% La Longue Marche — Volume 140-3 — mateo-canonical re-run
% Model: Gemini 3.1 Pro (medium thinking)
% Prompt: mateo-canonical (notation + structure block, Part I few-shot pool)
% Pages transcribed: 696/696
% Pages overlaid from diagram-tikzcd branch: 112
% Rebuilt: 2026-08-15

%% ===== Page 1 =====

\begin{center}
{\LARGE \textbf{La "Longue Marche" à travers la théorie de Galois}}
\vspace{1cm}

\textbf{Deuxième partie : variations sur l'équation des lacets [suite à La "Longue Marche" à travers la théorie de Galois\dots, pages 293 à 650, dont table des matières]\footnote{Lacunes entre les pages 371 et 379, la page 414 est manquante (ou erreur de pagination), la page 558 bis semble manquer. La page 518 est une photocopie.} : notes et copies de notes manuscrites (s.d.).}
\vspace{0.5cm}

Cote n$^\circ$ 140-3, 695 pages
\vspace{1cm}

Alexander GROTHENDIECK
\vspace{0.5cm}

s.d.
\end{center}

\newpage

%% ===== Page 2 =====

Pages 291-650 + table des matières

\underline{NB} Les pages 291, 292 figurent deux fois (erreur de numérotation)

\newpage

%% ===== Page 3 =====

\chapter*{2\textsuperscript{ème} partie : Variations sur l'équation des lacets}

\noindent \textbf{38} Détermination de $\hat{\Pi}_{0,3} \isommap \hat{\mathfrak{S}}_{0,3}^+$. \hfill 291

\vspace{0.5cm}
\noindent \textbf{39} Propriétés de congruence de l'opération de $\boldsymbol{\Gamma}$ sur $\hat{\Pi}_{0,3}$. Lien : l'hom $\mathcal{N}_1 \to \operatorname{Gl}(2, \hat{\mathbf{Z}})$. \hfill 299

\vspace{0.5cm}
\noindent \textbf{40} Paraphrases topologico-isotopiques de la théorie de Legendre \hfill 307'

\vspace{0.5cm}
\noindent \textbf{41} Conditions de commutation des $\hat{\Pi}_{0,3}^+ \isommap \hat{\mathfrak{S}}_{0,3}^+$ avec l'op. ext. de $\boldsymbol{\Gamma}$. \\
Calculs des autom. ext. dans $\hat{\Pi}_{0,3}$ et dans $\hat{\mathfrak{S}}_{0,3}^+$. \hfill 315

\vspace{0.5cm}
\noindent \textbf{42} Récapitulation sur le groupe $\hat{\mathfrak{S}}_{0,3}^+$ \\
\dots et sur les autom. à lacets de $\hat{\mathfrak{S}}_{0,3}^+$ et de $\hat{\Pi}_{0,3}$ : \hfill 331
\begin{itemize}
    \item[I] $\hat{\mathfrak{S}}_{0,3}^+$ \hfill 331
    \item[II] $\hat{\Pi}_{0,3}$ et $\hat{\mathfrak{S}}_{0,3}$ \hfill 332
    \item[III] relations avec $\hat{\Pi}_{0,3}$ et $\operatorname{Sl}(2, \hat{\mathbf{Z}})'$ \hfill 334
    \item[IV] relations avec $\hat{\mathfrak{S}}_{0,3}$ et $\operatorname{Sl}(2, \hat{\mathbf{Z}})'$ \hfill 335
    \item[V] les avatars de $\hat{\mathfrak{S}}_{0,3}$ \hfill 336
    \item[VI] Autom. à lacets de $\hat{\Pi}_{0,3}$ et $\hat{\mathfrak{S}}_{0,3}^+$ \hfill 337'
    \item[VII] Calculs dans $\operatorname{Sl}(2, \hat{\mathbf{Z}})'$ : cocycles normaux \hfill 342
    \item[VIII] Calculs dans $\operatorname{Sl}(2, \hat{\mathbf{Z}})'$ : cocycles sur $\mathfrak{S}_3$ \hfill 346
\end{itemize}
\marginpar{NB les calculs des sections VII, VIII en fait [unclear: s'appliquent au cas] d'un groupe qcq! [unclear: dans] des § 40 à 44 notations complexes...}

\vspace{0.5cm}
\noindent \textbf{43} Digressions sur les 1-cocycles des groupes cycliques \hfill 350 \\
(cas [unclear: pro-$p$, fini, $\dots$])

\vspace{0.5cm}
\noindent \textbf{44} Retour aux notations [unclear: initiales]. Mise en équations dans $\operatorname{Sl}(2, \hat{\mathbf{Z}})'$. \hfill 356

\vspace{0.5cm}
\noindent \textbf{45} L'hom. $\mathcal{M}_{0,3} \to \operatorname{Gl}(2, \hat{\mathbf{Z}})'$, et construction de l'extension $\mathcal{N}_{0,3}$ de $\boldsymbol{\Gamma}$ par $\operatorname{Sl}(2, \hat{\mathbf{Z}})'$. \hfill 384

\vspace{0.5cm}
\noindent \textbf{46} Notations définitives (?) -- récapitulation des calculs sur les autom. à lacets de $\hat{\Pi}_{0,3}$ et de $\hat{\mathfrak{S}}_{0,3}^+$. \hfill 390

\newpage

%% ===== Page 4 =====

\noindent I Les groupes $M_{0,3}^{\sim}$ et $N_{0,3}^{\sim \wedge}$ \hfill 390\\
II Calcul des éléments de $N_{0,3}^{\sim \wedge}$ \hfill 392\\
\hspace*{1em} 1) Voie $f, g, h$ \hfill 392\\
\hspace*{1em} 2) Voie $\alpha, \beta, \gamma$ \hfill 393\\
\hspace*{1em} 3) Voie $\alpha, \beta, f$ \hfill 393'\\
\hspace*{1em} 4) Voie $x, y, \xi$ \hfill 393'\\
\hspace*{1em} 5) Congruences au niveau 12, relat.\ aux cocycles \hfill 394\\
III Calculs dans $\operatorname{Sl}(2,\hat{\mathbf{Z}})'$ et cocycles dans $\operatorname{Sl}(2,\hat{\mathbf{Z}})$ \hfill 395\\
IV Les trois cas fondamentaux ($m=1, \rho=1, f=1$), et les relèvements $K_{1/2}, K_{2}, K_{\infty}$. Les sous-groupes $M_{0,3}^{\sim \wedge}(q)$, $q \in \{1, 2, 1/2\}$ \hfill 402\\
V Récapitulation en termes de $M_{0,3}^{\sim \wedge}[Q]$ \marginpar{calculs num. Les sous-groupes $M_{0,3}^{\sim \wedge}(q)$, $q \in \{1, 2, 1/2, \infty\}$} \hfill 404\\
VI Retour sur l'introduction de $\alpha$ et $\beta$ \hfill 414\\

\vspace{0.5em}
\noindent \underline{47.} Retour sur l'homomorphisme $\widehat{\Pi}_{0,3} \isommap \mathfrak{S}^{+\wedge}_{0,3}$, et les groupes $M_{0,3}^\wedge, N_{0,3}^\wedge, M_{0,3}^{\wedge(2)}$ \hfill 415\\
\underline{48.} Construction d'un schéma en groupes \hfill 420\\
\hspace*{1em} I) Etude du morphisme [unclear: $\alpha \mapsto \alpha p(\alpha)^{-1} f(p)$] \hfill 420\\
\hspace*{1em} II) L'équation [unclear: $\zeta = y \varepsilon_y (\xi \varepsilon - 1)$], et le groupe $B \tilde{\times} P_r$ \hfill 434'\\
\hspace*{1em} III) L'équation [unclear: $\alpha p(\alpha)^\sim = \tilde{p}(\beta)\varepsilon(\tilde{p}-1)$] \hfill 443'\\
\hspace*{1em} (IV) Un épinglage canonique associé, pour $\mathfrak{S}_{0,3}^\sim$ \hfill 445\\
\hspace*{1em} (V) L'équation $\alpha p(\alpha)^\sim = \varepsilon(\tilde{p} \sigma p(\sigma)^{-1}) \varepsilon_p(\mu - 1)$ et son étude au $2^e$ ordre. Le groupe $M_{0,3}^\sim$, la condition $Tr \sigma = 0$ \hfill 459\\
\hspace*{1em} (VI) Caractérisation du triplet $(E, \sigma, p)$ modulaire \hfill 462\\

\vspace{1em}
\noindent \underline{VII} Récapitulation sur les épinglages. \hfill 475\\
\underline{VIII} Récapitulation sur l'équation. Le groupe $S\mathcal{G}_{g,\nu}$ et ses variantes \hfill 485\\
\underline{IX} Les groupes $\mathcal{M}_{g,\nu}$ et $S\mathcal{M}_{g,\nu}$, et les homomorphismes [unclear: $\mathcal{S}\mathcal{M}_{0,3}^{\sim \wedge} \to \mathcal{M}_{g,\nu}^{(l)}/\mathcal{Z}_{g,\nu}$] \hfill 490'\\
\underline{X} Structure de $\mathcal{M}_{g,\nu}$, et l'homomorphisme canonique $\mathcal{M}_{g,\nu} \to \operatorname{Gl}(\nu) \simeq \operatorname{Gl}(2)$ \hfill 499\\
\underline{XI} Récapitulation sur $\mathcal{M}_{g,\nu}$ \hfill 508\\
\hspace*{1em} Reprise des (3°) $M_{g,\nu}$, $\mathcal{M}_{g,\nu}$, calculs de pts critiques \hfill 511\\
\hspace*{1em} 4°) L'hom $\theta: \mathcal{M}_{g,\nu} \to G$, $\bar{\theta}: \mathcal{M}_{g,\nu} \to P G$ \hfill 512, 513\\
\hspace*{1em} 5°) $R_{g,\nu}$, $S R_{g,\nu}$ \hfill 513'\\
\hspace*{1em} 6°) (Nouveau) - Le centre du radical $R_{g,\nu}$. Le sous-groupe $\Sigma \simeq \mathbf{Z}/2\mathbf{Z} \times \mathbf{Z}/2\mathbf{Z}$ \hfill 517\\
\hspace*{1em} 7°) Division de $R_{g,\nu}$ et de $M_{g,\nu}$ \hfill 526'\\
\hspace*{1em} 8°) Les paramètres $U, z, k, \mu, \varepsilon, e$ et $x, f, \alpha, \beta, \infty, \rho, (i, j, q, h)$ \hfill 529\\
\hspace*{1em} 9°) Les sous-groupes $M_{g,\nu}(\alpha)$, $M_{g,\nu}[\alpha]$, et l'algèbre $S\mathcal{M}_{g,\nu, \text{lisse}}$. L'équa $\sigma = S \Sigma_{\rho, \sigma}^\alpha$ \hfill 542\\
\hspace*{1em} 10°) Les sous-groupes (nouveau) $\mathcal{M}_{g,\nu}(A)$, $\mathcal{M}_{g,\nu}(p)$ \hfill 550'\\
\underline{XII} Rétrospective. \hfill 592

\newpage

%% ===== Page 5 =====

\section*{49. Les groupes $\mathcal{M}_{p,\sigma}$ généraux \hfill 550}

\noindent [ I Préliminaires heuristiques ]

\noindent II Les groupes $\mathcal{M}_{p,\sigma}$ ``sans signes'' -- nouvelles variantes sur la relation des lacets \hfill 553

\noindent III Passage de $\mathcal{M}_{0,3}^{\sim}$ à $\mathcal{N}_{1,1}^{\sim}$ -- et de $\operatorname{Gl}(2, \mathbf{Z})^{\wedge} / \pm 1$ à $\operatorname{Gl}(2, \mathbf{Z})^{\wedge}$ \hfill 565

\noindent [ IV Où on se bat avec ``les signes'' $\varepsilon_2, \varepsilon_3$ dans la relation des lacets -- calculs -- brouillons\dots ] \hfill 568

\noindent V Les groupes $\mathcal{M}_{p,\sigma}$ généraux (contexte pro-fini) \hfill 577
\begin{enumerate}
\item[1º)] Données \hfill 577
\item[2º)] $Z_{p,\sigma}, \mathfrak{S}_{p,\sigma}$ \hfill 577'
\item[3º)] $S_0 G_{p,\sigma}, G_{p,\sigma}$ \hfill 579
\item[4º)] $Z_p, Z_\sigma, \mathcal{N}_p, \mathcal{N}_\sigma, \mathcal{N}_p^\wedge, \mathcal{N}_\sigma^\wedge, S\mathcal{N}_p^\wedge, S\mathcal{N}_\sigma^\wedge$ \hfill 580
\item[5º)] $S_0 \mathcal{M}_{p,\sigma}, \mathcal{M}_{p,\sigma}, S_0 \Gamma_{p,\sigma}, \Gamma_{p,\sigma}$ \hfill 585
\item[6º)] $\Xi_{p,\sigma}, S\mathfrak{G}_{p,\sigma}, \Xi_{p,\sigma}^{\text{eff}}, S\mathfrak{G}_{p,\sigma}^{\text{eff}}$ \hfill 586
\item[7º)] Fonctorialités. \hfill 588
\item[8º)] \marginpar{rectifications pour la ``relation du cas pur''} Le cas universel $\mathfrak{G}^\wedge \simeq \operatorname{Gl}(2, \mathbf{Z})^\wedge$. Le sous-groupe $K_0$. Étude de l'extension $S_0 \Gamma_{p,\sigma}$ de $Z_{p,\sigma} = M(\Gamma_0)$ par $L_p \times L_\sigma$, et $S_0^{\text{eff}} \Gamma_{p,\sigma}^1$ de $Z_{p,\sigma}$ par $\mu = \pm 1$. \hfill 590
\item[9º)] Récapitulation sur le cas universel. Structure de $\mathcal{N}_\sigma, \mathcal{N}_\sigma^\wedge$. L'hom. can. $\mathcal{M} \to \mathcal{M}_{p,\sigma} / \mathcal{M}_p \isom \mathcal{M}_{p,\sigma} / \mathcal{N}_p$. \hfill 614
\end{enumerate}

\noindent VI Les groupes $\mathcal{M}_{p,\sigma}$ généraux : récapitulations et raffinement de la construction ci-dessus. \hfill 628
\begin{enumerate}
\item[1º)] Données \hfill 628
\item[2º)] $Z_{p,\sigma}, \mathfrak{S}_{p,\sigma}$ \hfill 629
\item[3º)] $G_{p,\sigma}, G_{p,\sigma} \dots$ \hfill 632
\item[4º)] $\mathcal{N}_p, \mathcal{N}_\sigma, \mathcal{N}_p^\wedge, \mathcal{N}_\sigma^\wedge \dots$ \hfill 635
\item[5º)] $\mathcal{M}_{p,\sigma}, \mathcal{M}_{p,\sigma} \dots$ \hfill 639
\item[6º)] Fonctorialités \hfill 640
\item[7º)] $\mathcal{M}_{p,\sigma}[p]$, $\mathcal{M}_{p,\sigma}[s]$ \hfill 642
\item[8º)] Autom. virtuels : $\mathcal{M}_p^{\text{virt}}$ \hfill 643
\item[9º)] Digression sur les schémas en groupes \hfill 648
\item[10º)] Lien avec §47
\end{enumerate}

\newpage

%% ===== Page 6 =====

\section*{29. (bis). Détermination explicite de l'hom. $\Pi'_{0,3} \to \underline{C}_{0,3}^+$}
(et de $\Pi'_{0,3}$ ([unclear: Gr. pro-catégorique] $\to \underline{C}_{0,3}$)).

J'ai envie de revenir sur la prop. du \S\ précédent (p. 262) avec surfaces $X$, groupe discret $G$ opérant proprement, les stabilisateurs des $x \in X$ conservent l'orientation et opèrent fidèlement --- sur un voisinage de $x$ ---
on prend $Y = X/G$, muni de la ``signature'' $\underline{d} = d(X/Y)$ définie par la ramification ---
d'où un iso
$$
\Pi = \Pi_1((X,G), x) \isommap \Pi_1((Y,\underline{d}), y) = \Pi' \leqno{(1)}
$$
($x \in X$, $y$ image de $x$ dans $Y$)
que je voudrais étudier de plus près.

On a dans $\Pi_1((Y,\underline{d}), y)$, pour tout $b \in Y!$ (supp.\ $d_b>1$)
des ``sous-groupes locaux-disques'', plus précisément
une classe de conjugaison de ss-groupes
quotients de $\mathbf{Z}/d_b\mathbf{Z}$, et plus précisément encore,
un hom.\ canonique
$$
T_b \otimes_{\mathbf{Z}} \mathbf{Z}/d_b\mathbf{Z} \to \Pi' \leqno{(2)}
$$
déduit par passage au quotient de l'hom.
$$
T_b \overset{\xi_b}{\to} \Pi_1(Y \setminus Y!) \quad (\to \Pi_1((Y,\underline{d}))) \leqno{(3)}
$$
(où $T_b$ est le $\mathbf{Z}$-module libre de rang $1$ des ``entiers tordus'' en $b$ sur $Y$). D'autre part, si $a \in X$ est au-dessus de $b$, le stabilisateur $G_a$ est un groupe cyclique d'ordre $d_a=d_b$, on a un isom.\ can.\ évident
$$
G_a \simeq T_a \otimes_{\mathbf{Z}} \mathbf{Z}/d_a\mathbf{Z} \simeq T_b \otimes \mathbf{Z}/d_b\mathbf{Z} \leqno{(4)}
$$
(NB $T_a \simeq T_b$ canoniquement).

\newpage

%% ===== Page 7 =====

D'autre part, à $a$ est associé une classe de conjugaison de scindages de l'extension
$$
1 \to \Pi_1(X,x) \to \Pi_1((X,G),x) \to G \to 1 \leqno{(5)}
$$
(cf § 15)
$$
k_a : G_a \to \Pi_1((X,G),x) \leqno{(6)}
$$
(à un $\Pi_1(X,x)$-conjugué près).
Ces classes (ou scindages (si $X \neq S^2$)) correspondent biunivoquement aux points fixes de $G$ opérant sur $X$. Si on prend l'homomorphisme extérieur défini par (6), celui-ci est défini à conjugaison près non seulement par les $g \in \Pi_1(X)$, mais par ceux de $\Pi_1(X,G)$ tout entier) on trouve une correspondance 1-1 entre l'ensemble des homomorphismes extérieurs (6) qui relèvent l'inclusion
$$
G_a \subset G
$$
(considérés comme homomorphismes extérieurs), et l'ensemble des orbites de $G$ dans $X!$

donc entre les points de $Y!$ En résumé -

\underline{Proposition} On a une bijection entre $Y!$ et l'ensemble des homomorphismes extérieurs "maximaux" ([unclear: $\varphi = \varphi_\Gamma : \Gamma \to \Pi_1$]) de sous-groupes $\Gamma \subset G$, (où $\Gamma \neq 1$, dans $\Pi_1(X,G)$ qui relèvent l'inclusion $\Gamma \hookrightarrow G$ (les dits $\Gamma$ sont nécessairement cycliques, les $\Gamma$

\newpage

%% ===== Page 8 =====

292 (bis)

en question correspondent à un $b \in Y!$ et à
les stabilisateurs $G_a$ dans $G$ des
$a \in X_b = X!_b$. De plus, le composé des
homo.
$$
T_b \otimes \mathbf{Z}/\nu\mathbf{Z} \xrightarrow{\sim} G_a \xrightarrow{\varphi_a=\varphi_b} \pi_1(X, G) \xrightarrow{\sim} \pi_1(Y, \underline{d}) \leqno{(7)}
$$
n'est autre (on s'en serait douté !) que l'hom.
"lacets tronqué" (2).

Supposons maintenant que $X$ soit de
la forme $\hat{X} \setminus S$, $\hat{X}$ surface compacte,
$S$ partie finie de $\hat{X}$, donc
$$
\begin{cases}
X = \hat{X} \setminus S \quad \text{et} \quad Y = \hat{Y} \setminus T \\
\hat{Y} = \hat{X} / G
\end{cases} \leqno{(8)}
$$
Considérons alors, pour tout $a \in S$, l'hom.
extérieur
$$
k_a : T_a \longrightarrow \pi_1(X) \hookrightarrow \pi_1(X, G) \leqno{(9)}
$$
que je vais comparer au composé
$$
T_b \xrightarrow{k_b} \pi_1(Y) \longrightarrow \pi_1(Y, \underline{d}) \leqno{(10)}
$$
pour $b = \hat{f}(a)$, moyennant l'isom. can.
$$
T_a \simeq T_b \leqno{(11)}
$$
Si $d_a$ désigne encore l'indice de ramifica-
tion en $a$, on aura un diagramme
commutatif de compatibilité :

\newpage

%% ===== Page 9 =====

(12)
\begin{tikzcd}
T_a \ar[rr, "K_a"] \ar[d, "\text{[unclear]} \text{ can}"'] & & \pi_1(X) \ar[r] & \pi_1(X, G) \ar[dd, "\wr"] \\
\tilde{T}_b \ar[d, "d_a \cdot id"'] & & & \\
T_b \ar[rrr, "K_b"'] & & & \pi_1(Y, \underline{d})
\end{tikzcd}

Nous allons appliquer ceci à la situation
où $X = U_{0,3}$ , $G = \mathfrak{S}_3$ , $\hat{X} = \mathbb{P}^{1 \text{ an}}$, $\hat{Y} \simeq \mathbb{P}^{1 \text{ an}}$,
où l'application de passage au quotient est
donnée par (p. 288' (58))

(13) $$ f(z) = \frac{3^3}{2^2} \frac{z^2(z-1)^2}{(z^2-z+1)^3} $$
$$ f : \begin{matrix} 0 \\ 1 \end{matrix} \mapsto 0 \qquad \begin{matrix} 1/2 \\ 2 \\ -1 \end{matrix} \mapsto 1 \qquad \begin{matrix} j \\ \bar{j} \end{matrix} \mapsto \infty $$

On voit donc, avec les notations de ,
que les générateurs $\ell'_0, \ell'_1, \ell'_\infty$ de $\Pi'_{0,3}$
(satisfaisant
(14) $$ \ell'_\infty \ell'_1 \ell'_0 = 1 \qquad {\ell'_1}^2 = {\ell'_\infty}^3 = 1 \quad ) $$
correspondent respectivement, par l'isomor-
phisme
(15) $$ \overline{\mathfrak{S}}^+_{0,3} = \Pi_1(U_{0,3}, \mathfrak{S}_3) \stackrel{\sim}{\longrightarrow} \Pi'_{0,3} = \overline{C}^+_{1,1} $$
[unclear: groupe cartographique bi-parti pondéré or.] \quad \quad \quad [unclear: groupe cartographique triangulé orienté $\simeq \operatorname{Sl}(2,\mathbb{Z})$]

aux éléments notés

(16)
$$ \begin{array}{rcccl}
\varepsilon_0, \varepsilon_1, \varepsilon_\infty & & & & \\
\text{(tels que } \varepsilon_i^2 = 1 \text{)} & & \longmapsto & \ell'_0 & \\
& & & & \\
\sigma_0, \sigma_1, \sigma_\infty & & \longmapsto & \ell'_1 & \quad \text{mod conjugaison} \\
& & & & \\
\text{[MARGIN: ou } \rho^{-1} \text{ ? ]} \quad \quad \rho & & \longmapsto & \ell'_\infty &
\end{array} $$

\newpage

%% ===== Page 10 =====

293

[MARGIN: de $\pi = \overline{C}_{0,3}^+$]
Ceci montre que l'on peut trouver dans 
un conjugué $\varepsilon_0'$ de $\varepsilon_0$, $\sigma_0'$ de $\sigma_0$, tels qu'on ait

(1) $\rho \sigma_0' \varepsilon_0' = 1$, en plus de $\rho^3 = {\sigma_0'}^2 = 1$

et de telle façon que ces relations constituent
une présentation de $\pi$ (donc en particulier,
que $\varepsilon_0', \sigma_0', \rho$ engendrent $\pi$). S'inspirant d'
ailleurs de l'isomorphisme plus gros qui
prolonge (15) et s'obtient :

(16) $$\overline{C}_{0,3} \simeq \pi_1 ( U_{0,3}, \underset{\overline{C}_{0,3}}{\underbrace{\mathfrak{S}_3 \times \mathbb{Z}/2}} ) \longrightarrow \overline{C}_{1,1}'$$
[MARGIN: groupe cartographique triangulaire étendu $\simeq \operatorname{Gl}(2, \mathbb{Z})$]

on devrait trouver dans le premier
trois générateurs $\tau_0', \tau_1', \tau_\infty'$ (correspondant aux
gén. analogues $\tau_0, \tau_1, \tau_\infty$ du groupe cartogra-
phique correspondant) satisfaisant

(17) $${\tau_0'}^2 = {\tau_1'}^2 = {\tau_\infty'}^2 = 1, \quad (\underbrace{\tau_0' \tau_1'}_{\text{corr. à } l_1'})^2 = (\underbrace{\tau_1' \tau_\infty'}_{\text{corr. à } l_\infty'})^3 = 1$$

de telle façon que ce soient là des
relations de présentation de $\overline{C}_{0,3}$. Il
faudrait que $\tau_\infty'$ soit conjugué de $\sigma_0$.

(18) $$
\begin{matrix}
\tau_\infty' \tau_1' = \varepsilon_0' & ; & \tau_0' \tau_\infty' = \sigma_0' & , & \tau_1' \tau_0' = \rho \\
\downarrow & & \downarrow & & \downarrow \\
l_0' & & l_1' & & l_\infty'
\end{matrix}
$$

\newpage

%% ===== Page 11 =====

$X$ surface topologique

$\mathcal{G}$ groupe discret opérant sur $X$

$G$ sous-groupe invariant dans $\mathcal{G}$, opérant proprement sur $X$, de sorte que $\forall x \in X$ les stabilisateurs $G_x$ ($x \in X$) opèrent sans renverser l'orientation locale en $x$.

$Y = X/G$ (surface topologique), $f : X \to Y$

$X'$ partie discrète de $X$, stable sous $\mathcal{G}$,

$d_X$ \og signature \fg{} sur $X$ de support $X'$, inv. par $\mathcal{G}$
$$ d_X : X' \to \mathbf{N} $$

$Y' = X'/G \subset Y$ image de $X'$ dans $Y$

On suppose $x \in X, \quad G_x \neq \{1\} \Rightarrow x \in X'$

$d_Y : Y' \to \mathbf{N}, \quad d_Y(y) = \nu_x d_X(x)$
si $y = f(x), \quad \nu_x = [G_x]$

$\operatorname{Rev}(X, d_X, \mathcal{G})$ catégorie des $\mathcal{G}$-revêtements ramifiés de $X$, subordonnés à $d_X$

$\operatorname{Rev}(Y, d_Y, \mathcal{H})$ catégorie des $\mathcal{H}$-revêtements ramifiés de $Y$, subordonnés à $d_Y$

Pour $X'' \in Ob \operatorname{Rev}(X, d_X, \mathcal{G})$, considérons
$X''/G$ comme surface sur laquelle opère
$\mathcal{H} = \mathcal{G}/G$, au-dessus de $Y = X/G$. Alors

\newpage

%% ===== Page 12 =====

\begin{enumerate}
\item[a)] $X'/G \in Ob \mathcal{R}ev(Y, d_Y, \mathcal{H})$
\item[b)] on trouve ainsi un foncteur
$$
\mathcal{R}ev(X, d_X, G) \to \mathcal{R}ev(Y, d_Y, \mathcal{H}) \leqno{(19)}
$$
\item[c)] ce foncteur est une équivalence de catégories.
\item[d)] Description d'un foncteur quasi-inverse par image inverse ``normalisée''.
\end{enumerate}

Remarques. OPS $d_X : \text{valeurs} > 0$ i.e.\ $d_Y : \text{valeurs} > 0$
(sinon, si $S = \{x \in X' \mid d_X(x) = 0\}$, remplacer $X$ par $X \setminus S$, $Y$ par $Y \setminus T$ ($T = S/G$)).

OPS $d_X(x) = 1 \implies [G_x] \neq 1$ (sinon, [unclear: enlever de] $X'$ [unclear: les points avec] $d_X(x) = 1$ et $G_x = 1$).

Le cas le plus intéressant pour [unclear: nous] est celui où $d_X(x) = 1$ $\forall x \in X'$, et où $G_x \neq \{1\}$ (donc $X' = \{x \in X \mid G_x \neq \{1\}\}$), dans ce cas on a
$$
\mathcal{R}ev(X, d_X, G) = \mathcal{R}ev(X, G) \leqno{(20)}
$$
et on [unclear: utilise] $\mathcal{R}ev(Y, d_Y, \mathcal{H})$ pour s'enquérir [unclear: sur] $\mathcal{R}ev(X, G)$
(sans [unclear: données] de ramification)

\newpage

%% ===== Page 13 =====

\noindent
\underline{Corollaire} Soit $x \in X \setminus X'$, $y$ son image dans $Y \setminus Y'$. On a un isomorphisme canonique
$$
\Pi_1((X, d_X, G)_x) \isommap \Pi_1((Y, d_Y, \mathcal{H})_y) \leqno{(21)}
$$
Nous allons expliciter cet isomorphisme en termes de classes de lacets classiques.

Un élément $\alpha$ de $\Pi_1((X, d_X, G)_x)$ provient d'un couple $(\gamma, \ell)$, où
$$
\gamma \in G, \quad \ell \text{ classe de chemins dans } X \setminus X' \text{ de } x \text{ vers } \gamma^{-1}x
$$
(avec $(\gamma, \ell)$ et $(\gamma', \ell')$ définissant le même élément ssi $\gamma' = \gamma$, et $\ell'^{-1} \ell$ de $\pi_1(X \setminus X', x)$ est dans le sous-groupe invariant engendré dans $\pi_1(X \setminus X', x) \longrightarrow \pi_1((X, d_X), x)$ --- par les $k \lambda_a^{d_X(a)} k^{-1}$, pour $a \in X'$, $\lambda_a$ désigne la classe dans $\pi_1(X \setminus X', x)$ d'un "lacet de $x$ entourant $a$").

Soit alors $j$ l'image de $\gamma$ dans $\mathcal{H}$ --- et soit $p_U : X \setminus X' \defeq U \longrightarrow Y \setminus Y' \defeq V \isom U/G$ l'application.

\newpage

%% ===== Page 14 =====

canoniques, d'où l'image $\beta$ \marginpar{de $\alpha \in \Pi_1(X, d_X, \mathcal{G}; x)$ défini par $(\gamma, l)$}
dans $\Pi_1(Y, d_Y, \mathcal{H}; y)$ est $(\dot{\gamma}, f_u(l))$

(NB $f_u(l)$ est une classe de chemins dans $V \defeq Y \setminus Y'$ de $y = f_u(x)$ vers $\dot{\gamma}^{-1}(y) = f_u(\gamma^{-1}x)$).

\marginpar{NB.}On voit tout de suite que si
$(\gamma, l)$ et $(\gamma', l')$ définissent le même $\alpha$, alors
ils définissent le même $\beta$, ce qui
provient du fait que
$$ \pi_1(f_u) \colon \pi_1(X \setminus X', x) \to \pi_1(Y \setminus Y', y) $$
applique un $\lambda_a$ dans $\lambda_b^{\nu(a)}$
($b = f(a)$, $\nu(a) = [\mathcal{G}_a] =$ degré de ramification de $f$ en $a$)
donc $\lambda_a^{d_X(a)}$ dans $\lambda_b^{d_X(a)\nu(a) = d_Y(b)}$

Revenons à la situation de (13) \marginpar{$X' = \{2,-1, 1/2 \}$}
$$ X = \mathcal{U}_{0,3}, \quad \mathcal{G} = \mathfrak{T}_{0,3} \simeq \mathfrak{S}_3 \times \mathbf{Z}/2, \quad \mathcal{G}^+ = \mathfrak{T}_{0,3}^+ = \mathfrak{S}_3 $$
$$ Y \simeq \mathbf{P}^1 \setminus \{0\}, \quad Y' = \{1, \infty\}, \quad \mathcal{H} = \mathbf{Z}/2 \mathbf{Z} $$
[unclear: prolongé aux quotients] de $\widehat{X} = \mathbf{P}^1$ et $\widehat{Y} = \mathbf{P}^1$
donnés par (13)
$$ f(z) = \frac{3^3}{2^2} \frac{z^2(z-1)^2}{(z^2-z+1)^3} $$

\newpage

%% ===== Page 15 =====

Pour calculer $\pi_1(Y, d_Y, \mathbf{Z}/2\mathbf{Z}) = \widetilde{\mathfrak{T}}_{0,3}$ (groupe cartographique trigénéré non orienté), on a pris le point base $y = -\bar{J}$ , et choisi un $x \in U_{0,3}$ au dessus de $y$, tel que
$$
f(x) = y = -\bar{J}
$$
- il y a un choix [unclear: d'un] tel dans un $f^{-1}(\{y\})$ à six éléments, dont trois dans l'hémisphère nord (et les trois autres dans l'hémisphère sud), ceux-ci sont permutés transitivement par $\mathfrak{S}_3$.
Ils se trouvent respectivement dans les triangles sphériques $(-\bar{J}, 0, 1/2)$, $(-\bar{J}, 1, 2)$, $(-\bar{J}, \infty, -1)$ (appliqués par $f$ homéomorphiquement sur le triangle sphérique $(0, 1, \infty)$ supérieur (ou inférieur resp.) de centre $-\bar{J}$ ). Le choix de $x$ revient donc à un choix de l'un des trois segments sphériques $(0, 1/2)$, $(1, 2)$, $(\infty, -1)$ — j'ai envie de prendre $(0, 1/2)$, correspondant au triangle $(-\bar{J}, 0, 1/2)$.

Les trois générateurs $\tilde{\tau}_0, \tilde{\tau}_1, \tilde{\tau}_\infty$ dans $\pi_1(Y, d_Y, \mathbf{Z}/2\mathbf{Z})$ au dessus de $\tau \in \mathbf{Z}/2\mathbf{Z}$ s'obtiennent en prenant des chemins $l_0', l_1', l_\infty'$ partant du point base $y = -\bar{J}$ vers $\tau y = -J$, qui sont les arcs de grand cercle coupant l'équateur resp. en $2, -1, 1/2$.

\newpage

%% ===== Page 16 =====

Ces chemins sont les images par $f$ de chemins sur $X$, soient $\lambda_0: x \xrightarrow{\lambda_0} g_0 x$, $\lambda_1: x \xrightarrow{\lambda_1} g_1 x$, $\lambda_\infty: x \xrightarrow{\lambda_\infty} g_\infty x$, qui coupent la subdivision barycentrique des dits triangles sphériques $(0,1,\infty)$ en exactement un point, situés respectivement sur les côtés $(1/2, \bar{J}), (\bar{J}, 0), (0, 1/2)$ du triangle sphérique $(0, 1/2, \bar{J})$.\marginpar{NB Les points $2_0, 2_1, 2_\infty$ sont les points [unclear: fixés] par $\tau_0, \tau_1, \tau_\infty$ sur $X$ (au dessus de [unclear: $\bar{J}$ ou $\tau \bar{J}$]) dans les triangles sphériques $(\bar{J}, 1/2, -1)$, $(0, \bar{J}, -1)$, $(\bar{J}, 1/2, 1)$.} Ces points (les $2_0, 2_1, 2_\infty$) sont [unclear: d'abscisses] et extrémités respectives de la forme $g_0 x = \sigma_0 \tau x, g_1 x = \sigma_1 \tau x, g_\infty x = \sigma_\infty \tau x$ (i.e. $g_0 = \sigma_0 \tau, g_1 = \sigma_1 \tau, g_\infty = \sigma_\infty \tau$). Ainsi les générateurs involutifs $\tau_0', \tau_1', \tau_\infty'$ de $\widetilde{\Pi}_{0,3} = \Pi_1(Y, dy, \mathbf{Z}/2\mathbf{Z}; y)$ sont les images des éléments (notés par le même signe)
$$ \tau_0' = (\sigma_0 \tau, \lambda_0) ; \quad \tau_1' = (\sigma_1 \tau, \lambda_1), \quad \tau_\infty' = (\sigma_\infty \tau, \lambda_\infty) $$
de $\Pi_1(X, \underline{G} ; x)$. Nous allons utiliser le bout de grand cercle $x \to -\bar{J} = p$ (origine habituelle pour le calcul de $\Pi_1(X, \underline{g})$) pour identifier $\Pi_1(X, \underline{G} ; x)$ à $\Pi_1(X, \underline{g} ; p)$.

On peut alors prendre
$$ \tau_0' = (\sigma_0 \tau, \lambda_0'), \quad \tau_1' = (\sigma_1 \tau, \lambda_1'), \quad \tau_\infty' = (\sigma_\infty \tau, \lambda_\infty') \leqno{(21)} $$
où $\lambda_0', \lambda_1', \lambda_\infty'$ sont les bouts de chemins de $p = -\bar{J}$ vers $g_0 p, g_1 p, g_\infty p$ [unclear: qui s'en déduisent].

\newpage

%% ===== Page 17 =====

[unclear: tien des demi-cercles en arêtes] pour le graphe barycentrique des doubles triangles sphériques à [unclear: sommet] $P=-j$, [unclear] les points, [unclear: situés] sur les segments sphériques $(0, 1/2)$ ! Les points $z_0, z_1, z_\infty$ sont construits respectivement comme les points $z$ tels que $f(z)=-j$ et qui se trouvent sur les triangles sphériques

$$([unclear: 1/2], j, 1), (-j, 0, -1), (0, 1/2, -j).$$

donnés par

$$\lambda'_i = [unclear](\lambda_i) l$$

On a donc

$$\lambda'_0 = (\sigma_0 \tau)(l') \cdot \lambda_0 \cdot l \quad \text{[MARGIN: homotope par un lacet trivial dans } U_{0s} \text{]}$$
[unclear: symétrie par rapport équateur] $(\infty, -j, 1/2, -j)$

$$\lambda'_1 = \sigma_1 \tau(l') \cdot \lambda_1 \cdot l \quad \text{[MARGIN: homotope en lacet trivial dans } U_{0s} \text{]}$$
[unclear: symétrie par l'équateur] $(0, -j, 1, -j)$

$$\lambda'_\infty = (\tau(l')) \cdot \lambda_\infty \cdot l \quad \text{[MARGIN: homotope au chemin du grand cercle } (-j, 1/2, j) \text{ de } P=\bar{j} \text{ vers } \tau P = -j \text{]}$$

\newpage

%% ===== Page 18 =====

En résumant, avec les notations de [C] (4b) p. 7 ff, on aura donc
$$
\begin{cases}
\tau'_0 = \widetilde{\sigma}_\infty & (\text{ou } \dot{\sigma}_1) \\
\tau'_1 = \widetilde{\sigma}_0 & (\text{ou } \dot{\sigma}_\infty) \\
\tau'_\infty = \tau_\infty & (\text{ou } \dot{\tau})
\end{cases} \qquad \text{NB } \tau_\infty = \widetilde{\sigma}_\infty \widetilde{\sigma}_0 = \widetilde{\sigma}_0 \widetilde{\sigma}_\infty \leqno{(22)}
$$

Il faut vérifier les relations essentielles
$$
\tau'_0{}^2 = \tau'_1{}^2 = \tau'_\infty{}^2 = (\underbrace{\tau'_\infty \tau'_0}_{\widetilde{\sigma}_\infty})^2 = (\underbrace{\tau'_1 \tau'_0}_{\rho^{-1}})^3 = 1 \leqno{(23)}
$$

Les quatre premières sont claires, la dernière s'écrit
$$
(\widetilde{\sigma}_0 \widetilde{\sigma}_\infty)^3 = 1
$$
ou
$$
\widetilde{\sigma}_0 \widetilde{\sigma}_\infty = \rho^{-1} \quad (\text{[C], formule 13,})
$$
d'où ce qu'on veut.

\marginpar{\raggedright On exprime les $l_0, l_1, l_\infty, \tau_0, \tau_1, \tau_\infty, \rho, \sigma_0, \sigma_1, \sigma_\infty \in \Pi_{0,3}(V)$ en fct de $\tau'_0, \tau'_1, \tau'_\infty$ ainsi : $\rho = \tau'_0 \tau'_1$, $\tau_\infty = \tau'_\infty$, $\tau_0 = \rho \tau_\infty \rho^{-1} = \tau'_0 \tau'_1 \tau'_\infty \tau'_1 \tau'_0$, $\tau_1 = \rho^{-1} \tau_\infty \rho = \tau'_1 \tau'_0 \tau'_\infty \tau'_0 \tau'_1$, $\sigma_\infty = \tau'_0$, $\sigma_0 = \tau'_1$, $\sigma_1 = \tau'_0 \tau'_1 \tau'_0 = \tau'_1 \tau'_0 \tau'_1$, $l_0 = \tau_\infty \sigma_1 = \tau'_\infty \tau'_0 \tau'_1 \tau'_0$, $l_1 = \tau_0 \sigma_0 = \tau'_0 \tau'_1 \tau'_\infty \tau'_1 \tau'_0 \tau'_1 = \tau'_0 \tau'_1 \tau'_\infty \tau'_0$, $l_\infty = \tau_1 \sigma_\infty = \tau'_1 \tau'_0 \tau'_\infty \tau'_0 \tau'_1 \tau'_0 = \tau'_1 \tau'_0 \tau'_\infty \tau'_1$.}
Calculons (cf notations de (18)) $\tau'_\infty \tau'_1$, $\tau'_0 \tau'_\infty$, $\tau'_1 \tau'_0$ -- qui devaient correspondre à $l'_0, l'_1, l'_\infty$ -- 
$$
l'_0 = \tau'_\infty \tau'_1 = \tau_\infty \widetilde{\sigma}_0 = \widetilde{\sigma}_\infty \underbrace{\widetilde{\sigma}_0 \widetilde{\sigma}_0}_{1} = \widetilde{\sigma}_\infty
$$
$$
l'_1 = \tau'_0 \tau'_\infty = \widetilde{\sigma}_0
$$
$$
l'_\infty = \tau'_1 \tau'_0 = \rho^{-1}
$$

i.e.
$$
l'_0 = \widetilde{\sigma}_\infty , \quad l'_1 = \widetilde{\sigma}_0 , \quad l'_\infty = \rho^{-1} \leqno{(24)}
$$

avec les relations caractéristiques
$$
l'_0 l'_1 l'_\infty = 1 \quad , \quad l'_1{}^2 = 1, \quad l'_\infty{}^3 = 1 \leqno{(25)}
$$

Tout est OK !

\newpage

%% ===== Page 19 =====

Je voudrais encore expliciter la relation
$$ \Pi_{0,3} \simeq \mathfrak{S}_{0,3} $$
\marginpar{qui revient à $\operatorname{Aut}_{ext}(\Pi_{0,3})$}
avec $GP(1,\mathbf{Z}) \simeq \operatorname{Gl}(2,\mathbf{Z})$.

i.e.\ $\operatorname{Gl}(2,\mathbf{Z}) / \{\pm 1\}$. Rappelons d'abord l'opération
de
$$ GP(1,\mathbf{R}) \simeq \operatorname{Gl}(2,\mathbf{R})/\mathbf{R}^* \simeq \operatorname{Gl}(2,\mathbf{R})^{\pm 1} / \{\pm 1\} \leqno{(26)} $$
\footnote{NB $GP(1,\mathbf{R})$ opère fidèlement sur $D$, et s'identifie au groupe des automorphismes de $D$, $\operatorname{Sl}(2,\mathbf{R})/\pm 1$ s'identifiant aux automorphismes analytiques.}dans le demi-plan de Poincaré $D$
$$ D = \{ z \in \mathbf{C} \ | \ \operatorname{Im}(z) > 0 \} = \{ x + i y \ | \ x \in \mathbf{R}, y \in \mathbf{R}^{*+} \} \leqno{(27)} $$
faisant opérer la matrice $u = \begin{pmatrix} a & b \\ c & d \end{pmatrix}$ par
$$ f_u(z) = \begin{cases} \frac{az+b}{cz+d} & \text{si } \det u > 0 \\ \frac{a\bar{z}+b}{c\bar{z}+d} & \text{si } \det u < 0 \end{cases} \leqno{(28)} $$
De cette façon, $GP(1,\mathbf{R})$ est isomorphe au groupe des automorphismes conformes de $D$.

D'ailleurs, la donnée d'une courbe elliptique\footnote{$E \simeq \mathbf{C} / \pi$} équivaut : à la donnée d'un espace vectoriel $V$ complexe de dimension 1 sur $\mathbf{C}$, et d'un "réseau" $\pi \subset V$ (sous-groupe discret de rang 2).
L'orientation canonique de $V$ correspond à un générateur pour $\bigwedge_{\mathbf{Z}}^2 \pi \simeq \mathbf{Z}$ dit "orienté". Le choix d'une rigidification

\newpage

%% ===== Page 20 =====

de base uni-modulaire de $\pi$, équivaut à (comme $\mathbf{Z}$-mod.) celui d'une base $(e_1, e_2)$ de $\pi$ telle que (comme base de $V$ sur $\mathbf{R}$) $e_1 \wedge e_2 = \eta$. L'élément $e_1$ permet d'ailleurs d'identifier $V$ à $\mathbf{C}$, $e_1$ devenant $1$. Ainsi la donnée de $E$ munie de sa rigidification de base uni-modulaire, équivaut à celle d'un élément $e_2 = \tau$ de $\mathbf{C}$, tel que
\begin{enumerate}
    \item[1)] $(1, \tau)$ engendre un sous-groupe discret de rang $2$, i.e.\ $\tau \notin \mathbf{R}$ i.e.\ $\Im \tau \neq 0$.
    \item[2)] L'orientation induite sur [unclear: le $\mathbf{R}$-e.v.] par la base $(1, \tau)$ est « directe », i.e.\ $\Im \tau > 0$, i.e.\ $\tau \in D$.
\end{enumerate}

Quand on regarde l'opération de $\operatorname{Sl}(2, \mathbf{Z})$ sur (les classes d') isomorphie de courbes ell.\ rigidifiées, ainsi identifiées à $D$, on trouve que $u = \begin{pmatrix} a & b \\ c & d \end{pmatrix}$ opère par $\tau \mapsto \frac{a\tau + b}{c\tau + d}$, c'est l'opération habituelle de $\operatorname{Sl}(2, \mathbf{Z})$ sur $D$, à travers $\operatorname{Sl}(2, \mathbf{Z}) \to \operatorname{GP}(1, \mathbf{Z})^+$.
Ce groupe opère proprement, il a exactement deux trajectoires critiques, qui corres-

\newpage

%% ===== Page 21 =====

-pondent respectivement aux $\tau$ qui avec $1$ engendrent le réseau des entiers alg. de $\mathbf{Q}(i)$ (soit $\mathbf{Z} + i\mathbf{Z}$), et le réseau des entiers de $\mathbf{Q}(j)$, formé des nbs de la forme $a + b j$ ($j = \frac{-1+i\sqrt{3}}{2} = \exp \frac{2 i \pi}{3}$) ; les [unclear: ordres] de fixité correspondantes st $2$ et $3$.

On a un isom. canonique (analytique complexe)
$$
D / \operatorname{Gl}(2,\mathbf{Z})^+ \isom \mathbf{P}^1_{\mathbf{C}} \setminus \{\infty\} = Y , \leqno{(29)}
$$
caractérisé par la condition que l'orbite critique de ramification $3$ s'envoie sur le point $0$, et celle de ramification $2$ au point $1$. On trouve ainsi un isom.
(posant $\mathcal{G} = \operatorname{Gl}(2,\mathbf{Z})$, $G = \operatorname{Gl}(2,\mathbf{Z})^+ = \operatorname{Sl}(2,\mathbf{Z})$)
\marginpar{$(x \in \overset{\circ}{D}!)$}
$$
\begin{tikzcd}
\pi_1(D, G, x) \arrow[r, "\sim"] \arrow[d, equal] & \pi_1(Y, \underline{d} ; y) \\
G &
\end{tikzcd} \leqno{(30)}
$$
que je voudrais expliciter, en identifiant les $\tau'_0, \tau'_1, \tau'_\infty$ comme éléments de $\operatorname{Sl}(2,\mathbf{Z})$.

\newpage

%% ===== Page 22 =====

39 Les conditions de congruence pour l'opération extérieure de $\Gamma_{\mathbb{Q}}$ sur $\hat{\Pi}_{0,3}$.

Considérons le diagramme

(1)
\begin{tikzcd}
1 \ar[r] & \underset{\simeq \operatorname{Sl}(2, \mathbb{Z})^\wedge}{\hat{\mathcal{T}}_{1,1}^+} \ar[r] \ar[d] & \mathcal{N}_{1,1} \ar[r] \ar[d] & \Gamma_{\mathbb{Q}} \ar[r] \ar[d, "\chi"] & 1 \\
1 \ar[r] & \operatorname{Sl}(2, \hat{\mathbb{Z}}) \ar[r] & \operatorname{Gl}(2, \hat{\mathbb{Z}}) \ar[r, "det"] & \hat{\mathbb{Z}}^* \ar[r] & 1
\end{tikzcd}

Il montre que l'opération extérieure de $\Gamma_{\mathbb{Q}}$ sur $\hat{\mathcal{T}}_{1,1}^+ \simeq \operatorname{Sl}(2, \mathbb{Z})^\wedge$ passe au quotient en une opération extérieure sur $\operatorname{Sl}(2, \hat{\mathbb{Z}})$, et cette opération extérieure, de plus, peut s'expliciter entièrement à l'aide du caractère cyclotomique

(2)
$$ \chi : \Gamma_{\mathbb{Q}} \longrightarrow \hat{\mathbb{Z}}^* $$

Pour préciser ce dernier point, notons que la deuxième ligne de (1) donne un homomorphisme

(3)
$$ \hat{\mathbb{Z}}^* \longrightarrow Aut ext (\operatorname{Sl}(2, \hat{\mathbb{Z}})) , $$

et l'hom canonique

(4)
$$ \Gamma_{\mathbb{Q}} \longrightarrow Aut ext (\operatorname{Sl}(2, \hat{\mathbb{Z}})) $$

\newpage

%% ===== Page 23 =====

n'est autre que la composée de (2) et (3).
Par ailleurs, l'hom.\ (4) se factorise
en
$$
\hat{\mathbf{Z}}^* / \hat{\mathbf{Z}}^{*2} \longrightarrow \Autext(\operatorname{Sl}(2, \hat{\mathbf{Z}})) \leqno{(5)}
$$
(provient du fait que les éléments de
$\operatorname{Gl}(2, \hat{\mathbf{Z}})$ qui sont dans le centre
$\hat{\mathbf{Z}}^*$ opèrent trivialement sur $\operatorname{Sl}(2, \hat{\mathbf{Z}})$,
et que la composée
$$ \hat{\mathbf{Z}}^* \underset{\text{(centre)}}{\hookrightarrow} \operatorname{Gl}(2, \hat{\mathbf{Z}}) \overset{\det}{\longrightarrow} \hat{\mathbf{Z}}^* $$
n'est autre que $\lambda \mapsto \lambda^2$). Plus
précisément, on voit que pour que l'
automorphisme de $\operatorname{Sl}(2, \hat{\mathbf{Z}})$ induit par un
$u \in \operatorname{Gl}(2, \hat{\mathbf{Z}})$ soit intérieur, i.e.\ de la forme
$\operatorname{int}(u_0)$ avec $u_0 \in \operatorname{Sl}(2, \hat{\mathbf{Z}})$, il f.\ et i.s.\
qu'il existe $u_0 \in \operatorname{Sl}(2, \hat{\mathbf{Z}})$ tel que
$u u_0^{-1}$ commute à $\operatorname{Sl}(2, \hat{\mathbf{Z}})$, et cela signifie
que c'est un scalaire - i.e.\
$u \in \operatorname{Sl}(2, \hat{\mathbf{Z}}) \cdot \hat{\mathbf{Z}}^*$. En d'autres termes,
l'hom.\ (5) est injectif, et on [unclear: remarquera]

\newpage

%% ===== Page 24 =====

le noyau de (4) n'est autre que le noyau du composé
$$ 
\Gamma_{\mathbf{Q}} \xrightarrow{\chi} \hat{\mathbf{Z}}^* \to \hat{\mathbf{Z}}^*/\hat{\mathbf{Z}}^{*2} \leqno{(6)} 
$$
i.e.\ (par la théorie des corps de classes pour $\mathbf{Q}$) l'opération ext.\footnote{NB L'op.\ ext.\ de $\Gamma_{\mathbf{Q}}$ sur $\operatorname{Sl}(2,\hat{\mathbf{Z}})$ se fait par l'intermédiaire de $\hat{\mathbf{Z}}^*/\hat{\mathbf{Z}}^{*2}$.} de $\Gamma_{\mathbf{Q}}$ sur $\operatorname{Sl}(2,\hat{\mathbf{Z}})$ se factorise par une opération ext.\ du $(\Gamma_{\mathbf{Q}_{ab}})$ -- le quotient "abélien" de $\Gamma_{\mathbf{Q}}$ -- et cette dernière est fidèle.

Ainsi, si on identifie $\Gamma_{\mathbf{Q}}$ [unclear: par ex.] à un sous-groupe de $\Autext(\operatorname{Sl}(2,\mathbf{Z})^\wedge)$, ce sous-groupe satisfait à des conditions assez fortes : préserver le sous-groupe distingué noyau de
$$ 
\operatorname{Sl}(2,\mathbf{Z})^\wedge \to \operatorname{Sl}(2,\hat{\mathbf{Z}}) \leqno{(7)} 
$$
et induire sur $\operatorname{Sl}(2,\hat{\mathbf{Z}})$ l'opération extérieure provenant d'un hom.\ $\Gamma \to \hat{\mathbf{Z}}^*/\hat{\mathbf{Z}}^{*2}$.

On peut d'ailleurs préciser cet hom.\ ainsi. Considérons dans $\operatorname{Sl}(2,\hat{\mathbf{Z}})$ le "sous-groupe de Cartan" provenant des points :

\newpage

%% ===== Page 25 =====

l'infini de $\mathcal{M}_{1,1\mathbf{Q}}$, il correspond :

une classe d'homomorphisme extérieur
$$ T = T_\infty(\overline{\mathbf{Q}}) \longrightarrow \operatorname{Sl}(2, \hat{\mathbf{Z}})^{\wedge} \leqno{(8)} $$
dont l'image, disons, est (à conjugaison près)
un sous-groupe "unipotent"
$$ \begin{pmatrix} 1 & 0 \\ n & 1 \end{pmatrix} , \quad n \in \hat{\mathbf{Z}} \leqno{(9)} $$
Cet homomorphisme extérieur est respecté par l'opération de $\boldsymbol{\Gamma}$, d'où une opération bien déterminée de $\boldsymbol{\Gamma}$ sur $T$
$$ \boldsymbol{\Gamma} \xrightarrow{\chi} \operatorname{Aut}(T) \simeq \hat{\mathbf{Z}}^* \leqno{(10)} $$
qui sera justement le caractère cyclotomique. De façon plus précise, pour \marginpar{[unclear: faisant l'identification]}
tout automorphisme extérieur $u$ de $\operatorname{Sl}(2, \mathbf{Z})^{\wedge}$ (ou seulement de $\operatorname{Sl}(2, \hat{\mathbf{Z}})$)
qui invariante la classe de "conjugaison" de l'image de (8) [ou seulement de son composé $T \to \operatorname{Sl}(2, \hat{\mathbf{Z}})$], il existe un $\chi(u) \in \operatorname{Aut}(T) \simeq \hat{\mathbf{Z}}^*$ unique, tel qu'on ait commutativité dans...

\newpage

%% ===== Page 26 =====

(11)
\begin{tikzcd}
T \ar[r, "k"] \ar[d, "\chi(u)"'] & \operatorname{Sl}(2, \mathbb{Z})^\wedge \ar[d, "u"] \\
T \ar[r] & \operatorname{Sl}(2, \mathbb{Z})^\wedge
\end{tikzcd}

- ou mieux, seulement dans

(11)
\begin{tikzcd}
T \ar[r, "k"] \ar[d, "\chi(u)"'] & \operatorname{Sl}(2, \hat{\mathbb{Z}}) \ar[d, "u"] \\
T \ar[r] & \operatorname{Sl}(2, \hat{\mathbb{Z}})
\end{tikzcd}

- et l'on a une action $\chi$ (10) sur le
groupe de tous les autom. ext. $u$ de
$\operatorname{Sl}(2, \hat{\mathbb{Z}})$ qui ont la propriété envisagée.

Ceci dit, $\Gamma$ est contenu dans ce
sous-groupe $\Gamma'$, et de plus dans le sous-
-groupe des autom. ext. qui de plus
préservent le noyau de (7), et qui
opèrent (dans le quotient $\operatorname{Sl}(2, \hat{\mathbb{Z}})$)
suivant le composé
$$ \Gamma \xrightarrow{\chi} \hat{\mathbb{Z}}^* \to \hat{\mathbb{Z}}^*/\hat{\mathbb{Z}}^{*2} \to Aut ext_{loc!} (\operatorname{Sl}(2, \hat{\mathbb{Z}})) $$

Car d'ailleurs notons que le centre $\{\pm 1\}$
de $\operatorname{Sl}(2, \mathbb{Z})^\wedge$ n'étant pas trivial, l'op. ext.

\newpage

%% ===== Page 27 =====

de $\Gamma_{\mathbf{Q}}$ sur $\hat{\mathfrak{T}}_{1,1}^+$ ne permet pas à priori de reconstituer (à isomorphisme unique près) l'extension $\mathcal{N}_{1,1}$ de $\Gamma_{\mathbf{Q}}$ par $\hat{\mathfrak{T}}_{1,1}^+$ de (1). Elle donne seulement l'extension $\mathcal{N}_{1,1}'$ de $\Gamma_{\mathbf{Q}}$ par $\hat{\mathfrak{T}}_{1,1}^{+\prime} \simeq \operatorname{Sl}(2, \mathbf{Z})^\wedge / \{\pm 1\} \marginpar{\defeq \operatorname{Sl}(2, \hat{\mathbf{Z}})'}$ qui est l'image inverse de l'extension

$$ 1 \to (\operatorname{Sl}(2, \hat{\mathbf{Z}})') \to \operatorname{Aut}_{lac}(\operatorname{Sl}(2, \hat{\mathbf{Z}})') \to \Autext_{lac}(\operatorname{Sl}(2, \hat{\mathbf{Z}})') \leqno{(13)} $$

Cette extension $\mathcal{N}_{1,1}'$ a été aussi interprétée comme l'extension $\mathcal{M}_{0,3} \simeq \mathcal{N}_{0,4}'$ de $\Gamma_{\mathbf{Q}}$ par $\hat{\mathfrak{S}}_{0,3} \simeq \hat{\mathfrak{T}}_{0,4}'$ (groupe qui est lui-même extension de $\mathfrak{S}_3 \simeq \operatorname{Sl}(2, \mathbf{Z}/2\mathbf{Z})$ par $\hat{\pi}_{0,3}$).

Nous savons déjà que l'opération extérieure de $\Gamma_{\mathbf{Q}}$ sur le noyau $\hat{\pi}_{0,3}$ de $\operatorname{Sl}(2, \hat{\mathbf{Z}})' \to \operatorname{Sl}(2, \mathbf{Z}/2\mathbf{Z})$ est fidèle, et concordante à l'opération extérieure de $\mathfrak{S}_3 \simeq \operatorname{Sl}(2, \mathbf{Z}/2\mathbf{Z}) \simeq \operatorname{Sl}(2, \hat{\mathbf{Z}})' / \hat{\pi}_{0,3}$.

Mais cette concordance permet...

\newpage

%% ===== Page 28 =====

nant s'interpréter comme conséquence du fait que l'opération ext. de $\Gamma_{\mathbf{Q}}$ sur $\mathfrak{S}_3 \simeq \operatorname{Sl}(2, \mathbf{Z}/2\mathbf{Z})$ est triviale, elle-même (dans l'optique présente) conséquence du fait qu'elle se fait par l'intermédiaire d'une op. (fidèle) de $(\mathbf{Z}/n\mathbf{Z})^\ast_{/ \pm}$ , où pour $n=2$ $(\mathbf{Z}/n\mathbf{Z})^\ast_{/\pm} = 1$ (assertion qui est par contre fausse pour les valeurs $n \neq 1,2$ i.e. $n \geq 3$). Cette relation de commutativité, de nature assez triviale, apparait donc à présent comme un cas très particulier des relations congruentielles nettement plus subtiles engendrées par les symétries galoisiennes. On peut l'exprimer de façon directe en termes de l'opération ext. de $\Gamma_{\mathbf{Q}}$ sur $\widehat{\pi}_{0,3}$, en introduisant les noyaux des hom. $\widehat{\mathfrak{T}}_{1,1}^+ \to \operatorname{Sl}(2, m)$, qui sont des sous-groupes ouverts
$$
\pi(2m) \subset \widehat{\pi}_{0,3} \leqno{(14)}
$$
distingués et [unclear: isomorphes] au pro $\mathfrak{S}_3$, noyaux

\newpage

%% ===== Page 29 =====

d'homomorphismes
$$ \phi_{2m} : \widehat{\Pi}_{0,3} \to \operatorname{Sl}(2, \mathbf{Z}/2m\mathbf{Z})' \leqno{(15)} $$
qu'il ne doit pas être bien difficile d'expliciter, provenant par "réduction mod $2m$" d'un hom.
$$ \phi_\infty : \widehat{\Pi}_{0,3} \to \operatorname{Sl}(2, \hat{\mathbf{Z}})' \leqno{(16)} $$
dont il a été question déjà précédemment - hom. dont l'image est le sous-groupe noyau de $\operatorname{Sl}(2, \hat{\mathbf{Z}}) \to \operatorname{Sl}(2, \mathbf{Z}/2\mathbf{Z})$, i.e. on a des suites exactes
$$ 1 \to \Pi(\infty) \to \widehat{\Pi}_{0,3} \to \operatorname{Sl}(2, \hat{\mathbf{Z}}) \to \operatorname{Sl}(2, \mathbf{Z}/2\mathbf{Z}) \to 1 \leqno{(17)} $$
$$ 1 \to \Pi(2m) \to \widehat{\Pi}_{0,3} \to \operatorname{Sl}(2, \mathbf{Z}/2m\mathbf{Z}) \to \operatorname{Sl}(2, \mathbf{Z}/2\mathbf{Z}) \to 1 \leqno{(18)} $$
Ceci posé, lesdites suites exactes permettent de façon naturelle de définir une opération ext. de $\hat{\mathbf{Z}}^\times / \hat{\mathbf{Z}}^{\times 2}$ sur $\widehat{\Pi}_{0,3} / \Pi(\infty)$ , ou encore des $(\mathbf{Z}/2m\mathbf{Z})^\times_2$ sur les $\widehat{\Pi}_{0,3} / \Pi(2m)$ , et la "coalition des congruences" s'exprimera en disant que la condition cyclotomique s'exprime par la condition que l'action $\chi$ de $\boldsymbol{\Gamma}_{\mathbf{Q}}$ (déduite de l'opération ext. de

\newpage

%% ===== Page 30 =====

$\Gamma_{\mathbf{Q}}$ sur $\hat{\mathfrak{T}}_{0,3}$ (ou $\hat{\mathfrak{S}}_{0,3}$ : [unclear: base] $\hat{\Pi}_{0,3}$), [unclear: provient] d'un
$$
\Gamma_{\mathbf{Q}} \xrightarrow{\chi} \hat{\mathbf{Z}}^\times \xrightarrow{\mathrm{can.}} \hat{\mathbf{Z}}^\times / \hat{\mathbf{Z}}^{\times 2} \leqno{(19)}
$$
et que cette opération extérieure de $\Gamma_{\mathbf{Q}}$ sur
$\hat{\Pi}_{0,3} / \Pi(\infty)$, est compatible avec
l'opération extérieure de $\Gamma_{\mathbf{Q}}$ sur $\hat{\Pi}_{0,3}$ lui-même.

Pour chaque entier $m$, cette condition
s'explicite par le fait que $\Pi(2m)$ est
invariant par l'opération extérieure de $\Gamma_{\mathbf{Q}}$, et que
l'opération extérieure induite sur
$\hat{\Pi}_{0,3} / \Pi(2m)$ est celle déduite du
caractère composé de (19) avec
$\hat{\mathbf{Z}}^\times / \hat{\mathbf{Z}}^{\times 2} \to (\mathbf{Z}/2m\mathbf{Z})_2^\times$ :
$$
\Gamma_{\mathbf{Q}} \to (\mathbf{Z}/2m\mathbf{Z})_2^\times \leqno{(20)}
$$
On devrait regarder si cette formulation
est vraiment équivalente à celle des
conditions de congruences dans $\hat{\mathfrak{T}}_{1,1}^{+ '}$
les autres, étant de simples
conséquences ; exercice ! Par contre, le

\newpage

%% ===== Page 31 =====

question de savoir si ces conditions
constituent bel et bien une restriction
effective, sur un autom.\ ext.\ du
gr.\ $\pi$ : [unclear: tout] $\hat{\pi}_{0,3}$, contrairement
à l'action de $\mathfrak{S}_3$, risque fort d'être
hors de portée, faute de savoir
construire effectivement de tels autom.\footnote{Ou ne fût-ce des éléments [unclear: mystérieux] de $\hat{\mathfrak{S}}_{0,3}^*$.}
extérieurs de $\pi$ - autrement que
justement par voie arithmétique, à
l'aide de l'opération de $\boldsymbol{\Gamma}_{\mathbf{Q}}$ sur
$\hat{\pi}_{0,3}$ ! (Il faudrait quand même
refaire à l'occasion un effort en ce
sens...)

Je me pose aussi la question, dans
quelle mesure on peut reconstituer, en
termes de la théorie des groupes profinis
à lacets et de leurs autom.\ ext.\ de
$\hat{\pi}_{0,3}$, ou de $\hat{\mathfrak{T}}_{0,3} \simeq \hat{\mathfrak{T}}_{1,1}^+$, l'homo-

\newpage

%% ===== Page 32 =====

\begin{enumerate}
    \item[a)] l'extension plus fine $\mathcal{N}_{1,1}$ de $\boldsymbol{\Gamma}_{\mathbf{Q}}$ par $\hat{\mathfrak{T}}_{1,1}^+$, qui est donc une ext.\ de $\mathcal{N}_{1,1}'$
    (censé être connue) par le gr.\ $\simeq \{\pm 1\}$
    \item[b)] l'hom.\ $\mathcal{N}_{1,1} \to \operatorname{Gl}(2, \hat{\mathbf{Z}})$ de (21),
    et à défaut, l'hom.
    $$ \mathcal{N}_{1,1}' \simeq \hat{\mathcal{M}}_{0,3}^{\text{paire}} \to \operatorname{Gl}(2, \hat{\mathbf{Z}})/\{\pm 1\} \leqno{(22)} $$
\end{enumerate}

Je commence par la question b).
Si on dispose de l'extension $\mathcal{N}_{1,1}$ de $\boldsymbol{\Gamma}_{\mathbf{Q}}$ par $\hat{\mathfrak{T}}_{1,1}^+ \simeq \operatorname{Sl}(2, \hat{\mathbf{Z}})^\sim$, on a un hom.
$$ \mathcal{N}_{1,1} \to \Autext(\underset{\operatorname{Sl}(2, \hat{\mathbf{Z}})^\sim}{\hat{\mathfrak{T}}_{1,1}^+}) \leqno{(23)} $$
qui passe au quotient en un hom.
$$ \mathcal{N}_{1,1} \to \Autext(\operatorname{Sl}(2, \hat{\mathbf{Z}})) \leqno{(24)} $$
on a un hom.\ évident
$$ \operatorname{Gl}(2, \hat{\mathbf{Z}}) \to \Autext(\operatorname{Sl}(2, \hat{\mathbf{Z}})) $$
qui se factorise en
$$ \operatorname{Gl}(2, \hat{\mathbf{Z}}) \to \operatorname{Gl}(2, \hat{\mathbf{Z}})/\hat{\mathbf{Z}}^\times \simeq \mathrm{PGL}(2, \hat{\mathbf{Z}}) \hookrightarrow \Autext(\operatorname{Sl}(2, \hat{\mathbf{Z}})) \leqno{(25)} $$
et il se trouve que (24) se factorise en
$$ \mathcal{N}_{1,1} \to \underset{\simeq \operatorname{Gl}(2, \hat{\mathbf{Z}})/\hat{\mathbf{Z}}^\times}{\mathrm{PGL}(2, \hat{\mathbf{Z}})} \hookrightarrow \Autext(\operatorname{Sl}(2, \hat{\mathbf{Z}})) \leqno{(26)} $$

\newpage

%% ===== Page 33 =====

ce qui permet donc de reconstituer, à partir de l'extension $\mathcal{N}_{1,1}$, un hom
$$
\mathcal{N}_{1,1} \overset{\varphi}{\longrightarrow} \operatorname{PGL}(1, \hat{\mathbf{Z}}) \simeq \operatorname{Gl}(2, \hat{\mathbf{Z}}) / \hat{\mathbf{Z}}^* \leqno{(27)}
$$
un peu plus faible que l'hom (21) qu'on demande : [unclear: à] construire. Mais notons que l'hom.
$$
\operatorname{Gl}(2, \hat{\mathbf{Z}}) \overset{u \mapsto (\bar{u}, \det(u))}{\longrightarrow} \operatorname{PGL}(1, \hat{\mathbf{Z}}) \times \hat{\mathbf{Z}}^* \leqno{(28)}
$$
a comme noyau l'ens. des $\lambda \in \hat{\mathbf{Z}}^*$ tels que $\lambda^2 = 1$ i.e. le groupe
$$
{_2\hat{\mathbf{Z}}^*} \simeq (\mathbf{Z} / 2 \mathbf{Z})^P \quad \text{ P ens. des nbs premiers} \leqno{(29)}
$$
Ainsi, on a l'hom
$$
\mathcal{N}_{1,1} \overset{u \mapsto (\varphi(u), \chi(u))}{\longrightarrow} \operatorname{PGL}(1, \hat{\mathbf{Z}}) \times \hat{\mathbf{Z}}^* \leqno{(30)}
$$
(déduit de $(27)$ et de $\chi$)\footnote{c.à.d. un scindage de l'ext. $\mathcal{N}_{1,1}$ de la cat. ou plutôt un relev. de l'op. de $\boldsymbol{\Gamma}_{\mathbf{Q}}$ sur $\hat{\mathfrak{T}}_{1,1}^\pm$} se factorise par l'image de $(28)$, et définit donc un hom
$$
\mathcal{N}_{1,1} \overset{\psi}{\longrightarrow} \operatorname{Gl}(2, \hat{\mathbf{Z}}) / {_2\hat{\mathbf{Z}}^*} \leqno{(31)}
$$
qui est encore une version affaiblie de $(21)$.

\newpage

%% ===== Page 34 =====

À vrai dire, l'homomorphisme (31) se factorise en
$$
[\mathcal{N}_{1,1} \to] \underset{\substack{|| \\ \mathcal{N}_{1,1}/\pm 1}}{\mathcal{N}'_{1,1}} \overset{\Psi'}{\longrightarrow} \operatorname{Gl}(2, \hat{\mathbf{Z}})/_2\hat{\mathbf{Z}}^* \leqno{(32)}
$$
et l'hom. $\Psi'$ peut sûrement se construire directement à partir de l'extension $\mathcal{N}'_{1,1}$ de $\Gamma_{\mathbf{Q}}$ par $\overline{\mathfrak{T}}_{1,1} \simeq \hat{\mathfrak{S}}_{0,3}^+ \simeq \operatorname{Sl}(2, \hat{\mathbf{Z}})'$, tout comme tantôt $\varphi$ : il suffit de vérifier que l'hom. canonique
$$
\begin{tikzcd}
\mathcal{G}P(1, \hat{\mathbf{Z}}) \arrow[r] & \Autext(\operatorname{Sl}(2, \hat{\mathbf{Z}})') \\
\operatorname{Gl}(2, \hat{\mathbf{Z}})/\hat{\mathbf{Z}}^* \arrow[u, "\simeq"']
\end{tikzcd}
\leqno{(33)}
$$
est injectif. La question (question du reste [unclear: des plus futiles]...) est comment relier $\Psi'$ de (31) avec l'hom. canonique
$$
\mathcal{N}'_{1,1} \longrightarrow \operatorname{Gl}(2, \hat{\mathbf{Z}})' \defeq \operatorname{Gl}(2, \hat{\mathbf{Z}})/\pm 1 \leqno{(34)}
$$
--- une fois cet hom. obtenu, on en déduit l'ext. $\mathcal{N}_{1,1}$ par image inverse de la structure d'ext.
$$
1 \to \{\pm 1\} \to \operatorname{Gl}(2, \hat{\mathbf{Z}}) \to \operatorname{Gl}(2, \hat{\mathbf{Z}})' \to 1 \leqno{(35)}
$$
ainsi que l'hom. (21) $\mathcal{N}_{1,1} \to \operatorname{Gl}(2, \hat{\mathbf{Z}})$.

\newpage

%% ===== Page 35 =====

Donc en résumé, la question essentielle semble ici la détermination d'un relèvement (34) de (32).

Reprenons nos arguments. On se donne un sous-groupe $\Gamma$ de $\Autext(\hat{\pi}_{0,3})$, commutant à $\mathfrak{S}_3$, d'où une extension $\mathcal{M}$ (notée $\mathcal{M}_{0,3}$ ou $\mathcal{N}_{1,1}'$) de $\Gamma$ par $\hat{\mathfrak{S}}_{0,3}^+$, ou encore de $\mathfrak{S}_3 \times \Gamma$ par $\hat{\pi}_{0,3}$, d'où un homomorphisme
$$
\mathcal{M}_{0,3} (\simeq \mathcal{N}_{1,1}') \to \operatorname{Aut}(\hat{\mathfrak{S}}_{0,3}^+ \simeq \operatorname{Sl}(2,\hat{\mathbf{Z}})'^\wedge) \leqno{(36)}
$$
induisant sur $\hat{\mathfrak{S}}_{0,3}^+$ l'opération intérieure de ce groupe sur lui-même.
On a par ailleurs un hom. injectif
$$
\operatorname{Gl}(2,\hat{\mathbf{Z}}) \hookrightarrow \operatorname{Aut}(\operatorname{Sl}(2,\hat{\mathbf{Z}})') \leqno{(37)}
$$
et $\Gamma \subset \Autext(\hat{\pi}_{0,3})$ est astreint (si ce n'est déjà une conséquence de l'hyp. de commutation à $\mathfrak{S}_3$) à la condition que l'action ext.
$$
\Gamma \to \Autext(\hat{\mathfrak{S}}_{0,3}^+ \simeq \operatorname{Sl}(2,\hat{\mathbf{Z}})'^\wedge) \leqno{(38)}
$$

\newpage

%% ===== Page 36 =====

soit compatible aux actions ext. de $\Gamma$, sur $\hat{\mathfrak{S}}_{0,3}^+$ par tautologie (du fait que $\Gamma$ est formé d'él. de $\Autext(\hat{\Pi}_{0,3})$ qui commutent à $\mathfrak{S}_3$), sur $\operatorname{Sl}(2,\hat{\mathbf{Z}})'$ via $\hat{\mathbf{Z}}^* / \hat{\mathbf{Z}}^{*2}$ et le caractère "cyclotomique"
$$
\chi: \Gamma \to \hat{\mathbf{Z}}^* \leqno{(39)}
$$
déduit de l'op. ext. tautologique de $\Gamma$ sur $\hat{\Pi}_{0,3}$ (respectant la structure à lacets).
Cette "condition de congruence" signifie que (36) se factorise par (37), en un hom. [unclear: disons] déterminant
$$
\varphi': \mathcal{M}_{0,3} \to \operatorname{Gl}(1, \hat{\mathbf{Z}}) \leqno{(40)}
$$
d'où, en utilisant de plus ledit caractère cyclotomique, un hom.
$$
\mathcal{M}_{0,3} \xrightarrow{u \mapsto (\varphi'(u), \chi(u))} \operatorname{Gl}(1, \hat{\mathbf{Z}}) \times \hat{\mathbf{Z}}^* \leqno{(41)}
$$
et la condition géométrique nous assure en fait, plus précisément, que ce dernier se factorise en un hom.
$$
\mathcal{M}_{0,3} \xrightarrow{\Psi'} \operatorname{Gl}(2, \hat{\mathbf{Z}}) / \hat{\mathbf{Z}}^*_2 \leqno{(42)}
$$
compte tenu de l'injection

\newpage

%% ===== Page 37 =====

$$
\operatorname{Gl}(2, \hat{\mathbf{Z}})/\hat{\mathbf{Z}}^{*2} \hookrightarrow P\operatorname{Gl}(2, \hat{\mathbf{Z}}) \times \hat{\mathbf{Z}}^* \leqno{(43)}
$$
$$
u \bmod \hat{\mathbf{Z}}^{*2} \longmapsto (u \bmod \hat{\mathbf{Z}}^*, \det u)
$$

Considérons l'hom. can. (surjectif)
$$
\underset{\operatorname{Gl}(2, \hat{\mathbf{Z}})/\pm 1}{\underset{\defeq}{\operatorname{Gl}(2, \hat{\mathbf{Z}})'}} \longrightarrow \operatorname{Gl}(2, \hat{\mathbf{Z}})/\hat{\mathbf{Z}}^{*2} \leqno{(44)}
$$

dont le noyau est donc :
$$
\nu = (\mathbf{Z}/2\mathbf{Z})^P/\text{diag} \qquad (P = \text{ens. des nb. premiers}) \leqno{(45)}
$$

L'hom. canonique $\Psi'$ (42), déduit
des conditions de congruences, définit
par image inverse une extension centrale
$$
1 \to \nu \to \overline{\mathcal{M}}_{0,3} \to \mathcal{M}_{0,3} \to 1 \leqno{(46)}
$$
et la donnée d'un hom. $\tilde{\varphi}_1$
$$
\mathcal{M}_{0,3} \to \operatorname{Gl}(2, \hat{\mathbf{Z}})' \leqno{(47)}
$$
qui relèverait (42), équivaut à celle d'un
scindage de cette extension (46). La
question à élucider à présent : est-il
possible d'expliciter un tel relèvement canonique
(42), ou scindage de (46), en termes de la

\newpage

%% ===== Page 38 =====

seule donnée de $\Gamma \subset \Autext(\widehat{\Pi}_{0,3})$, satisfaisant aux conditions congruentielles énoncées plus haut ?

Notons que sur le sous-groupe $\hat{\mathfrak{T}}_{0,3}^+ \simeq \operatorname{Sl}(2, \hat{\mathbf{Z}})/\pm 1$ de $\mathcal{M}_{0,3}$ , le relèvement cherché (47) est déjà connu, comme étant $\operatorname{Sl}(2, \hat{\mathbf{Z}})/\pm 1 \to \operatorname{Sl}(2, \hat{\mathbf{Z}})' \subset \operatorname{Gl}(2, \hat{\mathbf{Z}})'$. En d'autres termes, le scindage cherché de (46) est déjà connu au dessus du sous-groupe $\hat{\mathfrak{T}}_{0,3}^+$ de $\mathcal{M}_{0,3}$.

L'indétermination du relèvement ou scindage est a priori dans $\operatorname{Hom}(\mathcal{M}_{0,3}, \nu)$ , tenant compte de la condition précédente elle sera dans
$$
\operatorname{Ker}( \operatorname{Hom}(\mathcal{M}_{0,3}, \nu) \to \operatorname{Hom}(\hat{\mathfrak{T}}_{0,3}^+, \nu) ) \leqno{(48)}
$$
$$
\simeq \operatorname{Hom}(\Gamma, \nu)
$$
si on suppose que $\Gamma / \Ker \chi \xrightarrow{\chi} \hat{\mathbf{Z}}^*$ est iso, l'indétermination serait encore dans $\operatorname{Hom}( (\hat{\mathbf{Z}}^*)_2, \nu)$ - c'est gros ! Il

\newpage

%% ===== Page 39 =====

faut un élément un peu astucieux
pour justifier ce relèvement --- on devine
qu'il est lié justement à l'interprétation
d'un $M_{0,3}$ comme un $M_{1,1}$ --- i.e.\ à
l'introduction, dans une ``théorie profinie
de courbes rationnelles'', de la ``théorie
profinie des courbes elliptiques''.

\bigskip

\noindent \textbf{40} \marginpar{40} Paraphrases topologiques-isotopiques
(ou en termes des groupes fondamentaux)
de la théorie de Legendre.

Pour tâter le terrain, j'ai envie d'abord
de partir de la situation des groupes
discrets, qu'on a mieux en main, c'est
le cas de le dire !

Soit $\mathcal{C}_{0,3}$ la catégorie des groupes
(discrets) à lacets de type $0,3$, \marginpar{(avec les isomorphies ext. comme flèches)} dont le
représentant-type est $\Pi_{0,3}$ --- cette catégorie
est donc équivalente canoniquement à : celle
des torseurs sous le groupe
$$
\Autext(\Pi_{0,3}) \simeq \mathfrak{S}_{0,3} \simeq \operatorname{Gl}(2, \mathbf{Z}) / \pm 1 \simeq \operatorname{PGL}(2, \mathbf{Z}) \leqno (49)
$$

Soit $\mathcal{C}_{1,1}^{ext}$ la catégorie des groupes à lacets
de type $(1,1)$, dont le représentant-type est
$$
\Pi_{1,1} = \{x,y,l \mid [x,y]l = 1\} \ ; 
$$
cette catégorie est

\newpage

%% ===== Page 40 =====

équivalente à celle des torseurs sous le groupe
$C_{1,1}^{ext} \simeq$
$\operatorname{Aut}\, ext(\Pi_{1,1}) \simeq \operatorname{Gl}(2,\mathbb{Z})$ - l'isomorphisme

(50) $$\operatorname{Aut}\, ext(\Pi_{1,1}) \longrightarrow \operatorname{Gl}(2,\mathbb{Z})$$

il s'obtient en associant à un automorphisme
extérieur de $\Pi_{1,1}$ son action sur

(51) $$(\Pi_{1,1})_{ab} \simeq \mathbb{Z}^2$$

On a donc un foncteur canonique
[MARGIN: $\mathcal{L}$, initiale de "Legendre"]
(52) $$\mathcal{L} : C_{1,1}^{ext} \longrightarrow C_{0,3}$$

défini (à isomorphisme près) par la condition
de rendre commutatif le diagramme

(53)
\begin{tikzcd}
C_{1,1}^{ext} \ar[r] \ar[d, "\wr"'] & C_{0,3} \ar[d] \\
Tors(\operatorname{Gl}(2,\mathbb{Z})) \ar[r] & Tors(GPl(1,\mathbb{Z}))
\end{tikzcd}

où
$$Tors(\operatorname{Gl}(2,\mathbb{Z})) \longrightarrow Tors(GPl(1,\mathbb{Z}))$$
est le foncteur de chgt de groupe d'opérateurs
relatif à :
(54) $$\operatorname{Gl}(2,\mathbb{Z}) \xrightarrow{can} GPl(1,\mathbb{Z}) \simeq \operatorname{Gl}(2,\mathbb{Z}) / \pm 1$$

\newpage

%% ===== Page 41 =====

On peut [unclear: aussi] introduire le groupoïde
$C_{1,1}^{ext \ \prime}$, déduit de $C_{1,1}^{ext}$ en prenant le
quotient par l'automorphisme universel $-id$
- correspondant au centre $(\pm 1)$ (des groupes
fondamentaux de ce groupoïde - on trouve
un groupoïde connexe équivalent canon.
[MARGIN: groupoïde] à celui des torseurs sous $\operatorname{Gl}(2,\mathbb{Z})'= \operatorname{Gl}(2,\mathbb{Z})/(\pm 1)$ ,

dont on tire maintenant une
équivalence de catégories

(55) $$C_{1,1}^{ext \ \prime} \overset{\simeq}{\longrightarrow} C_{0,3}$$

rendant commutatif

(56)
\begin{tikzcd}
C_{1,1}^{ext \ \prime} \ar[r, "\simeq"] \ar[d] & C_{0,3} \ar[d] \\
Tors(\operatorname{Gl}(2,\mathbb{Z})') \ar[r, "\simeq"] & Tors(GPl(1,\mathbb{Z}))
\end{tikzcd}
,

où $Tors(\operatorname{Gl}(2,\mathbb{Z})') \overset{\simeq}{\longrightarrow} Tors(GPl(1,\mathbb{Z}))$ est
le foncteur tautologique associé à l'iso

(57) $$\operatorname{Gl}(2,\mathbb{Z})' \overset{\simeq}{\longrightarrow} GPl(1,\mathbb{Z})$$ .

\newpage

%% ===== Page 42 =====

En somme, je suis en train de faire une paraphrase topologique tautologique, en termes des groupes fondamentaux à lacets, de la théorie de Legendre des relations entre courbes elliptiques et courbes rationnelles. J'aimerais expliciter maintenant l'équivalence de Legendre (55), qui avait été obtenue grâce à la connaissance des groupes de Teichmüller discrets relatifs aux deux membres et d'un isomorphisme ad-hoc entre ces deux groupes, directement en termes des groupes fondamentaux de type $\Pi_{0,3} - \Pi_{1,1}$.

Si $\pi'$ est un groupe à lacets de type $1,1$, il convient d'abord de noter qu'il y a dans $\Autext(\pi')$\marginpar{NB $\mathfrak{S}_{1,1}$ conserve l'orientation, i.e.\ est sous-groupe d'indice 2 de $\Autext(\pi')$.} un élément canonique d'ordre $2$ --- en fait l'unique élément non trivial du centre, noté commodément $(-1)_{\pi'}$.

\newpage

%% ===== Page 43 =====

Ainsi, on peut considérer $\pi'$ comme muni\footnote{can!}
d'une opération extérieure de
$G = \mathbf{Z}/2\mathbf{Z}$ sur lui, d'où une extension
canonique
$$ 1 \to \pi' \to \mathcal{E} \to \mathbf{Z}/2\mathbf{Z} \to 1 \leqno (58) $$

Le modèle topologique nous apprend
qu'il y a exactement trois classes de $\pi'$-
-conjugaison de scindages de cette
extension, correspondant aux trois
points d'ordre $2$ d'un groupe topologique
$E \simeq \mathbf{U} \times \mathbf{U}$ -- les trois points fixes de
$x \mapsto -x$ dans $E \setminus \{0\}$. La théorie des

"franges et des scellements", relative à une
situation de groupes finis d'opérations exté-
rieures sur des groupes à lacets (théorie qui\marginpar{mais rien n'a encore été écrit sur le fond d'un tel groupe des franges, des conjectures...}
dans le cas discret n'a vraiment
été écrite -- il y a pour le moment des
conjectures...) permet alors

\newpage

%% ===== Page 44 =====

de construire canoniquement un groupe (noté $\tilde{\pi}'$) de type $\pi_{1,4}$, avec une opération extérieure de $G = \mathbf{Z}/2\mathbf{Z}$ sur ce groupe, et une partie $J$ de cardinal 3 \marginpar{en fait $G$ opère trivialement sur $I = I(\tilde{\pi}')$} de $I = I(\tilde{\pi}')$ (qui est de cardinal 4), $J$ stable par $G$ i.e. de pt fixe $i_0$, et un iso. du groupe $\tilde{\pi}'(J)$ (déduit par "bouchage des trous dans $J$") avec $\pi'$
$$
\tilde{\pi}'(J) \isom \pi' \leqno{(59)}
$$
qui soit compatible avec les opérations de $G$. Dans le modèle topologique, $\tilde{\pi}'$ est le groupe fond. de l'espace $E \setminus {}_2E$, sur lequel $G$ opère librement.

L'extension $\tilde{\mathcal{E}}$ de $G$ par $\tilde{\pi}'$ correspondante
$$
1 \to \tilde{\pi}' \to \tilde{\mathcal{E}} \to \mathbf{Z}/2\mathbf{Z} \to 1 , \leqno{(60)}
$$
s'identifie donc au groupe fondamental mixte de $(E \setminus {}_2E, G)$, s'identifie de [unclear: même]

\newpage

%% ===== Page 45 =====

au groupe fondamental de
$$
V = (E \setminus {}_2E)/\{\pm 1\} \leqno (61)
$$
qui est une surface top.\ de type $(0,4)$.

En d'autres termes, $\widetilde{\mathcal{E}}$ peut être considéré comme un groupe à lacets de type $(0,4)$, pour une structure à lacets facile à décrire, se déduisant de celle de $\widetilde{\pi'}$, et pour laquelle on a un iso.\ canonique
$$
I(\widetilde{\mathcal{E}}) = I(\widetilde{\pi'})/G \simeq I(\widetilde{\pi'}) \leqno (62)
$$

Ainsi on dispose d'un $i_0 \in I(\widetilde{\mathcal{E}})$, ce qui permet de "boucher le trou en $i_0$" pour obtenir un groupe à lacets effectif de type $(0,3)$, soit
$$
\pi = \mathcal{L}(\pi') \leqno (63)
$$
Reprenant un peu encore la relation entre $\pi$ et $\pi'$, on peut procéder ainsi : d'un groupe à lacets effectif de type $(0,3)$, on passe par

\newpage

%% ===== Page 46 =====

le foncteur "forage d'un trou" ($\pi \mapsto \widetilde{\pi}$) un groupe
à lacets extérieur de type (0,4), muni d'
un élément $i_0 \in I(\widetilde{\pi})$ tel que
$$ I(\widetilde{\pi}) \isom \{i_0\} \amalg \underset{\text{(cardinal 3)}}{I(\pi)} \leqno{(64)} $$
On a donc ($T$ étant le module d'orientation
de $\pi$)
$$ \widetilde{\pi}_{ab} \isom T^{I(\widetilde{\pi})} \qquad ((\widetilde{\pi}_{ab})_2 \isom (\mathbf{Z}/2\mathbf{Z})^{I(\widetilde{\pi})}) \leqno{(65)} $$
[unclear: il s'en déduit] un homomorphisme (extérieur)
$$ \widetilde{\pi} \longrightarrow \mathbf{Z}/2\mathbf{Z} \leqno{(66)} $$
canonique, défini par le fait d'être
"ramifié" en tous les $i \in I(\widetilde{\pi})$ - i.e. correspond
par l'homomorphisme
$$ (\widetilde{\pi}_{ab})_2 \isom (\mathbf{Z}/2\mathbf{Z})^{I(\widetilde{\pi})} \longrightarrow \mathbf{Z}/2\mathbf{Z} \leqno{(67)} $$
somme des coefficients. Le noyau $\widetilde{\pi}'$ de
(66) est un groupe à lacets de type (1,4)
- mais il n'est pas tout à fait un groupe
extérieur - il est par définition "à un signe
près", défini seulement modulo (au plus)
l'opération extérieure du quotient $\mathbf{Z}/2\mathbf{Z}$

\newpage

%% ===== Page 47 =====

dessus. Si on [unclear: dévisse] le groupe $\widetilde{\pi}$
(dans sa "classe extérieure") - on a donc une suite exacte
$$
1 \to \widetilde{\pi}' \to \widetilde{\pi} \to \mathbf{Z}/2\mathbf{Z} \to 1 \leqno{(68)}
$$
où $\widetilde{\pi}'$ est un groupe à lacets, avec
$$
I(\widetilde{\pi}') \simeq I(\widetilde{\pi}) \leqno{(69)}
$$
En bouchant les trous relatifs à 
$I(\infty) \subset I(\widetilde{\pi}) \simeq I(\widetilde{\pi}')$ , qui est de cardinal $3$,
on trouve un groupe extérieur de type $(1,1)$, avec une opération de $\mathbf{Z}/2\mathbf{Z}$
qui n'est autre que l'op. canonique ci-dessus. Mais [unclear: la justification des deux faits], c'est un peu [unclear: abusif] pour un groupe extérieur de type
$(1,1)$ i.e. un objet de $\mathcal{C}_{1,1}^{ext}$ - i.e. une [unclear: telle] classe "modulo symétrie" - i.e. un
objet de $\mathcal{C}_{1,1}^{ext}$. Pour avoir un objet
de $\mathcal{C}_{1,1}^{ext}$, il faut avoir "choisi un [unclear: noyau] $\widetilde{\pi}'$"
pour l'hom. extérieur $\widetilde{\pi} \to \mathbf{Z}/2\mathbf{Z}$, outre qu'un
groupe extérieur c'est la classe d'isomorphie.

\newpage

%% ===== Page 48 =====

dans une $(\mathbf{Z}/2\mathbf{Z})$-gerbe, i.e. un groupoïde connexe admettant $\mathbf{Z}/2\mathbf{Z}$ comme groupe fondamental.

Modulo une rédaction plus soignée, il me semble que ces réflexions incomplètes prouvent bel et bien un isomorphisme canonique (aux automorphismes intérieurs près)
$$
\mathcal{C}_{1,1}/\{\pm 1\} \xrightarrow{\sim} \mathfrak{S}_{0,3} \leqno{(70)}
$$
\marginpar{[unclear: en fait, justifié par un isomorphisme prouvant $\simeq \Pi_{0,3}$]}
— la détermination [unclear: de ce dernier] provenant du choix d'isomorphisme entre $\mathcal{L}(\Pi_{1,1})$ et $\Pi_{0,3}$.

Ce fait est tout à fait indépendant de celui que l'hom. naturel 
$$
\mathcal{C}_{1,1} \xrightarrow{\sim} \operatorname{Gl}(2, \mathbf{Z})
$$
soit un isomorphisme. Je n'y vois aucun inconvénient à identifier purement et simplement $\mathcal{C}_{1,1}$ à $\operatorname{Gl}(2, \mathbf{Z})$ ; par contre il est plus sage de continuer d'identifier $\mathfrak{S}_{0,3}$ à $\mathcal{C}_{1,1}/\{\pm 1\}$.

\newpage

%% ===== Page 49 =====

Si ce n'est seulement comme groupe
extérieur, à cause de l'indétermination
provenant du signe dans (70).

D'un point de vue des espaces de
modules analytiques, on a un
isomorphisme canonique \marginpar{avec un peu plus de soin, i.e. en [unclear: permettant] d'inclure dans la [unclear: famille] des complexes, (au lieu de les banir), ça donnerait un}
$$ M_{1,1}^{\text{an}} \longrightarrow \underline{M}_{0,3}^{\text{an}}\footnote{courbe analytique complexe de type (0,3) ``universelle''} \simeq M_{0,4}^{\text{an}} = (\widetilde{U}_{0,3} / \mathfrak{S}_3) \leqno (71) $$

qui induit un isomorphisme
$$ M_{1,1}^{\sim} \isommap M_{0,3}^{\sim} \leqno (72) $$

et l'effet de (71) sur les groupes fondamen-
-taux extérieurs n'est autre que l'hom.
$$ \mathfrak{T}_{1,1}^+ \longrightarrow \mathfrak{S}_{0,3}^+ \leqno (73) $$

induit par (70) --- est-il possible de
le préciser de façon utile, en un hom.
effectif, non seulement extérieur ? Il y
a un revêtement universel canonique
$\widetilde{M}_{1,1}$ de $M_{1,1}$, fourni par l'espace des

\newpage

%% ===== Page 50 =====

Teichmüller (le demi-plan de Poin\-caré $D$ du \S~38) - relatif à ce revêtement
universel - bel et bien un iso.\ can.
$$
\pi_1(\widetilde{M}_{1,1}) \simeq \operatorname{Sl}(2, \mathbf{Z}) \leqno{(74)}
$$
(pas seulement extérieur), qui n'est autre
d'ailleurs, essentiellement, que (50) -
(moyennant un sous-groupe d'indice 2
$\mathfrak{T}_{1,1}'$ de $\mathfrak{T}_{1,1}$). \marginpar{Pour la définition exacte (qui va suivre), y a} On a bel et bien un
morphisme canonique
$$
D \longrightarrow U_{0,3}^{\mathrm{an}} \quad (\simeq M_{1,1}[2]') \leqno{(75)}
$$
compatible avec un hom.\ sur les groupes
d'opérateurs
$$
\operatorname{Sl}(2, \mathbf{Z})' \longrightarrow \operatorname{Sl}(2, \mathbf{Z}/2\mathbf{Z}) \simeq \mathfrak{S}_3 \leqno{(76)}
$$
\footnote{c'est ça qu'il convient de rendre explicite dans \S~38.}
de sorte que $D$ apparaît comme un
rev.\ universel de $U_{0,3}^{\mathrm{an}}$, donc aussi
de $(U_{0,3}^{\sim}, \mathfrak{S}_3)$, et $\operatorname{Sl}(2, \mathbf{Z})'$ comme
le $\pi_1$ mixte relatif à ce rev.\ universel-là.
Pour l'identifier à $\pi_1( (U_{0,3}^{\sim}, \mathfrak{S}_3), P=\{ \dots \} )$, il

\newpage

%% ===== Page 51 =====

faut donc choisir un point dans $D$ au dessus
de $-\bar{j}$ (qui permettra alors d'identifier
$D$ au rev. universel pointé en $P$ de
$U_{0,3}$) - ce qui introduit justement
l'indétermination par un élément du
noyau de $\operatorname{Sl}(2, \mathbf{Z})' \to \operatorname{Sl}(2, \mathbf{Z}/2) \simeq \mathfrak{S}_3$,
lui-m. isomorphe à $\Pi_{0,3}$. Mais
on peut prendre p. ex.
$$
j = \exp \frac{2 i \pi}{3} = \frac{-1 + i \sqrt{3}}{2} \leqno{(77)}
$$
comme représentant pt. en dessus de
$P=-\bar{j}$ - en trichant un brin pour
qu'il soit bien en dessus de $P$, et non
de $\bar{P} = -j$, sinon c'est $-\bar{j} = \frac{1+i\sqrt{3}}{2} = \exp \frac{2 i \pi}{6}$
qu'on choisirait. [ NB On trouve que
les trois éléments de $\operatorname{Sl}(2, \mathbf{Z})'$
qui transforment $j$ en $j$ sont
$\begin{pmatrix} 1 & -1 \\ 1 & 0 \end{pmatrix}$, $\begin{pmatrix} 1 & 0 \\ 1 & 1 \end{pmatrix}$, $\begin{pmatrix} 0 & -1 \\ 1 & 0 \end{pmatrix}$, dont aucun
n'est $\equiv \begin{pmatrix} 1 & 0 \\ 0 & 1 \end{pmatrix} \pmod 2$, ce qui montre que
$j$ et $-\bar{j}$ [unclear: ont des images distinctes dans un rev. par congruences] ]

\newpage

%% ===== Page 52 =====

mod le noyau $\Gamma(2)$ de $\operatorname{Sl}(2,\mathbf{Z}) \to \operatorname{Sl}(2,\mathbf{Z}/2\mathbf{Z})$, de sorte qu'ils n'ont pas même image dans $U_{0,3} \simeq D/\Gamma(2)$, d'où leurs images respectives $P = - \bar{j}$ et $\bar{P} = -j$ (ou inversement) ] Le choix en question est alors caractérisé par le fait que $\tau \in D$ au-dessus de $P$ est alors une racine de l'unité, ce qui est quand même bien naturel. On peut dire donc qu'on a "normalisé" l'hom. (73) en un hom. effectif --- et de même pour (70).

On aurait envie de reprendre les constructions précédentes dans le contexte des groupes profinis --- malheureusement je manque des fondements nécessaires, et surtout d'une théorie des "tresses" --- "groupes de tresses" au niveau groupes à lacets profinis --- même des [unclear: espaces de configuration]

\newpage

%% ===== Page 53 =====

turale. Pourtant on sent que c'est
bien dans cette direction que se trouve
une compréhension "gruppentheoretisch"
de l'homomorphisme (40)
$$
\hat{\mathcal{M}}_{0,3} \to \operatorname{Gl}(2, \hat{\mathbf{Z}})/\pm 1 \simeq \text{[unclear: P\operatorname{Gl}(2, \hat{\mathbf{Z}})]} ,
$$
qui est la motivation principale des
réflexions du présent $\S$, et du précédent.
Ce devrait être une incitation pour une
réflexion renouvelée sur la question du
tressage dans le contexte profini... Mais
je me rends compte dès maintenant
que la question - s'il y a une
extension des constructions du
présent $\S$, au contexte profini, indépen-
damment de l'introduction d'éléments
de structure "arithmétiques" (tels que
la donnée d'un sous-groupe $\Gamma \subset \Autext(\hat{\Pi}_{0,3})$
satisfaisant à un certain nombre de
conditions $\pm$ semblables, genre relations
b-géométriques) est subordonnée à une

\newpage

%% ===== Page 54 =====

si ces relations congruentielles sont bel et
bien des relations de nature arithmétique
spécifiques aux autom. ext. provenant
de $\boldsymbol{\Gamma}_{\mathbf{Q}}$, ou au contraire des
conséquences $\pm$ automatiques d'un
[unclear: pur] (purement "algébrique") qui serait
formalisme des groupes profinis, calqué
sur le cas discret. Sans préjuger aucunement
de la réponse et du degré de
plausibilité de une telle formalisation, je
pourrais à titre expérimental [unclear: essayer]
d'en dégager un, pour tester où il
mènera.

\chapter*{\S \space 41. --- CONDITIONS DE COMPATIBILITÉ SUR $\widehat{\Pi}_{0,3} \isommap \widehat{\mathfrak{S}}_{0,3}^+$ AUX OPÉRATIONS DE $\boldsymbol{\Gamma}_{\mathbf{Q}}$. CALCUL D'AUTOMORPHISMES EXTÉRIEURS DE $\widehat{\Pi}_{0,3}$ ET $\widehat{\mathfrak{S}}_{0,3}^+$.}\thispagestyle{empty}
\addcontentsline{toc}{section}{41. Conditions de compatibilité sur $\widehat{\Pi}_{0,3} \isommap \widehat{\mathfrak{S}}_{0,3}^+$ aux opérations de $\boldsymbol{\Gamma}_{\mathbf{Q}}$. Calcul d'automorphismes extérieurs de $\widehat{\Pi}_{0,3}$ et $\widehat{\mathfrak{S}}_{0,3}^+$}
\label{sec:41}
\section*{}

Avant de me lancer dans cette
(tentative), je voudrais encore élucider
une question, concernant d'autres
éventuelles conditions nécessaires (en plus
des conditions congruentielles du $§ 38$)
pour qu'un autom. ext. de $\widehat{\Pi}_{0,3}$ provienne

\newpage

%% ===== Page 55 =====

de $\Gamma_{\mathbf{Q}}$. Pour ceci, reprenons l'hom. entièrement canonique $\varphi$
$$
\begin{tikzcd}
\Pi_{0,3} \arrow[r, "\text{épi}"] \arrow[rr, bend left=20, "\varphi"] & \Pi_{0,3}' \arrow[r, "\sim"] & \mathfrak{S}_{0,3}^+ \simeq \pi_1(\dots)
\end{tikzcd} \leqno{(1)}
$$
$$ \Pi_{0,3}' = \Pi_{0,3}/\{l_1^2, l_\infty^3\} \simeq \text{[unclear: gr. des autom. d'un polygone triangulé orienté]} $$

défini par
$$
\begin{cases}
\varphi(l_0) = l_0' = \sigma_\infty \rho \\
\varphi(l_1) = l_1' = \sigma_\infty \\
\varphi(l_\infty) = l_\infty' = \rho^{-1}
\end{cases} \leqno{(2)}
$$

(avec les notations de [C, \dots] - ),

qui s'étend par continuité en un hom. profini surjectif
$$
\begin{tikzcd}
\hat{\varphi}: \hat{\Pi}_{0,3} \arrow[r, "\text{épi}"] & \widehat{\mathfrak{S}}_{0,3}^+ \simeq \pi_1((\mathcal{M}_{0,3})_{\overline{\mathbf{Q}}}, P) \arrow[d, hook] \\
& \widehat{\mathfrak{T}}_{0,3} = \Autext(\hat{\Pi}_{0,3})
\end{tikzcd} \leqno{(3)}
$$

Cet homomorphisme est défini via les «~formules brutales~» (2), sans référence à un sous-groupe particulier $\Gamma$ de $\Autext(\hat{\Pi}_{0,3})$. Si on se donne un tel sous-groupe, contenant $\mathfrak{S}_3$, celui-ci opère automatiquement sur le groupe extérieur $\widehat{\mathfrak{S}}_{0,3}^+$, et il

\newpage

%% ===== Page 56 =====

318

s'insère dans une suite exacte

(4) $1 \longrightarrow \widehat{C}_{0,3}^+ \longrightarrow M_{0,3} \longrightarrow \Gamma \longrightarrow 1$ .

comparons à la suite ex. de

(5)
\begin{tikzcd}
1 \ar[r] & \widehat{\Pi}_{0,3} \ar[r] \ar[d] & \mathcal{E}_{0,3} \ar[r] \ar[d, "p"] & \Gamma \ar[r] \ar[d, equal] & 1 \\
1 \ar[r] & \widehat{C}_{0,3}^+ \ar[r] & M_{0,3} \ar[r] & \Gamma \ar[r] & 1
\end{tikzcd}

Le fait que l'homomorphisme
(3) commute à l'action de $\Gamma_{\mathbb{Q}} = \Gamma$
me semble bien a priori une condi-
tion non triviale sur $\Gamma$ - plus
précisément, sur les éléments de
$\Gamma$ séparément, qu'il faudrait bien
arriver à expliciter ([unclear: c a d par les]
conditions congruentielles), en
termes des calculs $\pm$ explicites
de [C]. Cette condition équivaut
à celle que l'hom (3) s'étende
en un homomorphisme d'extensions

\newpage

%% ===== Page 57 =====

[MARGIN: NB. $\mathcal{E}_{0,3}$ est le groupe noté $\mathcal{M}_{0,3}$ dans le § 46]

(6)
\begin{tikzcd}
1 \ar[r] & \hat{\Pi}_{0,3} \ar[r] \ar[d, "\hat{p}"] & \mathcal{E}_{0,3} \ar[r] \ar[d, "\underline{\Phi}"] & \Gamma \ar[r] \ar[d, equal] & 1 \\
1 \ar[r] & \hat{C}_{0,3}^+ \ar[r] & \mathcal{M}_{0,3} \ar[r] & \Gamma \ar[r] & 1 \quad ,
\end{tikzcd}

qui s'identifie à
moyennant le fait que $\mathcal{M}_{0,3}$ (interprété comme
groupe fondamental "à signatures" de
$(\mathbb{P}_{\overline{\mathbb{Q}}}^1 - \{0\}, 2 \{1\} + 3 \{\infty\} ) )$ s'identifie au
quotient de $\hat{\Pi}_{0,3}$ par le sous-groupe
fermé
invariant engendré par $l_1^2, l_\infty^3$.

Je commence à soupçonner
qu'il y aura toute une kyrielle
de conditions de même nature sur
$\Gamma$, qu'il faudrait bien arriver
à dégager de façon plus générale.

\newpage

%% ===== Page 58 =====

Je vais reprendre les calculs faits dans [C], concernant les autom. extérieurs de $\hat{\Pi}_{0,3}$ qui commutent à $\mathfrak{S}_3$, qui peuvent se réaliser de façon unique comme automorphismes \emph{effectifs} qui commutent à $\sigma_{\infty}$ (et qui de plus commutent extérieurement à $\rho$). Ils correspondent biunivoquement aux couples
$$
(g_0, p) \quad \begin{cases} g_0 \in \hat{\Pi}_{0,3} \mod L_0 \\ p \in \hat{\mathbf{Z}}^* \text{ (multiplicateur)} \end{cases} \leqno{(7)}
$$
satisfaisant les conditions suivantes (où $g_0 \in \hat{\Pi}_{0,3}$ est déterminé de façon unique par $u$ (ou encore, par l'autom. extérieur correspondant)), où on a posé
$$
h_0 = \sigma_{\infty}(g_0)^{-1} g_0 = [\sigma_{\infty}, g_0^{-1}] : \leqno{(8)}
$$

\footnote{NB l'autom. $A$ de $\hat{\Pi}_{0,3}$ (où $\sigma_{\infty} = \text{id}$) prov. par $g_0(p)$ comm. ssi le fait $h_0 \dots$}
\footnote{NB fausse et rétablie [unclear: plus simplement par :]}
$$
\operatorname{int}(h_0)(l_0^p) \operatorname{int}(l_1^p \rho(h_0)^{-1})(l_{\infty}^p) = 1 \leqno{(9)}
$$
$$
h_0 l_0^{-p} \rho(h_0) = l_1^p \rho(h_0)^{-1} [unclear: l_{\infty}^{-p}] \leqno{(10)}
$$

L'hom. $u$ est donné par :\footnote{NB (10) s'écrit aussi : $$ (l_0^p h_0^{-1}) \rho(l_0^p h_0^{-1}) \rho^2(l_0^p h_0^{-1}) = 1 \leqno{(10 \text{ bis})} $$}

\newpage

%% ===== Page 59 =====

$$
\begin{cases}
u(l_0) = \operatorname{int}(g_0)(l_0^\rho) \\
u(l_1) = \operatorname{int}(\sigma_\infty(g_0))(l_1^\rho) \\
u(l_\infty) = \operatorname{int}(g_0 l_0^{-\sigma} \rho(h_0))(l_\infty^\rho)
\end{cases} \leqno (11)
$$

\marginpar{(11 bis)}et il faut en toute rigueur, en plus des conditions $(9),(10)$, exiger que les deux éléments donnés par $(11)$ $u(l_0), u(l_1)$ forment des générateurs du groupe profini $\hat{\Pi}_{0,3}$.

Je vais calculer l'action de $u$ sur $\hat{\mathfrak{S}}_{0,3}$, en considérant ce groupe comme engendré par les éléments
$$
l'_1 = \sigma_\infty, \quad l'_\infty = \rho^{-1} \leqno (12)
$$

Il faut calculer simplement $u \sigma_\infty u^{-1}$ et $u \rho u^{-1}$. On a par construction ($u$ commutant à $\sigma_\infty$)
$$
u \sigma_\infty u^{-1} = \sigma_\infty \leqno (13)
$$

et d'autre part on vérifie que
$$
\begin{cases}
\rho^{-1} u \rho = \operatorname{int}(h) u \quad \text{i.e.} \quad u \rho u^{-1} = \operatorname{int}(\rho(h)) \rho \\
\text{avec} \\
h = (\rho^{-1}\sigma_\infty)(g_0) l_0^\rho g_0^{-1} \quad \text{i.e.} \quad \rho(h) = \sigma_\infty(g_0) l_1^\rho \rho(g_0)^{-1}
\end{cases} \leqno (14)
$$

[unclear: donc on trouve, en termes de] l'opération sur le groupe $\hat{\mathfrak{S}}_{0,3}^+$ engendré par $\sigma_\infty$ et $\rho^{-1}$
\marginpar{(cas $\rho$)}
$$
\begin{cases}
u(\sigma_\infty) = \sigma_\infty \\
u(\rho) = g \cdot \rho \quad \text{avec} \quad g = \sigma_\infty(g_0) l_1^\rho \rho(g_0)^{-1}
\end{cases} \leqno (15)
$$

\newpage

%% ===== Page 60 =====

NB La relation $u(\rho)^3 = u(\sigma_\infty)^2 = 1$ est triviale,
la relation $u(\rho)^3 = 1$ se réduit à
$$
g \rho(g) \rho^2(g) = 1 \leqno (16)
$$
qui est conséquence de l'expression explicite
(15) de $g$, et de la formule (10).

Notons qu'on \marginpar{devrait avoir} \footnote{[unclear: à retenir]}
$$
l'_0 = \sigma_\infty \rho = \varepsilon_0 \quad (\text{d'où } l'_0{}^2 = l_0)
$$
\marginpar{?}
$$
u(l'_0) = \operatorname{int}(g_0)(l'_0{}^p) \leqno (17)
$$
-- cette relation entraînant à son tour que
la première relation (11)
$$
u(l_0) = \operatorname{int}(g_0)(l_0{}^p) .
$$
[La relation (17) en est une conséquence,
si on admet que le centralisateur
de $l_0 = l_0{}^{\hat{\mathbf{Z}}}$ dans $\hat{\Pi}_{0,3}$ est réduit à $l_0$... ]

On devrait avoir aussi \marginpar{compte tenu des} les relations
$$
u(\rho) \sim \rho^p \quad \text{dans } \hat{\Pi}_{0,3} \leqno (18)
$$
\marginpar{[unclear: plus tard]} relations sur lesquelles nous reviendrons.

Il n'est pas clair pour moi qu'elles
soient conséquences des conditions envisagées
jusqu'à présent.

\newpage

%% ===== Page 61 =====

Considérons maintenant l'hom.\ (13)
$$
\begin{cases}
\hat{\varphi} : \hat{\Pi}_{0,3} \to \hat{\mathfrak{S}}_{0,3}^+ \\
l_0 \mapsto l'_0 = \sigma_\infty \rho = \varepsilon_0 \\
l_1 \mapsto l'_1 = \sigma_\infty \\
l_\infty \mapsto l'_\infty = \rho^{-1}
\end{cases} \leqno{(18)}
$$
et écrivons qu'il commute, i.e.\ à un autom.\ int.\ près, à $u$. Cela signifie donc qu'il existe un
$$
\lambda \in \hat{\mathfrak{S}}_{0,3}^+ \leqno{(19)}
$$
tel que l'on ait
$$
\hat{\varphi}(u(x)) = \lambda u(\hat{\varphi}(x)) \lambda^{-1} \qquad \forall x \in \hat{\Pi}_{0,3} \leqno{(20)}
$$
et il suffit d'exprimer cette condition pour $x$ égal à l'un des deux gén.\ $l_0, l_1$ de $\hat{\Pi}_{0,3}$.

Pour $x = l_0$ la relation (20) s'écrit
$$
u(l'_0) = \operatorname{int}(\lambda^{-1} \hat{\varphi}(g_0))(l'_0{}^p) \leqno{(21)}
$$
ce qui, comparé avec (17), donne
$$
\operatorname{int}(g_0'^{-1} \lambda^{-1} \hat{\varphi}(g_0))(l'_0{}^p) = l'_0{}^p \quad \text{d'où}
$$
$$
g_0'^{-1} \lambda^{-1} \hat{\varphi}(g_0) = l'_0{}^\beta \qquad (\beta \in \hat{\mathbf{Z}})
$$
ou encore
$$
\lambda = \hat{\varphi}(g_0) l'_0{}^{-\beta} g_0'^{-1} \quad [= \hat{\varphi}(g_0 l_0^{-\beta}) = \hat{\varphi}([unclear: \sigma_\infty \dots])] \leqno{(22)}
$$
\marginpar{cf plus bas (24 bis)}

Pour $x = l_1$ la relation (20) s'écrit

\newpage

%% ===== Page 62 =====

$$
u(l'_1) = \operatorname{int}(\lambda^{-1} \hat{\varphi}(q_1)) (l'_1 \rho) \leqno{(23)}
$$
\marginpar{où $q_1 = q_0 \rho$}

on aura 
$$
\sigma_\infty = \operatorname{int}(\lambda^{-1} \hat{\varphi}(q_1)) (\sigma_\infty)
$$

vu (22) :\footnote{Garde? Il faut bien sûr exprimer que ça [unclear: coincide] avec l'expression de $\lambda$ obtenue en (22) \dots}
$$
\lambda^{-1} \hat{\varphi}(q_1) = q_0 {l'_0}^\beta \hat{\varphi} (\underbrace{q_0^{-1} \sigma_\infty (q_0)}_{h_0^{-1}}) = q_0 {l'_0}^\beta \hat{\varphi}(h_0^{-1})
$$

on arrive à la relation :
$$
\sigma_\infty \big( \hat{\varphi}(q_0) . {l'_0}^{-\beta} . q_0^{-1} \big) = \hat{\varphi}(q_0) . {l'_0}^{-\beta} . q_0^{-1}
$$

qui s'écrit aussi :
$$
\sigma_\infty(\hat{\varphi}(q_0)) . \overset{\varepsilon_1^{-\beta}}{\sigma_\infty({l'_0}^{-\beta})} . \underset{(q_0 h_0)^{-1} = h_0^{-1} q_0^{-1}}{\underbrace{\sigma_\infty(q_0)^{-1}}} = \hat{\varphi}(q_0) {l'_0}^{-\beta} q_0^{-1}
$$
\marginpar{NB. si $q_0$ est symétrique, i.e. commute à l'involution [unclear: hyperelliptique], on a $\sigma_\infty q_0 \sigma_\infty = q_0^{-1}$, d'où $h_0 = q_0^2$. Si de plus $\beta = 1$ (cf. [unclear: bas] de la page) et $q_0 = \dots$}

$$
\sigma_\infty(\hat{\varphi}(q_0)) \varepsilon_1^{-\beta} h_0^{-1} q_0^{-1} = \hat{\varphi}(q_0) \varepsilon_0^{-\beta} q_0^{-1}
$$

$$
\varepsilon_1^{-\beta} h_0 \varepsilon_0^\beta = \sigma_\infty( (\hat{\varphi}(q_0))^{-1} ) \hat{\varphi}(q_0) \leqno{(24)}
$$
\marginpar{NB. Cette équation se déduit de (24 bis) [unclear: \dots]}

où $h_0$ a [unclear: l'avantage] [unclear: d'être] \dots
$$
\begin{cases}
\varepsilon_0 = l'_0 = \sigma_0 \rho \quad (\varepsilon_0^2 = l_0) \\
\varepsilon_1 = \rho \sigma_\infty = \sigma_1 \rho \quad (\varepsilon_1^2 = l_1)
\end{cases}
$$

Posons, pour $\mu \in \widehat{\mathfrak{S}}_{0,3}^+$ ,
$$
\tilde{\mu} = \sigma_\infty \mu^{-1} \sigma_\infty \mu = [\sigma_\infty, \mu] \leqno{(25)}
$$

(de sorte que $h_0 = \tilde{q}_0$) la formule (24) prend la forme\footnote{Il faudrait pouvoir [unclear: afficher] \dots}
$$
\varepsilon_1^{-\beta} \tilde{q}_0 \varepsilon_0^\beta = (\hat{\varphi}(q_0))^\sim \leqno{(27)}
$$

\newpage

%% ===== Page 63 =====

En résumé :

Proposition. Considérons un autom. ext. $u$
du groupe à lacets $\hat{\Pi}_{0,3}$, réalisé par
[unclear: via $l_0$] $g_0 \in \hat{\mathfrak{T}}_{0,3}$ (et $\beta \in \hat{\mathbf{Z}}^*$, et $\gamma \in \hat{\mathbf{Z}}^*$ uniquement
déterminé (en fonction de $g_0$ et $\beta$) par les
formules (11), au moyen d'un autom. effectif
(noté encore $u$)
qui commute à $\sigma_\infty$ — les données
$g_0, \beta$ (et $\gamma$) étant assujetties à vérifier les
conditions (9) et (10) ci-dessus, plus la
condition d'engendrement (11 bis). [unclear: (On prolongera un autom. de $\mathfrak{S}_{0,3}^{+\wedge}$ cf (15)).] Pour que
l'hom. extérieur $\hat{\varphi} : \hat{\mathfrak{T}}_{0,3} \to \mathfrak{S}_{0,3}^{+\wedge}$
défini par (18) commute à l'action
ext. de $u$, il faut et il suffit que l'on
puisse trouver un $\beta \in \hat{\mathbf{Z}}^*$, tel qu'on
ait (24)\footnote{Mieux : (24 bis).}, ou ce qui revient au même,
(27).

Cette condition (pour $\beta$ fixé) équivaut à la formule
(20) pour $x=l_\infty$ ($\lambda$ étant toujours donné par (22)),
c.à.d. à la formule
$$
u(l_\infty') = \operatorname{int}(\lambda^{-1} \hat{\varphi}(g_\infty))(l_\infty'^\beta) \leqno{(28)}
$$
où
$$
g_\infty = g_0 l_0^{-\beta} \hat{\varphi}(h_0) \quad \text{(cf. [unclear: la définition] (22))} \leqno{(29)}
$$

\newpage

%% ===== Page 64 =====

formule qui figure dans la formule présomptive (18). Elle s'écrit aussi, plus explicitement
$$
\rho^\gamma \rho^{-1} = \operatorname{int} \left( g_0 \varepsilon_0^\beta \hat{\varphi} \left( l_0^{-\gamma} p^{-1}(h_0) \right) \right) (\rho^{-\gamma} \rho)
$$
i.e. (en prenant les inverses et explicitant $g_0$)
\marginpar{NB On pose $\alpha = \beta + \gamma$. L'existence de $\beta \in \hat{\mathbf{Z}}$ vérifiant (24), équivaut à celle de $\alpha \in \hat{\mathbf{Z}}$ vérifiant (30).}
$$
\rho \rho^{-\gamma} = \operatorname{int} \left( \varepsilon_0^{\beta+\gamma} \hat{\varphi}(p^{-1}(h_0)) \right) (\rho^\gamma \rho^{-1}) \leqno{(30)}
$$

L'équivalence de (24) et (30) peut paraître un peu étrange, du fait qu'en (24), l'exposant $p$ n'intervient pas explicitement, mais seulement $h_0$, tandis que la formule (30) fait intervenir la valeur de $p \in \hat{\mathbf{Z}}^*$ mod $2$ (qui donne $\rho^p = \rho$ ou $\rho^p = \rho^{-1}$). Mais en tout état de cause, qu'on en soit sur la forme (24) ou sous la forme (30), il semble bien que la condition en question soit hautement tautologique !

Pour bien faire, il faudrait s'efforcer d'expliciter les conditions supplémentaires du § 39 en fonction de $g_0, p$.

\newpage

%% ===== Page 65 =====

Je voudrais comparer les autom. ext. de $\hat{\Pi}_{0,3}$, et ceux de $\mathfrak{S}^{+\wedge}_{0,3}$, tenant compte de la suite exacte
$$
1 \to \hat{\Pi}_{0,3} \to \mathfrak{S}^{+\wedge}_{0,3} \to \mathfrak{S}_3 \to 1 \leqno{(31)}
$$

Notons que tout autom. ext. $u$ de $\mathfrak{S}^{+\wedge}_{0,3}$ qui invarie le sous-groupe distingué $\hat{\Pi}_{0,3}$, et induit dans le quotient $\mathfrak{S}_3$ un autom. extérieur trivial, peut s'écrire comme déduit d'un autom. $u_0$ de $\mathfrak{S}^{+\wedge}_{0,3}$, stabilisant $\hat{\Pi}_{0,3}$, et induisant l'identité sur le quotient $\mathfrak{S}_3$ -- et $u_0$ est alors déterminé modulo composition avec $\operatorname{int}(g)$, où $g \in \mathfrak{S}^{+\wedge}_{0,3}$ -- tel que son image dans $\mathfrak{S}_3$ soit central, i.e. $= 1$, i.e. $g \in \hat{\Pi}_{0,3}$. Donc

$$
\begin{array}{lcl}
\begin{array}{l}
\text{Groupe des automorphismes} \\
\text{ext. de } \mathfrak{S}^{+\wedge}_{0,3} \text{ qui} \\
\text{induisent dans } \mathfrak{S}_3 \text{ un autom.} \\
\text{ext. trivial (donc invarient} \\
\hat{\Pi}_{0,3})
\end{array}
& \simeq &
\begin{array}{l}
\text{gr. de tous les} \\
\text{autom. de} \\
\mathfrak{S}^{+\wedge}_{0,3} \text{ qui} \\
\text{induisent} \\
\text{l'identité dans} \\
\text{le quotient } \mathfrak{S}_3 \\
\text{mod le sous-} \\
\text{groupe des} \\
\text{autom. int.} \\
\text{par des } g \in \hat{\Pi}_{0,3}
\end{array}
\end{array}
\leqno{(32)}
$$

\newpage

%% ===== Page 66 =====

Un autom. $u_0$ est-il connu par sa restrict. à $\hat{\pi}_{0,3}$ ? En d'autres termes, un autom. de $\widehat{\mathfrak{S}}_{0,3}^+$ (qui $\simeq \operatorname{Sl}(2, \hat{\mathbf{Z}})^{\wedge}$) qui induit l'identité sur $\hat{\pi}_{0,3}$ et sur $\mathfrak{S}_3$, est-il trivial ? Il suffit de prouver que sa restriction à l'image inverse de $\{ \sigma_\infty \} \subset \mathfrak{S}_3$ est trivial (car par symétrie on aura la [unclear: même chose] en remplaçant $\sigma_\infty$ par $\sigma_0, \sigma_1$, or ces trois sous-groupes engendrent $\widehat{\mathfrak{S}}_{0,3}^+$). On a 
$$
f(\sigma_\infty) = \sigma_\infty g, \quad \text{avec } g \in \hat{\pi}_{0,3}
$$ 
[et on doit avoir $\sigma_\infty g \sigma_\infty g = 1$ i.e. $g . \sigma_\infty(g) = 1$] de plus on doit avoir que $\operatorname{int}(\sigma_\infty g) | \hat{\pi}_{0,3} = \operatorname{int}(\sigma_\infty) | \hat{\pi}_{0,3}$, d'où $\operatorname{int}(g) | \hat{\pi}_{0,3} = \mathrm{id}$ i.e. $g \in \operatorname{Centr}(\hat{\pi}_{0,3})$ i.e. $g=1$, ok. 

Mais : quand un autom. de $\hat{\pi}_{0,3}$ peut-il se prolonger en un autom. de $\widehat{\mathfrak{S}}_{0,3}^+$ qui induise l'identité sur $\mathfrak{S}_3$ ? Bien sûr on voit aussitôt qu'il faut pour cela que l'autom. ext. de $\hat{\pi}_{0,3}$ qu'il induit commute à l'action de $\mathfrak{S}_3$.

\newpage

%% ===== Page 67 =====

C'est suffisant, soit $\Gamma \subset \Autext(\hat{\pi}_{0,3})$ le groupe des
automorphismes extérieurs de $\hat{\pi}_{0,3}$ qui a cette propriété (on ne
se préoccupe pas pour le moment de la
structure : [unclear: l'auto] de $\hat{\pi}_{0,3}$), on a
$\Gamma \cap \mathfrak{S}_3 = (1)$ (en identifiant $\mathfrak{S}_3$ à un
sous-groupe de $\Autext(\hat{\pi}_{0,3})$ car le centre
de $\mathfrak{S}_3$ est trivial, donc \marginpar{ceci est inexact}
$$ \mathfrak{S}_3 \times \Gamma \hookrightarrow \Autext(\hat{\pi}_{0,3}) \ , $$
et on a donc une structure d'extension
$$ \hat{\pi}_{0,3} \to \mathcal{E} \to \mathfrak{S}_3 \times \Gamma \to 1 \leqno{(38)} $$
(car le centre de $\hat{\pi}_{0,3}$ est trivial), on a [unclear: alors]
une structure d'extension
$$ 1 \to \hat{\mathfrak{S}}_{0,3}^+ \to \mathcal{E} \to \Gamma \to 1 \leqno{(39)} $$
qui définit une action extérieure de $\Gamma$
sur $\hat{\mathfrak{S}}_{0,3}^+$, qui est l'action naturelle
cherchée. Ainsi

Prop. Le groupe des autom. ext. de $\hat{\pi}_{0,3}$ qui \marginpar{Mais tout autom. de $\mathfrak{S}_3$ est intérieur.}
[unclear: commutent à] $\mathfrak{S}_3$ est isom. can. au groupe
des autom. ext. de $\hat{\mathfrak{S}}_{0,3}^+$ qui induisent
sur $\mathfrak{S}_3$ l'autom. ext. trivial.

\newpage

%% ===== Page 68 =====

Corollaire : Le groupe des autom. de $\hat{\Pi}_{0,3}$ qui
commutent aux extérieurs dans $\mathfrak{S}_3$, est isom.
au gr. des autom. de $\hat{\mathfrak{S}}_{0,3}^+$ qui
induisent sur $\mathfrak{S}_3$ l'autom. trivial (par
restriction de ces derniers).

(35)
\begin{tikzcd}
1 \ar[r] & \hat{\Pi}_{0,3} \ar[r] \ar[d, equal] & \operatorname{Aut}(\hat{\Pi}_{0,3} | \substack{ \text{comm.} \\ \text{à ext.} \\ \mathfrak{S}_3 }) \ar[r] \ar[d, "\simeq \text{ resp. } \hat{\Pi}_{0,3}"] & Aut ext(\hat{\Pi}_{0,3} | = \text{ dans } \mathfrak{S}_3) \ar[r] \ar[d, "\simeq"] & 1 \\
1 \ar[r] & \hat{\Pi}_{0,3} \ar[r] & \underbrace{\operatorname{Aut}(\hat{\mathfrak{S}}_{0,3}^+ | \substack{ \text{id. sur le} \\ \text{quot. } \mathfrak{S}_3 })}_{\mathcal{E}'} \ar[r] & Aut ext(\hat{\mathfrak{S}}_{0,3}^+ | \substack{ \text{triv.} \\ \text{sur quot.} \\ \mathfrak{S}_3 }) \ar[r] & 1
\end{tikzcd}

Cette extension $\mathcal{E}'$ de $\Gamma$ par $\hat{\Pi}_{0,3}$ est
en fait scindée, les [unclear] opèrent [unclear]
de $\hat{\Pi}_{0,3}$ qui commutent aux
[unclear] de $\hat{\mathfrak{S}}_{0,3}^+$ qui [unclear]
(par ex. : $\sigma_{\infty L}$ et [unclear] $\hat{\Pi}_{0,3}$)

(36)
$$ \begin{cases} u(\sigma_{\infty L}) = \sigma_{\infty L} \\ u(p) = g \cdot p \end{cases} \qquad (g \in \hat{\Pi}_{0,3}) $$

[MARGIN: provient au fait qu'il y a exact. un elt d'ordre 2 au dessus de $\sigma_\infty$...]

\newpage

%% ===== Page 69 =====

323

ici on a que $s$ est un homom. Car un "splitting" pour $\hat{\Pi}_{0,3}$ est défini par les gén. $\sigma_\infty, \rho$ avec les relations
$$
\sigma_\infty^2 = \rho^3 = 1 \leqno (37)
$$
\marginpar{[unclear: de plus il faut pour que $\sigma_\infty$ et $g\rho$ ($\in \hat{\Pi}_{0,3}$) engendrent $\hat{\Pi}_{0,3}$, val. $\Rightarrow (g\rho)^3=1$]}%
et on trouve la condition
$$
(g\rho)^3 = 1 \quad \text{i.e.} \quad g \rho(g) \rho^2(g) = 1 \leqno (38)
$$

Proposition. Soit $u \in \mathcal{E}^!$ un élément qui induit l'identité sur $\mathfrak{S}_3$. Conditions équivalentes :

\begin{enumerate}
\item[a)] $\exists p \in \hat{\mathbf{Z}}^*$ tel que $u(l_0)$ soit conjugué dans $\hat{\mathfrak{S}}_{0,3}^+$ à $l_0^p$ (ce $p$ est alors unique, il est en fait conjugué par un élém. de $\hat{\Pi}_{0,3}$)
\item[b)] $\exists p \in \hat{\mathbf{Z}}^*$ tel que $u(l_0)$ soit conjugué dans $\hat{\Pi}_{0,3}$ à $l_0^p$ (ce $p$ est alors unique)
\item[c)] $u_0 = u | \hat{\Pi}_{0,3}$ respecte la structure d'[unclear: inertie] de $\hat{\Pi}_{0,3}$.
\end{enumerate}

De plus, le $p$ dans a) et b) est le même.

Dém. Si $u(l_0) = \operatorname{int}(g_0)(l_0^p)$, avec $g_0 \in \hat{\mathfrak{S}}_{0,3}^+$, [unclear: comme $u$ induit] on voit que l'image de $g_0$ dans $\mathfrak{S}_3$ est $1$ ou $\sigma_0$, donc si on a (a $\Rightarrow$ b), si on a $g_0 = h \sigma_0$ ($h \in \hat{\Pi}_{0,3}$), on a 
$$u(l_0) = \operatorname{int}(h)(\sigma_0 l_0^p \sigma_0^{-1})$$
or on a $\sigma_0 l_0 \sigma_0^{-1} = l_0^{-1}$ (car $\operatorname{int}(\sigma_0)(l_0) = l_0^{-1}$)
donc $\sigma_0 l_0^p \sigma_0^{-1} = l_0^{-p}$, d'où $u(l_0) = \operatorname{int}(h)(l_0^{-p})$, d'où b) avec $-p$.

\newpage

%% ===== Page 70 =====

Comme on sait l'unicité de $p$ dans b),
[unclear]

(39) $$u(l_0) = int(g_0)(l_0^p) \quad (g_0 \in \hat{\Pi}_{0,3})$$
on déduit, écrivant
[crossed out: $\rho u_0 \rho^{-1}(l_0)$] $\rho u_0 \rho^{-1} = int(a) u_0 \quad (a \in \hat{\Pi}_{0,3})$

$$u_0(l_1) = int(a) \underbrace{\rho \underbrace{u_0 \underbrace{\rho^{-1} (l_1)}_{l_0}}_{int(g_0)(l_0^p)}}_{int(\rho(g_0))(l_1^p)} = int(a \rho(g_0)) (l_1^p)$$
i.e.
(40) $$u_0(l_1) = int(g_1)(l_1^p) \quad \dots \quad g_1 = a \rho(g_0)$$
et de même
$$u_0(l_\infty) = int(g_\infty)(l_\infty^p) \quad \dots \quad g_\infty = a \rho(g_1)$$
(41) $$= a \rho(a \rho(g_0))$$

donc b) $\implies$ c). Enfin prouvons c) $\implies$ a).

On a par hypothèse
$$u(l_0) = int(g) l_0^p$$
$$u(\varepsilon_0)^2 = int(g)((\varepsilon_0^p)^2) = (int(g) \varepsilon_0^p)^2$$

Je dis qu'on a alors
$$u(\varepsilon_0) = int(g) (\varepsilon_0^p)$$

Il suffit de [MARGIN: vérifier] sur le s-g de [crossed out: $\hat{\Gamma}_{0,3}^+$] [unclear] $l_0$ où $\varepsilon_0$ \dots OK

\newpage

%% ===== Page 71 =====

Désormais nous nous bornerons aux autom.
(eff. [unclear: et] extérieurs) de $\hat{\pi}_{0,3}$ qui respectent la str.
de lacets, et commutent à $\mathfrak{S}_3$ (extérieurement)
(eff. ou ext.),
- qui correspondent aux autom. (de
$\hat{\mathfrak{S}}_{0,3}^+$) qui induisent l'identité sur $\mathfrak{S}_3$,
et satisfont a) ci-dessus. Ils sont
Les autom. extérieurs de cette nature
correspondent à des autom. effectifs $u$
de $\hat{\pi}_{0,3}$ qui
fixent $\sigma_\infty$, et sont déterminés par
$g \in \hat{\pi}_{0,3}$ satisfaisant à (38), et tel que
de plus il existe $g_0 \in \hat{\pi}_{0,3}$ satisfaisant
$$ u(l_1') = \operatorname{int}(g_0)(l_1'^p) \quad \text{i.e.} \quad u(l_0) = \operatorname{int}(g_0)(l_0^p) \leqno{(39)} $$
(où $g_0$ est déterminé mod. $L_{00} = l_0^{\hat{\mathbf{Z}}}$ - ). Des
calculs faits par ailleurs montrent qu'on
doit avoir
$$ g = \sigma_\infty(g_0) l_1^r \rho(g_0)^{-1} \leqno{(40)} $$
avec $r \in \hat{\mathbf{Z}}$ bien déterminé (indépendant du choix de $g_0$)
i.e.
$$ \sigma_\infty(g_0)^{-1} g \rho(g_0) \in L_1 = \operatorname{Centr}_{\hat{\pi}_{0,3}}(l_1) \leqno{(41)} $$

\newpage

%% ===== Page 72 =====

En effet, on a par hyp. sur $g_0$ (cf. (39))
$$ u(L_0) = \operatorname{int}(g_0)(L_0) \leqno{(42)} $$
d'où
$$ u \sigma_\infty (L_0) = u(L_1) = \operatorname{int}(g)(L_1) $$
$$ \parallel \text{ (car } u \sigma_\infty = \sigma_\infty u \text{)} $$
$$ \sigma_\infty u(L_0) = \sigma_\infty(\operatorname{int}(g_0)(L_0)) = \operatorname{int}(\sigma_\infty(g_0)) (\sigma_\infty L_0) = \operatorname{int}(\sigma_\infty(g_0))(L_1) $$
d'où
$$ \operatorname{int}(\sigma_\infty(g_0)^{-1} g) (L_1) = L_1 $$
donc
$$ \sigma_\infty(g_0)^{-1} g p(g_0) \in \operatorname{Norm}_{\hat{\pi}_{0,3}} (L_1) = L_1 $$
d'où
$$ \sigma_\infty(g_0)^{-1} g p(g_0) = l_1^{\gamma} \quad (\gamma \in \hat{\mathbf{Z}}) $$
d'où (40). On notera, pour [unclear: s'en assurer], que l'élément $\gamma$ [unclear: et donc] $g$ défini au lemme est indépendant en effet du choix de $g_0 \pmod{L_0 \cdot \operatorname{Centr}_{\dots}}$ satisfaisant (42).

Ainsi, la donnée de $u$ [unclear: automorphisme] de $\hat{\pi}_{0,3}$ extériorisant [unclear: un élément] de $\hat{\mathfrak{S}}_3$, ou encore d'un [unclear: élément] de $\hat{\mathfrak{S}}_{0,3}^+$ induisant l'identité sur [unclear: \dots] et transformant $L_0$ en un conjugué $L_0' = z_0 \dots$, équivaut à la donnée d'un automorphisme $u_0$ de $\hat{\pi}_{0,3}$ induisant l'identité sur $\sigma_\infty$ et d'un [unclear: \dots], équivaut à la donnée de $g_0 \in \hat{\pi}_{0,3} \pmod{L_0}$, et de $\gamma \in \hat{\mathbf{Z}}$, satisfaisant la relation

\newpage

%% ===== Page 73 =====

(38) $g \rho g \rho g^2 \rho = 1$, où $g$ est donné par (40), et où
$u$ s'exprime en termes de $g$ de $g_0, \gamma$ par
(36)
$$
\begin{cases}
u(\sigma_{\infty}) = \sigma_{\infty} \\
u(g_1) = g\rho \qquad (= \sigma_{\infty}(g_0) l_1^\gamma \rho (g_0)^{-1} \rho)
\end{cases}
$$

la relation (42), s'écrit alors, [unclear] $L_0$ [unclear]
[unclear: $l_0 = l_0^{2\rho} = \dots$]
$(u(\sigma_{\infty}\rho))^3 \in int(g_0)(L_0')$
$\sigma_{\infty} g\rho = g_0 \sigma_{\infty} l_1^\gamma \rho g_0^{-1} = int(g_0) (\sigma_{\infty} l_1^\gamma \rho)$

donc (42) s'écrit
$(\sigma_{\infty} l_1^\gamma \rho)^2 \in L_0$

NB [unclear]
$(int(g_0)^{-1})(L_0) = (\sigma_{\infty} l_1^\gamma \rho)^2 = l_0^{2\gamma+1}$

soit $l_0^{2\gamma+1} \in L_0$

de [unclear] $u(L_0) =$
la relation (42), - la relation [unclear] plus forte
(43) $u(L_0') \subset int(g_0)(L_0')$

s'écrit
$int(g_0^{-1}) u(\sigma_{\infty}\rho) \in L_0'$
$\sigma_{\infty} g\rho = g_0 \sigma_{\infty} l_1^\gamma \rho g_0^{-1}$
$\sigma_{\infty} l_1^\gamma \rho = \sigma_{\infty} l_1^\gamma \sigma_{\infty} \sigma_{\infty} \rho$
$= \sigma_{\infty}(l_1)^\gamma \sigma_{\infty} \rho = l_0^{2\gamma} \varepsilon_0$

NB on a
$int(g_0^{-1}) u(l_0') = \varepsilon_0^{2\gamma+1}$

(44) $u(l_0') = u(\varepsilon_0) = int(g_0) (\varepsilon_0^{2\gamma+1})$

\newpage

%% ===== Page 74 =====

$$
u(l'_0) = int(g_0)(l_0^{' \rho}) \qquad \text{avec } p = 2\gamma + 1 \leqno{(44)}
$$
de même (39), donne [unclear: comme on s'en assure] :
$$
u(l_0) = int(g_0)(l_0^{\rho})
$$

\marginpar{NB} On notera que si $\rho \in \hat{\mathbf{Z}}^*$, alors $\rho - 1 \in 2\hat{\mathbf{Z}}$
i.e. $\frac{\rho-1}{2} \in \hat{\mathbf{Z}}$, car il suffit de vérifier la
relation analogue pour un $\rho \in \mathbf{Z}_\ell^*$, pour $\ell \neq 2$ on a $2\mathbf{Z}_\ell = \mathbf{Z}_\ell$,
dans le cas $\ell = 2$ $\rho \in \mathbf{Z}_2^*$ donne que l'image de
$\rho$ dans $\mathbf{Z}/4\mathbf{Z}$ est $1$ ou $3$ donc
$\rho - 1 \equiv 0 \text{ ou } 2 \mod 4$ et donc $\frac{\rho-1}{2} \in \mathbf{Z}_2$, OK.

Il reste donc à exiger la condition
(38), et on l'écrit sous la forme
$$
\sigma_\infty(q_0) l_1^r p(q_0^{-1}) p \sigma_\infty(q_0^{-1}) l_0^\gamma p(q_0) p \sigma_\infty(q_0) l_0^\gamma q_0^{-1} = 1
$$
soit, en appliquant $int(q_0^{-1})$ et posant
(cf (8))
$$
q_0^{-1} \sigma_\infty(q_0) = h_0^{-1} \leqno{(45)}
$$
i.e. $h_0 = \sigma_\infty(q_0^{-1}) q_0 = [\sigma_\infty^{-1}, q_0]$

$$
h_0^{-1} l_1^r p(h_0 l_1^r) p^2(h_0^{-1}) l_0^\gamma = 1
$$
[unclear: ou encore, par $int(l_0^\gamma)$]
$$
(h_0^{-1} l_1^r) \cdot p( l_1^r h_0^{-1} ) \cdot p^2( l_0^\gamma h_0^{-1} ) = 1 \qquad (\text{où } \gamma = \frac{\rho-1}{2}) \leqno{(46)}
$$

ce qui n'est autre que (10). (Quant à la
condition 17), on a vu plus haut qu'elle est automatiquement
vérifiée).

\newpage

%% ===== Page 75 =====

32bis

\textbf{Théorème} Un autom. ext. de $\hat{\Pi}_{0,3}$, commuant à
$\mathfrak{S}_3$, se relève de façon unique en un autom.
du groupe à lacet $\hat{\Pi}_{0,3}$, commuant à $\sigma_\varpi$ (et comm-
-utant à $\rho$ mod. autom. intérieurs), ou encore
comme restriction d'un autom. $u$ de $\hat{\mathfrak{S}}_{0,3}^+$
qui fixe $\sigma_\varpi$, laisse $\rho$ mod. $\hat{\Pi}_{0,3}$ (i.e. fixe $\bar{\rho}$ dans $\mathfrak{S}_3$),
et qui laisse la classe de conjugaison du \marginpar{fermé} s-groupe
$L'_0 = \langle\langle l'_0 \rangle\rangle$ engendré par $l'_0 = \varepsilon_0 = \sigma_\varpi \rho$. (ou encore
du sous-groupe $L_0 = \langle\langle l_0 \rangle\rangle$ engendré par $l_0$). La
donnée d'un tel $u$ équivaut à celle de son
multiplicateur $p \in \hat{\mathbf{Z}}^*$, et d'un $g_0 \in \hat{\Pi}_{0,3}/L_0$
satisfaisant les conditions suivantes :

\begin{enumerate}
\item[a)] Posant $h_0 = \sigma_\varpi(g_0)^{-1} g_0 = [\sigma_\varpi, g_0^{-1}]$ ($\in \hat{\Pi}_{0,3}$)
$$ \gamma = \frac{p-1}{2} \in \hat{\mathbf{Z}} \leqno (47) $$
on a la relation (46) ci-dessus (invariante
par remplacement de $g_0$ par $g_0 l_0^\lambda$, $\lambda \in \hat{\mathbf{Z}}$)
$$ (l_0^\gamma h_0^\gamma) \rho(l_0^{p\gamma} h_0^{-\gamma}) \rho^2(l_0^{-\gamma} h_0^{-\gamma}) = 1 \leqno (48) $$

\item[b)] $\sigma_\varpi$ et $g\rho$ engendrent top. $\hat{\mathfrak{S}}_{0,3}^+$, où $g$
est donné en fonction de $g_0$ par la
$$ g = \sigma_\varpi(g_0) l_1^\gamma \rho(g_0)^{-1} = \sigma_\varpi(g_0 l_0^\gamma) \rho(g_0 l_0^\gamma)^{-1} \leqno [(40)] $$
\end{enumerate}

(NB la relation (46) équivaut à (38) $g \rho(g) \rho^2(g) = 1$,
exprimée directement en termes de $g_0$, $\gamma = \frac{p-1}{2}$),
vu les relations génératrices $\sigma_\varpi^2 = (g\rho)^3 = 1$ i.e. :
Et l'automorphisme $u$ est donné en termes de $g_0$ et de $p$ par :

\newpage

%% ===== Page 76 =====

$$
\begin{cases}
u(\sigma_{\infty}) = \sigma_{\infty} \\
u(\rho) = g \rho
\end{cases}
\leqno{([36])}
$$
(où $g$ donné par $[(40)]$)

On aura alors
$$
\begin{matrix}
\begin{cases}
u(l_0) = \operatorname{int}(g_0) l_0^p \\
u(l_1) = \operatorname{int}(g_1) l_1^p = \operatorname{int}(\sigma_{\infty}(g_0))(l_1^p) \\
u(l_{\infty}) = \operatorname{int}(g_{\infty}) (l_{\infty}^p)
\end{cases}
&
\begin{aligned}
g_1 &= g \rho(g_0) = \sigma_{\infty}(g_0) l_1^\gamma \\
    &= \sigma_{\infty}(g_0 l_0^\gamma) \\
g_{\infty} &= g \rho(g_1) = \rho^2(g^{-1} g_0) \\
           &= \rho(g_0 l_0^\gamma) h_0
\end{aligned}
\end{matrix}
\leqno{(49)}
$$

\newpage

%% ===== Page 77 =====

Scholie \quad Dans le th. 327, le $g \in \Pi_{0,3}$\footnote{$\in \Pi_{0,3}$} n'intervient directe-
ment (dans $u$) s'exprime par l'intermédiaire
de $q_0$ et $\gamma = \frac{p-1}{2}$, \quad ($q_0$ défini par $L_0$)

- la condition (46) étant en fait une condition
qui s'exprime directement en termes de $q_0$,
mieux encore, en termes de
$$
\xi = l_0^\gamma h_0^{-1} \in \Pi_{0,3} \quad (d'où \quad h_0 = \xi^{-1} l_0^\gamma ) \leqno{(50)}
$$
comme
$$
\xi \rho(\xi) \rho^2(\xi) = 1 \quad , \leqno{(51)}
$$
qui équivaut à :
$$
(\xi \rho)^3 = 1 \quad . \leqno{(52)}
$$
Les solutions de l'équation (51), -- ou
ce qui revient au même de (52), correspondent aux
sous-groupes\footnote{dans $\mathfrak{S}_{0,3}^+$} isomorphes à $\mathbf{Z}/3\mathbf{Z}$ dans
$\hat{\mathfrak{T}}_{0,3}$, engendrés par $\rho \in \mathfrak{S}_3$, ou i.e. aux
éléments d'ordre 3 de $\mathfrak{S}_{0,3}^+$ au-dessus de $\rho$.

Sur l'ens. de ces sous-groupes resp. de ces
éléments, $\Pi_{0,3}$ opère par [unclear: autom. intérieurs].
$$
\operatorname{int}(\lambda) (\xi \rho) = \lambda \xi \rho \lambda^{-1} = \lambda \xi \rho \lambda^{-1} \rho^{-1} \rho = (\lambda \xi \rho(\lambda)^{-1} ) \rho \leqno{(53)}
$$
i.e. il opère sur l'ens. des solutions de (51)
par
\marginpar{NB $\lambda \mapsto \lambda \rho(\lambda)^{-1}$ \\ est injectif}
$$
\xi \mapsto \lambda \xi \rho(\lambda)^{-1} \leqno{(54)}
$$

\newpage

%% ===== Page 78 =====

Admettons (quitte à le vérifier par la suite)
qu'il y a exactement deux classes de conjugaison :
celle de $\{\rho\}$, correspondant au sous-groupe tautologique,
de sorte que pour de tels $\xi$ on a
$$
\xi = \lambda \rho(\lambda)^{-1} \quad (\lambda \in \hat{\Pi}_{0,3}) \leqno{(55)}
$$
et d'autre part la classe de conjugaison
des él.$^\text{ts}$ d'ordre 3 $\sigma_\infty \rho^{-1} \sigma_\infty^{-1} = \sigma_\infty \rho^2 \sigma_\infty$ de
$\hat{\mathfrak{S}}$ , correspondant à $\xi_0 \rho$ , avec
$$
\begin{cases}
\xi_0 &= \sigma_\infty \rho^{-1} \sigma_\infty \rho^{-1} = \sigma_\infty (\rho^{-1} \sigma_\infty) (\rho^{-1} \sigma_\infty) \sigma_\infty \\
&= \sigma_\infty \varepsilon_0^{-2} \sigma_\infty = \overset{\varepsilon_0^{-1}}{\sigma_\infty} (l_0^{-1}) = l_1^{-1}
\end{cases}
$$
(et on a bien
$$
l_1^{-1} \rho(l_1^{-1}) \rho^2(l_1^{-1}) = l_1^{-1} l_\infty^{-1} l_0^{-1} = (l_0 l_\infty l_1)^{-1} = 1 \ )
$$
Cette orbite correspond aux solutions de
la forme 
$$
\xi = \lambda l_1^{-1} \rho(\lambda)^{-1} \quad (\lambda \in \hat{\Pi}_{0,3}) \leqno{(56)}
$$
[On notera que $l_0^{-1}$ et $l_1^{-1}$ sont sur la m.
orbite
$$
l_1^{-1} = l_0 (l_0^{-1}) \rho(l_0)^{-1} ]
$$
i.e. avec $\lambda = l_0$.

A) Cas où $\xi$ est de la forme (55), donc
$$
h_0 = \rho(\lambda) \lambda^{-1} l_0^\gamma \leqno{(57)}
$$

\newpage

%% ===== Page 79 =====

Il reste donc à résoudre l'équation en $g_0$ (47)
$$
\sigma_\infty(g_0)^{-1} g_0 = \underbrace{\rho(\lambda)\lambda^{-1} l_0^\gamma}_{h_0}
\leqno{(58)}
$$
Pour $\gamma$ et $h_0$ fixés, il y a au plus un seul $g_0$\footnote{Résultant de ce que $\hat{\Pi}_{0,3}^{\sigma_\infty} = \{1\}$.} qui satisfait à cette condition, et un tel $g_0$ existe ssi $h_0 \sigma_\infty(h_0) = 1$, ce qui s'écrit,
$$
\sigma_\infty(\rho(\lambda)) \sigma_\infty(\lambda^{-1}) l_1^\gamma \rho(\lambda)\lambda^{-1} l_0^\gamma = 1
\leqno{(59)}
$$
qui s'écrit aussi, en reprenant l'expression $h_0 = \xi l_0^\gamma$
$$
\sigma_\infty(\xi) = l_1^{-\gamma} \xi^{-1} l_0^{-\gamma}
\leqno{(60)}
$$

En résumé

Corollaire. Les $u$ qui [unclear: se distinguent] par la classe de $\hat{\Pi}_{0,3}$-conjugaison de $\rho$ correspondent 1-1 aux systèmes $(\lambda, \rho) \in \hat{\Pi}_{0,3} \times \hat{\mathbf{Z}}^*$, tels que, en posant $\xi = \lambda \rho(\lambda)^{-1}$, (où $\gamma = \frac{\rho-1}{2}$), on ait la relation (60). On notera que pour $\lambda$ de $\hat{\Pi}_{0,3}$ fixé, il existe au plus un $\gamma \in \hat{\mathbf{Z}}$ tel que la relation (60) soit valable, donc $\rho$ est déterminé par $\lambda$ ; mieux : la condition suppose que $2\gamma+1$ de $\hat{\mathbf{Z}}$ soit dans $\hat{\mathbf{Z}}^*$ (i.e. $\exists \gamma' \in \hat{\mathbf{Z}}$ $2\gamma\gamma' + \gamma + \gamma' = 0$).

\newpage

%% ===== Page 80 =====

NB Si $\lambda$ est remplacé par $\lambda' = l_0^\alpha \lambda$,
$\zeta = \lambda \rho(\lambda)^{-1}$ est remplacé par $\zeta' = l_0^\alpha \zeta \rho(l_0)^{-\alpha}$,
et $h_0 = \zeta^{-1} l_0^\gamma$ est remplacé par
$h'_0 = \rho(l_0)^\alpha h_0 l_0^{-\alpha}$, enfin $g_0$ est remplacé
par $g'_0 = g_0 l_0^{-\alpha}$, OK. (L'équation (60) pour $\zeta$
et pour $\zeta'$ sont équivalentes)

Finalement, que se donne-t-on si on ne change
pas la classe de $\widehat{\Pi}_{0,3}$-conj. de $\rho$, la
démarche la plus claire pour déterminer
les $u$ en question me semble celle-ci :

1º) On choisit un $\zeta \in \widehat{\Pi}_{0,3}$ satisfaisant à l'équation
$$
\zeta \rho(\zeta) \rho^2(\zeta) = 1 \leqno (A)
$$
(on le présume soit sous la forme (55), soit
sous la forme (56) :
$$
\zeta = \lambda \rho(\lambda)^{-1} \quad \text{ou} \quad \zeta = \lambda l_0^{-1} \rho(\lambda)^{-1} \quad \text{, avec } \lambda \in \widehat{\Pi}_{0,3} ) \leqno (B)
$$

2º) On choisit $\gamma \in \widehat{\mathbf{Z}}$ tel que\marginpar{ici, $\zeta$ est bien entendu déterminé de $\lambda$ par les éq (B)}
$$
\sigma_\infty(\zeta) = l_1^{\gamma} \zeta^{-1} l_0^{\gamma} \leqno (C)
$$
ce qui détermine $\gamma$ de façon unique. On
[unclear: dégage] donc des $u$ en se bornant aux $\zeta$
pour lesquels un tel $\gamma$ existe, et tel de plus,
que $p = 2\gamma + 1$ soit $\in \widehat{\mathbf{Z}}^\ast$.

\newpage

%% ===== Page 81 =====

3°) On construit alors
$$
h_0 = \xi^{-1} l_0^\gamma \in \hat{\Pi}_{0,3}
$$
La relation de 2°) assure que $h_0 \sigma_\infty(h_0) = 1$ i.e.\ que $h_0$ peut se mettre sous la forme
$$
h_0 = \sigma_\infty(g_0)^{-1} g_0
$$
$g_0 \in \hat{\Pi}_{0,3}$.\footnote{NB Il faudrait s'assurer cependant que la condition sur $h_0$ détermine uniquement $g_0$ mod.\ le s-g engendré par $l_0$, i.e.\ que le commutant de $\sigma_\infty$ dans $\hat{\Pi}_{0,3}$ soit [unclear: engendré par] $l_0$ ? À vérifier !}

4°) On peut maintenant poser
$$
g = \sigma_\infty(g_0) l_1^\gamma \rho(g_0)^{-1} = (g_0 h_0^{-1}) \xi \rho(g_0 l_0^{-\gamma})^{-1}
$$
définissant ce qu'il faut pour définir
$$
\begin{cases} 
u(\sigma_\infty) = \sigma_\infty \\ 
u(\rho) = g \rho 
\end{cases}
$$

Ici, la seule équation délicate, pour laquelle l'existence (non l'unicité, on l'a vu) d'une solution, est problématique, est l'équation (C) (sur $\gamma \in \hat{\mathbf{Z}}$) --- plus la condition annexe que $2\gamma + 1 \in \hat{\mathbf{Z}}^\times$. C'est donc une condition sur $\lambda \in \hat{\Pi}_{0,3}$ --- ou plutôt deux types de conditions, suivant que $\xi$ est défini en termes de $\lambda$ par l'une ou l'autre formule (B).

Je prévois qu'on est dans l'un ou

\newpage

%% ===== Page 82 =====

l'autre un des formules (B), suivant
que $p \bmod 3$ est congru à $1$, ou à $-1$
(NB. il ne doit pas être congru à $0$,
puisque $p \bmod 3 \in (\mathbf{Z}/3\mathbf{Z})^*$) i.e. suivant
que
$$
2\gamma \equiv 0 \text{ ou } -2 \pmod{3} \leqno{(3)}
$$
i.e.
$$
\gamma \equiv 0 \text{ ou } -1 \pmod{3} \leqno{(3)}
$$
(NB. il ne faut pas
que $\gamma$ soit $\equiv 1 \pmod{3}$, sinon $2\gamma+1=p$
serait congru à $0 \pmod{3}$). C'est
en tous cas ce que donne le cas où
$u \in \Gamma_{\mathbf{Q}_0}$.

Pour le tester, réduisons le groupe $\hat{\mathfrak{S}}_{0,3}^+$
mod les équations $l_0^3=1$ i.e. $(\sigma_{\infty}\rho)^6=1$
de sorte que $\hat{\mathfrak{S}}_{0,3}^+$ est remplacé par une
extension de $\mathfrak{S}_3$ par le groupe
défini par les générateurs $l_0,l_1,l_{\infty}$
satisfaisant
$$
l_0^3 = l_1^3 = l_{\infty}^3 = (l_{\infty} l_0 l_0) = 1
$$
- ou mieux, utilisons
$$
\hat{\mathfrak{S}}_{0,3}^+ \simeq \operatorname{Sl}(2,\hat{\mathbf{Z}})'
$$
et réduisons au niveau de
$$
\operatorname{Sl}(2,\mathbf{Z}/6\mathbf{Z})' \simeq \operatorname{Sl}(2,\mathbf{Z}/2\mathbf{Z}) \times \operatorname{Sl}(2,\mathbf{Z}/3\mathbf{Z})'.
$$

\newpage

%% ===== Page 83 =====

Mais pour faire ces calculs, il faut pouvoir identifier $\sigma_\infty, \rho$ comme éléments de $\operatorname{Sl}(2, \hat{\mathbf{Z}})$ — explicitation que nous ferons tantôt.

Revenant à (24 bis), et y portant $h_0 = \xi l_0^\gamma$, il vient
$$
g_0 = \hat{\rho}(\xi) \varepsilon_0^{\gamma-\beta} \leqno{(61)}
$$

On trouve donc que $u$ correspond à $(\xi, g_0, \gamma, \beta)$, avec $\xi, g_0 \in \widehat{\Pi}_{0,3}$ et $\gamma, \beta \in \hat{\mathbf{Z}}$ satisfaisant
$$
\begin{cases} 
\xi \rho(\xi) \rho^2(\xi) = 1 \\ 
\sigma_\infty(g_0)^{-1} g_0 = l_0^\gamma \\ 
g_0 = \hat{\rho}(\xi) \varepsilon_0^{\gamma-\beta} \\ 
2\gamma+1 \in \hat{\mathbf{Z}}^* 
\end{cases} \leqno{(62)}
$$
d'où on peut éliminer $g_0$, et on trouve les équations en $(\xi, r, s)$ $(r = \gamma-\beta, s = \gamma+\beta)$
$$
\begin{cases} 
\xi \rho(\xi) \rho^2(\xi) = 1 \quad (\text{d'où } \xi = \lambda \rho(\lambda)^{-1}) \\ 
\sigma_\infty(\hat{\rho}(\xi))^{-1} \hat{\rho}(\xi) = \varepsilon_1^r \varepsilon_0^{-s} \\ 
r+s+1 \in \hat{\mathbf{Z}}^* 
\end{cases} \leqno{(63)}
$$

Ici, $\lambda$ est déterminé mod transformation $\lambda \mapsto l_0^\alpha \lambda$, qui donne $\xi' = l_0^\alpha \xi l_1^{-\alpha}$.

\newpage

%% ===== Page 84 =====

\section*{42. Récapitulation sur le groupe}
$(\mathfrak{S}_{0,3}\footnote{On a défini $\mathfrak{S}_{0,3}$ comme un [unclear: certain relèvement] (En fait il y a deux choix différents \dots) et $\mathfrak{S}_{0,3}$ \dots} \simeq \underline{\Pi}_1(U_{0,3}^{\mathrm{an}}, \mathfrak{S}_3 ; -\vec{j}) \simeq (\Pi_{0,3}' \simeq \pi_1((\mathbb{P}^{1\mathrm{an}} \setminus \{0\} + 2\{1\} + 3\{\infty\}), \{-\vec{j}\} ))$
$\simeq (\underline{\underline{\operatorname{Gl}(2, \mathbf{Z})'}} = \operatorname{Gl}(2, \mathbf{Z})/\pm 1 \simeq \operatorname{PGL}(2, \mathbf{Z})) \simeq \mathfrak{T}_{1,1}' = \mathfrak{T}_{1,1}^+ \ltimes B$

et sur les automorphismes "de lacets" de $\hat{\mathfrak{S}}_{0,3}$ et de $\hat{\Pi}_{0,3}$
(Formulaire fondamental).

I) Le groupe $\mathfrak{S}_{0,3}$ est présenté comme extension
$$
1 \to \underset{\Pi_1(U_{0,3} ; -\vec{j})}{\Pi_{0,3}} \to \mathfrak{S}_{0,3} \to \underset{\mathfrak{S}_3 \times \mathbf{Z}/2\mathbf{Z}}{\mathfrak{T}_{0,3}} \to 1 \leqno{(1)}
$$
où $\mathfrak{S}_3 \times \mathbf{Z}/2\mathbf{Z}$ est considéré comme groupe
opérant sur $U_{0,3} = \mathbb{P}^1_{\mathrm{an}} \setminus \{0, 1, \infty\} = \mathbb{C} \setminus \{0, 1\}$,
$\mathfrak{S}_3$ opérant par permutations --
c-à-d le groupe des automorphismes conformes de $\mathbb{P}^1$ qui stabilisent $\{0, 1, \infty\}$. Le groupe $\mathfrak{S}_3$ est le sous-groupe des homographies, il s'identifie au groupe des permutations de $\{0, 1, \infty\}$, ces six éléments sont
a) $\sigma_0, \sigma_1, \sigma_\infty$ (transpositions), $\rho, \rho^{-1}$ (permutations circulaires)
donnés par les formules explicites
$$
\begin{cases} \sigma_0(z) = \frac{z}{z-1}, \sigma_1(z) = \frac{1}{z}, \sigma_\infty(z) = 1-z \\ \rho(z) = \frac{1}{1-z}, \rho^{-1}(z) = 1 - \frac{1}{z} = \frac{z-1}{z} \end{cases} \leqno{(2)}
$$

$$
\begin{cases} \sigma_0 \sigma_1 = \sigma_1 \sigma_\infty = \sigma_\infty \sigma_0 = \rho \\ \rho \sigma_0 \rho^{-1} = \sigma_1, \rho \sigma_1 \rho^{-1} = \sigma_\infty, \rho \sigma_\infty \rho^{-1} = \sigma_0 \end{cases} \leqno{(3)}
$$

$$
\sigma_0^2 = \sigma_1^2 = \sigma_\infty^2 = \rho^3 = 1 \leqno{(4)}
$$

Le groupe $T_{0,3} \simeq \mathbf{Z}/2\mathbf{Z} \subset \mathfrak{T}_{0,3}$ apparaît comme le centre,

\newpage

%% ===== Page 85 =====

il est engendré par la conjugaison complexe $\tau$
(5) $$ \tau(z) = \bar{z} $$
les cinq autres éléments de $\mathcal{T}_{0,3}$ sont
(6) $$ \tau \sigma_0 = \tau'_0 \quad , \quad \tau \sigma_1 = \tau'_1 \quad , \quad \tau \sigma_\infty = \tau'_\infty \quad , \quad \tau \rho \quad , \quad \tau \rho^{-1} = (\tau \rho)^{-1} , $$
les anti-involutions $\tau'_i$ ($i = 0, 1, \infty$) en particulier --
sont des involutions anti-holomorphes, données par
(7) $$ \tau'_0(z) = \frac{\bar{z}}{\bar{z}-1} \quad , \quad \tau'_1(z) = \frac{1}{\bar{z}} \quad , \quad \tau'_\infty(z) = 1-\bar{z} . $$

Les points fixes des autom. continus de $(\mathbb{P}^{1 \text{an}} \setminus \{0,1,\infty\})$
sont les suivants

(8) $$
\left\{
\begin{aligned}
& (\mathbb{P}^1_\mathbb{C})^{\sigma_0} = \{0, 2\} \quad , \quad (\mathbb{P}^1_\mathbb{C})^{\sigma_1} = \{1, -1\} \quad , \quad (\mathbb{P}^1_\mathbb{C})^{\sigma_\infty} = \{\infty, 1/2\} \\
& (\mathbb{P}^1_\mathbb{C})^\rho = (\mathbb{P}^1_\mathbb{C})^{\rho^{-1}} = \{-j, -j^2\} \qquad \text{où } j = \exp 2i\pi/3 = \frac{-1+i\sqrt{3}}{2} \text{ racine primitive 3ème de 1} \\
& \qquad \qquad \qquad \qquad \qquad \qquad \quad -j = \exp 2i\pi/6 = \frac{1+i\sqrt{3}}{2} \text{ racine primitive 6ème de 1}
\end{aligned}
\right.
$$

[MARGIN: NB $(\mathbb{P}^1_\mathbb{C})^\sigma \cap (\mathbb{P}^1_\mathbb{C})^{\sigma'} = \emptyset$]

(9) $$
\left\{
\begin{aligned}
& (\mathbb{P}^1_\mathbb{C})^{\tau} = \mathbb{P}^1_\mathbb{R} \\
& (\mathbb{P}^1_\mathbb{C})^{\tau'_0} = \{ z \mid z + \bar{z} = z \bar{z} \} = \left\{ x+iy \mid (x-1)^2 + y^2 = 1 \right\} \\
& (\mathbb{P}^1_\mathbb{C})^{\tau'_1} = \{ z \mid z \bar{z} = 1 \} = \{ x+iy \mid x^2+y^2=1 \} = U \\
& (\mathbb{P}^1_\mathbb{C})^{\tau'_\infty} = \{ z \mid z + \bar{z} = 1 \} = \{ x+iy \mid x = 1/2 \} = 1/2 + i\mathbb{R}
\end{aligned}
\right.
$$

L'image la plus suggestive de tout cela est la
sphère de Riemann, où $\tau, \tau'_0, \tau'_1, \tau'_\infty$ apparaissent comme
des réflexions par r. aux "équateurs" correspondants, qui
sont 4 "lignes" (cercles ou droites)
passant toutes par les pôles $-j, -j^2$, qui
divisent le dit [unclear] de l'équateur en six
segments circulaires, par les pts de
subdivision $0, 1/2, 1, 2, \infty, -1$ .

Notons aussi les formules
(10) $$
\left\{
\begin{aligned}
& \rho \tau'_0 \rho^{-1} = \tau'_1 \quad , \quad \rho \tau'_1 \rho^{-1} = \tau'_\infty \quad , \quad \rho \tau'_\infty \rho^{-1} = \tau'_0 \\
& \tau'_0 \tau'_1 = \tau'_1 \tau'_\infty = \tau'_\infty \tau'_0 = \rho
\end{aligned}
\right.
$$

\newpage

%% ===== Page 86 =====

Les quatre grands cercles (9) découpent la sphère en 12 triangles sphériques, qui forment des domaines fondamentaux pour le groupe $\mathfrak{T}_{0,3} \simeq \mathfrak{S}_3 \times \mathbf{Z}/2$ d'ordre 12, et il les permute entre eux de façon simple-ment transitive. (Il s'agit de la sous-division barycentrique des doubles triangles sphériques déterminés par l'équateur $\mathbb{P}^1_{\mathbf{R}}$, et ses trois points de subdivision : $0, 1, \infty$). Prend un de ces triangles sphériques, disons $(-\bar{j}, 0, 1/2)$, les réflexions par rapport à ses trois côtés, i.e. $\tau, \tau'_0, \tau'_1$, forment un système de générateurs de $\mathfrak{T}_{0,3}$ , avec les relations

$$
\tau^2 = {\tau'_0}^2 = {\tau'_1}^2 = (\underset{\sigma'_0}{\tau \tau'_0})^2 = (\underset{\rho^{-1}}{\tau'_0 \tau'_1})^3 = (\underset{\sigma_0}{\tau'_1 \tau})^2 = 1 \leqno (11)
$$

quant à $\mathfrak{S}_3 = \mathfrak{T}_{0,3}^+$, il est engendré par $\tau \tau'_0 = \sigma'_0$, $\tau'_0 \tau'_1 = \rho^{-1}$ et $\tau'_1 \tau = \sigma_0$, avec les relations

$$
\sigma'_0 \rho^{-1} \sigma_0 = {\sigma'_0}^2 = \rho^3 = {\sigma_0}^2 = 1 \leqno (12)
$$

II) Le groupe $\Pi_{0,3}$ est engendré par $l_0, l_1, l_\infty$, liés par la relation génératrice

$$
l_\infty l_1 l_0 = 1 \leqno (13)
$$

Le sous-groupe d'indice 2 $\mathfrak{S}'_3 = (\mathfrak{T}_{0,3})_{-\bar{j}} \subset \mathfrak{T}_{0,3}$ , stabilisateur de $-\bar{j}$, formé des $(g, \tau^\varepsilon) \in \mathfrak{S}_3 \times \mathbf{Z}/2 \mathbf{Z}$ avec $\operatorname{sgn}(g) = (-1)^\varepsilon$ , opère sur $\Pi_{0,3}$. Ce sous-groupe est formé de $1, \tau'_0, \tau'_1, \tau'_\infty, \rho, \rho^{-1}$, et leurs opérations sur $\Pi_{0,3}$ sont notées respectivement $1, \tau'_0, \tau'_1, \tau'_\infty, \rho, \rho^{-1}$.\marginpar{NB $\tau'_0, \tau'_1, \tau'_\infty$ étant des [unclear: anti-automorphismes] (dans (C)) \dots opèrent sur $\Pi_{0,3}$ \dots}

\newpage

%% ===== Page 87 =====

On identifie ces éléments à des éléments du groupe $\operatorname{Aut}_{\text{loc}}(\Pi_{0,3}) \simeq \mathfrak{S}_{0,3}$
par $g(x) = g x g^{-1}$ ; et on trouve pour $g \in \mathfrak{S}'_J, x \in \Pi_{0,3}$, --
$$
\begin{array}{lll}
\rho(l_0) = \rho l_0 \rho^{-1} = l_1, & \rho(l_1) = \rho l_1 \rho^{-1} = l_\infty, & \rho(l_\infty) = \rho l_\infty \rho^{-1} = l_0 \\
\tau'_0(l_0) = l_0^{-1}, & \tau'_0(l_1) = l_\infty^{-1}, & \tau'_0(l_\infty) = l_1^{-1} \\
\tau'_1(l_1) = l_1^{-1}, & \tau'_1(l_0) = l_\infty^{-1}, & \tau'_1(l_\infty) = l_0^{-1} \\
\tau'_\infty(l_\infty) = l_\infty^{-1}, & \tau'_\infty(l_0) = l_1^{-1}, & \tau'_\infty(l_1) = l_0^{-1}
\end{array}
\leqno{(14)}
$$

$$
\tau'_0 \tau'_1 = \tau'_1 \tau'_\infty = \tau'_\infty \tau'_0 = \rho \qquad (\text{on a bien sûr } \tau'^2_0 = \tau'^2_1 = \tau'^2_\infty = \rho^3 = 1)
\leqno{(15)}
$$

Dans le groupe $\mathfrak{S}_{0,3}$ on définit des éléments $\sigma_0, \sigma_1, \sigma_\infty$ au-dessus de $\bar{\sigma}_0, \bar{\sigma}_1, \bar{\sigma}_\infty$, par
$$
\begin{array}{lll}
\sigma_0(l_1) = l_\infty, & \sigma_0(l_\infty) = l_1, & \sigma_0(l_0) = (l_1 l_\infty)^{-1} \\
\sigma_1(l_\infty) = l_0, & \sigma_1(l_0) = l_\infty, & \sigma_1(l_1) = (l_\infty l_0)^{-1} \\
\sigma_\infty(l_0) = l_1, & \sigma_\infty(l_1) = l_0, & \sigma_\infty(l_\infty) = (l_0 l_1)^{-1}
\end{array}
\leqno{(16)}
$$

On a
$$
\begin{array}{lll}
\rho \sigma_0 \rho^{-1} = \sigma_1, & \rho \sigma_1 \rho^{-1} = \sigma_\infty, & \rho \sigma_\infty \rho^{-1} = \sigma_0 \\
\rho \tau'_0 \rho^{-1} = \tau'_1, & \rho \tau'_1 \rho^{-1} = \tau'_\infty, & \rho \tau'_\infty \rho^{-1} = \tau'_0 \\
\sigma_0 \sigma_1 = \rho l_0 = l_1 \rho, & \sigma_1 \sigma_\infty = \rho l_1 = l_\infty \rho, & \sigma_\infty \sigma_0 = \rho l_\infty = l_0 \rho
\end{array}
\leqno{(17)}
$$

\marginpar{NB $\sigma_0, \sigma_1, \sigma_\infty \in \mathfrak{S}_{0,3}^+$ \\ $\tau'_0, \tau'_1, \tau'_\infty \in \mathfrak{S}_{0,3}$ \\ $\tau_0, \tau_1, \tau_\infty \in \mathfrak{S}_{0,3}$}
On a ainsi des éléments de $\mathfrak{S}_{0,3}$ au-dessus des éléments $1, \bar{\sigma}_0, \bar{\sigma}_1, \bar{\sigma}_\infty, \rho, \rho^{-1}, \bar{\tau}'_0, \bar{\tau}'_1, \bar{\tau}'_\infty$, mais pas encore au-dessus de $\bar{\tau} \in \bar{\mathfrak{S}}_{0,3}$ (qui permet d'en avoir un [unclear: ex $\tau_0, \tau_1^{-1}$]), on posera\marginpar{[unclear: rappelons que pour ... l'identité]}
$$ \tau_i = \sigma_i \tau'_i \quad (= \tau'_i \sigma_i) $$
(qui est sur $\bar{\tau}$, donc de carré 1), $\tau_i$ étant caractérisé par la condition d'être sur $\bar{\tau}$ et de commuter à $\sigma_i$ (ou encore, à $\tau'_i$).

$$
[\tau_0, \sigma_0] = [\tau_1, \sigma_1] = [\tau_\infty, \sigma_\infty] = 1
\qquad
\begin{array}{l}
\tau_0 = \sigma_0 \tau'_0 = \tau'_0 \sigma_0, \\
\tau_1 = \sigma_1 \tau'_1 = \tau'_1 \sigma_1, \\
\tau_\infty = \sigma_\infty \tau'_\infty = \tau'_\infty \sigma_\infty
\end{array}
\leqno{(19)}
$$
$$
\tau_0^2 = \tau_1^2 = \tau_\infty^2 = 1 \qquad \text{et} \qquad \tau_0, \tau_1, \tau_\infty \text{ sur } \bar{\tau}
\leqno{(20)}
$$
$$
\text{On trouve} \quad \tau_0 \tau_1 = l_\infty \ , \ \tau_0 \tau_\infty = l_1 \ , \ \tau_1 \tau_\infty = l_0
\leqno{(21)}
$$

Pour bien faire, il faudrait expliciter la description de Malgoire (faite dans [unclear: (C)]) de $\sigma_0, \sigma_1, \sigma_\infty, \rho, \tau'_0, \tau'_1, \tau'_\infty, \tau_0, \tau_1, \tau_\infty$.

\newpage

%% ===== Page 88 =====

On a
$$
\begin{cases}
(\sigma_\infty\rho)^2 = (\rho\sigma_1)^2 = l_0 = \tau_0^2, & \xi_0 \defeq \sigma_\infty \rho = \rho \sigma_1 \\
(\sigma_0\rho)^2 = (\rho\sigma_\infty)^2 = l_1 = \tau_1^2, & \xi_1 \defeq \sigma_0 \rho = \rho \sigma_\infty \\
(\sigma_1\rho)^2 = (\rho\sigma_0)^2 = l_\infty = \tau_\infty^2, & \xi_\infty \defeq \sigma_1 \rho = \rho \sigma_0
\end{cases}
\leqno{(22)}
$$

L'élément $\xi_i \in \mathfrak{S}_{0,3}$ (en fait $\xi_i \in \mathfrak{S}_{0,3}^+$) au dessus de $\dot{\sigma}_i$, il est caractérisé par la condition $\xi_i^2 = l_i$.

Dans le groupe $\bar{\mathfrak{T}}_{0,3}$ d'ordre 12, il y a sept éléments involutifs $\dot{\sigma}_0, \dot{\sigma}_1, \dot{\sigma}_\infty, \dot{\tau}'_0, \dot{\tau}'_1, \dot{\tau}'_\infty$, (dont un central), d'ordre 3 $\dot{\rho}, \dot{\rho}^{-1}$ (et $\dot{\rho} \varepsilon, \dot{\rho}^{-1} \varepsilon$ qui sont des éléments d'ordre 6).

[unclear: Nous allons préciser le choix de nos] relèvements en des éléments [unclear: d'ordre fini] de $\mathfrak{S}_{0,3}$. On sait que les possibilités, mod $\Pi_{0,3}$-conjugaison, de tels relèvements, correspondent aux composantes connexes des pts fixes dans $\tilde{U}_{0,3}$ des actions topologiques correspondantes. - pour $\dot{\sigma}_0, \dot{\sigma}_1, \dot{\sigma}_\infty, \dot{\tau}'_0, \dot{\tau}'_1, \dot{\tau}'_\infty$ il n'y a qu'une possibilité mod conjugaison (on a choisi $\sigma_0, \sigma_1, \sigma_\infty$, $\tau_0, \tau_1, \tau_\infty$ de la façon qui a semblé la plus commode pour les calculs), pour $\varepsilon$ il y en a trois - on a eu les trois choix $\tau_0, \tau_1, \tau_\infty$ - pour $\dot{\rho}, \dot{\rho}^{-1}$ il y en a deux, correspondantes aux points $-j, -j^2$ - nous avons choisi celle qui correspond au point $-j$. Les autres relèvements correspondants à $-j^2$ peuvent s'en déduire en conjuguant par un $\tau_i$ :

\newpage

%% ===== Page 89 =====

$$
\dots = \tilde{l}_0 \tilde{\rho} \circ \tilde{\rho} \tilde{l}_1 \tilde{\rho}^{-1} = \tilde{l}_0 \tilde{\rho}^2 \tilde{l}_1 \tilde{\rho}^{-1} = \tilde{l}_1 \tilde{\rho} = \tilde{\rho} \tilde{l}_1 \tilde{\rho}^{-1} \dots \leqno{(23)}
$$
\marginpar{* NB\\ $\rho \rho_0 \rho^{-1} = \rho_1$\\ $\rho \rho_1 \rho^{-1} = \rho_\infty$\\ $\rho \rho_\infty \rho^{-1} = \rho_0$}

Notons aussi que dans les formules usuelles de la théorie générale il suffit que les composantes [unclear: connexes] soient des [unclear: points] pour que [unclear: l'invariant] soit [unclear: trivial].

$$
\Pi_{0,3}^{\sigma_1} = \Pi_{0,3}^{\rho} = \Pi_{0,3}^{\tau_1} = \Pi_{0,3}^{\tau_1'} = \{1\} \leqno{(24)}
$$

Si on regarde le petit [unclear: relèvement] dans $\mathfrak{S}_{0,3}$, satisfaisant l'équation de [unclear: conjugaison], sur les points [unclear: fixes] $(2, -1, 1/2)$ et $(-j, -j)$, le stabilisateur d'un point de ces orbites est d'ordre 4 resp.\ 6. Pour le point 2 (et pour $\infty$), ce stabilisateur est le groupe $D_2 \simeq \mathbf{Z}/2\mathbf{Z} \times \mathbf{Z}/2\mathbf{Z}$, engendré par $\sigma_1, \tau_0$ (ou $\sigma_0, \tau_0$) \dots qui a bel et bien relevé [unclear: multiplicativement] dans $\mathfrak{S}_{0,3}$.
Il en est de même pour le point $-j$, dans $\mathfrak{S}_{0,3}'$, qui a été relevé également. Pour le [unclear: diviseur] général, on sait que les classes de [unclear: conjugaison] d'une relation correspondent aux points fixes. 
Le [unclear: groupe] [unclear: invariant] des [unclear: pointements] (sur $\mathcal{M}_{0,3}$), pour $\mathfrak{D}_3$ il n'y a qu'un seul point fixe, d'ordre 2, pour $\mathfrak{S}_3 = D_3$ il y en a deux, $-j$ et $-\bar{j}$, \dots [unclear: avant] d'aller [unclear: chercher] \dots des [unclear: conjugués].

\newpage

%% ===== Page 90 =====

en conjuguant par l'une quelconque des $\tau_i$ \\
$(\tau_0, \tau_1, \tau_\infty)$.

\vspace{1em}

III. Relations avec le gr. [unclear: cartographique] triangulaire non orienté,
$$ \Pi'_{0,3} = \pi_1((\mathbb{P}^1_{\mathbf{C}}, \{0\} + 2\{1\} + 3\{\infty\}) ; -j) $$
et
$$ = (\Pi_{0,3} \cdot \{1, \tau\}) / (l_1^2, l_\infty^3) \marginpar{avec $S\ell(2,\mathbf{Z})$} $$

On a des générateurs involutifs $\overline{\sigma}_0, \overline{\sigma}_1, \overline{\sigma}_2$ de $\Pi'_{0,3}$,
( $\sigma_i$ assoc. à un repère du cart. [unclear: fixant] la ième [unclear: arête] par symétrie par rapport au repère, donc $\dots$ changement d'orientation [unclear: tri-rectangles]
$\sigma_0$ et $\sigma_2$ sont resp. le passage au triangle [unclear: adjacent] resp. [unclear: transversal], $\sigma_1$ le passage au triangle [unclear: adjacent] )
avec les relations génératrices\footnote{On pourrait d'ailleurs utiliser $\overline{\sigma}_\infty$ au lieu de $\overline{\sigma}_2$ (qui en fait est $\overline{\sigma}_2 \overline{\sigma}_0$). On aurait une covariance parfaite des rôles (par conjugaison) entre $(\overline{\sigma}_0, \overline{\sigma}_1, \overline{\sigma}_\infty)$ et $(\tau_0, \tau_1, \tau_\infty)$.}
$$
\begin{cases}
\overline{\sigma}_0^2 = \overline{\sigma}_1^2 = \overline{\sigma}_2^2 = 1 \\
(\overline{\sigma}_0 \overline{\sigma}_1)^2 = 1 \\
(\overline{\sigma}_1 \overline{\sigma}_2)^3 = 1
\end{cases} \leqno{(25)}
$$

Le sous-groupe $\Pi''_{0,3}$ d'indice 2 [unclear: (groupe cart. dir. gal.-strict)] (qui est engendré
par $\overline{\sigma}_0 \overline{\sigma}_1 = \rho_f$ (associé à $l'_\infty$), $\overline{\sigma}_2 \overline{\sigma}_1 = \rho_s$ (associé à $l'_0$), $\overline{\sigma}_0 \overline{\sigma}_2 = \sigma$ (associé à $l'_1$), avec les
relations génératrices
$$
\rho_f \sigma \rho_s = 1 \quad \sigma^2 = \rho_f^3 = 1 \leqno{(26)}
$$
C'est donc le quotient de $\Pi_{0,3}$ par $l_1^2, l_\infty^3$
$$
\begin{tikzcd}
\Pi_{0,3} \arrow[r] & \Pi''_{0,3} \\
l_0 \arrow[r, mapsto] & l'_0 = \rho_s \\
l_1 \arrow[r, mapsto] & l'_1 = \sigma \\
l_\infty \arrow[r, mapsto] & l'_\infty = \rho_f
\end{tikzcd} \leqno{(27)}
$$
Cet homomorphisme se prolonge [unclear: en un prolongement]

\newpage

%% ===== Page 91 =====

$$
\tilde{\Pi}_{0,3} = \Pi_{0,3} \cdot \{t_1, t_2\} \longrightarrow \tilde{\Pi}'_{0,3} \leqno{(28)}
$$
$$
\tau_0, \tau_1, \tau_\infty \longmapsto \sigma_0, \sigma_1, \sigma_\infty
$$

Ceci posé, on définit un isomorphisme\footnote{Cet iso. est l'unique homomorphisme de $\tilde{\Pi}'_{0,3}$ (resp. de $\tilde{\Pi}_{0,3}$) sur $\tilde{\mathfrak{S}}_{0,3}$ qui induit l'iso. $\Pi'_{0,3} \isommap \mathfrak{S}^+_{0,3}$ (resp. [unclear: ...]).}
$$
\tilde{\Pi}'_{0,3} \isommap \tilde{\mathfrak{S}}_{0,3} \leqno{(29)}
$$
$$
(\sigma_0, \sigma_1, \sigma_\infty) \longmapsto (\tau'_0, \tau'_1, \tau'_\infty) \quad \text{notés aussi } \tau_0, \tau_1, \tau_\infty \text{ (réflexions fondamentales)}
$$
qui induit un iso.
$$
\Pi'_{0,3} \isommap \mathfrak{S}^+_{0,3} \leqno{(30)}
$$
$$
\begin{cases}
\rho_s = l'_0 \longmapsto \varepsilon_0 = \sigma_0 \rho \\
\sigma = l'_1 \longmapsto \sigma_\infty \\
\rho_f = l'_\infty \longmapsto \rho^{-1}
\end{cases}
\quad \text{NB } 
\begin{cases}
\tau'_0 \tau'_1 = \varepsilon_0 = l'_0 \\
\tau'_1 \tau'_\infty = \sigma_\infty = l'_1 \\
\tau'_0 \tau'_\infty = \rho^{-1} = l'_\infty
\end{cases}
$$
$$
\text{i.e. } \begin{cases}
\tau_0 \tau_1 = l'_0 (= \varepsilon_0) \\
\tau_1 \tau_\infty = l'_1 (= \sigma_\infty) \\
\tau_0 \tau_\infty = l'_\infty (= \rho^{-1})
\end{cases}
$$

On trouve donc que $\tilde{\mathfrak{S}}_{0,3}$ a des relations génératrices
$$
{\tau'_0}^2 = {\tau'_1}^2 = {\tau'_\infty}^2 = \underset{\varepsilon_0}{(\tau'_0 \tau'_1)^2} = \underset{\rho^{-1}}{(\tau'_0 \tau'_\infty)^3} = 1 \leqno{(31)}
$$
et $\mathfrak{S}^+_{0,3}$ a des relations génératrices
$$
l'_0 l'_1 l'_\infty = {l'_1}^2 = {l'_\infty}^3 = 1 \leqno{(32)}
$$
ou aussi simplement, en éliminant $l'_0$ et
posant $l'_1 = \sigma_\infty$ et $l'_\infty = \rho$
$$
\sigma_\infty^2 = \rho^3 = 1 \leqno{(33)}
$$

Pour mémoriser ces expressions des éléments remarquables de $\mathfrak{S}^+_{0,3}$ en termes de $\sigma_\infty, \rho$ :

\newpage

%% ===== Page 92 =====

(34)
$$ \begin{cases}
\varepsilon_0 = \sigma_\infty \rho & l_0 = \sigma_\infty \rho \sigma_\infty \rho = (\sigma_\infty \rho)^2 \\
\varepsilon_1 = \rho \sigma_\infty & l_1 = \rho \sigma_\infty \rho \sigma_\infty = (\rho \sigma_\infty)^2 \\
\varepsilon_\infty = \rho^{-1} \sigma_\infty \rho^{-1} & l_\infty = \rho^{-1} \sigma_\infty \rho^{-1} \rho^{-1} \sigma_\infty \rho^{-1} = (\rho^{-1} \sigma_\infty \rho^{-1})^2
\end{cases} $$

(35)
$$ \begin{cases}
\sigma_0 = \rho \sigma_\infty \rho^{-1} & (= \varepsilon_1 \rho^{-1} = \rho \varepsilon_1^{-1}) \\
\sigma_1 = \rho^{-1} \sigma_\infty \rho & (= \rho^{-1} \varepsilon_0 = \varepsilon_\infty \rho^{-1})
\end{cases} $$

et dans $\widehat{G}_{0,3}$, en fonction de $\tau_\infty', \tau_0', \tau_\infty$
les expressions de $\tau_0, \tau_1, \tau_\infty'$ :
(36)
$$ \begin{cases}
\tau_0 = \rho \tau_\infty' \rho^{-1} = \tau_\infty' \tau_0' \tau_\infty' \tau_0' \tau_\infty' & (= r_0 r_1 r_\infty r_1 r_0) \\
\tau_1 = \rho^{-1} \tau_\infty' \rho = \tau_0' \tau_\infty' \tau_0' \tau_\infty' \tau_0' & (= r_1 r_0 r_\infty r_0 r_1) \\
\tau_\infty' = \rho \tau_0' \rho^{-1} = \tau_\infty' \tau_0' \tau_\infty' & (= r_0 r_1 r_0)
\end{cases} \qquad \begin{cases} \tau_\infty' = r_0 \\ \tau_0' = r_1 \\ \tau_\infty = r_\infty \end{cases} $$

## IV) Relations avec $\mathcal{G}l(2, \mathbb{Z}) = \operatorname{Gl}(2, \mathbb{Z})/\pm 1 \simeq \operatorname{Aut}(\mathbb{P}^1_{\mathbb{Z}})$.

Ce groupe opère sur le demi-plan de Poincaré
$$D = \{z \in \mathbb{C} \mid \Im z > 0\}$$

par
(37)
$$ \begin{cases}
u_g(z) = \frac{az+b}{cz+d} & \text{si } ad-bc = +1 \\
u_g(z) = \frac{a\bar{z}+b}{c\bar{z}+d} & \text{si } ad-bc = -1
\end{cases} \qquad g = \begin{pmatrix} a & b \\ c & d \end{pmatrix} \quad a,b,c,d \in \mathbb{Z}, \quad ad-bc \in \pm 1 $$

On a dans $D$ un domaine fondamental bien connu, qui correspond au triangle sphérique $(-J, 1, 0)$ pour l'homéomorphisme du demi-plan naturel à $\Delta$ envoyant $\infty$ à $1$, $\rho$ à $-J$, et le point $i$ à $0$. Ainsi [unclear: on a les symétries] par rapport aux trois côtés de ce triangle sont donnés par

(38)
$$ r_0(z) = -\bar{z}, \quad r_1(z) = 1-\bar{z}, \quad r_\infty(z) = 1/\bar{z} $$

Ils décrivent donc [unclear: le groupe de symétries]
et les produits deux à deux donnent les [unclear: rotations affines]

\newpage

%% ===== Page 93 =====

$$
\begin{cases}
r_\infty r_1(z) = \ell_0'(z) = z-1 \\
r_0 r_\infty(z) = \ell_1'(z) = -\frac{1}{z} \\
r_1 r_0(z) = \ell_\infty'(z) = 1-\frac{1}{z} = \frac{z-1}{z}
\end{cases} \leqno (39)
$$

On trouve ainsi un isomorphisme
$$
\mathfrak{S}_{0,3} \isommap \operatorname{Gl}(2,\mathbf{Z})' \simeq \operatorname{Gl}(2,\mathbf{Z})/\pm 1 \simeq GP(1,\mathbf{Z}) \leqno (40)
$$
$$
\left.
\begin{aligned}
r_0 &= \tau_0' \longmapsto r_0 = \begin{pmatrix} 0 & 1 \\ 1 & 0 \end{pmatrix} \\
r_1 &= \tau_1' \longmapsto r_1 = \begin{pmatrix} -1 & 1 \\ 0 & 1 \end{pmatrix} \\
r_\infty &= \tau_\infty \longmapsto r_\infty = \begin{pmatrix} 1 & 0 \\ 0 & -1 \end{pmatrix}
\end{aligned}
\right\} \footnote{il vaut mieux des signes changés cf plus loin}
$$

et une présentation de $GP(1,\mathbf{Z})$ par
$$
r_0^2 = r_1^2 = r_\infty^2 = (r_0 r_\infty)^2 = (r_1 r_0)^3 = 1. \leqno (41)
$$

Ceci donne un isomorphisme
$$
\mathfrak{S}_{0,3}^+ \isommap \operatorname{Sl}(2,\mathbf{Z})' \defeq \operatorname{Sl}(2,\mathbf{Z})/\pm 1 \simeq GP(1,\mathbf{Z})^+ \leqno (42)
$$
$$
\begin{cases}
\ell_1' = \sigma_\infty \longmapsto \begin{pmatrix} 0 & -1 \\ 1 & 0 \end{pmatrix} \\
\ell_\infty' = \rho^{-1} \longmapsto \begin{pmatrix} 1 & -1 \\ 1 & 0 \end{pmatrix}
\end{cases}
$$
(d'où
$$
\rho \longmapsto \begin{pmatrix} 0 & 1 \\ -1 & 1 \end{pmatrix} \quad ).
$$

On identifiera par la suite les éléments de $\mathfrak{S}_{0,3}$ avec leurs images dans $\operatorname{Gl}(2,\mathbf{Z})' \simeq GP(1,\mathbf{Z})$.

Notons la formule [unclear: suivante bien] utile pour
pour une matrice $u = \begin{pmatrix} a & b \\ c & d \end{pmatrix} \in \operatorname{Gl}(2,\mathbf{Z})$

\newpage

%% ===== Page 94 =====

\marginpar{NB 2 : Si $u \in \operatorname{Gl}(2, A)$, $A$ anneau commutatif, on a pour $u = \begin{pmatrix} a & b \\ c & d \end{pmatrix}$, $\det(u) \cdot u^{-1} = \begin{pmatrix} d & -b \\ -c & a \end{pmatrix}$}
$$
u^{-1} = \begin{pmatrix} d & -b \\ -c & a \end{pmatrix} \leqno{(43)}
$$
\marginpar{cf ``formulaire'' pour l'étude des relations d'incidence plus complètes}
$$
\begin{aligned}
\varepsilon_0 &= \sigma_{\infty} \rho = \begin{pmatrix} 1 & -1 \\ 0 & 1 \end{pmatrix} \quad , \quad \varepsilon_0^{-1} = \begin{pmatrix} 1 & 1 \\ 0 & 1 \end{pmatrix} \\
\varepsilon_1 &= \rho \sigma_{\infty} = \begin{pmatrix} 1 & 0 \\ 1 & 1 \end{pmatrix} \quad , \quad \varepsilon_1^{-1} = \begin{pmatrix} 1 & 0 \\ -1 & 1 \end{pmatrix} \\
\varepsilon_{\infty} &= \rho^{-1} \sigma_{\infty} \rho^{-1} = \begin{pmatrix} 2 & -1 \\ 1 & 0 \end{pmatrix} \quad , \quad \varepsilon_{\infty}^{-1} = \begin{pmatrix} 0 & 1 \\ -1 & 2 \end{pmatrix} \\
\ell_0 &= \varepsilon_0^2 = \begin{pmatrix} 1 & -2 \\ 0 & 1 \end{pmatrix} \qquad \qquad \qquad \qquad \sigma_0 = \begin{pmatrix} 1 & -1 \\ 2 & -1 \end{pmatrix} \\
\ell_1 &= \varepsilon_1^2 = \begin{pmatrix} 1 & 0 \\ 2 & 1 \end{pmatrix} \qquad \qquad \qquad \qquad \sigma_1 = \begin{pmatrix} 1 & -2 \\ 1 & -1 \end{pmatrix} \\
\ell_{\infty} &= \varepsilon_{\infty}^2 = \begin{pmatrix} 3 & -2 \\ 2 & -1 \end{pmatrix} \qquad \qquad \qquad \qquad \sigma_{\infty} = \begin{pmatrix} 0 & -1 \\ 1 & 0 \end{pmatrix}
\end{aligned} \leqno{(44)}
$$

$$
\rho = \begin{pmatrix} 0 & 1 \\ -1 & 1 \end{pmatrix} \quad , \quad \rho^{-1} = \begin{pmatrix} 1 & -1 \\ 1 & 0 \end{pmatrix} \leqno{(45)}
$$

$$
\ell_{\infty} \ell_1 \ell_0 = \begin{pmatrix} -1 & 0 \\ 0 & -1 \end{pmatrix} \quad \begin{array}{l} \text{l'élément non trivial} \\ \text{du centre de } \operatorname{Sl}(2,\mathbf{Z}) \\ \text{qui (heureusement !) donne} \\ 1 \text{ dans } \operatorname{Sl}(2,\mathbf{Z})' \end{array} \leqno{(46)}
$$

Le groupe $\Pi_{0,3}$ se récupère comme le noyau de l'homomorphisme
$$
\begin{tikzcd}
\operatorname{Sl}(2,\mathbf{Z})' \arrow[r] & \operatorname{Sl}(2, \mathbf{Z}/2\mathbf{Z}) \\
& \mathfrak{S}_3 \arrow[u, "\simeq"']
\end{tikzcd}
$$
la suite exacte
$$
\begin{tikzcd}[row sep=small]
1 \arrow[r] & \Pi_{0,3} \arrow[r] & \mathfrak{S}_{0,3}^+ \arrow[r] & \mathfrak{T}_{0,3}^+ \arrow[r] & 1 \\
& & & \mathfrak{S}_3 \arrow[u, "\simeq"'] &
\end{tikzcd} \leqno{(47)}
$$
se récupère comme
$$
\begin{tikzcd}[row sep=small]
1 \arrow[r] & \pi(2)' \arrow[r] & \operatorname{Sl}(2,\mathbf{Z})' \arrow[r] & \operatorname{Sl}(2, \mathbf{Z}/2\mathbf{Z}) \arrow[r] & 1 \\
& \Pi_{0,3} \arrow[u, "\parallel"'] & & \mathfrak{S}_3 \arrow[u, "\simeq"'] &
\end{tikzcd} \leqno{(48)}
$$

\newpage

%% ===== Page 95 =====

Ce [unclear: noeud] formelle dans les mains de [unclear: Teichmüller] dans $\operatorname{Sl}(2,\mathbf{Z})$ lui-même, extension centrale de $\mathfrak{S}_{0,3}$ - en fait, cette extension s'interprète comme $\mathfrak{T}_{1,1}$, mais nous y reviendrons ultérieurement...

V) Résumons les avatars des divers groupes, que nous avons rencontrés jusqu'à présent :

\begin{enumerate}
    \item[1)] $\mathfrak{S}_{0,3} \defeq \operatorname{Aut}_{\mathrm{loc}}(\Pi_{0,3})$ \quad avec sous-groupe $\mathfrak{S}_{0,3}^+$ (respectant orientation)
    \item[2)] $\mathfrak{T}_{0,4}' \simeq \Autext_{\mathrm{loc}}(\Pi_{0,4})'$ \quad automorphisme extérieur de $\Pi_{0,4}$ respectant la structure à lacets et fixant la classe de lacets $L_4$ (et permutant $L_1, L_2, L_3$)
    
    [unclear: $\mathfrak{S}_{0,4}^{+'}$ composantes orientées]
    \item[3)] $A_{0,4}' / A_{0,4}^0$ groupe des composantes connexes de $A_{0,4}' \subset \operatorname{Aut}_{\mathrm{top}}(U_{0,4}^{\mathrm{top}})$ sous-groupe des automorphismes qui fixent le point : $x_4=1/2$
    
    avec $A_{0,4}' / A_{0,4}^0$ (composantes orientées) $\simeq \mathfrak{S}_{0,3}$
    \item[4)] $\pi_1( (U_{0,3}^{\mathrm{top}}, \mathfrak{S}_{0,3}, -J) )$ (décrit à la J. Malgoire)
    
    avec le sous-groupe $\pi_1( (U_{0,3}^{\mathrm{top}}, \mathfrak{S}_{0,3}^+) )$ \quad $(\mathfrak{S}_3 = \mathfrak{T}_{0,3}^+)$
    \marginpar{[unclear: cette subtile groupe] $\pi_{0,3} \simeq \Pi_{0,3} \cdot \{1, \varepsilon \}$}
    \item[5)] $\pi_1( (\mathbb{P}_{\mathbf{C}}^{1 \wedge}, \infty \cdot \{0\} + 2 \cdot \{1\} + 3 \cdot \{\infty\} ), -J )$ (à la J. Malgoire)
    
    avec le sous-groupe $\pi_1( (\mathbb{P}_{\mathbf{C}}^{1 \wedge}, \infty \cdot \{0\} + 2 \cdot \{1\} + 3 \cdot \{\infty\} ) )$
    \item[6)] Groupe cartographique triangulaire $\bar{\Gamma}_{0,3}$
    
    $$ \bar{\Gamma}_{0,3} = \{ \tau_0, \tau_1, \tau_\infty \mid \tau_0^2 = \tau_1^2 = \tau_\infty^2 = (\tau_1\tau_\infty) = (\tau_\infty\tau_0)^2 = (\tau_0\tau_1)^3 = 1 \} $$
    
    (avec sous-groupe des triangles orientés
    
    $$ \Gamma_{0,3} = \{ \rho_0, \rho_1, \rho_\infty \mid \rho_\infty\rho_1\rho_0 = \rho_0^2 = \rho_1^3 = 1 \} \quad ) $$
    \vspace{-0.5cm}
    $$ \hspace{2cm} \underset{\tau_1\tau_\infty}{\uparrow} \hspace{0.3cm} \underset{\tau_\infty\tau_0}{\uparrow} \hspace{0.3cm} \underset{\tau_0\tau_1}{\uparrow} $$

    \marginpar{NB $\mathfrak{A}_{0,3}^{(loc)} \simeq \pi_1((U_{0,3}, \mathfrak{T}_{0,3}))$ ce qui permet d'identifier 7) \& 4)}
    \item[7)] $\pi_1( U_{0,3}^{\mathrm{conf}}, pt(U_{0,3}, -J) ) \quad \text{avec sous-groupe} \quad \pi_1( U_{0,3}, pt(U_{0,3}, -J) )$
    
    \begin{tabular}{ll}
    domaine de [unclear: structures] & sur les [unclear: holomorphes] \\
    universelles de & universelles de type \\
    type $(0,3)$ & $(0,3)$
    \end{tabular}

    \item[8)] $\mathfrak{T}_{1,1}' = \mathfrak{T}_{1,1} / \underset{\text{centre}}{\pm 1}$ \quad avec \quad $\mathfrak{T}_{1,1} \simeq \Autext_{\mathrm{loc}}(\Pi_{1,1}) \simeq A_{1,1} / A_{1,1}^0$
    
    ($A_{1,1} \simeq$ automorphisme topologique de $U_{1,1}^{\mathrm{top}}$) sous-groupe orienté $\mathfrak{T}_{1,1}^+ \subset \mathfrak{T}_{1,1}'$
\end{enumerate}

\newpage

%% ===== Page 96 =====

9) $\pi_1(M_{1,1}^{\text{anf}}, \text{pt})' \simeq \pi_1(\ )/\operatorname{centre}\{\pm 1\} \simeq \pi_1(M_{1,1}^{\prime \text{anf}}, \text{pt})$

où $\underline{M}_{1,1}^{\text{an}}$ est la multiplicité modulaire [unclear: conforme] analytique complexe des courbes elliptiques (i.e. courbes [unclear: conformes] de type 1,1),

$\underline{M}_{1,1}^{\prime \text{an}}$ celle des courbes elliptiques analytiques "mod. symétrie", avec sous-groupes

$\pi_1(M_{1,1}^{\text{an}}, \text{pt})^+ \simeq \pi_1(\ )/\operatorname{centre}\{\pm 1\} \simeq \pi_1(M_{1,1}^{\prime \text{an}}, \text{pt})^+$

où le pt base est la courbe ell. conforme [unclear: an.], quotient de $\mathbf{C}$ par $\mathbf{Z} + \mathbf{Z}(-J)$.

NB Les multiplicités modulaires sont $(D, \operatorname{Gl}(2, \mathbf{Z}))$, $(D, \operatorname{Sl}(2, \mathbf{Z}))$ dans le cadre conforme resp. analytique pour $M_{1,1}$,

et $(D, \operatorname{Gl}(2, \mathbf{Z})' = \operatorname{PGL}(2, \mathbf{Z}))$ resp. $(D, \operatorname{Sl}(2, \mathbf{Z})' = \operatorname{PSL}(2, \mathbf{Z}))$ dans le cadre conforme resp. analytique pour $M_{1,1}'$.

10) $\operatorname{Gl}(2, \mathbf{Z})' \simeq \operatorname{PGL}(2, \mathbf{Z})$, avec le sous-groupe $\operatorname{Sl}(2, \mathbf{Z})' \simeq \operatorname{PSL}(2, \mathbf{Z})$

\newpage

%% ===== Page 97 =====

(VI) Les automorphismes à l'ext. de $\widehat{\Pi}_{0,3}$ et de $\widehat{\mathbb{C}}_{0,3}^+$

On a un diagramme commutatif canonique :

(49)
\begin{tikzcd}
\operatorname{Aut}\,ext(\widehat{\mathbb{C}}_{0,3}^+)' \ar[r, "\text{restriction}", "\sim"'] & \operatorname{Aut}\,ext^*(\widehat{\Pi}_{0,3})' \\
\operatorname{Aut}(\widehat{\mathbb{C}}_{0,3}^+, \sigma_\infty)' \ar[u, "\wr"] \ar[r, "\text{restriction}", "\sim"'] & \operatorname{Aut}(\widehat{\Pi}_{0,3}, \sigma_\infty) \ar[u, "\wr"']
\end{tikzcd}

où on a :
$\operatorname{Aut}\,ext(\widehat{\mathbb{C}}_{0,3}^+)' = \text{groupe des autom. ext. induisant l'identité sur le quotient } \mathfrak{S}_3 \simeq \widehat{\mathbb{C}}_{0,3}^+/\widehat{\Pi}_{0,3}$

$\operatorname{Aut}\,ext^*(\widehat{\Pi}_{0,3})' = \text{groupe des autom. extérieurs qui commutent à l'action nat. de } \mathfrak{S}_3$

$\operatorname{Aut}(\widehat{\mathbb{C}}_{0,3}^+, \sigma_\infty)' = \text{groupe des autom. de } \widehat{\mathbb{C}}_{0,3}^+ \text{ qui fixent } \sigma_\infty, \text{ et qui fixent } \rho \bmod \widehat{\Pi}_{0,3}$

$\operatorname{Aut}(\widehat{\Pi}_{0,3}, \sigma_\infty) = \text{groupe des aut. qui commutent à } \sigma_\infty \text{ et qui } [unclear] \text{ ext. } id$

Les $u \in \operatorname{Aut}(\widehat{\mathbb{C}}_{0,3}^+, \sigma_\infty)'$ sont de la forme

[MARGIN: On écrit les formules dans $\operatorname{Aut}(\widehat{\mathbb{C}}_{0,3}^+)' = \mathcal{G}$ groupe des [unclear: classes d'] autom. extérieurs qui ind. l'identité dans $\mathcal{G}/\widehat{\Pi}_{0,3} \simeq \mathfrak{S}_3$, et $\operatorname{Aut}(\widehat{\Pi}_{0,3})' = gr. des [unclear]$ ext. de $\dots$ NB. On admet ici que le centre de $\widehat{\mathbb{C}}_{0,3}^+$ est trivial.]

(50) 
$$
\begin{cases}
u(\sigma_\infty) = \sigma_\infty \\
u(\rho) = g \rho \qquad (i.e.\ u \rho u^{-1} = g)
\end{cases}
$$

où $g \in \widehat{\Pi}_{0,3}$ est assujetti à la condition de
$$ \boxed{g \rho g \rho g \rho = 1} $$
(51) exprimant la relation
(52) $$(g \rho)^3 = 1$$

On trouve ainsi une corresp. 1-1 entre les éléments de $\operatorname{Aut}(\widehat{\mathbb{C}}_{0,3}^+, \sigma_\infty)'$ et les $g \in \widehat{\Pi}_{0,3}$ satisfaisant (51), et tels de plus que

(53) $\sigma_\infty, g \rho \text{ avec les relations } \sigma_\infty^2 = (g \rho)^3 = 1 \text{ présentent } \widehat{\mathbb{C}}_{0,3}^+$

\newpage

%% ===== Page 98 =====

(que je vais appeler la ``condition liminaire'').

Pour $u \in \mathcal{G} = \operatorname{Aut}(\widehat{\mathfrak{S}}_{0,3}^+)'$ (ne fixant pas nécessairement $\sigma_\infty$), on a des conditions équivalentes
$$
\begin{array}{ll}
\mathrm{(i)} & u(l_0) \in \text{conjugué de } l_0^{\hat{\mathbf{Z}}} = L_0 \quad (\text{dans } \widehat{\mathfrak{S}}_{0,3}^+) \\
\mathrm{(ii)} & u(l_0) \in \text{conjugué de } l_0^{\hat{\mathbf{Z}}} = L_0 \quad (\text{dans } \hat{\Pi}_{0,3}) \\
\mathrm{(iii)} & u(l_i) \in \text{conjugué de } l_i \text{ pour } i = 0, 1, \infty \\
\mathrm{(iv)} & \exists p \in \hat{\mathbf{Z}}^*, \ g_0 \in \hat{\Pi}_{0,3} \qquad \boxed{u(l_0) = \operatorname{int}(g_0)(l_0^p)} \\
\mathrm{(v)} & \exists p \in \hat{\mathbf{Z}}^*, \ g_\infty \in \hat{\Pi}_{0,3} \qquad \boxed{u(l_\infty) = \operatorname{int}(g_\infty)(l_\infty^p)}
\end{array}
\leqno{(54)}
$$

On aura alors aussi
$$
\begin{cases}
u(l_1) = \operatorname{int}(g_1)(l_1^p) & g_1 = g p(g_0) \\
u(l_\infty) = \operatorname{int}(g_\infty)(l_\infty^p) & g_\infty = g p(g_1) = g p(g) p^2(g_0)
\end{cases}
$$
$$
p \in \hat{\mathbf{Z}}^* \qquad p = \chi(u) \quad (\text{où } u \text{ donné par } (50)), \leqno{(55)}
$$
est bien déterminé par $u$, donc par $g$, on l'appelle le \underline{multiplicateur} de $u$.
$$
g_0 \in \hat{\Pi}_{0,3} \leqno{(56)}
$$
est déterminé à [unclear: un facteur droit] $L_0 = l_0^{\hat{\mathbf{Z}}}$ (changeant
$$
g_0 \mapsto g_0 l_0^\alpha \qquad ). \leqno{(57)}
$$

\marginpar{Notations :\\ $\Autext(\widehat{\mathfrak{S}}_{0,3}^+)'$\\ $\operatorname{Aut}(\widehat{\mathfrak{S}}_{0,3}^+)'$\\ $\Autext(\hat{\Pi}_{0,3}, \sigma_\infty)'$\\ $\operatorname{Aut}(\hat{\Pi}_{0,3}, \sigma_\infty)'$}
Si les conditions (54) sont satisfaites, on dit que $u$ \underline{respecte la structure à lacets} de $\widehat{\mathfrak{S}}_{0,3}^+$ (ou de $\hat{\Pi}_{0,3}$) -- en sous-entendant pour $\widehat{\mathfrak{S}}_{0,3}^+$ le choix des ``générateurs'' $l_0^{\hat{\mathbf{Z}}} = L_0$. On aura beaucoup, par la suite à [unclear: faire] à de tels autom.. NB Ces conditions ne dépendent que de l'autom. ext. défini par $u$, dans $\hat{\Pi}_{0,3}$ ou dans $\widehat{\mathfrak{S}}_{0,3}^+$.

\newpage

%% ===== Page 99 =====

on s'intéresse aux relèvements canoniques de tels autom.\ extérieurs donnés par (49), (50), en des autom.\ effectifs.

La formule de définition (50) de $u$ ne fait intervenir que $g_0 \in \hat{\Pi}_{0,3}$ -- en principe $g_0 \pmod{L_0}$ dans $\hat{\Pi}_{0,3}/L_0$ se déduit de $g$. Mais en pratique $g_0$ joue un rôle aussi important que $g$, on verra que dans un cas au moins il détermine $g$ (formule (62)). S'il n'y avait pas lieu d'introduire $g_0$, pour exprimer les conditions (54), on pourrait aisément paramétrer les $g$ satisfaisant à (53) et donc correspondants à des relèvements de $\bar{g} \in \mathfrak{S}_3$ en des éléments d'ordre 3 de $\hat{\mathfrak{S}}_{0,3}^+$ -- en notant qu'il y a exactement deux classes de conjugaison de tels relèvements, (correspondants aux pts fixes $j$ et $-j$ de $\bar{g}$ opérant sur $U_{0,3}^{an}$), qui se paramétrisent par\marginpar{cette représentat\textsuperscript{ion} se trouvant d'ailleurs mal appropriée par la suite.}
$$
\begin{cases}
g = \lambda_0 \rho \lambda_0^{-1} & \text{(classe principale, correspond\textsuperscript{t}} \\
& \text{au ``pôle nord'' } j) \\
g = \lambda_0' \rho^{-1} {\lambda_0'}^{-1} & \text{(classe adjointe, correspond\textsuperscript{t}} \\
& \text{au ``pôle sud'' } -j)
\end{cases} \leqno{(58)}
$$
où $\lambda_0, \lambda_0' \in \hat{\Pi}_{0,3}$. On verra plus bas que ces cas correspondent respectivement aux cas
$$
p \equiv 1 \pmod 3 \quad \text{et} \quad p \equiv -1 \pmod 3 \leqno{(59)}
$$

\newpage

%% ===== Page 100 =====

Mais vu la réduction à ces conditions par la
nécessité de tenir compte des conditions (59)
(par sur les $f_i \to (iv)$) -- d'où l'introduction
conjointement $g_0, \rho$. Il y a lieu alors d'introduire
la quantité
$$
T = \frac{\rho-1}{2} \in \hat{\mathbf{Z}} \leqno{(60)}
$$
(NB on suppose $\rho \in \hat{\mathbf{Z}}^*$, se réduit dans $\hat{\mathbf{Z}}/2\hat{\mathbf{Z}}$ à $1$,
i.e.\ $\rho - 1 \in 2\hat{\mathbf{Z}}$, d'où la possibilité de trouver $T$),
$$
h_0 = g_0^{-1} \sigma_\infty(g_0)
$$
(que nous noterons, puisque nous l'avons pas
fait $l_0 = \sigma_\infty(g_0) g_0$, pour [unclear: prendre] l'inversion par
rapport aux notations précédentes\footnote{NB. rectif: l'aut.\ $\sigma_\infty$ d'ordre $2$, que j'ai considéré au § 27, 28 est le symétrique par rapport au centre de celui que j'ai pris dans ces notations.}). On vérifie alors que
$g$ s'exprime à l'aide de $g_0$ et de $T$ (donc de $\rho$)
à l'aide de la formule (équivalente à (54) (iv) ou (v))
\marginpar{NB l'aut.\ expr.\ d'ordre 2 $g_0 \mapsto \sigma_\infty(g_0) = g_0 h_0$}
$$
g = \sigma_\infty(g_0) l_0^T \rho(g_0)^{-1} = (g_0 h_0) (h_0^T h_0) \rho(g_0 h_0)^{-1} \leqno{(62)}
$$
Dans la suite, ce qui nous intéresse se décrit
par les couples $(g_0, T) \quad (g_0 \in \widehat{\mathfrak{S}}_{0,3}, T \in \hat{\mathbf{Z}})$
satisfaisant la condition
$$
2T+1 \in \hat{\mathbf{Z}}^* \quad (\text{NB } 2T+1 = \rho = \chi(\dots)) \leqno{(63)}
$$
et la condition (51), où $g$ est exprimé en
fonction de $g_0$ et de $T$ par (62). Cette
relation (51) équivaut à son tour à la relation
$$
[ \{ \xi \} , \{ \eta \} ] = 1 \leqno{(64)}
$$
$$
\xi = l_0^T h_0 \quad \text{avec (66)} \quad h_0 = g_0^{-1} \sigma_\infty(g_0) \dots \leqno{(65)}
$$
\marginpar{NB. On rappelle que $g_0 \in \widehat{\mathfrak{S}}_{0,3}$ dont on déduit $g$ satisfait à ...}

\newpage

%% ===== Page 101 =====

Ainsi, les $u$ sont paramétrisés par les couples $(q_0 \in \hat{\Pi}_{0,3}/L_0, \gamma \in \hat{\mathbf{Z}})$, satisfaisant les conditions (63) et (64), qui moyennant (65) s'écrit sous forme explicite\footnote{Il vaudrait mieux considérer cette équation comme l'« équation en $h_0$ » (symétrique en $h_0$).} :
$$
(l_0^\gamma \underbrace{q_0^{-1}\sigma_\infty(q_0)}_{h_0}) \rho(l_0^\gamma \underbrace{q_0^{-1}\sigma_\infty(q_0)}_{h_0}) \rho^2(l_0^\gamma \underbrace{q_0^{-1}\sigma_\infty(q_0)}_{h_0}) = 1 \leqno{(65 \text{ bis})}
$$
(« équation en $q_0$ » - par opposition avec l'« équation en $\xi$ » plus bas). (De plus, il faut ajouter la condition liminaire, pour mémoire...)

On peut aussi considérer que $u$ (via $g$) est défini par $(\xi \in \hat{\Pi}_{0,3}, \gamma \in \hat{\mathbf{Z}})$, qui permettent en effet de retrouver
$$
h_0 = l_0^{-\gamma} \xi
$$
par (65), puis $q_0$ par (66) :
$$
q_0^{-1} \sigma_\infty(q_0) = h_0
$$
qui détermine en effet $q_0$ de façon unique, grâce au fait (qu'on admet provisoirement) que $\hat{\Pi}_{0,3}^{\sigma_\infty} = \{1\}$. La possibilité de trouver $q_0$ satisfaisant à cette condition équivaut (on l'admettra [unclear: aussi] ici) : la condition
$$
\sigma_\infty(h_0) h_0 = 1 \leqno{(67)}
$$
qui, en termes de $\xi, \gamma$, s'écrit aussi sous la forme

\newpage

%% ===== Page 102 =====

$$
\sigma_\infty(\xi) = l_1^\gamma \xi^{-1} l_0^{-\gamma} \leqno{(68)}
$$

Les équations (64), (68) sont appelées les "équations en $\xi$"\footnote{Plutôt "équations en $h_0$".} - par opposition à "l'équation en $q_0$" (66), qui à première vue a l'avantage d'être une seule, et le désavantage d'être plus compliquée. En résumé, les $u$ qui nous intéressent ($\in \operatorname{Aut}(\widehat{\mathfrak{S}}_{0,3}^+, \sigma_\infty)'$) sont donnés par les couples $(\xi \in \widehat{\Pi}_{0,3}, \gamma \in \hat{\mathbf{Z}})$, satisfaisant les conditions

\begin{enumerate}
\item[a)] $2\gamma + 1 \in \hat{\mathbf{Z}}^\times \hfill (63)$
\item[b)] $\xi \rho(\xi)\rho^2(\xi) = 1 \hfill (64)$
\item[c)] $\sigma_\infty(\xi) = l_1^\gamma \xi^{-1} l_0^{-\gamma} \hfill (68)$
\item[d)] Condition liminaire pour assurer (53) (qui ne s'exprime apparemment pas directement de façon commode en termes de $\xi, \gamma$).
\end{enumerate}

Pour que deux couples $(\xi, \gamma)$, $(\xi', \gamma')$ définissent le même $u$, il faut et il suffit qu'on ait :
$$
\gamma = \gamma', \quad \xi' = l_0^{-\alpha} \xi l_1^\alpha \quad (\alpha \in \hat{\mathbf{Z}}) . \leqno{(69)}
$$

La condition (64) peut encore se remplacer par une représentation paramétrique
$$
\begin{cases}
\xi = \lambda \rho(\lambda)^{-1} \\
\xi = \lambda l_0 \rho(\lambda)^{-1}
\end{cases} \text{ ou } \leqno{(70)}
$$
avec $\lambda \in \widehat{\Pi}_{0,3}$ --- auquel cas l'équation (68)

\newpage

%% ===== Page 103 =====

peut s'interpréter comme une "équation en $\lambda$" qui suivant les deux cas prend la forme
$$
\begin{cases}
\sigma_\infty(\lambda \rho(\lambda)^{-1}) = l_1^\gamma (\rho(\lambda) \lambda^{-1}) l_0^{-\gamma} \\
\text{ou} \\
\sigma_\infty(\lambda l_0^{-1} \rho(\lambda)^{-1}) = l_1^\gamma (\rho(\lambda) l_0 \lambda^{-1}) l_0^{-\gamma}
\end{cases} \leqno{(71)}
$$
Pour que $(\lambda, \gamma)$ et $(\lambda', \gamma')$ définissent le même $u$, il faut et il suffit qu'on ait
$$
\lambda' = l_0^{-\alpha} \lambda \quad , \quad \gamma = \gamma' \leqno{(72)}
$$

Remarque 1) Au facteur "anodin" $l_0^\gamma$ près, qui différencie $h_0$ et $\xi$ (65), on peut dire que la recherche des $u \in \operatorname{Aut}_{loc}(\widehat{\mathfrak{S}}_{0,3}^+, \sigma_\infty)$ équivaut à celle des solutions simultanées des deux équations de cocycles
$$
\begin{cases}
\xi \rho(\xi) \rho^2(\xi) = 1 \\
h_0 \sigma_\infty(h_0) = 1
\end{cases} \leqno{(73)}
$$
relatifs respectivement aux groupes $\mathbf{Z}/3\mathbf{Z}$, opérant sur $\widehat{\Pi}_{0,3}$ via $\rho$, et $\mathbf{Z}/2\mathbf{Z}$, opérant via $\sigma_\infty$, dont les solutions se donnent paramétriquement respectivement par
$$
\begin{cases}
\xi = \lambda \rho(\lambda)^{-1} \quad \text{ou} \quad \xi = \lambda l_0^{-1} \rho(\lambda)^{-1} \\
\qquad (\text{suivant que } \gamma \equiv 1 \ (3) \text{ ou } \gamma \equiv -1 \ (3) \\
\qquad \text{i.e.\ } \gamma \equiv 1 \ (3) \text{ ou } \gamma \equiv -1 \ (3)) \\
h_0 = g_0 \sigma_\infty(g_0)^{-1} \quad \text{[unclear: (NB on a pris } g_0 \dots \text{)]}
\end{cases} \leqno{(74)}
$$
\footnote{NB l'expression (62) pour $\xi$ met aussi en évidence l'intervention de $l_0$ dans la formule (74).}

\newpage

%% ===== Page 104 =====

341

3) La relation entre le $\lambda_0$ dans la représentation paramétrique de $g$ (58), et le $\lambda$ dans celle de $\xi$ (74 a) est implicite dans la deuxième expression (62) de $g$, savoir

(75) $$ \lambda_0 = g_0 l_0^{-\gamma} \lambda $$

4) Notons ici [unclear: que] les quantités [unclear: précédemment considérées] (avec les indéterminations
$$g, \gamma, g_0, h_0, \xi, \lambda_0, \lambda$$
$$\in \quad \in \quad \in \quad \in \qquad \in$$
$$\hat{\mathbb{Z}}^\times \quad \hat{\mathbb{Z}} \quad \hat{\Pi}_{0,3} \quad \hat{\Pi}_{0,3}/L_0 \quad \hat{\Pi}_{0,3}$$

spécifiques à $g_0, h_0, \xi, \lambda$) peuvent être considérées, [unclear: pour un repérage] fixé, comme des fonctions sur le groupe 
$$Autext_{Loc}(\widehat{\mathfrak{G}}_{0,3}^+)' \simeq Autext_{Loc} (\hat{\Pi}_{0,3})'$$ (sous-groupe des autom. ext. continus du groupe $\hat{\Pi}_{0,3}$ qui commutent jusqu'à $\mathfrak{S}_3$), à valeurs dans les ensembles $\hat{\mathbb{Z}}^\times, \hat{\mathbb{Z}}, \hat{\Pi}_{0,3}, \hat{\Pi}_{0,3}/L_0, \, \cdot \, , \hat{\Pi}_{0,3}$.

En [unclear: particularisant], on trouve ainsi $8$ fonctions sur $\widehat{\Gamma}_{0,3}$ : à valeurs dans les espaces appropriés - [unclear: ... pour la connexion cyclotomique.]

[MARGIN: On verra [unclear: temps!] non ah ! que pour un [unclear: autom] $g$ et [unclear: par l'él...] $h_0, \xi, \lambda_0, \lambda$ seront définis de façon canon. et rigoureuses dans 2 paragraphes...]

\newpage

%% ===== Page 105 =====

Notons à ce propos les conditions équivalentes
$$
\begin{cases}
\text{i)} \quad p=1 \\
\text{ii)} \quad \gamma=0 \\
\text{iii)} \quad g = \sigma_\infty(g_0) \rho(g_0)^{-1} \\
\text{iv)} \quad h_0 = \xi \\
\text{v)} \quad \sigma_\infty(\xi) = \xi^{-1} \quad \text{("équation en $\xi$ simplifiée !")} \\
\text{vi)} \quad h_0 \rho(h_0) \rho^2(h_0) = 1
\end{cases} \leqno{(75 \text{bis})}
$$
signifiant toutes que $\chi(u)=1$. Donc la recherche des $u \in \Autext(\widehat{\mathfrak{S}}_{0,3}^+)'$ de multiplicateur $1$ ($u \in \Autext_{\text{loc}}(\widehat{\mathfrak{S}}_{0,3}^+)^{\prime \circ}$) équivaut bel et bien à celle des solutions simultanées des équations $(73)$, ou des équations paramétriques ("de cobord") $(74)$.

2) Revenons au cas général (multiplicateurs quelconques), le seul en fait qu'on sache écrire explicitement (cas $u = c$) correspond à $\tau_\infty \in \widehat{\mathfrak{S}}_{0,3}$ qui a les caractéristiques :
$$
\begin{cases}
u = \tau_\infty \ , \ \chi(u) = p = -1 \ , \ \gamma = -1 \\
g = l_1^{-1} \\
g_0 = 1 \ , \ h_0 = 1 \ , \ \xi = l_0^{-1} \\
\lambda' = 1 \\
[\lambda_0 = \tau_\infty \ , \ \lambda = l_0^{-1} \tau_\infty] \quad \leftarrow \text{cf. ci-dessous}
\end{cases} \leqno{(75 \text{ter})}
$$

Ce serait délicieux par un beau dimanche d'écrire un paquet d'autres $u$, pour se rendre compte notamment si les foules de conditions nécessaires qu'on est amené à dégager sur les éléments provenant de $\boldsymbol{\Gamma}_{\mathbf{Q}}$, ne sont pas automatiquement satisfaites !

\newpage

%% ===== Page 106 =====

\footnote{Calculs \og par conjug. \fg~de (VII) ne serviraient peut-être plus !} \noindent (VII) Calculs dans $\operatorname{Sl}(2,\hat{\mathbf{Z}})' = \operatorname{Sl}(2,\hat{\mathbf{Z}})/\pm 1$

\noindent ($6$ cocycles sous $p$).

Ce groupe est un (tout petit !) quotient de $\operatorname{Sl}(2,\hat{\mathbf{Z}})$, qui a l'avantage de se prêter à des calculs explicites commodes, en termes de matrices. Désignons par les mêmes lettres $g, g_0, h_0, \xi, \lambda_0, \lambda$ des éléments de $\widehat{\pi}_{0,3}$ (ou de $\mathfrak{S}_{0,3}^{+\wedge}$ plus généralement) et leurs images dans $\bar{\Gamma}(2)' \subset \operatorname{Sl}(2,\hat{\mathbf{Z}})'$ (ou dans $\operatorname{Gl}(2,\hat{\mathbf{Z}})' \defeq \operatorname{Gl}(2,\hat{\mathbf{Z}})/\pm 1$). De façon générale, considérons une matrice
$$
g = \begin{pmatrix} a & b \\ c & d \end{pmatrix} \in \operatorname{Gl}(2,\hat{\mathbf{Z}}) \leqno{(76)}
$$
On a alors
$$
\begin{aligned}
p(g) &= p g p^{-1} = \begin{pmatrix} c+d & -c \\ c+d-a-b & a-c \end{pmatrix}, & \det(g) p(g^{-1}) &= \begin{pmatrix} a-c & c \\ a+b-c-d & c+d \end{pmatrix} \\
p^2(g) &= p^{-1} g p = \begin{pmatrix} d-b & a+b-c-d \\ -b & a+b \end{pmatrix}, & \det(g) p^2(g^{-1}) &= \begin{pmatrix} a+b & c+d-a-b \\ b & d-b \end{pmatrix}
\end{aligned} \leqno{(77)}
$$
$$
[g, p] = g p g^{-1} p^{-1} = \begin{pmatrix} A & B \\ C & D \end{pmatrix} \leqno{(78)}
$$
\marginpar{NB $C+D-B =$ [rayé]\\ $\delta = \det g = ad-bc$}
$$
\begin{aligned}
\delta A &= (a^2+b^2+ab) - (ac+bd+bc) \\
\delta B &= ac+bd+bc \\
\delta C &= -(c^2+d^2+cd) + (ac+bd+ad) \\
\delta D &= c^2+d^2+cd
\end{aligned}
$$
Nous écrirons l'équation
$$
g p(g) p^2(g) = 1 \leqno{(78\text{bis})}
$$
sous la forme équivalente
$$
g p(g) = p^{-1}(g^{-1}) \quad (= p^2(g^{-1})) \quad \text{dans } \operatorname{Gl}(2,\hat{\mathbf{Z}})' \leqno{(79)}
$$
où
$$
g p(g) = \begin{pmatrix} P & Q \\ R & S \end{pmatrix} \leqno{(80)}
$$

\newpage

%% ===== Page 107 =====

$$
\begin{aligned}
P &= (a+b)(c+d-b) & Q &= -a(c-b)-bc = \delta - a(c+d-b) \\
R &= (c+d)^2 - d(a+b) & S &= ad-c(c+d) = \delta - c(c+d-b) \\
  &= (c+d)(c+d-b) - \delta & &
\end{aligned} \leqno{(81)}
$$
où on a posé $\delta = ad-bc = \det g$.

Comparant (81) et l'expression (77) de $\rho^2(g^{-1})$, l'équation (79) -- qui s'écrit aussi
\marginpar{Chercher les $g$ pour l'éq. (79) se ramène à l'équation}
$$ g \rho(g) = \pm \rho^2(g^{-1}) $$
a pour solutions les solutions suivantes :
$$ g \rho(g) \rho^2(g) = \pm 1 \quad \text{i.e.} \quad g \rho(g) = \rho^{-1}(g^{-1}) \quad \text{dans } \operatorname{Gl}(2, \hat{\mathbf{Z}})' \leqno{(82)} $$
On étudiera les solutions suivant les facteurs du produit
\marginpar{NB: l'on utilisera cette étude dans le cas $\delta = 1$ (pour $g \in \SL$)}
$$ \operatorname{Gl}(2, \hat{\mathbf{Z}})' \isom \prod_{p \text{ premier}} \operatorname{Gl}(2, \mathbf{Z}_p)' $$
La relation
$$ \frac{1}{\delta} (a+b) = P = (a+b)(c+d-b) $$
donne les deux solutions $a+b=0$ ou $c+d-b = \frac{1}{\delta}$ dans $\mathbf{Z}_\ell$ (car $\mathbf{Z}_\ell$ est intègre).

1º) $(a+b) \neq 0$ (dans $\mathbf{Z}_\ell$) d'où $c+d-b = \frac{1}{\delta}$
$$ c+d-b = \frac{1}{\delta} \leqno{(83)} $$
Moyennant (83), on vérifie sur les formules $g \rho(g) = \rho^2(g^{-1})$ -- qu'il reste à vérifier...
\marginpar{La nécessité de cette formule (83) était évidente a priori.}
$$ P = \frac{1}{\delta}(a+b) \quad Q = \frac{1}{\delta}(\delta - a) \quad R = \frac{1}{\delta}(c+d) - \delta \quad S = \frac{1}{\delta}(\delta - c) $$
équivaut à

2º) $a+b = 0$ i.e. $b=-a$ i.e. $g = \begin{pmatrix} a & -a \\ c & d \end{pmatrix}$ et $\delta = a(c+d)$
$$ \begin{cases} P = 0 \\ Q = \delta + b(c+d-b) \\ R = (c+d)^2 - \delta \\ S = \delta - c(c+d+a) \end{cases} \leqno{(84)} $$

\newpage

%% ===== Page 108 =====

1°) $a+b \neq 0$ (dans $\mathbf{Z}_\ell$) d'où
$$ c+d-b = \frac{1}{\delta} \leqno{(83)} $$

Moyennant (83)\footnote{[unclear: NB les conditions $c+d-b=\frac{1}{\delta}$, $\delta^3=1$ sont surabondantes (résultent de l'équation (81) avec $a+b \neq 0$) : elles découlent de l'équation $\rho^2(g) = \rho^{-1}(g^{-1})$, qui donne aussi le cas 2° ci-dessous.]}, on vérifie que la formule (82) -- où il reste à vérifier
$$ Q = \frac{c+d}{\delta} = \left(\frac{1}{\delta^2} - \frac{a}{\delta}\right) \quad , \quad R = \frac{b}{\delta} $$
$$ S = \frac{d-b}{\delta} = \left(\frac{1}{\delta^2} - \frac{c}{\delta}\right) $$
équivaut à :
$$ \delta^3 = 1 \leqno{(84)} $$
(dont la nécessité était évidente \emph{a priori}).

2°) $a+b = 0$ (dans $\mathbf{Z}_\ell$) d'où
$$ a = -b \quad \text{i.e.} \quad g = \begin{pmatrix} -b & b \\ c & d \end{pmatrix} \leqno{(85)} $$
$$ \det g = -b(c+d) = \delta \quad \text{d'où} \quad d = -\frac{\delta}{b} - c \leqno{(86)} $$
i.e.
$$ g = \begin{pmatrix} -b & b \\ c & -\frac{\delta}{b} - c \end{pmatrix} \quad , $$
et
$$ \begin{cases} P = 0 \quad , \quad Q = \delta + b(c+d-b) = -b^2 \\ R = (c+d)^2 = \frac{\delta^2}{b^2} \quad , \quad S = \delta + c\left(\frac{\delta}{b} + b\right) \quad , \end{cases} $$
L'équation (82) s'écrit alors, [unclear: comme on a]
$$ \rho^2(g^{-1}) = \begin{pmatrix} 0 & -1/b \\ b/\delta & -\frac{1}{b}\left(\frac{\delta}{b} + b + c\right) \end{pmatrix} $$
et l'équation (82) s'écrit alors
$$ \begin{cases} \text{a)} & b^3 = \delta^3 = 1 \\ \text{b)} & c\left( (b\delta)^2 + (b\delta) + 1 \right) + \left( \delta^2 + b\delta + b \right) = 0 \end{cases} \leqno{(87)} $$

\newpage

%% ===== Page 109 =====

On distinguera encore deux cas :

a) $b\delta = 1$ i.e.\ $b = \frac{1}{\delta}$, alors l'équation (* b) devient
$$ c = -\frac{1}{3} (2\delta^2 + 1) \leqno{(88)} $$
Or $\delta \in \hat{\mathbf{Z}}^\times \implies \delta^2 \equiv 1 \ (3)$ donc $2\delta^2+1 \equiv 0 \ (3)$, donc l'équation admet une solution dans $\hat{\mathbf{Z}}$ :
correspondant à :
$$ g = \begin{pmatrix} [unclear: -\delta^5] & \delta^4 \\ [unclear: -c] & 0 \end{pmatrix} \leqno{(89)} $$
où, je le rappelle, $\delta$ est assujetti à la seule condition $\delta^3=1$.

b) $b\delta \neq 1$ donc (comme $(b\delta)^3=1$) $(b\delta)^2+b\delta+1=0$, et l'équation (* b) devient
$$ b^2 \delta + b = 0 \leqno{(90)} $$
On pose $b = \frac{1}{\delta} u$ avec $u^3=1, u \neq 1$ de sorte que $u^2+u+1=0$
et l'équation se réduit à :
[unclear: u^2 + \dots = 0]
ce qui montre que $u$ et $\delta$ sont des solutions distinctes de l'équation cyclotomique $X^2+X+1=0$. Or $u \neq 1$ signifie $u \neq \delta$ i.e.\ $b \neq \delta^2$ ce qui exige i.e.\ $b=\delta$.
$$ g = \begin{pmatrix} -\delta & \delta \\ c & [unclear: -1-c] \end{pmatrix} \leqno{(91)} $$

\newpage

%% ===== Page 110 =====

Proposition : \marginpar{344} Soit $g \in \operatorname{Gl}(2, \hat{\mathbf{Z}})$, $g = \begin{pmatrix} a & b \\ c & d \end{pmatrix}$, $a,b,c,d \in \hat{\mathbf{Z}}$, $ad-bc = \delta \in \hat{\mathbf{Z}}^*$. Pour qu'on ait
$$ g \rho(g) \rho^2(g) = 1 \leqno{(81)} $$
il faut et il suffit qu'on ait $\delta^3 = 1$, et qu'on soit dans l'un des deux cas suivants, qui s'excluent mutuellement :
\begin{enumerate}
\item[$1^\circ)$] $c+d-b = \delta'$
\item[$2^\circ)$] $g = \begin{pmatrix} 0 & \delta' \\ -\delta' & -\delta' \end{pmatrix} \quad (\text{où } \delta' = \frac{1}{\delta} = \delta^2)$
\end{enumerate}

De plus, dans le cas $1^\circ)$, on a $a+b=0$ [unclear: ssi]

Corollaire : Les $g$ pour lesquels on a $1^\circ)$ et $a+b \neq 0$ sont ceux de la forme
$$ g = \begin{pmatrix} a & b \\ c & b-c+\frac{1}{\delta} \end{pmatrix} $$
avec $c = \frac{1}{a+b} ( a(b+\delta') - \delta )$,

ceux pour lesquels on a $1^\circ)$ et $a+b=0$ sont ceux de la forme
$$ g = \begin{pmatrix} -u\delta' & u\delta' \\ c & -\delta'u' - c \end{pmatrix} $$
où $\delta' = \frac{1}{\delta} = \delta^2$, $u' = \frac{1}{u} = u^2$, $u$ étant une solution de l'équation cyclotomique $u^2+u+1=0$ i.e.\ racine primitive $3^e$ de l'unité.

\newpage

%% ===== Page 111 =====

et on a en fait alors
$$ g\rho(g)\rho^2(g) = g^3 = \begin{pmatrix} -1 & 0 \\ 0 & -1 \end{pmatrix} = -1 $$

Remarques

1) Evidemment $-\rho$ est une solution de (87)
(tout comme $\rho^{-1}$) (tout comme $-\rho^{-1}$), mais ils sont justiciables
respectivement des cas 1).

2) Pour déterminer les solutions dans $\operatorname{Gl}(2, \hat{\mathbf{Z}})$
de $g\rho(g)\rho^2(g)=1$, il suffit de choisir pour
tout $\ell$ une telle solution dans $\operatorname{Gl}(2, \mathbf{Z}_\ell)$, en
utilisant les résultats (86) ; on notera que
les "$a$" obtenus 1) à 3) peuvent varier d'un
$\ell$ à un autre.

On fera attention que les cas 1) 2) 3)
correspondent resp. à des familles de
solutions dépendant de 2, 1 et 0 paramètres.

3) J'ignore si une $\bar{g} \in \operatorname{Sl}(2, \hat{\mathbf{Z}})$, $\bar{g} = \left(\begin{smallmatrix} A & B \\ C & D \end{smallmatrix}\right)$ satisfaisant :
$$ \bar{g} \rho(\bar{g}) \rho^2(\bar{g}) = 1 $$
admet une représentation
\marginpar{[unclear: NB $g \mapsto \lambda_0^{-1} g \lambda_0$ conserve les relations du centre, l'équation $g' = -\lambda_0^{-1} g \lambda_0$]}paramétrique de la forme (88)
$(= \lambda_0 ( - \delta' ) \lambda_0^{-1} \dots)$
(il suffit
de le prouver dans le $\operatorname{Sl}(2, \mathbf{Z}_\ell)$).
Je l'ignore même [unclear: en supposant] $g \equiv \mathrm{id} \ (2)$
(ce qui est une condition sur le $1$-cocycle
dans $\operatorname{Sl}(2, \mathbf{Z}_2)$ seulement). On voit en

\newpage

%% ===== Page 112 =====

tous cas sur les formules (78) que pour $\xi$ satisfaisant la condition $\xi\rho(\xi)\rho^2(\xi)=1$, de la forme $\lambda_0\rho(\lambda_0)^{-1}$ ([unclear: on verra bien plus loin]), on aura $C+D-B=1$ --- i.e.\ on est dans le cas ``général'' à deux paramètres.

\marginpar{OK, cf.\ (98)\\ ci-dessous}Si on prend
$$
g = \lambda_0 \rho (\lambda_0)^{-1},
$$
on sera dans le cas B), avec $C+D-B=-1$
i.e.\ le cas $1^\circ$) encore à deux paramètres. Je n'ai pas eu le temps de l'éclaircir -
(sans pouvoir à l'avance exclure le cas où on aurait en même temps
$A+B=0$, ce qui en fait nous ramènerait au cas
$$
\xi = \begin{pmatrix} A & B \\ C & D \end{pmatrix}
$$
i.e.\ le cas à un paramètre
$$
\xi = \begin{pmatrix} -u & u \\ C & -u-C \end{pmatrix}
$$
avec $u$ racine primitive $3^{\text{ème}}$ de $1$,
$u' = u^{-1} = u^2$ l'autre racine primitive.)

\newpage

%% ===== Page 113 =====

(VIII) --- Calculs dans $\operatorname{Sl}(2, \hat{\mathbf{Z}})$ : cocycles sous $\sigma_{\infty}$.

Reprenons $g = \begin{pmatrix} a & b \\ c & d \end{pmatrix} \in \operatorname{Sl}(2, \hat{\mathbf{Z}})$,
et calculons $\sigma_{\infty} g \sigma_{\infty}^{-1}$, on trouve, grâce à
$$
\sigma_{\infty} = \begin{pmatrix} 0 & -1 \\ 1 & 0 \end{pmatrix} \leqno{(48)}
$$
\footnote{NB On a $\sigma_{\infty}^{-1} = \begin{pmatrix} 0 & 1 \\ -1 & 0 \end{pmatrix} = -\sigma_{\infty}$ et $\sigma_{\infty}^2 = \begin{pmatrix} -1 & 0 \\ 0 & -1 \end{pmatrix} \neq 1$ dans $\operatorname{Sl}(2, \hat{\mathbf{Z}})$ !}les formules\footnote{NB Si $g \in \operatorname{Gl}(2, A)$, $A$ anneau quelconque (ici $A=\hat{\mathbf{Z}}$), on a $\sigma_{\infty} g \sigma_{\infty}^{-1} = (\det g) {}^t(g^{-1})$. $\sigma_{\infty} g \sigma_{\infty}^{-1} = \begin{pmatrix} d & -c \\ -b & a \end{pmatrix}$}
$$
\sigma_{\infty} g \sigma_{\infty}^{-1} = \begin{pmatrix} d & -c \\ -b & a \end{pmatrix} = {}^t(g^{-1})
$$
la contragrédiente de $g$ !
i.e.
$$
\sigma_{\infty} g^{-1} \sigma_{\infty}^{-1} = {}^t g \leqno{(90)}
$$
$$
[\sigma_{\infty}, g^{-1}] = ({}^t g) (g) = \begin{pmatrix} a^2+b^2 & ac+bd \\ ac+bd & c^2+d^2 \end{pmatrix}
$$
$$
= \sigma_{\infty} g^{-1} \sigma_{\infty}^{-1} g
$$
donc, en prenant l'inverse des deux membres dans la dernière formule
$$
g^{-1} \sigma_{\infty} (g) = \begin{pmatrix} c^2+d^2 & -(ac+bd) \\ -(ac+bd) & a^2+b^2 \end{pmatrix} \leqno{(91)}
$$
$$
= ({}^t g \cdot g)^{-1}
$$

D'autre part, soit $h_0 \in \operatorname{Sl}(2, \hat{\mathbf{Z}})$
on a
$$
h_0 \sigma_{\infty}(h_0) = h_0 ({}^t h_0^{-1})
$$
donc
$$
h_0 \sigma_{\infty}(h_0) = 1 \iff h_0 = {}^t h_0 \leqno{(*)}
$$

Soit enfin
$$
\xi = \begin{pmatrix} A & B \\ C & D \end{pmatrix} \leqno{(92)}
$$
et considérons l'équation (68)
$$
\sigma_{\infty}(\xi) = l_1 \xi^{-1} l_0 \leqno{(**)}
$$
qui équivaut aussi à (*) où $h_0$ donné par (67) pour $l_0, l_1$ (49)

\newpage

%% ===== Page 114 =====

$$
\begin{aligned}
h_0 = l_0^{-\gamma} \xi &= \begin{pmatrix} 1 & + 2\gamma \\ 0 & 1 \end{pmatrix} \begin{pmatrix} A & B \\ C & D \end{pmatrix} \\
&= \begin{pmatrix} A + 2\gamma C & B + 2\gamma D \\ C & D \end{pmatrix}
\end{aligned} \leqno{(92)}
$$
donc "l'équation en $\xi$" $(**)$ (décrite dans la
question $\operatorname{Sl}(2, \hat{\mathbf{Z}})'$ dans $\hat{\mathfrak{S}}_{0,3}^+$), s'écrit $B + 2\gamma D = C$
i.e.
\marginpar{où on a rappelé que $2\gamma + 1 = p$}
$$ C - B = 2\gamma D \quad \text{i.e.} \quad p D = 1 \leqno{(93)} $$

(On fera attention que l'équation $(**)$ \underline{dans}
$\operatorname{Sl}(2, \hat{\mathbf{Z}})'$ signifie $h_0 \sigma_\infty(h_0) = \pm 1$ i.e.\ $h_0 = \pm {^t h_0}$,
mais le signe $-$ impliquerait $D = -D$ i.e.\ $2D = 0$ donc $D = 0$, ce qui est impossible
si on suppose que $\xi = \begin{pmatrix} A & B \\ C & D \end{pmatrix}$ provient
de $\hat{\mathfrak{T}}_{0,3}$ donc est dans $\overline{\Gamma}'(2)$ i.e.\
$\equiv 1 \pmod 2$ (de sorte que $D \equiv 1 \pmod 2$).
Donc on devra bien avoir (93).

Prenons maintenant la représentation
paramétrique (70) de $\xi$ (tenant compte de
l'équation $\xi \sigma_\infty {^t \xi} \sigma_\infty^{-1} = 1$ (64))
$$ \xi = \lambda {^t \lambda}^{-1} \quad (\lambda = \lambda' \sigma_\infty, \lambda' \in \operatorname{Sl}(2, \hat{\mathbf{Z}})) \leqno{(94)} $$
(en fait $\lambda \equiv 1 \pmod 2$ dans le cas
qui nous intéresse, mais peu importe...)

\newpage

%% ===== Page 115 =====

alors la formule $(78)$ donne $A, B, C, D$ comme
fonctions quadratiques en $a,b,c,d$ (en faisant
attention que $\delta = -1$ dans le cas 2))
$$
\begin{cases} B+C = \text{[raturé]} D-1 \\ D = \delta (c^2+d^2+cd) \end{cases} \leqno (95)
$$
donc $(93)$ s'écrit
$$
\text{[raturé]} = 2\gamma D \quad \text{i.e.}
$$
$$
(2\gamma+1)D = 1 \quad \text{i.e.} \quad pD = 1 \quad | \quad D = \delta (c^2+d^2+cd) \leqno (96)
$$
i.e. le multiplicateur $p$ s'obtient en fonction
de $\xi \in \widehat{\Pi}_{0,3}$ en prenant l'image
de $\xi$ dans $\operatorname{Sl}(2, \hat{\mathbf{Z}})'$, sous forme $\begin{pmatrix} A & B \\ C & D \end{pmatrix}$
(il y a une ambiguité de signe, levée par le
choix de $\lambda$ tel que $\xi = \lambda \rho(\lambda)^{-1}$, [unclear: l'image]
de $\lambda$ dans $\operatorname{Gl}(2, \hat{\mathbf{Z}})'$ [unclear: détermine une
matrice f. unité] de $\operatorname{Gl}(2, \hat{\mathbf{Z}})$ qu'au
signe près, mais $\lambda \rho(\lambda)^{-1}$ n'en dépend plus du
choix de ce signe \dots), par la formule
$$
p = \frac{1}{D} \qquad \left( D = \delta (c^2+d^2+cd) \right) \leqno (97)
$$
Notons maintenant qu'en réduisant $\mod 3$,
la forme quadratique $c^2+d^2+cd = (c+d)^2 - cd$
ne prend que les valeurs $0$ et $1$ (clair si $c=0$,
si $d=0$, \text{[raturé]} et dans le troisième cas, sinon
$c, d = \pm 1$, et $f(c,d) = 1+1 \pm 1 \equiv 1 \text{ ou } 0 \pmod 3$. Or ici
[unclear: $p$ doit trouver une unité dans $\hat{\mathbf{Z}}$, le \dots à cacher]

\newpage

%% ===== Page 116 =====

\marginpar{NB On a $D \equiv p \equiv \delta \pmod 3$ (la cong. $D \equiv \delta$ est d'ailleurs démontrée ci-contre, puisqu'on a supposé ici que $\xi = \lambda \rho \lambda^{-1}$, $\lambda = \lambda_{t_0}$, $\det \lambda = 1 = \delta$ \\ NB $C+D-B = -1$}

$$
p \equiv \delta \pmod 3 \leqno{(98)}
$$
(car $\delta = \det \lambda \in \{\pm 1\}$, $\xi = \lambda \rho (\lambda^{-1})$)
$$
\begin{cases}
\text{Cas } \xi = \lambda \rho \lambda^{-1} \ (\lambda \in \operatorname{Sl}(2, \hat{\mathbf{Z}})) \iff p \equiv 1 \pmod 3 \iff \gamma \equiv 0 \pmod 3 \\
\text{Cas } \xi = \lambda' l_1 \rho (l_1)^{-1} \ (\lambda' \in \operatorname{Sl}(2, \hat{\mathbf{Z}})) \iff p \equiv -1 \pmod 3 \iff \gamma \equiv -1 \pmod 3
\end{cases}
$$

B) Cas $\xi = \lambda' l_1 \rho (\lambda')^{-1} \ (= (\lambda' l_1) \rho (\lambda' l_1)^{-1})$.

Utilisons la deuxième formule (87) -- on trouve
$$
\xi = \begin{pmatrix} a & b \\ c & d \end{pmatrix} \begin{pmatrix} 1 & 2 \\ 0 & 1 \end{pmatrix} \begin{pmatrix} a-c & c \\ a+b-c-d & c+d \end{pmatrix} = \begin{pmatrix} A & B \\ C & D \end{pmatrix}
$$
avec
$$
\begin{cases}
A = 3a^2 + b^2 + 3ab - 3ac - 2ad - bc - bd \\
C = - (3c^2 + 3cd + d^2) + 3ac + ad + 2bc + bd
\end{cases}
\quad
\begin{cases}
B = 3ac + 2ad + bc + bd \\
D = 3c^2 + 3cd + d^2
\end{cases} \leqno{(99)}
$$

\marginpar{on aurait pu s'en convaincre sur l'équation de trace $\operatorname{Tr} \xi = 0$ \\ si on suppose $p \equiv -1$ (i.e. $\gamma \equiv -1$) \\ pour $D=1 \pmod 2$.}

Comme tantôt la formule (93) (qui exprime "l'égalité à un signe près" (91), ou celle en $th.$ de [unclear: as]) s'écrit, compte tenu de la relation
$$
B - C = -1 + D
$$
sous la forme
$$
2\gamma D = -1 - D \quad \text{i.e.}
$$
$$
(2\gamma+1)D = -1 \quad \text{i.e.} \quad p D = -1 \leqno{(100)}
$$

donc ici le multiplicateur $p$ s'exprime en fonction de l'image de $\xi$ dans $\operatorname{Sl}(2, \hat{\mathbf{Z}})'$, sous sa forme $\begin{pmatrix} A & B \\ C & D \end{pmatrix}$ où $A,B,C,D \in \hat{\mathbf{Z}}$ -- et on lève l'ambiguïté de signe comme tantôt par la formule
$$
p = - \frac{1}{D} \qquad (D = 3c^2 + 3cd + d^2) \leqno{(101)}
$$

\newpage

%% ===== Page 117 =====

Ici il est clair que les seules valeurs mod $3$
représentées par la forme quadratique
$3c^2 + 3cd + d^2 \equiv d^2 \pmod{3}$ sont $0$ et $1$, d'où
$$
p \equiv -1 \pmod{3} \leqno{(102)}
$$
(i.e.\ $\gamma \equiv -1 \pmod{3}$) si $\xi = \lambda \tau_0^{-1} \rho(\lambda)^{-1}$
($\Leftrightarrow \xi = \lambda_0 \tau_0^{-1} \rho(\lambda_0)^{-1}$)

Cette étude montre donc bien que l'on est
dans le cas $\xi = \lambda \rho(\lambda)^{-1}$, resp.\ $\xi = \lambda \tau_0^{-1} \rho(\lambda)^{-1}$,
suivant que $p = \chi(u) \equiv 1$ ou $\equiv -1 \pmod{3}$,
i.e.\ suivant que $\gamma \equiv 0$ ou $\gamma \equiv -1 \pmod{3}$.

Remarques. 1) Se plaçant dans $\operatorname{Sl}(2,\hat{\mathbf{Z}})'$ (ou plutôt, dans son
image $\overline{\Gamma}$ sous-groupe $\Pi(2)' = \operatorname{Ker}(\operatorname{Sl}(2,\hat{\mathbf{Z}})' \to \operatorname{Sl}(2,\mathbf{Z}/2\mathbf{Z}))$
au lieu de $\hat{\Pi}_{0,3}$), pour résoudre l'équa-
solutions de « l'équation en $\xi$ » $(**)$ (cf.\ $(68_3)$),
quand $\xi$ est mis sous forme paramétrique
(94), on voit qu'en tant que condition sur
$\lambda = \begin{pmatrix} a & b \\ c & d \end{pmatrix}$, cette équation est relativement
anodine et très simple - elle dit
simplement que $p D = 1$, où
$D = c^2+d^2+cd$ dans le premier cas,
$D = 3c^2+3cd+d^2$ dans le deuxième. Si on
ne fixe pas $p \in \hat{\mathbf{Z}}^*$ d'avance, cette
condition dit simplement que l'élément

\newpage

%% ===== Page 118 =====

$D$ de $\hat{\mathbf{Z}}$ soit une unité, i.e.
$$
D \in \hat{\mathbf{Z}}^* \leqno{(103)}
$$
\marginpar{NB On a $\rho$}et on détermine alors de façon unique en termes de $D$, donc de $c$ et $d$ et du signe $\delta$ par $D = \delta(c^2 + d^2 + cd)$ ($a$ et $b$ s'en calculant) -- donc aussi en termes de $\lambda \in \widehat{\Pi}_{0,3}^{\sim, (1, 2)}$ ; $\rho$ lui-même est uniquement déterminé par $\xi$.

Ainsi la connaissance de $\xi \in \widehat{\Pi}_{0,3}$ suffit à déterminer le multiplicateur $p = 2\gamma + 1 \in \hat{\mathbf{Z}}^*$, donc aussi $\gamma \in \hat{\mathbf{Z}}$, et $h_0 = l_0^{-\gamma} \xi$, donc aussi $g_0$ (par la formule (66), $g_0^{-1} \sigma_a g_0 = h_0$), et par suite $g$ par la formule (62) $g = \sigma_a (g_0) l_1^\gamma \rho(g_0)^{-1}$, et $u$ par (50). Il est intéressant de s'assurer que cela ne change pas quand on fait le changement (69)
$$
\xi \longmapsto \xi' = l_0^{-\alpha} \xi l_1^\alpha
$$
en termes matriciels, on a en effet
$$
l_0^{-\alpha} \xi l_1^\alpha = \begin{pmatrix} 1 & 2\alpha \\ 0 & 1 \end{pmatrix} \begin{pmatrix} A & B \\ C & D \end{pmatrix} \begin{pmatrix} 1 & 0 \\ 2\alpha & 1 \end{pmatrix} = \begin{pmatrix} A + 2\alpha(B + C) + 4\alpha^2 D & B + 2\alpha D \\ C + 2\alpha D & D \end{pmatrix} \leqno{(104)}
$$

2) Mais on fera attention qu'à toute rigueur la seule connaissance de l'image de $\xi$ dans $\operatorname{Sl}(2, \hat{\mathbf{Z}})'$ ne permet de déterminer le multiplicateur $p$ qu'au signe près -- par contre la connaissance de l'image de $\lambda$ dans $\operatorname{Gl}(2, \hat{\mathbf{Z}})'$ c-à-d d'une matrice $\begin{pmatrix} a & b \\ c & d \end{pmatrix} \in \operatorname{Sl}(2, \hat{\mathbf{Z}})$ définie au signe près -- permet de déterminer

\newpage

%% ===== Page 119 =====

$p$ comme l'inverse de $D$, donné [unclear: mieux]
par $D = \delta(c^2 + cd + d^2)$ où $\delta$ est déterminé par la
condition $D \equiv \delta \quad (3)$.

3) On sait par voie arithmétique (comme
$$
\Gamma_{\mathbf{Q}} \to \hat{\mathbf{Z}}^* \quad \text{est surjectif)}
$$
qu'on doit trouver n'importe quel $p \in \hat{\mathbf{Z}}^*$ sous
la forme $1/D$ (ou encore, sous la forme $D$),
i.e. $D = c^2 + cd + d^2$ si $p \equiv 1 \quad (3)$, et
$D = -3(c^2 + cd) + d^2$ si $p \equiv -1 \quad (3)$, avec $c, d \in \hat{\mathbf{Z}}$,
et même $c \equiv 0 \quad (2)$, $d \equiv 1 \quad (2)$. (On est ramené à le
vérifier séparément pour les différents $\mathbf{Z}_{\ell}$, i.e.
à prouver le

\underline{Corollaire} : Si $\ell$ est premier $\neq 3$, alors l'-
él. $p \in \mathbf{Z}_{\ell}^*$ s'écrit indifféremment sous les formes
$c^2 + cd + d^2$, ou $-3(c^2 + cd) + d^2$\footnote{m'en fiche de $-3(c^2+cd)+d^2$ ! $p$} (avec $c,d \in \mathbf{Z}_{\ell}$). Si
$\ell = 2$, on peut supposer $c \equiv 0$, $d \equiv 1 \quad (2)$. Pour
$p \in \mathbf{Z}_3^*$, $\ell = 3$, il se met sous la forme
$c^2 + cd + d^2$ (resp. $-3(c^2 + cd) + d^2$) ssi $p \equiv 1 \quad (3)$
(resp. $p \equiv -1 \quad (3)$), ce qui est mutuellement
exclusif.

4) Revenant à la réduction "anodine" dans
$\operatorname{Sl}(2, \hat{\mathbf{Z}})'$ de "l'équation en $\lambda$" comme
signifiant simplement que $c^2 + cd + d^2 \in \hat{\mathbf{Z}}^*$
(resp. $-3(c^2 + cd) + d^2 \in \hat{\mathbf{Z}}^*$), on pourrait espérer

\newpage

%% ===== Page 120 =====

[unclear: trouver] une méthode pour construire [unclear: au moins] dans
$\hat{\Pi}_{0,3}$ des gros paquets de solutions de l'équa-
tion en question. Or je n'y suis absolu-
ment pas arrivé $-$ pas [unclear: à même] d'en construire
explicitement une seule, en dehors de celles
qui correspondent à $u=1$, et $u=t_0$.

Une question : Soit $\lambda_0 \in \operatorname{Sl}(2, \hat{\mathbf{Z}})'$, $\lambda_0 \equiv 1$ (2), tel que
$\xi_0 = \lambda_0 \rho(\lambda_0)^{-1}$ (disons, $-$ $\xi = \lambda_0(-\rho(\lambda_0)^{-1})$) soit tel qu'il
satisfasse dans $\operatorname{Sl}(2, \hat{\mathbf{Z}})'$ à une "équation en $\xi$"
pour $\gamma = p^{-1}$ convenable ($p \in \hat{\mathbf{Z}}^\times$) $-$ i.e.
$c^2 + d^2 + cd \in \hat{\mathbf{Z}}^\times$ (resp. $3(c^2+cd)+d^2 \in \hat{\mathbf{Z}}^\times$).
Existe-t-il alors un $\tilde{\lambda} \in \hat{\Pi}_{0,3}$ au-dessus de
$\lambda_0$, tel que $\tilde{\xi} = \tilde{\lambda} \rho(\tilde{\lambda})^{-1}$ (resp. ...) satis-
fasse aussi une équation en $\tilde{\xi}$ (nécessairement
avec le même $\gamma$) ?

\newpage

%% ===== Page 121 =====

\section*{43. Retour sur les [unclear: relations] --- Digression sur les $1$-cocycles des groupes cycliques}

Soit $\pi$ un groupe, $\rho$ un automorphisme de $\pi$ tel que
$$
\rho^n = \operatorname{id} \leqno{(1)}
$$
où $n \in \mathbf{N}^*$. On a donc une opération du groupe cyclique $\mathbf{Z}/n\mathbf{Z}$ sur $\pi$, d'où un produit semi-direct
$$
\tilde{\pi} = \pi \cdot \mathbf{Z}/n\mathbf{Z} \leqno{(2)}
$$
et une structure d'extension
$$
1 \to \pi \to \tilde{\pi} \to \mathbf{Z}/n\mathbf{Z} \to 1 \leqno{(3)}
$$
On identifie $\rho$ avec le générateur de $\mathbf{Z}/n\mathbf{Z}$, et on [unclear: l'identifie] à un élément de $\tilde{\pi}$, par le scindage canonique
$$
\mathbf{Z}/n\mathbf{Z} \longrightarrow \tilde{\pi} \qquad (1 \bmod n \longmapsto \rho) \leqno{(4)}
$$
On a donc
$$
\rho(g) = \rho g \rho^{-1} \qquad (g \in \pi) \leqno{(5)}
$$
Les scindages de l'extension (3), correspondent biunivoquement aux éléments $\rho'$ de $\tilde{\pi}$ sur $\rho$, tels que
$$
\rho'^n = 1 \leqno{(6)}
$$
Dire que $\rho'$ est sur $\rho$ signifie que l'on a
$$
\rho' = g \rho \qquad \text{avec } g \in \pi , \leqno{(7)}
$$
et en termes d'un tel $g \in \pi$, l'équation (6) s'écrit aussi
$$
g \rho(g) \rho^2(g) \cdots \rho^{n-1}(g) = 1 \leqno{(8)}
$$

\newpage

%% ===== Page 122 =====

Le groupe $\pi$ opère par conjugaison sur l'ensemble
des scindages de l'extension (3), c'est-à-dire,
ce qui revient au même, sur l'ensemble des $\rho'$
au dessus de $\rho$, satisfaisant (6). On a
d'ailleurs, pour $\alpha \in \pi$
$$
\alpha \rho'_g \alpha^{-1} = \alpha g \rho \alpha^{-1} = \alpha g \rho \alpha^{-1} \rho^{-1} \rho = (\alpha g \rho(\alpha)^{-1}) \rho = \rho'_{\alpha g \rho(\alpha)^{-1}} \leqno{(9)}
$$
i.e.\ l'opération par conjugaison de $\pi$ sur l'ensemble
des scindages de (3), s'interprète comme
l'opération de $\pi$ sur l'ensemble des solutions de (8)
donnée (pour $\alpha \in \pi$) par
$$
g \mapsto \alpha * g \defeq \alpha g \rho(\alpha)^{-1} \leqno{(10)}
$$
(on dira que $\alpha g \rho(\alpha)^{-1}$ est le $\rho$-conjugué de
$g$ par l'élément $\alpha \in \pi$). Les classes de
$\pi$-conjugaison de scindages de (3) correspon-
dent donc aux classes de $\pi$-$\rho$-conjugaison
de solutions de (8).

Parmi ces classes, la classe "triviale" (du scindage
canonique de (3), i.e.\ la classe de $g=1$), est
constituée exactement par l'ensemble des
$$
g = \alpha \rho(\alpha)^{-1} \quad (\alpha \in \pi) \leqno{(11)}
$$

Soient $g, g'$ deux solutions congruentes de

\newpage

%% ===== Page 123 =====

l'équation (8), de sorte que
$$
g' = \alpha g \rho(\alpha)^{-1} \quad \text{i.e.} \quad g' = \theta(\alpha).g \leqno{(12)}
$$
pour un $\alpha \in \pi$ convenable. Quand est-on assuré de l'unicité de $\alpha$ --- plus généralement, quel degré d'incertitude, pour $g, g'$ fixés, de l'ens. des $\alpha \in \pi$ qui tels qu'on ait (12) (équiv. pour « relever » la $\rho$-transgression de $g$ dans $g'$, dans le groupe $\pi$). C'est évidemment une classe à droite modulo le sous-groupe $\pi(\rho,g)$ de $\pi$ qui stabilise $g$ :
$$
\begin{aligned}
\pi(\rho,g) &= \{ \beta \in \pi \mid \theta_\rho(\beta).g = g \} \\
&= \{ \beta \in \pi \mid \beta g \rho(\beta)^{-1} = g \} \\
&= \{ \beta \in \pi \mid \operatorname{int}(\beta)(g\rho) = g\rho \} \\
&= \pi^{g\rho} = \{ \beta \in \pi \mid \operatorname{int}(g\rho)(\beta) = \beta \} \\
&= \text{Centralisateur de } g\rho \text{ dans } \pi
\end{aligned} \leqno{(13)}
$$
En d'autres termes, si $\alpha$ satisfait (12), l'ens. des $\alpha'$ qui satisfont la même relation est l'ens. $\alpha \pi(\rho,g)$ :
$$
\begin{gathered}
\alpha' g \rho(\alpha')^{-1} = \alpha g \rho(\alpha)^{-1} \iff \alpha' \in \alpha \pi(\rho,g) \\
\text{i.e. } \alpha^{-1}\alpha' \in \pi^{g\rho}
\end{gathered} \leqno{(14)}
$$
Donc l'unicité de $\alpha$ satisfaisant (12) $\iff$
$$
\pi^{g\rho} = 1 \leqno{(15)}
$$
condition qui ne dépend que de la classe de $\pi$-conjugaison du sous-groupe envisagé de (3) défini par $g$, i.e. $g\rho$.

\newpage

%% ===== Page 124 =====

\begin{itemize}
\item[ ]
En fait, non seulement le groupe $\pi$, mais le groupe
$\widetilde{\pi}$ lui-même opère par conjugaison sur l'ensemble des
scindages de l'extension (3). Il en est en
particulier (de $p$) de $\rho$, et on a :
$$ ^\rho(g \rho) = \rho g \rho \rho^{-1} = \rho(g\rho)\rho^{-1} = (\rho g \rho^{-1}) . \rho = (^\rho g) \rho \leqno (16) $$
i.e. l'opération de $\rho$ sur l'ensemble des scindages
en question, s'identifie à l'opération induite
par $\rho$ (par $g \mapsto ^\rho g$) sur l'ensemble des
solutions de (8) ($g \in \pi$). Notons d'ailleurs
que $g$ et $p(g)$ sont $\pi$-équivalents, car
$$ p(g) = \alpha g p(\alpha)^{-1} \quad \text{avec} \quad \alpha = g^{-1} \leqno (17) $$
Notons que pour l'opération de $\widetilde{\pi}$ sur l'ensemble
des scindages de (3), ou encore sur l'ensemble
des solutions de (8), est donnée par
$$ \Theta_\alpha : g \mapsto \alpha g p(\alpha)^{-1} \leqno (14) $$
où maintenant $\alpha \in \widetilde{\pi}$ (et plus
nécessairement $\alpha \in \pi$). Les classes de $\widetilde{\pi}$-conjugaison
(comme il résulte
en se réduisant aux classes de $\pi$-conjugaison) des
solutions de (8). En particulier, si $\alpha \in \widetilde{\pi}$, l'élément
$$ \alpha p(\alpha)^{-1} \quad (\alpha \in \widetilde{\pi}) \leqno (18) $$
de $\widetilde{\pi}$ est dans $\pi$, et il en est ainsi de (8) si
\end{itemize}

\newpage

%% ===== Page 125 =====

se trouve dans les deux [unclear: classes principales] \dots

$$
\alpha \rho(\alpha)^{-1} = \alpha_0 \rho(\alpha_0)^{-1} \qquad \text{pour } \alpha = \alpha_0 \pi \text{ avec } \alpha_0 \in \tilde{\pi}, \pi \in Z \leqno{(19)}
$$

En fait, on peut écrire $\alpha = \alpha_0 p^i$ (d'où $\rho(\alpha)^{-1} = p^{-i} \rho(\alpha_0)^{-1}$), d'où (19). Plus généralement, on aura, si
$$
\alpha = \alpha_0 p^i \qquad (\text{avec } \alpha_0 \in \tilde{\pi}, i \in \mathbf{Z})
$$
la relation

[unclear: $\alpha g \rho(\alpha)^{-1} = \alpha_0 p^i g \rho(p)^{-i} \rho(\alpha_0)^{-1}$]

avec $\qquad p^i g \rho(p^{-i}) = p^i g (\rho(p)^i)^{-1} \qquad \beta = (g \rho(g) \cdots \rho^{i-1}(g))^{-1}$

$$
\begin{cases}
\alpha g \rho(\alpha)^{-1} = (\alpha_0 p^i) g \rho(\alpha_0 p^i)^{-1} \\
p^i = (g \rho(g) \cdots \rho^{i-1}(g))^{-1}
\end{cases} \leqno{(19)}
$$
(pour $\alpha = \alpha_0 p^i$, $i \in \mathbf{N}$).

De même, la représentation (14) d'une solution de l'équation (8), [unclear: fournie] par un élément $\alpha \in \tilde{\pi}$, n'est plus [unclear: univoque] de nouveau, puisque pour des $p$ [unclear: fixés], la [unclear: condition] \dots $p = \rho(p)^{-1} = 1$.

Supposons plus généralement qu'on dispose d'un groupe
$$
E \supset \tilde{\pi}
$$
laissant par conjugaison $\pi = \pi_{0}$ (resp. $\tilde{\pi}$) invariant, i.e. un hom. de $E \to \operatorname{Aut}(\pi_{0})$ (resp. $E \to \operatorname{Aut}(\tilde{\pi})$), dans lequel $\tilde{\pi}$ s'envoie par [unclear: $\operatorname{int}$], de sorte que $E$ opère par [unclear: automorphismes] \dots

\newpage

%% ===== Page 126 =====

sur l'ensemble des $p'$ dans $\mathcal{E}$ tels que
$$
1 \to \pi \to \mathcal{E} \to G \to 1 \leqno{(20)}
$$

\marginpar{NB l'homomorphisme est déterminé par l'image de $1 \in \mathbf{Z}/n\mathbf{Z}$}
$$
\mathbf{Z}/n\mathbf{Z} \to G, \quad 1 \bmod n \mapsto \dot{p} \in G \quad (\text{élément d'ordre } n) \leqno{(21)}
$$

On ne peut dire en général, si $p'$ est un relèvement de $\dot{p}$ dans $\mathcal{E}$ ( $=$ s'il satisfait $p'^n=1$ ), que pour tout $\alpha \in \mathcal{E}$, $\operatorname{int}(\alpha)(p')$ est encore un relèvement au dessus de $\dot{p}$ : pour qu'il en soit ainsi, il faut et il suffit que $\dot{\alpha} \in G$ centralise $\dot{p}$. Soit donc
$$
\mathcal{E}_{\dot{p}} = \{ \alpha \in \mathcal{E} \mid [\dot{\alpha}, \dot{p}] = 1 \} \leqno{(22)}
$$
C'est un sous-groupe de $\mathcal{E}$ contenant $\pi$, qui opère sur l'ensemble des $p'$ d'ordre $n$ au dessus de $\dot{p}$, et ce par les opérations
$$
\alpha p' \alpha^{-1} = p'_{\alpha p' \alpha^{-1}} \leqno{(23)}
$$
de sorte qu'on [unclear: a un homomorphisme en considérant]
$$
\begin{cases}
\Theta_{\dot{p}} : \mathcal{E}_{\dot{p}} \to \operatorname{Aut}(\pi) \\
\alpha \mapsto \Theta(\alpha) : g \mapsto \alpha g \alpha^{-1} \quad (\alpha \in \mathcal{E}_{\dot{p}})
\end{cases}
\leqno{(24)}
$$
Cette [unclear: fois], une opération $\Theta(\alpha)$ pourra permuter les classes de $\pi$-conjugaison des solutions de (8) dans $\pi$. On [unclear: reviendra] sur $\mathcal{E}_{\dot{p}}$.

\newpage

%% ===== Page 127 =====

353

une structure d'extension
$$
1 \to \pi \to \mathcal{E}_\rho \to G^\rho \to 1 \leqno{(25)}
$$
$$ G^\rho = \operatorname{Centr}_G(\rho) = \{g \in G \mid \rho g \rho^{-1} = g\} = \{g \in G \mid [\rho, g] = 1\} $$
et $\mathcal{E}_\rho$ opère sur l'ensemble des classes de $\pi$-conjugaison de scindages de (1), (ou des classes de $\pi$-conjugaison de solutions de (8)) via une opération de $G^\rho$ dessus.

Soit $\alpha \in \mathcal{E}$, considérons l'élément
$$
\alpha \rho \alpha^{-1} \rho^{-1} = \alpha (\rho \alpha \rho^{-1})^{-1} = [\alpha, \rho]
$$
Il est dans $\pi$ si et seulement si
$$
[\dot{\alpha}, \rho] = 1 \quad \text{i.e.} \quad \dot{\alpha} \in G^\rho \quad \text{i.e.} \quad \alpha \in \mathcal{E}_\rho.
$$

Pour une solution $g$ de l'équation (8), il pourra arriver qu'on ne puisse trouver aucun $\alpha$ pour la mettre sous la forme $\alpha \rho' \alpha^{-1}$ avec $\alpha \in \pi$, mais qu'on y arrive avec $\alpha \in \mathcal{E}_\rho$. Pour qu'il existe un tel $\alpha$, il faut et il suffit que la classe de $\pi$-conjugaison $\Pi_g$ de $g$ (parmi les solutions de (8)) soit dans l'orbite de la classe $\Pi_{\rho'}$ de $\rho'$ par l'action sus-définie de $G^\rho$ sur cet ensemble de classes de conjugaison ; et les $\alpha \in \mathcal{E}_\rho$ qui satisfont
$$
g = \alpha \rho' \alpha^{-1} \leqno{(26)}
$$

\newpage

%% ===== Page 128 =====

satisfait à : la relation
$$
\dot{\alpha} [1]_\rho = [g]_\rho \leqno{(27)}
$$
et pour tout $\dot{\alpha} \in G^\rho$ satisfaisant à cette relation
(27), il existe au moins un représentant $\alpha \in \mathcal{E}_\rho$
satisfaisant (26) (ce qui est pure tautologie).

À quelle condition a-t-on unicité de
$\alpha \in \mathcal{E}_\rho$ satisfaisant à (26) ? L'ensemble de ces
$\alpha$ est une classe à droite sur le sous-groupe
de $\mathcal{E}_\rho$ formé des $\alpha$ tels que
$$
\alpha \rho(\alpha)^{-1} = 1 \quad \text{i.e.} \quad \alpha = \rho(\alpha) \quad \text{i.e.} \quad [\alpha, \rho] = 1
$$
ce groupe n'est autre que
$$
\mathcal{E}^\rho = \operatorname{Centr}_{\mathcal{E}}(\rho) \subset \mathcal{E}_\rho \leqno{(28)}
$$
groupe qui contient $\rho$ (donc le sous-groupe de $\mathcal{E}$
$\simeq \mathbf{Z}/n\mathbf{Z}$ engendré par $\rho$ — il n'est donc
jamais trivial — i.e.\ on n'a jamais unicité
(sauf dans le cas trivial où $g=1$ i.e.\ $\rho=1$,
mais alors $\mathcal{E}^\rho = \mathcal{E}$ et il faudrait $\mathcal{E}=1$ ... !)

Mais si on suppose $\pi^\rho = \{1\}$ i.e.
$$
\mathcal{E}^\rho \cap \pi = \{1\} \leqno{(29)}
$$
alors
\begin{tikzcd}
\mathcal{E}^\rho \arrow[r, hook, "\text{injectif}"'] & G^\rho \subset G
\end{tikzcd}
alors les $\alpha \in \mathcal{E}_\rho$ solutions de (26) sont en [unclear: correspon-]

\newpage

%% ===== Page 129 =====

dce biunivoque avec leurs images $\alpha$ conjuguées dans $G$, i.e. avec la trans-formée de $[1]_\rho$ sous $[g]_\rho$ : si on choisit un élément de ce trans-formant, on a défini par lui-même un $\alpha$, solution de (26).

Dans le cas de $\pi = \pi_{0,3}$\footnote{$= \pi_1(\mathcal{M}_{0,3})$}, $\mathcal{E} = \mathfrak{S}_{0,3}$\footnote{$= \pi_1(U_{0,3}, \mathfrak{S}_{0,3})$}, $G = \mathfrak{S}_{0,3} \simeq \mathfrak{S}_3 \times \mathbf{Z}/2\mathbf{Z}$, avec $\rho, \tau$ comme d'accou-tumée, l'ens. des classes de $G$-conjugaison de ss-groupes de (3) i.e. des solutions de (8) s'identifie à $\mathcal{M}_{0,3}$, qui a deux éléments, les deux "pôles" $-J, J$ --- dont l'un, choisi comme point base, correspond à la classe triviale.

$G^\rho = \mathfrak{S}_\rho \simeq \mathbf{Z}/3\mathbf{Z} \times \mathbf{Z}/2\mathbf{Z}$ est le groupe engendré par $\rho$ et $\tau$, et $\rho$ opère trivialement sur l'un des deux éléments, tandis que $\tau$ les permute (i.e. $\mathbf{Z}/2\mathbf{Z}$ engendré par $\tau$ opère de façon simplement transitive). Donc si on pose
$$
\mathcal{E}' = ( \pi . \{1, \tau_\infty\} ) = \text{image inverse de } \{1, \tau\} \subset G \text{ dans } \mathcal{E} = \mathfrak{S}_{0,3} \leqno{(30)}
$$
\marginpar{[unclear: (groupe cartographique tri-valent pondéré non orienté...)]}
on voit que
$$
\mathcal{E}'^\rho = \{1\} \leqno{(31)}
$$
Conséquence immédiate de la relation $\pi^\rho = \{1\}$
et l'ens. des solutions dans $\pi_{0,3}$ de l'équation
$$
g \rho(g) \rho^2(g) = 1 \leqno{(32)}
$$

\newpage

%% ===== Page 130 =====

est en correspondance 1-1 avec l'ensemble des éléments
de $\widetilde{\Pi}_{0,3} = \mathcal{E}'$ (le groupe cartographique triangulaire pondéré non
orienté, ayant comme générateurs les éléments
$\tau_0, \tau_1, \tau_\infty$ reliés par les relations $\tau_0^2 = \tau_1^2 = \tau_\infty^2 = 1$ sans
plus, avec $l_0 = \tau_\infty \tau_1$, $l_1 = \tau_0 \tau_\infty$, $l_\infty = \tau_1 \tau_0$ \dots) Il doit
encore en être de même quand on se place
dans les groupes $\Pi_{0,3}$ et $\hat{\Pi}_{0,3} = \widetilde{\Pi}_{0,3} \cdot \{1, \tau_\infty\}$,
si on admet (donc à vérifier ultérieure-
ment)
$$
(\hat{\Pi}_{0,3})^\rho = \{1\} \leqno{(33)}
$$
et que les éléments d'ordre fini de $\mathfrak{S}_{0,3}$ sont conjugués d'éléments d'ordre fini de $\Pi_{0,3}$ \dots

Notons que l'on a :
$$
\tau_\infty \rho (\tau_\infty)^{-1} = \tau_\infty \rho \tau_\infty = l_1^{-1} , \leqno{(34)}
$$
en vertu des formules (23) du § 42 :
$$
\rho_\infty \defeq \tau_\infty \rho \tau_\infty^{-1} = l_1^{-1} \rho = \rho l_0^{-1}
$$

La situation est encore plus simple pour
l'élément $\tau_\infty$ de $\mathfrak{S}_{0,3}$, opérant sur $\Pi$ par
un automorphisme d'ordre $n=2$, car ici toute
solution de l'équation correspondante
$$
h \tau_\infty(h) = 1 \quad (h \in \Pi) \leqno{(35)}
$$
se met sous la forme évidente
$$
h = \beta \tau_\infty(\beta)^{-1} \quad (\beta \in \Pi) \leqno{(36)}
$$
et ceci de façon unique. Cela résulte
en effet de la théorie générale \dots , compte
tenu que
$$
U_{0,3}^{\tau_\infty} = \left\{ \frac{1}{2} \right\} \leqno{(37)}
$$

\newpage

%% ===== Page 131 =====

l'ens.\ des pts fixes de $U_{0,3}$ par $\sigma_{\infty}$ est
réduit à un seul élément, ce qui implique

$$
\Pi_{0,3}^{\sigma_{\infty}} = \{1\} \ .
$$

Nous admettrons même que ces faits s'étendent
aux éléments de $\hat{\Pi}_{0,3}$ -- ce qui
sera le cas si (en plus de ce qui a été
admis à propos de $\rho$) on admet provisoirement
$$
\hat{\Pi}_{0,3}^{\sigma_{\infty}} = \{1\} \ . \leqno{(38)}
$$

Notons qu'on a ici
$$
G^{\sigma_{\infty}} = \mathfrak{T}_{0,3}^{\sigma_{\infty}} = \mathbf{Z}/2\mathbf{Z} \times \mathbf{Z}/2\mathbf{Z} \ , \leqno{(39)}
$$
sous-groupe engendré par $\sigma_{\infty}, \tau_{\infty}$ , dont les
éléments sont $\sigma_{\infty}, \tau_{\infty}, \tau_{\infty}'$ - qui se
relient par $\sigma_{\infty} \tau_{\infty} = \tau_{\infty}'$ ; et on a
$$
\mathcal{E}^{\sigma_{\infty}} = \mathfrak{S}_{0,3}^{\sigma_{\infty}} \simeq \mathbf{Z}/2\mathbf{Z} \times \mathbf{Z}/2\mathbf{Z} \ , \leqno{(40)}
$$
groupe engendré par $\sigma_{\infty}, \tau_{\infty}$ (l'autre él.\
non trivial étant $\tau_{\infty}'$). Si donc on écrit
une solution $h \in \Pi$ de (35) sous la forme (36),
avec $\beta \in \mathfrak{S}_{0,3}$ modulo $\mathfrak{S}_{0,3}^{\sigma_{\infty}}$ (par ex.\ en prenant
$\beta \in \pi$), alors l'indétermination est de
la forme (40), d'où de façon plus
précise, si $\beta_0$ est une solution dans $\pi$, alors
les autres solutions sont de la forme
$$
\beta = \beta_0 \sigma_{\infty}^{\varepsilon} \tau_{\infty}^{\varepsilon'} \qquad \varepsilon, \varepsilon' \in \mathbf{Z}/2\mathbf{Z} \leqno{(41)}
$$

\newpage

%% ===== Page 132 =====

(avec $\varepsilon, \varepsilon'$ uniquement déterminés). Pour qu'on ait $\beta \in \widehat{\mathfrak{S}}_{0,3}^+$, il faut et il suffit que $\varepsilon'=0$.

\newpage

%% ===== Page 133 =====

\marginpar{Ces notations seront révisées encore une fois au \S~46}

\chapter*{\S \space 44. --- RETOUR SUR LES NOTATIONS. MISE EN ÉQUATIONS DANS $\operatorname{Sl}(2, \hat{\mathbf{Z}})$.}
\thispagestyle{empty}
\addcontentsline{toc}{section}{44. Retour sur les notations. Mise en équations dans $\operatorname{Sl}(2, \hat{\mathbf{Z}})$.}
\label{sec:44}
\section*{}

Soit $u$ un autom. de $\hat{\Pi}_{0,3}$ (— ou ce qui revient au même, un autom. de $\widehat{\mathfrak{S}}_{0,3}^+$ qui induise l'identité dans $\bar{\mathfrak{S}}_{0,3}^+ = \mathfrak{S}_3$. Alors $u$ est connu quand on connaît ses valeurs sur les deux générateurs $\sigma_{\infty}, \rho$ de $\widehat{\mathfrak{S}}_{0,3}^+$, qui seront resp. deux éléments de $\widehat{\mathfrak{S}}_{0,3}^+$ au-dessus de $\bar{\sigma}_{\infty}$ et de $\bar{\rho}$, soit
$$
u(\sigma_{\infty}) = h \sigma_{\infty} \quad , \quad u(\rho) = g \rho \quad (h, g \in \hat{\Pi}_{0,3}) \leqno{(1)}
$$
où $h, g$ sont assujettis aux conditions
$$
h \sigma_{\infty}(h) = 1 \quad , \quad g \rho(g) \rho^2(g) = 1 \leqno{(2)}
$$
exprimant respectivement les relations
$$
(h \sigma_{\infty})^2 = (g \rho)^3 = 1 \leqno{(3)}
$$

(Il y a en plus à vérifier la condition [unclear: limitative] — à savoir que (3) sont des relations génératrices de $\widehat{\mathfrak{S}}_{0,3}^+$...). Pour le moment, il n'y a pas lieu d'imposer aucune relation entre $g$ et $h$ — sauf les relations (2) sur chacun de ces éléments — puisque $\widehat{\mathfrak{S}}_{0,3}^+$ est le groupe ayant les relations génératrices
$$
\sigma_{\infty}^2 = \rho^3 = 1
$$

\newpage

%% ===== Page 134 =====

Les éléments $g$ et $h$, du point de vue
des automorphismes de $\hat{\pi}_{0,3}$, expriment les
relations de commutation dans ce groupe avec
$\rho$ et $\sigma_\infty$ respectivement :
$$ [u,\rho] = \underset{u(\rho)}{u\rho u^{-1}}\rho^{-1} = \operatorname{int}(g) \quad , \quad [u,\sigma_\infty] = u\sigma_\infty u^{-1}\sigma_\infty^{-1} = \operatorname{int}(h) \leqno{(4)} $$
--- c'est une chose remarquable que la
connaissance des éléments $g, h$ de $\hat{\pi}$
qui vérifient (4), implique déjà la
connaissance de $u$.

Écrivons maintenant la compatibilité de
$u$ avec la structure à lacets, on l'avait
formulé précédemment par une relation
de la forme
$$ u(l_0) = \operatorname{int}(g_0)(l_0^\lambda) \quad \text{ou} \quad u(l_0) = \operatorname{int}(g_0) l_0^\lambda , $$
pour $\lambda \in \hat{\mathbf{Z}}^*, g_0 \in \hat{\pi}_{0,3}$
incommode. La notation $g_0$ n'est pas très
heureuse, puisque $g$ et $g_0$ jouent
des rôles tout à fait dissemblables. L'indice
$0$ dans $g_0$ était destiné à pouvoir écrire
plus généralement
$$ u(l_i) = \operatorname{int}(g_i) l_i^\lambda \quad i \in \{0, 1, \infty\} . $$

\newpage

%% ===== Page 135 =====

Je préfère utiliser\marginpar{35 1/2} cette lettre $f_0$ - quitte à supprimer l'indice $0$ s'il n'y a pas de confusion possible. La lettre $p$ pour un multiplicateur n'est pas un choix très heureux, car par la suite je vais avoir sûrement besoin de cette lettre pour désigner une caractéristique. Je préfère la lettre $\mu$ (initiale grecque de "multiplicateur"), et la relation de lacets prendrait la forme
$$
u(\varepsilon_0) = \operatorname{int}(f_0)\varepsilon_0^\mu
$$
En fait, je préfère y remplacer $f_0$ par $f_0^{-1}$, pour obtenir par la suite des formules où la symétrie des rôles joués par $\rho$ et par $\sigma_\infty$ apparaisse plus nettement, - [unclear: moyennant] que la formule de définition de $h_0$ en termes de $g_0$ (§ 13, formule (65))
$$
h_0 = g_0^{-1} \sigma_{\infty}(g_0)
$$
devienne plutôt
$$
h_0 = g_0 \sigma_{\infty}(g_0)^{-1}
$$
[unclear: un lapsus $\xi = \lambda(\rho\lambda)^{-1}$ \dots]

Donc la "relation de lacets" prend la forme
$$
u(\varepsilon_0) = \operatorname{int}(f_0^{-1}) \varepsilon_0^\mu \qquad \text{pour } f_0 \in \hat{\Pi}_{0,3}, \mu \in \hat{\mathbf{Z}}^*, \leqno{(5)}
$$
on en déduit aussitôt

\newpage

%% ===== Page 136 =====

$$
\begin{cases}
u(\varepsilon_1) = \operatorname{int}(f_1^{-1})(\varepsilon_1^\mu), & u(h_1) = \operatorname{int}(f_1^{-1})(h_1^\mu) \\
u(\varepsilon_\infty) = \operatorname{int}(f_\infty^{-1})(\varepsilon_\infty^\mu), & u(h_\infty) = \operatorname{int}(f_\infty^{-1})(h_\infty^\mu)
\end{cases}
\leqno{(7)}
$$

$$
f_1 = \rho(f_0) g^{-1}, \quad f_\infty = \rho(f_1) g^{-1} = \rho^2(f_0) \rho(g^{-1}) g^{-1} . \leqno{(8)}
$$

Dans la suite, on écrira simplement $f$ au lieu de $f_0$, et on va transformer l'équation
$$
\operatorname{int}(f)u(\varepsilon_0) = \varepsilon_0^\mu \leqno{(5)}
$$
en explicitant le premier membre

$$
\operatorname{int}(f)u(\varepsilon_0) = f u(\sigma_\infty\rho) f^{-1} = f u(\sigma_\infty) u(\rho) f^{-1} = f h_\infty \sigma_\infty g \rho f^{-1}
$$
$$
= f h_\infty \sigma_\infty g \sigma_\infty^{-1} \underbrace{\sigma_\infty \rho}_{\varepsilon_0} f^{-1} = f h_\infty \sigma_\infty(g) \varepsilon_0 f^{-1} \varepsilon_0^{-1} \varepsilon_0
$$
i.e.
$$
\operatorname{int}(f)u(\varepsilon_0) = \left[ f h_\infty \sigma_\infty(g) \varepsilon_0(f)^{-1} \right] \varepsilon_0 = \left[ f h_\infty \sigma_\infty(g \rho(f)^{-1}) \right] \varepsilon_0 \leqno{(9)}
$$
où les trois facteurs $f, h_\infty, \sigma_\infty(g \rho(f)^{-1})$ dans $[ \quad ]$ sont $\in \hat{\Pi}_{0,3}$. Ainsi, (5) prend la forme

$$
f h_\infty \sigma_\infty(g \rho(f)^{-1}) = \varepsilon_0^{\mu-1}
$$
ce qui donne, en posant
$$
\mu-1 = 2 \nu \quad (\nu \in \hat{\mathbf{Z}}) \leqno{(10)}
$$
(on a noté donc $\nu$ ce qui était désigné par $\gamma$ précédemment) on trouve la relation
$$
f h_\infty \sigma_\infty(g \rho(f)^{-1}) = e_0^\nu
$$
ou encore, en transformant par $\sigma_\infty$
$$
\sigma_\infty(f) \sigma_\infty(h_\infty) g \rho(f)^{-1} = e_1^\nu \quad , \text{ soit}
$$

\newpage

%% ===== Page 137 =====

$$
\sigma_\infty(h) \cdot g = \sigma_\infty(f)^{-1} h_1^\nu \rho(f)^{-1} \leqno (11)
$$

Cette relation, qui relie $h$ et $g$ via $f$, peut à juste titre s'appeler la "relation [unclear: de tresses]".
On voit, par cette relation, que non seulement la connaissance de $h$ et $g$ détermine $f$ et $\mu$\footnote{$f$ déterminé mod. le seul s-g $\simeq \hat{\mathbf{Z}}$, i.e. $f$ et $\nu$.} (i.e. $f$ et $\nu$), mais aussi inversement, une fois $f$ et $\nu$ connus, $g$ et $h$ se déterminent mutuellement via (11).

Notons que pour $u$ correspondant à $g, h, f$, si $\gamma \in \hat{\Pi}_{0,3}$, l'autom. $\operatorname{int}(\gamma) u = u'$ correspond à\marginpar{Cela vaut plus généralement si $\gamma \in \hat{\mathfrak{S}}_{0,3}$, (12) valant même pour $\gamma \in \hat{\mathfrak{T}}_{0,3}$}
$$
g' = \gamma g \rho(\gamma)^{-1}, \quad h' = \gamma h \sigma_\infty(\gamma)^{-1}, \quad f' = f \gamma^{-1} \leqno (12)
$$

Les équations (2) "en $h$" et "en $g$" se résolvent paramétriquement de façon unique, par
$$
g = \alpha \rho(\alpha)^{-1}, \quad h = \beta \sigma_\infty(\beta)^{-1} \leqno (13)
$$
$$
\begin{cases} 
\alpha \in \tilde{\Pi}_{0,3} = \hat{\Pi}_{0,3} \cdot \{ 1, \varepsilon_0 \} \\ 
\beta \in \hat{\Pi}_{0,3} 
\end{cases} \leqno (14)
$$
de façon précise, la donnée des $g, h$ satisfaisant à (2) équivaut resp. à celle de $\alpha, \beta$ dans $\tilde{\Pi}_{0,3}$ et $\hat{\Pi}_{0,3}$, via les formules (13).

\newpage

%% ===== Page 138 =====

Les éléments $g$ et $h$ étant ainsi explicités en termes des lacets $\alpha$ et $\beta$, la relation (11) peut être considérée comme une relation entre $\alpha, \beta, f, \nu$

$$
\sigma_{\infty}(\beta \sigma_{\infty}(\beta)^{-1}) (\alpha \rho(\alpha)^{-1}) = \sigma_{\infty}(f)^{-1} h^{\nu} \rho(f) \leqno{(15)}
$$

où $\nu$ est déterminé de façon unique par $\alpha, \beta$, et $f$ est déterminé dans $h^{\hat{\mathbf{Z}}} \backslash \hat{\Pi}_{0,3}$. (On notera que le deuxième membre de (15) ne change pas par le remplacement de $f$ par $h^n f$ (avec $n \in \hat{\mathbf{Z}}$).

La dissymétrie principale me semble-t-il, entre $\sigma_{\infty}$ et $\rho$ dans cette théorie, c'est que l'on a toujours $\beta \in \hat{\Pi}_{0,3}$, tandis que $\alpha$ n'est pas nécessairement dans $\hat{\Pi}_{0,3}$ - il l'est, en fait, ssi $\nu \equiv 0$ (3) i.e.\ $\mu \equiv 1$ (3).

L'autom.
$$
u' = \operatorname{int}(\beta^{-1}) u
$$
sera donc congru à $u$ mod. hom. intérieurs, et correspondra à :
$$
\begin{cases}
g' = \beta^{-1} g \rho(\beta) = (\beta^{-1} \alpha) \rho(\beta^{-1} \alpha)^{-1} \\
h' = \beta^{-1} h \sigma_{\infty}(\beta) = 1 \\
f' = f \beta
\end{cases} \leqno{(16)}
$$

ainsi $u$ est congru mod. hom. intérieurs à un [unclear: automorphisme], i.e.\ [unclear: symétrique].

\newpage

%% ===== Page 139 =====

à un unique $u' \in M$ tel que $u'(\sigma_{\infty}) = \sigma_{\infty}$ i.e.\
un $u'$ qui commute à $\sigma_{\infty}$. Mais
[unclear: pour des raisons] symétriques on prendra
$$
u'' = \operatorname{int}(\alpha^{-1}) u
$$
correspondant à
$$
\begin{cases}
g'' = \alpha^{-1} g \rho(\alpha) = 1 \\
h'' = \alpha^{-1} h \sigma_{\infty}(\alpha) = (\alpha^{-1} \beta) \sigma_{\infty}(\alpha^{-1} \beta)^{-1} \\
f'' = f \alpha
\end{cases} \leqno{(17)}
$$
mais un tel $u''$ sera $\hat{\Pi}_{0,3}$-conjugué de $u'$ - il
sera conjugué à un élément de $M$ qui
commute à $\rho$ ssi $g$ est $\hat{\Pi}_{0,3}-\rho$-conjugué
à $1$ i.e.\ ssi $\alpha \in \hat{\Pi}_{0,3}$ i.e.\ ssi $\mu \equiv 1 \; (3)$
ou encore $\nu \equiv 0 \; (3)$.

Or tous les calculs des \S \S \, précédents
étaient relatifs au cas $h=1$, i.e.\ $u$
commute à $\sigma_{\infty}$ -- alors il reste les
quantités $g$ et $f$,
\marginpar{On verra que $u$ est déterminé de façon unique et d'une façon plus rigide par $f$ plus bas}
$g$ et $f$ se déterminent mutuellement par
(11) ($f$ donnant un élément de $\operatorname{Loc}(\hat{\Pi}_{0,3})$).
Pour $f$ fixé, l'équation en $g$ peut s'interpréter
comme une équation en $f$, ou mieux, en\footnote{Cette notation n'est utile que si $\beta = 1$.}\footnote{c'est essentiellement l'élément noté $h_0$ précédemment}
$$
\varphi = f \sigma_{\infty}(f)^{-1} \leqno{(18)}
$$

\newpage

%% ===== Page 140 =====

une "équation en $\varphi$"
$$
(l_0^\nu \varphi) p(l_0^\nu \varphi) p^2(l_0^\nu \varphi) = 1 , \leqno{(13)}
$$
on pourra poser
$$
\zeta = l_0^\nu \varphi \text{ i.e. } \varphi = l_0^{-\nu} \zeta \leqno{(14)}
$$
ce qui la présente comme "équation en $\zeta$"
$$
\zeta p(\zeta) p^2(\zeta) = 1 . \leqno{(15)}
$$

Ainsi, les automorphismes à lacets de $\widehat{\Pi}_{0,3}$\footnote{avec multiplicateur $\mu=2\nu+1$,} qui
commutent à $\sigma_{\infty}$ correspondent aux
couples de solutions $(\zeta, p)$ des
équations
\marginpar{[unclear: On trouve que $\zeta$ et $p$ satisfont...]}
$$
\boxed{ p \sigma_{\infty}(p) = 1 \qquad \zeta p(\zeta) p^2(\zeta) = 1 } \leqno{(16)}
$$
satisfaisant de plus la condition de
compatibilité (14) (particulièrement simple
si $\mu=1$ i.e. $\nu=0$, puisqu'alors $\varphi=\zeta$ !). On
récupère $g$ (donc $u$) en fonction
de $(\varphi,f)$ par la formule du bas de (11), qui
ici prend la forme :
$$
g = \sigma_{\infty}(f)^{-1} l_1^{-\nu} p(f) = (l_0^{2\nu} f)^{-1} (l_0^{-\nu} \varphi) p(l_0^\nu f) \leqno{(17)}
$$
$f \in \widehat{\Pi}_{0,3}$ étant déterminé de façon unique en
fonction de $\varphi$ par la formule (12).

De façon symétrique, on va s'intéresser
aux automorphismes à lacets $u$ de $\widehat{\Pi}_{0,3}$ qui commu-

\newpage

%% ===== Page 141 =====

tant à $\rho$, i.e. ceux pour lesquels $g=1$, décrits par suite en termes de $h, f$,
$$ u(\sigma_\infty) = h \sigma_\infty, \quad u(\rho) = \rho \leqno (18) $$
$$ u(\varepsilon_0) = \operatorname{int}(f^{-1})(\varepsilon_0^\nu) \leqno (19) $$
cette dernière relation prenant la forme (11),
$$ \sigma_\infty(h) = \sigma_\infty f^{-1} l_1^\nu p(f) \leqno (20) $$
qui s'écrit aussi, [unclear: en multipliant par] $\sigma_\infty$
$$ h = f^{-1} l_0^\nu (\sigma_\infty p)(f) \leqno (21) $$
"L'équation en $h$" (2) :
$$ h \sigma_\infty(h) = 1 \leqno (*) $$
s'écrit donc, comme une "équation en $f$" :
$$ h \sigma_\infty(h) = f^{-1} l_0^\nu \underbrace{(\sigma_\infty p)(f) \sigma_\infty(f^{-1})}_{\substack{\sigma_\infty(p(f) f^{-1}) \\ \psi^{-1}}} l_1^\nu p(f) $$
qui s'écrit, en posant $\psi = f p(f)^{-1}$
$$ l_0^\nu \sigma_\infty(\psi^{-1}) l_1^\nu \underbrace{p(f) f^{-1}}_{\psi^{-1}} = 1 $$
soit (en passant à l'inverse)
$$ \psi l_1^{-\nu} \underbrace{\sigma_\infty(\psi) l_0^{-\nu}}_{\sigma_\infty(\psi l_1^{-\nu})} = 1 \quad \text{soit } \eta \sigma_\infty(\eta) = 1 $$
Donc, posant
$$ \psi \defeq f p(f)^{-1} \leqno (22) $$
$$ \eta = l_0^{-\nu} \psi \quad \text{i.e.} \quad \psi = l_0^\nu \eta \leqno (23) $$
on trouve que les automorphismes à lacets $u$ de $\hat{\mathfrak{T}}_{0,3}$ qui commutent à $\rho$ correspondent aux triplets $(\eta, \psi, \nu)$, ($\eta, \psi \in \hat{\pi}_{0,3}$, $\nu \in \hat{\mathbf{Z}}$ (avec $2\nu+1 \in \hat{\mathbf{Z}}^*$)), satisfaisant

\newpage

%% ===== Page 142 =====

\marginpar{[unclear: On trouve d'ailleurs, d'où $\nu=1$ ou $\dots$]}
$$
\psi \rho(\psi) \rho^2(\psi) = 1 \, , \quad \eta \, \sigma_{\infty}(\eta) = 1 \leqno{(24)}
$$
et en plus la relation de compatibilité (23) --- particulièrement simple si $\mu=1$ (i.e.\ si $\nu=0$, cas dans leq.\ elle signifie $\psi=\eta$ ! Mais en plus de ces relations, il faut supposer que la solution $\psi$ de l'équation (24) "en $\psi$" [unclear: soit dans la feuille] principale, i.e.\ que, quand on écrit $\psi$ sous forme
$$
\psi = f \rho(f)^{-1} \leqno{(12)}
$$
avec $\displaystyle f \in \hat{\Pi}_{0,3}^{\wedge} = \hat{\Pi}_{0,3} \cdot \{1, \tau_0\}$, on ait bien $f \in \hat{\Pi}_{0,3}$ (donc que l'image de $f$ dans $\hat{\Pi}_{0,3}^{\wedge} / \hat{\Pi}_{0,3} \simeq \{1, \tau_0\} \simeq \mathbf{Z}/2$ soit nulle...). Ayant ainsi déterminé $f$ en fonction de $\psi$, on déterminera $h$ par (21), d'où $u$.

Remarques 1) Si cette condition supplémentaire sur $\psi$ n'est pas satisfaite, la formule (21) garde un sens et définit $h \in \hat{\Pi}_{0,3}^{\wedge}$, satisfaisant "l'équation en $h$" $\quad h \sigma_{\infty}(h) = 1 \quad$ donc définissant un automorphisme de $\hat{\Pi}_{0,3}^{\wedge}$, commutant à $\rho$.

On aura alors
$$
u(l_0) = \operatorname{int}(f^{-1})(l_0^\mu)
$$
ou encore, posant
$$
f = \tau_0 f_0 \quad \text{avec} \quad f_0 \in \hat{\Pi}_{0,3}
$$
$$
u(l_0) = \operatorname{int}(f_0^{-1})(\tau_0(l_0^\mu)) = \operatorname{int}(f_0^{-1})(l_0^{-\mu})
$$

\newpage

%% ===== Page 143 =====

i.e.\ $u$ est un automorphisme [unclear: de $\hat{\pi}$], de multiplicateur $\mu = -(2\nu+1)$.

2) Notons que l'indétermination de $f$ (pour $\mu$ fixé, i.e.\ $g, h$ fixés) est dans
$$ f \mapsto f' = l_0^n f \quad (n \in \hat{\mathbf{Z}}) \leqno{(25)} $$
ce qui donne un
\marginpar{NB. On a de même pour $\alpha_0 \mapsto \alpha'_0 = l_0^n \alpha_0$, $\beta_0 \mapsto \beta'_0 = l_0^n \beta_0$ (26)}
$$ \begin{cases} \varphi \mapsto \varphi' = l_0^n \varphi l_0^{-n}, & \psi \mapsto \psi' = l_0^n \psi l_0^{-n} \\ \xi \mapsto \xi' = l_0^n \xi l_0^{-n}, & \eta \mapsto \eta' = l_0^n \eta l_0^{-n} \end{cases} $$

(NB -- $l_1 = \sigma_\infty(l_0) = p(l_0)$, donc $\xi'$ et $\eta'$ sont les $p$-transformés de $\xi, \eta$ par $l_0^{-n}$, et $\varphi', \psi'$ les $\sigma_\infty$-transformés de $\varphi, \psi$ par $l_0^{-n}$ -- en particulier les équations (16) (en $\varphi$ et $\psi$) et (24) (en $\xi$ et $\eta$) sont bien invariantes par ces transformations, ainsi d'ailleurs que les relations de compatibilité (14) resp.\ (23) (réécriture de (11)) ).

3) On a les formules suivantes :
$$ \sigma_\infty(h) g = (l_0^\nu f)^{-1} \xi p(l_0^\nu f) = \sigma_\infty(f^{-1} l_0^{-\nu}) \eta (f^{-1} l_0^{-\nu})^{-1} \leqno{(27)} $$
(cas $h=1$) - expriment $g$ comme le $p$-conjugué de $\xi$ par $(l_0^\nu f)^{-1}$, et de façon symétrique dans le cas $g=1$, comme le $\sigma_\infty$-conjugué de $\eta$ par $\sigma_\infty(f^{-1} l_0^{-\nu})$. Si $h=1$, c'est la première expression qui est intéressante (indiquant l'équation en $f$), pour $g=1$ c'est la seconde, qui s'écrit précisément (aux [unclear: notations] près)
$$ h = (f^{-1} l_0^{-\nu}) (\eta^{-1}) \sigma_\infty(f^{-1} l_0^{-\nu})^{-1} \leqno{(28)} $$

\newpage

%% ===== Page 144 =====

exprimant $h$ comme $\sigma_{\infty}$-transformé de $\eta^{-1}$ par $f^{-1} l_0^{-n}$. Si on veut l'exprimer comme $\sigma_{\infty}$-transformé de $\eta$ (au lieu de $\eta^{-1}$), on écrit :
$$ y \eta^{-1} \sigma_{\infty}(y)^{-1} = y \eta^{-1} \cdot \eta \cdot \sigma_{\infty}(y \eta^{-1})^{-1} $$
d'où
$$ h = (f^{-1} l_0^{-n} \eta^{-1}) \eta \sigma_{\infty}(f^{-1} l_0^{-n} \eta^{-1})^{-1} \ . \leqno{(28\text{bis})} $$

Ces formules permettent de mettre en relation, dans le cas $h=1$, les solutions des deux équations
$$ g = \alpha p(\alpha)^{-1} \quad \text{et} \quad \xi = \alpha_0 p(\alpha_0)^{-1} \leqno{(29)} $$
par
$$ \alpha = (l_0^n f)^{-1} \alpha_0 = f^{-1} l_0^{-n} \alpha_0 \ , \leqno{(30)} $$
et dans le cas $g=1$, les solutions des deux équations
$$ h = \beta \sigma_{\infty}(\beta)^{-1} \quad , \quad \eta = \beta_0 \sigma_{\infty}(\beta_0)^{-1} \leqno{(31)} $$
par
$$ \beta = (f^{-1} l_0^{-n} \eta^{-1}) \beta_0 \ . \leqno{(32)} $$

La dissymétrie entre ces deux formules ( (32) étant un peu plus compliquée que (30) par le facteur supplémentaire $\eta^{-1}$ du coefficient de $\beta_0$ ) provient du fait que $f$ a été défini via le formalisme des Cartes (6), première formule, où $l_0$ est défini comme le carré de $\sigma_{0\infty} = \varepsilon_0$, et non de $p \sigma_{\infty} = \varepsilon_1$.

\newpage

%% ===== Page 145 =====

Si on identifie maintenant $\bar{\mathfrak{S}}_{0,3}^{\pm \wedge}$ à $\mathrm{Gl}(2, \hat{\mathbf{Z}})^\wedge$,
et qu'on regarde les images des éléments $g, h, f$ etc dans le quotient
$$ \mathrm{Gl}(2, \hat{\mathbf{Z}})^{\pm 1} = \mathrm{Gl}(2, \hat{\mathbf{Z}})^\wedge / \pm 1 $$
- qu'on désignera par les mêmes lettres, par abus de langage - on sera amené à calculer
avec des éléments de $\mathrm{Gl}(2, \hat{\mathbf{Z}})$, plutôt que de $\mathrm{Gl}(2, \hat{\mathbf{Z}})^\wedge$. On choisira $\alpha, \beta$ dans $\Pi_{0,3}^\wedge$ et $\bar{\Pi}_{0,3}^\wedge$ conformément à (13), d'où
$$ \alpha, \beta \in \mathrm{Gl}(2, \hat{\mathbf{Z}})^{\pm 1} \quad \text{modulo signes} \leqno (33) $$
mais les éléments correspondants
$$ g = \alpha \rho(\alpha)^{-1}, \quad h = \beta \sigma_\infty(\beta)^{-1} \quad (\in \operatorname{Sl}(2, \hat{\mathbf{Z}})) \leqno (34) $$
sont déterminés dans $\mathrm{Gl}(2, \hat{\mathbf{Z}})$ sans ambiguïté de signe.
De même
$$ \alpha_0 \ (\text{cas } h=1), \quad \beta_0 \ (\text{cas } g=1) \ \in \mathrm{Gl}(2, \hat{\mathbf{Z}})^{\pm 1} \quad \text{modulo signe}, \leqno (35) $$
et alors
$$ \xi = \alpha_0 \rho(\alpha_0)^{-1}, \quad \eta = \beta_0 \sigma_\infty(\beta_0)^{-1} \quad (\in \operatorname{Sl}(2, \hat{\mathbf{Z}})) \leqno (36) $$
sont déterminés sans ambiguïté de signe. On a
$$ f \in \operatorname{Sl}(2, \hat{\mathbf{Z}}) \quad \text{modulo signe} \quad (\text{pour } f \in \bar{\mathfrak{S}}_{0,3} \text{ fixé}) \leqno (37) $$
et on en déduit
$$ \varphi = f \sigma_0(f)^{-1}, \quad \psi = f \rho(f)^{-1} \in \operatorname{Sl}(2, \hat{\mathbf{Z}}) \leqno (38) $$
sans ambiguïté de signe.
Dans le choix des représentations...

\newpage

%% ===== Page 146 =====

de $\varepsilon_0 \in \operatorname{Gl}(2,\hat{\mathbf{Z}})^{\pm 1}$ dans $\operatorname{Sl}(2,\hat{\mathbf{Z}})$ il y a
ambiguïté de signe\footnote{[unclear: non que ce signe en serait qu'un mauvais choix traditionnel à aucun moment un élément d'ordre 2 engendrant un s-groupe à 1 paramètre...]}
$$
\varepsilon_0 = \pm \begin{pmatrix} 1 & -1 \\ 0 & 1 \end{pmatrix}
$$
mais pour
$$
l_0 = \varepsilon_0^2 = \begin{pmatrix} 1 & -2 \\ 0 & 1 \end{pmatrix} \leqno{(39)}
$$
il n'y a pas d'ambiguïté de signe. Ces
étant, les relations
$$
\xi = l_0^2 \varphi \quad (\text{cas } h=1) \leqno{(40)}
$$
et
$$
\eta = l_0^{-2} \psi \quad (\text{cas } g=1) \leqno{(41)}
$$
sont elles vraies sans signes correcteurs ??
Le cas $h=1$ correspond aux données
(dans $\widehat{\Pi}_{0,3}^\wedge$) de $\alpha_0$ et de $f$, satisfaisant
$$
\xi \rho(\xi) \rho^2(\xi) = 1, \quad \xi = l_0^2 (f \varpi_0(f)^{-1})
$$
ou mieux aux données $\alpha_0 \in \widehat{\Pi}_{0,3}^\wedge$, $f \in \widehat{\Pi}_{0,3}^\wedge$, $\nu \in \hat{\mathbf{Z}}$
satisfaisant les relations
$$
\alpha_0 \rho(\alpha_0) = l_0^\nu f \varpi_0(f)^{-1} \leqno{(42)}
$$
en termes de ces données on définit
$$
\begin{cases}
\xi = \alpha_0 \rho(\alpha_0) \\
\varphi = f \varpi_0(f)^{-1} \\
g = \varpi_0(f)^{-1} l_0^\nu \rho(f) = (l_0^\nu f)^{-1} (l_0^\nu \varphi) \rho(l_0^\nu f)
\end{cases} \leqno{(43)}
$$
Ceci vu, [unclear: on] remonte $\alpha_0, f$ dans $\operatorname{Gl}(2,\hat{\mathbf{Z}})^{\pm 1}$

\newpage

%% ===== Page 147 =====

sans ambiguïté de signes, et $\xi, \varphi, g$
sont bien déterminés dans $\operatorname{Sl}(2,\hat{\mathbf{Z}})$, indépendamment
des choix de signes, dans $\alpha_0$ et $f$ et de $g$.
On a donc dans $\operatorname{Sl}(2,\hat{\mathbf{Z}})$
$$
\xi = \varepsilon l_0^v \varphi \qquad \text{i.e.} \qquad \alpha_0 \rho(\alpha_0) = \varepsilon l_0^v f v_\infty f^{-1} \qquad \text{dans } \operatorname{Sl}(2,\hat{\mathbf{Z}})
\leqno{(44)}
$$
avec un signe $\varepsilon \in \{\pm 1\}$ bien déterminé, et
la question qui se pose est de savoir si $\varepsilon = +1$.

Notons que la formule définitive $\rho(l_0) = l_1$
est vraie dans $\operatorname{Sl}(2,\hat{\mathbf{Z}})$ (sans signe correcteur)
donc l'expression de $g$ dans (43) est valable dans
$\operatorname{Sl}(2,\hat{\mathbf{Z}})$ (sans signe), ce qui montre que
$g \in \operatorname{Sl}(2,\hat{\mathbf{Z}})$ est un $\rho$-transformé
de $l_0^v \varphi$, donc son cube donnera le $\rho$-transformé
du cube de $g$. En fait, $g$ satisfait (comme on l'a vu) la
relation
$$
g \rho(g) \rho^2(g) = \varepsilon_0 \cdot 1, \qquad \varepsilon_0 \in \{\pm 1\},
$$
et de plus on a l'identité, si $g' = \alpha g \rho(\alpha)^{-1}$
($\alpha \in \operatorname{Gl}(2,\hat{\mathbf{Z}})$), que
$$
g' \rho(g') \rho^2(g') = \varepsilon_0 \cdot 1
$$
donc en particulier pour $g' = l_0^v \varphi \overset{(44)}{=} \varepsilon \xi$
on a or $\xi$ donné par (42) satisfait à
$$
\xi \rho(\xi) \rho^2(\xi) = 1
$$
(NB on n'a pas $\rho^3 = 1$ dans $\operatorname{Sl}(2,\hat{\mathbf{Z}})$, mais
$\rho^3 = -1$, ce qui indique que l'évaluation [unclear: de l'expression]
de $\xi$ par $\xi \dots$ [unclear: donne $1$]), d'où
$$
\varepsilon_0 = \varepsilon
$$

\newpage

%% ===== Page 148 =====

i.e. la question de la validité de la formule
$$ \xi = l_0^\vee \varphi $$
dans $\operatorname{Sl}(2,\hat{\mathbf{Z}})$ équivaut à celle de la validité
de la formule
$$ g \rho(g) \rho^2(g) = 1 , $$
que nous laissons en suspens pour le
moment.\footnote{Il me semble maintenant pratiquement sûr (d'après les p. [unclear: suivantes]) que par définition de $g$, on a bien $g \rho(g) \rho^2(g) = 1$, et par suite $\xi = l_0^\vee \varphi$ dans $\operatorname{Sl}(2,\hat{\mathbf{Z}})$. La lettre suggère plutôt de poser $g = [unclear: \sigma_\infty (l_0^\vee f)^{-1} \sigma_\infty \dots]$. NB On aura peut-être $\varepsilon = \varepsilon_0$.} Notons pour mémoire les
formules suivantes dans $\operatorname{Sl}(2,\hat{\mathbf{Z}})$

$$
\begin{cases}
\alpha_0, f \in \operatorname{Gl}(2,\hat{\mathbf{Z}})^{\pm 1} \quad (\text{dét } \alpha_0 = \pm 1, \; \text{dét } f = 1, \\
\qquad \qquad \qquad \qquad \quad \text{[unclear: car on en choisit]}) \\
\xi \defeq \alpha_0 \rho(\alpha_0)^{-1} \qquad \Big| \in \operatorname{Sl}(2,\hat{\mathbf{Z}}) \\
\varphi \defeq f \alpha_0 f^{-1} \\
g \defeq [unclear: \sigma_\infty f^{-1} l_0^{\vee -1} \varphi l_0^\vee f] = (l_0^\vee f)^{-1} \cdot (l_0^\vee \varphi) \cdot \rho(l_0^\vee f) \\
\text{où } \xi = \varepsilon \cdot l_0^\vee \varphi \qquad \varepsilon \in \{\pm 1\} \\
[ g \rho(g) \rho^2(g) = \varepsilon \cdot 1 ] \\
g = \varepsilon \cdot \alpha \rho(\alpha)^{-1} \qquad \alpha = (l_0^\vee f)^{-1} \alpha_0 = f^{-1} l_0^{\vee -1} \alpha_0
\end{cases} \leqno{(45) \quad (\text{Cas } h=1)}
$$

De façon symétrique

$$
\begin{cases}
\beta_0, f \in \operatorname{Sl}(2,\hat{\mathbf{Z}}) \quad (\text{[unclear: signes du log], on les choisit}) \\
\eta \defeq \beta_0 \sigma_\infty \rho(\beta_0)^{-1} \\
\psi \defeq f \beta_0 f^{-1} \\
h \defeq f^{-1} l_0^{\vee -1} \sigma_\infty \rho(l_0^\vee f) = \dots
\end{cases} \leqno{(46) \quad (\text{Cas } g=1)}
$$

\newpage

%% ===== Page 149 =====

Je vais regarder le cas $h=1$ (cf (45)), en posant\footnote{NB Dans le cas <<~arithmétique~>> (i.e.\ provenant de $\widehat{\Pi}_{0,3}$) $\alpha_0, f \in \widehat{\Pi}_{0,3} \subset \operatorname{Sl}(2,\hat{\mathbf{Z}})$ avec $\alpha_0 \equiv 1\ (2)$, $f \equiv 1\ (2)$, i.e.\ $a,d,r,u \equiv 1\ (2)$, $b,c,s,t \equiv 0\ (2)$}\footnote{changements de signes dans $\alpha_0, f$ est immatériel}
$$ 
\begin{cases} 
\alpha_0 = \begin{pmatrix} a & b \\ c & d \end{pmatrix} \\ 
f = \begin{pmatrix} r & s \\ t & u \end{pmatrix} 
\end{cases} \leqno{(47)} 
$$
où $a,b,c,d,r,s,t,u \in \hat{\mathbf{Z}}$, avec
$$ ad-bc = \delta \in \{ \pm 1 \} \leqno{(48)} $$
$$ ru-st = 1 \leqno{(49)} $$

On pose
$$ \xi = \alpha_0 p (\alpha_0)^{-1} = \begin{pmatrix} A & B \\ C & D \end{pmatrix} \leqno{(50)} $$
$$ \varphi = f \sigma_\infty f^{-1} = \begin{pmatrix} R & S \\ T & U \end{pmatrix} = \begin{pmatrix} r^2+s^2 & rt+su \\ rt+su & t^2+u^2 \end{pmatrix} \leqno{(51)} $$
de sorte qu'on aura (cf § 42, formules 78)
\marginpar{NB on a $\delta=\pm 1$}
$$ 
\begin{cases} 
\delta A = [unclear: (a^2+ab+b^2)] - (ac+bd+bc) & \delta B = ac+bd+bc \\ 
\delta C = ac+bd+ad - (c^2+d^2+cd) & \delta D = c^2+d^2+cd 
\end{cases} \leqno{(52)} 
$$
et (§ 42, (90) et 91)
$$ R = r^2+s^2, \quad S = T = rt+su, \quad U = t^2+u^2 \leqno{(53)} $$

On aura les formules
$$ A+D = 1 \leqno{(54)} $$
$$ S = T \leqno{(55)} $$

La relation essentielle reliant $\alpha_0$ et $f$, i.e.\ reliant les <<~variables~>> $a,b,c,d$ aux $r,s,t,u$, est la relation
\marginpar{<<~relation des lacets~>>}
$$ \xi = \varepsilon l_0^\nu \varphi \qquad \varepsilon \in \{ \pm 1 \} \leqno{(56)} $$
qui s'écrit aussi
$$ l_0^{-\nu} \xi = \varepsilon \varphi, \quad \text{i.e.} $$

\newpage

%% ===== Page 150 =====

$$ 
\begin{pmatrix} A+2\nu C & B+2\nu D \\ C & D \end{pmatrix} = \varepsilon \begin{pmatrix} R & S \\ T & U \end{pmatrix} 
\leqno (57) 
$$

Mais cette condition implique d'abord que la matrice du premier membre est symétrique (puisque $S=T$) i.e.
$$ 
2\nu D = C - B 
$$
condition qui se réduit --- par (54), à
$$ 
2\nu D = 1 - D \quad \text{i.e.} \quad \mu D = 1 \quad \text{i.e.} \quad D = 1/\mu \, , \, \mu = 1/D 
\leqno (58) 
$$
où $\mu = 2\nu + 1$,
ce qui s'explique déjà pour $D$ inversible.

En écrivant (56) sous la forme
$$ 
\varepsilon \begin{pmatrix} A & B \\ C & D \end{pmatrix} = \begin{pmatrix} R - 2\nu T & S - 2\nu U \\ T & U \end{pmatrix} 
$$
La condition $C+D-B=1$ devient, sur les éléments de la matrice $\varphi$
$$ 
T+U-(S-2\nu U) = \varepsilon 
$$
i.e. \quad $(2\nu + 1) U = \varepsilon \quad \text{i.e.}$
$$ 
\mu U = \varepsilon \quad \text{i.e.} \quad U = \varepsilon/\mu \, , \, \mu = \varepsilon/U 
\leqno (59) 
$$
ce qui équivaut --- tenant compte de (58), à la relation
$$ 
D = \varepsilon U \quad - \quad U = \varepsilon D 
\leqno (60) 
$$
qui n'est autre que l'égalité des termes aux coins inférieurs droits des matrices (57).

\newpage

%% ===== Page 151 =====

On peut expliciter les relations (57), comme quatre relations sur les $9$ variables $a,b,c,d, r,s,t,u, \mu$ en plus des relations (48), (49) :
\marginpar{$A,B$ calculés en fonction des autres compte tenu de $C+D-B=1$ et $AD-BC=1$}
$$
\begin{cases}
D = \varepsilon U \\
\mu D = 1 \quad \text{i.e. } \mu U = \varepsilon \\
C = \varepsilon T \\
A = \varepsilon R - \frac{(\mu-1)}{2\nu} C
\end{cases}
\leqno (61)
$$

soit -
\marginpar{Conséquence des précédentes}
$$
\begin{cases}
c^2 + d^2 + cd = \varepsilon \delta (t^2 + u^2) \\
ac + bd + ad - (c^2 + d^2 + cd) = \delta \varepsilon (rt + su) \\
a^2 + b^2 + ab - (ac + bd + bc) = \delta \varepsilon \left(r^2 + s^2 - (\mu-1)(rt + su)\right)
\end{cases}
\leqno (62)
$$

Quand les signes $\varepsilon, \delta$ sont fixés, on trouve avec (48), (49) un système de 6 équations en 9 variables $a,b,c,d,r,s,t,u,\mu$ --- on doit pouvoir vérifier que le schéma (sur $\operatorname{Spec} \mathbf{Z}$) des solutions en question est de dimension\marginpar{non, il y a 4 dim. de relations} relative 3, et de dimension relative 1 sur les "facteurs" $(a,b,c,d)$ resp. $(r,s,t,u)$. On a donc un point adélique de ce schéma absolu --- un pt à coefficients dans $\hat{\mathbf{Z}}$...

\textbf{Théorème} Le multiplicateur $\mu$ est déterminé en fonction de $\zeta = \begin{pmatrix} A & B \\ C & D \end{pmatrix}$ (considéré comme élém. de $\operatorname{Sl}(2,\hat{\mathbf{Z}})$) par la relation
$$ 
\mu D = 1 \leqno (58) 
$$
et il est déterminé au signe près par $\varphi = \begin{pmatrix} R & S \\ T & U \end{pmatrix}$ par la relation
$$ 
\mu U = \varepsilon \leqno (59) 
$$
où $\varepsilon \in \{\pm 1\}$,

\newpage

%% ===== Page 152 =====

et où $\varepsilon$ est le signe qui intervient dans les formules
$$
\begin{cases}
f = \varepsilon l_0^4 \varphi & (56) \\
g \rho g \rho g^{-1} = \varepsilon & (45) \\
g = \varepsilon \alpha \rho \alpha \rho^{-1}
\end{cases}
$$

On a\footnote{en d'autres termes, on a $\mu \equiv \varepsilon \pmod 4$. Car $U \equiv 1 \pmod 4$ et $\varepsilon \equiv \varepsilon \pmod 4$, donc $\mu = \varepsilon / U \equiv \varepsilon \pmod 4$. Ce serait [unclear: dommage] d'oublier pour la démonstration [unclear: la même] idée.} \quad $\varepsilon = 1$ ssi $\mu \equiv 1 \pmod 4$ (i.e.\ $\nu \equiv 0 \pmod 2$), \quad $\varepsilon = -1$ ssi $\mu \equiv -1 \pmod 4$ (i.e.\ $\nu \equiv 1 \pmod 2$).

Pour la prouver, on se rappelle que
$$ U = t^2+u^2 $$
on voit que les seules valeurs mod 4 prises par $t^2$ et $u^2$ sont 0 et 1, donc celles prises par $t^2+u^2$ sont 0, 1 et 2, et la seule unité parmi elles est 1. Donc [unclear: $\varepsilon/\mu$ prend sa valeur dans le groupe] $(\mathbf{Z}/4\mathbf{Z})^\times$ mod 4 ssi $\varepsilon/\mu \equiv 1 \pmod 4$
ou encore $\mu \equiv \varepsilon \pmod 4$ cqfd.

Ainsi on aurait vu de façon toute analogue, que : la formule $\mu D = 1 \quad (58)$ et $D = \delta(c^2+d^2+cd)$ donne $\mu \equiv \delta \pmod 3$ i.e.\ $\delta = 1 \ (\text{resp.\ } -1)$ ssi $\mu \equiv 1 \pmod 3 \quad (\text{resp.\ } \mu \equiv -1 \pmod 3)$. 

notons, pour $p$ un premier, $\varepsilon_p$ la valeur du ``caractère quadratique'' relatif à $p$, égal à 1 ssi $\mu \in (\mathbf{Z}_p^\times)^2$. On aura donc
$$
\begin{cases}
\delta = ad-bc = \chi_3(\mu) = \varepsilon_3 \\
\varepsilon = [unclear: D/U] = \chi_2(\mu) = \varepsilon_2
\end{cases}
\leqno{(63)}
$$
\marginpar{$(\mathbf{Z}/12\mathbf{Z})^\times \simeq (\mathbf{Z}/3\mathbf{Z})^\times \times (\mathbf{Z}/4\mathbf{Z})^\times \simeq \{\pm 1\} \times \{\pm 1\}$}
et la connaissance de $\varepsilon_3$ (resp.\ $\varepsilon_2$, resp.\ des deux) équivaut à celle de $\mu \pmod 3$ (resp.\ $\pmod 4$, resp.\ $\pmod{12}$).

\newpage

%% ===== Page 153 =====

\noindent
Nous allons aussi calculer \sout{sur ces (cf formules (45))}
$$
\alpha = (l_0^\nu f)^{-1} \alpha_0 = f^{-1} l_0^{-\nu} \alpha_0
$$
$$
g = \varepsilon \alpha \rho(\alpha)^{-1} = \sigma_0 \underbrace{f^{-1}}_{t_f} l_1^\nu \rho(f) = (l_0^\nu f)^{-1} f \rho(l_0^\nu f)
$$
On trouve\footnote{NB Contrairement à $\alpha, f$, et à (64) qui s'en déduisent, l'évaluation qu'on va trouver pour $g$ est indépendante des ambiguïtés de signes dont on a pu les affecter (i.e. du choix du signe de $\mu$ etc...). Ainsi cette éval. reste valable dans tous les cas.} :
$$
\alpha = \underbrace{\begin{pmatrix} u & -s \\ -t & r \end{pmatrix}}_{f^{-1}} \begin{pmatrix} 1 & 2\nu \\ 0 & 1 \end{pmatrix} \begin{pmatrix} a & b \\ c & d \end{pmatrix} = \begin{pmatrix} u(a+2\nu c)-cs & u(b+2\nu d)-ds \\ -t(a+2\nu c)+cr & -t(b+2\nu d)+dr \end{pmatrix} \leqno{(64)}
$$
$$
g = \underbrace{\begin{pmatrix} r & t \\ s & u \end{pmatrix}}_{\sigma_0(f^{-1}) = t_f} \underbrace{\begin{pmatrix} 1 & 0 \\ 2\nu & 1 \end{pmatrix}}_{l_1^\nu} \underbrace{\begin{pmatrix} 1+\mu t(t+u) & -\mu t^2 \\ -1+\mu u(t+u) & 1-\mu ut \end{pmatrix}}_{\rho(f) \text{ (cf § 43 formules (77))}} \quad \text{, soit}
$$
$$
g = \begin{pmatrix} 1+\mu t(t+u) & -\mu t^2 \\ -1+\mu u(t+u) & 1-\mu u t \end{pmatrix} \leqno{(65)}
$$

Il apparaît que $g$ s'exprime simplement en termes de $f$ i.e. de $r,s,t,u$ directement, et non en termes de $\alpha_0$ ou $\alpha$ i.e. via les $a,b,c,d - A,B,C,D$. Par contre, pour exprimer $\alpha$ (64), il faut faire intervenir à la fois $\underbrace{r,s,t,u}_{f}$ et les $a,b,c,d$.

Il était naturel, en relation avec l'hom. canonique
$$
N_{1,1} \to \operatorname{Gl}(2, \hat{\mathbf{Z}})
$$
(cf § 39, (1)) de chercher une matrice [unclear: $\hat{g}$] qui se [unclear: superpose] sur ce $g$ - $g$ je vais désigner par la même lettre, $g$, l'automorphisme envisagé de $\hat{\Pi}_{0,3}$ dont on était parti. (Tout comme je désigne par les mêmes lettres les éléments $f,g$ etc de $\hat{\Pi}_{0,3}$ ou de $\hat{\mathcal{G}}_{0,3}$, et leurs images dans $\operatorname{Gl}(2, \hat{\mathbf{Z}})/\pm 1$, voire même des éléments correspondants dans $\operatorname{Gl}(2, \hat{\mathbf{Z}})$ lui-même).

\newpage

%% ===== Page 154 =====

$$
u \in \operatorname{Gl}(2, \hat{\mathbf{Z}})
$$
[unclear: soit] de déterminant égal à $\mu$, et induisant sur $\operatorname{Sl}(2, \hat{\mathbf{Z}})'$, par autom. intérieur, le même autom. ([unclear: à l'intime] près) que [unclear: notre] pour $f \mapsto g$ en question), ce qui s'explicite par les formules
$$
\det u = \mu \leqno{(66)}
$$
$$
u \sigma_\infty u^{-1} = \varepsilon' \sigma_\infty \qquad \varepsilon' \in \{\pm 1\} \leqno{(67)}
$$
$$
u \rho u^{-1} = \varepsilon'' g \qquad \varepsilon'' \in \{\pm 1\} \leqno{(68)}
$$

En fait, (68) s'écrit
$$
[u, \rho] = u \rho u^{-1} = \varepsilon'' g
$$
et appliquant aux deux membres l'opération $x \mapsto x \rho(x) \rho^2(x)$, et tenant compte que $g \rho(g) \rho^2(g) = \varepsilon_2$, or comme on a $x \rho(x) \rho^2(x) = 1$ pour $x = u \rho u^{-1}$, on trouve
$$
\varepsilon'' = \varepsilon_2 \leqno{(68 \text{ bis})}
$$
et (68) équivaut donc à
$$
u \rho u^{-1} = \varepsilon_2 g = \alpha \rho(\alpha)^{-1}
$$
d'où i.e.\ à
$$
u = \alpha u_0 \qquad \text{avec} \qquad [u_0, \rho] = 1 \qquad (\text{i.e. } u_0 \rho u_0^{-1} = 1). \leqno{(69)}
$$

Nous avons déjà remarqué que (67), (68) déterminent $u$ seulement modulo multiplication par un scalaire de $\hat{\mathbf{Z}}^*$, et compte tenu de (66), $u$ est déterminé seulement modulo $\lambda$ satisfaisant $\lambda^2 = 1$ i.e.\ $\lambda \in \{\pm 1\}^P$ (où $P$ est l'ens. des nbs premiers), mais on espère néanmoins tomber sur un choix canonique.

Dans mes premiers calculs, j'avais pris comme départ (69), où la condition $u_0$ commutant à $\rho$ s'exprimait sous la forme
$$
u_0 = p \mathbf{1} + q \rho = \begin{pmatrix} p & q \\ -q & p-q \end{pmatrix} \quad (p, q \in \hat{\mathbf{Z}}) \leqno{(70)}
$$

\newpage

%% ===== Page 155 =====

de sorte que
$$
\det u_0 = p^2 + p q + q^2 \leqno{(71)}
$$
et que (66) s'écrit (compte tenu de (69), où $\det \alpha = \delta = \varepsilon_3$)
$$
p^2 + p q + q^2 = \mu \varepsilon_3 \leqno{(72)}
$$
ce qui, compte tenu de la relation $c^2 + c d + d^2 = \varepsilon_3 / \mu \quad (51 \text{ et } 58)$
suggérait des relations de la forme
$$
p = \mu c, \quad q = \mu d, \quad \text{ou} \quad p = \mu d, \quad q = \mu c
$$
En fait, exprimant (67) comme équation en $p, q$ (s'ajoutant à 72), je trouvais (au signe près bien sûr) que [unclear: l'ordre] est inversé)
$$
p = \mu d, \quad q = \mu c \quad \text{i.e.} \quad u_0 = \mu(c\rho + d) \leqno{(73)}
$$
mais en me plaçant dans le cas $\varepsilon' = +1$ exclusivement.
(j'avais cru conclure, un peu trop hâtivement, que le signe $\varepsilon'$ était toujours $+1$, avant de m'apercevoir que ce n'était guère compatible avec le corps des classes, et de subodorer que l'on devait avoir $\varepsilon' = +1$ si $p \equiv 1 \ (4)$.) En faisant ce détour assez pesant par les $(a,b,c,d)$, je suis quand même arrivé, après bien des péripéties (erreurs de calcul etc) à la formule où il n'y avait plus de $a,b,c,d$ du tout, d'une simplicité déroutante :

\marginpar{Je souligne la lettre $\underline{u}$ pour la distinguer de $u$}
$$
u_0 = \mu(t \sigma_\infty + \underline{u}) = \mu \begin{pmatrix} \underline{u} & t \\ -t & \underline{u} \end{pmatrix} \qquad (\Longleftrightarrow \text{sans doute, } \varepsilon_2 = 1) \leqno{(74)}
$$
qui correspond manifestement au cas $\varepsilon'_1 = 1$, i.e. celui où $u$ commute à $\sigma_\infty$. Comme...

\newpage

%% ===== Page 156 =====

$$
\det(t v_\infty + u) = t^2 + u^2 = \varepsilon_2 / \mu \quad (= 1 / \mu \text{ si } \varepsilon_2 = +1) ,
$$
on aura bien, dans le cas $\varepsilon_2 = 1$, donc
$$
\det u_0 = \mu^2 \det(t v_\infty + u) = \mu^2 1/\mu = \mu
$$
i.e.\ (66). D'autre part, on aura\footnote{On pourrait utiliser les formules § 93 (78).}
$$
[u_0, \rho] = [t v_\infty + u, \rho] = \begin{pmatrix} u & t \\ -t & u \end{pmatrix} \begin{pmatrix} 1 & -1 \\ 1 & 0 \end{pmatrix} \mu \begin{pmatrix} u & -t \\ t & u \end{pmatrix} \begin{pmatrix} 0 & 1 \\ -1 & 1 \end{pmatrix}
$$
$$
= \begin{pmatrix} -t & u+t \\ -u & u-t \end{pmatrix} \begin{pmatrix} u-t & -u \\ u+t & -t \end{pmatrix}
$$
$$
= \mu \begin{pmatrix} \frac{1}{\mu} + t(u+t) & -t^2 \\ t(u-t) & \frac{1}{\mu} - ut \end{pmatrix}
$$
$$
= \begin{pmatrix} 1 + \mu t(u+t) & -\mu t^2 \\ \mu t(u-t) & 1 - \mu ut \end{pmatrix} ,
$$
et en comparant avec (65) on trouve bien
$$
[u_0, \rho] = g
$$
compte tenu que
$$
\mu(u^2+t^2) = \varepsilon_2 = 1
$$
donc
$$
-1 + \mu u(t+u) = \underset{-\mu t^2}{-1 + \mu u^2} + \mu ut = \mu t(u-t) !
$$
J'aimerais maintenant, sans me lancer dans des nouveaux calculs, trouver "au pif" une expression pour $u_0$ dans le cas où $\varepsilon_2 = -1$, et où on s'attend à trouver dans (67) le signe $\varepsilon' = -1$ -- i.e.\ que $u_0$ soit une antisimilitude
$$
{}^t u u = - \det u
$$
Mais la matrice
$$
\tau_\infty = \begin{pmatrix} -1 & 0 \\ 0 & 1 \end{pmatrix} \quad (\text{c'est une réflexion orthogonale})
$$
est le cas-type d'une telle antisimilitude, et je prévois que la matrice
$$
u_0 \overset{\text{pif}}{=} \mu \begin{pmatrix} -1 & 0 \\ 0 & 1 \end{pmatrix} \begin{pmatrix} u & t \\ -t & u \end{pmatrix} = \mu \begin{pmatrix} -u & -t \\ -t & u \end{pmatrix}
$$

\newpage

%% ===== Page 157 =====

conviendra. On aura, bien sûr
$$ u {}^t u = \mu^2 \begin{pmatrix} -u & -t \\ -t & u \end{pmatrix} \begin{pmatrix} -u & -t \\ -t & u \end{pmatrix} = \mu^2 \begin{pmatrix} u^2+t^2 & 0 \\ 0 & u^2+t^2 \end{pmatrix} = \mu^2 \begin{pmatrix} -\frac{1}{\mu} & 0 \\ 0 & -\frac{1}{\mu} \end{pmatrix} $$
i.e.
$$ u {}^t u = -\mu \cdot I_d $$
et
$$ \det u = \mu^2 (-u^2 - t^2) = \mu^2 \left(\frac{1}{\mu}\right) = \mu $$
compte tenu que maintenant
$$ u^2 + t^2 = \varepsilon_2 / \mu = -\frac{1}{\mu} \qquad \text{puisqu'} \quad \varepsilon_2 = -1 . $$

Enfin pour vérifier (68) calculons
$$
\begin{aligned}
[u, \rho] &= \left[ \mu \begin{pmatrix} -u & -t \\ -t & u \end{pmatrix} , \rho \right] = \begin{pmatrix} u & t \\ t & -u \end{pmatrix} \begin{pmatrix} 0 & 1 \\ -1 & 1 \end{pmatrix} \mu \begin{pmatrix} -u & -t \\ -t & u \end{pmatrix} \begin{pmatrix} 1 & -1 \\ 1 & 0 \end{pmatrix} \\
&= \begin{pmatrix} -t & u+t \\ -u & u-t \end{pmatrix} \mu \begin{pmatrix} -u-t & u \\ u-t & t \end{pmatrix} \\
&= \mu \begin{pmatrix} u(u+t) & t^2 \\ -(u+t)^2 - u(u-t) & u^2+t^2-ut \end{pmatrix} = \begin{pmatrix} \mu u(u+t) & \mu t^2 \\ 1-\mu u(u+t) & -1-\mu ut \end{pmatrix}
\end{aligned}
$$
(compte tenu que $\mu(u^2+t^2) = -1$), d'où a-t-on bien
$$ [u, \rho] = -g \quad \text{i.e.} $$
compte de la formule (65) qui explicite $g$ en
termes de $r,s,t,u$, on trouve que cette formule
pourrie équivaut à
$$ 2\mu ut = 0 \quad \text{i.e.} \quad ut = 0 $$
décidément, ça déconne! Mais je vois que si on
remplace dans la définition de $u$ $t$ par $-t$,
i.e. si on définit
$$ u = \mu (\tau_\infty + u)\tau_\infty = \mu \begin{pmatrix} u & t \\ -t & u \end{pmatrix} \begin{pmatrix} 1 & 0 \\ 0 & -1 \end{pmatrix} = \mu \begin{pmatrix} u & -t \\ -t & -u \end{pmatrix} $$
(au lieu de $u = \mu \begin{pmatrix} -u & -t \\ -t & u \end{pmatrix}$), alors $u$ est toujours
une antisimilitude de déterminant $\mu$, et de
plus maintenant
$$ [u, \rho] = \begin{pmatrix} \mu u(u-t) & \mu t^2 \\ 1-\mu u(u+t) & -1+\mu ut \end{pmatrix} = -g $$

\newpage

%% ===== Page 158 =====

ce qui se réduit en effet à vérifier
$$ \mu(u-t) = -(1+\mu t(t-u)) $$
i.e.
$$ \mu[u - t + t^2 - tu] = -1 \qquad \text{O.K.} $$

En résumé, on trouve

Théorème Soit $\alpha_2 \in \mathbf{Z}/2\mathbf{Z}$ défini par $(-1)^{\alpha_2} = \varepsilon_2$
(donc $\alpha_2 \equiv 0 \pmod 2$ ssi $\mu \equiv 1 \pmod 4$, $\alpha_2 \equiv 1 \pmod 2$ ssi $\mu \equiv -1 \pmod 4$). Posons\footnote{Dans la marge, l'auteur a explicité cette matrice : $= \begin{pmatrix} u & \varepsilon_2 t \\ -t & \varepsilon_2 u \end{pmatrix} = \begin{cases} \mu \begin{pmatrix} u & t \\ -t & u \end{pmatrix} & \alpha_2 = 0 \\ \mu \begin{pmatrix} u & -t \\ -t & -u \end{pmatrix} & \alpha_2 = 1 \end{cases}$}
$$ \underline{u} = \mu(t\overline{\infty}+u)\overline{\infty}^{\alpha_2} = \mu\begin{pmatrix} u & t \\ -t & u \end{pmatrix} \begin{pmatrix} 1 & 0 \\ 0 & \varepsilon_2 \end{pmatrix} \leqno{(75)} $$
On a alors les formules
$$ \det \underline{u} = \mu \leqno{(76)} $$
$$ [\underline{u}, \overline{\infty}] = -\varepsilon_2\cdot \operatorname{id} \quad \text{i.e. } \underline{u}(\overline{\infty}) = \varepsilon_2 \overline{\infty} \leqno{(77)} $$
$$ [\underline{u}, \rho] = \varepsilon_2\cdot g \quad \text{i.e. } \underline{u}(\rho) = (\varepsilon_2 g)\cdot \rho \leqno{(78)} $$

Corollaire Soit $u$ un autom. du groupe
à lacets $\hat{\Pi}_{0,3}$, invariant relative-
ment à $\mathfrak{S}_3$ (i.e. à $\overline{\infty}$ et à $\rho$), $\mu$
son multiplicateur, désignons aussi par
$u$ l'extension canonique de $u$ à
$\mathfrak{S}_{0,3}^{+\wedge} \simeq \operatorname{Sl}(2,\hat{\mathbf{Z}})^{\wedge'}$, induisant l'identité sur $\mathfrak{S}_3 \simeq \mathfrak{S}_{0,3}^{+\wedge}/\hat{\Pi}_{0,3}$,
et considérons le quotient $\operatorname{Sl}(2,\hat{\mathbf{Z}})'$ de
$\operatorname{Sl}(2,\hat{\mathbf{Z}})^{\wedge'}$. Alors l'action de $u$ passe
en une action dans le quotient $\operatorname{Sl}(2,\hat{\mathbf{Z}})'$,
et cette action est induite par autom.
intérieur de $\operatorname{Gl}(2,\hat{\mathbf{Z}})$, définie par une
matrice $\underline{u} \in \operatorname{Gl}(2,\hat{\mathbf{Z}})$ telle que $\det \underline{u} = \mu$.

\newpage

%% ===== Page 159 =====

(ce qui décrit l'action extérieure de $u$ sur $\operatorname{Sl}(2, \hat{\mathbf{Z}})'$ en termes du multiplicateur $\mu$).

\marginpar{NB Pour [unclear: vérifier] ce corollaire, j'ai pris soin de choisir un $u$ qui laisse $\sigma_\infty$ et $\rho$ fixes. On a en fait [unclear: supposé que] ceci se passe [unclear: ainsi].}
En fait, on a beaucoup mieux encore modulo signe, puisqu'on a défini canoniquement en termes de $u$, la matrice $\underline{u}$ par la formule (75), où $t,u \in \hat{\mathbf{Z}}$ sont définis ainsi: la donnée de l'autom. $u$ définit $f \in \hat{\Pi}_{0,3}$, modulo opération $f \mapsto l_0^n f$ (pour $n \in \hat{\mathbf{Z}}$), et on désigne encore par $f$ l'image de $f$ dans $\operatorname{Sl}(2, \hat{\mathbf{Z}})/\pm 1$, qui est une matrice $\begin{pmatrix} r & s \\ t & u \end{pmatrix}$ définie mod signe.

On note que dans $\operatorname{Sl}(2, \hat{\mathbf{Z}})$ on a
$$
l_0^n \begin{pmatrix} r & s \\ t & u \end{pmatrix} = \begin{pmatrix} 1 & -2n \\ 0 & 1 \end{pmatrix} \begin{pmatrix} r & s \\ t & u \end{pmatrix} = \begin{pmatrix} r - 2nt & s - 2nu \\ t & u \end{pmatrix} \leqno{(79)}
$$
($n \in \hat{\mathbf{Z}}$), donc $(t,u)$ ne change pas (mod signe) quand on remplace $f$ par $l_0^n f$ --- donc $(t,u)$ est déterminé mod signe par l'autom. $u$ de $\hat{\Pi}_{0,3}$. L'observation analogue s'applique d'ailleurs à $\alpha_0 = \begin{pmatrix} a & b \\ c & d \end{pmatrix}$, qui est déterminé mod signe à pre-multiplication à gauche par un $l_0^n$ --- le vecteur $(c,d)$ est déterminé mod signe, donc l'élément
$$
c\rho + d
$$

\newpage

%% ===== Page 160 =====

est déterminé mod signes, et on doit avoir une relation
$$
u \overset{?}{=} \mu \alpha(cp+d) = \mu f^{-1} l_0^{-\nu} \alpha_0(cp+d) \leqno{(80)}
$$
\footnote{NB cette formule est en accord avec cf (69), (73), pourvu au moins que $\varepsilon_2=1$ et $\delta=\varepsilon_3$, cf (82) \dots sinon il y a lieu peut-être de corriger par un facteur $\tau_\infty$ quelque part\dots}

$$
l_0^{-\nu} \alpha_0(cp+d) = \begin{pmatrix} 1 & 2\nu = \mu-1 \\ 0 & 1 \end{pmatrix} \begin{pmatrix} a & b \\ c & d \end{pmatrix} \begin{pmatrix} d & c \\ -c & c+d \end{pmatrix}
$$
$$
\begin{pmatrix} ad-bc = \delta & ac+bc+bd = \delta B \\ 0 & c^2+cd+d^2 = \delta D \end{pmatrix} = \delta \begin{pmatrix} 1 & B \\ 0 & D \end{pmatrix}
$$
$$
l_0^{-\nu} \alpha_0(cp+d) = \delta \begin{pmatrix} 1 & B+2\nu D = C \\ 0 & D \end{pmatrix}
$$

i.e.
$$
l_0^{-\nu} \alpha_0(cp+d) = \delta \begin{pmatrix} 1 & C \\ 0 & D \end{pmatrix} \leqno{(81)}
$$

d'où
$$
\mu \alpha(cp+d) = \mu f^{-1} (l_0^{-\nu} \alpha_0(cp+d)) = \mu \delta \begin{pmatrix} u & -s \\ -t & r \end{pmatrix} \begin{pmatrix} 1 & C \\ 0 & D \end{pmatrix}
$$
$$
= \mu \delta \begin{pmatrix} u & uC - sD \\ -t & -tC + rD \end{pmatrix}
$$
$$
= \delta \begin{pmatrix} \mu u & -s + \mu u C \\ -\mu t & r - \mu t C \end{pmatrix}
$$

or (61)
$$
C = \varepsilon_2 T = \varepsilon_2 (rt+su)
$$

d'où
$$
-s + \mu u C = -s + \varepsilon_2 \mu u r t + \varepsilon_2 \mu u^2 s = s \underbrace{(-1 + \varepsilon_2 \mu u^2)}_{= -\varepsilon_2 \mu t^2} + \varepsilon_2 \mu u r t = \varepsilon_2 \mu t \underbrace{(-st + ur)}_{= +1} = \varepsilon_2 \mu t
$$

\newpage

%% ===== Page 161 =====

$$r - \mu t C = r - \varepsilon_2 \mu r t^2 - \varepsilon_2 \mu s t u = \varepsilon_2 \mu u (ru-st) = \varepsilon_2 \mu u \quad , \quad d\text{'où}$$
[MARGIN: $\overbrace{r(1-\varepsilon_2 \mu t^2)}^{r(\varepsilon_2 \mu s u)}$]

$$ \mu \alpha(cp+d) = \delta \mu \begin{pmatrix} u & -\varepsilon_2 t \\ -t & \varepsilon_2 u \end{pmatrix} = \delta \underline{u} \quad , $$

i.e.
(82) $\qquad \underline{u} = \delta \mu \alpha(cp+d) = (\delta \mu) (f^{-1} l_0^{-2} \alpha_0) (cp+d)$ [MARGIN: où $\varepsilon_2=1$, et $u \equiv \mu (4)$]

Formule qui s'écrit aussi
i.e.
(83) $\qquad f (t_{0\infty} * \underline{u}) = \delta \alpha (cp+d)$
où $\alpha = f^{-1} l_0^{-2} \alpha_0 \quad , \quad \delta = \varepsilon_3 \in \{\pm 1\}, \delta = ad-bc, \delta \equiv \mu \pmod 3$

ou encore
(84) $\qquad f (t_{0 \infty} * \underline{u}) \tilde{t}_{\infty}^{\alpha_2} = \delta l_0^{-2} \alpha_0 (cp+d)$

$$ \begin{pmatrix} r & s \\ t & u \end{pmatrix} \begin{pmatrix} 1 & 0 \\ 0 & \varepsilon_2 \end{pmatrix} \bigg/ \begin{pmatrix} 1 & 2\nu \\ 0 & 1 \end{pmatrix} \begin{pmatrix} a & b \\ c & d \end{pmatrix} \bigg/ \begin{pmatrix} d & c \\ -c & c+d \end{pmatrix} $$

[unclear crossed out equations: $\begin{pmatrix} 1 & rt+su \\ 0 & \varepsilon_2 t^2 + u^2 \end{pmatrix} \dots = \begin{pmatrix} 1 & C \\ 0 & D \end{pmatrix}$ cf (81)]

et sous cette forme la formule [unclear: extraite] en fait
conserve, [unclear: comme] le contenu essentiel de
la fameuse relation des lacets, sans fioritures
matricielles (56) ou (57)

$$ \varepsilon_2 u = D \quad , \quad \varepsilon_2 t = C \quad ! $$

Remarque. En fait, sous une forme apparemment
d'arithmétique, on a fait par la bande des
[unclear: formules] d'algèbre linéaire, valables sur des
[unclear]

\newpage

%% ===== Page 162 =====

anneaux quelconques et formelles sans doute sous des conditions beaucoup plus générales et plus « compréhensives » concernant des équations de la forme
$$
[\alpha, \rho] = \xi
$$
où $\rho \in \operatorname{Gl}(2,A)$ est donné, ainsi que $\xi$, et $\alpha \in \operatorname{Gl}(2,A)$ est l'inconnue (assujettie éventuellement à des conditions du type $\det\alpha = \delta$, $\delta$ fixé), ainsi que des équations simultanées en $\alpha$ et $f$
$$
[\alpha, \rho] = \varepsilon \cdot \mathcal{E}_0(\mu-1) [f, \sigma] \quad (cf\ (56)) \leqno{(*)}
$$
où $\rho, \sigma \in \operatorname{Gl}(2,A)$ sont donnés ainsi qu'un scalaire $\varepsilon \in A^*$, et $\lambda \mapsto \mathcal{E}_0(\lambda)$ est un « sous-groupe à un paramètre », p. ex. on peut supposer
$$
\mathcal{E}_0(\lambda) = \begin{pmatrix} 1 & -\lambda \\ 0 & 1 \end{pmatrix} = \varepsilon_0^\lambda , \qquad \varepsilon_0 = \begin{pmatrix} 1 & -1 \\ 0 & 1 \end{pmatrix}
$$
— en supposant au besoin une relation de la forme $\sigma f = \varepsilon_0$.
Modulo des ajustements convenables, toutes les formules du présent $\S$ devraient pouvoir se développer dans ce contexte plus général (ainsi qu'engendrer des complications substantielles) — mais je ne sais si je vais expliciter

\newpage

%% ===== Page 163 =====

ici ces questions délectables d'algèbre linéaire, voire d'interprétations schématiques du genre que le schéma des solutions de $(*)$ (pour $f$, $\alpha$, $\rho$, $\mu$ variables) est un schéma régulier lisse de dimension relative 4, qui sait. Je vais seulement expliciter, dans le contexte présent adélique, trop particulier manifestement, la

\textbf{Proposition} Soient $c, d, t, u \in \hat{\mathbf{Z}}$, pour qu'ils correspondent à des données $\alpha_0 = \begin{pmatrix} a & b \\ c & d \end{pmatrix}$, $f = \begin{pmatrix} r & s \\ t & u \end{pmatrix}$ satisfaisant
$$
ad - bc \in \{ \pm 1\}, \ ru - st = 1, \ \alpha_0 \rho (\alpha_0)^{-1} = \varepsilon_2 f_0^2(f \sigma_\alpha / f)^{-1} \defeq \varepsilon_2 \leqno (85)
$$
avec $\varepsilon_2 \in \{\pm 1\}$ et $f \in \hat{\mathbf{Z}}$ convenables (cf. (48)(49)(53)), il faut et il suffit qu'on ait
$$
c^2 + d^2 = \pm (t^2 + u^2) \in \hat{\mathbf{Z}}^* \leqno (85)
$$
Alors $(\varepsilon_2, \varepsilon_3, \mu) \in \{ \pm 1\} \times \{ \pm 1\} \times \hat{\mathbf{Z}}^*$ est déterminé modulo le groupe [unclear: engendre par] $(- \varepsilon_2, - \varepsilon_3, - \mu)$ par $c, d, t, u$, par les relations (où $\mu = 2 \gamma + 1$)
$$
\mu = \varepsilon_2 (t^2 + u^2) = \varepsilon_3 (c^2 + cd + u^2) \leqno (86)
$$
(qui détermine $\varepsilon_3, \mu$ quand on se donne $\varepsilon_2$)
et on a
$$
\begin{cases}
\varepsilon_2 \equiv \mu \pmod 4 \\
\varepsilon_3 \equiv \mu \pmod 3 \\
\gamma = (\mu - 1)/2
\end{cases} \leqno (87)
$$

\newpage

%% ===== Page 164 =====

Les solutions s'obtiennent alors ainsi : on choisit $\varepsilon_2, \varepsilon_3, \mu$ comme dessus (deux possibilités), puis arbitrairement $r,s \in \hat{\mathbf{Z}}$ satisfaisant
$$ u r - t s = 1 $$
(il y en a toujours car $t^2+u^2 \in \hat{\mathbf{Z}}^*$ - donc $\forall p \in \mathbf{P}$, $t_p$ et $u_p$ ne sont pas tous deux dans $p\mathbf{Z}_p$)
et si $r_0, s_0$ est une solution, les autres s'obtiennent sous forme paramétrique par
$$ (r,s) = (r_0, s_0) + \lambda(t, u) = (r_0+\lambda t, s_0+\lambda u) \leqno{(87')} $$
Ayant choisi $r,s$ donc $f = \begin{pmatrix} r & s \\ t & u \end{pmatrix} \in \operatorname{Sl}(2, \hat{\mathbf{Z}})$,
on en déduit $f \cdot {}^t\!f = \begin{pmatrix} R & S \\ T & U \end{pmatrix}$ (où $R = r^2+s^2$, $S=T=rt+su$, $U=t^2+u^2$), et on détermine successivement les quantités
$$
\begin{aligned}
D &= \frac{1}{\mu} = \varepsilon_2(t^2+u^2) = \varepsilon_2 U \\
C &= \varepsilon_2 T \ (= \varepsilon_2 S) \\
B &= C+D-1 = \varepsilon_2(S+U)-1 \\
A &= \frac{1}{D}(1+BC) = \frac{1}{\varepsilon_2 U} (1+S(S+U)-\varepsilon_2 S)
\end{aligned} \leqno{(88)}
$$
d'où une matrice $\begin{pmatrix} A & B \\ C & D \end{pmatrix}$ satisfaisant $C+D-B=1$ et $D \in \hat{\mathbf{Z}}^*$ et satisfaisant la relation (57)
$$ \begin{pmatrix} A+2\nu C & B+2\nu D \\ C & D \end{pmatrix} = \varepsilon_2 \begin{pmatrix} R & S \\ T & U \end{pmatrix} $$
qui se réduit ici à la relation du coin sup. gauche
$$ A+2\nu C = \varepsilon_2 R \quad \text{i.e.} $$

\newpage

%% ===== Page 165 =====

il reste à choisir $a,b$ de façon à satisfaire aux
deux relations $\det \begin{pmatrix} a & b \\ c & d \end{pmatrix} = \varepsilon_3$ et $C = \varepsilon_2 T$ (cf (61) (62))
$$
\begin{cases}
d \underline{a} - c \underline{b} = \varepsilon_3 \\
(c+d)\underline{a} + d \underline{b} = \varepsilon_2 \varepsilon_3 (\underbrace{rt+su}_{S}) + (\underbrace{c^2+cd+d^2}_{1/\mu \varepsilon_3})
\end{cases} \leqno{(88)}
$$
qui est un système de deux équations linéaires
en $\underline{a}, \underline{b}$ , dont le déterminant est connu
par hasard $c^2+cd+d^2 = 1/\mu \varepsilon_3 \in \hat{\mathbf{Z}}^\times$ - ce
qui assure une solution unique !

Il semble plus naturel peut être de
partir du multiplicateur $\mu$, et de chercher
les systèmes $\begin{pmatrix} a & b \\ c & d \end{pmatrix} \begin{pmatrix} r & s \\ t & u \end{pmatrix}$ qui y
correspondent, mais en mettant l'accent
sur $(c,d) \ (t,u)$ en premier lieu.
A l'aide de $\mu$ on détermine $\varepsilon_2, \varepsilon_3$ par
$$
\mu \equiv \varepsilon_2 \ (2) \ , \ \mu \equiv \varepsilon_3 \ (3) \ ,
$$
[unclear: procédé] qui permet de représenter
$1/\mu$ sous la forme
$$
1/\mu = \varepsilon_2 (t^2+u^2) = \varepsilon_3 (c^2+cd+d^2)
$$
(avec $t,u, c,d \in \hat{\mathbf{Z}}$ , les systèmes $(t,u)$ comme
$(c,d)$ "dépendant de 1 paramètre", 
je serais très content qu'on [unclear: puisse y voir plus clair dans ces]
[unclear: précisions] \& y [unclear: introduire] avec des
paramètres $\sigma, \tau$ explicites \dots Une
fois $c,d$, et $t,u$ choisis, on choisit $r,s$ comme il a

\newpage

%% ===== Page 166 =====

été dit (avec un paramètre (linéaire))
ce qui détermine ensuite $r,s$ de façon
unique, d'une façon qui fait intervenir
ce paramètre de façon linéaire-affine.
En fait, pour $c,d,t,u$ fixés, et $\varepsilon_2, \varepsilon_3$,
satisfaisant : $\varepsilon_2(t^2+u^2) = \varepsilon_3(c^2+cd+d^2)$, les trois équations
qui relient $a,b,r,s$ sont linéaires-affines

\marginpar{on a souligné les "coefficients" inconnus $\underline{a}, \underline{b}, \underline{r}, \underline{s}$ ; le rang de ce système est $3$, l'ens. des solutions est un espace affine de dim $1$ sur $\hat{\mathbf{Z}}$...}

$$
\begin{cases}
d\underline{a} - c\underline{b} = \varepsilon_3 \\
u\underline{r} - t\underline{s} = 1 \\
(c+d)\underline{a} + d\underline{b} - \varepsilon_2 \varepsilon_3 (u\underline{s} + t\underline{r}) = c^2+cd+d^2
\end{cases}
\leqno (89)
$$

Remarque. On n'a pas [unclear: jusqu'à présent]
tenu compte des conditions de congruences
$$
\begin{pmatrix} a & b \\ c & d \end{pmatrix} \equiv 1 \quad (2) \qquad \begin{pmatrix} r & s \\ t & u \end{pmatrix} \equiv 1 \quad (3)
$$
qui expriment que ces matrices proviennent
à vrai dire par "réduction" d'éléments de
$\hat{\mathfrak{T}}_{0,3}$, $\hat{\mathfrak{T}}_{1,1}$. Pour les quantités $(c,d)$, $\varepsilon_3$ on
trouve les conséquences
$$
\begin{cases} d \equiv \pm 1 \ (4) \text{ donc } d^2 \equiv 1 \ (4) \\ c \equiv 0, 2 \ (4) \text{ donc } c^2 \equiv 0 \ (4) \end{cases} \implies c^2+d^2+cd \equiv 1+cd \ (4) \equiv 1+c \ (4)
$$
donc
$$
\mu = \varepsilon_3(c^2+cd+d^2) \equiv \varepsilon_3(1+c) \ (4)
$$
($NB$ $d$ est inversible $\bmod 4$ puisque $d \equiv 1 \ (2)$).

\newpage

%% ===== Page 167 =====

On voit donc que la connaissance de $c \bmod 4$
et celle du signe $ad-bc = \varepsilon_3$, permet de
trouver le signe $\varepsilon_2$ par
$$
\varepsilon_2 \equiv \varepsilon_3(1+c) \ (4) \qquad \text{i.e.} \qquad \varepsilon_2 \varepsilon_3 \equiv 1+c \ (4) \leqno(90)
$$

Si on se propose, pour $c,d,t,u$ fixés,
satisfaisant
$$
\begin{cases}
c \equiv t \equiv 0 \ (2) \qquad d \equiv u \equiv 1 \ (2) \\
(c^2+cd+d^2) = \varepsilon(t^2+tu+u^2) \in \hat{\mathbf{Z}}^*
\end{cases} \leqno(91)
$$
où $\varepsilon \in \{\pm 1\}$ est le signe défini par
$$
\varepsilon \equiv 1+c \ (4)
$$
($\varepsilon=1$ si $c \equiv 0 \ (4)$, $\varepsilon=-1$ si $c \equiv 2 \ (4)$), de trouver
les $a,b,r,s$ correspondants, satisfaisant
les conditions (89), et de plus les
conditions de congruence
$$
b \equiv s \equiv 0 \ (2) \qquad a \equiv r \equiv 1 \ (2), \leqno(92)
$$
on est ramené à relever une solution déjà
donnée mod $2$ de ce système d'équations,
en une solution. En termes d'un torseur
$T$ des solutions, ce signifie qu'on a
donné un élément du $\mathbf{Z}/2\mathbf{Z}$-torseur quotient
correspondant — l'un de ses relèvements s'identifie
à un torseur sous $2\hat{\mathbf{Z}} \simeq \hat{\mathbf{Z}}$ \dots Aucun
mystère !

\newpage

%% ===== Page 168 =====

Remarque Le groupe additif $\hat{\mathbf{Z}}$ opère sur l'ens. des solutions $(\alpha_0, f) = \left( \begin{pmatrix} a & b \\ c & d \end{pmatrix}, \begin{pmatrix} r & s \\ t & u \end{pmatrix} \right)$ des équations envisagées (48)(49)(56)
$$ \det \alpha_0 \in \{\pm 1\}, \det f = 1, \alpha_0 \rho(\alpha_0)^{-1} = \varepsilon \, l_0 \big( f v_0 \rho(f)^{-1} \big) \ , \ \varepsilon \in \{\pm 1\}, \leqno{(93)} $$
par
$$ (\alpha_0, f) \longmapsto (\alpha'_0, f') = (l_0^n \alpha_0, l_0^n f) \leqno{(94)} $$
$$ = \left( \begin{pmatrix} a - 2nc & b - 2nd \\ c & d \end{pmatrix}, \begin{pmatrix} r - 2nt & s - 2nu \\ t & u \end{pmatrix} \right) $$
$(n \in \hat{\mathbf{Z}})$.

Au facteur 2 près dans $2n$, c'est justement l'indétermination dans l'ensemble des solutions de (93) lorsque $c, d, t, u$ sont fixés (cf (87) et propriété que [unclear: la matrice solution est dans $\GL_2$]). On peut se demander s'il y a un moyen de déterminer intrinsèquement (en termes de l'action de $\widehat{\Pi}_{0,3}$ comme [unclear: groupe d'automorphisme]) :
\begin{enumerate}
\item[a)] le multiplicateur $\mu \in \hat{\mathbf{Z}}^*$
\item[b)] les quantités $c, d, t, u \in \hat{\mathbf{Z}}$
\end{enumerate}
satisfaisant la surprenante relation (86)
$$ \varepsilon_2 (t^2+u^2) = \varepsilon_2 (c^2+cd+d^2) = 1/\mu \leqno{(96)} $$
et les relats. de congruences
$$ c \equiv t \equiv 0 \pmod{2} \quad d \equiv u \equiv 1 \pmod{2} \leqno{(97)} $$

Il est vrai que l'on a vu qu'il y a un moyen (dans $\widehat{\Pi}_{0,3}$) de normaliser le choix de $f$, de façon à lever l'ambiguïté des $l_0^n f$ (cf § 91 où c'était effleuré - j'y reviendrai encore par la suite...) mais il me semble [unclear: rationnel]...

\newpage

%% ===== Page 169 =====

pour le moment d'ignorer cette possibilité de normalisation -- ce qui met en évidence, en effet, la conviction cruciale, dans tous ces calculs héroïques, du rôle joué par les quatre quantités $c,d,t,u \in \hat{\mathbf{Z}}$.

\newpage

%% ===== Page 170 =====

\chapter*{\S \space 45. --- L'HOMOMORPHISME $\tilde{\mathcal{M}}_{0,3} \to \operatorname{Gl}(2, \hat{\mathbf{Z}})$, CONSTRUCTION DE L'EXTENSION $\tilde{\mathcal{N}}_{1,1}$ DE $\boldsymbol{\Gamma}$ PAR $\operatorname{Sl}(2, \hat{\mathbf{Z}})^{\wedge}$}\thispagestyle{empty}
\addcontentsline{toc}{section}{45. L'homomorphisme $\tilde{\mathcal{M}}_{0,3} \to \operatorname{Gl}(2, \hat{\mathbf{Z}})$, construction de l'extension $\tilde{\mathcal{N}}_{1,1}$ de $\boldsymbol{\Gamma}$ par $\operatorname{Sl}(2, \hat{\mathbf{Z}})^{\wedge}$}
\label{sec:45}
\section*{}

On désigne par $\tilde{\mathcal{M}}_{0,3}$\footnote{On désignera par $\tilde{\mathcal{M}}_{0,3}$ le sous-groupe de $\mathcal{M}_{0,3}$ (lequel s'identifie à l'extension [unclear: ( )] de $\mathfrak{S}_3$ par $\hat{\Pi}_{0,3}$) formé de ceux qui opèrent extérieurement à lacets sur $\hat{\Pi}_{0,3}$ et induisent l'identité sur son centre (qui est $\hat{\mathbf{Z}}$).} le groupe des automorphismes à lacets de $\hat{\mathfrak{S}}_{0,3}^+$ qui induisent sur le quotient $\mathfrak{S}_3 \simeq \hat{\mathfrak{S}}_{0,3}^+/\hat{\Pi}_{0,3}^+$ l'action triviale --- on a l'extension
$$
1 \to \hat{\Pi}_{0,3}^+ \to \tilde{\mathcal{M}}_{0,3} \to \tilde{\mathcal{N}}_{0,3} \to 1 , \leqno{(1)}
$$
$\tilde{\mathcal{N}}_{0,3}$ désigne le groupe des automorphismes à lacets de $\hat{\mathfrak{S}}_{0,3}^+$ qui induisent l'identité dans $\mathfrak{S}_3$. On notera que
$$
\boldsymbol{\Gamma} = \Gamma_{\mathbf{Q}} \hookrightarrow \tilde{\boldsymbol{\Gamma}} ,
$$
dont l'image inverse de $\Gamma_{\mathbf{Q}}$ dans l'ext. (1) s'identifie à $\mathcal{M}_{0,3}$ --- mais pour le moment je n'ai pas de caractérisation bien plausible à proposer du sous-groupe $\Gamma_{\mathbf{Q}}$ de $\tilde{\boldsymbol{\Gamma}}$. On a un scindage canonique de l'extension (1), via le sous-groupe $\Gamma_{\mathbf{Q}}$ de $\tilde{\mathcal{M}}_{0,3}$ des éléments de $\tilde{\mathcal{M}}_{0,3}$ qui commutent à $\mathfrak{S}_3$.

\newpage

%% ===== Page 171 =====

de sorte que $\mathcal{M}_{0,3}^{\sim}$ se présente comme un produit 1/2 direct

$$
\mathcal{M}_{0,3}^{\sim} \simeq \underset{\simeq \operatorname{Sl}(2,\hat{\mathbf{Z}})^{\wedge}}{\hat{\mathfrak{S}}_{0,3}^{+}} \cdot \tilde{\Gamma}_{\sigma} \leqno{(3)}
$$

Considérons l'homomorphisme canonique
$$
\hat{\mathfrak{S}}_{0,3}^{+} \simeq \operatorname{Sl}(2,\hat{\mathbf{Z}})^{\wedge} \xrightarrow{\varphi_0'} \operatorname{Sl}(2,\hat{\mathbf{Z}})' \subset \operatorname{Gl}(2,\hat{\mathbf{Z}})' \leqno{(4)}
$$
et [unclear: l'homomorphisme]
$$
\begin{aligned}
\tilde{\Gamma}_{\sigma} &\xrightarrow{\varphi_{\sigma}'} \operatorname{Gl}(2,\hat{\mathbf{Z}})' \\
u &\longmapsto \mu \begin{pmatrix} u & \varepsilon t \\ -t & \varepsilon u \end{pmatrix}
\end{aligned} \leqno{(5)}
$$
où $\mu \in \hat{\mathbf{Z}}^*$, $t, u \in \hat{\mathbf{Z}}$, $\varepsilon \in \{\pm 1\}$, $\alpha = \alpha_2 \in \mathbf{Z} / 2\mathbf{Z}$ sont les quantités [unclear: composantes] de $u \in \tilde{\Gamma}_{\sigma}$ envisagées au \S précédent, avec
$$
\varepsilon = (-1)^\alpha.
$$
S'il y a lieu, on écrira $t(u), u(u), \varepsilon(u), \alpha(u)$ pour indiquer la dépendance par rapport à $u$.

Considérons l'application (continue)
$$
\varphi' : \mathcal{M}_{0,3}^{\sim} \longrightarrow \operatorname{Gl}(2,\hat{\mathbf{Z}})' \leqno{(6)}
$$
définie par
$$
\varphi'(a.u) = \varphi_0'(a) \varphi_{\sigma}'(u) \qquad \begin{cases} a \in \hat{\mathfrak{S}}_{0,3}^{+} \\ u \in \tilde{\Gamma}_{\sigma} \end{cases} \leqno{(7)}
$$

\newpage

%% ===== Page 172 =====

\marginpar{385}Théorème. L'application $\varphi'$ est un homomorphisme de groupes.

Cela signifie, compte tenu que $\varphi'_0$ est déjà un hom. de groupes, le théorème signifie que

\begin{enumerate}
\item[a)] $\varphi'_\varpi$ est un hom. de groupes, i.e.
$$
\varphi'_\varpi(\underline{u}' \underline{u}) = \varphi'_\varpi(\underline{u}') \varphi'_\varpi(\underline{u}) \quad (\underline{u}', \underline{u} \in \widetilde{\Gamma}_\varpi) \leqno (8)
$$

\item[b)] on a, pour $a \in \widehat{\mathfrak{S}}_{0,3}^+$, $\underline{u} \in \widetilde{\Gamma}_\varpi$
$$
\varphi'_\varpi(\underbrace{a \underline{u} a^{-1}}_{\underline{u}(a)}) = \varphi'_0(a) \varphi'_\varpi(\underline{u}) \varphi'_0(a)^{-1} \leqno (9)
$$
\end{enumerate}

Mais la formule (9) est un fait déjà établi, puisque la construction même de la matrice
$$
\varphi'_\varpi(\underline{u}) = \mu (t \sigma_\varpi + u) \tau_\varpi^\varepsilon
$$
au $\S$ précédent était faite de façon à ce que cette formule soit vérifiée (cf (67)(68), et le cor. au th. (77)...). Pour vérifier (8), nous allons d'abord expliciter les quantités $g, h, f$ pour un produit de deux opérateurs $\underline{u}', \underline{u} \in \widetilde{\mathcal{M}}_{0,3}$, en termes de leurs valeurs pour $\underline{u}', \underline{u}$ eux-mêmes. Un calcul immédiat donne, en un sens, des "relations de cocycle"
$$
\begin{cases}
g(\underline{u}' \underline{u}) = g(\underline{u}') \cdot g(\underline{u}) \\
h(\underline{u}' \underline{u}) = \underline{u}'(h(\underline{u})) \cdot h(\underline{u}') \\
f(\underline{u}' \underline{u}) = f(\underline{u}') \cdot \underline{u}'(f(\underline{u}))
\end{cases} \leqno (10)
$$

\newpage

%% ===== Page 173 =====

Remarques

1) On aura d'ailleurs de même, avec les notations du $\S$ précédent
$$
\begin{cases}
\alpha(u'u) = u'(\alpha(u)). \alpha(u') \\
\beta(u'u) = u'(\beta(u)). \beta(u')
\end{cases} \leqno{(11)}
$$

2) La formule (10) sera un peu plus sympa pour $f$, que pour $g$ et $h$, si on avait envie de changer encore des notations, et de remplacer également $g$ et $h$ par $g^{-1}, h^{-1}$ (comme précédemment on [unclear: l'avait fait pour] $f^{-1}$ ...).

Désignons encore par les mêmes lettres $g(u), h(u)$ etc les images de ces quantités dans $\operatorname{Sl}(2,\hat{\mathbf{Z}})'$, les formules analogues dans $\operatorname{Sl}(2,\hat{\mathbf{Z}})'$ (ou $\operatorname{Gl}(2,\hat{\mathbf{Z}})'$) pour (10) et les (11) qui s'écrivent
$$
\begin{cases}
g(\underline{u'u}) = \operatorname{int}(\varphi'_0(\underline{u'}))(g(\underline{u})) . g(\underline{u'}) \\
h(\underline{u'u}) = \operatorname{int}(\varphi'_0(\underline{u'}))(h(\underline{u})) . h(\underline{u'}) \\
f(\underline{u'u}) = f(\underline{u'}) . \operatorname{int}(\varphi'_0(\underline{u'}))(f(\underline{u})) \\
\alpha(\underline{u'u}) = \operatorname{int}(\varphi'_0(\underline{u'}))(\alpha(\underline{u})) . \alpha(\underline{u'}) \\
\beta(\underline{u'u}) = \operatorname{int}(\varphi'_0(\underline{u'}))(\beta(\underline{u})) . \beta(\underline{u'})
\end{cases} \leqno{(12)}
$$
dans $\operatorname{Gl}(2,\hat{\mathbf{Z}})'$.

Ces formules dans $\operatorname{Gl}(2,\hat{\mathbf{Z}})'$ ont besoin en outre d'être explicitées du point de vue calcul effectif, contrairement

\newpage

%% ===== Page 174 =====

au cas des formules initiales dans $\pi_{0,3}^{\wedge\wedge}$ ou $\pi_{0,3}^\wedge$,
où des quantités comme $w(\underline{g}(\underline{u}))$ ne sont pas
explicitables à l'aide de formules rel. simples
en fonction des données $\underline{g}(\underline{u})$ etc qui déterminent $\underline{u}$ et $\underline{u}'$.
Ces formules se complètent, pour mémoire, par

$$ \mu(\underline{u}' \underline{u}) = \mu(\underline{u}') \mu(\underline{u}) \leqno{(13)} $$

impliquant notamment

$$ \varepsilon_2(\underline{u}' \underline{u}) = \varepsilon_2(\underline{u}') \varepsilon_2(\underline{u}) \leqno{(14)} $$

Pour faire le calcul de $\varphi_0(\underline{u}' \underline{u})$ et
vérifier qu'il s'accorde à $\varphi_0(\underline{u}') \varphi_0(\underline{u})$, on
introduira les coefficients $r, s, t, u$
de $f = f(\underline{u})$, et $r', s', t', u'$ de $f' = f(\underline{u}')$, et
$\mu, \mu', \mu''$, $r'', s'', t'', u''$ de $f'' = f(\underline{u}' \underline{u})$, avec $\varepsilon, \varepsilon', \varepsilon''$,
d'où

$$ \mu'' = \mu \mu' \; , \; \varepsilon'' = \varepsilon \varepsilon' \leqno{(15)} $$

et la formule en $f$ dans (12) s'écrit alors

$$ f'' = f' \underline{w}' f \underline{w}'^{-1} \leqno{(16)} $$

On écrit, comme au \S précédent,
$\underline{w}'$ au lieu de $\varphi_0(\underline{u}')$, soit en termes
explicites
\marginpar{NB $\det \underline{w}' = \mu'$}
$$ \begin{pmatrix} r'' & s'' \\ t'' & u'' \end{pmatrix} = \begin{pmatrix} r' & s' \\ t' & u' \end{pmatrix} \begin{pmatrix} \mu' & \varepsilon' t' \\ -t' & \varepsilon' \mu' \end{pmatrix} \begin{pmatrix} r & s \\ t & u \end{pmatrix} \begin{pmatrix} \mu' & -\varepsilon' t' \\ t' & \mu' \end{pmatrix} $$

\newpage

%% ===== Page 175 =====

Le produit des trois derniers termes du second membre n'est pas engageant :
$$
\underline{u}' \underline{f} \underline{u}'^{-1} = \mu' \begin{pmatrix} \varepsilon' r u'^2 + s t' u' + t t' u' + \varepsilon' u t'^2 & -\varepsilon' r t' u' + s u'^2 - t t'^2 + \varepsilon' u' u t' \\ -\varepsilon' r u' t' - s t'^2 + t u'^2 + \varepsilon' u' u t' & -\varepsilon' r t'^2 - s u' t' - t u' t' + \varepsilon' u u'^2 \end{pmatrix}
$$
mais le produit par $\underline{f}'$ se simplifie miraculeusement, grâce à la relation :
$$
t'^2 + u'^2 = \varepsilon' / \mu'
$$
et on trouve
$$
\underline{f}'' = \begin{pmatrix} r'' & s'' \\ t'' & u'' \end{pmatrix} = \frac{1}{\mu'} \begin{pmatrix} * \text{bordel} & \text{bordel} \\ \varepsilon' t u' + u t' & -\varepsilon' t t' + u u' \end{pmatrix}
$$
d'où
$$
t'' = \frac{1}{\mu'} (\varepsilon' t u' + u t') \quad \quad u'' = \frac{1}{\mu'} (-\varepsilon' t t' + u u') \leqno (18)
$$
(avec $\mu'' = \mu \mu'$) et par suite (tenant compte de $\varepsilon'' = \varepsilon \varepsilon'$)
$$
\underline{u}'' = \mu \mu' \begin{pmatrix} -\varepsilon' t t' + u u' & \varepsilon t u' + \varepsilon \varepsilon' u t' \\ - \varepsilon' t u' - u t' & -\varepsilon \varepsilon' t t' + \varepsilon \varepsilon' u u' \end{pmatrix}
$$
et la formule à vérifier
$$
\underline{u}'' = \underline{u}' \underline{u}
$$
se réduit à
$$
\mu \mu' \begin{pmatrix} - \varepsilon' t t' + u u' & \varepsilon t u' + \varepsilon \varepsilon' u t' \\ - \varepsilon' t u' - u t' & -\varepsilon \varepsilon' t t' + \varepsilon \varepsilon' u u' \end{pmatrix} = \mu' \begin{pmatrix} u' & \varepsilon' t' \\ -t' & \varepsilon' u' \end{pmatrix} \mu \begin{pmatrix} u & \varepsilon t \\ -t & \varepsilon u \end{pmatrix}
$$
qui est bien vraie, OK !

Ici encore, on sent que ces calculs devraient s'éliminer d'eux-mêmes, [unclear: tapis] d'algèbre linéaire que j'ai la flemme de dérouler in extenso. Il faudrait paraphraser la description de $M_{0,3}$ (ou plutôt $\hat{\mathfrak{T}}_{0,3}$) décrit par un système $(\rho_1, \rho_2)$ de relations [unclear: du groupe]...

\newpage

%% ===== Page 176 =====

dans $\operatorname{Gl}(2, \hat{\mathbb{Z}})/_{\pm 1}$, puis on remplacera $\hat{\mathbb{Z}}$ par
un anneau quelconque, et bien sûr
$\operatorname{Gl}(2, \hat{\mathbb{Z}})'$ par $\operatorname{Gl}(2, \hat{\mathbb{Z}})$ lui-même. Peut-
être y reviendrai-je finalement ....

Ayant ainsi construit
$$ \varphi : \widetilde{M}_{0,3} \longrightarrow \operatorname{Gl}(2, \hat{\mathbb{Z}})' = \operatorname{Gl}(2, \hat{\mathbb{Z}})/\pm 1 $$
on déduit, par image inverse de l'extension
centrale $\operatorname{Gl}(2, \hat{\mathbb{Z}})$ de $\operatorname{Gl}(2, \hat{\mathbb{Z}})'$ par $\{\pm 1\}$, une
extension centrale de $\widetilde{M}_{0,3}$ par $\{\pm 1\}$,
que je note $\widetilde{N}_{1,1}$, et on a
un hom d'extensions

(19)
\begin{tikzcd}
1 \ar[r] & \{\pm 1\} \ar[r] \ar[d, equal] & \widetilde{N}_{1,1} \ar[r] \ar[d, "\varphi"] & \widetilde{M}_{0,3} \ar[r] \ar[d, "\varphi'"] & 1 \\
1 \ar[r] & \{\pm 1\} \ar[r] & \operatorname{Gl}(2, \hat{\mathbb{Z}}) \ar[r] & \operatorname{Gl}(2, \hat{\mathbb{Z}})' \ar[r] & 1
\end{tikzcd}

[Paragraphe raturé au bas de la page]

\newpage

%% ===== Page 177 =====

Il faudrait vérifier que l'hom $\varphi'$
qu'on vient de construire ci-dessus sur
$\mathcal{M}_{0,3} \subset \tilde{\mathcal{M}}_{0,3}$ ... est l'hom. "modulaire"
bien connu du § 39 (?) - il suffit
de le vérifier sur le sous-groupe
$\Gamma_{Q,\sigma}$ de $\tilde{\Gamma}_Q$ correspondant à $\Gamma_Q$ -
ce ne doit pas être bien difficile, et
je vais l'admettre ici. Ceci étant
admis, on trouve que l'image
inverse de $\mathcal{M}_{0,3}$ dans $\tilde{\mathcal{N}}_{1,1}$ n'est
autre (à isom. can. près ) que
$\mathcal{N}_{1,1}$ lui-m. - ou tout au moins l'
image inverse du [unclear: crossed out] n'est
autre que $\widehat{C}_{1,1}^+ \simeq \operatorname{Sl}(2,\mathbb{Z})^\wedge$. On aura alors
(avec cette identification)

(20)
\begin{tikzcd}
1 \ar[r] & \widehat{C}_{1,1}^+ \ar[r] \ar[d] & \tilde{\mathcal{N}}_{1,1} \ar[r] \ar[d, "\varphi"] & \tilde{\Gamma} \ar[r] \ar[d, "det"] & 1 \\
1 \ar[r] & \operatorname{Sl}(2,\hat{\mathbb{Z}}) \ar[r] & \operatorname{Gl}(2,\hat{\mathbb{Z}}) \ar[r] & \hat{\mathbb{Z}}^\times \ar[r] & 1
\end{tikzcd}

\newpage

%% ===== Page 178 =====

donc on a reconstruit en termes d'autom[orphisme]s
à l'aide de $\hat{\Pi}_{0,3}$ le diagramme (1) du
§ 39, - plutôt un diagramme un
peu plus général, dans lequel se plonge
celui du § 39 - il resterait encore à
caractériser $\Gamma_Q$ comme sous-groupe de $\tilde{\Gamma}$.

Le diagramme déduit directement du
th. précédent

(21)
\begin{tikzcd}
1 \ar[r] & \hat{C}_{0,3}^+ \ar[r] \ar[d] & \tilde{\mathcal{M}}_{0,3} \ar[r] \ar[d, "\phi'"] & \tilde{\Gamma} \ar[r] \ar[d, "\bar{\chi}"] & 1 \\
1 \ar[r] & \operatorname{Sl}(2,\hat{\mathbb{Z}})' \ar[r] & \operatorname{Gl}(2,\hat{\mathbb{Z}})' \ar[r, "\det"'] & \hat{\mathbb{Z}}^\times \ar[r] & 1
\end{tikzcd}

peut être considéré comme étant déduit
du précédent par réduction mod les
[unclear: sous-groupes centraux] $\{\pm 1\}$ dans les deux lignes.

Que devient dans (20) le scindage
(13) de $\tilde{\mathcal{M}}_{0,3}$ ? Y a-t-il une façon
privilégiée, pour les éléments de
$\tilde{\mathcal{N}}_{1,1}$ - i.e. un couple d'un
$u = (g,h,f) \in \tilde{\mathcal{M}}_{0,3}$, et d'un $u_0 \in \operatorname{Gl}(2,\hat{\mathbb{Z}})$

\newpage

%% ===== Page 179 =====

qui relève $\underline{u}_\infty = \varphi'_\infty(u)$ de $\operatorname{Gl}(2,\hat{\mathbf{Z}})'$,
de le ``normaliser'' mod. composition
avec un élément de $\operatorname{Sl}(2,\hat{\mathbf{Z}})^\sim$, de
façon à ce qu'il ait une (disons...)
propriété de commutation avec $\sigma_\infty$, du
type $u(\sigma_\infty) = \sigma_\infty$, ou $u(\sigma_\infty) = \pm \sigma_\infty$ ?

Une autre façon de poser la question,
en termes du sous-groupe $\boldsymbol{\Gamma}^\sim_{\mathbf{Q}} \subset \mathcal{M}_{0,3}^\sim$,
est la suivante : y-a-t-il un relève-
ment canonique $\boldsymbol{\Gamma}^\sim_{\mathbf{Q}} \to \mathcal{N}_{1,1}^\sim$,
associant à tout $u = (g,f) \in \boldsymbol{\Gamma}^\sim_{\mathbf{Q}}$
(d'où un [unclear: $\underline{u} = \dots \in \operatorname{Gl}(2,\hat{\mathbf{Z}})'$],
soit i.e.\ un $\underline{u}$ défini mod. signe)
un représentant $u_0 \in \operatorname{Gl}(2,\hat{\mathbf{Z}})$ - i.e.
une façon de lever cette ambiguïté
de signe ? En fait, si on remplace
$\boldsymbol{\Gamma}^\sim_{\mathbf{Q}}$ par $\Gamma_{\mathbf{Q}, \infty}$, le yoga
standard suggère que les relèvements en

\newpage

%% ===== Page 180 =====

question doivent correspondre aux courbes
elliptiques sur $\mathbf{Q}$, munies d'une rigidifica-
tion des points d'échelon 2, correspondant
à l'isomorphisme $\mathcal{M}_{1,1,\mathbf{Q}} \simeq \mathcal{M}_{0,3,\mathbf{Q}}$
au point $1/2$ de $\mathcal{M}_{0,3}(\mathbf{Q})$. L'indétermina-
tion pour le choix d'une telle courbe
elliptique est justement dans $H^1(\mathbf{Q}, \{\pm 1\})$
(c-à-d un [unclear: choix] d'ext. quadratiques de $\mathbf{Q}$)
qui s'identifie à $\operatorname{Hom}(\boldsymbol{\Gamma}_{\mathbf{Q}}, \{\pm 1\}) -$
ce qui donne bien la même indétermina-
tion pour un relèvement. Ceci prouve
essentiellement que l'idée suggérée par
le yoga général est bel et bien
correcte - et cela montre en même temps
qu'il n'y a pas (même sur $\boldsymbol{\Gamma}_{\mathbf{Q}, S}$) de
choix vraiment privilégié. Qu'il y ait
(sur $\boldsymbol{\Gamma}_{\mathbf{Q}, S}$)
de tels relèvements provient du fait
qu'il y a bel et bien des courbes,

\newpage

%% ===== Page 181 =====

elliptiques répondant à la question, la plus simple étant donnée par l'équation
$$
y^2 - x(x-1)(x-1/2) = 0 \leqno{(22)}
$$
et voilà le choix fait ! Il resterait à essayer de l'expliciter un peu dans un formulaire calculatoire, si ce n'est pas sans espoir. Il faudrait que j'examine cette question de plus près, à l'occasion.

\newpage

%% ===== Page 182 =====

\section*{45 Notations définitives (?) - nouvelle récapitulation}
(pour les autom. à lacets de $\hat{\Pi}_{0,3}$ et de $\mathfrak{S}_{0,3}^{+\wedge} \simeq \operatorname{Sl}(2, \hat{\mathbf{Z}})^{\sim}$)

\subsection*{I Les groupes $\tilde{\mathcal{M}}_{0,3}$ et $\tilde{\mathcal{N}}_{0,3}$}

On pose\footnote{NB $\tilde{\mathcal{M}}_{0,3}$ est aussi le normalisateur de $\mathfrak{S}_{0,3}^{+\wedge}$ dans $\hat{\mathfrak{S}}_{0,3}$ et de $\hat{\Pi}_{0,3}$ (le normalisateur de $\hat{\Pi}_{0,3}$ dans $\hat{\mathfrak{S}}_{0,3}$ [unclear: est $\tilde{\mathcal{N}}_{0,3}$]).}
$$
\tilde{\mathcal{M}}_{0,3} = \text{groupe des autom. de } \mathfrak{S}_{0,3}^{+\wedge} \simeq \operatorname{Sl}(2, \hat{\mathbf{Z}})^{\sim} \leqno{(1)}
$$
qui invariant $\hat{\Pi}_{0,3}$, et qui respectent la « structure à lacets » i.e.\ satisfont $u(x_i) = \operatorname{int}(f_i)(x_i^{k_i})$, avec $f_i \in \mathfrak{S}_{0,3}^{+\wedge}$, $i \in \{0,1,\infty\}$, $k \in \hat{\mathbf{Z}}^*$ (où $f_i \in \hat{\Pi}_{0,3}$, et $f_i$ est déterminé alors dans $x_i \backslash \hat{\Pi}_{0,3}$)

$\simeq$ groupe des autom. ext. à lacets de $\hat{\Pi}_{0,3}$ qui commutent extérieurement à l'action extérieure naturelle de $\mathfrak{S}_3$.

On a une structure d'extension
$$
1 \to \hat{\Pi}_{0,3} \to \tilde{\mathcal{M}}_{0,3} \to \tilde{\mathcal{N}}_{0,3} \to 1 \leqno{(2)}
$$

$$
\tilde{\mathcal{N}}_{0,3} \defeq \tilde{\mathcal{M}}_{0,3} / \hat{\Pi}_{0,3} \leqno{(3)}
$$
$\simeq$ groupe des autom. ext. à lacets de $\hat{\Pi}_{0,3}$, commutant à $\mathfrak{S}_3$.

L'homomorphisme interne : $\mathfrak{S}_3 \to \operatorname{Aut}(\mathfrak{S}_3)$ est un isomorphisme
$$
\begin{tikzcd}
\mathfrak{S}_3 \arrow[r, "\sim", "x \mapsto \operatorname{int}(x)"'] & \operatorname{Aut}(\mathfrak{S}_3)
\end{tikzcd} \leqno{(4)}
$$

et l'homomorphisme naturel
$$
\tilde{\mathcal{M}}_{0,3} \to \operatorname{Aut}(\mathfrak{S}_3 = \mathfrak{S}_{0,3}^{+\wedge} / \hat{\Pi}_{0,3})
$$
définit alors un homomorphisme
$$
\tilde{\mathcal{M}}_{0,3} \to \mathfrak{S}_3 \leqno{(5)}
$$
qui peut aussi s'interpréter comme la composition
$$
\begin{tikzcd}
\tilde{\mathcal{M}}_{0,3} \arrow[r, hook] & \underset{\substack{\text{tous les autom.} \\ \text{à lacets de } \hat{\Pi}_{0,3}}}{\hat{\mathfrak{S}}_{0,3}} \arrow[r, "\text{action sur les}"] & \underset{\text{classes de lacets}}{\mathfrak{S}_3}
\end{tikzcd} \leqno{(6)}
$$

\newpage

%% ===== Page 183 =====

et qui est trivial sur $\Pi_{0,3}^\wedge$, donc se factorise en

$$(7) \quad \mathcal{N}_{0,3}^\sim \longrightarrow \mathfrak{S}_3$$

Les noyaux de (5), (7) sont notés $\mathcal{M}_{0,3}^{\sim !}$, $\mathcal{N}_{0,3}^{\sim !}$ respectivement, de sorte qu'on a un diagramme de suites exactes :

\begin{tikzcd}
& & \Pi_{0,3}^\wedge \ar[d] & & & \\
(8) & 1 \ar[r] & \mathcal{M}_{0,3}^{\sim !} \ar[r] \ar[d] & \mathcal{M}_{0,3}^\sim \ar[r] \ar[d] & \mathfrak{S}_3 \ar[r] \ar[d, equal] & 1 \\
(9) & 1 \ar[r] & \mathcal{N}_{0,3}^{\sim !} \ar[r] \ar[d] & \mathcal{N}_{0,3}^\sim \ar[r] & \mathfrak{S}_3 \ar[r] & 1 \\
& & 1 & & &
\end{tikzcd}

(9) pourra à son tour être considéré comme quotient de (8) par $\Pi_{0,3}^\wedge$ \quad Notons aussi la suite exacte correspondante

$$(9 \text{ bis}) \quad 1 \longrightarrow \Pi_{0,3}^\wedge \longrightarrow \mathcal{M}_{0,3}^\sim \longrightarrow \mathcal{N}_{0,3}^\sim \longrightarrow 1$$

On a un hom. canonique

$$(10) \quad \mathfrak{S}_3 \simeq \widehat{G_{0,3}^+} / \Pi_{0,3}^\wedge \longrightarrow \mathcal{N}_{0,3}^\sim$$

dont le composé avec (7) est l'identité, d'où une scission

$$(11) \quad \mathcal{N}_{0,3}^\sim \simeq \mathfrak{S}_3 \times \mathcal{N}_{0,3}^{\sim !}$$

D'autre part, on a un isomorphisme

\begin{tikzcd}
(12) & \Gamma_{0,3} \ar[r, "\sim"] \ar[d, "\simeq"'] & \mathcal{N}_{0,3}^\sim \\
& \mathfrak{S}_3 \times \mathbb{Z}/2\mathbb{Z} & 
\end{tikzcd}

qui sur le facteur $\mathfrak{S}_3$ ( $= \Gamma_{0,3}^+$ ) coïncide avec (10), et qui en outre applique l'élément non trivial $\tau$ de $\mathbb{Z}/2\mathbb{Z}$ sur l'élément, également noté $\tau$, du deuxième facteur $\mathcal{N}_{0,3}^{\sim !}$ de (11).

\newpage

%% ===== Page 184 =====

L'homomorphisme composé
(13) $$ \mathcal{M}_{0,3}^\sim \to \hat{\mathfrak{S}}_{0,3} \xrightarrow{\text{multiplicateur } \chi} \hat{\mathbb{Z}}^* $$
donne un hom.
(14) $$ \mathcal{M}_{0,3}^\sim \xrightarrow{\chi} \hat{\mathbb{Z}}^* $$
surjectif, se factorise par $\mathcal{N}_{0,3}^\sim$ et même $\mathcal{N}_{0,3}^\sim / \mathfrak{S}_3$
(15) 
\begin{tikzcd}
\mathcal{N}_{0,3}^\sim \ar[rr, "\chi"] \ar[dr] & & \hat{\mathbb{Z}}^* \\
& \mathcal{N}_{0,3}^\sim / \mathfrak{S}_3 \ar[ur] & \\
& = \mathcal{N}_{0,3}^1 &
\end{tikzcd}

et donne lieu à des structures d'extensions :
(16)
\begin{tikzcd}
1 \ar[r] & \mathcal{M}_{0,3}^{\sim +} \ar[r] \ar[d] & \mathcal{M}_{0,3}^\sim \ar[r, "\chi"] \ar[d] & \hat{\mathbb{Z}}^* \ar[r] \ar[d, equal] & 1 \\
1 \ar[r] & \mathcal{N}_{0,3}^{\sim +} \ar[r] & \mathcal{N}_{0,3}^\sim \ar[r, "\chi"] & \hat{\mathbb{Z}}^* \ar[r] & 1
\end{tikzcd}

où on a posé :
(17) $$ \begin{cases} \mathcal{M}_{0,3}^{\sim +} = \operatorname{Ker}(\mathcal{M}_{0,3}^\sim \to \hat{\mathbb{Z}}^*) \\ \mathcal{N}_{0,3}^{\sim +} = \operatorname{Ker}(\mathcal{N}_{0,3}^\sim \to \hat{\mathbb{Z}}^*) \end{cases} $$

- la deuxième ligne de (16) s'interprète comme quotient de la première ligne par $\hat{\Pi}_{0,3}$ .

On a un plongement canonique
(18) $$ \Gamma_{\mathbb{Q}} \stackrel{df}{=} \operatorname{Gal}(\bar{\mathbb{Q}}/\mathbb{Q}) \hookrightarrow \mathcal{N}_{0,3}^\sim $$
[MARGIN: clôture alg. de $\mathbb{Q}$ dans $\mathbb{C}$]
unique à [unclear: un autom. int. de conjugaison près]

\newpage

%% ===== Page 185 =====

complexe $\tau$ ou l'élément déjà noté $(\tau)$
(cf (11)). L'image inverse du sous-
groupe $\Gamma_{\mathbb{Q}}$ par $\mathcal{M}_{0,3}^{\sim} \to \mathcal{N}_{0,3}^{\sim} \simeq \mathcal{N}_{0,3}^{1 \sim} \rtimes \hat{\mathfrak{S}}_3$
i.e l'image inverse de $\hat{\mathfrak{S}}_3 \times \Gamma_{\mathbb{Q}} \subset$
(19) $$ \mathcal{N}_{0,3} \overset{dfn}{=} \hat{\mathfrak{S}}_3 \times \Gamma_{\mathbb{Q}} \subset \mathcal{N}_{0,3}^{\sim} = \hat{\mathfrak{S}}_3 \times \mathcal{N}_{0,3}^{1 \sim} $$
est [unclear] notée $\mathcal{M}_{0,3}$ par ailleurs :

[MARGIN: On a une [unclear] de $\mathcal{M}_{0,3}^{\sim}$ de (7)]
(20)
\begin{tikzcd}
1 \ar[r] & \Pi_{0,3}^{\sim} \ar[r] \ar[d] & \mathcal{M}_{0,3} \ar[r] \ar[d] & \hat{\mathfrak{S}}_3 \times \Gamma_{\mathbb{Q}} \ar[r] \ar[d, "(\alpha, \beta) \mapsto \beta"] & 1 \\
1 \ar[r] & \hat{\mathfrak{S}}_{0,3}^+ \ar[r] & \mathcal{M}_{0,3} \ar[r] & \Gamma_{\mathbb{Q}} \ar[r] & 1
\end{tikzcd}

Bien entendu, c'est par $\mathcal{M}_{0,3}$ et $\mathcal{N}_{0,3}$ plutôt que
par $\mathcal{M}_{0,3}^{\sim}$ et $\mathcal{N}_{0,3}^{\sim}$ que nous serons intéressés -
donc par le sous-groupe remarquable $\Gamma_{\mathbb{Q}}$ de
$\mathcal{N}_{0,3}^{1 \sim}$. Nous allons "débroussailler" un peu
le groupe $\mathcal{N}_{0,3}^{1 \sim}$ - et il s'agira dans la suite
du travail de trouver un maximum
de conditions nécessaires sur des éléments
de $\mathcal{N}_{0,3}^{1 \sim}$ pour provenir de $\Gamma_{\mathbb{Q}}$ - jusqu'à
trouver, si possible, une caractérisation (au
moins plausible, sinon prouvée?) des éléments
de $\Gamma_{\mathbb{Q}}$ - i.e. de $\Gamma_{\mathbb{Q}}$ lui-même.

\newpage

%% ===== Page 186 =====

II Détermination combinatoire des éléments de $\mathcal{N}_{0,3}^{! \sim}$.

1) Détermination par $h, g, f$.

Rappelons que
$$ \mathcal{N}_{0,3}^{! \sim} = \text{autom. de } \widehat{\mathfrak{S}_{0,3}^+} \simeq \operatorname{Sl}(2, \hat{\mathbf{Z}})^\sim \text{ qui stabilisent } \leqno{(21)} $$
$\hat{\pi}_{0,3}$, induisent l'identité sur $\mathfrak{S}_3 = \widehat{\mathfrak{S}_{0,3}^+} / \hat{\pi}_{0,3}$
et respectent la structure à bouts (définie par la classe de $\varepsilon_0$).

Un $u \in \mathcal{N}_{0,3}^{! \sim}$ est déterminé par un système
$$ (h, g, f) \qquad \left( h, g, f \in \hat{\pi}_{0,3} \atop f \text{ déterminé mod } \text{[unclear: } \varepsilon_0^n \text{]}, \ n \in \hat{\mathbf{Z}} \right) \leqno{(22)} $$
via les formules
$$ \begin{cases} u(\sigma_\infty) = h \sigma_\infty \\ u(p) = g p \end{cases} \leqno{(23)} $$
La condition que ces formules définissent un automorphisme (ou dumoins un endomorphisme) de $\widehat{\mathfrak{S}_{0,3}^+}$ (qui induise nécess. l'identité sur $\mathfrak{S}_3$) est exprimée par les formules
$$ (h \sigma_\infty)^2 = (g p)^3 = 1 \leqno{(24)} $$
qui s'écrivent aussi
$$ \begin{cases} h \sigma_\infty(h) = 1 \\ g p(g) p^2(g) = 1 \end{cases} \leqno{(25)} $$
et la condition qu'il respecte les bouts : l'un [unclear: d'eux] s'exprimant par\marginpar{Relations des bouts dans $\widehat{\mathfrak{S}_{0,3}^+}$ (forme brute)}
$$ u(\varepsilon_0) = f (\varepsilon_0^\mu) f^{-1} \qquad \mu \in \hat{\mathbf{Z}}^* \text{ convenable}, \leqno{(26)} $$
[unclear: équivaut à la condition :]

\newpage

%% ===== Page 187 =====

\marginpar{relations des \\ lacets dans $\hat{\Pi}_{0,3}$ \\ (forme brute)}
$$ u(l_0) = \operatorname{int}(f^{-1})(l_0^\mu) \leqno{(27)} $$

elle s'explicite aisément en termes de $h,g,f$ (sous forme (26) de préférence) compte tenu que $l_0 = \sigma_\infty p$

$$ (h \sigma_\infty) (\underline{g p}) = \operatorname{int}(f^{-1}) \varepsilon_0^\mu \quad \text{, i.e.} $$
$$ h \sigma_\infty(g) \underbrace{\sigma_\infty p}_{l_0} = f^{-1} \varepsilon_0^\mu f $$
$$ f \, h \, \sigma_\infty(g) \, l_0 \, f^{-1} = \varepsilon_0^\mu \quad \text{soit} $$
$$ \text{[unclear: } \varepsilon_0 (f) \varepsilon_0 = \sigma_\infty(p(f)^{-1}) \varepsilon_0 \text{]} $$
$$ f \, h \, \sigma_\infty(g \, p \, f \, p^{-1}) = \varepsilon_0^{\mu-1} \quad \text{soit, posant} $$
\marginpar{possible, car \\ $\mu \in \hat{\mathbf{Z}}^*$ donc \\ $\mu \equiv 1 \, (2)$}
$$ \nu \defeq \frac{\mu-1}{2} \in \hat{\mathbf{Z}} \leqno{(28)} $$

la relation (compte tenu que $\varepsilon_0^2 = l_0$)
$$ f \, h \, \sigma_\infty(g \, p \, f \, p^{-1}) = l_0^\nu \quad , $$
on trouve, appliquant $\sigma_\infty$, comme $\sigma_\infty(l_0) = l_1$
$$ \sigma_\infty(f) \sigma_\infty(h) g \, p \, f \, p^{-1} = l_1^\nu $$
soit enfin
\marginpar{$\nu \in \hat{\mathbf{Z}}$ \\ NB $\nu$ est uniquement \\ déterminé par $g, h, f$ et $\mu \in \hat{\mathbf{Z}}^*$ \\ car $2\nu+1$ est le multiplicateur $\dots$}
$$ \sigma_\infty(h) g = \sigma_\infty(f)^{-1} l_1^\nu p(f) \leqno{(29)} $$
(relation des lacets en $h, g, f$).

NB Il faudra de plus exiger, pour que $h, g, f$ définissent bien un autom., que [unclear: (24) est une partie] génératrice de $\hat{\mathfrak{S}}_{0,3}^+$ --- par la suite, nous sous-entendrons tjrs tacitement cette condition liminaire.

\newpage

%% ===== Page 188 =====

2) Détermination par $\alpha_0, \beta_0, f$.

Il y a correspondance biunivoque entre les couples $(g,h)$ satisfaisant (25), et les couples
$$
(\alpha_0, \beta_0) \quad (\alpha_0 \in \hat{\Pi}_{0,3}^{\sim} = \hat{\Pi}_{0,3} \setminus \{1, [unclear: \tau_\infty]\} , \beta_0 \in \hat{\Pi}_{0,3} )
\leqno{(31)}
$$
$$
\begin{cases} 
g = \alpha_0 \rho(\alpha_0)^{-1} = [\alpha_0, \rho] \\ 
h = \beta_0 \sigma_\infty(\beta_0)^{-1} = [\beta_0, \sigma_\infty] 
\end{cases}
\leqno{(32)}
$$

Ainsi, les $u \in \tilde{\mathcal{M}}_{0,3}^{1 \sim}$ correspondent aux systèmes
$$
(\alpha_0, \beta_0, f) \quad \alpha_0 \in \hat{\Pi}_{0,3}^{\sim} \text{, } \beta_0, f \in \hat{\Pi}_{0,3} \text{, } f \text{ déterm. mod } f \mapsto l_1^\nu f
\leqno{(33)}
$$
satisfaisant la condition (29), considérée comme condition en $\alpha_0, \beta_0, f$, qu'on y substitue à $g, h$ les valeurs données par (32). La relation s'écrit
$$
\sigma_\infty(\beta_0) \beta_0^{-1} \alpha_0 \rho(\alpha_0)^{-1} = \sigma_\infty(f)^{-1} l_1^\nu \rho(f)
$$
ou encore
$$
\sigma_\infty(f \beta_0) \underbrace{\beta_0^{-1} \alpha_0}_{(f \beta_0)^{-1} (f \alpha_0)} \rho(f \alpha_0)^{-1} = l_1^\nu
$$
ou encore
$$
(f \alpha_0) \rho(f \alpha_0)^{-1} = \left( (f \beta_0) \sigma_\infty(f \beta_0)^{-1} \right) l_1^\nu
\leqno{(34)}
$$
(relation des lacets en $\alpha_0, \beta_0, f$). Comme pour (29), $\nu \in \hat{\mathbf{Z}}$ est uniquement déterminé en termes de $\alpha_0, \beta_0$, par la possibilité de trouver $f$ satisfaisant à la relation (34).

L'avantage de l'introduction de $\alpha_0, \beta_0$ au lieu

\newpage

%% ===== Page 189 =====

de $g, h$ est de rendre inutiles (car automatiques)
les relations (25), - il reste la seule relation
(34) à vérifier. Les relations (23) s'écrivent au demeurant :
$$ u(\sigma_\infty) = \operatorname{int}(\beta_0)(\sigma_\infty), \quad u(\rho) = \operatorname{int}(\alpha_0)(\rho) \leqno{(35)} $$

3) Détermination par $(\alpha, \beta, f)$.

La forme particulière de (34), où le triple $(\alpha_0, \beta_0, f)$
n'intervient que par l'intermédiaire du
couple correspondant $(f\alpha_0, f\beta_0)$, conduit à poser
$$ \alpha = f\alpha_0 , \quad \beta = f\beta_0 , \leqno{(35')} $$
on trouve que les $u \in \mathcal{M}_{0,3}^{1 \sim}$ sont déterminés
par les triples
$$ (\alpha, \beta, f)\footnote{Une dissymétrie est l'indétermination d'un tel $f$ (pour des $\alpha, \beta$ donnés), mais indétermination levée plus bas.} \qquad \alpha \in \hat{\Pi}_{0,3}, \beta \in \hat{\Pi}_{0,3}, f \in \hat{\Pi}_{0,3} \text{ [unclear: définis à mult près par } \alpha \mapsto l\alpha, \beta \mapsto l\beta, f \mapsto lf] \leqno{(36)} $$
satisfaisant à la condition
$$ (\alpha \rho(\alpha)^{-1}) = (\beta \sigma_\infty(\beta)^{-1}) \cdot l_1^\nu \leqno{(37)} $$
(relation des lacets en $\alpha, \beta$). L'avantage de
cette relation est que $f$ n'y figure plus.

On retrouve $\alpha_0, \beta_0, g, h$ en fonction de $\alpha, \beta, f$
par les formules
$$ \alpha_0 = f^{-1} \alpha , \quad \beta_0 = f^{-1} \beta \leqno{(38)} $$
$$ g = \alpha_0 \rho(\alpha_0)^{-1} = f^{-1}(\alpha \rho(\alpha)^{-1}) \rho(f) = f^{-1} \xi \rho(f) \leqno{(39)} $$
$$ h = \beta_0 \sigma_\infty(\beta_0)^{-1} = f^{-1}(\beta \sigma_\infty(\beta)^{-1}) \sigma_\infty(f) = f^{-1} \eta \sigma_\infty(f) , \leqno{(40)} $$
où on a posé
$$ \xi = \alpha \rho(\alpha)^{-1} , \quad \eta = \beta \sigma_\infty(\beta)^{-1} \qquad (\xi, \eta \in \hat{\Pi}_{0,3}) \leqno{(41)} $$

\newpage

%% ===== Page 190 =====

Les formules (23) à (35) déterminant $u$ en termes des données calculatoires, deviennent ici :
$$ u(\sigma_\infty) = \operatorname{int}(f^{-1} \beta)(\sigma_\infty) \qquad u(\rho) = \operatorname{int}(f^{-1} \alpha)(\rho) \leqno{(42)} $$
(42 bis) i.e.\ $u = \operatorname{int}(f^{-1}) u_{\alpha,\beta}$, où $u_{\alpha,\beta}$ ne dépend que de $\alpha, \beta$ :\footnote{La formule (42 ter) qui suit est raturée dans le manuscrit : $u_{\alpha,\beta}(\sigma_\infty) = \operatorname{int}(\beta)(\sigma_\infty)$, $u_{\alpha,\beta}(\rho) = \operatorname{int}(\alpha)(\rho)$.}

La formule (37) se traduit par la relation en $\xi, \eta$
$$ \xi = \eta h^\nu \leqno{(43)} $$

4) Détermination par $(\xi, \eta, f)$.

La donnée du couple $(\alpha, \beta) \in \hat{\Pi}_{0,3} \times \hat{\Pi}_{0,3}$
équivaut à : celle du couple $(\xi, \eta)$ qui s'en
déduit par (41), assujetti aux équations
$$ \begin{cases} \xi \rho(\xi) \rho^2(\xi) = 1 \\ \eta \sigma_\infty(\eta) = 1 \end{cases} \leqno{(44)} $$
Ainsi, la donnée de $u \in \tilde{\mathcal{M}}_{0,3}$ équivaut à celle d'un triplet
$$ (\xi, \eta, f) \qquad \begin{cases} \xi, \eta \in \hat{\Pi}_{0,3} \text{ modulo} \\ \xi \mapsto l^{-1} \xi \rho(l), \eta \mapsto l^{-1} \eta \sigma_\infty(l), f \mapsto l^{-1} f \end{cases} \leqno{(45)} $$
satisfaisant les relations (44), et la relation
$$ \boxed{\xi = \eta h^\nu} \leqno{(46) = (43)} $$
qui joue le rôle de "relation des lacets en $\xi, \eta$". L'avantage de cette présentation est

\newpage

%% ===== Page 191 =====

la simplicité maximale des équations (44), (46) -
et le fait que $f$ n'y intervienne pas ; en outre
les deux premières équations sont des équations
en $\xi$ seul et en $\eta$ seul, et la troisième relie
$\xi$ et $\eta$ de façon particulièrement simple - en
fait si $\mu=1$, i.e.\ $\nu=0$, c'est la relation $\xi=\eta$ !

L'inconvénient est qu'il y a trois équations
au lieu d'une, et qu'on n'a pas d'expression\footnote{mais si, on a [unclear: juste vu] l'expression explicite (42 ter) : $u(\sigma_\infty) = \operatorname{int}(f^{-1})(\xi \sigma_\infty)$, $u(\rho) = \operatorname{int}(f^{-1})(\eta \rho)$ !!}
explicite de $u(\sigma_\infty)$ et de $u(\rho)$ en termes
de $\xi, \eta, f$ : il faut d'abord, pour cela,
résoudre les équations :
$$
\begin{cases}
\alpha \rho(\alpha)^{-1} = \xi & (\alpha \in \widetilde{\Pi}_{0,3}^\wedge) \\
\beta \sigma_\infty(\beta)^{-1} = \eta & (\beta \in \widehat{\Pi}_{0,3}^\wedge)
\end{cases}
\leqno{(47)}
$$
en $\alpha, \beta$, pour pouvoir ensuite appliquer (42).

Un autre\marginpar{(!)} inconvénient, d'importance, c'est
que lorsqu'on travaille avec les images des
éléments envisagés dans $\operatorname{Gl}(2, \hat{\mathbf{Z}})^{\pm}$ la connaissance
(même dans $\operatorname{Sl}(2, \hat{\mathbf{Z}})$ lui-même)
de ces images de $\xi, \eta$ est
assez faible, qui ne permet nullement de
retrouver $\alpha, \beta$ comme éléments de $\operatorname{Gl}(2, \hat{\mathbf{Z}})^{\pm}$,
cf plus bas.

\newpage

%% ===== Page 192 =====

5) Composition avec un autom.\ intérieur. Relations des cocycles pour un composé $u'' = u'u$.

Soit $\gamma \in \widetilde{\Pi}_{0,3}$, et considérons le composé
$$
u' = \operatorname{int}(\gamma) \circ u
\leqno{(48)}
$$
il correspond à des invariants
$$
f', g', h', \alpha_0', \beta_0', \alpha', \beta', \zeta', \eta'
$$
donnés par
$$
\begin{cases}
f' = \gamma f \\
g' = \gamma g u(\gamma)^{-1}, \quad h' = \gamma h u(\gamma)^{-1} \\
\alpha_0' = \gamma \alpha_0, \quad \beta_0' = \gamma \beta_0 \\
\alpha' = \alpha, \quad \beta' = \beta \\
\zeta' = \zeta, \quad \eta' = \eta
\end{cases}
\leqno{(49)}
$$
Ainsi $\alpha$ et $\beta$ ne changent pas quand on remplace $u$ par $\operatorname{int}(\gamma) u$, ils ne dépendent que de l'autom.\ ext.\ à lacets défini par $u$, i.e.\ de l'image de $u$ dans $M_{0,3}^\sim$. C'est là un avantage substantiel de l'utilisation des $(\alpha, \beta, f)$ ou des $(\zeta, \eta, f)$, de préférence à $(\alpha_0, \beta_0, f)$ ou des $(g, h, f)$.

Soit maintenant
$$
u, u' \in M_{0,3}^\sim \quad , \quad \text{soit } u'' = u' u
\leqno{(50)}
$$
$$
\text{soient} \quad
\begin{cases}
f, g, h, \alpha_0, \beta_0, \alpha, \beta, \zeta, \eta \\
f', g', h', \alpha_0', \beta_0', \alpha', \beta', \zeta', \eta' \\
f'', g'', h'', \alpha_0'', \beta_0'', \alpha'', \beta'', \zeta'', \eta''
\end{cases}
\quad \text{et}
\leqno{(51)}
$$

\newpage

%% ===== Page 193 =====

les invariants relatifs à $u, u'; u''$ respectivement, on les exprime (dans la mesure du possible) en fonction les uns des autres,
$$
\begin{matrix}
g'' = u'(g)g' & \alpha_0'' = u'(\alpha_0) \alpha_0'^{\ast} & \alpha'' = (f'u')(\alpha) . \alpha'^{\ast} \\
h'' = u'(h)h' & \beta_0'' = u'(\beta_0) \beta_0' & \beta'' = (f'u')(\beta) . \beta' \\
\text{mais } f'' = f' u'(f) & &
\end{matrix} \leqno{(52)}
$$
\footnote{$\ast$ $u'(\alpha_0)$ et $(f'u')(\alpha)$ n'ont de sens à priori que si $\alpha_0, \alpha \in \Pi_{0,3}^\wedge$ (c.a.d. $M = \{1,2,3\}$) bien qu'on puisse convenir de $u'(\alpha_0) = u'\alpha_0 u'^{-1}$ d'où il résulte un sens pour $u''(\alpha_0)$ à une $\operatorname{int} \gamma$ près ds $\Pi_{0,3}^\wedge$, mais ce $\alpha''$ ne dépendra pas fonctoriellement de $u$.}

On voit donc que si on considère (pour $u$ variable dans $\mathcal{M}_{0,3}^{1\sim}$) $f, g, h, \alpha_0, \beta_0, \alpha, \beta$ comme des fonctions
$$
\mathcal{M}_{0,3}^{1\sim} \to \Pi_{0,3}^\wedge \text{ ou } \mathfrak{T}_{0,3}^{\sim\wedge}
$$
alors $f, g^{-1}, h^{-1}, \alpha_0^{-1}, \beta_0^{-1}$ satisfont les relations des cocycles, et $\alpha^{-1}, \beta^{-1}$ aussi quand on se limite au sous-groupe pour lequel $f=1$, i.e. les automorphismes $u \in \mathcal{M}_{0,3}^{1\sim}$ qui immobilisent $L_0 = \{ l \delta^n \mid n \in \hat{\mathbf{Z}} \}$. A vrai dire, cet énoncé pour $f, \alpha, \beta$ est un peu boiteux, à cause des indéterminations pour $f$, impliquant une indétermination correspondante pour $\alpha, \beta$, --- [unclear: un $f \mapsto b f$ amène $\alpha \mapsto b \alpha, \beta \mapsto b \beta$]. Pour obtenir des cocycles canoniques correspondants à $f, \alpha, \beta$, il faudrait introduire une vraie extension de $\mathcal{M}_{0,3}^{1\sim}$ par $\hat{\mathbf{Z}}$, sur le modèle des groupes de Teichmüller spéciaux, sur lequel $f, \alpha, \beta$ soient définis sans ambiguïté. [unclear: On n'entrera pas dans ces détails...]

\newpage

%% ===== Page 194 =====

\section*{III \quad Calculs dans $\operatorname{Gl}(2, \hat{\mathbf{Z}})'$ et dans $\operatorname{Gl}(2, \hat{\mathbf{Z}})$.}

Pour $x \in \mathcal{M}_{0,3}^{\sim \wedge}$ donné, d'où les quantités\footnote{NB : faut déterminer sous forme explicite le système des indéterminations de la p. précé. d'où la levée d'indétermination de $\alpha_0, \beta_0, \alpha, \beta$ ... elles logent alors dans $\operatorname{Gl}(2, \hat{\mathbf{Z}})/...$ au lieu de loger dans $\operatorname{Gl}(2, \hat{\mathbf{Z}})'$ comme $f, ...$ cf p. 197 en relisant formules plus bas.}
$$
f, g, h, \alpha_0, \beta_0, \alpha, \beta, \xi, \eta  \leqno{(53)}
$$
éléments de $\Pi_{0,\hat{3}}^\wedge$, sauf $\alpha_0, \alpha$ qui n'en sont pas (dans $\Pi_{0,\tilde{3}}^\wedge \simeq \hat{\mathfrak{T}}_{\hat{1},\hat{3}} \cdot \{1, \tau_0\}$) on désigne par les mêmes lettres les images dans $\operatorname{Gl}(2, \hat{\mathbf{Z}})' \defeq \operatorname{Gl}(2, \hat{\mathbf{Z}})/\pm 1$, en fait dans
$$
\operatorname{Sl}(2, \hat{\mathbf{Z}})' = \operatorname{Sl}(2, \hat{\mathbf{Z}})/\pm 1 ,
$$
sauf $\alpha_0, \alpha$ qui sont seulement dans $\operatorname{Gl}(2, \hat{\mathbf{Z}})'^{\pm 1}$ :
$$
\det \alpha_0 = \det \alpha = \varepsilon_3 \in \{\pm 1\} \leqno{(54)}
$$
où $\varepsilon_3$ est défini par
$$
\varepsilon_3 \equiv \mu \pmod 3 \leqno{(55)}
$$
à part cela :
$$
f, g, h, \beta_0, \beta, \xi, \eta \in \operatorname{Sl}(2, \hat{\mathbf{Z}})' \quad \text{i.e.} \leqno{(56)}
$$
$$
\det f = \det g = \det h = \det \beta_0 = \det \beta = \det \xi = \det \eta = 1
$$

Car on a $\det(-1)=1$, donc $\det$ se factorise
$$
\operatorname{Gl}(2, \hat{\mathbf{Z}})' \xrightarrow{\det} \hat{\mathbf{Z}}^*
$$
pour donner une suite exacte
$$
1 \to \operatorname{Sl}(2, \hat{\mathbf{Z}})' \to \operatorname{Gl}(2, \hat{\mathbf{Z}})' \to \hat{\mathbf{Z}}^* \to 1 . \leqno{(57)}
$$

On choisira des représentants
$$
f, \alpha, \beta \in \operatorname{Gl}(2, \hat{\mathbf{Z}}) \leqno{(58)}
$$
ce qui équivaut au choix de trois signes,

\newpage

%% ===== Page 195 =====

et en fonction de ces choix on définit les autres quantités par

$$
\begin{cases}
\xi = \alpha p(\alpha)^{-1} \quad , \quad \eta = \beta p(\beta)^{-1} \\
\alpha_0 = f^{-1}\alpha \quad , \quad \beta_0 = f^{-1}\beta \\
g = \alpha_0 p(\alpha_0)^{-1} = f^{-1}\xi p(f) \quad , \quad h = \beta_0 p(\beta_0)^{-1} = f^{-1}\eta p(f)
\end{cases} \leqno{(59)}
$$

On notera que $\xi$ et $\eta$, $g, h$ sont indépendants des choix des signes qui ont été faits. La relation essentielle qui relie $(f, \alpha_0, \beta_0)$, ou plutôt $\alpha, \beta$ (en plus des relations (59), qui jouent le rôle de relations de définition) est la "relation des lacets", sous forme matricielle, qui prend la forme correspondante : (46) (ou : (43))

$$
\xi = \varepsilon \eta [unclear: h^\nu] \qquad \varepsilon \in \{\pm 1\} \leqno{(60)}
$$

où $\varepsilon$ est un signe qui ne dépend plus d'aucun choix — on a (cf.\ aussi formules plus bas)

$$
\varepsilon = \varepsilon_2 \quad \text{(défini univoquement par } \mu \equiv \varepsilon_2 \pmod{4}\text{)} \leqno{(61)}
$$

\newpage

%% ===== Page 196 =====

et où $l_1$ est défini matriciellement par
$$
l_1 = \begin{pmatrix} 1 & 0 \\ 2 & 1 \end{pmatrix}
$$
d'où
$$
l_1^\nu = \begin{pmatrix} 1 & 0 \\ 2\nu & 1 \end{pmatrix} = \begin{pmatrix} 1 & 0 \\ \mu-1 & 1 \end{pmatrix}. \leqno{(62)}
$$

\underline{Remarques.}
(Relation avec les notations du § 41 (p.\ 356 ff, notamment formulaire récap.\ p.\ 364). Dans les groupes $\hat{\mathfrak{T}}_{0,4}$ et $\hat{\mathfrak{T}}_{0,3}$, les lettres $f, g, h$ de loc.\ cit.\ ont le même sens qu'ici, mais non les lettres $\alpha, \beta$ etc. Je vais distinguer par des soulignés, dans le cas d'ambiguïté avec la signification attachée dans loc.\ cit., on trouve :
$$
\underline{\alpha} = \alpha_0, \quad \underline{\beta} = \beta_0 \leqno{(i)}
$$
à partir de la p.\ 357 (i.e.\ $\underline{u}(\sigma_\infty) = \sigma_\infty$) on suppose pratiquement tjrs
$$
\underline{\beta} = f \beta_0 = f \leqno{(ii)}
$$
et loc.\ cit.\ (11) donne
$$
\underline{\varphi} = f \sigma_\infty(f)^{-1} = \beta \sigma_\infty(\beta)^{-1} = \eta \leqno{(iii)}
$$
tandis que (loc.\ cit.\ (44))
$$
\underline{\xi} = l_0^\nu \underline{\varphi} = l_0^\nu \eta
$$
qui comparé avec
$$
\xi = \eta l_1^\nu \leqno{(45)}
$$
donne
$$
\underline{\xi} = l_0^\nu \xi l_1^{-\nu} = l_0^\nu \xi \rho(l_0^\nu)^{-1} \leqno{(iv)}
$$
i.e.\ $\underline{\xi}$ est le $\rho$-transformé de $l_0^\nu$.
La relation (loc.\ cit.\ (49))
$$
\underline{\xi} = \alpha_0 \rho(\alpha_0)^{-1}
$$
comparée à $\xi = \alpha \rho(\alpha)^{-1}$ et à (iv), donne

\newpage

%% ===== Page 197 =====

(v) $\alpha_0 = l_0^\nu \alpha$.

Les formules (i) à (v) permettent de
traduire le formulaire de loc. cit. en termes
du formulaire que je vais expliciter.

On explicite $\alpha$ et $\beta$ par leurs matrices

(63) $\alpha = \begin{pmatrix} a & b \\ c & d \end{pmatrix}$, $\beta = \begin{pmatrix} r & s \\ t & u \end{pmatrix}$

où $a, b, c, d, r, s, t, u \in \hat{\mathbb{Z}}$, avec

(64) $ad-bc = \delta$, $ru-st = 1$.

On aura

(65) $\xi = \alpha \rho(\alpha)^{-1} = \begin{pmatrix} A & B \\ C & D \end{pmatrix}$

avec (cf. p. 342, 08)

(66) $\begin{cases} \delta A = (a^2+ab+b^2) - (ac+bd+bc) & \delta B = ac+bd+bc \\ \delta C = -(c^2+cd+d^2) + (ac+bd+ad) & \delta D = c^2+cd+d^2 \end{cases}$

(67) $\eta = \beta \sigma_0(\beta)^{-1} = \begin{pmatrix} R & S \\ T & U \end{pmatrix}$, avec

(68) $R = r^2+s^2$, $S = T = rt+su$, $U = t^2+u^2$.

La relation des lacets (60) prend la forme

(69) $\begin{pmatrix} A & B \\ C & D \end{pmatrix} = \varepsilon \begin{pmatrix} R & S \\ T & U \end{pmatrix} \begin{pmatrix} 1 & 0 \\ 2\nu & 1 \end{pmatrix}$

$= \varepsilon \begin{pmatrix} R+2\nu S & S \\ T+2\nu U & U \end{pmatrix}$,

\newpage

%% ===== Page 198 =====

\begin{flushright}
398
\end{flushright}

Notons les relations (conséquences de (66), (67)) :
$$
AD - BC = 1 \qquad RU - ST = 1 \leqno{(70)}
$$
$$
C + D - B = 1 \qquad T - S = 0 \text{ i.e. } T = S \leqno{(71)}
$$
Cette relation des lacets (69) est équivalente au système de relations
$$
\begin{cases}
D = \varepsilon U \\
\mu D = 1 \qquad (\text{où } \mu = 2\nu+1) \\
B = \varepsilon S
\end{cases} \leqno{(72)}
$$
i.e.\ au système d'équations
$$
\begin{cases}
c^2+cd+d^2 = \varepsilon(t^2+u^2) = \delta/\mu \\
ac+bd+bc = \varepsilon(rt+su)
\end{cases} \leqno{(73)}
$$

[unclear: C'est de la première de ces relations (73) que l'on] tire aussitôt
$$
\delta \equiv \mu \pmod{3}, \quad \varepsilon \equiv \mu \pmod{4} \leqno{(74)}
$$
qui caractérisent les signes $\delta, \varepsilon$ comme
$$
\delta = \varepsilon_3(\mu) = \varepsilon_3, \quad \varepsilon = \varepsilon_4(\mu) = \varepsilon_4 , \leqno{(75)}
$$
cf.\ th.\ p.\ 365.

La formule (83) de la page 370 prend ici la forme très simple (valable sans conditions autres que celles précisées jusqu'ici sur $\alpha, \beta, f$ --- $f$ n'intervient d'ailleurs pas dans le second membre de la formule) :

\newpage

%% ===== Page 199 =====

$$ \underset{(= \varepsilon_3)}{(\delta)} \alpha(\varphi + d) = \beta (t \tau_\infty + u)\tau_\infty^{\alpha_2} \leqno (76) \quad (\alpha = \alpha_{\sim}) $$

où $\alpha_2$ est défini par la condition 
$$ (-1)^{\alpha_2} = \varepsilon \quad (= \varepsilon_2) \quad . $$

Explicitant (76) par le calcul matriciel, on a
$$ \delta \alpha(\varphi + d) = \delta \begin{pmatrix} a & b \\ c & d \end{pmatrix} \begin{pmatrix} d & c \\ -c & c+d \end{pmatrix} = \delta \begin{pmatrix} 1 & B \\ 0 & D \end{pmatrix} $$
$$ \beta(t \tau_\infty + u)\tau_\infty^{\alpha_2} = \begin{pmatrix} r & s \\ t & u \end{pmatrix} \begin{pmatrix} u & \varepsilon t \\ -t & \varepsilon u \end{pmatrix} \begin{pmatrix} 1 & \alpha_2 \\ 0 & 1 \end{pmatrix} = \begin{pmatrix} 1 & \varepsilon S \\ 0 & \varepsilon U \end{pmatrix} $$
de sorte que (76) se réduit aux équations
$$ D = \varepsilon U, \quad B = \varepsilon S \quad , $$
donc la relation des lacets (68, 69) est équivalente à la conjonction de (76) et la relation
$$ \mu D = 1 \quad . \leqno (77) $$

Ceci posé, on posera (pour $(\alpha, \beta)$ fixés de façon\footnote{en fait, $\mu, \rho$ ne dépendent que de $\xi = \alpha(c), \eta = \beta(c)$, de même $\varepsilon, \delta$, donc la seule cause de dépendance de $u_{\alpha, \beta}$ tient à $B$, qui, lui, n'est pas déterminé par $\alpha \dots$ voir plus loin p. 410} à satisfaire la formule des lacets)
$$ u_{\alpha, \beta} \defeq (\mu \delta) \cdot \alpha(\varphi + d) = \mu \beta (t \tau_\infty + u) \tau_\infty^{\alpha_2} = H \begin{pmatrix} 1 & B \\ 0 & D \end{pmatrix} = \begin{pmatrix} H & HB \\ 0 & 1 \end{pmatrix} $$
de façon que $u_{\alpha, \beta}$ ne dépend en fait que de $\beta$ et des signes $\alpha_2$ qui, lui, n'est pas déterminé par $\beta$, mais par la valeur de $\delta(c^2+cd+d^2) \pmod 2 \dots$ ), et on posera, pour un triple $(\alpha, \beta, f)$ d'éléments de $\operatorname{Gl}(2, \hat{\mathbf{Z}})$, $\det f = 1$, $\alpha$ et $\beta$ reliés par la formule des lacets

\newpage

%% ===== Page 200 =====

$$
u_{\alpha,\beta,f} \defeq f^{-1} u_{\alpha,\beta} = (\mu\delta) \overbrace{f^{-1}\alpha}^{\alpha_0} (c\rho+d) = \mu \underbrace{f^{-1}\beta}_{\beta_0} (t\sigma_\infty+u)\tau_\infty^{\alpha_2} \leqno{(79)}
$$

et on trouve, comme il se doit\footnote{NB si $u,\sigma$ sont des éléments d'un groupe $G$, on pose bien sûr $u(\sigma) = u \sigma u^{-1}$.}\footnote{Ceci s'explique par le fait que $(u_{\alpha,\beta,f})$ a même effet sur $\hat{\pi}_1$ que l'homomorphisme de $\hat{\pi}_{0,3}$ déduit (par symétrie) des formules [unclear: pour] $V_1, V_2$ (générateurs de $\operatorname{Sl}(2,\mathbf{Z})$).}
$$
\begin{cases}
u_{\alpha,\beta,f}(\sigma_\infty) = \operatorname{int}(\beta_0)(\sigma_\infty) \\
u_{\alpha,\beta,f}(\rho) = \operatorname{int}(\alpha_0)(\rho)
\end{cases} \leqno{(80)}
$$
formules qui se réduisent à des formules où $f$ a disparu, [unclear: dont on dispose]
$$
\begin{cases}
u_{\alpha,\beta}(\sigma_\infty) = \operatorname{int}(\beta)(\sigma_\infty) \\
u_{\alpha,\beta}(\rho) = \operatorname{int}(\alpha)(\rho)
\end{cases} \leqno{(81)}
$$
qui résultent trivialement des deux expressions de (78) de $u_{\alpha,\beta}$ !

Ces formules (80) se complètent de façon naturelle en
$$
u_{\alpha,\beta,f}(\varepsilon_0) = \operatorname{int}(f^{-1})(\varepsilon^\mu) \quad (\text{où } \varepsilon = \varepsilon_0 = \pm 1) \leqno{(82)}
$$
qui à son tour se réduit à
$$
u_{\alpha,\beta}(\varepsilon_0) = \varepsilon^\mu \leqno{(83)}
$$
où $f$ ne figure plus. Comme $\varepsilon_0 = \sigma_\infty \rho$, cette formule se réduit à :
$$
u_{\alpha,\beta}(\sigma_\infty) u_{\alpha,\beta}(\rho) = \varepsilon^\mu
$$
qui compte tenu de (81) se réduit à :
$$
\operatorname{int}(\beta)(\sigma_\infty) \operatorname{int}(\alpha)(\rho) = \varepsilon^\mu \marginpar{juste}
$$

\newpage

%% ===== Page 201 =====

or le premier membre s'écrit
$$
\beta \sigma_\infty \rho^{-1} \alpha \rho \alpha^{-1} = (\beta \sigma_\infty \beta^{-1} \sigma_\infty) (\sigma_\infty \rho) (\rho^{-1} \alpha \rho \alpha^{-1})
$$
$$
= \eta \cdot \varepsilon_0 \cdot \rho^{-1}(\xi) = \eta \cdot ( \underbrace{\varepsilon_0 \rho^{-1}(\xi) \varepsilon_0}_{\sigma_\infty(\xi)} ) \cdot \varepsilon_0 = \eta \sigma_\infty(\xi) \varepsilon_0
$$
donc la relation (83) s'écrit aussi
$$
\eta \sigma_\infty(\xi) = \varepsilon \varepsilon_0^{\mu-1} \quad (= \varepsilon l_0^\nu)
$$
or on a $\xi = \varepsilon \eta l_0^\nu$ (relation des lacets (60))
donc
$$
\eta \sigma_\infty(\xi) = \varepsilon \eta \sigma_\infty(\eta) l_0^\nu = \varepsilon l_0^\nu \quad , \quad \text{OK} \, .
$$

\underline{Remarques}. Toutes ces formules sont valables dans $\operatorname{Sl}(2, \hat{\mathbf{Z}})$, sans ambiguité de signes, dès que $\alpha, \beta, f$ sont choisis. Si on change les signes de $\alpha$ et de $\beta$ --- les deux, la relation des lacets n'est pas modifiée - et $u_{\alpha,\beta}$ n'est pas changé non plus. Par contre ni $u_{\alpha,\beta,f}$ l'est [unclear: au signe près : $f$]. Par contre, un changt. de signe sur $f$ implique un changt de signes sur $u_{\alpha,\beta,f}$. Mais un tel changement de signes ne change bien sûr rien à la validité des formules (80), (82). Quand on fait

\newpage

%% ===== Page 202 =====

sur $(\alpha, \beta, f)$ le changement banal
$$
\alpha' = \ell_0^n \alpha, \quad \beta' = \ell_0^{-n} \beta, \quad f' = \ell_0^{-n} f \quad (n \in \mathbf{Z}) \leqno{(84)}
$$
alors on a
$$
c' = c, d' = d \qquad t' = t, u' = u \qquad [unclear: p'_0 = p_0]
$$
$$
\begin{cases}
f'^{-1} \alpha' = f^{-1} \alpha \\
\alpha'_0 = \alpha_0
\end{cases}
\quad
\begin{cases}
f'^{-1} \beta' = f^{-1} \beta \\
\beta'_0 = \beta_0
\end{cases}
\leqno{(85)}
$$
enfin
$\delta' (= \det \alpha') = \delta (= \det \alpha)$ (car $\det \ell_0^n = 1$)
$$
\text{d'où} \quad \mu' (= \delta' / (c'^2 + c'd' + d'^2)) = \mu (= \delta / (c^2 + cd + d^2)) \leqno{(86)}
$$
d'où $\varepsilon'_2 = \varepsilon_2, \alpha'_2 = \alpha_2 \dots$ (79)

c'est au-delà de ce qu'il faut pour vérifier que
$$
u_{\alpha', \beta', f'} = u_{\alpha, \beta, f}. \leqno{(87)}
$$

Rappelons aussi (cf. page pr. 384)

Proposition. Soient $\mu \in \hat{\mathbf{Z}}^\ast$, d'où les signes $\varepsilon_2 = \varepsilon_2(\mu)$ et $\varepsilon_3 = \varepsilon_3(\mu)$ (définis par $\mu \equiv \varepsilon_2 \, (4)$, $\mu \equiv \varepsilon_3 \, (3)$).\footnote{NB Pour le strict énoncé on se fixe $c, d, t, u \in \hat{\mathbf{Z}}$. Pour qu'il existe $a, b, r, s \in \hat{\mathbf{Z}}$ donnant $\alpha, \beta$ satisfaisant (88) !}
Pour qu'il existe $a, b, r, s \in \hat{\mathbf{Z}}$ des matrices
$$
\alpha = \begin{pmatrix} a & b \\ c & d \end{pmatrix}, \quad \beta = \begin{pmatrix} r & s \\ t & u \end{pmatrix}
$$
satisfaisant (64)
$$
ad - bc = \varepsilon_3 \qquad ru - st = 1
$$
et la relation des lacets
$$
\alpha \rho \ell_0 \varepsilon_1^2 (\beta \sigma_{\infty} \ell_0 \beta^{-1}) \ell_1^{\nu} \qquad \text{où } \nu = (\mu - 1)/2
$$
(i.e. les relations (73)) il faut et il suffit qu'on ait
$$
\varepsilon_3(c^2 + cd + d^2) = \varepsilon_2(t^2 + tu + u^2) = 1/\mu \leqno{(88)}
$$
et l'ensemble des couples $(\alpha, \beta)$ satisfaisant à ces conditions est un torseur sous $\hat{\mathbf{Z}}$, opérant...

\newpage

%% ===== Page 203 =====

par
$$
(\alpha,\beta) \longmapsto (l_0^n \alpha, l_0^{-n} \beta).
$$

Corollaire. Pour qu'on puisse trouver $\alpha, \beta$ tels qu'on ait de plus $\alpha \equiv \beta \equiv 1 \quad (2)$, il faut et il suffit qu'on ait
$$
c \equiv t \equiv 0 \ (2) \quad d \equiv u \equiv 1 \ (2) \leqno{(89)}
$$
et alors l'ensemble des solutions $(\alpha,\beta)$ satisfaisant à ces conditions supplémentaires est un torseur sous $\mathbf{Z}$, opérant par
$$
(\alpha,\beta) \longmapsto (l_0^n \alpha, l_0^{-n} \beta)
$$

Scholie. On peut donc dire que pour un élément $\underline{\mu}$ de $\mathcal{N}_{0,3}^{\wedge \sim}$ qui définit deux\footnote{modulo la relation d'équivalence $(\alpha, \beta) \mapsto (l_0^n \alpha, l_0^{-n} \beta)$} $\alpha, \beta$ de $\Pi_{0,3}^\wedge, \Pi_{0,3}^\wedge$ (satisfaisant la relation des lacets), il lui est attaché un "invariant adélique"
$$
\Theta = (\mu, c, d, t, u) \leqno{(90)}
$$
($(c, d, t, u)$ déterminés aux ambiguïtés de signes près, sans plus) satisfaisant (88) --- et c'est ce qu'on peut faire de mieux dans le contexte actuel : les $\alpha, \beta$ n'ont pas un caractère invariant. Quand on a un élément de $\mathcal{N}_{0,3}^{\wedge \sim}$ --- i.e. un [unclear: relèvement absolu] celà --- on a de plus un $f \in \Pi_{0,3}^\wedge$, donnant un $\bar{f} \in \operatorname{Sl}(2, \hat{\mathbf{Z}})'$, d'où
$$
f = \begin{pmatrix} j & k \\ l & m \end{pmatrix} \leqno{(91)}
$$
défini mod $\pm 1$ et mod $f \mapsto l_0^n f$.

\newpage

%% ===== Page 204 =====

Mais cette donnée supplémentaire est un peu trop faible pour $\underline{\underline{u}} = \underline{\underline{u}}_{\alpha,\beta,f}$ l'élément de $\operatorname{Sl}(2,\hat{\mathbf{Z}})'$ associé à $u$.

Il vaut mieux dire que pour $\mu, c, d, t, u$ fixés on a un $\hat{\mathbf{Z}}$-torseur des $(\alpha, \beta) \equiv (1,1) \ (2)$ satisfaisant aux conditions dites dans la proposition 8 et le corollaire, soit $T = T_{\Theta}$, on fait opérer alors $\hat{\mathbf{Z}}$ diagonalement sur $T \times \operatorname{Sl}(2,\hat{\mathbf{Z}})$ par
$$
(\alpha, \beta, f) \longmapsto (l_0^n \alpha, l_0^{-n} \beta, l_0^{-n} f)
$$
et on regarde le quotient - c'est dans le torseur par $\check{T}$ (torseur opposé : $\check{T}$) de $\operatorname{Sl}(2,\hat{\mathbf{Z}})$, sur lequel $\hat{\mathbf{Z}}$ opère par $f \mapsto l_0^{-n} f$. c'est la donnée d'un point de
$$
\check{T} \wedge^{\hat{\mathbf{Z}}} \operatorname{Sl}(2,\hat{\mathbf{Z}}) \quad \text{(c'est un } \operatorname{Sl}(2,\hat{\mathbf{Z}})\text{-torseur à droite)}
$$
qui permet de définir un élément
$$
\underline{\underline{u}} = \underline{\underline{u}}_{\alpha, \beta, f} ,
$$
qui changera de signe si $f$ est remplacé par $(-f)$.

Décidément, ces formulations [unclear: sont un] peu enthousiasmantes, et il devient clair qu'il faudra introduire explicite-

\newpage

%% ===== Page 205 =====

ment et étudier de façon un peu minutieuse [unclear: décrire] le groupe abélien formé par les systèmes $(\alpha, \beta, f)$ satisfaisant aux relations des lacets, et envoyer l'homomorphisme canonique de $\widetilde{\mathcal{M}}_{0,3}$ dans ledit groupe. Mais je voudrais avant cela terminer les gammes notationnelles.

\newpage

%% ===== Page 206 =====

\noindent \underline{IV. Les trois cas particuliers fondamentaux.}

On revient aux éléments $g, h, f$ etc dans $\hat{\Pi}_{0,3}$ ou $\hat{\tilde{\Pi}}_{0,3}$ qui correspondent à $u$, pouvant être décrits de quatre façons différentes par l'un ou l'autre des quatre systèmes 
$(g, h, f)$, $(\alpha_0, \beta_0, f)$, $(\alpha, \beta, f)$, $(\xi, \eta, f)$.

\underline{1\textsuperscript{er} Cas} :
$$
h = 1, \text{ i.e. } \beta_0 = 1, \text{ i.e. } \beta = f, \text{ i.e. } \eta = f y f y^{-1}, \leqno{(93)}
$$
i.e.
$$
[u, \tau_\infty] = 1 \leqno{(94)}
$$
Dans un système $(g, h, f)$ etc quelconque, il y a un autom.\ intérieur et un seul $\operatorname{int}(v)$ tel que $u' = \operatorname{int}(v)(u)$ satisfasse à la propriété précédente i.e.\ commute à $\tau_\infty$, ce sera pour $v = \beta_0^{-1}$ (cf formules (49)).

Les $u$ satisfaisant (93), s'exprimant en termes des triples $(\xi, \eta, f)$ satisfaisant les trois relations (44) (46) (en faisant pour $\eta$ la substitution [unclear: indiquée]), se correspondent ainsi. La relation $\eta = f y f^{-1}$ de (93) amène aux couples $(\xi, f)$ d'éléments de $\hat{\Pi}_{0,3}$ satisfaisant :
\marginpar{[unclear: NB on avait $\xi = f y f y^{-1}$ au lieu de $\dots$ il est plus judicieux d'écrire $\dots$ mais $\dots$]}
$$
\xi \rho(\xi) \rho^2(\xi) = 1 \leqno{(95)}
$$
$$
\xi = (f y f^{-1}) l_1^2 \quad \text{avec } \dots \leqno{(96)}
$$

\newpage

%% ===== Page 207 =====

\marginpar{\alpha, \beta \text{ astreints à la relation (cf. § 4, (37))}}
ou aux couples $(\alpha, \beta)$ d'éléments de $\widehat{\Pi}_{0,3}$ satisfaisant la relation des lacets (37)
$$ (\alpha \rho(\alpha)^{-1}) = (\beta \tau_\infty(\beta)^{-1}) b^\alpha \leqno{(97)=(37)} $$
(qu'on complétera dans un moment par $f = \beta$).
Au changement près de $\alpha$ en $b^\alpha \alpha = \alpha_0$, c'étaient les conventions du § 4, où la formule prend encore la forme
$$ \underbrace{\alpha_0 \rho(\alpha_0)^{-1}}_{\xi} = b^\alpha \underbrace{(f \tau_\infty(f)^{-1})}_{[unclear: f=\beta]} $$

Compte tenu de (79), les conditions (97) (94) prennent aussi la forme
$$ u_{\alpha,\beta,f} = \mu (b \tau_\infty + u) \tau_\infty^{\alpha_2} \leqno{(98)} $$
et impliquent que $u = u_{\alpha,\beta,f}$ normalise la sous-algèbre de l'algèbre des matrices $M_2(\hat{\mathbf{Z}})$ engendrée par $\sigma_\infty$ (compte tenu de la relation
$$ \tau_\infty (\sigma_\infty) = - \sigma_\infty \qquad \text{dans } \operatorname{Gl}(2, \hat{\mathbf{Z}})). \leqno{(99)} $$

Si on désigne le sous-groupe de $\widehat{\mathcal{M}}_{0,3}^\sim$ formé des $u$ satisfaisant aux conditions équivalentes (93) (94) (98)
$$ \widehat{\mathcal{M}}_{0,3}^\sim(1/2) = \left\{ u \in \widehat{\mathcal{M}}_{0,3}^\sim \mid [u, \tau_\infty] = 1 \right\} \isom \widehat{\mathcal{N}}_{0,3}^\sim \leqno{(100)} $$
on trouve que l'hom. composé

\newpage

%% ===== Page 208 =====

\footnote{NB $\mathcal{M}_{0,3}^{1 \sim}$ est stable par $x \mapsto \sigma_\infty x \sigma_\infty^{-1}$ et par $x \mapsto \tau_\infty x \tau_\infty^{-1}$ (puisque $\sigma_\infty$ et $\tau_\infty$ normalisent $\sigma_\infty$)}
$$
\begin{tikzcd}
\mathcal{N}_{0,3}^{1 \sim} \simeq \mathcal{M}_{0,3}(1/2)^{\sim} \arrow[r, hook] & \mathcal{M}_{0,3}^{1 \sim} \arrow[r] & \operatorname{Gl}(2, \hat{\mathbf{Z}})' \\
& u \arrow[r, maps to] & \underline{\underline{u}} = u_{\alpha, \beta, f}
\end{tikzcd} \leqno{(101)}
$$
qui (préféré le voir en termes du groupe $\operatorname{Gl}(2, \hat{\mathbf{Z}}) \simeq M_2(\hat{\mathbf{Z}})^*$, plutôt que de l'algèbre des matrices, elle se factorise par le normalisateur du tore maximal
$$ \underline{T}_{\hat{\mathbf{Z}}} = \left\{ t \sigma_\infty + u \mid (t, u) \in \hat{\mathbf{Z}}^2, \, \det(t \sigma_\infty + u) \in \hat{\mathbf{Z}}^* \right\} \leqno{(102)} $$

\marginpar{sûrement inutile! il y a ainsi l'isomorphisme spect. qu'il n'y a que 2 classes de courbes elliptiques d'invariant $j=\infty$ (concerne l'étude de l'orbite $1/2, 2, -1$ des $\mathcal{N}_{0,3}$ dans ...)}
l'homomorphisme $k_{1/2}$ de (101) de l'extension $\mathcal{M}_{0,3}^{1 \sim}$ (de $\mathcal{N}_{0,3}^{1 \sim}$ par $\hat{\pi}_{0,3}$, duquel on déduit que $\mathcal{N}_{0,3}^{1 \sim}$ pour $\widehat{\mathfrak{T}}_{0,3}^+$) que se factorise par un abélianisateur.

2\textsuperscript{o} \underline{Cas} \quad C'est le cas
$$ g = 1 \text{ i.e. } \beta_0 = 1 \text{ i.e. } \alpha = f \text{ i.e. } \xi = f \rho f \rho^{-1} \text{, i.e.} \leqno{(103)} $$
$$ [\alpha, \rho] = 1 . \leqno{(104)} $$

Pour $u \in \widehat{\Pi}_{0,3}^\sim = \widehat{\Pi}_{0,3} \cdot (1, \tau_\infty)$ quelconques, il y a un $\gamma \in \operatorname{Centr}_{\widehat{\Pi}_{0,3}}(\tau_\infty)$ tel que $u! = \operatorname{int}(\alpha)(\gamma) \cdot u = \gamma u$
commute ainsi à $\rho$, savoir $\gamma = \beta_0^{-1}$.

Mais de ce fait attention que $\gamma$ n'est pas toujours dans $\widehat{\Pi}_{0,3}$, en fait
$$ \gamma \in \alpha \beta_0 \widehat{\Pi}_{0,3} \iff \varepsilon_3(u) = 1 \text{ i.e. } f(u) = 1 \leqno{(105)} $$

\newpage

%% ===== Page 209 =====

Si on désigne par $\mathcal{N}_{0,3}^{!\sim}(-J)$ le sous-groupe\footnote{NB $\mathcal{M}_{0,3}^{!\sim}(-J)$ est stable par $u \mapsto \rho u \rho^{-1}$.}
de $\mathcal{M}_{0,3}^{!\sim}$ formé des automorphismes
dont l'image dans $\mathcal{N}_{0,3}^{!\sim}$ est dans
le sous-groupe d'indice $2$ noyau de
$$
u \mapsto \varepsilon_3(\mu(u)) = \varepsilon_3(u)
$$
i.e.
$$
\mathcal{N}_{0,3}^{!\sim}(-J) = \left\{ u \in \mathcal{N}_{0,3}^{!\sim} \mid \varepsilon_3(u) = 1 \text{ i.e. } \bar{j}(u) \equiv 1 \ (3) \right\}
\leqno{(105)}
$$
et l'on a
$$
\mathcal{M}_{0,3}^{!\sim}[-J] \isommap \mathcal{N}_{0,3}^{!\sim}[-J] \ .
\leqno{(107)}
$$
En l'incorporant on trouve donc un
homomorphisme composé
$$
\mathcal{N}_{0,3}^{!\sim}(-J) \isommap \mathcal{M}_{0,3}^{!\sim}(-J) \hookrightarrow \mathcal{M}_{0,3}^{!\sim}
\leqno{(108)}
$$
dont le composé avec $k(-J) : \mathcal{M}_{0,3}^{!\sim} \to \dots$ est un
homomorphisme qui se factorise par le tore
maximal $T_{\rho}$ de $\operatorname{Gl}(2, \hat{\mathbf{Z}})'$, formé des
$$
\left\{ c \rho + d \mid c, d \in \hat{\mathbf{Z}}, c^2 - c d + d^2 \in \hat{\mathbf{Z}}^{\times} \right\} \ .
\leqno{(109)}
$$

Je suppute\footnote{canulé, pas [unclear: plus de] raison que tout à l'heure...} (c'est le seul sous-groupe de $\mathcal{M}_{0,3}^{!\sim}$) qui ait cette propriété. Celle-ci résulte de la forme équivalente des conditions (103) (101), formulée en (79) :

\newpage

%% ===== Page 210 =====

$$
u = u_{\alpha, \beta, f} = \mu(c\rho+d) \leqno (110)
$$
(NB $\delta = \varepsilon_3(\mu)$ est égal à $1$ dans ce cas).

Les $u$ en question correspondent aux couples $(\eta, f)$ satisfaisant\marginpar{On pourra - si on y tient - écrire (111) sous la forme d'un couple $(\alpha, \beta)$ it. pour (112)}
$$
\mathfrak{d}_{03}(\eta) = 1 \leqno (111)
$$
$$
(f \rho f j)^{-1} = \eta \rho^\nu \leqno (112)
$$
ou encore aux couples $(\alpha, \beta)$ satisfaisant la relation en lacets [unclear: $(g f) = (\beta f)$], qui se complète alors par $f = \alpha$.

3° Cas C'est celui où
$$
f \equiv 1 \ (L_0) \text{ i.e. } \alpha \equiv \alpha_0 \ (L_0 \text{ à g.}) \text{ i.e. } \beta \equiv \beta_0 \ (L_0 \text{ à g.}) \leqno (113)
$$
$$
u(\varepsilon_0) = \varepsilon_0^\nu \ (\nu \in \hat{\mathbf{Z}}^\times) \text{ i.e. } \leqno (114)
$$
$$
[u, \varepsilon_0] = l_0^\nu \quad (\nu \in \hat{\mathbf{Z}} \text{ - ou } \nu = \pm 1)\footnote{NB $L_0 = l_0^{\hat{\mathbf{Z}}}$ agissant sur $\hat{\pi}_{0,3}$ i.e. $u$ normalise $L_0 = \{ l_0^n \mid n \in \hat{\mathbf{Z}} \}$ - ou encore dans $\hat{\mathbf{Z}}$.}
$$
Dans ce cas, la substitution $u \mapsto (l_0^\nu \alpha, l_0^\nu \beta, l_0^\nu f)$ va exiger
$$
f = 1 \leqno (115)
$$
ce qui détermine $\alpha, \beta$ de façon unique.

Le sous-groupe $\mathcal{M}_{0,3}^{\wedge \sim}[0]$\footnote{On ne remonte pas dans $\mathcal{M}_{0,3}^{\wedge \sim}(0)$ par $\mathcal{M}_{0,3}^{\wedge \sim}$ car il n'y a pas de section de l'ext. $\hat{\pi}_{0,3} \to \mathcal{M}_{0,3}^{\wedge \sim}$ qui normalise $L_0$.} des $u \in \mathcal{M}_{0,3}^{\wedge \sim}$ qui normalisent $L_0$ est en bijection avec l'ensemble des couples $(\alpha, \beta) \ (\alpha \in \hat{\pi}_{0,3} = \hat{\pi}_{0,3} \cdot \{1, \tau_0\}, \beta \in \hat{\pi}_{0,3})$ satisfaisant

\newpage

%% ===== Page 211 =====

les sempiternelles relations des lacets\footnote{Dans tout ceci, je laisse de côté la "condition lointaine" sur les relations dans $\hat{\Pi}_{0,3}$.}
$$
(\alpha p l_0 \alpha^{-1}) = (l_0 \beta l_0^{-1}) l_1^\nu \quad (\nu \in \hat{\mathbf{Z}} \text{ conv.})
$$
ou encore (aux couples $(f, g)$ d'éléments de $\hat{\Pi}_{0,3}$ satisfaisant
$$
f \rho l_0 \rho f^{-1} = 1, \quad g \tau l_0 \tau g^{-1} = 1, \quad f = g l_1^\nu \quad (\nu \in \hat{\mathbf{Z}} \text{ conv.}) \leqno{(116)}
$$
(=(44)+(45))
au couple $(\alpha, \beta)$ correspond... l'automorphisme $u$ défini par
$$
u(v_0) = \alpha v_0, \quad u(p) = g p \leqno{(117)}
$$
= (23)
ou encore par
$$
u(v_0) = \operatorname{int}(\alpha) v_0, \quad u(p) = \operatorname{int}(\beta) p \leqno{(118)}
$$

Dans le groupe $\mathcal{M}_{0,3}^{\sim}[0]$, il y a le sous-groupe
$$
\mathcal{M}_{0,3}^{\sim}[0] \cap \hat{\Pi}_{0,3} \simeq L_0 \leqno{(119)}
$$
L'inclusion
$$
\hat{\mathbf{Z}} \xrightarrow{\sim} L_0 \hookrightarrow \mathcal{M}_{0,3}^{\sim} \leqno{(120)}
$$
étant décrite par
$$
\lambda = l_0^n \longmapsto (\alpha = l_0^n, \beta = l_0^n) \quad \text{i.e. } \alpha=\lambda, \beta=\lambda \leqno{(121)}
$$
ou encore
$$
\lambda = l_0^n \longmapsto (f = l_0^n l_1^{-n}, g = l_0^{-n} l_1^n) \leqno{(122)}
$$
la multiplication par $\lambda \in L_0$ ($\lambda = l_0^n$) dans $\mathcal{M}_{0,3}^{\sim}[0]$ étant explicitée par
$$
(\alpha, \beta) \longmapsto (l_0^n \alpha, l_0^{-n} \beta) \leqno{(123)}
$$
$$
(f, g) \longmapsto (l_0^n f l_1^{-n}, l_0^{-n} g l_1^n) \leqno{(124)}
$$

\newpage

%% ===== Page 212 =====

Comme on a un iso
$$
\mathcal{M}_{0,3}^{\sim}[0] / L_0 \isommap \mathcal{N}_{0,3}^{\sim}
$$
on trouve le

Corollaire. Le groupe $\mathcal{N}_{0,3}^{\sim}$ est en bijection
canonique avec l'ens. des couples
$$
(\alpha, \beta) \in \hat{\Pi}_{0,3} \times \hat{\Pi}_{0,3}
$$
satisfaisant la relation
des cartes (97)
$$
\alpha \rho(\alpha)^{-1} = (\beta \tau_0(\beta)^{-1}) l_1^v \quad (\text{pour } v \in \hat{\mathbf{Z}} \text{ conv.})
$$
modulo les transformations $(\alpha, \beta) \mapsto (l_0 \alpha, l_0^{-1} \beta)$
de $\mathcal{M}_{0,3}^{\sim}[0]$ ou de son quotient $\mathcal{N}_{0,3}^{\sim}$.

La loi de groupe s'explicite en termes
de $(g, h)$, ou de $(\alpha, \beta)$, par les "formules
de cocycles" (52) : où $u'' = u'u$, si $u', u$ sont donnés
par $(g', h')$, $(g, h)$, $(g'', h'')$, resp. par
$(\alpha', \beta')$, $(\alpha, \beta)$ et $(\alpha'', \beta'')$, on a :
$$
\begin{cases}
g'' = u'(g) g' & h'' = u'(h) h' \\
\alpha'' = u'(\alpha) \alpha' & \beta'' = u'(\beta) \beta'
\end{cases} \leqno{(125)}
$$

Remarques sur les notations $\mathcal{M}_{0,3}^{\sim}(q)$, pour $q = \frac{1}{2}, -1, 0$.
En fait pour tout point algébrique [unclear: considérant un $K_q \subset \mathbf{C}$]
$$
q \in \mathcal{M}_{0,3}(\mathbf{C}) = \mathbf{C} \setminus \{0, 1\} = \mathbb{P}^1(\mathbf{C}) \setminus \{0, 1, \infty\} \leqno{(126)}
$$
\footnote{En fait, il faut plus précisément [unclear: lire (127) qu'il s'agit d'un scindage de l'extension]} on obtient un [unclear: homomorphisme injectif partiel]
$$
K_q : \mathcal{N}_{0,3}^{\sim} \hookrightarrow \mathcal{M}_{0,3}^{\sim} \quad , \quad \operatorname{Im} K_q = \mathcal{M}_{0,3}^{\sim}(q)
$$
d'où [unclear: le scindage] de l'extension $\mathcal{M}_{0,3}^{\sim}$ de $\mathcal{N}_{0,3}^{\sim}$ par $\hat{\Pi}_{0,3}$

\newpage

%% ===== Page 213 =====

donne au moins un scindage partiel au dessus
de $\boldsymbol{\Gamma}_{\mathbf{Q}}(q)$, qu'on s'efforcera de déterminer ;
$\mathcal{N}_{0,3}^{\prime \sim}(q)$ désignera un sous-groupe
topologique de $\mathcal{N}_{0,3}^{\prime \sim}$ qui l'\footnote{La définition importe peu lorsque $q$ est élément d'une extension cyclotomique de $\mathbf{Q}$, de telle façon que $\boldsymbol{\Gamma}_{\mathbf{Q}}(q)$ contienne un [unclear: sous-groupe distingué dont l'image dans \dots]} [unclear: envoie]
: déterminer, s'il y a lieu
pour tous les $q$ (enfin $\dots$) , l'indice dans
$\boldsymbol{\Gamma}_{\mathbf{Q}}$ du sous-groupe $\boldsymbol{\Gamma}_{\mathbf{Q}}(q)$ de $\boldsymbol{\Gamma}_{\mathbf{Q}}$
stabilisateur de $q$.

L'analyse des cas
$1^o$ et $2^o$ ci-dessus reviendrait en fait à
la détermination de tels scindages pour
$q = \frac{1}{2}$ et pour $q = \overline{-1}$ - dans le premier
cas $\boldsymbol{\Gamma}_{\mathbf{Q}}(q)$ est $\boldsymbol{\Gamma}_{\mathbf{Q}}$ tout entier, dans le
deuxième c'est le sous-groupe d'indice 2
noyau de $\varepsilon_2 : \boldsymbol{\Gamma}_{\mathbf{Q}} \to \{ \pm 1 \}$, et idem pour
$\mathcal{M}_{0,3}^{\prime \sim}(q)$. En transformant par des $\rho$
on déduirait des scindages $k_{1/2}$ les scindages
$k_2$ et $k_{-1}$ - en transformant par des élé-
ments quelconques des $\sigma_i$ [unclear: engendrant] $\dots$ on
déduirait du scindage $k_{\overline{-1}}$ un scindage de
type $k_{\overline{-1}}$.

Quant à $\mathcal{M}_{0,3}^{\prime \sim}(\underline{o})$, il correspond [unclear: un point]

\newpage

%% ===== Page 214 =====

$0$ de $\hat{U}_{0,3} = \mathbb{P}^1_{\mathbf{Q}}$, il y correspond deux autres sous-groupes analogues $\mathcal{M}_{0,3}^{\sim \wedge}[q]$, où $q=1, q=\infty$, correspondant aux deux autres points "à l'infini" de $U_{0,3}$ ; et qui se déduisent du précédent en appliquant $\rho, \rho^2$. J'ai dit des bêtises quand au commencement j'ai cherché des sections partielles de $\mathcal{M}_{0,3}^\sim$ sur $\mathcal{N}_{0,3}^{\sim \wedge}$, puisque l'intersection $\hat{\Pi}_{0,3}$ n'est pas réduite à $1$, mais est l'un des trois sous-groupes-lacets type $L_0, L_1, L_\infty$. Mais on sait que le choix d'une uniformisante de l'anneau local de $\hat{U}_{0,3}$ en $0$ (disons) détermine un scindage de $\mathcal{M}_{0,3}^\sim(0)$ au dessus de son quotient extérieur $\mathcal{N}_{0,3}^{\sim \wedge}$ par $L_0$ — et il faudra dans la suite essayer de saisir un tel scindage (associé à une uniformisante) dans l'optique de notre calcul, comme une façon privilégiée de fixer les $(\alpha, \beta)$ satisfaisant aux relations des lacets, de "normaliser" le choix de $(\alpha, \beta)$ mod $L_0$. J'y reviendrai par la suite, sûrement.

\newpage

%% ===== Page 215 =====

\section*{I Récapitulation en termes de $\mathcal{M}_{0,3}^{1 \sim}$}

Il apparait finalement que la façon la plus simple d'étudier les autom.\ [unclear: de lacet] de $\hat{\pi}_{0,3}$ (correspondant exact.\ à $\mathfrak{S}_3$, i.e.\ les éléments de $\mathcal{M}_{0,3}^{1 \sim}$) est de commencer par l'étude dans $\mathcal{M}_{0,3}^{1 \sim}$ du sous-groupe $\mathcal{M}_{0,3}^{1 \sim} [\sigma]$, normalisateur de $L_0$.

Posant pour $u \in \mathcal{M}_{0,3}^{1 \sim} [\sigma]$
$$
u(\sigma_0) = \eta \sigma_0 \quad , \quad u(p) = \xi p \quad (\eta, \xi \in \hat{\pi}_{0,3})
\leqno{(128)}
$$
i.e.
$$
[u, \sigma_0] = \eta \quad , \quad [u, p] = \xi
\leqno{(129)}
$$
qui expriment la commutation ext.\ de $u$ à $\mathfrak{S}_3$ (engendré en effet par $\sigma_0$ et $p$), la condition de normalisation de $L_0$ s'exprime par une relation de la forme\marginpar{(130 bis) équivaut à dire: \\ $u(e_1) = \operatorname{int} p (e_1^\mu)$ \\ $u(e_0) = e_0^\mu$ \\ $u(e_\infty) = \operatorname{int} p^{-1} (e_\infty^\mu)$}
$$
u(l_0) = l_0^\mu \quad \text{ou} \quad \varepsilon(u) = \mu \quad (\mu \in \hat{\mathbf{Z}}^*).
\leqno{(130)}
$$

Or on [unclear: calcule] (compte tenu de $l_0 = \sigma_0 p \sigma_0 p$, par un calcul immédiat) la formule
$$
u(l_0) = u(\sigma_0 p \sigma_0 p) = \eta \sigma_0 \xi p \eta \sigma_0 \xi p = \eta \sigma_0(\xi) l_0
$$
donc (130) s'écrit
$$
\eta \sigma_0(\xi) = l_0^{\mu - 1} \quad \text{soit, en posant } \mu - 1 = 2 \nu \ (\nu \in \hat{\mathbf{Z}})
$$
$$
\eta \sigma_0(\xi) = l_0^{2\nu} \quad \text{soit, appliquant } \sigma_0 \text{ : }
$$
$$
\sigma_0(\eta) \xi = l_0^{2\nu} .
\leqno{(131)}
$$

Or les relations
$$
(\eta \sigma_0)^2 = 1 \quad , \quad (\xi p)^3 = 1
\leqno{(132)}
$$
s'écrivent aussi
$$
\eta \sigma_0(\eta) = 1 \quad , \quad \xi p(\xi) p^2(\xi) = 1 \ ,
\leqno{(133)}
$$
d'où $\sigma_0(\eta) = \eta^{-1}$, de sorte que (131) s'écrit aussi de façon plus symétrique
$$
\xi = \eta l_0^{2\nu} \qquad (\text{relation des lacets})
\leqno{(134)}
$$

\newpage

%% ===== Page 216 =====

On trouve ainsi le

\underline{Théorème}. Les éléments $u$ de $M_{0,3}^{!\sim}[\vec{0}]$ --- i.e.\ les
autom. ext. de $\widehat{\Pi}_{0,3}$ qui commutent ext. à
$\mathfrak{S}_3$ et qui normalisent $L_0$ --- sont en corr.
biunivoque\footnote{via les formules (128) ou (129)} avec les couples $(\xi, \eta)$ ($\xi, \eta \in \widehat{\Pi}_{0,3}$),
satisfaisant les relations (133) (134), et tels que
de plus ("condition lumineuse")
les relations (132) (équivalentes à (133)) soient généra-
trices pour $\widehat{\Pi}_{0,3}$. Dans (134) $\nu$ est un élément
de $\widehat{\mathbf{Z}}$, évidemment uniquement déterminé,
et [unclear: inversible] : la condition $2\nu+1 \in \widehat{\mathbf{Z}}^*$\footnote{cette condition $2\nu+1 \in \widehat{\mathbf{Z}}^*$ est une conséquence de (132), (133), (134)), cf plus bas p. 409!}
(qui est impliquée par la condition lumineuse, et
peut-être l'implique ...)

Si l'inclusion évidente
$$ L_0 \hookrightarrow M_{0,3}^{!\sim}(\omega) \leqno (135) $$
$$ l_0^n = \lambda \mapsto u(\lambda) $$
est donnée par
$$ l_0^n = \lambda \mapsto (\xi = \lambda \rho(\lambda)^{-1} = l_0^n l_1^{-n}, \eta = \lambda \sigma(\lambda)^{-1} = l_0^n l_1^{-n}) \leqno (136) $$
et plus généralement l'action à gauche de $L_0$
sur $M_{0,3}^{!\sim}(\omega)$ est donnée, pour $\lambda = l_0^n \in L_0$, par
$$ (\xi, \eta) \mapsto (\xi' = l_0^n \xi l_1^{-n}, \eta' = l_0^n \eta l_1^{-n}) \leqno (137) $$
donc
$$ M_{0,3}^{!\sim} / L_0 = N_{0,3}^{!\sim} $$
est isomorphe au quotient de l'ensemble des
couples $(\xi, \eta)$ satisfaisant à (133) (134) [et la condition
lumineuse], par les opérations (137) --- compte tenu
d'une évidente structure de groupe naturelle ; d'ailleurs
cet ens. de couples $(\xi, \eta)$, donné par

\newpage

%% ===== Page 217 =====

$$
(\xi', \eta') (\xi, \eta) = (\xi'', \eta'')
\leqno{(138)}
$$
avec
$$
\xi'' = u'(\xi) \xi' \quad , \quad \eta'' = u'(\eta) \eta'
\leqno{(139)}
$$
$u'$ étant l'automorphisme défini par $(\xi', \eta')$.

Corollaire. Tout $u \in \mathcal{M}_{0,3}^{\sim !}$ s'écrit sous la forme
$$
u = f^{-1} u_{\xi, \eta} \quad , \quad = \operatorname{int}(f^{-1}) \circ u_{\xi, \eta}
\leqno{(140)}
$$
avec $f \in \hat{\Pi}_{0,3}$ et $(\xi, \eta)$ satisfaisant aux conditions du théorème précédent, de sorte qu'on aura
$$
u(\sigma_\infty) = h \sigma_\infty \quad , \quad u(\rho) = g \rho \quad \text{avec}
\leqno{(141)}
$$
$$
h = f^{-1} \eta \sigma_\infty(f) \quad , \quad g = f^{-1} \xi \rho(f)
\leqno{(142)}
$$
Pour qu'on ait
$$
f^{-1} u_{\xi, \eta} = f'^{-1} u_{\xi', \eta'}
$$
il faut et il suffit qu'il existe $n \in \hat{\mathbf{Z}}$ tel que
$$
\xi' = l_0^n \xi l_1^{-n} \quad , \quad \eta' = l_0^n \eta l_1^{-n} \quad , \quad f' = l_0^n f \quad .
\leqno{(144)}
$$

NB Comme la formule des lacets, une fois fixé le multiplicateur $\mu$ i.e.\ $\nu$ (cet un fait $\mu$ est déjà déterminé de façon unique par la seule connaissance de $\xi$, comme le montrent les calculs analogues dans $\operatorname{Sl}(2, \hat{\mathbf{Z}})'$ \dots), par (134) $\xi$ et $\eta$ se déterminent mutuellement, donc on peut éliminer l'un ou l'autre du système d'équations (133),(134), qui devient soit
$$
\xi \rho(\xi) \rho^2(\xi) = 1 \quad , \quad (\xi l_1^{-\nu}) \sigma_\infty(\xi l_1^{-\nu}) = 1 \quad \text{i.e.}
\leqno{(145)}
$$
$$
\sigma_\infty(\xi) = l_1^\nu \xi^{-1} l_0^{-\nu}
$$
soit
$$
\eta \sigma_\infty(\eta) = 1 \quad , \quad (\eta l_1^\nu) \rho(\eta l_1^\nu) \rho^2(\eta l_1^\nu) = 1 \quad .
\leqno{(146)}
$$

\newpage

%% ===== Page 218 =====

Le calcul adélique. On désigne encore par $f, g$ les images de $f, g$ dans $\operatorname{Sl}(2, \hat{\mathbf{Z}})'$, qui se relèvent en des éléments, notés encore $f, g$, de $\operatorname{Sl}(2, \hat{\mathbf{Z}})$. Il y a une ambiguïté de signes\footnote{vrai ?!... pour faire le choix de signes convenables pour $g$ et $g \rho$ (mais alors pourquoi la convention $y = \rho \sigma(g)$ ?...) par ailleurs $y = \rho_0(g)^{-1}$}, qu'on lève pour $f$ en imposant
$$
f \rho f \rho^2(f) = + 1 \quad (\text{dans } \operatorname{Sl}(2, \hat{\mathbf{Z}})) \leqno{(147)}
$$

pour $g$ : il n'y a pas de façon simple de lever l'ambiguïté pour le moment, et on admettra la possibilité d'écrire $g \in \hat{\Pi}_{0,3}$ sous la forme $\beta \sigma_\infty(\beta)^{-1}$ dont on aura a priori seulement
$$
g \sigma_\infty(g) = \pm 1 \leqno{(148)}
$$
avec possibilité de modifier le signe de $g$ [unclear: mais par un changement] $g \mapsto -g$ ! Posons
$$
f = \begin{pmatrix} A & B \\ C & D \end{pmatrix} , \quad g = \begin{pmatrix} R & S \\ T & U \end{pmatrix} \leqno{(149)}
$$

On aura
$$
AD - BC = 1 \qquad RU - ST = 1 \leqno{(149')}
$$
et la relation (147) s'exprimera alors :
$$
C + D - B = 1 \leqno{(150)}
$$

tandis (du moins dans chaque composante $\ell$-adique on devrait avoir soit cette relation, soit
$$
f_\ell = \begin{pmatrix} A_\ell & B_\ell \\ C_\ell & D_\ell \end{pmatrix} = -\rho^{-1} = \begin{pmatrix} 0 & -1 \\ 1 & -1 \end{pmatrix} ,
$$
satisfaisant prop. de la p.~394 - mais cette solution est incompatible avec la condition $f \equiv 1 \pmod{2}$, i.e.\ $A \equiv D \equiv 1 \pmod{2}$ et $B \equiv C \equiv 0 \pmod{2}$ $\dots$) donc on a (150) vraie dans $\hat{\mathbf{Z}}_\ell$ pour tout $\ell$, donc dans $\hat{\mathbf{Z}}$ $\dots$).

\newpage

%% ===== Page 219 =====

La relation (148) s'écrit encore
$$
\eta = \pm \sigma_{\infty}(\eta^{\sim})
$$
or $\sigma_{\infty}(\eta^{\sim}) = {}^t\eta$, donc (148) s'écrit
$$
\eta = \pm {}^t\eta \ ,
$$
mais si on avait le signe $-$ , on aurait $U = -U$
donc $2U=0$ , donc $U=0$ , ce qui est incompatible
avec $\eta \equiv 1 \ (2)$
qui implique $R, U \equiv 1 \ (2) \ , \ S, T \equiv 0 \ (2)$ . On a
donc
$$
\eta = {}^t\eta \quad \text{i.e.}
$$
$$
\eta \sigma_{\infty}(\eta) = + 1_2 \ , \quad \text{i.e.} \leqno{(151)}
$$
$$
S = T \leqno{(152)}
$$

La relation des lacets (138) s'écrit
$$
\begin{pmatrix} A & B \\ C & D \end{pmatrix} = \varepsilon \begin{pmatrix} R & S \\ T & U \end{pmatrix} \begin{pmatrix} 1 & 0 \\ 2\nu & 1 \end{pmatrix} \leqno{(153)}
$$
$$
= \varepsilon \begin{pmatrix} R+2\nu S & S \\ T+2\nu U & U \end{pmatrix} \qquad \text{avec } \varepsilon \in \{\pm 1\}\footnote{Obs. si $\varepsilon = 1$, en choisissant la base convenable (cf (155)) (la notation $\nu$ au lieu de $2\nu$ coïncide par hasard avec celle introduite dans le par. précédent où on prenait de plus symétriques les $\alpha,\beta$, elle parait plus logique).}
$$
et elle équivaut aux relations
$$
B = \varepsilon S \ , \quad D = \varepsilon U \leqno{(154)}
$$
$$
\mu D = 1 \qquad \text{où } \mu = 2\nu + 1 \leqno{(155)}
$$
équivalente aussi moyennant (154) , à
$$
\mu U = \varepsilon \qquad \text{(compte tenu de (152))} \leqno{(155 \text{ bis})}
$$
et cette dernière en fait qui exprime la
relation
$$
\overbrace{\varepsilon(T+2\nu U)}^{A'} + \overbrace{\varepsilon U}^{B'} - \overbrace{\varepsilon S}^{C'} = 1 \qquad \text{paraphrasant}
$$
$$
A + B - C = 1 \qquad \text{etc}
$$

\newpage

%% ===== Page 220 =====

elle ne fait qu'exprimer (compte tenu de (154) et de (150)) la relation
$$
C = \overbrace{\varepsilon (T+2\nu U)}^{C'} .
$$
Ainsi les relations (154) (155) expriment le fait que les trois coefficients sont égaux, c.à.d. $B=B'$, $C=C'$, $D=D'$
et l'égalité pour le dernier
$$
A = \overbrace{\varepsilon (R+2\nu S)}^{A'}
$$
s'en déduit alors grâce aux relations (149), qui donnent $AD-BC=1$, $A'D'-B'C'=1$, qui donnent
$$
A = \frac{1+BC}{D} , \quad A' = \frac{1+B'C'}{D'}
$$
(NB on sait que $D=D'$ est inversible, par (155))

donc $A=A'$.

On trouve ainsi le

Théorème. Pour une solution $(x,y)$ dans $\hat{\Pi}_{0,3} \times \hat{\Pi}_{0,3}$ du système d'équations (133)(134), désignant encore par $x,y$ des éléments de $\operatorname{Sl}(2,\hat{\mathbf{Z}})$ qui leur correspondent, l'ambiguïté de signe pour $y$ étant levée par la relation (147), on aura les relations (154) et (155), qui impliquent les conséquences suivantes :

\begin{enumerate}
\item[a)] $\mu \defeq 2\nu+1$ est inversible, et il est entièrement déterminé par $x$ (il est déterminé par $y$ seulement au signe près).
\end{enumerate}

\newpage

%% ===== Page 221 =====

Compte tenu des relations (149) (150)
$$
A D - B C = 1 \quad , \quad A + D - B = 1
$$
et du fait que $D$ est inversible,
les matrices $\xi$, $\eta$ sont entièrement déterminées par
les coeffs $B, D$ — les couples $(\xi, \eta)$
de matrices $\in \operatorname{Sl}(2, \hat{\mathbf{Z}})$, satisfaisant

\marginpar{Bien sûr, l'intro-\\duction des congruences\\ici est tout à fait postiche\\et accessoire -\\l'énoncé algébrique\\vaut indépendamment de\\telles conditions...}
$$
\begin{cases}
[unclear: \{ (\xi \eta^2)^2 = 1, \eta \xi \eta = 1, \eta = \xi l_\nu^{-1} \quad (\nu \in \hat{\mathbf{Z}} \text{ convenable}) \}] \\
\xi \equiv 1 \; (2), \quad \eta \equiv 1 \; (2)
\end{cases} \leqno (156)
$$

correspondent biunivoquement aux couples $(B,D)$, $B \in \hat{\mathbf{Z}}$, $D \in \hat{\mathbf{Z}}^\times$,\footnote{dont on déduit $B \equiv 0 \; (2)$} par
$$
\xi = \begin{pmatrix} A & B \\ C & D \end{pmatrix} , \quad \eta = \begin{pmatrix} R & S \\ T & U \end{pmatrix} , \quad \text{et } \nu \in \hat{\mathbf{Z}}
$$

$$
C = B - D + 1 \quad , \quad A = \frac{1 + BC}{D} = \frac{1}{D}(1 + B^2 - BD + B) \leqno (157)
$$
$$
= \mu(1 + B + B^2) - B
$$

$$
\mu = \frac{1}{D} \quad , \quad \nu = \frac{\mu - 1}{2} \leqno (158)
$$

\marginpar{on a calculé\\$\eta$ en f. de $\xi$\\et $\nu$ dans (156)}
$$
\eta = \xi l_\nu^{-1} = \begin{pmatrix} A & B \\ C & D \end{pmatrix} \begin{pmatrix} 1 & 0 \\ -2\nu & 1 \end{pmatrix} = \begin{pmatrix} A - 2\nu B & B \\ C - 2\nu D & D \end{pmatrix} \leqno (159)
$$

$$
U = D, S = B, R = A - 2\nu B = \mu(1 + B^2) = \frac{1 + B^2}{D}
$$

\marginpar{N.B. On déduit\\de la rel. que\\$\det u_\xi = \mu = \frac{1}{D}$\\et $\dots$\\ (161)}
Posons (comme p. 398' (78))
$$
u_\xi = \mu \begin{pmatrix} 1 & B \\ 0 & D \end{pmatrix} = \begin{pmatrix} \mu & \mu B \\ 0 & 1 \end{pmatrix} = \begin{pmatrix} \frac{1}{D} & B/D \\ 0 & 1 \end{pmatrix} \in \operatorname{Gl}(2, \hat{\mathbf{Z}}) \leqno (160)
$$

de sorte que la connaissance de $B, D$, i.e. de $u_\xi$,
équivaut à celle de $\xi \in \operatorname{Sl}(2, \hat{\mathbf{Z}})$, ou encore
à celle du couple $(\xi, \eta)$ ($\eta$ normalisé par
$\eta = \xi l_\nu^{-1}$). (Dans loc. cit. cet élément était
noté $u_{\alpha \beta}$ \dots je crois il ne dépend que de $\xi \in \hat{\pi}_{0,3}$,
et [unclear: l'enveloppe] de $\xi \in \operatorname{Sl}(2, \hat{\mathbf{Z}})$ --- et sa
connaissance équivaut à celle de $\xi \in \operatorname{Sl}(2, \hat{\mathbf{Z}})$. ---)

\newpage

%% ===== Page 222 =====

\footnote{Les formules (163) s'écrivent aussi $u_{\underline{\xi}}(\sigma_\infty) = \eta \sigma_\infty$, $u_{\underline{\xi}}(\rho) = \xi \rho \leqno{(163\text{ bis})}$, elles se déduisent de façon évidente de la condition à l'aide de $u_{\underline{\xi}}$.}
$$
u_{\underline{\xi}}(\sigma_\infty) = \eta \sigma_\infty = (\xi \rho_1^{-1}) \sigma_\infty \quad , \quad u_{\underline{\xi}}(\rho) = \xi \rho \leqno{(163)}
$$
$$
u_{\underline{\xi}}(\varepsilon_0) = \varepsilon_0^\mu \quad \text{et} \quad u_{\underline{\xi}}(l_0) = l_0^\mu \leqno{(164)}
$$

Si on considère une matrice inversible quelconque
$$
\underline{u} = \begin{pmatrix} L & M \\ N & P \end{pmatrix} \in \operatorname{Gl}(2, \hat{\mathbf{Z}})
$$
la relation
$$
\underline{u}(\varepsilon_0) = \varepsilon_0^\mu \quad \text{i.e.} \quad \underline{u} \varepsilon_0 = \varepsilon_0^\mu \quad (\text{où } \mu \in \hat{\mathbf{Z}}^*)
$$
s'écrit
$$
\begin{pmatrix} L & -L+M \\ N & -N+P \end{pmatrix} = \begin{pmatrix} L - \mu N & M - \mu P \\ N & P \end{pmatrix}
$$
et équivaut à :
$$
N = 0 \quad , \quad L = \mu P
$$
i.e.
$$
\underline{u} = \begin{pmatrix} \mu P & M \\ 0 & P \end{pmatrix} \quad , \quad \text{d'où } \det \underline{u} = \mu P^2
$$

la condition supplémentaire
$$
\det \underline{u} = \mu
$$
donne la condition
$$
P^2 = 1
$$
la solution la plus évidente est $P=1$, ce qui donne juste les matrices
$$
\underline{u} = \begin{pmatrix} \mu & M \\ 0 & 1 \end{pmatrix} = \underline{u}_{\underline{\xi}}
$$

\newpage

%% ===== Page 223 =====

où $\xi = \begin{pmatrix} A & B \\ C & D \end{pmatrix}$ satisfaisant à (411) est donné en termes de $B, D$ par $D = \frac{1}{\mu}$, $B = D M = M/\mu$.

Théorème. L'application\footnote{NB Du fait qu'ils sont naturels ! (168) (où on trouve $\mathcal{M}_{0,3}$ bien dans $\operatorname{Gl}(2, \hat{\mathbf{Z}})/\pm 1$, pas seulement dans $\operatorname{Gl}(2, \hat{\mathbf{Z}})$)}\footnote{NB Cette relation s'obtient plus bas par B ci-dessous.}
$$
\mathcal{M}_{0,3}^{\sim \wedge}[\infty] \isommap \operatorname{Gl}(2, \hat{\mathbf{Z}}) \leqno{(169)}
$$
$$
u_{\xi, \eta} \mapsto \underline{u}_{\xi} \quad, \quad \text{notée aussi } u \mapsto \underline{u}
$$
où $\underline{u}_{\xi}$ est donné par (160) est un hom. de groupes.

\underline{Dém.}

Tout d'abord $u$ a une image dans $\operatorname{Gl}(2, \hat{\mathbf{Z}})$, et si $x \in \hat{\Pi}_{0,3}$, $\underline{x}$ son image dans $\operatorname{Sl}(2, \hat{\mathbf{Z}})$ (dans $\hat{\mathfrak{T}}_{0,3}$) : on a 
$$
\underline{u(x)} = \underline{u}(\underline{x}) \quad \left( \defeq \underline{u} \, \underline{x} \, \underline{u}^{-1} \right) \leqno{(166)}
$$
car il suffit de le voir pour $x = \varpi_\infty$ et pour $x = \rho$, où cette formule se réduit à (161). Ceci noté,
si $u = u_{\xi, \eta}$, $u' = u_{\xi', \eta'}$ sont deux éléments de $\mathcal{M}_{0,3}^{\sim \wedge}(\infty)$, $u'' = u' u$ leur produit, on aura \leqno{(139)}
$$
\eta'' = u'(\eta) \eta'
$$
d'où
$$
\xi'' = \underline{u}'(\xi) \xi' \defeq \underline{u}'(\xi) \xi' \quad \text{i.e.}
$$
$$
\xi'' = \operatorname{int}(\underline{u}_{\xi'})(\xi) \xi' \leqno{(167)}
$$
et il s'agit de vérifier que si trois solutions $\xi, \xi', \xi''$ des équations
$$
\xi = \begin{pmatrix} A & B \\ C & D \end{pmatrix} , \quad \xi' = \begin{pmatrix} A' & B' \\ C' & D' \end{pmatrix} , \quad \xi'' = \begin{pmatrix} A'' & B'' \\ C'' & D'' \end{pmatrix}
$$
des équations $AD-BC=1$, $C+D B = 1$, $D \in \hat{\mathbf{Z}}^\times$,
sont reliées par (167), alors on a
$$
\underline{u}_{\xi''} = \underline{u}_{\xi'} \cdot \underline{u}_{\xi} \leqno{(168)}
$$

\newpage

%% ===== Page 224 =====

Comme $\xi$ et $\underline{u} = \underline{u}_{\xi}$ se déduisent mutuellement (162), il revient au m. de dire que si, sur le groupe
$$ \mathcal{B} \subset \operatorname{Gl}(2, \hat{\mathbf{Z}}) $$
$$ \mathcal{B} = \left\{ \begin{pmatrix} \mu & M \\ 0 & 1 \end{pmatrix} \mid \mu, M \in \hat{\mathbf{Z}}, \mu \text{ inv.} \right\} \leqno (169) $$
on définit (cf (162))
$$
\begin{cases}
F : \mathcal{B} \to \operatorname{Sl}(2, \hat{\mathbf{Z}}) \\
F_{\xi}(\underline{u}) = (\underbrace{\det \underline{u}}_{\mu})^{-1} [\underline{u}, \xi]
\end{cases} \leqno (170)
$$
alors on a la relation des 1-cocycles
$$ F_{\xi}(\underline{u}' \underline{u}) = \operatorname{int}(\underline{u}')(F_{\xi}(\underline{u})) \cdot F_{\xi}(\underline{u}') \leqno (171) $$
Par multiplicativité de $\det$, on en voit que
$$ \underline{u} \mapsto [\underline{u}, \xi] = \underline{u} \xi \underline{u}^{-1} \xi^{-1} $$
satisfait la relation des cocycles, ce qui est trivial.

Corollaire\footnote{Attention, ça ne dit rien du tout sur $\mathcal{M}_{0,3}(\underline{0})$ puisqu'on a supposé $L_0 = 1$ !! il faudrait donc refaire le calcul avec $L_0 \neq 1$ (i.e. $L_0$ qq) pour voir ce que ces relations donnent.} : L'application
$$ \mathcal{M}_{0,3} = \tilde{\mathcal{M}}_{0,3}[\underline{0}] \underset{\text{produit 1/2 directs}}{\cdot} \widehat{\mathfrak{S}}_{0,3} \to \operatorname{Gl}(2, \hat{\mathbf{Z}})' \leqno (172) $$
$$ \underline{u} \cdot a \mapsto (\underline{u}_{\xi} \cdot \underline{a} \pmod{\pm 1}) $$
et où $\underline{a}$ désigne l'image de $a$ dans $\operatorname{Gl}(2, \hat{\mathbf{Z}})'$, est un homomorphisme de groupes.

Résulte aussitôt du th. et des formules (163).

\newpage

%% ===== Page 225 =====

Dorénavant je vais repérer le couple $(\xi, \eta)$ d'un point de vue adélique (dans $\operatorname{Gl}(2, \hat{\mathbf{Z}})$) à l'aide de $\underline{u}_f$, qui est donc de la forme
$$ \underline{u} = \left(\begin{array}{cc} \mu & M \\ 0 & 1 \end{array}\right) \quad \text{avec } M \equiv 0 (2) \text{ si } \underline{u} \text{ provient de } \mathcal{M}_{0,3}^{\sim [\mathbf{0}]} \text{ (163 bis)} \leqno (173) $$
dont on tire $\xi, \eta$ par les formules
$$ \xi = [\underline{u}, \rho], \quad \eta = [\underline{u}, \sigma_{\infty}] \leqno (174) $$

Pour $\underline{u} \in \mathcal{M}_{0,3}^{\sim [\mathbf{0}]}$, défini par $(\xi, \eta)$, tel qu'il ne [unclear: faut pas] entrevoir la possibilité d'écrire $\xi$ et $\eta$ de façon unique comme $\xi = \alpha \rho(\alpha)^{-1}$, $\eta = \beta \sigma_{\infty}(\beta)^{-1}$ avec $\alpha \in \Pi_{0,3}^{\sim \wedge}$, $\beta \in \Pi_{0,3}^{\sim \wedge}$, les seuls invariants adéliques qu'on puisse lui associer directement de façon naturelle est le couple $(\xi, \eta)$, ou encore $\underline{u}_f$ qui revient au même, mais est plus joli. Les éléments $\lambda = \underline{l}_0^{-n} \in L_0 \subset \mathcal{M}_{0,3}^{\sim [\mathbf{0}]}$ correspondent à :

$$ \underline{\lambda} = \underline{\underline{l}}_0^{-n} = \left(\begin{array}{cc} 1 & -2n \\ 0 & 1 \end{array}\right) \leqno (175) $$

d'ailleurs on a
$$ \underbrace{\left(\begin{array}{cc} 1 & -2n \\ 0 & 1 \end{array}\right)}_{\underline{l}_0^{-n}} \underbrace{\left(\begin{array}{cc} \mu & M \\ 0 & 1 \end{array}\right)}_{\underline{u}_f} = \left(\begin{array}{cc} \mu & M-2n \\ 0 & 1 \end{array}\right) \leqno (177) $$

ce qui montre que la question des sous-groupes $B_i$ de $B$ (169) pose des

\newpage

%% ===== Page 226 =====

matrices $u$ telles que $u \equiv 1\ (2)$ i.e.\ $M \equiv 0\ (2)$
par le ss-groupe fermé $L_0$ engendré par $l_0 = \begin{pmatrix} 1 & -2 \\ 0 & 1 \end{pmatrix}$
s'identifie au groupe des matrices $\begin{pmatrix} \mu & 0 \\ 0 & 1 \end{pmatrix}$
tore maximal $T_0$ de $B_1$, de façon plus
précise on a une structure de produit semi-direct
$$
\underset{\displaystyle \left\{ \begin{pmatrix} \mu & M \\ 0 & 1 \end{pmatrix} \;\middle|\; \substack{\mu \in \hat{\mathbf{Z}}^\ast \\ M \in 2\hat{\mathbf{Z}}} \right\}}{B_1} \isommap \underset{\displaystyle \left\{ \begin{pmatrix} \mu & 0 \\ 0 & 1 \end{pmatrix} \;\middle|\; \mu \in \hat{\mathbf{Z}}^\ast \right\}}{T_0} \cdot \underset{\displaystyle \substack{ \left\{ l_0^n \mid n \in \hat{\mathbf{Z}} \right\} \\ \rotatebox{90}{$=$} \\ \left\{ \begin{pmatrix} 1 & -2n \\ 0 & 1 \end{pmatrix} \right\} }}{L_0} \leqno{(178)}
$$

Cela montre que si on s'intéresse aux
quotients $\mathcal{M}_{0,3}^{\sim \wedge} (0) / L_0 \isommap \mathcal{N}_{0,3}^{\sim \wedge}$, alors
le seul invariant adélique qu'on puisse lui
associer dans le contexte présent (quand on
ne travaille pas avec $\alpha, \beta$) est le multiplica-
teur $\mu$. Mais il y a mieux, on trouve grâce
à la décomposition (178) dans $\operatorname{Gl}(2, \hat{\mathbf{Z}})$
un scindage canonique de l'extension
$\mathcal{M}_{0,3}^{\sim \wedge} (0)$ de $\mathcal{N}_{0,3}^{\sim \wedge}$ par $L_0$ :

Proposition. Pour $\underline{u} \in \mathcal{N}_{0,3}^{\sim \wedge}$, il existe
un représentant unique $u \in \mathcal{M}_{0,3}^{\sim \wedge} [0]$ tel
que l'on ait $u \in T_0$ i.e.\ $u$ de la forme $\begin{pmatrix} \mu & 0 \\ 0 & 1 \end{pmatrix}$.

\newpage

%% ===== Page 227 =====

$$
\mathcal{M}_{0,3}^{!\sim}(0) \defeq \left\{ u \in \mathcal{M}_{0,3}^{!\sim}[0] \mid \underline{u} \in T_0 \text{ i.e. } \underline{u} = \begin{pmatrix} \mu & 0 \\ 0 & 1 \end{pmatrix} \right\} \leqno{(179)}
$$
est un sous-groupe section de l'extension $\mathcal{M}_{0,3}^{!\sim}[0]$ de $\mathcal{N}_{0,3}^{!\sim}$ par $L_0$, donc aussi de l'extension $\mathcal{M}_{0,3}^{!\sim}$ de $\mathcal{N}_{0,3}^{!\sim}$ par $\hat{\pi}_{0,3}$, ou de l'extension $\widetilde{\mathcal{M}}_{0,3}$ de $\widetilde{\mathcal{N}}_{0,3}$ par $\widehat{\mathfrak{S}}_{0,3}$.

Cela provient du fait que si $u \in \mathcal{M}_{0,3}^{!\sim}[0]$, $\underline{u} = \begin{pmatrix} \mu & M \\ 0 & 1 \end{pmatrix}$, on a donc $M \equiv 0 \pmod 2$, et si $M \in 2\hat{\mathbf{Z}}$, il existe un unique $n \in \hat{\mathbf{Z}}$ tel que $l_0^n \underline{u} = \begin{pmatrix} \mu & M-2n \\ 0 & 1 \end{pmatrix}$ soit dans $T_0$ i.e. tel que $l_0^n u \in \mathcal{M}_{0,3}^{!\sim}(0)$, savoir $n = \frac{M}{2}$.

Si $(\xi,\eta)$ satisfont aux conditions (133) (134) i.e. définissent un $u = u_{\xi,\eta}$ (cf [unclear: loc.] p. 407), lors on dira que $(\xi,\eta)$ est normalisé si on a $u_{\xi,\eta} \in \mathcal{M}_{0,3}^{!\sim}(0)$. De même, si nous prenons un élément quelconque de $\mathcal{M}_{0,3}^{!\sim}$ et l'écrit sous la forme
$$
u = f^{-1} \cdot u_{\xi,\eta} = f^{-1} u_f
$$
\footnote{c'est-à-dire de la condition univoque (136)} on dit que cette représentation est normalisée si $u_{\xi,\eta} = u_f$ est normalisé [i.e. si $\underline{u}_f = \begin{pmatrix} A & B \\ C & D \end{pmatrix}$, on a $B=0$]. Ainsi, tout élément $u$ de $\mathcal{M}_{0,3}^{!\sim}$ s'écrit de façon unique sous forme
$$
u = f^{-1} u_f \leqno{(180)}
$$

\newpage

%% ===== Page 228 =====

avec $f$ normalisé. En fait, ceci correspond
à la décomposition en produit $1/2$-direct
$$
\mathcal{M}_{0,3}^{!\sim \wedge} \simeq \mathcal{M}_{0,3}^{!\sim \wedge}(0) \cdot \Pi_{0,3}^\wedge
\leqno{(180)}
$$
- et de même on a
$$
\mathcal{M}_{0,3}^{\sim \wedge} \simeq \mathcal{M}_{0,3}^{\sim \wedge}(0) \cdot \widehat{\mathfrak{S}}_{0,3}^+
\leqno{(181)}
$$

[unclear: Rmq en passant : l'occasion moyen] de relier ces scindages canoniques en relation avec le choix d'une uniformisante en $0$ de $\mathbf{P}_{\mathbf{Q}}^1 = \overline{U}_{0,3 \mathbf{Q}}$.

Bien entendu, la même question se pose pour les "points à l'infini" $1, \infty$ de $U_{0,3 \mathbf{Q}}$, relativement aux sous-groupes sections
$$
\mathcal{M}_{0,3}^{!\sim \wedge}(q) \qquad (q \in \{1, \infty\})
\leqno{(182)}
$$
correspondant à ces points à l'infini - la construction d'un groupe-section $\mathcal{M}_{0,3}^{!\sim \wedge}(0)$ dans le centralisateur de $L_0$, s'applique évidemment [unclear: de même] pour $L_1, L_{\infty}$.

On a de même
$$
\rho(\mathcal{M}_{0,3}^{!\sim \wedge}(q)) = \mathcal{M}_{0,3}^{!\sim \wedge}(\rho(q)) \quad \text{pour } \rho \in \mathfrak{S}_3
\leqno{(183)}
$$
où les \underline{trois} $\mathcal{M}_{0,3}^{!\sim \wedge}(q)$ se déduisent de l'un quelconque d'entre eux en appliquant $\rho, \rho^2$.

\newpage

%% ===== Page 229 =====

47. | Retour sur l'homomorphisme $\widehat{\Pi}_{0,3} \stackrel{\hat{\varphi}}{\longrightarrow} \widehat{\mathfrak{S}}_{0,3}^{+}$
et la condition de commutation extérieure à celui-ci. Les groupes $\mathcal{M}_{0,3}^{!\sim}$, $\mathcal{M}_{0,3}^{\sim}$ et l'hom $\varphi \colon \mathcal{M}_{0,3}^{\sim} \to \mathcal{M}_{0,3}$.

Rappelons (§ 41) que $\hat{\varphi}$ est décrit par
$$
\hat{\varphi}(e_0) = \varepsilon_0 = \sigma_0 \rho \quad , \quad \hat{\varphi}(e_1) = \sigma_{\infty} \quad , \quad \hat{\varphi}(e_{\infty}) = \rho^{-1} \leqno (1)
$$

Soit
$$
u = u_{\alpha,\beta,f} \in \mathcal{M}_{0,3}^{!\sim}
$$
un autom. à lacets de $\widehat{\mathfrak{S}}_3$, noté encore $u$, défini par
$$
u(\sigma_{\infty}) = h \sigma_{\infty} = \operatorname{int}(\beta_0) \sigma_{\infty} \quad (h = \beta_0 \sigma_{\infty} \beta_0^{-1}, \beta_0 = f^{-1} \beta) \leqno (2)
$$
$$
u(\rho) = g \rho = \operatorname{int}(\alpha_0) \rho \quad (g = \alpha_0 \rho (\alpha_0)^{-1}, \alpha_0 = f^{-1} \alpha) \leqno (3)
$$

et écrivons que $u$ commute extérieurement à $\hat{\varphi}$, i.e. qu'il existe
$$
\lambda \in \widehat{\Pi}_{0,3} \leqno (4)
$$
tel que l'on ait
$$
\hat{\varphi} u(x) = \operatorname{int}(\lambda) u \hat{\varphi}(x) \quad \forall x \in \widehat{\Pi}_{0,3} \leqno (5)
$$

Il suffit de vérifier (5) pour $x = e_0$ et $x = e_1$.

1) Cas $x = e_0$, d'où
$$
u(x) = \operatorname{int}(f^{-1})(e_0^\mu)
$$
$$
\hat{\varphi}(u(x)) = \operatorname{int}(\hat{\varphi}(f)^{-1}) \hat{\varphi}(e_0^\mu) = \operatorname{int}(\hat{\varphi}(f)^{-1}) (\varepsilon_0^\mu)
$$
$$
u \hat{\varphi}(x) = u(\varepsilon_0) = \operatorname{int}(f^{-1}) (\varepsilon_0^\mu)
$$
donc la relation (5) s'écrit
$$
\operatorname{int}( [unclear: \lambda f \hat{\varphi}(f)^{-1}] ) (\varepsilon_0^\mu) = \varepsilon_0^\mu
$$

\newpage

%% ===== Page 230 =====

équivaut\marginpar{NB on n'a pas supposé a priori $\lambda \in \hat{\Pi}_{0,3}$ i.e.\ que $u \dots$} à $\hat{\varphi}(f) \lambda f^{-1} \in L_0 \defeq \operatorname{Centr}_{\hat{\Pi}_{0,3}}(\varepsilon_0^\mu)$
i.e.\ à l'existence de $r \in \hat{\mathbf{Z}}$ tel que l'on
ait $\hat{\varphi}(f) \lambda f^{-1} = \varepsilon_0^r$, i.e.
$$
\lambda = \hat{\varphi}(f)^{-1} \varepsilon_0^r f \leqno{(6)}
$$
Ceci permet donc d'éliminer $\lambda$ de
la deuxième relation :

2) Cas $x = l_1$, d'où
$$ u(x) = \operatorname{int}(f_1^{-1})(l_1^\mu) \quad (f_1 = g \rho(f^{-1})) $$
$$ \hat{\varphi}(u(x)) = \operatorname{int}(\hat{\varphi}(f_1)^{-1}) \sigma_\infty^\mu = \operatorname{int}(\hat{\varphi}(f_1)^{-1}) \sigma_\infty $$
(compte tenu que $\sigma_\infty^\mu = \sigma_\infty$ car $\mu \equiv 1 \, (2)$).
et $u \hat{\varphi}(x) = u(\sigma_\infty) = \operatorname{int}(\beta_0)(\sigma_\infty)$,

Donc la relation (5) s'écrit
$$ \operatorname{int}( \hat{\varphi}(f_1) \lambda \beta_0 ) \sigma_\infty = \sigma_\infty $$
ce qui signifie aussi qu'il existe $i \in \mathbf{Z}/2\mathbf{Z}$ tel
qu'on ait
$$ \hat{\varphi}(f_1) \lambda \beta_0 = \sigma_\infty^i \quad \text{soit} \quad \hat{\varphi}(f_1)^{-1} = \lambda \beta_0 \sigma_\infty^i $$
et remplaçant $\lambda$ par sa valeur (6), et
$f_1$ par sa valeur $g \rho(f^{-1})$, on trouve la
condition
$$ \hat{\varphi}(g \rho f^{-1})^{-1} = \hat{\varphi}(f)^{-1} \varepsilon_0^r f \beta_0 \sigma_\infty^i \quad \text{soit} $$
$$ \hat{\varphi}(f g \rho f^{-1}) = \varepsilon_0^r \beta_0 \sigma_\infty^i $$
$$ \xi = \alpha \rho(\alpha)^{-1} $$

\newpage

%% ===== Page 231 =====

Ainsi :

Pour que\footnote{NB [unclear: $\pi_A$ tout bêtement n'a de grand intérêt qd on le prend com. aucun des diviseurs].}
$$
\hat{\varphi} : \Pi_{0,3} \longrightarrow \widehat{\mathfrak{S}}_{0,3}^+
$$
corresponde extérieurement à l'élément $u \in \mathcal{M}_{0,3}^\sim$, correspondant aux invariants $g, h, f, [unclear: \alpha, \beta, \alpha, \beta, \xi, \eta]$ (cf $§ 46$), il faut et s. qu'on ait
$$
\hat{\varphi}(\xi) = \varepsilon_0^r \beta \sigma_\infty^{-\alpha} \leqno{(7)}
$$
pour $r \in \hat{\mathbf{Z}}$, $\alpha \in \mathbf{Z} / 2 \mathbf{Z}$ convenables.

On aura alors (5), avec $\lambda$ donné par (6) ; $\lambda = \hat{\varphi}(f)^{-1} l_0^{-r} f$.

Notons que $u$ ne définit le triplet $(\alpha, \beta, f)$ qu'à une transformation près de la forme
$$
(\alpha, \beta, f) \longmapsto (\alpha' = l_0^{-n} \alpha, \beta' = l_0^{-n} \beta, f' = l_0^{-n} f ) \leqno{(8)}
$$
où $n \in \hat{\mathbf{Z}}$. On aura alors (pour $\xi' = \alpha' \beta' (\alpha')^{-1} = l_0^{-n} \xi l_0^n$)
$$
\hat{\varphi}(\xi') = \varepsilon_0^{r'} \beta' \sigma_\infty^{-\alpha'} \leqno{(*)}
$$
avec $r' \in \hat{\mathbf{Z}}$, $\alpha' \in \mathbf{Z} / 2 \mathbf{Z}$ également bien déterminés. Or on aura
$$
\hat{\varphi}(\xi') = \hat{\varphi}(l_0^{-n} \xi l_0^n) = \hat{\varphi}(l_0^{-n}) \hat{\varphi}(\xi) \hat{\varphi}(l_0^n) = \varepsilon_0^{-n} \hat{\varphi}(\xi) \varepsilon_0^n
$$
et la relation $(*)$ s'écrit
$$
\hat{\varphi}(\xi) = \varepsilon_0^{r'+n} \beta' \sigma_\infty^{-\alpha'} \varepsilon_0^{-n}
$$
ou encore, comme $\beta' = l_0^{-n} \beta = \varepsilon_0^{2n} \beta$
$$
\hat{\varphi}(\xi) = \varepsilon_0^{r'+n} \beta \sigma_\infty^{-\alpha' -n}
$$
ce qui, comparé à (7), donne
$$
r' = r - n, \quad \alpha' = \alpha + n \mod 2 \leqno{(9)}
$$

\newpage

%% ===== Page 232 =====

Cela implique le

Corollaire ! Supposons que $u$ commute exté-
rieurement à $\hat{\varphi} : \Pi_{0,3} \longrightarrow \widehat{\mathfrak{S}}_{0,3}^+$. Alors
il est possible de façon unique de
décrire $u$ par un triple $(\alpha, \beta, f)$
(a priori déterminé modulo une transfor-
mation du type (8)), de façon à ce qu'on
ait (posant $\xi = \alpha g(\alpha)^{-1}$, comme à l'accou-
tumée) la relation
$$ \hat{\varphi}(\xi) = \beta \sigma_\infty^i \quad \text{pour } i \in \mathbf{Z}/2\mathbf{Z} \text{ convenable.} \leqno{(10)} $$
Un tel choix sera appelé \emph{normalisé}\footnote{Il faudrait examiner si cette définition coïncide avec celle de la p. 40', (9'), formulée sans faire intervenir $\hat{\varphi}$.}.

Ainsi les $u \in \mathcal{M}_{0,3}^{\sim}$ qui commutent
extérieurement à $\hat{\varphi}$ correspondent
biunivoquement aux triples $(\alpha, \beta, f)$, $\alpha \in \Pi_{0,3}^{\sim}$, $\beta, f \in \Pi_{0,3}$
satisfaisant, d'une part la relation des
lacets
$$ \xi = \eta l_1^\nu \leqno{(11)} $$
où $\xi = \alpha g(\alpha)^{-1}$, $\eta = \beta \sigma_\infty g(\beta)^{-1}$,
et $\nu \in \hat{\mathbf{Z}}$ convenable
(càd que $2\nu+1 \in \hat{\mathbf{Z}}^\times$)
d'autre part la relation (10) (pour $i \in \mathbf{Z}/2\mathbf{Z}$ convenable)
— l'une et l'autre relation ne faisant
d'ailleurs pas intervenir $f$.

\newpage

%% ===== Page 233 =====

Notons que $i \in \mathbf{Z}/2\mathbf{Z}$ est défini en termes de $\xi$ (compte tenu de $\beta \in \Pi_{0,\hat{3}} = \operatorname{Ker}(\hat{\mathfrak{S}}_{0,3}^+ \to \mathfrak{S}_3)$) en réduisant (10) mod $\Pi_{0,\hat{3}}$. Désignons par
$$
\hat{\varphi}^\circ : \Pi_{0,\hat{3}} \to \mathfrak{S}_3 \leqno{(12)}
$$
la composée
$$
\Pi_{0,\hat{3}} \overset{\hat{\varphi}}{\to} \hat{\mathfrak{S}}_{0,3}^+ \overset{\mathrm{can}}{\to} \mathfrak{S}_3,
$$
la formule (10) donne alors
$$
\hat{\varphi}^\circ(\xi) = \sigma_\infty^i \quad (i \in \mathbf{Z}/2\mathbf{Z}) \leqno{(13)}
$$
ce qui détermine bien $i \in \mathbf{Z}/2\mathbf{Z}$ en termes de $\xi \in \Pi_{0,\hat{3}}$, et implique une condition intrinsèque sur l'élément $\xi$ de $\Pi_{0,\hat{3}}$, à savoir que son image dans $\mathfrak{S}_3$ par $\hat{\varphi}^\circ$ soit dans le sous-groupe $\{1, \sigma_\infty\}$ de $\mathfrak{S}_3$. En fait, le noyau de $\hat{\varphi}^\circ$ est le sous-groupe invariant de $\Pi_{0,\hat{3}}$ engendré par $\ell_0^2, \ell_1^2, \ell_\infty^3$, et la condition envisagée signifie donc que modulo ce sous-groupe, $\xi$ soit de la forme $\ell_\infty^i$.

Par la formule (10), on peut éliminer $\beta$ de la donnée $(\alpha, \beta, f)$, et on peut remplacer la donnée $\alpha$ par la donnée $\xi$. Ainsi :

\newpage

%% ===== Page 234 =====

Corollaire 2. --- Les éléments du sous-groupe de $\mathcal{M}_{0,3}^{\sim}$ formés des $u$ qui commutent extérieurement à $\hat{\rho}$, est en corresp. biunivoque avec l'ensemble des couples $(\xi, f)$ (où $f \in \Pi_{0,\hat{3}}$) satisfaisant les conditions (ne faisant intervenir que $\xi$)
\marginpar{il faut de plus la condition ``éliminant'' $i$ :}
$$ \xi \rho(\xi) \rho^2(\xi) = 1 \leqno{(14)} $$
$$ \xi = (\hat{\varphi}(\xi) \sigma_{\infty}(\hat{\varphi}(\xi))^{-1}) l_1^\nu \qquad (\nu \in \mathbf{Z} \text{ convenable, tel que } 2\nu+1 \in \hat{\mathbf{Z}}^*) \leqno{(15)} $$

Pour un tel $\xi$, $\hat{\varphi}^{\bullet}(\xi)$ dans $\mathfrak{S}_3$ est de la forme $\sigma_{\infty}^i$, avec $i \in \mathbf{Z}/2\mathbf{Z}$, et posant
$$
\begin{aligned}
\beta &= \hat{\varphi}(\xi) \sigma_{\infty}^{-i} \quad (\in \Pi_{0,\hat{3}}) \\
\alpha &\in \Pi_{0,\hat{3}} \text{ l'unique élém. de } \Pi_{0,\hat{3}}^\sim \text{ [unclear: } \simeq \dots \text{] tel que } \alpha \rho(\alpha)^{-1} = \xi
\end{aligned} \leqno{(16)}
$$
$u$ se déduit\footnote{en fait $u$ se déduit ainsi de $\alpha,\beta$} de $\xi, f$ par les formules
$$
\begin{cases}
u(\sigma_{\infty}) = \operatorname{int}(f^{-1}\beta)(\sigma_{\infty}) = \operatorname{int}(f^{-1})(\operatorname{int}(\beta)(\sigma_{\infty})) = \operatorname{int}(f^{-1})(\xi l_1^{-\nu} \sigma_{\infty}) \\
u(\rho) = \operatorname{int}(f^{-1}\alpha)(\rho) = \operatorname{int}(f^{-1})(\operatorname{int}(\alpha)(\rho)) = \operatorname{int}(f^{-1})(\xi \rho)
\end{cases} \leqno{(17)}
$$

On notera que dans ces formules, $f$ n'intervient que par le fait de composer avec $\operatorname{int}(f)$. Si on [unclear: passe] donc, pour [unclear: $\mathcal{G}_{\Pi_{0,\hat{3}}}$]

\newpage

%% ===== Page 235 =====

satisfaisant (14) (15), par $u_\xi$\footnote{cette notation a déjà été malencontreusement utilisée car je l'ai utilisée pour désigner un des automorphismes $u_{g, h}$ dans la prop. qui précède... un abus de notations que j'espère sans gravité} l'autom. correspondant, en termes du Corollaire 2, au couple $(\xi, 1)$, de sorte que
$$
\begin{cases}
u_\xi(\sigma_\infty) = (\xi l_1^{-\nu}) \sigma_\infty \\
u_\xi(\rho) = \xi \rho
\end{cases}
\leqno{(18)}
$$
Les formules (17) s'écrivent aussi, simplement
$$
u_{\xi, f} = \operatorname{int}(f^{-1}) u_\xi .
\leqno{(19)}
$$
D'autre part les formules (5), (6) donnent :
$$
\hat{\varphi} \circ u_\xi = u_\xi \circ \hat{\varphi}
\leqno{(20)}
$$
où, bien sûr, dans le premier membre $u_\xi$ désigne la restriction de $u_\xi$ à $\Pi_{0,3}^{\wedge}$.

Notons

Corollaire 3. Les automorphismes\footnote{au sens de p. 416} (géométriques) de $\mathfrak{S}^{+\wedge}_{0,3}$ de la forme $u_\xi$ (i.e. exactement ceux pour lesquels on a la propriété formelle (20), i.e. qui commutent à $\hat{\varphi}$) et qui de plus normalisent $L_0$.

\newpage

%% ===== Page 236 =====

En effet, en termes des triples ``spéciaux'' $(\alpha, \beta, f)$ ces
propriétés s'expriment par les conditions
que d'une part $f$ de la forme $l_0^n$
(qui exprime la condition de normaliser $L_0$)
et que $\lambda$ donné par (6) (où $r=0$, i.e.
$$
\lambda = \hat{\varphi}(f)^{-1} f \leqno (21)
$$
soit égal à $1$, i.e. la condition
$$
\hat{\varphi}(f) = f \quad , \leqno (22)
$$
qui en l'occurrence s'écrit
$$
\varepsilon_0^n = \underbrace{l_0^n}_{\varepsilon_0^{2n}} \quad \text{i.e.} \quad \varepsilon_0^n = 1
$$
ce qui implique $f = l_0^n = 1$, donc
$u$ de la forme $u_{(\alpha, \beta, 1)} = u_{\xi}$, avec $\xi = \alpha p(\alpha)^{-1}$.

Corollaire 4 Les $u_{\xi}$ forment un
sous-groupe normalisé de $\operatorname{Norm}_{M_{0,3}^{\sim}}(L_0)$,
soit $M_{0,3}^{\sim (\infty)}$, dont l'intersection avec
$\Pi_{0,3}^{\wedge}$ se réduit à $\{1\}$, et dont l'image dans
$\mathcal{N}_{0,3}^{\sim}$\footnote{sous-groupe de $\mathcal{N}_{0,3}^{\sim}$} est formé des automorphismes
ext. à lacets $^+$ de $\hat{\mathfrak{S}}_{0,3}^+$, stabilisant
$\Pi_{0,3}^{\wedge}$ (i.e. $u \in \mathcal{M}_{0,3}^{\sim +}$) tels que $\hat{\varphi} : \Pi_{0,3}^{\wedge} \to \hat{\mathfrak{S}}_{0,3}^+$
commute extérieurement à $u$, on a
donc un isomorphisme
$$
M_{0,3}^{\sim (\infty)} = \left\{ u_{\xi} \mid \xi \in \Pi_{0,3}^{\wedge} \text{ solution de } (14), (15) \right\} \isommap \mathcal{N}_{0,3}^{\sim (\infty)} \leqno (23)
$$

\newpage

%% ===== Page 237 =====

Il reste seulement à vérifier qu'un $u_\xi$ ne peut être un automomorphisme intérieur que si $u_\xi = 1$ i.e.\ $\xi = 1$ (donc $\nu = 0$). Or $u_\xi$ est un $u_{\alpha, \beta} = u_{\alpha, \beta, 1}$ ($\alpha, \beta$ définis par $\alpha \rho(\alpha)^{-1} = \xi$, $\rho\sigma_0(\beta)\beta^{-1} = \xi l_1^{-\nu}$), qui est un automomorphisme intérieur si et seulement si on a $\alpha = l_0^n$, $\beta = l_0^n$ (cf.\ [unclear: bas] de la page 404), i.e.
$$ \xi = l_0^n l_1^{-n}, \quad \eta = \rho\sigma_0(l_0^n) l_0^{-n} = l_0^{-n} l_1^{-n}, $$
et la relation pour $\xi = \eta l_1^\nu$ s'écrit $l_1^\nu = 1$ i.e.\ $\nu = 0$. La relation (15) devient donc (pour
$$ \hat{\varphi}(\xi) = \hat{\varphi}(l_0^n l_1^{-n}) = \varepsilon_0^{-n} \sigma_\infty^{-n} $$
$$ \sigma_\infty(\hat{\varphi}(\xi)) = \varepsilon_1^n \sigma_\infty^{-n} \quad ) \text{ donne la relation} $$
$$ l_0^n l_1^{-n} = \varepsilon_0^{-n} \sigma_\infty^{-n} \varepsilon_1^{-n} \sigma_\infty^n, $$
i.e.\ (comme $l_0 = \varepsilon_0^2$, $l_1 = \varepsilon_1^2$) si $n$ pair
$$ \varepsilon_0^n = \varepsilon_1^{3n} \quad \text{d'où } n=0, $$
et si $n$ impair
$$ \varepsilon_0^n = \sigma_\infty(\varepsilon_1)^{-n} = \varepsilon_0^{-n} $$
i.e.\ $\varepsilon_1^{2n} = 1$, i.e.\ encore $n=0$, cqfd.

On a donc trouvé une section partiellement canonique\footnote{NB Il y a un [unclear: scindage] de $\mathcal{M}_{0,3}^{!\simeq} \to \mathcal{N}_{0,3}^{!\simeq}$ tout entier, voir p.\ suivante.} de $\mathcal{M}_{0,3}^{!\simeq}$ au dessus du sous-groupe $\mathcal{N}_{0,3}^{!\simeq}$ de $\mathcal{N}_{0,3}^{\sim}$. Comme

\newpage

%% ===== Page 238 =====

ce sous-groupe contient $\boldsymbol{\Gamma}_{\mathbf{Q}}$, on trouve donc une section de l'extension $\mathcal{N}_{0,3}^1$ de $\boldsymbol{\Gamma}_{\mathbf{Q}}$ par $\hat{\Pi}_{0,3}$, qui est en fait une section du sous-groupe de $\mathcal{N}_{0,3}^1$ normalisateur de $l_0$. Il faudrait vérifier que cette section canonique est associée au choix d'une uniformisante sur le point $0$ de $\mathcal{M}_{0,3,\mathbf{Q}}$, (i.e.\ d'un générateur de $\mathcal{M}_0/\mathcal{M}_0^2$ sur $\mathbf{Q}$, où $\mathcal{M}_0$ désigne l'idéal maximal de l'anneau local en $0$ dans $\mathcal{M}_{0,3,\mathbf{Q}} = \mathbf{P}_{\mathbf{Q}}^1 \setminus \{0,1,\infty\}$) --- il faudrait déterminer cet élément de $\mathcal{M}_0/\mathcal{M}_0^2$. ---

On a envie de prendre un élément qui soit, non seulement une base sur $\mathbf{Q}$, mais (vu la situation sur les entiers) une base sur $\mathbf{Z}$ --- condition qui détermine un élément de $\mathcal{M}_0/\mathcal{M}_0^2$ au signe près --- ce sera en fait le choix correspondant à l'uniformisante $z$, ou $-z$. Mais est-ce bien le bon ? Il faudrait y revenir !

\newpage

%% ===== Page 239 =====

Travailons désormais dans le sous-groupe $\mathcal{M}_{0,3}^{1\simeq}$ de $\mathcal{M}_{0,3}^{1\sim}$, formé des $u$ qui
commutent rationnellement à $\widehat{\varphi}$,
i.e. images inverses de $\mathcal{N}_{0,3}^{1\simeq}$, de sorte
qu'on a
$$ 1 \to \Pi_{0,\hat{3}} \to \mathcal{M}_{0,3}^{1\simeq} \to \mathcal{N}_{0,3}^{1\simeq} \to 1 \leqno{(24)} $$

On a donc trouvé un sous-groupe section
$$ \mathcal{M}_{0,3}^{1\simeq}(0) \isommap \mathcal{N}_{0,3}^{1\simeq} \leqno{(25)} $$
$$ \defeq \{ u \in \mathcal{M}_{0,3}^{1\simeq} \text{ sol. de } (14), \text{ et } (15) \} $$

En termes de celui-ci on peut donc écrire $\mathcal{M}_{0,3}^{1\simeq}$ comme un produit $1/2$ direct
$$ \mathcal{M}_{0,3}^{1\simeq} \simeq \mathcal{M}_{0,3}^{1\simeq}(0) \cdot \Pi_{0,\hat{3}} \leqno{(26)} $$
en écrivant tout élément de $\mathcal{M}_{0,3}^{1\simeq}$ de façon unique comme
$$ u = [unclear: u_f] \cdot u_0 \leqno{(27)} $$
$$ ( u_0 \in \mathcal{M}_{0,3}^{1\simeq}(0) \text{ i.e. } u_0 \text{ sol. de (14), (15)} ) $$
$$ ( [unclear: u_f] \text{ de (19), soit, } f \in \Pi_{0,\hat{3}} \text{ bien déterminé} ) $$

Ce sera [unclear: justement le]
$$ [unclear: u] = u_{\alpha, \beta} \cdot f \leqno{(28)} $$
avec $\alpha, \beta$ déterminés en termes de $f$ par
$$ \alpha(x) = [unclear: f \rho(x) f^{-1}] \quad , \quad \beta = [unclear: \dots] \leqno{(29)} $$

\newpage

%% ===== Page 240 =====

c'est ainsi qu'on peut interpréter les
triples $(\alpha, \beta, f)$ normalisés du cor.\ 1.

(Récapitulation)
Rappelons que le groupe $\mathcal{N}_{0,3}^{!\simeq} \simeq \mathcal{M}_{0,3}^{!\simeq}(0)$ est
en corr.\ 1-1 avec l'ens.\ des $\{ \xi \in \widehat{\Pi}_{0,3} \}$
satisfaisant les deux relations
$$
\begin{cases}
\xi \rho(\xi) \rho^2(\xi) = 1 \\
\xi = (\hat{\varphi}(\xi) \sigma_{\infty}(\hat{\varphi}(\xi))^{-1}) l_1^V \quad \text{pour } V \in \hat{\mathbf{Z}} \text{ conv.}
\end{cases} \leqno{(30)}
$$
(on pose $\xi = \alpha \rho(\alpha^{-1})$ , $\alpha \in \widehat{\Pi}_{0,3}^{\wedge}$ )

ou encore l'ens.\ des $\alpha \in \widehat{\Pi}_{0,3}^{\wedge}$ satisfaisant les
deux mêmes relations (30), dans lesquelles on remplace
$\xi$ par $\alpha \rho(\alpha)^{-1}$. L'élt $u_{\xi}$ de $\mathcal{M}_{0,3}^{!\simeq}(0)$
défini par $\xi$ est défini par les formules
(18), i.e.\ par
$$
u_{\xi}(\sigma_{\infty}) = \xi l_1^{-V} \sigma_{\infty} \quad , \quad u_{\xi}(\rho) = \xi \rho \leqno{(31)}
$$
En plus, si on a $u_{\xi}$ et $u_{\xi'}$, alors
$$
u_{\xi'} u_{\xi} = u_{\xi''} \leqno{(32)}
$$
où $\xi''$ est donné grâce aux formules
p.\ 395 (52) --- en notant qu'ici $\xi=g$, $\xi'=g'$, $\xi''=g''$ :

\newpage

%% ===== Page 241 =====

$$
\xi'' = u'(\xi) . \xi'
\leqno{(33)}
$$

Posons maintenant
$$
\mathcal{N}_{0,3}^{\approx} = \mathfrak{S}_3 \times \mathcal{N}_{0,3}^{!\simeq} \subset \mathcal{N}_{0,3}^{\simeq} = \mathfrak{S}_3 \times \hat{\mathcal{N}}_{0,3}^{!\simeq}
\leqno{(34)}
$$
et soit $\mathcal{M}_{0,3}^{\approx} \subset \mathcal{M}_{0,3}^{\simeq}$ l'image inverse de $\mathcal{N}_{0,3}^{\approx}$ dans $\mathcal{N}_{0,3}^{\simeq}$, de sorte qu'on a une structure d'extension
$$
1 \to \hat{\Pi}_{0,3}^{\sim} \to \mathcal{M}_{0,3}^{\approx} \to \underset{\simeq \mathfrak{S}_3 \times \mathcal{N}_{0,3}^{!\simeq}}{\mathcal{N}_{0,3}^{\approx}} \to 1
\leqno{(35)}
$$
contenant l'extension (24) comme [unclear: sous-extension]

$$
1 \to \hat{\mathfrak{S}}_{0,3}^+ \to \mathcal{M}_{0,3}^{\approx} \to \mathcal{N}_{0,3}^{!\simeq} \to 1
\leqno{(36)}
$$
contenant (24) comme sous-extension
= quotient par un groupe plus petit.
Le sous-groupe $\mathcal{M}_{0,3}^{!\simeq}(\omega)$, qui en est une section dans (24), est donc aussi une section dans (36), d'où un isomorphisme canonique
$$
\mathcal{M}_{0,3}^{\approx} \simeq \underset{\simeq \mathcal{N}_{0,3}^{!\simeq}}{\mathcal{M}_{0,3}^{!\simeq}(\omega)} \ltimes \hat{\mathfrak{S}}_{0,3}^+
\leqno{(37)}
$$
(produit semi-direct).

\newpage

%% ===== Page 242 =====

Considérons maintenant l'application
$$
\begin{tikzcd}[row sep=small, column sep=large]
\Phi: \mathcal{M}_{0,3}^{!\approx} \arrow[r] \arrow[d, phantom, sloped, "\simeq"] & \mathcal{M}_{0,3}^{\approx} \arrow[d, phantom, sloped, "\simeq"] \\
\mathcal{M}_{0,3}^{!\approx}(\omega) . \hat{\Pi}_{0,3} & \mathcal{M}_{0,3}^{!\approx}(\omega) . \widehat{\mathfrak{S}}_{0,3}^+
\end{tikzcd} \leqno{(38)}
$$
qui est donné par
$$
\Phi(u.a) = u \hat{\varphi}(a) \leqno{(39)}
$$
$(u \in \mathcal{M}_{0,3}^{!\approx}(\omega), a \in \hat{\Pi}_{0,3})$. Il résulte de la construction de $\mathcal{M}_{0,3}^{!\approx}(\omega)$ (n° 3) - plus particulièrement du fait que les éléments de $\mathcal{M}_{0,3}^{!\approx}(\omega)$ commutent effectivement (pas seulement extérieurement) à $\hat{\varphi}$, que $\Phi$ est un homomorphisme de groupes. Ce n'est autre bien sûr que (à restriction près aux sous-groupes $\mathcal{M}_{0,3}^{!} = \mathfrak{S}_{0,3}$ et $\mathcal{M}_{0,3}$) l'hom. $\Phi$ envisagé dans le § 11 (p. 316) - dans loc. cit. cet hom. pouvait se définir d'ailleurs de façon effective - par scindements : $\hat{\Pi}_{0,3} - \text{conj.}$ puis par : $\widehat{\mathfrak{S}}_{0,3}^+ - \text{conj.}$ puis - et un peu de soin devrait permettre de vérifier l'égalité de ces homomorphismes, $\Phi$ du § 11 et $\Phi$ de ici, de façon effective.

\newpage

%% ===== Page 243 =====

Calcul de $x = l_\infty$. \hfill 422

Je vais revenir sur la détermination des
$u_{\rho} \in \mathcal{M}_{0,3}^{\prime \sim}[0] = \operatorname{Norm}_{\mathcal{M}_{0,3}^{\prime \sim}}(L_0)$ qui sont normalisés
i.e. qui commutent à $\hat{\varphi}$ i.e. satisfaisant (5)
avec $\lambda = 1$. Pour $x = l_0$ cette formule, en vertu
de la forme particulière de $u_{\rho}$, est tautologique
($\hat{\varphi} u(l_0) = \hat{\varphi}(l_0 u') = \varepsilon_0 u'$, $u \hat{\varphi}(l_0) = u(l_0) = \varepsilon_0 u'$), donc on
voit qu'il suffit de la poser pour
$x = l_1$, ou pour $x = l_\infty$. Pour $x = l_1$, la condition
prend la forme (10) -- remarquable du fait
qu'elle donne un calcul de la dérivation
"délicate" $\beta$ (dont l'existence en effet,
en termes de $\rho$, repose sur des fondements
qui n'ont pas encore été établis...), en termes
des paramètres "grossiers" $\rho$ -- compte tenu
du fait que $t \in \hat{\mathbf{Z}}^*$ est déterminé en fait
aisément en termes de $\rho$.

Nous allons donner à cette condition une
forme un peu plus efficace, en
écrivant la relation de commutation (5)
pour $x = l_\infty$, au lieu de $x = l_1$. Notons d'abord
qu'on a la relation
$$
\tau'_\infty(\rho) = \rho^{-1} \leqno{(40)}
$$
Car $\tau'_\infty(\rho) = \sigma_\infty(\underset{=\rho}{\tau_\infty(\rho)}) = \text{[unclear: crossed out]} = \underset{\sigma_\infty \rho \sigma_\infty^{-1}}{\sigma_\infty(\rho)} = \text{[unclear: $\varepsilon_0^{-1} \sigma_0 \varepsilon_0^{-1}$]} = \rho^{-1}$ ,

\newpage

%% ===== Page 244 =====

il en résulte que si $u \in \mathcal{M}_{0,3}^{\prime \sim}$ satisfait
$$
u(\rho) = \rho^\mu \qquad \mu \in \hat{\mathbf{Z}} / 3 \hat{\mathbf{Z}} \text{ inversible} \leqno{(41)}
$$
alors on a nécessairement $\mu \equiv 1 \text{ ou } -1 \pmod{3}$
et si $\alpha_3(\mu) \in \mathbf{Z}/2\mathbf{Z}$ est l'entier
$$
\alpha_3(\mu) \in \mathbf{Z}/2\mathbf{Z} \text{ caractérisé par } \mu \equiv (-1)^{\alpha_3(\mu)} \pmod{3} \leqno{(42)}
$$
on aura donc $u(\rho) = \rho$ ou $u(\rho) = \rho^{-1}$ suivant que $\alpha_3(\mu) = 0$ ou $\alpha_3(\mu) = 1$, et on trouve
$$
u(\rho) = \tau_{\infty}^{\alpha_3} (\rho) \leqno{(43)}
$$
Ceci posé, faisons $x = l_{\infty}$ dans $(5)$ :
$$
\hat{\varphi} u (l_{\infty}) = \hat{\varphi} ( \operatorname{int}(f_{\infty}^{-1}) l_{\infty}^\mu ) = \operatorname{int} ( \hat{\varphi} (f_{\infty}^{-1}) ) \rho^{-\mu}
$$
(car $\hat{\varphi}(l_{\infty}) = \rho^{-1}$), et
$$
u ( \hat{\varphi} (l_{\infty}) ) = u (\rho^{-1}) = u(\rho)^{-1} = (\tau_{\infty}^{\alpha_3} (\rho))^{-1} = \tau_{\infty}^{\alpha_3}(\rho^{-1})
$$
donc la formule $(5)$ pour $\lambda = 1$ devient
$$
\operatorname{int} (\hat{\varphi}(f_{\infty}) \alpha) (\rho^{-1}) = \rho^{-\mu} \qquad \text{soit}
$$
$$
\operatorname{int} (\hat{\varphi}(f_{\infty}) \alpha) \rho = \rho^{\mu}
$$
et en vertu de $(43)$ on a $\rho^\mu = \tau_{\infty}^{\alpha_3}(\rho) \defeq \operatorname{int}(\tau_{\infty}^{\alpha_3}) (\rho)$
donc la condition devient
$$
\operatorname{int} (\tau_{\infty}^{\alpha_3} \hat{\varphi}(f_{\infty}) \alpha) \rho = \rho
$$
i.e. que l'élément $\tau_{\infty}^{\alpha_3} \hat{\varphi}(f_{\infty}) \alpha$ de $\widehat{\mathcal{G}}_{0,3}^{\prime \sim} = \widehat{\mathcal{G}}_{0,3}^{\sim} \cdot \{1, \tau_{\infty}\}$ commute à $\rho$.

\newpage

%% ===== Page 245 =====

\marginpar{[unclear: il faudrait y voir de plus près d'ailleurs, à vrai dire, dans le cadre profini (les principes généraux n'ont pas encore été précisés !!!)]}
d'après des principes généraux (pas encore établis, à vrai dire, dans le cadre profini) le commutant de $\rho$ dans le groupe en question est réduit au groupe engendré par $\rho$, $\{1, \rho, \rho^2\}$, de sorte que la relation vue prend la forme
$$
\tau_\infty^{\alpha_3} \hat{\varphi}(f_\infty) \alpha = \rho^j \quad \text{pour } j \in \hat{\mathbf{Z}} / 3 \hat{\mathbf{Z}} \text{ conv.}
$$
(bien déterminé).

On l'explicite en utilisant la valeur de $f_\infty$
$$
f_\infty^{-1} = g \rho(g) \rho^2(f_0^{-1}) = g \rho(g) = \rho^2(g)^{-1} = \rho^{-1}(g^{-1}) ,
$$
(NB ici $f_0^{-1} = 1$)
donc la relation devient
$$
\tau_\infty^{\alpha_3} \hat{\varphi}(\rho^{-1}(g)) \alpha = \rho^j
$$
soit
$$
\hat{\varphi}(\rho^{-1}(g)) = \tau_\infty^{-\alpha_3} \rho^j \alpha^{-1}
$$
(où $\tau_\infty^{-\alpha_3} \rho^j = \tau_\infty^{-\alpha_3}(\rho^j) \tau_\infty^{-\alpha_3}$)
comme $\tau_\infty(\rho) = \rho^{-1}$ on a $\tau_\infty^{\alpha_3}(\rho^j) = \rho^{(-1)^{\alpha_3} j}$,
$$
\hat{\varphi}(\rho^{-1}(g)) = \rho^{j'} \tau_\infty^{-\alpha_3} \alpha^{-1} \leqno{(44)}
$$
où $\alpha_3 = \alpha_3(\mu) \in \hat{\mathbf{Z}} / 2 \hat{\mathbf{Z}}$
défini par $(-1)^{\alpha_3} = \mu$ (3)
et $j' \in \hat{\mathbf{Z}} / 3 \hat{\mathbf{Z}}$ (bien déterminé).

Nous allons encore transformer un peu cette relation, en écrivant
$$
\tau_\infty^{-\alpha_3} = (\sigma_\infty \tau_\infty)^{-\alpha_3} = \sigma_\infty^{-\alpha_3} \tau_\infty^{-\alpha_3} \quad d'où
$$
$$
\tau_\infty^{-\alpha_3} \alpha^{-1} = \sigma_\infty^{-\alpha_3} (\alpha \tau_\infty^{\alpha_3})^{-1} \leqno{(\ast)}
$$

\newpage

%% ===== Page 246 =====

où
$$ \tilde{\alpha} = \alpha \tau_\infty^{\alpha_3} \leqno{(45)} $$
a l'avantage sur $\alpha$ lui-même d'être dans $\hat{\pi}_{0,3}$ dans tous les cas (NB on a $\alpha \in \hat{\pi}_{0,3} \cdot \{ 1, \tau_\infty \}$, et $\alpha \in \hat{\pi}_{0,3}$ ssi $\alpha_3 = 0$).
Il est égal à $\alpha$ ssi $\alpha_3 = 0$, et dans le cas $\alpha_3 = 1$ il est relié directement à $\xi$ par la formule
$$ \xi = \tilde{\alpha} h_1^{-1} \rho(\tilde{\alpha})^{-1} \leqno{(46)} $$
(déduite de (45), de $\xi = \alpha \rho(\alpha)^{-1}$, et de
$$
\tau_\infty \rho(\tau_\infty)^{-1} = h_1^{-1} \ ). \leqno{(* *)}
$$

Ainsi (44) prend la forme définitive\footnote{toutes ces formules restent valables en désignant par $g \mapsto g^{\pm 1}$ l'op.\ de multiplication par $\pm 1$ dans $\mathcal{G}_{\mathbf{Q}}$}
$$ \hat{\varphi}(\rho^{-1}(g)) = \rho^{J'} \tau_\infty^{\alpha_3} \tilde{\alpha}^{-1} \quad (\tilde{\alpha} = \alpha \tau_\infty^{\alpha_3}) \leqno{(47)} $$
qui peut s'écrire, en passant aux inverses
$$ \hat{\varphi}(\rho^{-1}(g^{-1})) = \tilde{\alpha} \tau_\infty^{\alpha_3} \rho^{-J'} \leqno{(47^{\text{bis}})} $$
ou encore, en résolvant en $\tilde{\alpha}$
$$ \underset{\substack{\parallel \\ \alpha \tau_\infty^{\alpha_3}}}{\tilde{\alpha}} = \hat{\varphi}(\rho^{-1}(g^{-1})) \rho^{J'} \tau_\infty^{\alpha_3} \leqno{(48)} $$

On rapprochera cette formule de la formule (10), qui s'écrivait en posant $\beta = \alpha_\infty^{-1}$
$$ \hat{\varphi}(g) = \beta \sigma_\infty^I \leqno{(49)=(10)} $$
$$ \beta = \hat{\varphi}(g) \sigma_\infty^{-I} \leqno{(50)} $$

\newpage

%% ===== Page 247 =====

Le fait que les formules (47) sont bien
équivalentes à (10), ou encore à (49),
lorsque $\alpha \in \hat{\pi}_{0,3}^{\sim}$ et $\beta \in \hat{\pi}_{0,3}$ sont liés
par $\xi = \eta l_1^{-1}$ (où $\xi = \alpha \rho(a)^{-1}$, $\eta = \beta \sigma_3(y)^{-1}$),
mériterait une vérification directe. L'intérêt
de (47), --- de façon à pouvoir résoudre en $\alpha$ (48),
tout comme de (10) $\overset{(49)}{\longleftrightarrow}$ (50), consiste en ce
que ces formules donnent les paramètres
"délicats" $\alpha, \beta$, en termes des "paramètres
grossiers" $\xi, \eta$. L'une et l'autre sont des
formules dans $\widehat{\mathfrak{S}}_{0,3}^+$ --- réduites modulo
$\hat{\pi}_{0,3}$, elles donnent des formules dans
$\widehat{\mathfrak{S}}_3$, savoir (compte tenu que $\tilde{\alpha}, \beta \in \hat{\pi}_{0,3}$) :

$$
\begin{cases}
\hat{\varphi}(\rho^{-1}(g_1)) = \hat{\rho} j' \hat{\sigma}_{\infty}^{\alpha_3} \\
\hat{\varphi}(\xi) = \hat{\sigma}_{\infty}^i
\end{cases} \leqno{(51)}
$$

ce qui redonne exactement, en termes
des éléments $\hat{\varphi}(g_1)$ et $\hat{\varphi}(\rho^{-1}(y))$ de
$\widehat{\mathfrak{S}}_3$, les valeurs de $i, \alpha_3 \in \mathbf{Z}/2\mathbf{Z}$ et de
$j' \in \mathbf{Z}/3\mathbf{Z}$\footnote{Il faudrait voir si $i, j'$ (supposés connus $\alpha_3$), s'explicitent en termes du multiplicateur $\mu$, voire en termes de [unclear: $p, q$].}.

\newpage

%% ===== Page 248 =====

restes quadratiques $\epsilon_l(\mu)$). On peut les
compléter par une formule déduite de
l'une et de l'autre, compte tenu que
$$
\left.\begin{aligned}
\xi \rho(\xi) \rho^{-1}(\xi) &= 1 \quad \text{i.e.} \quad \rho^{-1}(\xi)\xi = \rho(\xi)^{-1} \\
\rho^{-1}(\xi) &
\end{aligned}\right.
$$
on trouve donc, en multipliant (47) et (10),
$$
\hat{\varphi}(\underbrace{\rho^{-1}(\xi)\xi}_{\rho(\xi)^{-1}}) = \rho^i j' \sigma_\infty^{\alpha_3} \tilde{\alpha}^{-1} \beta \sigma_\infty^i
$$
soit encore
$$
\hat{\varphi}(\rho(\xi)) = \sigma_\infty^{-i} (\beta^{-1}\tilde{\alpha}) \sigma_\infty^{-\alpha_3} \rho^{-j'} \leqno{(52)}
$$
soit, sous forme résolue en $\beta^{-1}\tilde{\alpha}$:
$$
\underbrace{\beta^{-1}\tilde{\alpha}}_{(\beta^{-1}\alpha)^{\tau_\infty^{\alpha_3}}} = \sigma_\infty^i \hat{\varphi}(\rho(\xi)) \sigma_\infty^{\alpha_3} \rho^{j'} \quad . \leqno{(53)}
$$
On peut considérer que les trois
"paramètres délicats" principaux sont
les quantités $\beta, \alpha$ et $\beta^{-1}\alpha$, qui s'expri-
ment directement, de façon remarquable-
ment simple, en termes des valeurs
de $\hat{\varphi}$ pour $\xi$, $\rho^{-1}(\xi)$ et $\rho^{-2}(\xi) = \rho(\xi)$, ce
qui dispenserait sans doute (en

\newpage

%% ===== Page 249 =====

écrivant les calculs en termes de $\xi, \rho$)
exclusivement, ce qui demanderait un
peu plus de soins...) de faire appel à des
résultats délicats non encore démontrés,
du type "$1 = \alpha \rho(\alpha)$", "$\gamma = \beta \sigma_{\infty}(\beta)$",
pour justifier l'introduction de ces paramè-
tres $\alpha, \beta$.

\marginpar{(récapitulation)}
$$
\begin{cases}
\hat{\varphi}(\xi) = \beta \sigma_{\infty}^i \\
\hat{\varphi}(\rho^{-1}(\xi)) = \rho' \sigma_{\infty}^{\alpha_3} \tilde{\alpha}^{-1} \qquad (\tilde{\alpha} = \alpha \tau_{\infty}^{\alpha_3}) \\
\hat{\varphi}(\rho(\xi)) = \sigma_{\infty}^{-i} (\beta^{-1} \tilde{\alpha}) \sigma_{\infty}^{\alpha_3} \rho'^{-1}
\end{cases} \leqno{(54)}
$$

Supposons pour simplifier
$$
\alpha_3 = 0 \quad i.e. \quad \alpha = \tilde{\alpha} \quad i.e. \quad \alpha \in \Pi_{0,\hat{S}} \leqno{(55)}
$$
Alors, tirée de la deuxième formule
(54), en prenant les inverses
$$
\hat{\varphi}(\rho^{-1}(\xi^{-1})) = \alpha \rho'^{-1}
$$
et [unclear: le rapprochant] des deux membres
$$
\xi = \hat{\varphi}(\rho^{-1}(\xi^{-1})) \rho' \hat{\varphi}(\rho^{-1}(\xi))^{-1} \leqno{(56)}
$$
qui implique déjà la première formule
(14) des les formules fond.\ (14) (15).

\newpage

%% ===== Page 250 =====

On trouve donc dans ce cas ($\mu \equiv 1 \ (3)$ i.e.\ $\nu \equiv 0 \ (3)$) que les deux relations (14) (15) équivalent en fait à l'une des deux relations

\marginpar{valable d'ailleurs dès $\nu \equiv 0 \ (3)$ i.e. $\mu \equiv 1 \ (3)$ i.e. $\partial_y(u) \equiv 1$}
$$ \xi = \hat{\varphi}(\rho(q^{-1})) \rho(\hat{\varphi}(\rho(q^{-1})))^{-1} \leqno{(57)} $$
$$ \xi = \hat{\varphi}(q) \sigma_{\infty}(\hat{\varphi}(q))^{-1} h \leqno{(58)} $$
\marginpar{pour $\nu \equiv 0 \ (3)$}

qui caractérisent donc les $u = u_q \in \mathcal{M}_{0,3}^{! \approx}(0)$, dont le multiplicateur $\mu$ est $\equiv 1 \ (3)$.

Rmq. Le seul élément de $\mathcal{M}_{0,3}^{! \approx}(0) \simeq \mathcal{N}_{0,3}^{\sim}$ qui se calcule explicitement, en dehors de l'élément unité (et l'unique élément des $\mathcal{N}_{0,3}^{\sim} \supset \mathcal{N}_{0,3}^{! \approx}$) est l'élément $\tau \in \mathcal{T}_{0,3} \subset \mathcal{N}_{0,3}^{\sim}$ qui vient des $\mathcal{M}_{0,3}^{! \approx}(0)$ correspondant à l'élément $\tau_{\infty}$, posons $u = u_{\tau_{\infty}}$.

\marginpar{Attention, ici $\alpha_3 = 1$ donc (57) n'est pas valable pour $u = u_{\tau_{\infty}}$}
$$ h = \beta_0 = 1, \ f = 1, \ \alpha = \alpha_0 = \tau_{\infty}, \ \xi = g = l_1^{-2}, \ \eta = 1 \leqno{(59)} $$
$$ \nu = -1, \ \mu = -1 $$

ce qui satisfait bien à [unclear: les deux relations $\xi \dots$ (54)]

\newpage

%% ===== Page 251 =====

avec
$$
i = 1,\ \alpha_3 = 1 \text{ i.e. } \varepsilon_3 = -1,\ \gamma' = -1 \leqno (60)
$$

\marginpar{on refait un p'tit calcul dont je fais grâce, qui nous montre que ces [unclear: éventualités] multiplicatives sont bien une vue de l'esprit (cf. p. 423) i.e. : pour qu'on ait ceci il faut que les coeff. de la ligne sup. de $\xi$ soient de parité diff.}
\emph{Rmq} On peut se poser la question, de façon générale, de déterminer les éléments $u_{\alpha,\beta} = u_{g,\eta}$ de $\mathcal{M}_{0,3}^{1\sim}[0]$ (pas nécessairement normalisés dans un sens ou dans l'autre) tels qu'on ait (non seulement $f=1$, mais aussi) $\beta=1$ i.e. $\gamma=1$ (NB On sait déjà que si $\alpha=1$ i.e. $\xi=1$, alors $u=1$). On aura alors $\xi_g = \eta h^\nu = l_1^{2\nu}$, et il reste à vérifier
$$
\rho(\xi)\rho(\xi') = 1 \text{ avec } \xi = h^\nu,
$$
[unclear: car dans] $\operatorname{Sl}(2, \hat{\mathbf{Z}})'$, [unclear: on a vu] que
$$
\xi = \begin{pmatrix} 1 & 0 \\ 2\nu & 1 \end{pmatrix}
$$
(p. 248) ce qui signifie que l'on ait soit
$$
C+D-B = \pm 1 \quad i.e. \quad 2\nu+1 = \mu = \pm 1 \quad i.e.
$$
$\nu = 0$ ou $-1$ -- ce sont les deux cas triviaux $u_g = id$ et $u_g = \tau_\infty$.

\newpage

%% ===== Page 252 =====

Notons la\footnote{NB. Cette proposition a l'air de se relier avec le § 46 VI, où on l'étudie d'un point de vue [unclear: purement géom.]} \underline{Proposition} : Le commutant de $\tau$ (élément non trivial de $\mathfrak{T}_{0,3}^{!} \simeq \mathbf{Z}/2\mathbf{Z} \subset \mathcal{M}_{0,3}^{!\sim}$) dans $\mathcal{M}_{0,3}^{!\sim}$ est réduit à $\{1, \tau\}$.

On fera le calcul dans le groupe d'automorphismes $\mathcal{M}_{0,3}^{!\sim}(0)$ (où $\tau$ correspond à $\tau_\infty$), et en fait il se trouve que le centralisateur de $\tau$ dans $\mathcal{M}_{0,3}^{!\sim}[0]$ tout entier (le groupe sous-jacent à [unclear: celui de] $\mathcal{M}_{0,3}^{!\sim}$, formé des $u_\xi$ sans conditions de normalisation) est réduit à $\{1, \tau\}$. En effet, il faut exprimer la commutation d'un $u = u_\xi$ avec $u_{\xi_1} = \tau_\infty$, \quad $\xi' = \ell_1^{-1}$, on a
$$
\begin{cases}
u_\xi u_{\xi_1} = u_{\xi''} \text{ avec } \xi'' = u_{\xi_1}(\xi)\xi' \\
u_{\xi_1} u_\xi = u_{\xi'''} \text{ avec } \xi''' = u_\xi(\xi')\xi
\end{cases}
$$
et il faut exprimer $\xi'' = \xi'''$, i.e.
$$
\tau_\infty(\xi) \ell_1^{-1} = u_\xi(\ell_1^{-1}) \xi \leqno{(61)}
$$
Or on a
$$
u_\xi(\varepsilon_1) = \operatorname{int}(\xi)(\varepsilon_1^H)
$$
$$
u_\xi(\underbrace{\ell_1^{-1}}_{\varepsilon_1^2}) = \operatorname{int}(\xi)(\ell_1^{-H})
$$
d'où et (61) s'écrit
$$
\tau_\infty(\xi) \ell_1^{-1} = \xi \ell_1^{-H}
$$
i.e.
$$
\xi^{-1} \tau_\infty(\xi) = \ell_1^{1-H} = \ell_1^{-2\nu} \leqno{(62)}
$$

\newpage

%% ===== Page 253 =====

Cette équation a la solution
$$
\zeta = l_1^\nu \leqno{(63)}
$$
dans $\Pi_{0,3}^\wedge$ (car $\tau_{\infty}(l_1) = l_1^{-1}$ donc $\tau_{\infty}(\zeta) = l_1^{-\nu}$)
et cette solution est unique si on admet
que les résultats discrets se prolongent au
cas profini, de façon à avoir
$$
(\Pi_{0,3}^\wedge)^{\tau_{\infty}} = \{1\} \leqno{(64)}
$$
Or il faut de plus que $\zeta$ satisfasse la
relation
$$
\zeta \rho(\zeta) \rho^2(\zeta) = 1
$$
et on voit de suite que cela implique
$\nu = 0$, i.e. $u \in \{1, \tau_{\infty}\}$.

\newpage

%% ===== Page 254 =====

\section*{48. Construction d'un schéma en groupes.}

I) Étude de l'application $\alpha \mapsto \alpha p(\alpha)^{-1} = [\alpha, p]$.

Soit $G$ un schéma en groupes sur un anneau $A$ (p. ex. $A = \mathbf{Z}$, $G = \operatorname{Sl}(2)_{\mathbf{Z}}$ ou $\operatorname{Gl}(2)_{\mathbf{Z}}$), $p$ une section de $G$, et considérons le morphisme
$$
F_p : \alpha \mapsto \alpha p(\alpha)^{-1} = [\alpha, p] : G \to G \leqno{(1)}
$$
Soit
$$
T_p = \operatorname{Centr}_G(p) = G^p \subset G
$$
alors
$$
F_p(\alpha') = F_p(\alpha) \iff \alpha^{-1} \alpha' \text{ section de } T_p \leqno{(2)}
$$
donc on a un monomorphisme canonique induit par $F_p$
$$
G/T_p \hookrightarrow G \leqno{(3)}
$$
dont il s'agit de caractériser l'image.

Dans le cas qui nous occupe, le groupe quotient $G_{ab}$ est représentable (c'est le groupe unité pour $G = \operatorname{Sl}(2)_{\mathbf{Z}}$, et c'est $\mathbf{G}_m$ si $G = \operatorname{Gl}(2)_{\mathbf{Z}}$), et une première condition pour une section $\xi \in G(S)$, d'un $A$-schéma $S$, pour qu'elle soit dans l'image, c'est qu'on ait
$$
\delta(\xi) = 1 \leqno{(4)}
$$

\newpage

%% ===== Page 255 =====

$$
\delta : G \to G_{ab} \leqno{(5)}
$$
est l'homomorphisme canonique. Une deuxième condition, c'est que $\xi \rho$, qui doit être (localement) de la forme $\alpha \rho \alpha^{-1}$, soit conjuguée localement à $\rho$ — cette condition est d'ailleurs nécessaire de façon évidente, tautologiquement, pour que $\xi \rho$ appartienne à l'image de $F_\rho$, et implique donc la condition (4). Mais cette dernière peut servir parfois d'ingrédient, cela dans un [unclear: ensemble] de conditions de nature "calculatoire", qui soient nécessaires et suffisantes pour qu'on ait $\xi \rho \sim \rho$. Dans le cas d'un groupe réductif $G$, la condition de [unclear: conjugaison] locale de sections est d'ailleurs assez bien comprise, depuis notamment les travaux de Steinberg. On a un "schéma des invariants" $I_G$, et un morphisme canonique
$$
G \to I_G \leqno{(6)}
$$
(sur $A$, et commutant aux changements de base)
qui soit universel pour les morphismes de $G$ dans des schémas, qui soient invariants par automorphismes intérieurs - i.e.\ $I_G$ est l'enveloppe représentable de $G/\operatorname{int}(G)$ (faisceau des classes de \dots)

\newpage

%% ===== Page 256 =====

de conjugaison
$$
G/\operatorname{int}(G) \longrightarrow I_G . \leqno{(7)}
$$
A vrai dire, ce morphisme est un épimorphisme (pour la topologie fppf), mais pas un isomorphisme i.e.\ pas un monomorphisme --- en d'autres termes, pour deux sections (en l'occurrence $\rho$, $\xi\rho$) de $G$, il est bien nécessaire, mais pas toujours suffisant, que leurs images dans $I_G$ soient égales, pour qu'elles soient conjuguées localement fppf. Mais dans le cas où les deux sections sont régulières au sens de Steinberg (ce qui est une condition sur les fibres) il est suffisant. Elle l'est aussi si les deux sections sont $1/2$ simples (ce qui est également une condition sur les fibres).

Dans le cas d'un groupe localement isomorphe à $\operatorname{Sl}(2)$, ou $\operatorname{Gl}(2)$, les sections régulières au sens de Steinberg sont d'ailleurs celles qui n'ont aucun point où elles soient centrales --- ce qui est bien le cas pour le $\rho$ ou le $\xi_\sigma \rho$ des \S\ précédents. D'autre part, $I_G$ est alors bien connu, dans le cas de ces groupes.

\newpage

%% ===== Page 257 =====

déployé, plus généralement de $\operatorname{Sl}(E)$ ou $\operatorname{Gl}(E)$,
$E$ loc. libre de rang 2, il est can. isomorphe
à $E^1$ ou à $E^{1*} \times E^1$ [unclear: (convenons que le rang est $\geq 2$)]

dans le cas d'une forme de $\operatorname{Sl}(2)$ c'est
donné par
$$
G \longrightarrow I_G \simeq E^1
\leqno{(7)}
$$
$$
x \longmapsto \operatorname{Tr} x
$$
dans le cas d'une forme de $\operatorname{Gl}(2)$ c'est
donné par
$$
G \longrightarrow I_G = E^{1*} \times E^1
\leqno{(8)}
$$
$$
x \longmapsto (\det x, \operatorname{tr} x)
$$

On trouve donc (pour mémoire)

\underline{Proposition} Soit $G$ une forme de $\operatorname{Sl}(2)$ ou
$\operatorname{Gl}(2)$ sur un schéma $S$, soit $\rho$ une section
de $G$ qui n'est centrale en aucune fibre.
Pour qu'une section $\xi$ de $G$ soit (dans
l'image de $F_\rho$ ((1) ou (2))), il faut et il suffit (localement),
qu'elle satisfasse les conditions suivantes :
\begin{enumerate}
\item[a)] $\rho \xi \rho^{-1}$ n'est centrale en aucune fibre, i.e.
en aucune fibre on n'a $\xi = \lambda \rho^{-1}$, où
$\lambda$ soit un élément central
\item[b)] $\operatorname{tr} \xi \rho = \operatorname{tr} \rho$
\item[c)] (si $G$ forme de $\operatorname{Gl}(2)$) $\det \xi = 1$
i.e. $\xi$ section de $\operatorname{Sl}(2)$. \marginpar{équivaut à : \\ $\det(\xi \rho) = \det \rho$}
\end{enumerate}

\newpage

%% ===== Page 258 =====

930

\underline{Corollaire}. Quand ces conditions sont satisfaites, le schéma des solutions en $x$ de l'équation
$$
[\alpha, \rho] = \xi
$$
est un torseur (à droite) sous $T_\rho$, qui est un schéma en groupes lisse commutatif sur $S$, de dimension relative 1 (si $G$ forme de $\operatorname{Sl}(2)$) ou 2 (si $G$ forme de $\operatorname{Gl}(2)$).

\underline{Exemples}. $G = \operatorname{Gl}(2)_\mathbf{Z}$, $\rho = \begin{pmatrix} 0 & 1 \\ -1 & 1 \end{pmatrix}$ est l'élément de $\operatorname{Sl}(2, \mathbf{Z})$ déjà noté ainsi (à partir de la p. 335). Il est bien régulier, et de même l'élément $\sigma = \sigma_\infty = \begin{pmatrix} 0 & -1 \\ 1 & 0 \end{pmatrix}$.

De façon générale, la donnée d'une forme de $\operatorname{Gl}(2)$ équivaut à celle d'une forme de $\operatorname{Sl}(2)$, ou de $\operatorname{PGL}(2)$, ou encore à celle d'une algèbre d'Azumaya de rang 4 sur $S$, soit $\mathcal{M}$ -- i.e.\ une $\mathcal{O}_S$-algèbre localement isomorphe à $M_2(\mathcal{O}_S)$ -- par
$$
\begin{cases}
G(\mathcal{M}) = \text{schéma des unités de } \mathcal{M} \\
\operatorname{SG}(\mathcal{M}) = \text{[unclear: id.] de dét.\ 1} \\
\operatorname{PG}(\mathcal{M}) = G(\mathcal{M}) / \operatorname{Centr} \simeq G(\mathcal{M}) / \mathbf{G}_m
\end{cases} \leqno{(9)}
$$
Dans ce [unclear: dictionnaire] la donnée d'un [unclear: point] de $G$ équivaut à celle d'une section de $\mathcal{M}$, à

\newpage

%% ===== Page 259 =====

la condition de régularité signifie que la sous-algèbre monogène sur $\mathcal{O}_S$ (et par suite commutative [unclear: et stable au changement de base]) engendrée par $\rho$ sur chaque fibre est de dimension $2$ (ce qui implique que $\rho$ annule un polynôme quadratique)
$$
\mathcal{O}_S + \mathcal{O}_S \rho = A_\rho \leqno{(10)}
$$
qui est [unclear: localement libre en tant que] $\mathcal{O}_S$-module, i.e. à formation commutante au chgt. de base... Et on a
$$
T_\rho = \text{schéma des unités de } A_\rho \quad (\text{cas de } G = G(M)) \leqno{(11)}
$$
$$
T_\rho = \text{schéma des unités de } A_\rho \text{ de norme égale à } 1 \quad (\text{cas de } G = S G(M)) \leqno{(11 bis)}
$$

On notera que $T_\rho$ est un tore\footnote{(au sens usuel)} pour $\rho$ semi-simple, i.e. semi-simple sur chaque fibre, ce qui équivaut, puisqu'il est régulier, à dire que dans chaque fibre ses deux valeurs propres sont distinctes. Pour
$$
\rho = \begin{pmatrix} 0 & 1 \\ -1 & 1 \end{pmatrix} \quad (\rho^6 = 1, \rho^3 = -1) \leqno{(12)}
$$
ce qui est le cas en car. $\neq 3$ --- au sens strict $T_\rho$ en tant que schéma sur $\operatorname{Spec}(\mathbf{Z})$ n'est pas un tore -- il

\newpage

%% ===== Page 260 =====

devient un produit d'un tore par un groupe unipotent en car. $3$. De même, pour
$$
\sigma = \sigma_0 = \begin{pmatrix} 0 & -1 \\ 1 & 0 \end{pmatrix} \quad , \quad \sigma^4 = 1, \sigma^2 = -1 \leqno{(13)}
$$
$\sigma$ est semi-simple en toutes les car. sauf en car. $2$, où il est unipotent - et $T_\sigma$ est un tore sauf en car. $2$, où il devient un produit d'un tore par un unipotent.

\marginpar{Revenons à un cas général}
Quand on prend le centralisateur dans $S\mathfrak{G}(\mathcal{M})$, on trouve
$$
S T_\rho = T_\rho \quad , \quad S T_\sigma = \Ker \left( T_\sigma \xrightarrow{\substack{\text{hom.} \\ \text{induit par le} \\ \text{déterminant}}} G_m \right) \leqno{(14)}
$$

L'avantage de $T_\rho$ sur $S T_\rho$, c'est que tout torseur sous $T_\rho = \text{schéma des unités de } A_\rho$ est localement scindé au sens Zariski, tandis qu'il n'en est plus ainsi pour un torseur sous $S T_\rho$. Ainsi, quand on se propose de résoudre l'équation en $\alpha$
$$
[\alpha, \rho] = \gamma
$$
pour $\gamma$ donné satisfaisant les conditions de congruences, c'est possible (théor. 90) dans $\mathfrak{G}(\mathcal{M})$, mais peut-être pas dans $S\mathfrak{G}(\mathcal{M})$.

\newpage

%% ===== Page 261 =====

Revenons au cas particulier
$$
\rho = \begin{pmatrix} 0 & 1 \\ -1 & 1 \end{pmatrix} , \quad A = \mathbf{Z}
$$

alors
$$
A_\rho = \text{anneau des entiers du corps } \leqno{(14)}
$$
$$
\mathbf{Q}(\sqrt[3]{-1}), \quad \rho \text{ correspondant à } \exp(2i\pi/6) = \frac{1+i\sqrt{3}}{2}
$$

Quand on se donne $\xi$
$$
\xi = \begin{pmatrix} A & B \\ C & D \end{pmatrix} , \text{ d'où } \xi\rho = \begin{pmatrix} -B & A+B \\ -D & C+D \end{pmatrix} \leqno{(15)}
$$

les conditions b), c) de la proposition
deviennent
$$
\det \xi = 1 , \quad \operatorname{Tr}(\xi\rho) = \operatorname{Tr}\rho
$$
d'où

$$
AD-BC = 1 \qquad C+D-B = 1 \leqno{(16)}
$$

-- conditions bien familières ! Elles impliquent
la possibilité (sur une base quelconque)
de trouver loc. Zaris. un $\alpha \in \operatorname{Gl}(2)$ \marginpar{[unclear: donnant par descente de Weil un $A_\rho$-module... (*)]}
qui
vérifie
$$
\alpha \rho(\alpha)^{-1} = [\alpha, \rho] = \xi
$$
et sur $\mathbf{Z}$ c'est possible globalement,
car l'obstruction est un $A_\rho$-Module inversible, or
$A_\rho$ est principal i.e.
$$
\operatorname{Pic}(A_\rho) = 1 \leqno{(17)}
$$

\newpage

%% ===== Page 262 =====

Mais si $\alpha$ est une solution, on a a priori
$$ \det \alpha \in \mathbf{Z}^* = \{\pm 1\} \leqno{(18)} $$
et toute autre solution sera de la forme
$\alpha' = \alpha f$, avec $f \in A_\rho^*$, d'où
$$ \det \alpha' = \det \alpha \cdot \det f = \det \alpha $$
car $\det f = f \overline{f} > 0$ donc $\det f = +1$,
ainsi le signe dans (18) ne dépend pas
du choix de $\alpha$ --- s'il est $-1$, on ne peut
trouver de solution de l'équation ($\ast$) dans
$\operatorname{Sl}(2, \mathbf{Z})$. Ce signe ne dépend que de $\xi$ i.e.\ que
de $A,B,C,D$ --- et en fait on [unclear: voit] que de $B,D$
(qui déterminent, au signe près, $\xi$, grâce aux
formules (16)) --- les calculs faits par ailleurs
suggèrent que ce signe est égal au
reste de $D \bmod 3$ : $+1$ si $D \equiv 1\ (3)$, $-1$ si
$D \equiv -1\ (3)$.

Revenant au cas général, si on a trouvé
une solution
$$ [\alpha, \rho] = \alpha \rho(\alpha)^{-1} = \xi \qquad , \qquad \alpha \in \operatorname{Gl}(\mathcal{M}) \leqno{(19)} $$
pour [unclear: pouvoir] trouver une solution $\alpha'$ qui
soit $\in \operatorname{Sl}(\mathcal{M})$ i.e.\ $\det \alpha' = 1$, on note que
les solutions $\alpha' \in \operatorname{Gl}(\mathcal{M})$ sont alors données par
$$ \alpha' = \alpha f \qquad f \in \Gamma \mathcal{O}_S^* \quad , $$

\newpage

%% ===== Page 263 =====

et la condition cherchée
$$
\det \alpha' = 1 \quad \text{i.e.} \quad \det (\alpha f) = 1 \quad \text{i.e.}
$$
$$
\det(f) = \det \alpha
$$
donc la condition néc. et suffisante pour l'existence d'une solution, est que la section $\det \alpha$ de $A^*$ soit ds l'image de l'homomorphisme norme\marginpar{NB $S = \text{Spec } A$}
$$
\mathcal{N}_{A_{\rho}/A} : A_{\rho}^* \to A^* \leqno{(20)}
$$
i.e.\ l'obstruction est dans le groupe
$$
A^* / \mathcal{N}_{A_{\rho}/A} A_{\rho}^* \leqno{(21)}
$$
qui est un invariant particulièrement bien connu dans la théorie des corps de classes !

\marginpar{NB $\mathcal{N}(c\rho+d) = c^2+d^2-cd$}Sauf erreur, il résulte de la théorie des corps de classes que pour $\rho$ donné par (11), et pour $A = \mathbf{Z}_{\ell}$ ($\ell$ nb premier), en entiers $\ell$-adiques, ce groupe quotient est nul dans tous les cas, sauf pour $\ell=3$, où il est isomorphe à :
$$
\mathbf{Z}/2\mathbf{Z}. \qquad \text{[unclear: A biffer]}
$$

\marginpar{[unclear: bon à retenir]}Notons que (hors ds le cas particulier (12))
$$
\text{si} \quad \alpha = \begin{pmatrix} a & b \\ c & d \end{pmatrix} \leqno{(22)}
$$
des calculs immédiats déjà faits montrent que

\newpage

%% ===== Page 264 =====

$$
(ad-bc) D = (c^2 + cd + d^2) \leqno{(23)}
$$

donc $D$ est inversible $\Leftrightarrow c^2+cd+d^2 = \mathcal{N}_{A_{\rho}/A}(c \rho + d)$ l'est i.e. ssi $c\rho+d$ est inversible, et alors
\marginpar{Cas $c^2+cd+d^2 \equiv 1 (3)$}
$$
ad-bc = (c^2+cd+d^2)/D
$$
ce qui montre que la classe de $\det \alpha = ad-bc$ est alors égale à la valeur de $D \pmod 3$.

(Mais est-ce que $D$ peut être égal à $0 \pmod 3$ ?
donc $\pmod 3$ la matrice $\xi$ dépendrait des seuls paramètres $A,B$
$$
\xi = \begin{pmatrix} A & B \\ 1+B & 0 \end{pmatrix} \qquad -B(1+B)=1 \text{ i.e. } B^2+B+1=0
$$
$$
\text{donc } B \equiv 1 \pmod 3
$$
donc $\xi = \begin{pmatrix} A & 1 \\ -1 & 0 \end{pmatrix}$, cette matrice est régulière ssi $A \neq -1$, ok.) Donc a priori on ne peut exclure le cas où $D$ prendrait une valeur nulle $\pmod 3$... )

Lorsqu'on prend la section $\sigma$ de $\operatorname{Sl}(2, \mathbf{Z})$ donnée par
$$
\sigma = \sigma_{0} = \begin{pmatrix} 0 & -1 \\ 1 & 0 \end{pmatrix} \qquad \sigma_{0}^4=1, \sigma_{0}^2=-1 \leqno{(24)}
$$

on a encore (sans peine)
\marginpar{Cas $A = \mathbf{Z}$.}
$$
Pic(A_{\sigma}) = 0 \leqno{(25)}
$$

La condition numérique sur
$$
\eta = \begin{pmatrix} R & S \\ T & U \end{pmatrix} \leqno{(26)}
$$
pour être locale dans l'image

\newpage

%% ===== Page 265 =====

de $F_\sigma$ (en plus de la condition de régularité -
sur $\eta\sigma$)
$$
\eta = \beta \sigma \beta \sigma^{-1} = [\beta, \sigma] \leqno{(27)}
$$
c'est que l'on ait
$$
\det \eta = 1, \quad \operatorname{Tr}(\eta\sigma) = \operatorname{Tr}\sigma \leqno{(28)}
$$
or
$$
\eta\sigma = \begin{pmatrix} R & S \\ T & U \end{pmatrix} \begin{pmatrix} 0 & -1 \\ 1 & 0 \end{pmatrix} = \begin{pmatrix} S & -R \\ U & -T \end{pmatrix}
$$
et (28) devient
$$
RU - ST = 1, \quad S - T = 0 \quad \text{i.e.} \quad S = T \leqno{(29)}
$$
i.e. $\eta$ doit être \underline{symétrique}. Si cette
condition est satisfaite, il s'ensuit que (NB
- quand la base est $\mathbf{Z}$ - que l'on peut
trouver une solution $\beta \in \operatorname{Gl}(2, \mathbf{Z})$ de (27),
l'obstruction à trouver une solution dans
$\operatorname{Sl}(2, \mathbf{Z})$ est encore un signe, savoir un
élément dans
$$
\mathbf{Z}^* / \mathcal{N}(A^*) \simeq \pm 1
$$
- comme on a pour
$$
\beta = \begin{pmatrix} r & s \\ t & u \end{pmatrix} \leqno{(30)}
$$
la matrice satisfaisant (27), la relation :
$$
(ru - st) U = r^2 + u^2 = \mathcal{N}(r+ui) \leqno{(31)}
$$

\newpage

%% ===== Page 266 =====

on sait comment tester que $U$ est dans l'im.\ de $\mathcal{N}$, et que dans ce cas, l'obstruction dans $\mathbf{Z}$ (dans $\mathbf{Z}$, ou $\mathbf{Z}_l$) est le reste de $U \pmod 4$ \marginpar{car $\mathcal{N}(u) \equiv 1 \pmod 4$} (qui est $\equiv \pm 1$)...

Résumons pour mémoire

Proposition. Soit $A = \mathbf{Z}$, ou $\mathbf{Z}_l$, ou $\hat{\mathbf{Z}}$, et considérons les équations
$$[\alpha, \rho] \defeq \alpha \rho \alpha^{-1} = \xi$$
$$[\beta, \sigma] \defeq \beta \sigma \beta^{-1} = \eta$$
où
$$\rho = \begin{pmatrix} 1 & 1 \\ 0 & 1 \end{pmatrix}, \quad \sigma = \begin{pmatrix} 0 & -1 \\ 1 & 0 \end{pmatrix},$$
$\xi$ et $\eta \in \operatorname{Sl}(2, A)$ donnés, et on cherche $\alpha$ et $\beta$ dans $\operatorname{Sl}(2, A)$. On suppose que
$$ \xi = \begin{pmatrix} A & B \\ C & D \end{pmatrix} \quad \text{resp. } \eta = \begin{pmatrix} R & S \\ T & U \end{pmatrix} $$
satisfont les conditions nécessaires pour l'existence de ces éléments :
$$AD - BC = 1 \qquad B+D-C = 1 \leqno{(16)}$$
$$RU - ST = 1 \qquad S-T = 0 \text{ i.e. } S=T \leqno{(29)}$$
.

Alors :

1) Pour qu'il existe une solution $\alpha$ (resp. $\beta$) dans $\operatorname{Sl}(2, A)$, il faut et il suffit que $\xi \rho = \begin{pmatrix} A & A+B \\ C & C+D \end{pmatrix}$ soit régulier, i.e. [unclear: ne soit pas scalaire] (resp. que $\eta \sigma = \begin{pmatrix} S & -R \\ U & -T=-S \end{pmatrix}$ le soit) - et pour

\newpage

%% ===== Page 267 =====

ceci il suffit resp. que sur aucune fibre, $\xi$ (resp. $\eta$)
ne soit un multiple scalaire de $\rho^{-1}$ (resp. $\sigma^{-1}$)
- et pour ceci il suffit que $D$ (resp. $U$) soit
inversible (car le coeff. correspondant dans
$\rho^{-1} = \begin{pmatrix} 1 & -1 \\ 1 & 0 \end{pmatrix}$ resp. $\sigma^{-1} = \begin{pmatrix} 0 & 1 \\ -1 & -1 \end{pmatrix}$ est zéro).

2) Lorsque $D$ (resp. $U$) est inversible, il y a
toujours une solution $\alpha$ resp. $\beta$, avec $\det\alpha \in \pm 1$
(resp. $\det\beta \in \{\pm 1\}$).

3) Pour qu'il y ait une solution dans
$\operatorname{Sl}(2, A)$, il faut et il suffit qu'on ait $D \equiv 1\ (3)$
(resp. $U \equiv 1\ (4)$) - condition qui est
[unclear: automatique] si $A = \mathbf{Z}_\ell$ avec $\ell \neq 3$ (resp.
avec $\ell \neq 2$).

\marginpar{Étude d'abord de l'équation}Étude de l'équation
$$
[\alpha,\rho] = [\beta,\sigma] \ \varepsilon \ (\in \{\pm 1\}) \leqno{(32)}
$$

On suppose par ex. que $G$ soit une forme
de $\operatorname{Gl}(2)$, associée à une algèbre d'Azumaya
$\mathcal{M}$ sur la base $S = \operatorname{Spec} A$, et qu'on
se donne des sections $\rho, \sigma$ régulières de $G$ i.e.
deux sections inversibles de $\mathcal{M}$. On se donne
de plus un ``sous-groupe à 1 paramètre''
(additif) de $G$, i.e. un homomorphisme
$$
\mathbb{G}_a \to G \ , \quad \lambda \mapsto \varepsilon(\lambda) \leqno{(33)}
$$

\newpage

%% ===== Page 268 =====

Le cas qui nous intéresse plus particulièrement
est celui où $G = \operatorname{Gl}(2)_{\mathbf{Z}}$, $\rho$ et $\sigma$ comme
dessus, et
$$ \varepsilon_1(\lambda) = \varepsilon_1^\lambda \equiv \begin{pmatrix} 1 & 0 \\ \lambda & 1 \end{pmatrix} \leqno{(34)} $$
On notera que l'on a
$$ \varepsilon_1 = \varepsilon_1(1) = \begin{pmatrix} 1 & 0 \\ 1 & 1 \end{pmatrix} = \rho \sigma^{-1} \leqno{(35)} $$

Nous considérons (32) comme des équations reliant les quantités variables $\alpha, \beta, \mu$ - $\alpha, \beta$ sections de $G$, $\mu$ section de [unclear: $\mathbf{G}_a$] - les hypothèses que nous allons faire impliqueront d'ailleurs que $\mu$ sera en fait une section de $\mathbf{G}_m$ (s'il [unclear: existe $\alpha, \beta$ vérifiant] la relation (32)).

On commencera par étudier les syst. $(\xi, \eta, \mu)$ satisfaisant à
$$ \xi = \eta \varepsilon_1(\mu - 1) \leqno{(33)} $$
$$ \det \xi = \det \eta = 1 \leqno{(34)} $$
\marginpar{NB On aura $\det \varepsilon_1(\mu-1)=1$ donc $\det \xi = \det \eta$ [unclear: automatiquement]}
$$ \xi \rho = \rho \xi, \quad \eta \sigma = \sigma \eta \leqno{(35)} $$

On peut éliminer $\eta$ de ces équations
grâce à (33), on trouve

\newpage

%% ===== Page 269 =====

(36) $$det \xi = 1, \quad Tr \xi \rho = 0, \quad Tr((\xi \varepsilon_1 (1-\mu) - 1)\sigma) = 0$$.

posons pour simplifier
$$i = Tr \rho \qquad j = Tr \sigma$$

La troisième relation (36), qui en développant par rapport aux coefficients de $\xi$, peut s'écrire
(37) $$P\mu - Q = 0$$

Les coefficients $P, Q$ sont des fonctions linéaires affines par rapport aux composantes de $\xi$, et données par l'identité (en $\xi, \mu$)
(38) $$Tr((\xi \varepsilon_1(1-\mu) - 1)\sigma) = P(\xi)\mu - Q(\xi)$$

Le terme constant $-Q(\xi)$ est égal (s'obtient en prenant la valeur du premier membre de (38) pour $\mu = 0$)
i.e.
(39) $$-Q(\xi) = Tr((\xi \varepsilon_1 - 1)\sigma)$$
$$= Tr(\xi \varepsilon_1 \sigma) - Tr \sigma$$

On suppose qu'on a [MARGIN: $\lambda$ section centrale de $G$ i.e. $\lambda$ scalaire inversible]
(40) $$\varepsilon_1 \sigma = - \lambda \rho$$
de sorte que
(40 bis) $$\varepsilon_1 = - \lambda \rho \sigma^{-1}$$ [MARGIN: NB Si $det \rho = det \sigma = 1$, alors on doit avoir $\lambda^2 = 1$ donc si $A$ intègre, $\lambda \in \{ \pm 1 \}$]
et (39) devient,
$$Tr(\xi \varepsilon_1 \sigma) = Tr(\xi(-\lambda \rho)) = - \lambda Tr \xi \rho$$ [unclear]

\newpage

%% ===== Page 270 =====

on trouve
$$
Q(\xi) = \underbrace{\lambda_0 \Tr\rho + \Tr\sigma}_{c} \quad (\text{constante, indépendante de } \xi) \leqno{(41)}
$$

et on suppose
$$
c = \lambda_0 \Tr\rho + \Tr\sigma \quad \text{scalaire inversible de } \underline{O}_S \leqno{(42)}
$$

(Dans le cas qui nous intéresse, on aura
$$
\sigma^{-1} = -\sigma, \quad \xi_1 = \rho\sigma = -\rho\sigma^{-1}, \quad \text{i.e. } \lambda_0 = 1
$$
Et $c = \lambda_0 \underbrace{\Tr\rho}_{1} + \underbrace{\Tr\sigma}_{0} = 1 \quad \text{ok.}$ )

Ceci posé, la troisième équation (36) devient (compte tenu des deux premières)
$$
P(\xi) \mu = c \leqno{(43)}
$$

ce qui signifie en soi que $P(\xi)$ est inversible, et
$$
\mu = c / P(\xi) \ . \leqno{(44)}
$$

Ainsi, $\mu$ est déterminé de façon unique par $\xi$, par l'équation (44) - et $\mu$ sera nécessairement inversible.

En résumé :

Proposition. Si on suppose satisfaite (40 bis) pour l'élément universel $\xi_1 = \xi_1(\underline{t})$, et que l'on ait (42) ($\lambda_0 \Tr\rho + \Tr\sigma$ inversible ), alors

\newpage

%% ===== Page 271 =====

l'ensemble des solutions en $(\xi, \eta, \mu)$ des
équations simultanées (33)(34)(35)
correspond 1-1 avec l'ensemble des $\xi$
dans $G = G(\mathcal{M})$ satisfaisant les
conditions
$$
\begin{cases}
\det \xi = 1, \Tr \xi \rho = \Tr \rho \text{, et} \\
P(\xi) \text{ section inversible de } \underline{\mathcal{O}}_S.
\end{cases}
\leqno{(45)}
$$
En termes de $\xi$, la solution correspondante du
système (33 à 35) est donnée par
$$
(\xi, \eta = \xi \xi_1(1-\mu), \mu = c/P(\xi))
$$
- et $\mu$ (le « multiplicateur ») est
nécessairement inversible.

NB. Si on avait résolu en $\eta$, on aurait
trouvé une équation de la forme
$$
P'(\eta)\mu - Q'(\eta) = 0
$$
avec
$$
Q'(\eta) = \Tr \eta \varepsilon_0 \rho - \Tr \rho
$$
et la condition naturelle serait
$$
\varepsilon_0^{-1} \rho = - \lambda_0' \sigma^- \quad (\lambda_0' \text{ scalaire inversible}) ,
$$
qui est équivalente à (40 bis) avec $\lambda_0' = 1/\lambda_0$ ,
et implique que
$$
Q'(\eta) = \lambda_0' \Tr \sigma - \Tr \rho = c' ,
$$

\newpage

%% ===== Page 272 =====

on voit tout de suite que
$$
c' = \lambda'_0 c
$$
donc la condition $c$ inv. équivaut à $c'$ inv., et on a alors le système d'équations en $\eta$
$$
\det \eta = 1,\ \operatorname{Tr} \eta \sigma = \operatorname{Tr} \sigma^{-},\ P'(\eta) \mu = c' \quad (\text{i.e. } \mu = c' / P'(\eta))\footnote{$P'(\eta)$ inv.}
$$

On voit donc que les hypothèses faites sont en fait symétriques en $\xi,\eta$, et permettent un traitement symétrique de la résolution, soit via $\xi$, soit via $\eta$.

Détermination de $P(\xi)$. On a
$$
\varepsilon_1(\lambda) = 1 + N_1 \lambda = (1+N_1)^\lambda
$$
$N_1$ une matrice de carré nul
$$
N_1^2 = 0,
$$
et
$$
\operatorname{Tr}(\xi \varepsilon_1(1-\mu) \sigma) - \operatorname{Tr} \sigma = \operatorname{Tr} \xi \sigma - \operatorname{Tr} \sigma + (1-\mu) \operatorname{Tr} (\xi N_1 \sigma)
$$
donc
$$
P(\xi) = \operatorname{Tr} (\xi (-N_1 \sigma)) \leqno (47)
$$

D'ailleurs la formule
$$
\varepsilon_1 = 1 + N_1 \leqno (48)
$$
et la formule (40 bis) s'écrit
$$
\sigma + N_1 \sigma = - \lambda_0 \rho^- \leqno (49)
$$
$$
N_1 = -(\lambda_0 \rho^- \sigma^{-1} + 1)
$$
soit [unclear: de nouveau]

\newpage

%% ===== Page 273 =====

$$
N_1\sigma = - (\lambda_0\rho + \sigma)
$$

donc on aura

$$
P(\xi) = \operatorname{Tr} \xi (\lambda_0\rho + \sigma) \leqno{(50)}
$$

\underline{Corollaire de la Proposition}. Le schéma des solutions des équ.\ simultanées (33), (35) en $\xi, \eta, \mu$ est un schéma \underline{relatif} lisse, de dimension relative 2 -- car l'équation $\operatorname{Tr}(\xi\rho) = \text{cte}$ définit toujours un sous-schéma lisse de dim 1 de $\operatorname{Sl}(\mathcal{M})$ (dès que $\rho$ est une section qui n'est nulle en aucune fibre), -- dont il faut prendre ensuite l'ouvert défini par "$P(\xi)$ inversible".

Dans le cas qui nous intéresse,
$$
\rho = \begin{pmatrix} 0 & 1 \\ -1 & 1 \end{pmatrix}, \quad \sigma = \begin{pmatrix} 0 & -1 \\ 1 & 0 \end{pmatrix},
$$
la condition $P(\xi)$ inv.\ s'écrit
pour
$$
\xi = \begin{pmatrix} A & B \\ C & D \end{pmatrix}, \quad \eta = \begin{pmatrix} R & S \\ T & U \end{pmatrix}
$$
on aura
$$
P(\xi) = D \quad , \quad P(\eta) = U
$$
$$
\eta \varepsilon_1(1-\mu) = \begin{pmatrix} R + (\mu-1)S & S \\ T + (\mu-1)U & U \end{pmatrix} \leqno{(51)}
$$

\newpage

%% ===== Page 274 =====

et l'équation
$$ \xi = \eta \varepsilon_1(\mu-1) $$
équivaut à
$$ B = S, \quad D = U \leqno{(52)} $$
et en fait elle équivaut à l'une de ces
deux équations et à
$$ \mu D = 1 \quad (\Rightarrow \mu U = 1) \leqno{(53)} $$
quand on suppose déjà que $\xi, \eta$ satisfont
à des équations de traces et de déterminants
$$ \begin{cases} A D - B C = 1 & R U - S T = 1 \\ A+D-B = 1 & T-S = 1 \end{cases} \leqno{(54)} $$

Le schéma des solutions des
équations simultanées (33), (35),
est celui des $(D,B)$ avec $D$
inversible, ou encore des $(S,U)$ avec
$U$ inversible. On peut l'identifier
au schéma des matrices
$$ \begin{pmatrix} \mu & M \\ 0 & 1 \end{pmatrix} = \begin{pmatrix} 1/D & B/D \\ 0 & 1 \end{pmatrix} \leqno{(55)} $$
(cf. p. 410 - ), qui normalisent le
sous-groupe à 1 paramètre
$$ \varepsilon_1(\lambda) = \begin{pmatrix} 1 & -\lambda \\ 0 & 1 \end{pmatrix} = \rho^{-1} \varepsilon_0(\lambda) \rho. $$

\newpage

%% ===== Page 275 =====

Revenons au cas général, où on suppose néanmoins que $\varepsilon_1$ donné par (40 bis),
$$
\varepsilon_0 = \rho^{-1} \varepsilon_1 \rho = -\lambda_0 \sigma^{-1} \rho = 1 + N_0 \leqno{(56)}
$$
$$
N_0 = \rho^{-1} N_1 \rho = -(\lambda_0 \sigma^{-1} \rho + 1) \leqno{(57)}
$$
de sorte que $N_0^2 = 0$, et
$$
\varepsilon_0 : \lambda \mapsto \varepsilon_0(\lambda) \defeq 1 + \lambda N_0 \leqno{(58)}
$$
est un sous-groupe à 1 paramètre de $\operatorname{Gl}(M)$, soit $L_0$, dont le normalisateur dans $\operatorname{Gl}(M)$ sera noté $\mathcal{B}_0$. On a deux homomorphismes
$$
\mathcal{B}_0 \xrightarrow{\chi(\underline{u}) = \det(\underline{u})} \mathbb{G}_m \leqno{(59)}
$$
$\chi'(\underline{u})$ donné par $\underline{u}(\varepsilon_0(\lambda)) = \varepsilon_0(\chi'(\underline{u})\lambda)$\marginpar{je souligne $\underline{u}$, pour ne pas le confondre avec des coefficients (comme $\gamma = \dots$) !}

(en termes linéaires dans $M$ :
$$
\underline{u} N_0 \underline{u}^{-1} = \chi'(\underline{u}) N_0 ) \leqno{(60)}
$$
et on considère le sous-groupe
$$
\mathcal{B}_0^* \subset \mathcal{B}_0 \text{ défini par } \chi(\underline{u}) = \chi'(\underline{u})
$$
$$
\text{i.e. par } \underline{u} N_0 \underline{u}^{-1} = \det(\underline{u}) N_0 \leqno{(61)}
$$

\footnote{Texte biffé en marge : il faudrait vérifier que $L_0 \subset \mathcal{B}_0^*$ c.-à-d. si pour $\underline{u} \in L_0$, $\dots$ c'est vrai, cf plus bas.}on a $L_0 \subset \mathcal{B}_0^*$ et une suite exacte de sch. en gr.
$$
1 \to L_0 \to \mathcal{B}_0^* \xrightarrow{\det} \mathbb{G}_m \to 1.
$$

\newpage

%% ===== Page 276 =====

Soit $\underline{u}$ une section de $\tilde{B}_0^*$, j'ai envie de prouver que si on pose
$$
\xi = [\underline{u}, \rho] \ , \quad \eta = [\underline{u}, \sigma] \leqno{(63)}
$$
on aura automatiquement la relation
$$
\xi = \eta \, \varepsilon_1(\mu-1) \qquad \text{où } \mu = \det \underline{u} . \leqno{(64)}
$$

Pour nous rassurer, prenons d'abord $\underline{u}$ dans $L_0$, i.e.
$$
\underline{u} = \varepsilon_0(\lambda)
$$
d'où
$$
[\underline{u}, \rho] = \varepsilon_0(\lambda) \rho \varepsilon_0(-\lambda) \rho^{-1} = \varepsilon_0(\lambda) \varepsilon_1(-\lambda)
$$
$$
[\underline{u}, \sigma] = \varepsilon_0(\lambda) \sigma \varepsilon_0(-\lambda) \sigma^{-1}
$$

Or
$$
\sigma \varepsilon_0(-\lambda) \sigma^{-1} = 1 + (\sigma N_0 \sigma^{-1}) \lambda 
$$
or en vertu de 57) on a
$$
\sigma N_0 \sigma^{-1} = -(\lambda_0 \rho \sigma^{-1} + 1) = N_1 \qquad (49)
$$
donc on a
$$
\begin{cases}
\sigma \varepsilon_0(\lambda) \sigma^{-1} = \varepsilon_1(\lambda) \\
\rho \varepsilon_0(\lambda) \rho^{-1} = \varepsilon_1(\lambda)
\end{cases} \leqno{(65)}
$$
[unclear: et par suite]

donc on trouve bien
$$
[\underline{u}, \rho] = [\underline{u}, \sigma] = \varepsilon_0(\lambda) \varepsilon_1(-\lambda)
$$
i.e. (64), puisque ici $\mu = \det \underline{u} = 1$, donc $\mu - 1 = 0$.

\newpage

%% ===== Page 277 =====

Nous avons vu que $\rho$ et $\sigma$ tous les deux transforment $N_0$ dans $N_1$, donc le sous-groupe à 1 paramètre additif
$$ L_0 = \{ \underset{= 1 + \lambda N_0}{\varepsilon_0(\lambda)} \mid \lambda \text{ dans } G_a \} $$
dans le sous-groupe analogue
$$ L_1 = \{ \underset{= 1 + \lambda N_1}{\varepsilon_1(\lambda)} \mid \lambda \text{ dans } G_a \} \ , $$
et donc $\widetilde{B}_0^\ast$ dans $\widetilde{B}_1^\ast$ (défini en termes de $L_1$, i.e.\ de $N_1$, comme $\widetilde{B}_0^\ast$ en termes de $L_0$). La formule (64) s'écrit alors
$$ \rho \underline{u}^{-1} \rho^{-1} = \sigma \underline{u}^{-1} \sigma^{-1} \varepsilon_1(\mu - 1) $$
ou plutôt, en remplaçant $\underline{u}$ par $\underline{u}^{-1}$ :
$$ \rho \underline{u} \rho^{-1} = \sigma \underline{u} \sigma^{-1} \ \varepsilon_1(\frac{1}{\mu} - 1) \leqno (65) $$
or
$$ \rho \underline{u} \rho^{-1} = \operatorname{int}(\lambda)(\sigma \underline{u} \sigma^{-1}) $$
où $\lambda = \rho \sigma^{-1}$, et la formule (65) s'écrit
$$ \operatorname{int}(\lambda)(v) = v \cdot \varepsilon_1(\frac{1}{\mu} - 1) \qquad \text{pour } v \text{ dans } \widetilde{B}_1^\ast \text{ et } \mu = \det v \ ; \leqno{(\ast\ast)} $$
et en vertu de (40 bis) $\operatorname{int} \lambda = \operatorname{int}(\varepsilon_1)$, donc la formule se réduit à
$$ \operatorname{int}(\varepsilon_1)(v) = v \cdot \varepsilon_1(\frac{1}{\mu} - 1) $$

\newpage

%% ===== Page 278 =====

$$
\varepsilon_1 v \varepsilon_1(-1) = v \varepsilon_1 \Big(\frac{1}{\mu}-1\Big) \quad \text{soit}
$$
$$
\underbrace{\varepsilon_1}_{1+N_1} v = v \underbrace{\varepsilon_1\Big(\frac{1}{\mu}\Big)}_{1+\frac{1}{\mu}N_1} \quad \text{soit} \quad v + N_1 v = v + \frac{1}{\mu} v N_1
$$
soit enfin
$$
v N_1 v^{-1} = \mu N_1 \quad (= \det(v). N_1)
$$
ce qui n'est autre que la définition de $\tilde{B}_1^\gamma$ !

On trouve donc un homomorphisme
$$
\begin{cases}
\tilde{B}_1^\gamma \longrightarrow \mathcal{F}_{\rho,\sigma} \\
\underline{u} \longmapsto ([\underline{u},\rho], [\underline{u},\sigma], \det \underline{u})
\end{cases}
\leqno (66)
$$
de $\tilde{B}_1^\gamma$ dans l'ensemble des solutions du système (33) à (35) , est-ce un isomorphisme ?

Notons que l'égalité
$$
[\underline{u},\rho] = [\underline{u}',\rho] \quad \text{signifie que} \quad \underline{u}' \underline{u}^{-1} \in T_\rho
$$
signifie simplement que
$$
T_\rho \cap \tilde{B}_1^\gamma = \{1\}
$$
en fait l'analyse faite tantôt montre

\newpage

%% ===== Page 279 =====

que cette intersection est égale à \dots
: $T_0 \cap \tilde{B}^?_0$, donc contenue dans
$T_0 \cap T_{\rho}$. Je vais supposer
$$
T_0 \cap T_{\rho} = \mathbb{G}_m \quad (= \operatorname{Centr}_{\mathfrak{S}}(\mathcal{M})) \leqno{(68)}
$$
-- il suffira pour ceci de supposer qu'en
degrés $\ge 1$, $s, \sigma, \rho$ sont linéairement indé-
pendants -- c'est le cas bien sûr dans le
cas qui nous intéresse. Donc
$$
T_0 \cap T_{\rho} \cap \tilde{B}^?_0 \subset \mathbb{G}_m \cap \tilde{B}^?_0,
$$
mais
$$
\mathbb{G}_m \cap \tilde{B}^?_0 = \mu_2
$$
car si $\lambda \in \mathbb{G}_m$, alors
$$
\lambda N_0 \lambda^{-1} = N_0 = 1 \cdot N_0
$$
mais $\det \lambda = \lambda^2$, et la condition que $\lambda$
soit dans $\tilde{B}^?_0$, c'est bien $\lambda^2 = 1$.

On trouve donc ainsi
$$
T_0 \cap \tilde{B}^?_0 = T_{\rho} \cap \tilde{B}^?_0 = \mu_2 \leqno{(69)}
$$
et l'homomorphisme (66) n'est pas un
monomorphisme, mais se factorise

\newpage

%% ===== Page 280 =====

en un monomorphisme
$$
\widetilde{B}_0 / \mu \hookrightarrow \widetilde{\mathcal{I}}_{p,\sigma} \leqno{(70)}
$$
À vrai dire, on voit que $\widetilde{B}_0$ n'est pas le "bon" groupe à introduire ici - dans le cas particulier $\varepsilon_0 = \begin{pmatrix} -1 & 0 \\ 0 & 1 \end{pmatrix}$, on a :
$$
B_0 = \left\{ \begin{pmatrix} p & \ell \\ 0 & q \end{pmatrix} \mid pq \text{ inv.} \right\}
$$
et la relation $\det(\underline{u}) = \chi(\underline{u})$ s'écrit, par un calcul immédiat
$$
p/q = pq \quad \text{i.e.} \quad q^2 = 1 \leqno{(71)}
$$
tandis que le "bon" groupe qui nous intéresse est défini par
$$
B'_0 = \left\{ \begin{pmatrix} p & \ell \\ 0 & 1 \end{pmatrix} \mid p \text{ inv.} \right\} ,
$$
donc avec $q=1$. Pour donner la description dans le cas général, on note que la fonction déterminant sur $B_0^*$ se décompose intrinsèquement en deux facteurs
$$
\det \underline{u} = p(\underline{u}) q(\underline{u}) , \leqno{(72)}
$$
qu'on peut distinguer l'un de l'autre par la condition que
$$
\underline{u} N_0 \underline{u}^{-1} = \frac{p(\underline{u})}{q(\underline{u})} N_0 \leqno{(73)}
$$
et la condition qui va nous donner le "bon" groupe n'est pas (71), mais bien
$$
q = q(\underline{u}) = 1 . \leqno{(74)}
$$

\newpage

%% ===== Page 281 =====

qui définit donc un groupe qui est libre
(sauf en caractéristique 2 !) et s'insère dans une suite exacte canonique
$$
1 \to \mu_2 \to B'_0 \xrightarrow{\det} \mathbb{G}_m \to 1 \leqno{(74)}
$$
On a trivialement
$$
B'_0 \cap \mathbb{G}_m = \{ 1 \} \leqno{(75)}
$$
ce qui implique que le morphisme
induit par (66)
$$
B'_0 \to \mathfrak{X}_{p,\sigma} \leqno{(76)}
$$
est un monomorphisme.

Pour montrer que c'est un isomorphisme, il suffit
maintenant, en vertu des principes généraux
de montrer que c'est bijectif sur les
points à valeurs dans un corps algébriquement clos de caractéristique $0$
(pour nous simplifier la vie). Je laisse
au lecteur le détail des vérifications.

On [unclear: prendra], pour la condition
$$
\text{que l'on a} \quad \varepsilon_i = 1 + \tilde{N}_i \quad (\tilde{N}_i^2 = 0) \leqno{(40 \text{bis})}
$$
dans un [unclear: sous-groupe] de $\operatorname{Sl}(\mathcal{M})$, et de même
pour [unclear: $p_0$] que $\sigma$ s'en déduit par [unclear: limite].

\newpage

%% ===== Page 282 =====

formules
$$ \sigma = - \varepsilon_1^{-1} \rho \quad (= - (1-N_1)\rho \text{ i.e. } \sigma = -\rho + N_1 \rho \text{, i.e. } \rho+\sigma = N_1\rho ) \leqno{(77)} $$
(je laisse tomber la constante $\lambda_0$, ou plutôt\footnote{je prends la constante qui intervient dans la condition d'anomalie (40 bis) de façon que $\varepsilon_1 = \rho \sigma$.} la situation n'est pas changée quand on remplace $\rho, \sigma$ par des multiples scalaires, car $\Gamma_\rho$ et $\Gamma_\sigma$ n'est pas changé ...) La
condition (41) s'écrit alors
$$ c = \Tr\rho + \Tr\sigma = \Tr (1 - \varepsilon_1^{-1}) \rho = \Tr N_1 \rho \in \Gamma_{\mathcal{O}_S}^* $$
(section inversible de $\mathcal{O}_S$). On suppose
que $1, \rho, \sigma$ sont linéairement indépendants dans
les fibres, il revient au même de dire que
$1, \rho$ et $N_1 \rho$ le sont, ou encore que
$\rho^{-1}, 1$ et $N_1$ le sont, ou que $1, \sigma, N_1 \sigma$ le sont,
ou $\sigma^{-1}, 1, N_1$. En résumé -

\marginpar{[unclear: NB. $\rho + \sigma = N_1 \rho = \rho N_1 = - N_1 \sigma = - \sigma N_1$ \\ (78 bis) quitte à multiplier $\rho, \sigma$ par des scalaires on a $\rho\sigma=1$ et $\varepsilon_1 = \rho\sigma$ (ça n'introduit d'ailleurs aucune restriction dans la position du problème)]}

Proposition. Soit $M$ une alg. d'Azumaya sur
le schéma de base $S$, $N_1$ une section de
trace nulle dans $M$, i.e. $\varepsilon_1 = 1+N_1$ une section
unipotente, $\rho, \sigma$ des sections reliées par
$$ \varepsilon_1 = - \rho \sigma^{-1} \text{ i.e. } \sigma = - \varepsilon_1^{-1} \rho \text{ i.e. } \rho + \sigma = N_1\rho \text{ i.e. } \rho + \sigma = - N_1\sigma \leqno{(78)} $$
(chacune détermine donc chacune des trois
sections $\varepsilon_1, \rho, \sigma$ - et $\varepsilon_1$ ou $N_1$) - l'une des
deux autres, et par ex. de façon arbitraire
$\rho$ et $N_1$ (i.e. $\rho$ et $N_1$ de trace nulle) -

\newpage

%% ===== Page 283 =====

supposer en plus
\leqno{(2^\circ)} $1, \rho, \sigma$ linéairement indépendants en chaque fibre i.e. $1, \rho^{-1}, N_1$ lin. indép. en chaque fibre i.e. $1, \sigma^{-1}, N_1$ lin. indép. en chaque fibre\footnote{NB $2^\circ)$ est corollaire du $3^\circ)$, [unclear: sauf peut-être pour les s-schémas fermés de $S$ où il s'annule].}
et \leqno{(3^\circ)} $\Tr N_1 \rho$ inversible.
(Ceci posé, on aura $N_1 \rho = - N_1 \sigma = \rho + \sigma$).
$\operatorname{int}(\rho^{-1}) \varepsilon_1 = \operatorname{int}(\sigma^{-1}) \varepsilon_1$, soit $\varepsilon_0$
(dans la notation de (74) on a $\varepsilon_0 = 1 + N_0$)

Soit $B_0$ le [unclear: sous-schéma en groupes normalisateur de la droite] $L_0 = S \cdot N_0$
$B'_0 \subset B_0$ le sous-schéma en groupes
défini par (74), de sorte qu'on a
$\varepsilon_0 = 1 + N_0$
posons
$$ \varepsilon_0(\lambda) = 1 + \lambda N_0 , \quad \varepsilon_1(\lambda) = 1 + \lambda N_1 $$
(sous-groupes à 1-paramètre additifs (isomorphes à $\mathbb{G}_a$) dans $\mathbb{G}(M)$)
et soient $L_0 =$ sous-schéma en groupes $\varepsilon_0(\mathbb{G}_a)$
$L_1 =$ sous-schéma en groupes $\varepsilon_1(\mathbb{G}_a)$,

(donc $\rho(L_0) = \sigma(L_0) = L_1$, )

Soit $B'_0$ le sous-schéma en groupes de $B_0$ défini par (74), de sorte que $B_0$
s'insère dans une suite exacte
$$ 1 \to L_0 \hookrightarrow B_0 \xrightarrow{\det} \mathbb{G}_m \to 1 $$
et que $u \in \Gamma(B_0)$ implique
$\underline{u}(N_0) = \det(u) N_0$, i.e.
$\underline{u}(\varepsilon_0(\lambda)) = \varepsilon_0(\det(u) \lambda)$.

Ceci posé, considérons le morphisme
$$ B'_0 \to \mathbb{G}(M) \times \mathbb{G}(M) \times \mathbb{G}_m $$
$$ u \mapsto ([\underline{u}, \rho], [\underline{u}, \sigma], \det(u)) . $$

\newpage

%% ===== Page 284 =====

Celui-ci prend ses valeurs dans $\mathcal{X}_{\rho, \sigma}$, le schéma des solutions des équations simultanées (33) à (35), et induit un iso
$$
B'_0 \isom \mathcal{X}_{\rho, \sigma} \cdot \leqno{(79)}
$$

Rappelons pour mémoire

\noindent \textbf{Corollaire.} La projection $\mathcal{X}_{\rho} \to \mathcal{G}(\mathcal{M})$
$$
(\xi, \eta, \nu) \mapsto \xi
$$
est un mono, et l'image est formée des $\xi$ satisfaisant
$$
\det \xi = 1, \quad \operatorname{Tr} \xi \rho = \operatorname{Tr} \rho, \quad \operatorname{Tr} \xi (-N_{1\rho}) \text{ inv.} \marginpar{[unclear: -N_{1\rho} = N_0 \rho = N_0 - \sigma N_0]}
$$
De même la projection $\mathcal{X}_{\rho} \to \mathcal{G}(\mathcal{M})$
$$
(\xi, \eta, \nu) \mapsto \eta
$$
est un mono, et l'image est formée des $\eta$ satisfaisant
$$
\det \eta = 1, \quad \operatorname{Tr} \eta \sigma = \operatorname{Tr} \sigma, \quad \operatorname{Tr} \eta (\rho + \sigma) \text{ inv.}
$$

Enfin, on a
$$
\operatorname{Tr} \xi (\rho + \sigma) = \operatorname{Tr} \eta (\rho + \sigma) = 1/\nu \cdot \leqno{(80)}
$$

\noindent \textbf{Scholie.} Par transport de structure, on trouve une structure de groupe sur $\mathcal{X}_{\rho, \sigma}$.

\newpage

%% ===== Page 285 =====

En l'explicitant, on trouve
$$
\begin{array}{clcl}
\text{si} & \underline{u} & \text{ correspond à } & \xi, \eta, \mu \\
& \underline{u}' & \text{---} & \xi', \eta', \mu' \\
& \underline{u}'' = \underline{u}' \underline{u} & \text{---} & \xi'', \eta'', \mu''
\end{array}
$$
on aura
$$
\mu'' = \mu \mu', \quad \xi'' = u'(\xi) \xi', \quad \eta'' = u'(\eta) \eta' \leqno{(81)}
$$
où on pose, comme d'habitude $u'(x) = u' x u'^{-1}$, pour deux éléments $x,u'$ d'un groupe.

\subsection*{III) Étude de l'équation \hfill (Suite)}

$$
\alpha \rho(\alpha)^{-1} = \beta \sigma(\beta)^{-1} \varepsilon_1(\mu-1) \leqno{(82)}
$$
$$
(\alpha, \beta \in G(\mathcal{M}), \mu \in \Gamma_{\mathfrak{S}'})
$$
i.e.
$$
[\alpha, \rho] = [\beta, \sigma] \varepsilon_1(\mu-1) \leqno{(82bis)}
$$

Soit $G_{\rho, \sigma}$ le schéma des solutions, on a donc un morphisme
$$
G_{\rho, \sigma} \to \mathfrak{X}_{\rho, \sigma} \leqno{(83)}
$$
$$
(\alpha, \beta, \mu) \mapsto (\alpha \rho(\alpha)^{-1}, \beta \sigma(\beta)^{-1}, \mu)
$$
et il résulte du théorème précédent que c'est un épimorphisme, plus précisément

\newpage

%% ===== Page 286 =====

On a une section canonique

(84)
\begin{tikzcd}
\mathcal{G}_{\rho,\sigma} \ar[r] & \mathcal{X}_{\rho,\sigma} \\
& B'_0 \ar[u]
\end{tikzcd}
$$(\underline{u}, \underline{u}, \det \underline{u}) \longleftarrow \underline{u} \longmapsto ([\underline{u},\rho], [\underline{u},\sigma], \det \underline{u})$$

Cette section prend ses valeurs dans la « diagonale », [unclear: formée par ce qui] correspond aux couples $(\underline{u}, \mu)$ (avec $\underline{u} \in \Gamma \mathcal{M}^*, \mu \in \Gamma E'$) tels que

(85) $$[\underline{u}, \rho] = [\underline{u}, \sigma] \varepsilon_1(\mu-1) \quad ,$$

qu'on tâchera de déterminer ainsi.

Conjuguant les résultats de (I) et (II) on voit que la solution générale de (82) sera fournie par

$$(\underline{u}\underline{a}, \underline{u}\underline{b}, \mu = \det \underline{u})$$

où $\underline{u}$ est une section de $B'_0$, $\underline{a}$ une section locale de $T_\rho$, $\underline{b}$ une section de $T_\sigma$.

De façon précise

Th. Considérons

(86)
\begin{tikzcd}
B'_0 \times T_\rho \times T_\sigma \ar[r] & G \times G \times E' \\
(\underline{u}, \underline{a}, \underline{b}) \ar[r, maps to] & (\underline{u}\underline{a}^{-1}, \underline{u}\underline{b}^{-1}, \det \underline{u})
\end{tikzcd}

\newpage

%% ===== Page 287 =====

prend ses valeurs dans $\mathcal{G}_{p, \sigma}$ ,
et induit un
isomorphisme
$$ B'_0 \times T_p \times T_\sigma \isommap \mathcal{G}_{p, \sigma} \leqno{(87)} $$
De façon plus précise, soit
$(\alpha, \beta, \mu)$ une section de $\mathcal{G}_{p, \sigma}$ , alors
il existe des uniques sections $\underline{u}$ de $B'_0$,
$\underline{a}$ de $T_p$, $\underline{b}$ de $T_\sigma$, telles que l'on ait
$$ \alpha \underline{a}, \beta \underline{b} \in B'_0 \leqno{(88)} $$
et on aura alors nécessairement
$$ \alpha \underline{a} = \beta \underline{b} \ \left( \defeq \underline{u} \right) $$
d'où l'on aura
$$ ([\alpha, p], [\beta, \sigma]) = ([\underline{u}, p], [\underline{u}, \sigma]) $$
et l'unique $\mu$ qui satisfasse à (82)
est donc $\det \underline{u}$ , d'où
$$ (\alpha, \beta, \mu) = (\underline{u}\underline{a}^{-1}, \underline{u}\underline{b}^{-1}, \det \underline{u}) \quad (\text{de façon unique}) \leqno{(89)} $$
Dém. Il reste seulement à expliciter
l'unicité de $\underline{a}, \underline{b}$ satisfaisant à
(88), ce qui est conséquence de
$$ T_p \cap B'_0 = T_\sigma \cap B'_0 = \{1\} \leqno{(90)} $$
puis [unclear: on achève] d'expliciter dans (II) ,

\newpage

%% ===== Page 288 =====

En termes de la donnée de $\alpha$, on aura
donc un $\underline{a}$ uniquement déterminé par
la condition $\alpha \underline{a} \in \Gamma B'_0$, et on aura

$$
\underline{a} = c_0 \rho + d_0 \qquad c_0, d_0 \in \Gamma(\mathcal{O}_S) \text{ déterminés par } \alpha
\leqno{(91)}
$$

de même on aura

$$
\underline{b} = t_0 \sigma + u_0 \qquad t_0, u_0 \in \Gamma(\mathcal{O}_S) \text{ déterminés par } \beta
\leqno{(92)}
$$

de telle façon donc qu'on aura

$$
\alpha \underline{a} = \beta \underline{b} \quad (\overset{\text{déf}}{=} \underline{u})
\leqno{(93)}
$$

Dans le cas numérique
$$
\rho = \begin{pmatrix} 0 & -1 \\ 1 & 1 \end{pmatrix}, \quad \sigma = \begin{pmatrix} 0 & -1 \\ 1 & 0 \end{pmatrix},
$$
posant
$$
\alpha = \begin{pmatrix} a & b \\ c & d \end{pmatrix}, \quad \beta = \begin{pmatrix} r & s \\ t & u \end{pmatrix}
$$
on aura (cf p.\ 398' (75)(78)) sauf erreur\footnote{$\delta = \det \alpha = ad-bc$, $\delta' = \det \beta = ru-st$. Il faudra préciser ça dans le cas général.}

$$
\underline{a} = \frac{1}{\delta D} (c \rho + d) = \frac{1}{c^2+cd+d^2} (c\rho+d)
\leqno{(94)}
$$
où $\delta D = c^2+cd+d^2$, $\delta = \det \alpha$, on aura besoin de $\xi = \alpha \rho \alpha^{-1} = \begin{pmatrix} A & B \\ C & D \end{pmatrix}$

$$
\underline{b} = \frac{1}{\delta' U} (t \sigma + u) = \frac{1}{t^2+u^2} (t\sigma+u)
\leqno{(95)}
$$
où $\delta' U = t^2+u^2$, on aura besoin de $\eta = \beta \sigma \beta^{-1} = \begin{pmatrix} R & S \\ T & U \end{pmatrix}$

Le seul ennui, pour faire le lien avec
les calculs ``bourrins'' de loc.\ cit., est dans
le facteur [unclear: $\tau_\infty$] de p.\ 398' (78), où
$\alpha = \Phi_2(\mu)$, donc là on a $\det \alpha = 1 \dots$

\newpage

%% ===== Page 289 =====

à première vue il semble y avoir incom-
patibilité entre la théorie (plus simple)
développée ici, et ce facteur $\tau_\infty^2$ dont je
suis persuadé qu'il a de très bonnes
raisons d'être là ! Et [unclear: d'ailleurs] dans loc.
cit. on écrivait l'équation des cartes avec
la probabilité d'un signe\footnote{vu que $\xi, \eta$ y sont définis comme $\alpha\rho(\alpha)^{-1}$, $\beta\sigma(\beta)^{-1}$}
$$ \xi = \varepsilon \eta h_1 \quad \varepsilon \in \{ \pm 1 \} $$
($h_1 = \varepsilon_1(\mu_{-1})$ avec les
notations actuelles)

c'est donc une généralisation p. r. à l'équation
pour [unclear: le début], où l'on a supposé $\varepsilon = +1$.
Il faudra donc tantôt reprendre la
théorie précédente, dans le contexte plus
général d'une équation
$$ \underset{\alpha\rho(\alpha)^{-1}}{\overset{\shortparallel}{\xi}} = \varepsilon . \underset{\beta\sigma(\beta)^{-1}}{\overset{\shortparallel}{\eta}} \varepsilon_1(\mu_{-1}) \quad \varepsilon \text{ [unclear: central dans] } \mathfrak{G} \leqno{(95)} $$
où l'on devra avoir (en prenant les
[unclear: dét.] des deux membres) 
$$ \varepsilon^2 = 1 $$
- mais il n'y a pas lieu sans doute de
supposer qu'on ait $\varepsilon \in \{\pm 1\}$. Mais je vais
continuer d'abord les calculs dans le cas
présent $\varepsilon = +1$.

\newpage

%% ===== Page 290 =====

\section*{IV \quad Un épinglage canonique}

Rappelons que la donnée\footnote{sur un schéma de base $S$} d'une forme $G$ de $\operatorname{Gl}(2)_S$, ou $SG$ de $\operatorname{Sl}(2)_S$ (ou encore d'un $PG$ de $P\operatorname{Gl}(2)_S$), d'une algèbre d'Azumaya $\mathcal{M}$ (forme de $M_2(\mathcal{O}_S)$), revient à la donnée d'un fibré en droites projectives $\mathcal{P}$ sur $S$. $\mathcal{P}$ peut s'interpréter comme le schéma des sous-groupes de Borel $B$ de $G$, ou $SB$ de $SG$, ou $PB$ de $PG$, ou comme le schéma des sous-groupes à 1-paramètre ``boréliens'' (i.e.\ identiques à la partie unipotente d'un Borel convenable de $G$, ou de $SG$, ou $PG$ (NB. Les morphismes canoniques $SG \to G \to PG$ induisent des bijections entre sous-groupes de Borel des trois groupes en question, et aussi entre les uns et les autres des sous-groupes à 1-paramètre ``boréliens'' de ces groupes). Également, $\mathcal{P}$ peut s'interpréter en termes des sections $N$ de $\mathcal{M}$ qui sont non nulles sur chaque fibre et telles que
$$ 
\det N = 0, \Tr N = 0 \qquad \text{d'où } N^2 = 0 \quad \text{(les sections } N \text{ sont régulières, c-à-d.\ non nulles sur ch.\ pt.)}, \leqno{(1)} 
$$
(``strictement nilpotentes''), à toute telle section correspond un sous-groupe à 1-paramètre ``borélien'' de $SG$, [unclear: pour]
$$ 
\mathbf{G}_a \isom U, \quad \lambda \mapsto \exp(\lambda N) = 1 + \lambda N \leqno{(2)} 
$$
Mais il est évident que si $N$ et $N'$ sont de telles

\newpage

%% ===== Page 291 =====

sections qui sont homothétiques (i.e.\ $N' = a N$
avec $a \in \Gamma(S, \mathcal{O}_S)$), alors elles définissent le même
sous-groupe à 1 paramètre, -- donc ces directions
correspondent en fait à des sections du
fibré projectif (de dim relative 3) $\mathbf{P}(\mathcal{M})$, --
si on tient compte de $Tr N = 0$, soit à des
sections du fibré projectif (de dim relative
2 -- donc un fibré en plans projectifs)
$\mathbf{P}(\mathcal{M}_0)$ (où $\mathcal{M}_0 \subset \mathcal{M}$ est défini justement
par la condition $Tr N = 0$) -- à savoir les
sections de la quadrique conique "définie"
par l'équation
$$ \det N = 0 $$

Les $N$ correspondent d'ailleurs aussi aux
opérateurs
$$ e = e_N = 1 + N $$
satisfaisant les conditions
$$ \det e = 1, \quad Tr e = 2 \leqno (3) $$
ce qui équivaut en fait, en posant
$e = 1 + N$, i.e.\ en définissant $N$ par
$$ N = e - 1 $$
aux conditions (1) [car $Tr$ et $\det$ de deux
opérateurs $N$ et $e$ reliés par $e = \lambda + N$
s'obtiennent mutuellement par
$$
\begin{cases}
Tr e = 2 \lambda + Tr N, & \det e = \lambda^2 + \lambda Tr N + \det N \\
Tr N = - 2 \lambda + Tr e, & \det N = \lambda^2 - \lambda Tr e + \det e
\end{cases}
$$
]

\newpage

%% ===== Page 292 =====

Les opérateurs sections de $G$ (en fait de $S G$)
satisfaisant à (3), (et qui avec $e=1$ engendrent un $P(M)$)
\marginpar{terminologie un peu lourdingue}
on les appellera les sections \underline{strictes}
unipotentes régulières, deux telles sections $e, e'$
sont considérées "équivalentes" si ...
$$
e' = e^{\lambda} \defeq (1 + \lambda N) \quad (\lambda \in \Gamma_{\mathbf{Q}} = \widehat{\mathbf{Z}}^\times) \leqno{(4)}
$$
$$
= 1 + \lambda(e-1) = (1-\lambda) + \lambda e
$$

On voit donc que les sections de $\mathcal{M}$
(strictement) unipotentes régulières correspondent
1-1
aux sections nilpotentes régulières,
et la relation d'équivalence d'homothétie
devient celle d'équivalence par $e \mapsto e^{\lambda}$ --
et les classes d'équivalence des unes
aux autres correspondent aux sous-groupes
additifs boréliens de $G$ (ou $S G$, ou $P G$
ou $S G \to P G$) -- ou encore (en
prenant les normalisateurs de ceux-ci dans
$G, S G \text{ ou } P G$) aux Borels de ces groupes,
i.e. aux sections de $\mathbb{P}$ sur $S$.

Supposons donnée une telle section,
soit $\mathcal{P}$, elle permet donc de [unclear: munir canoniquement]
$P$ de un fibré [unclear: affine]
$$
E (= \mathcal{P} \setminus \{ 1 \}) \leqno{(5)}
$$
, cette [unclear: structure vectorielle] autorise des translations...

\newpage

%% ===== Page 293 =====

et $x_\infty$ est donné [unclear: rigidement] par un sous-
groupe à $1$-paramètre borélien
$$L \subset SG \subset G \leqno{(6)}$$
(puis, par $SG \to PG$, donnera naissance à un
ss-groupe additif borélien $L \subset PG$ de
$PG$...). Localement, $L$ sera engendré par
une section nilpotente régulière
$$e = 1+N \qquad \text{($N$ nilpotente régulière)} \leqno{(7)}$$
formant le ss-groupe [unclear: à un paramètre]
$$e^\lambda = e(\lambda) = e_N(\lambda) = 1 + \lambda N \leqno{(8)}$$
mais comme on a dit, $e$, $N$ ne sont pas
déterminés intrinsèquement en termes du
$L \subset SG$. De façon précise, $L$ est munie
canoniquement d'une structure de fibré
vectoriel (respectée par ses normalisateurs
dans $G$), à savoir justement celle
déduite par transport de structure de
celle de $G_a$ par
$$e : G_a \isommap L \qquad e(\lambda) = e^\lambda = 1 + \lambda N$$
et le groupe $L$ s'identifie ainsi (fibré
vectoriel canoniquement associé au fibré en
droites $\mathcal{N} \subset \mathcal{M}$ sous-tendu par les $N$)
$$L = \check{V}(\mathcal{N}) \leqno{(9)}$$

\newpage

%% ===== Page 294 =====

le choix de $e$, i.e.\ $N$, correspond simple-
ment à celui d'une base $N_0$ de $\mathcal{N}$, ou
d'une section de $L$ partout $\neq 1$ (i.e.\ d'une sect.\ $e$ de $L$\footnote{N.B.\ un tel $L$ est aussi ``fibré principal'' sous un fibré vectoriel des translations, i.e.\ un fibré en droites affines $\mathbb{P}(\mathcal{N})$.}
qui soit $\neq 1$ sur les fibres).

Revenons à l'équation (3), et notons
que $\rho$ et $\sigma$ n'y interviennent que modulo
$(\lambda, \lambda' \in \Gamma_{\max})$,
les transformations $\rho \mapsto \lambda \rho$, $\sigma \mapsto \lambda' \sigma$ (qui ne
changent pas les commutateurs
$[\alpha, \rho]$ et $[\beta, \sigma]$) ; d'autre part, on peut l'exprimer en
disant que $\rho, \sigma$ n'interviennent pas
[unclear: par leur composante]
dans les sous-groupes additifs [unclear: booléens]
de $G$ mais par leurs images dans $\Gamma(PG)$, d'
ailleurs :
$$
PG \simeq \operatorname{Aut}_{S\text{-gr}}(G) \simeq \operatorname{Aut}_{S\text{-gr}}(SG) \simeq \operatorname{Aut}_{\text{s-gr}}(P\mathcal{G}) \leqno{(4)}
$$
$$
\simeq \operatorname{Aut}_{S\text{-Alg}}(\mathcal{M}) \simeq \operatorname{Aut}_S(P)
$$
et cette remarque revient à dire que
$\rho$ et $\sigma$ n'interviennent que par leur action
sur $G$ (ou sur $SG$), ce qui est clair
aussi sur les formules
$$
[\alpha, \rho] = \alpha \rho(\alpha^{-1}), \quad [\beta, \sigma] = \beta \sigma(\beta^{-1}) .
$$

D'autre part, le facteur $\varepsilon_1(H-1)$ de (32)
peut être considéré comme une section
d'un sous-fibré en s-groupes additifs $L_1$ de
$G$ (en fait, de $SG$), qui s'exprime bien

\newpage

%% ===== Page 295 =====

entendu abélien --- l'équation (32)
s'écrit donc sous la forme (où je n'écris
même plus $\sim$, seulement $\alpha$ et $\rho$)
$$
[\beta, \sigma]^{-1} [\alpha, \rho] \in \Gamma L_1 \leqno{(11)}
$$
ou
$$
[\alpha, \rho] = [\beta, \sigma] \cdot \gamma \qquad \text{avec} \quad \gamma \in \Gamma L_1 . \leqno{(12)}
$$

Sous la forme (12), on peut la regarder
comme une équation reliant $\alpha, \beta, \gamma$.

Considérons l'image
$$
\widetilde{L}_1 \subset P G \leqno{(13)}
$$
de $L_1$ par $G \to P G$, la relation (40 bis)
s'écrit alors sous la forme
$$
\widetilde{\varepsilon}_1 = \tilde{\rho} \tilde{\sigma}^{-1} \in \Gamma \widetilde{L}_1 \quad (\equiv \Gamma L_1) . \leqno{(14)}
$$

Cette relation provient donc d'une section
bien déterminée exact au-dessus de $\tilde{\rho}, \tilde{\sigma}$
$$
\varepsilon_1 \in \Gamma L_1 \leqno{(15)}
$$
décrite en termes de $\rho, \sigma$ d'ailleurs
par la constatation qu'on ait une relation de
la forme\footnote{NB $\lambda_0$ est uniquement déterminé par $\rho, \sigma$ si $2$ est inversible dans $\hat{\mathbf{Z}}$ (car l'équation à résoudre est de la forme $\dots$). Mais pour l'application qui nous intéresse le plus $\dots$}
$$
\varepsilon_1 = - \lambda_0 \rho \sigma^{-1} \qquad (\lambda_0 \in \Gamma L_1 ) \leqno{(16) = (40 bis)}
$$
par commodité (à cause de l'application au cas
abélien qui nous intéresse le plus $\dots$) .

\newpage

%% ===== Page 296 =====

Revenons sur la situation générale d'une forme $G$ de $\P\operatorname{Gl}(2)_S$, et supposons qu'elle puisse s'écrire sous la forme
$$ G = \operatorname{Aut}_{\mathcal{O}_S\text{-Mod}}(F) \leqno{(17)} $$
où $F$ est un $\mathcal{O}_S$-Module localement libre de rang $2$,
C'est le cas dès qu'on a une section $x$ de $\mathcal{P}$ sur $S$, i.e.\ un sous-groupe [unclear: fermé] borélien $L$ de $G$ --- ou de $\SL_2$ (en fait on peut prendre alors pour $F$ l'extension de $\mathcal{O}_S$ par $\mathcal{N}$ qui détermine le fibré en droites affine $\mathcal{P} \setminus \{x\}$).
Mais le $\mathcal{P}$ se réalise
$$ \mathcal{P} = \mathbf{P}(F) \simeq \check{\mathbf{P}}(\check{F}) \leqno{(18)} $$
le fibré projectif associé à $F$, et les sections de $\mathcal{P}$ s'identifient aux ``droites'' $D \subset F$.

De même on a
$$ \mathcal{M} = \operatorname{End}_{\mathcal{O}_S\text{-Mod}}(F) \simeq \check{F} \otimes_{\mathcal{O}_S} F \leqno{(19)} $$
et le sous-Module $\mathcal{N}$ associé à la section $x$ de $\mathcal{P}$ n'est autre que
$$ \mathcal{N} = D^\circ \otimes D \quad (D^\circ = (F/D)^{\check{}} \subset \check{F}, \text{ l'orthogonal de } D \text{ dans } \check{F}) \leqno{(20)} $$
ses sections sont les $u: F \to F$ telles que
$$ u(F) \subset D, \quad u(D) = \{0\} \leqno{(21)} $$

\newpage

%% ===== Page 297 =====

\begin{flushleft}
Proposition. Soit $G$ forme de $\operatorname{Gl}(2)$, $\rho, \sigma$ deux sections
de $G$ telles que $\rho \sigma^{-1} = \varepsilon_1$ soit sect. unipot. régulière de $G$, (de droite
inv. $L_1$). Alors les conditions suivantes sont équivalentes : \marginpar{NB On a désigné par $x_1$ la section de $P$ associée à $L_1$}

(i) $\varepsilon_1$ est section régulière (i.e. $\neq 1$ dans chaque fibre)
i.e. $\varepsilon_1 = 1 + N_1$, avec $N_1$ une base de $\mathcal{N}_1$) et
$\tilde{\rho}^{-1}(x_1)$ est une section de $P$ partout $\neq x_1$ dans
chaque fibre. \marginpar{(i bis) Idem, avec $\tilde{\rho}$ au lieu de $\tilde{\rho}^{-1}$}

(ii) Idem, avec $\tilde{\sigma}^{-1}$ au lieu de $\tilde{\rho}^{-1}$.

(iii) On a (42), i.e. $c = \lambda_0 \Tr \rho + \Tr \sigma$ est inversible.

Dém. On a, grâce au fait que $\rho \sigma^{-1}$ [unclear: opère sur] $L_1$,
normalise $L_1$, (la relation
$\tilde{\rho}^{-1}(L_1) = \tilde{\sigma}^{-1}(L_1)$
qui revient au même), on a
$$ \tilde{\rho}^{-1}(x_1) = \tilde{\sigma}^{-1}(x_1) \leqno{(22)} $$
ce qui prouve l'équivalence de (i) et (ii).

Montrons que (i) $\Leftrightarrow$ (iii). OPS que $S$ est le
spectre d'un [unclear: anneau local], $L_1$ sa
[unclear: direction], et OPS $\lambda_0 = 1$ (quitte à remplacer
$\rho$ par $\lambda_0 \rho$). Si on avait $\varepsilon_1 = 1$ i.e.
$\rho = \sigma$, on aurait [unclear: Faux]
$\Tr(\rho + \sigma) = \Tr \rho + \Tr \sigma = 2c$, or (iii) implique la
régularité de $\varepsilon_1$, i.e. de $N_1$, OPS donc que
$\varepsilon_1$ est régulier. Il faut montrer que
(si $\varepsilon_1 = 1+N_1$ est unipot. régulier $\varepsilon_1$)
[unclear: la direction $D$ dans $F$, associé à $D \subset F$,]
\end{flushleft}

\newpage

%% ===== Page 298 =====

et une involution $\sigma$ de $F$, d'où
$\rho = - \varepsilon_1 \sigma$, et donc
$c = \Tr \sigma - \sigma + \Tr \rho = \Tr \sigma - \Tr \varepsilon_1 \sigma - \Tr N_1 \sigma$ est inversible ssi $D + \sigma \cdot D = F$
i.e.\ $\sigma D + D = F$. Prenons une base
$e_0, e_1$ de $F$ telle que $e_1$ engendre $D$, et
que $N_1$ s'écrive par la matrice $\left(\begin{matrix} 0 & 0 \\ 1 & 0 \end{matrix}\right)$,
soit $\sigma = \left(\begin{matrix} a & b \\ c & d \end{matrix}\right)$, on a $N_1 \sigma = \left(\begin{matrix} 0 & 0 \\ a & b \end{matrix}\right)$,
$\Tr N_1 \sigma = b$, $\sigma e_1 = a e_0 + b e_1$, et $D + \sigma \cdot D = D + b e_0$
est égale à $F = \mathbf{Q} s^2$ ssi $b$ inv., cqfd.

Supposons vérifiée ces conditions
équivalentes, pour $\tilde{\rho}$ qui s'exprime aussi
en disant que $L_1$ et
$$
L_0 = \tilde{\rho}^{-1}(L_{-1}) = \tilde{\sigma}^{-1}(L_{-1}) \leqno{(33)}
$$
sont disjoints sur chaque fibre, i.e.\ définissent
des Borels $B_1, B_0 \subset G$ tels que
$$
T = B_1 \cap B_0 \leqno{(34)}
$$
soit un tore (maximal) de $G$,
ce qui revient à dire l'intersection
des stabilisateurs $B_0, B_1$ dans $G$ des
sections $X_0, X_1$ de $T$ définis correspondant à $L_0, L_1$
(et d'ailleurs $X_1 = \tilde{\rho} X_0 = \tilde{\sigma} X_0 \leqno{(25)}$),
et le rang est dans le sous-groupe (fourni par $\sigma$)

\newpage

%% ===== Page 299 =====

dont $T$ dans $PGL_2 = \operatorname{Aut}_{\overline{\mathbf{Q}}}(\mathbb{P}^1)$. On peut s'en
assurer par le [unclear: choix d'une section]
[unclear: rationnelle] $x_0, x_1$ de $\mathbb{P}^1$ [unclear: projectifs]
...
...
...
\marginpar{le couple $(B_1, T)$ s'appelle une [unclear: paire de Borel]}
$$
B_1 = T \cdot L_1 \quad \text{ (produit semi-direct)} \leqno (26)
$$
qui [unclear: prend] la [unclear: suite] ... $x_0$ ...
...
[unclear: Borel] $B_0$ tel que
$$
T = B_0 \cap B_1 \ . \quad (\text{Les Borels } B_0 \text{ tels que}
$$
$B_0 \cap B_1 = T$, ou plus généralement $B_0 \cdot B_1$ soit [unclear: dense]
$B_0 \neq B_1$, sur chaque fibre, [unclear: correspondent aux] ...
[unclear: de l'espace] ... $T \to B_1$,
[unclear: formés] ... $T = B_0 \cap B_1$ .)

Cela nous amène à : la [unclear: donnée] ...
[unclear: de] $B_1$ et une [unclear: section] ...
... $U_1 = \mathbb{P}^1 \setminus \{x_1\}$, ...
...
...
... $L_1$ ...
$$
U_1 = \mathbb{P}^1 \setminus \{x_1\} \simeq L_1 \leqno (27)
$$
...

\newpage

%% ===== Page 300 =====

de réduire $\mathcal{P}$ à un fibré projectif (type $\mathbf{P}^1_S$)
il faut encore se donner une base
de $L_1$ i.e.\ une section $x_{01}$ de
$$ L_1^* = U_1 \setminus \{x_0\} \simeq \mathcal{P} \setminus \{x_1, x_0\} . \leqno{(28)} $$
Or la section $\varepsilon_1$ de $L_1$ est justement
une section de $L_1^*$ --- on a donc un
épinglage complet de la situation. En
termes de $(F, D=D_1)$ correspondant à $L_1$,

$$ \begin{cases} p^{-1}(D_1) = \sigma^{-1}(D_1) \stackrel{\text{déf}}{=} D_0 \\ F = D_0 \oplus D_1 \end{cases} \leqno{(29)} $$
$$ N_1 | D_0 : D_0 \xrightarrow{\sim} D_1 \leqno{(30)} $$
et l'épinglage : (tensoriant par $D_1^{\otimes -1}$) on a
$D_1 \simeq \mathcal{O}_S$, d'où
$$ \begin{cases} F \simeq \mathcal{O}_S^2 = \mathcal{O}_S e_0 + \mathcal{O}_S e_1 \\ D_0 = \mathcal{O}_S e_0, \quad D_1 = \mathcal{O}_S e_1 \\ N_1 e_1 = 0, \quad N_1 e_0 = e_1 \quad \text{i.e. } N_1 = \begin{pmatrix} 0 & 1 \\ 0 & 0 \end{pmatrix} \end{cases} \leqno{(31)} $$

\newpage

%% ===== Page 301 =====

i.e.
$$ 
N_1 = \begin{pmatrix} 0 & 0 \\ 1 & 0 \end{pmatrix} \quad , \quad \varepsilon_1 = \begin{pmatrix} 1 & 1 \\ 0 & 1 \end{pmatrix} \ . \leqno (32) 
$$

D'autre part nous choisissons $\rho, \sigma$ relevées dans $\operatorname{Gl}(2)$ au dessus de $\bar{\rho}, \bar{\sigma}$ (sections de $P\operatorname{Gl}(2) \simeq \operatorname{Gl}(2)/\mathbb{G}_m$) satisfaisant $\bar{\rho} \bar{\sigma}^{-1} = \bar{\varepsilon}_1$, et on supposera $\rho, \sigma$ choisis de telle façon qu'on ait $\rho \sigma^{-1} | D_1 = -1$, ce qui compte tenu que $\overline{\rho \sigma^{-1}} = \bar{\varepsilon}_1$ implique que $\rho \sigma^{-1} | (F/D_1) = -1$ i.e. de façon à avoir
$$ 
\rho \sigma^{-1} = - \varepsilon_1 \leqno (33) 
$$
ce qui détermine chacune des deux sections $\rho, \sigma$ en fonction de l'autre (déterminées [unclear: modulo un scalaire, i.e. fonction de $\mathcal{O}_S^*$] près). Comme $\rho D_0 = D_1$,
$$ 
\rho = \begin{pmatrix} 0 & u \\ v & w \end{pmatrix} \quad , \quad \det \rho = -uv \text{ inv.} \leqno (*) 
$$
et on peut lever l'ambiguïté scalaire sur $\rho$ en imposant par ex.
$$ 
\rho = \begin{pmatrix} 0 & 1 \\ v & w \end{pmatrix} \quad , \quad v \in \mathbb{G}_m \ . \leqno (34) 
$$

\newpage

%% ===== Page 302 =====

Notons que (33) s'écrit $\rho = -(1+N_1)\sigma$ i.e.
$$
\rho + \sigma = -N_1 \sigma \leqno (35)
$$
multipliant par $N_1$ à gauche, donne alors
$$
N_1 \rho + N_1 \sigma = 0
$$
(puisque $N_1^2 = 0$), d'ailleurs
$$
N_1 \rho = \rho N_0, \quad N_1 \sigma = \sigma N_0
$$
(car $N_1 = \rho N_0 \rho^{-1} = \sigma N_0 \sigma^{-1}$), donc (35) s'écrit
$$
\rho + \sigma = -N_1 \sigma = + N_1 \rho = - \sigma N_0 = + \rho N_0 \defeq \mathcal{N} \leqno (36)
$$
d'ailleurs quand $\rho$ est mis sous la forme
générale $\rho = \begin{pmatrix} 0 & u \\ v & w \end{pmatrix}$ on trouve
$$
\mathcal{N} = \begin{pmatrix} 0 & 0 \\ 0 & u \end{pmatrix} , \quad N^2 = \begin{pmatrix} 0 & 0 \\ 0 & u^2 \end{pmatrix}
$$
donc la normalisation (34) $u=1$ signifie
$$
\rho + \sigma = \begin{pmatrix} 0 & 0 \\ 0 & 1 \end{pmatrix} \leqno (37)
$$
ou encore
$$
N^2 = \mathcal{N} \quad (\text{où } \mathcal{N} = \rho + \sigma) \leqno (37 \text{ bis})
$$
ce qui donne pour $\sigma$
\marginpar{[unclear: note marginale rayée : NB On voit aussitôt que les conditions envisagées sur $\rho, \sigma$ impliquent que l'on peut se ramener à $u=1$ i.e.]}
$$
\sigma = \begin{pmatrix} 0 & -1 \\ -v & 1-w \end{pmatrix}
$$

D'ailleurs, $N_0 = \rho^{-1} N_1 \rho$ sera de la forme $\begin{pmatrix} 0 & \alpha \\ 0 & 0 \end{pmatrix}$,
un calcul immédiat montre que
$\alpha$ est donné par la formule $\alpha = u/v$, i.e. (vu la
normalisation $u=1$)

\newpage

%% ===== Page 303 =====

$$ N_0 = \begin{pmatrix} 0 & 1/\mu \\ 0 & 0 \end{pmatrix}, \quad \varepsilon_0 = \begin{pmatrix} 1 & 1/\mu \\ 0 & 1 \end{pmatrix} \leqno{(39)} $$

Ces normalisations canoniques étant faites\footnote{NB On aura (39bis) $\sigma^{-1} \rho = -\varepsilon_0$ car $\sigma^{-1} \rho = \sigma^{-1} (\rho \sigma^{-1}) \sigma = \sigma^{-1} (-\varepsilon_1) \sigma = -\sigma^{-1} (\varepsilon_1) = -\varepsilon_0$.}, revenons à la proposition de la page 436 ; comme $\varepsilon_1$ a été choisi (à [unclear: part] $\pm$ ce qu'il y a de canonique) il n'y a plus d'inconvénient à laisser l'équation (32), (mise s.\ la forme où on a plus l'habitude (4) -- (11)) sous sa forme [unclear: paramétrisée], a priori on s'attendait à introduire le paramètre $\gamma$ ou
$$ [\alpha, \rho] = \rho [\beta, \rho] \varepsilon_1(\gamma) $$
mais il apparaît que c'est $\gamma + 1 = \mu$ qui a une signification particulièrement remarquable dans la théorie - d'où on gardera la forme (32), ou le facteur $\varepsilon_1(\mu-1)$.

L'expression (47) de $P(\xi)$, grâce à (37), s'écrit
$$ P(\xi) = D \quad \text{pour} \quad \xi = \begin{pmatrix} A & B \\ C & D \end{pmatrix} \leqno{(40)} $$

\newpage

%% ===== Page 304 =====

Donc la proposition 453 nous dit que le schéma $\mathcal{X}_{\rho, \sigma}$ des solutions simultanées $(\xi, \eta, \mu)$ de
$$
\det\xi = \det\eta = 1, \quad \xi\eta = \varepsilon_1(\mu-1), \quad \operatorname{Tr}(\xi\rho) = \operatorname{Tr}\rho, \quad \operatorname{Tr}\eta\sigma = \operatorname{Tr}\sigma \leqno{(41)}
$$
est isomorphe, via $(\xi, \eta, \mu) \mapsto \xi$, à celui des
$$
\xi = \begin{pmatrix} A & B \\ C & D \end{pmatrix}
$$
satisfaisant les conditions
$$
\underbrace{\det\xi = 1}_{AD-BC}, \quad \underbrace{\operatorname{Tr}(\xi\rho)}_{C+wD+vB} = \underbrace{\operatorname{Tr}\rho}_{w}, \quad \text{et } D \text{ entier inversible dans } \hat{\mathbf{Z}} \leqno{(42)}
$$
et on retrouve $\mu$ en termes de $\xi$ par
$$
\mu = 1/D \leqno{(43)}
$$
(car $c = \operatorname{Tr}(\rho\varepsilon_1^{-1}) = \operatorname{Tr}\begin{pmatrix} 0 & 0 \\ 0 & 1 \end{pmatrix} = 1$), d'où
et $\eta$ par $\eta = \xi^{-1}\varepsilon_1(\mu-1)$.

Le sous-schéma en groupes $B'_0$ de la page 442 n'est autre que le schéma des matrices
$$
B'_0 = \left\{ \left. \begin{pmatrix} \mu & \lambda \\ 0 & 1 \end{pmatrix} \, \right| \, \mu \in \mathbf{G}_m, \lambda \in \mathbf{G}_a \right\} \leqno{(44)}
$$

Nous allons expliciter l'isomorphisme (39)
$$
u \longmapsto ([u,\rho], [u,\sigma], \det u) \leqno{(45)}
$$
en termes des coefficients $\mu, \lambda$ de $B'_0$, à l'aide des paramètres $v,w$ dans $\rho = \begin{pmatrix} 0 & 1 \\ v & w \end{pmatrix}$.

\newpage

%% ===== Page 305 =====

On trouve

$$
[u,\rho] = \left(\begin{matrix} \mu & \lambda \\ 0 & 1 \end{matrix}\right) \left(\begin{matrix} 0 & 1 \\ v & w \end{matrix}\right) \frac{1}{\mu} \left(\begin{matrix} 1 & -\lambda \\ 0 & \mu \end{matrix}\right) \frac{1}{-v} \left(\begin{matrix} w & -1 \\ -v & 0 \end{matrix}\right)
$$
$$
= \frac{1}{-\mu v} \left(\begin{matrix} \lambda v & \mu + \lambda w \\ v & w \end{matrix}\right) \left(\begin{matrix} w+\lambda v & -1 \\ -\mu v & 0 \end{matrix}\right)
$$
$$
= \frac{1}{\mu} \left(\begin{matrix} -v \lambda^2 + w \lambda \mu + \mu^2 - w \lambda & \lambda \\ -v \lambda + w \mu - w & 1 \end{matrix}\right)
$$

\marginpar{on pose $\sigma = \left(\begin{matrix} 0 & -1 \\ v' & w' \end{matrix}\right)$ \\ avec $v'=-v$ \\ $w'=1-w$}
$$
[u,\sigma] = \left(\begin{matrix} \mu & \lambda \\ 0 & 1 \end{matrix}\right) \left(\begin{matrix} 0 & -1 \\ v' & w' \end{matrix}\right) \frac{1}{\mu} \left(\begin{matrix} 1 & -\lambda \\ 0 & \mu \end{matrix}\right) \frac{1}{v'} \left(\begin{matrix} w' & 1 \\ -v' & 0 \end{matrix}\right)
$$
$$
= \frac{1}{\mu v'} \left(\begin{matrix} \lambda v' & -\mu + \lambda w' \\ v' & w' \end{matrix}\right) \left(\begin{matrix} w' + \lambda v' & 1 \\ -\mu v' & 0 \end{matrix}\right)
$$
$$
= \frac{1}{\mu} \left(\begin{matrix} v' \lambda^2 - w' \lambda \mu + \mu^2 + w' \lambda & \lambda \\ v' \lambda - w' \mu + w' & 1 \end{matrix}\right)
$$
$$
= \frac{1}{\mu} \left(\begin{matrix} -v \lambda^2 + (w-1) \lambda \mu + \mu^2 - (w-1) \lambda & \lambda \\ -v \lambda + (w-1) \mu - (w-1) & 1 \end{matrix}\right)
$$

En résumé\marginpar{abréviations}
$$
\begin{cases}
\xi = [u,\rho] = \frac{1}{\mu} \left(\begin{matrix} \mu^2 + (-v \lambda + w(\mu-1)) \lambda & \lambda \\ -v \lambda + w(\mu-1) & 1 \end{matrix}\right) \\ 
\eta = [u,\sigma] = \frac{1}{\mu} \left(\begin{matrix} \mu^2 + (-v \lambda + (w-1)(\mu-1)) \lambda & \lambda \\ -v \lambda + (w-1)(\mu-1) & 1 \end{matrix}\right)
\end{cases}
\leqno(46)
$$

$\mu = \det u$

donc si $\xi = \left(\begin{matrix} A & B \\ C & D \end{matrix}\right)$
$$
B = \lambda / \mu, \quad D = 1 / \mu \leqno(47)
$$
$$
\mu = 1 / D, \quad \lambda = B / D, \quad \text{donc } u = \left(\begin{matrix} 1/D & B/D \\ 0 & 1 \end{matrix}\right) \leqno(48)
$$

\newpage

%% ===== Page 306 =====

De même, si $\eta = \begin{pmatrix} R & S \\ T & U \end{pmatrix}$, on a $S=B, U=D$
$$ S = B = \lambda/\mu \ , \ U = D = 1/\mu \quad \text{d'où} \leqno{(47 \text{ bis})} $$
et $\mu = 1/U$, $\lambda = S/U$

donc
$$ u = \begin{pmatrix} 1/U & S/U \\ 0 & 1 \end{pmatrix} \leqno{(48 \text{ bis})} $$

Considérons maintenant l'équation
$$ [\alpha, \beta] = [\beta, \sigma] \varepsilon_1(11-1) \ , \ \text{avec} \leqno{(49)} $$
$$ \alpha = \begin{pmatrix} a & b \\ c & d \end{pmatrix} \ , \ \beta = \begin{pmatrix} r & s \\ t & u \end{pmatrix} \ , \leqno{(50)} $$
et posons, pour des calculs plus élégants
\marginpar{on ne pose $v_1=1, v'_1=1$ par la suite pour un calcul plus [unclear: simple] de $[\beta, \sigma]$ qu'après $[\alpha, \rho]$ !}
$$ \rho = \begin{pmatrix} 0 & v_1 \\ v & w \end{pmatrix} \ , \ \sigma = \begin{pmatrix} 0 & v'_1 \\ v' & w' \end{pmatrix} \leqno{(51)} $$
(i.e.\ on a écrit $v_1, v'_1$ au lieu de $u, u'$ pour ne pas interférer avec la notation $u$ pour un des coefficients de $\beta$), avec
$$ v' = -v \ , \ v'_1 = -v_1 \ , \ w+w' = v_1 \ , \ \text{posons} \leqno{(52)} $$
$$ \delta_\alpha = \det \alpha = ad-bc \ , \ \delta_\beta = \det \beta = ru-st \ , \leqno{(53)} $$

donnons alors
$$ \zeta = [\alpha, \rho] = \begin{pmatrix} a & b \\ c & d \end{pmatrix} \begin{pmatrix} 0 & v_1 \\ v & w \end{pmatrix} \frac{1}{\delta_\alpha} \begin{pmatrix} d & -b \\ -c & a \end{pmatrix} \frac{1}{-v v_1} \begin{pmatrix} w & -v_1 \\ -v & 0 \end{pmatrix} $$
$$ = \frac{1}{-\delta_\alpha v v_1} \begin{pmatrix} b v & a v_1 + b w \\ d v & c v_1 + d w \end{pmatrix} \begin{pmatrix} d w + b v & -d v_1 \\ -c w - a v & c v_1 \end{pmatrix} = \begin{pmatrix} A & B \\ C & D \end{pmatrix} $$

\marginpar{NB la formule pour $B/D$ [unclear: se simplifie]}
$$
\begin{cases}
(- \delta_\alpha v v_1) A = \dots \text{(barré)} \ , \quad (-\delta_\alpha v v_1) B = -v v_1 bd + v_1^2 ac + v_1 w bc \\
(- \delta_\alpha v v_1) C = \dots \text{(barré)} \ , \quad (-\delta_\alpha v v_1) D = -v v_1 d^2 + v_1 w cd + v_1^2 c^2
\end{cases} \leqno{(54)}
$$

\newpage

%% ===== Page 307 =====

et on trouve de même
$$
y = (\beta, \sigma) = \begin{pmatrix} R & S \\ T & U \end{pmatrix}
$$
$$
\begin{cases}
(-v v_1 \delta_\beta) R = \dots & (-v v_1 \delta_\beta) S = (v_1 v) su + v_1^2 tr - v_1 \omega' st \\
(-v v_1 \delta_\beta) T = \dots & (-v v_1 \delta_\beta) U = -(v_1 v) u^2 - v_1 \omega' tu + v_1^2 t^2
\end{cases}
\marginpar{\underline{NB} on peut diviser par $v_1$}
\leqno{(55)}
$$

On en tire en particulier
$$
\begin{cases}
D = \frac{-v v_1 d^2 + v_1 \omega cd + v_1^2 c^2}{-v v_1 (ad-bc)} \\
\mu = 1/D = \frac{-v v_1 (ad-bc)}{-v v_1 d^2 + v_1 \omega cd + v_1^2 c^2}
\end{cases}
\leqno{(56)}
$$

$$
B/D = \frac{-v_1 v bd + v_1^2 ac + v_1 \omega bc}{-v_1 v d^2 + v_1 \omega cd + v_1^2 c^2}
$$

ce qui permet donc d'expliciter
$$
\mathfrak{y} = \begin{pmatrix} 1/D = \mu & B/D = \lambda \\ 0 & 1 \end{pmatrix}
$$
en termes de $a, b, c, d$ (et des paramètres $v, v_1, \omega$ qui déterminent $\rho$).

Pour $v_1=1$, la formule pour $D$ de (56), s'écrit
$$
D = \det(d\rho+c)/-v \delta_\alpha = \det(d\rho+c)/\det \rho \det \alpha
\leqno{(57)}
$$

la condition de "Divis."
$$
D(\alpha) \text{ inversible } \iff d\rho+c \text{ inversible.}
\leqno{(58)}
$$

mais je m'aperçois que c'est une fausse piste,

\newpage

%% ===== Page 308 =====

\noindent
\underline{V \quad Étude (de l'équation 45) dans $GP(1)$, et la condition $\operatorname{Tr} \sigma = 0$.}
\vspace{0.5em}

Jusqu'à présent nous avons étudié l'équation (32)
$$
[\alpha, \rho] = [\beta, \sigma] \varepsilon_1(\mu_{-1}) \leqno{(*)}
$$
dans $G (\simeq \operatorname{Gl}(2))$, ou ce qui revient au même, en tant que relation dans $\operatorname{Sl}(2)$. Ici $\alpha, \beta$ sont des sections variables dans $G \simeq \operatorname{Gl}(2)$, mais $[\alpha, \rho]$ et $[\beta, \sigma]$ ne dépendent que des sections correspondantes $\tilde{\alpha}, \tilde{\beta}$ dans $PG (\simeq GP(1)) \simeq G/\mathbb{G}_m \simeq \operatorname{Aut}(C)$\footnote{centre de $G$} ([unclear: car ne dépendent des paramètres de $\rho, \sigma$ que par $\tilde{\rho}, \tilde{\sigma}$]). On peut donc considérer dans $SG \simeq \operatorname{Sl}(2)$ que c'est une équation liant $\tilde{\alpha}, \tilde{\beta}, \rho$.

On peut se proposer d'étudier l'équation correspondante dans $PG \simeq GP(1)$,
$$
[\tilde{\alpha}, \tilde{\rho}] = [\tilde{\beta}, \tilde{\sigma}] \tilde{\varepsilon}_1(\mu_{-1}) \leqno{(1)}
$$
ce qui s'écrit aussi, dans $\operatorname{Gl}(2)$, sous la forme équivalente
$$
[\alpha, \rho] = z [\beta, \sigma] \tilde{\varepsilon}_1(\mu_{-1}) \leqno{(2)}
$$
où $z$ est une section de $\mathbb{G}_m$, d'ailleurs en fait [unclear: un scalaire] uniquement déterminé par $\alpha, \beta, \mu$ (et ne dépend que de $\tilde{\alpha}, \tilde{\beta}, \mu$). On [unclear: verra]...

\newpage

%% ===== Page 309 =====

$$
\varepsilon^2 = 1 \leqno (3)
$$
(prendre le dét.\ des deux membres). On
considère (2) comme une équation
reliant les variables $\alpha, \beta, \mu, \varepsilon$ ($\alpha, \beta$ dans
$\operatorname{Gl}(2)$, $\mu$ [unclear: est] dans $\hat{\mathbf{Z}}^*$ -- et en fait [unclear: nécessairement]
dans $\mathbb{G}_m$, et $\varepsilon$ dans $\mu_2 = {}_2(\mathbb{G}_m)$), ou
[unclear: bien] $\tilde{\alpha}, \tilde{\beta}, \mu, \varepsilon$. Posons, en analogie avec (II),
$$
\xi = [\alpha, \rho] \quad , \quad \eta = \varepsilon [\beta, \sigma] \leqno (4)
$$
de sorte que (2) prend la forme
$$
\xi = \eta \varepsilon_1(\mu-1) \leqno (5)
$$
comme dans (II). Pour pouvoir appliquer
(II) à cette situation, il faudra s'assurer
que $\eta = \varepsilon [\beta, \sigma]$ satisfait aux [unclear: mêmes] conditions
préliminaires que 
$$
y_0 = [\beta, \sigma]
$$
lui-même, savoir $\det \eta = 1$ ($\Longleftrightarrow \varepsilon^2 = 1$) et
$$
\Tr \eta = \Tr \sigma
$$
$$
\varepsilon \Tr y_0 = \varepsilon \Tr \sigma
$$
i.e.
$$
\varepsilon \Tr \sigma = \Tr \sigma
$$

\newpage

%% ===== Page 310 =====

Lemme : Soit $\sigma^-$ une section de $G$ (dans ou $\operatorname{Gl}(2)$, associé à une alg. d'Azumaya $M$) -- i.e. section de $M$ telle que $\det \sigma^-$ inversible. Soit $A_{\sigma^-}$ le schéma des sections $y$ de $G$ satisfaisant aux conditions
$\det y = 1$ (i.e. $y$ dans $SL$), $\operatorname{Tr} y \sigma^- = \operatorname{Tr} \sigma^-$.

Conditions équivalentes a) b) c)\marginpar{d) On a $(\sigma^-)^2 \in \mathbb{G}_m$, i.e. $(\sigma^-)^2 = -1$} :
\begin{enumerate}
\item[a)] $A_{\sigma^-}$ est stable par translations par $\mu_2 \subset \mathbb{G}_m$ :
$$y \mapsto \varepsilon y \quad (\varepsilon^2 = 1)$$
\item[b)] $\operatorname{Tr} \sigma^- = 0$
\item[c)] (Dans le contexte du n° précédent, avec les normalisations (34) (38) pour $\rho, \sigma^-$) on a
$$ \sigma^- = \begin{pmatrix} 0 & -1 \\ \omega & 0 \end{pmatrix} = \begin{pmatrix} 0 & -1 \\ 1 & 0 \end{pmatrix} \leqno{(6)} $$
i.e. (avec les notations du loc. cit.) on a $\omega=1$, i.e.
$$ \rho = \begin{pmatrix} 0 & 1 \\ 1 & 1 \end{pmatrix} \leqno{(7)} $$
\end{enumerate}

Enfin ces conditions impliquent la [unclear: nullité] (où d'ailleurs l'équivalence purement [unclear: locale] a) $\iff$ b) s'impose [unclear: de façon évidente]) que $E'^{\natural 2} = (0)$, qui est donc une de ces implications immédiates. L'équivalence de b) et d) est immédiate, reste à prouver b) $\implies$ c). Or par le théorème de Hamilton-Cayley appliqué à $\sigma^-$ :
$$(\sigma^-)^2 - (\operatorname{Tr} \sigma^-) \sigma^- + \det \sigma^- = 0.$$

\newpage

%% ===== Page 311 =====

qui montre que $\mathrm{Tr}\sigma=0 \implies \sigma^2 = -\det\sigma \cdot 1$,
et inversement si $\sigma^2$ est un scalaire,
on aura que $(\mathrm{Tr}\sigma)\sigma = \sigma^2 + \det\sigma \cdot 1$ est
un scalaire, ce qui n'est possible (si $\sigma$ n'est
pas un scalaire sur aucune fibre) que
si $\mathrm{Tr}\sigma=0$.

Supposons par la suite ces conditions
satisfaites. Soit alors le
schéma des $(\alpha,\beta,\mu,\varepsilon)$ satisfaisant
$$ \mathcal{G}^*_{\rho,\sigma} \subset G \times G \times E' \times E' \leqno{(8)} $$
le schéma des solutions en $(\alpha,\beta,\mu,\varepsilon)$
de l'équation (2), qui est ainsi l'image
inverse dans $G \times G \times E' \times E'$ d'un sous-schéma
$$ \tilde{\mathcal{G}}^*_{\rho,\sigma} \subset PG \times PG \times E' \leqno{(9)} $$
formé des solutions en $(\alpha,\beta,\mu)$ de (4). La
projection sur le dernier facteur $E'$
fournit un morphisme
$$ \mathcal{G}^*_{\rho,\sigma} \to E' \leqno{(10)} $$

\newpage

%% ===== Page 312 =====

qui se factorise d'ailleurs par $\tilde{\mathcal{G}}^*_{\rho,\sigma}$ :
$$ \tilde{\varepsilon} : \tilde{\mathcal{G}}^*_{\rho,\sigma} \to \mu_2 . \leqno{(11)} $$
Notons qu'on a
$$ \mathcal{G}_{\rho,\sigma} \subset \mathcal{G}^*_{\rho,\sigma} $$
plus précisément
$$ \mathcal{G}_{\rho,\sigma} = \varepsilon^{-1}(1) \leqno{(12)} $$
ce qui conduit à introduire encore
$$ \tilde{\mathcal{G}}_{\rho,\sigma} \defeq \tilde{\varepsilon}^{-1}(1) , \leqno{(13)} $$
de sorte que $\mathcal{G}_{\rho,\sigma}$ est l'image
inverse dans $G \times G \times E'$ de $\tilde{\mathcal{G}}_{\rho,\sigma}$.

On a un morphisme
$$ (\alpha, \beta, \mu, \varepsilon) \longmapsto \left( \underset{\in \mathcal{H}_{\rho,\sigma}}{(\eta = \alpha \rho(\alpha)^{-1} = \varepsilon \beta \sigma \rho(\beta)^{-1}, \mu)}, \underset{\in \mu_2}{\varepsilon} \right) \leqno{(14)} $$
de
$$ \mathcal{G}^*_{\rho,\sigma} \to \mathcal{H}_{\rho,\sigma} \times \mu_2 $$
et l'image inverse d'une section
$( (\eta, \mu), \varepsilon )$ du second membre est le
schéma des solutions en $\alpha, \beta$ des équations
$$ (a)\quad \alpha \rho(\alpha)^{-1} = \eta, \qquad (b)\quad \beta \sigma \rho(\beta)^{-1} = \varepsilon \eta \leqno{(15)} $$

\newpage

%% ===== Page 313 =====

Or on a, via l'isom.
$$ \mathcal{B}_0 \isommap \mathcal{X}_{\rho,\sigma} $$
une solution privilégiée $\alpha = \underline{u}$ de la
première équation (15), qui est
caractérisée comme l'unique solution
qui soit une section de $\mathcal{B}_0$ (en vertu
de $T_\alpha \cap \mathcal{B}_0 = \{ 1 \}$ ).

Mais on a, pour ce $u$
$$ \underbrace{u \sigma(u)^{-1}}_{[u,\sigma]} = \eta \quad (\text{et non } \varepsilon \eta \ !) $$
pour trouver au moyen de ce $u$ une solution
de l'équation (15 b)
$$ \beta \sigma(\beta)^{-1} = \varepsilon \eta , $$
on est conduit : on l'essaye sous la forme
$$ \beta = u \tau \quad \text{d'où} \quad \beta \sigma(\beta)^{-1} = u \tau \sigma(\tau)^{-1} \sigma(u)^{-1} , \leqno{(16)} $$
en essayant de déterminer $\tau$ de façon
qu'on ait
$$ \underbrace{\tau \sigma(\tau)^{-1}}_{[\tau,\sigma]} = \varepsilon \leqno{(17)} $$
ce qui implique bien
$$ \beta \sigma(\beta)^{-1} = u \varepsilon \sigma(u)^{-1} = \varepsilon (u \sigma(u)^{-1}) = \varepsilon \eta $$
i.e.\ (15 b).

\newpage

%% ===== Page 314 =====

La relation (17) s'écrit aussi (en prenant l'inverse)
$$
\tau \sigma \tau^{-1} \sigma^{-1} = \varepsilon \quad \text{ou} \quad \sigma \tau \sigma^{-1} \tau^{-1} = \varepsilon \leqno{(18)}
$$
$$
\sigma = \varepsilon \tau(\sigma) \quad \text{ou} \quad \tau(\sigma) = \varepsilon \sigma \leqno{(19)}
$$
ce qui implique que $\tau$ normalise $T_\sigma$
$$
\tau(T_\sigma) = T_\sigma \text{ i.e. } \tau \in \mathcal{N}_\sigma \leqno{(20)}
$$
où
$$
\mathcal{N}_\sigma \defeq \operatorname{Norm}_G(T_\sigma). \leqno{(21)}
$$

Corollaire au lemme (p. 455) [unclear: les conditions algébriques] suivantes sont équivalentes :
\begin{enumerate}
\item[e)] [unclear: $\forall$ signe $\varepsilon$ de $\pm 1$], existe section $\tau$ de $G$ telle qu'on ait (17) (ou encore (19))
\item[f)] (si $2$ inv. dans $S$) idem, avec $\varepsilon = -1$.
\end{enumerate}

Si on a e), alors on voit que pour tout $\varepsilon$, on doit avoir $\operatorname{Tr} \varepsilon \sigma = \operatorname{Tr} \sigma$ ([unclear: grâce à l'invariance de la trace]), or $\operatorname{Tr} \varepsilon \sigma = \varepsilon \operatorname{Tr} \sigma$, donc ($\operatorname{Tr} \sigma(1-\varepsilon) = 0$). Si on la suppose satisfaite pour $\varepsilon = -1$, on aura $\operatorname{Tr} \sigma = - \operatorname{Tr} \sigma$ i.e. $2 \operatorname{Tr} \sigma = 0$, i.e. $\operatorname{Tr} \sigma = 0$ si $2$ inv. Inversement, si $\operatorname{Tr} \sigma = 0$, alors $\operatorname{Tr} \varepsilon \sigma = 0$, d'ailleurs $\det(\varepsilon \sigma) = \varepsilon^2 = 1$, donc $\varepsilon \sigma$ satisfait aux ``conditions numériques'' pour être de la forme $[\tau, \sigma]$ pour $\tau$ convenable. Puis $\dots$

\newpage

%% ===== Page 315 =====

assurer l'existence locale de $\tau$, il reste à s'assurer que $\varepsilon \sigma-$ est section régulière au sens de Steinberg i.e. n'est nulle part un scalaire, ce qui provient du fait qu'il en est ainsi de $\sigma-$.

Ainsi, sous les conditions actuelles, soit $\mathcal{N}'_{\sigma-} \subset \mathcal{N}_{\sigma-}$ le sous-schéma [unclear: de $G$] défini par
$$
\mathcal{N}'_{\sigma-} = \{ \tau \mid [\tau, \sigma-] \in \Gamma \mu_2 \} \leqno{(22)}
$$
alors $\mathcal{N}'_{\sigma-}$ est stable par translations à droite par $T_{\sigma-}$ — mieux, c'est un sous-schéma (contenant $T_{\sigma-}$) en groupes de $\mathcal{N}_{\sigma-}$, car si $\tau, \tau'$ sont tels que
$$
\tau(\sigma-) = \varepsilon \sigma- \quad , \quad \tau'(\sigma-) = \varepsilon' \sigma- ,
$$
$$
\tau'\tau(\sigma-) = \varepsilon' \varepsilon . \sigma- ,
$$
en même temps on voit que
$$
\tau \mapsto \varepsilon
$$
est un homomorphisme de groupes $\mathcal{N}'_{\sigma-} \to \mu_2$, et on a une suite exacte de schémas en groupes
$$
1 \to T_{\sigma-} \to \mathcal{N}'_{\sigma-} \xrightarrow{\varepsilon} \mu_2 \to 1 . \leqno{(23)}
$$
On notera que si $\tau \in \Gamma G$ normalise $T_{\sigma-}$, on a
$$
\tau(\sigma-) = a \sigma- + b \quad a,b \in \Gamma \mathcal{O}_S
$$

\newpage

%% ===== Page 316 =====

et prenant les traces des deux membres on
trouve (compte tenu que $\operatorname{Tr} \tau(\sigma) = \operatorname{Tr} \sigma = 0$)
$\operatorname{Tr} b \cdot 1 = 0$ i.e. $2b = 0$, donc si $2$ est inversible,
$b=0$, donc $\tau(\sigma) = a \sigma$ et nécessairement
on a $a^2 = 1$ i.e. $a$ est un $\varepsilon$. Donc
$$
N'_\sigma = N_\sigma \quad \text{en car. résiduelles } \neq 2. \leqno (24)
$$

Plus génér., sur une base $S$ telle que
$2$ ne soit pas diviseur de zéro (i.e. régulier)
dans $\Gamma(S, \mathcal{O}_S)$, on voit que toute
section de $N'_\sigma$ est en fait une section de
$N_\sigma$. Mais il n'en est plus ainsi en
car $2$, où
$$ \sigma' = a \sigma + b \quad (a \in \Gamma(\mathcal{O}_S), b \in \Gamma(\mathcal{O}_S)) $$
dès qu'il satisfait
$\operatorname{Tr} \sigma' = 0$, $\det \sigma' = \det \sigma$
est conjugué loc! à $\sigma$ - et il suffit pour
s'en assurer (comme $\det(a\sigma+b) = a^2 \det \sigma + a b \operatorname{Tr} \sigma + b^2$)
[unclear: d'écrire]
$$
\begin{cases}
-v a^2 + b^2 = -v \\
2b = 0 \quad (\text{automatique en car. 2})
\end{cases}
$$
qui n'impliquent nullement que l'on ait $b=0$ \dots

Considérons maintenant le
groupe produit
$$
B'_0 \times A_p \times N'_\sigma
$$

\newpage

%% ===== Page 317 =====

et le morphisme de schémas
$$
\begin{cases} B'_0 \times A_\rho \times N'_0 \longrightarrow G \times G \times E' \times E' \\ (\underline{u}, \underline{a}, \underline{b}) \longmapsto (\alpha = \underline{u}\underline{a}^{-1}, \beta = \underline{u}\underline{b}^{-1}, \mu = \det \underline{u}, \varepsilon = \varepsilon_\rho(\underline{b})) \end{cases} \leqno{(25)}
$$
je dis que l'image est dans $\mathcal{G}_{\rho,\sigma}^*$, i.e. qu'on a bien
$$
\alpha \rho(\alpha)^{-1} = \varepsilon (\beta \sigma(\beta)^{-1}) \varepsilon_1(\mu^{-1})
$$
En effet
$$
\alpha \rho(\alpha)^{-1} = \underline{u} \rho(\underline{u})^{-1}
$$
$$
\beta \sigma(\beta)^{-1} = \underline{u} \underbrace{\underline{b}^{-1} \sigma(\underline{b})}_{\varepsilon_\rho(\underline{b})^{-1} = \varepsilon_\rho(\underline{b})} \sigma(\underline{u})^{-1} = \varepsilon \underline{u} \sigma(\underline{u})^{-1}
$$
et on sait (par II) que
$$
\underline{u} \rho(\underline{u})^{-1} = \underline{u} \sigma(\underline{u})^{-1} \varepsilon_1(\mu^{-1})
$$
d'où ça, puisque $\varepsilon^2 = 1$. Mais de plus le morphisme (25) est un isomorphisme, car en termes de $(\alpha, \beta, \mu, \varepsilon)$, $\underline{u}$ est déterminé (en termes de $\alpha$ déjà), comme l'unique section de $B'_0$ qui soit de la forme $\alpha \underline{a}$, où $\underline{a} \in \Gamma_{T_\rho}$ (car $\Gamma_{T_\rho} \cap B'_0 = \{1\}$), et $\underline{b}$ est alors déterminé par la relation $\beta = \underline{u}\underline{b}^{-1}$. On trouve donc que (25) est un iso
$$
\begin{cases} B'_0 \times A_\rho \times N'_0 \xrightarrow{\sim} \mathcal{G}_{\rho,\sigma}^* \\ (\underline{u}, \underline{a}, \underline{b}) \longmapsto (\alpha = \underline{u}\underline{a}^{-1}, \beta = \underline{u}\underline{b}^{-1}, \mu = \det \underline{u}, \varepsilon = \varepsilon_\rho(\underline{b})) \end{cases} \leqno{(26)}
$$

\newpage

%% ===== Page 318 =====

Si on a deux sections $(\underline{u}, \underline{a}, \underline{b})$ et $(\underline{u}', \underline{a}', \underline{b}')$ dans $\mathcal{B}'_0 \times \mathcal{A}_{\overline{p}} \times \mathcal{N}'_0$, à quelle condition a-t-on
$$
\widetilde{\alpha} = \widetilde{\alpha}', \quad \widetilde{\beta} = \widetilde{\beta}' \ ? \leqno{(27)}
$$
Il faut évidemment qu'on ait
$$
\underline{a}' = \lambda_1 \underline{a} , \quad \underline{b}' = \lambda_2 \underline{b} \quad \lambda_1, \lambda_2 \in \Gamma_{G} \leqno{(28)}
$$
et je dis que ces conditions sont aussi nécessaires. Cela provient du fait que $\mathcal{B}'_0 \to PG \simeq G P(N)$ est mono (soit $\overline{\mathcal{B}}'_0$ l'image de $\mathcal{B}'_0$), que $\mathcal{A}_{\overline{p}} \to PG$ et $\mathcal{N}'_0 \to PG$ ont [unclear: $\Gamma$ comme noyau] (soient $\overline{\mathcal{T}}_{\overline{p}}$ et $\overline{\mathcal{N}}'_0$ les images), et que dans $PG$
$$
\overline{\mathcal{T}}_{\overline{p}} \cap \overline{\mathcal{B}}'_0 = \{1\} \leqno{(29)}
$$
de sorte que $\overline{u} \in \overline{\mathcal{B}}'_0$ et $\overline{a} \in \overline{\mathcal{T}}_{\overline{p}}$ sont définis de façon unique en termes de $\widetilde{\alpha}$ par la relation
$$
\widetilde{\alpha} = \overline{u} \overline{a}^{-1}
$$
-- et qu'alors $\widetilde{\underline{b}}$ est défini par la relation
$$
\widetilde{\beta} = \underline{u} \widetilde{\underline{b}}^{-1} .
$$

Ceci montre que l'on a un isomorphisme
$$
\begin{cases}
\mathcal{B}'_0 \times \mathcal{A}_{\overline{p}} \times \mathcal{N}'_0 \isommap \mathcal{G}_{p,\sigma}^* \subset R_{1} \times P_1 \times E' \\
(\underline{u}, \widetilde{\underline{a}}, \widetilde{\underline{b}}) \mapsto (\underline{u} \widetilde{\underline{a}}^{-1}, \underline{u} \widetilde{\underline{b}}^{-1}, \varepsilon(\underline{u}))
\end{cases} \leqno{(30)}
$$

\newpage

%% ===== Page 319 =====

où
$$
\delta(\tilde{u}) = \det(\underline{u}) \leqno{(31)}
$$
quand $\tilde{u}$ provient de $\underline{u} \in \Gamma B_0'$ -- ce qui
revient au même (cf p. 441) $\delta(\tilde{u})$ est
défini par
$$
\tilde{u}(N_0) = \delta(\tilde{u})N_0 \ . \leqno{(32)}
$$

Rmq. Notons qu'une section $(\tilde{\alpha}, \tilde{\beta}, \mu)$ \footnote{NB le centralisateur de $\overline{\Delta}_p$ dans $G \simeq P\Gamma(\dots)$ est $\overline{T}_0$, une fois [unclear: $\mu = \pm 1$] exclue} de
$\mathcal{G}_{\rho, \sigma}$ est \marginpar{mod op. de [unclear: $N_0'$ sur $\beta$]} déterminée, par
\footnote{NB $\tilde{\alpha}$ dépend du trio par [unclear: sa projection $\alpha$]}
[unclear: ce qui] précède, en termes de $\tilde{\alpha}$ --
condition à satisfaire sur $\tilde{\alpha}$, pour
qu'il existe $\tilde{\beta}, \mu$ qui la complètent
en une section de $\mathcal{G}_{\rho, \sigma}^*$, étant la
suivante\footnote{NB On a [unclear: d'ailleurs] le relèvement exact [unclear: de] $\alpha \in \Gamma G$, qui est [unclear: défini indépendamment] du choix du relèvement} :
Le coefficient $D$ de $\xi = \alpha \rho(\alpha)^{-1} = \begin{pmatrix} A & B \\ C & D \end{pmatrix}$ est inversible (ce qui permet
de déterminer $\mu$ par $\mu = 1/D$), puis
on détermine $\tilde{\gamma}$ par $\tilde{\gamma} = \tilde{\xi} (\mu - 1)$, et
enfin on trouve $\tilde{\beta}$ par la condition
$$
\tilde{\rho}^{\tilde{\alpha}} \tilde{\beta}^{-1} = \tilde{\gamma} \ .
$$
De même, une section de $\mathcal{G}_{\rho, \sigma}^*$
(ou $\alpha, \beta, \mu, \varepsilon$) est déterminée par [unclear: $\alpha$] seul, mod
[unclear: $B_0 \times \dots \times N_0$] on trouve d'abord $\underline{u}$, en termes

\newpage

%% ===== Page 320 =====

de $\alpha$ comme dessus, d'où $\mu = \det \alpha$, d'où $\varepsilon$, et il reste ensuite à déterminer $\beta$ comme $\beta = u b^{-1}$, avec $b \in \Gamma/\mathcal{N}_{0'}$, arbitraire, $\varepsilon$ étant défini alors en termes du choix de $b$ comme cl. de $\varepsilon(b)$. La condition pour que $\alpha$ fixé puisse se compléter par $(\beta, \mu, \varepsilon)$ en $(\alpha, \beta, \mu, \varepsilon) \in \Gamma_{\hat{\mathfrak{S}}_{\hat{T}^{0'}}}$ est [unclear: tirée] la condition que le coeff. $D$ de 
$$ \zeta = \alpha \rho (\alpha)^{-1} = \begin{pmatrix} A & B \\ C & D \end{pmatrix} \leqno (33) $$
soit inversible, ce qui permet de déterminer $\mu$ par
$$ \mu = 1/D \leqno (34) $$

Posons
$$ \alpha = \begin{pmatrix} a & b \\ c & d \end{pmatrix} \leqno (35) $$
et explicitons $\alpha \rho (\alpha)^{-1}$ ; où $\rho = \begin{pmatrix} 0 & 1 \\ 1 & 1 \end{pmatrix}$ \marginpar{\delta = \det \alpha = ad-bc}
$$ \alpha \rho (\alpha)^{-1} = \begin{pmatrix} a & b \\ c & d \end{pmatrix} \begin{pmatrix} 0 & 1 \\ 1 & 1 \end{pmatrix} \frac{1}{\delta} \begin{pmatrix} d & -b \\ -c & a \end{pmatrix} (-1) \begin{pmatrix} 1 & -1 \\ -1 & 0 \end{pmatrix} $$
$$ = - \frac{1}{\delta} \begin{pmatrix} b & a+b \\ d & c+d \end{pmatrix} \begin{pmatrix} d+b & -d \\ -c-a & c \end{pmatrix} \quad , \quad d'où $$
$$ 
\begin{aligned} 
(-\delta) A &= b(d+b) - (a+b)(c+a) & (-\delta) B &= -bd + c(a+b) \\ 
(-\delta) C &= d(d+b) - (c+d)(c+a) & (-\delta) D &= -d^2 + cd + c^2 
\end{aligned} \leqno (36) 
$$
La quantité la plus essentielle étant donc
$$ (-\delta) D = -d^2 + cd + c^2 \quad , \leqno (37) $$

\newpage

%% ===== Page 321 =====

$$
(- \delta v) B = a c + b c - v b d  \leqno{(38)}
$$

ce qui redonne $u$ en termes de $a, b, c, d$,
grâce à la formule $\underline{\text{IV}} (48)$ (p. 453')
$$
u = \begin{pmatrix} 1/D & B/D \\ 0 & 1 \end{pmatrix} = \begin{pmatrix} \frac{-v(ad-bc)}{-vd^2+cd+c^2} & \frac{ac+bc-vbd}{-vd^2+cd+c^2} \\ 0 & 1 \end{pmatrix} \leqno{(39)}
$$

Remarques. Par contre, la donnée de $\bar{\beta}$ seule (sans le "signe" $\varepsilon$) ne permet pas de récupérer $\mu$. Mais la donnée de $(\bar{\beta}, \varepsilon)$ le permet, car on en déduit la [unclear: matrice] 
$$
\eta = \varepsilon \beta v \beta^{-1} \defeq \begin{pmatrix} R & S \\ T & U \end{pmatrix} , \leqno{(40)}
$$
qui sera relié à $\zeta = \alpha \rho \alpha^{-1}$ par
$$
\zeta = \eta \ e_1(\mu^{-1})
$$
et par le [unclear: théorème] de (II), on a
$$
\mu = 1/U \leqno{(41)}
$$
(mais en termes de $\beta v \beta^{-1} \defeq \begin{pmatrix} R_0 & S_0 \\ T_0 & U_0 \end{pmatrix}$)
on aura
$$
\mu = \frac{\varepsilon}{U_0} \leqno{(42)}
$$
i.e. il n'est pas [unclear: connu] par $U_0$ seul, mais seulement par $U_0$ et $\varepsilon$. 

Pour préciser ces points, on va regarder quand est-ce que, pour $(\bar{\beta}, \bar{\gamma})$ fixés, on peut compléter par $\alpha, \mu$ de façon à ...

\newpage

%% ===== Page 322 =====

\marginpar{461}
section $(\alpha, \beta, \mu, \varepsilon)$ de $\mathcal{G}_{\beta,\sigma}$. Une condition nécessaire est que l'on ait $U_0$ inversible (elle ne dépend que de $\beta$, pas de $\varepsilon$) --- et alors par le choix de $\varepsilon$ on déduit un élément de $\Gamma_{g,\nu}$, soit $\eta = \varepsilon \beta \sigma \gamma \beta^{-1}$ satisfaisant dans l'[unclear: anneau] de $\mathcal{G}$ aux relations de $\det \eta = 1, \operatorname{Tr} \eta \sigma = \operatorname{Tr} \sigma$ (cf.\ lemme p.~455), et par le théorème de (II), pour compléter par $\mu = 1/U = \frac{\varepsilon}{U_0}$ la section, est définie
$$ \xi = \eta \, \varepsilon_{c(\mu-1)} , \qquad \text{loc.cit.} $$
(pour que le tout s'écrive dans la forme $\alpha \rho(\alpha)^{-1}$, l'indétermination de $\alpha$ d'ailleurs dans $\mathfrak{T}_g$ via $\alpha \mapsto \alpha' = \alpha b, b \in \Gamma_g$).

Ceci montre bien que $\mu$ ne peut prendre, pour un $\beta$ fixé, toutes les valeurs
$$ \frac{\varepsilon}{U_0} \quad \left( \beta \sigma \gamma \beta^{-1} = \begin{pmatrix} R_0 & S_0 \\ T_0 & U_0 \end{pmatrix} \right) \dots $$

\newpage

%% ===== Page 323 =====

\noindent
(VI) \underline{Caractérisation des triplets $(E, \rho, \sigma)$ modulaires.}

1) Soient $G, SG, PG$ des formes associées de $\operatorname{Gl}(2), \operatorname{Sl}(2), P\operatorname{Gl}(2)$ sur une base $S$ --- correspondant donc à un fibré en coniques $\mathcal{P}$ (i.e.\ torseur sous $P\operatorname{Gl}(2)_S$), ou une Algèbre d'Azumaya $\mathcal{A}$ de rang $4$, et considérons deux sections $\tilde{\rho}, \tilde{\sigma}$ de $PG$, telles que
$$
\tilde{\rho}^2 = \tilde{\sigma}^3 = 1 \leqno (1)
$$
soit la section $\tilde{\tau}_0 \defeq \tilde{\rho} \tilde{\sigma}$\footnote{NB Si $\tilde{\tau}_0 = \tilde{\rho} \tilde{\sigma} = 1$, on a par (1) $\tilde{\rho} = \tilde{\sigma}^{-1}$, donc $\tilde{\sigma}^2 = \tilde{\rho}^2 = 1$, d'où $1 = \tilde{\sigma}^3 = \tilde{\sigma}$, i.e.\ $\tilde{\rho} = \tilde{\sigma} = 1$, et $x_0 = x_1$ --- ce cas est donc écarté (sur toute fibre).},
dont on suppose 
équivalente à : $\tilde{\tau}_0$ a d.\ p.\ f., i.e.\ que les sections $x_0, x_1$ de $\mathcal{P}$ qui correspondent à ses points fixes $x_0, x_1$, soient partout distinctes dans chaque fibre).
Alors il existe un isomorphisme unique\footnote{NB L'unicité de l'isomorphisme (2) vient de ce que (par $\tilde{\tau}_0 \neq 1$) $\rho, \sigma$ engendrent l'algèbre $M_2(\mathcal{O}_S)$.}
$$
\operatorname{Gl}(2)_S \isommap G \leqno (2)
$$
et des sections $\rho, \sigma, w$ de $G$, vérifiant
$$
\rho^2 = w, \quad \sigma^3 = w \leqno (3)
$$
$$
w^2 = 1, \quad w \in \Gamma(\mathcal{O}_S^\times) \leqno (4)
$$
telles que $\tilde{\rho}, \tilde{\sigma}$ soient les images, par le morphisme $\operatorname{Gl}(2)_S \isom G \to PG$, des sections
$$
\rho = \begin{pmatrix} 0 & 1 \\ w & 0 \end{pmatrix}, \quad \sigma = \begin{pmatrix} 0 & -1 \\ w & -1 \end{pmatrix} \leqno (5)
$$

\newpage

%% ===== Page 324 =====

de $\operatorname{Gl}(2)_S$ - et inversement, pour un épinglage $y$ et des sections $v, w, v', w'$ satisfaisant (4) - i.e.
$$
v' = -v , \quad w' = 1-w , \quad \begin{matrix} v \in \Gamma(S, \underline{O}_S^\times) \\ w \in \Gamma(S, \underline{O}_S) \end{matrix} \leqno{(4\text{bis})}
$$
les images $\tilde{\rho}, \tilde{\sigma}$ par la composition $\operatorname{Gl}(2)_S \to PG \simeq G/G_m$ des sections $\rho, \sigma$ de $\operatorname{Gl}(2)_S$ données par (5) satisfont les conditions plus haut ($\tilde{\rho} \tilde{\sigma}^{-1} \defeq \tilde{\varepsilon}_1$ est unipotente régulière, et partout non équivalente à $\tilde{\sigma}^{-1}(\tilde{\varepsilon}_1) = \tilde{\rho}^{-1}(\tilde{\varepsilon}_1) \defeq \tilde{\varepsilon}_0$).

Posons alors
$$
\varepsilon_1 = \begin{pmatrix} 1 & 0 \\ 1 & 1 \end{pmatrix} , \quad \varepsilon_0 = \begin{pmatrix} 1 & 1/v \\ 0 & 1 \end{pmatrix} \leqno{(6)}
$$

on a
$$
\rho \sigma^{-1} = - \varepsilon_1 , \quad \sigma^{-1} \rho = - \varepsilon_0 \leqno{(7)}
$$
$$
\varepsilon_0 = \rho(\varepsilon_1) = \sigma^{-1}(\varepsilon_1) , \quad \varepsilon_1 = \rho(\varepsilon_0) = \sigma(\varepsilon_0) \leqno{(7\text{bis})}
$$
dont les images par $\operatorname{Gl}(2)_S \to PG$ des sections $\varepsilon_1, \varepsilon_0$ sont $\tilde{\varepsilon}_1, \tilde{\varepsilon}_0$. On a les formules
$$
\operatorname{Tr} \sigma = w' , \quad \operatorname{Tr} \rho = w , \quad \operatorname{Tr} (\sigma^{-1} \rho) = 1 \leqno{(8)}
$$
$$
\sigma + \rho = \begin{pmatrix} 0 & 0 \\ 0 & 1 \end{pmatrix} \defeq N \leqno{(9)}
$$
$$
\operatorname{Tr} \xi N = D \quad \text{si} \quad \xi = \begin{pmatrix} A & B \\ C & D \end{pmatrix} , \leqno{(10)}
$$
$$
\det \sigma = \det \rho = -v = -v' . \leqno{(11)}
$$

\newpage

%% ===== Page 325 =====

2°) Les conditions suivantes sont équivalentes

(12)
\begin{tikzcd}
\boxed{Tr \sigma = 0} \ar[r, "\Leftrightarrow"] \ar[d, "\Updownarrow"] & Tr \rho = 1 \ar[r, "\Leftrightarrow"] \ar[d] & \boxed{\tilde{\sigma}^2 = 1} \ar[d] \\
\boxed{w' = 0} & \boxed{w = 1} & \sigma^2 \in \mathbb{G}_m \ar[d] \\
& & \boxed{\sigma^2 = v \text{ id.}}
\end{tikzcd}

Elles équivalent encore à ceci : désignant par $N'_\sigma$ le sous-schéma en groupes de $\operatorname{Gl}(2)_S$ défini par la condition
(13) $\tau(\sigma) = \varepsilon \sigma$ i.e. $\tau(\sigma) = \varepsilon \sigma$, $\varepsilon \in \mathbb{G}_m$
(avec nécessairement $\varepsilon^2 = 1$ i.e. $\varepsilon \in \mu_2$)
le morphisme

(14)
\begin{tikzcd}
N'_\sigma \ar[r] & \mu_2 \\
\tau \ar[r, |->] & \varepsilon
\end{tikzcd}

est un épimorphisme, de façon qu'on a une suite exacte

(15) $$1 \longrightarrow T_\sigma \longrightarrow N'_\sigma \longrightarrow \mu_2 \longrightarrow 1$$

où
(16) $T_\sigma = Cent_{\operatorname{Gl}(2)_S}(\sigma) = \{ \dots \}$
($N_\sigma \subset N'_\sigma$ est $\operatorname{Norm}_{\operatorname{Gl}(2)_S}(T_\sigma)$, les
[unclear]
une section de $N'_\sigma$ [unclear] vers $\mu_2$ est une section de
$N_\sigma$ [unclear] en fait. On a [unclear]
(18) $N'_\sigma = N_\sigma \quad (\text{car } \dots)$

La condition ci-dessus signifie que $\sigma, \rho$ prennent la forme

\newpage

%% ===== Page 326 =====

$$
\sigma = \begin{pmatrix} 0 & -1 \\ v' & w' \end{pmatrix}, \quad \rho = \begin{pmatrix} 0 & 1 \\ v & w \end{pmatrix} \leqno{(19)}
$$
avec $v+v'=0$.

Ces résultats de $1^\circ)$, $2^\circ)$ ne sont qu'une [unclear: reformulation] d'une partie des résultats des sections IV, V précédentes. On vient d'étudier des conditions qui reviennent à la [unclear: détermination] de $w$ (ou $w, w'$) de façon très explicite. On va se tourner vers des conditions qui font revenir à fixer $v$ (ou encore $v'$) de façon explicite.

$3^\circ)$ Conditions équivalentes
$$
\begin{cases}
\det \sigma = 1 \iff \det \rho = 1 \iff v = -1 \iff v' = 1 \\
\iff \varepsilon_0 = \begin{pmatrix} 1 & -1 \\ 0 & 1 \end{pmatrix}
\end{cases} \leqno{(20)}
$$

Ces équivalences résultent aussitôt de (6) et (11). Les conditions signifient que $\sigma$ et $\rho$ sont de la forme
$$
\sigma = \begin{pmatrix} 0 & -1 \\ 1 & w' \end{pmatrix}, \quad \rho = \begin{pmatrix} 0 & 1 \\ -1 & w \end{pmatrix} \leqno{(21)}
$$
$$
(w+w'=1)
$$

$4^\circ)$ Conjonction de $2^\circ)$ et $3^\circ)$.
Conditions équivalentes :

\newpage

%% ===== Page 327 =====

$$
\begin{matrix}
\text{Conjonction des} \\
\text{conditions de } 2^\circ \\
\text{et des conditions} \\
\text{de } 3^\circ
\end{matrix}
\iff \sigma^2 = -1 \quad \text{i.e. } \sigma^2+1=0 \iff \tilde{\sigma}^2 = \tilde{\rho}^3 = 1 \implies \begin{bmatrix} \rho^3 = -1 \\ \rho^2-\rho+1=0 \end{bmatrix} \text{ (cf. rem. plus bas)}
$$

$$
\begin{cases}
v = -1, \omega = 1 \iff \rho = \begin{pmatrix} 0 & 1 \\ -1 & 1 \end{pmatrix} \iff \Tr \rho = \det \rho = 1 \\
v' = 1, \omega' = 0 \iff \sigma = \begin{pmatrix} 0 & -1 \\ 1 & 0 \end{pmatrix} \iff (\Tr \sigma = 0, \det \sigma = 1)
\end{cases} \leqno{(22)}
$$

Dém. Le fait que la conjonction des conditions de $2^\circ)$ et $3^\circ)$ soit équivalente à chacune des conditions en $2^\text{e}$ et $3^\text{e}$ lignes, est trivial. L'équivalence avec $\sigma^2=-1$ provient du fait que la condition $2^\circ$ signifie que $\sigma^2=v \, \text{id}$, tandis que $3^\circ$ signifie $v=-1$, d'où $\sigma^2=-1$ équivaut bien à la conjonction.

Notons que\footnote{Pour une autre démonstration de l'équivalence de $2^\circ + 3^\circ)$ avec $\tilde{\sigma}^2=\tilde{\rho}^3=1$, cf. Remarque ci-dessous.} $2^\circ)$ s'écrit aussi $\tilde{\sigma}^2=1$, et équivaut à : $\omega'=0$, soit $\omega=1$, et -- d'autre part, en écrivant $\rho$
$$
\rho = \begin{pmatrix} 0 & 1 \\ v & 1 \end{pmatrix}
$$
$$
\rho^3 = \begin{pmatrix} v & v+1 \\ v(v+1) & 2v+1 \end{pmatrix} \leqno{(23)}
$$
donc $\rho^3$ ne peut être un scalaire i.e. $\tilde{\rho}^3=1$ que si les deux termes diagonaux sont égaux, i.e. leur diff. est $(2v+1)-v=v+1$, nulle, donc $v=-1$,

\newpage

%% ===== Page 328 =====

donc la condition s'écrit $\upsilon = -1$, i.e. c'est
la condition $3^\circ)$, mais alors l'expression
(23) de $\rho^3$ nous montre que $\rho^3 = -1$.

Remarques. Des conditions qui précèdent
impliquent les formules
$$ \varepsilon_0 = \sigma \rho \leqno{(24)} $$
$$ \varepsilon_1 = \rho \sigma \leqno{(25)} $$
$$ \rho \varepsilon_1 = \varepsilon_0^{-1} \leqno{(26)} $$
$$ \rho^2 \sigma = \varepsilon_0^{-1} \leqno{(27)} $$
dont chacune\footnote{sauf la dernière (27)} les implique à son tour

- p. ex. on a $\varepsilon_0 = - \sigma^{-1} \rho$, donc (24)
équivaut à $\sigma = - \sigma^{-1}$, i.e. $\sigma^2 = -1$, et (26) équivaut à
(25) en multipliant par $\rho$ aux deux membres,

On a d'ailleurs
$$ \rho \varepsilon_1 = \begin{pmatrix} 0 & 1 \\ \upsilon & w \end{pmatrix} \begin{pmatrix} 1 & 1 \\ 0 & 1 \end{pmatrix} = \begin{pmatrix} 0 & 1 \\ \upsilon & \upsilon+w \end{pmatrix} , \quad \varepsilon_0^{-1} = \begin{pmatrix} 1 & -1 \\ 0 & 1 \end{pmatrix} $$
qui montre que $\rho \varepsilon_1 = \varepsilon_0^{-1}$ équivaut à
$w=1, \upsilon=-1$ i.e. $\rho = \begin{pmatrix} 0 & 1 \\ -1 & 1 \end{pmatrix}$,
(d'où la forme $\rho \varepsilon_1 = \rho^2 \sigma$) OK.

Comprendre (25) et (26) --
comme conséquences des conditions
$2^\circ) + 3^\circ)$, est un calcul brutal :

$$ \rho^2 \sigma = \begin{pmatrix} -vw & w - (v+w^2) \\ -w(v+w^2) & -w(w^2+v)+vw^2 \end{pmatrix} \quad \text{[rayé]} $$

L'égalité $\rho^2 \sigma = \varepsilon_0^{-1}$ [unclear: qui donne $-v = -1$ i.e. $v=1$]
implique $w = 1, v = 1$, et $v+w^2 = 0$, (i.e. $2=0$ !)

\newpage

%% ===== Page 329 =====

$$
\rho^2 \sigma = \begin{pmatrix} w^3 & w \\ 0 & -w^3 \end{pmatrix}
$$

cela montre que
$$
\rho^2 \sigma = \varepsilon_0^{-1} \iff v+w^2=0, w^3=1 \quad (\implies v = -w^2 = -1/w, \text{ i.e. } vw = -1, \text{ OK}) \leqno{(27 \text{ bis})}
$$
Mais $(27)$, en multipliant par $\rho$, équivaut à :
$$
\rho^3 \sigma = \rho \varepsilon_0^{-1} \leqno{(27')}
$$
or on sait que $\varepsilon_0 = -\sigma^{-1} \rho$ i.e. $\varepsilon_0^{-1} = -\rho^{-1} \sigma$
i.e. $\rho \varepsilon_0^{-1} = -\sigma$ ; donc la relation précédente s'écrit
$$
\rho^3 \sigma = -\sigma \quad \text{ou encore}
$$
$$
\rho^3 = -1 \leqno{(28)}
$$
\marginpar{\dots et (28) est la condition (27 bis)}
On peut voir cette équivalence directement en calculant $\rho^3$ en fonction de $v, w, -$
on trouve :
$$
\rho^3 = \begin{pmatrix} vw & vw^2 \\ v(v+w^2) & w(2v+w^2) \end{pmatrix} \leqno{(29)}
$$
d'où on ne peut avoir $\tilde{\rho}^3 = 1$ i.e. $\rho^3 \sim$ scalaire que si $v+w^2=0$ i.e. $v = -w^2$
ce qui implique
$$
\rho^3 = \begin{pmatrix} -w^3 & 0 \\ 0 & -w^3 \end{pmatrix}
$$
i.e. on trouve

Lemme. On a $\tilde{\rho}^3 = 1 \iff v+w^2=0$ i.e. $v = -w^2$, i.e. on a
$$
\rho = \begin{pmatrix} 0 & 1 \\ -w^2 & w \end{pmatrix} \leqno{(30)}
$$
et alors on a
$$
\rho^3 = -w^3 \cdot 1 \leqno{(31)}
$$
donc on a
$$
\rho^3 = -1 \iff w^3=1 \text{ i.e. } w \in \mu_3 \text{ et } v=-w^2 (=-1/w) \leqno{(32)}
$$

\newpage

%% ===== Page 330 =====

Corollaire ! Les conditions équivalentes (32)
équivalent aussi à
$$
\rho^3 = -1 \text{ et } \det \rho = 1 \leqno{(33)}
$$
ou encore :
$$
\rho^3 = -1 \text{ et } \operatorname{Tr} \rho = 1 \leqno{(34)}
$$
en d'autres termes, considérant les
trois relations
$$
\det \rho = 1, \quad \operatorname{Tr} \rho = 1, \quad \rho^3 = -1 \leqno{(35)}
$$
la conjonction de deux de celles-ci impliquent la troisième,
et équivalent aux conditions équivalentes
de (32).

Notons une autre façon de trouver ce
corollaire, en notant que
$$
\rho^3 + 1 = (\rho + 1)(\rho^2 - \rho + 1) \quad \text{(identité algébrique)}
$$
\marginpar{faux}or $\rho$ n'est pas un scalaire sur aucun [unclear: corps $k$],
donc $\rho+1$ est inversible, d'où $\rho^3 = -1$ i.e. $\rho^3 + 1 = 0$
équivaut à la relation $\rho^2 - \rho + 1 = 0$ :
$$
\rho^2 - \rho + 1 = 0 \quad (\iff \rho^3 = -1), \leqno{(36)}
$$
qu'on compare avec Hamilton-Cayley
$$
\rho^2 - (\operatorname{Tr} \rho)\rho + \det \rho = 0 \leqno{(37)}
$$
donc $\operatorname{Tr} \rho = \det \rho = 1$ impliquent (36),
d'autre part si on a (36), on conclut
via 37
$$
(\operatorname{Tr} \rho - 1) \rho = (\det \rho - 1) 1
$$

\newpage

%% ===== Page 331 =====

Une autre façon de trouver certaines
équivalences est via les équations de Hamilton-Cayley pour $\sigma$ et $\rho$
$$ \sigma^2 - (\underbrace{\Tr \sigma}_{[unclear: w'=0]})\sigma + \underbrace{\det \sigma}_{[unclear: v'=1]} = 0 \leqno{(36)} $$
$$ \rho^2 - (\underbrace{\Tr \rho}_{w})\rho + \underbrace{\det \rho}_{v} = 0 \leqno{(37)} $$
L'équation (36) redonne que la conjonction
de $\Tr \sigma=0$, $\det \sigma=1$ implique $\sigma^2+1=0$ i.e.
$\sigma^2=-1$, et inversement, compte tenu que
$1$ et $\sigma$ sont [unclear: linéairement] indépendantes en $M_2([unclear: \mathcal{O}_S])$.
De même (37) montre
$$ (\Tr \rho=1 \text{ et } \det \rho=1) \Longleftrightarrow (\rho^2-\rho+1 = 0) \leqno{(38)} $$
et l'on a [unclear: de plus] l'implication
$$ \rho^2-\rho+1=0 \implies \rho^3+1=0 \text{ i.e. } \rho^3=-1 \leqno{(39)} $$
comme [unclear: suite de] l'identité algébrique
$$ \rho^3+1 = (\rho+1)(\rho^2-\rho+1). $$
Donc

Corollaire 2. Les conditions équivalentes de
(22) équivalent aussi à la relation
$$ \rho^2-\rho+1=0 \leqno{(40)} $$

\newpage

%% ===== Page 332 =====

5º) Théorème d'épinglage fuchsien.

\underline{Théorème} Soient $G, SG, PG$ des [unclear: foncteurs] (sous-entendus sur un anneau) de $\operatorname{Gl}(2), \operatorname{Sl}(2), \P\operatorname{Gl}(2)$ (associées), et $\tilde{\sigma}, \tilde{\rho}$ deux sections de $PG$.

5º) Un retour sur 1º), relatif à la situation
$$
\begin{cases}
1^\circ) & \tilde{\varepsilon}_0 = \tilde{\sigma}^{-1} \tilde{\rho} \text{ unipotent régulier } (\text{i.e.} \\
& \tilde{\varepsilon}_1 = \tilde{\rho} \tilde{\sigma}^{-1} = \tilde{\rho}(\tilde{\varepsilon}_0) = \tilde{\sigma}(\tilde{\varepsilon}_0) \text{ unipotent régulier) et} \\
2^\circ) & \tilde{\varepsilon}_0 \text{ et } \tilde{\varepsilon}_1 \text{ sont équivalentes ou d. gen. fib.}
\end{cases} \leqno{(41)}
$$

Soient $\varepsilon_0, \varepsilon_1$ les relèvements unipotents réguliers dans $G$ (identifiés à des sections de $\mathcal{M}$
$$
\varepsilon_0 = 1 + N_0\footnote{$N_0, N_1$ strictement nilpotents i.e. leur det et leur tr sont nuls.}, \quad \varepsilon_1 = 1 + N_1 \leqno{(42)}
$$
$$
\text{avec} \quad \tilde{p}(\varepsilon_0) = \tilde{\varepsilon}_0, \quad [unclear: \tilde{p}(\varepsilon_1) = \tilde{\varepsilon}_1]. \leqno{(43)}
$$

On peut alors relever $\tilde{\sigma}$ et $\tilde{\rho}$ de façons uniques en des sections $\sigma, \rho$ de $G$ (i.e. de $\mathcal{M}$) satisfaisant les conditions suivantes :
$$
1^\circ) \quad -\sigma^{-1} \rho = \varepsilon_0 \quad \left( \iff -\rho \sigma^{-1} = \varepsilon_1 \right) \leqno{(44)}
$$

De plus, il existe un unique triplet $(\varphi, \rho_0, \sigma_0)$
$$
\begin{cases}
\varphi : G_0 \simeq \operatorname{Gl}(2) \to G \quad (\text{épimorphisme}) \\
\rho_0 \in \Gamma(G_0), \rho_0 = \begin{pmatrix} 0 & 1 \\ -1 & \omega \end{pmatrix}, \sigma_0 = \begin{pmatrix} 0 & -1 \\ 1 & 0 \end{pmatrix} \\
(\det(\sigma_0)=1, \text{tr}(\sigma_0)=0, \omega \in \Gamma(G_m))
\end{cases} \leqno{(45)}
$$
tel que l'on ait
$$
\varphi(\rho_0) = \rho, \quad \varphi(\sigma_0) = \sigma \leqno{(46)}
$$
qui impliquent
$$
\varphi([unclear: \varepsilon_{0}]) = \varepsilon_0, \quad \varphi([unclear: \varepsilon_{1}]) = \varepsilon_1 \leqno{(47)}
$$

\newpage

%% ===== Page 333 =====

où \quad $(\varepsilon_1)_0 = -\rho_0\sigma_0^{-1} =$
$$
(\varepsilon_0)_0 = -\sigma_0^{-1}\rho_0 = \begin{pmatrix} 1 & 1 \\ 0 & 1 \end{pmatrix} , \quad \begin{pmatrix} 1 & 0 \\ 1 & 1 \end{pmatrix} \leqno{(48)}
$$

-- et inversement, pour $\rho_0, \sigma_0$ comme dessus,
et tout hom. $\varphi: G_0 = \operatorname{Gl}(2) \to G$, les images
$\tilde{\rho}, \tilde{\sigma}$ de $\rho_0, \sigma_0$ dans $PG$ satisfont (41), et les
images $\rho, \sigma$ de $\rho_0, \sigma_0$ par $\varphi$ sont des relèvements
de $\tilde{\rho}, \tilde{\sigma}$ satisfaisant (44).

Théorème d'épinglage modulaire. Soient $G, SG, PG$
des formes (sur schéma $S$) de $\operatorname{Gl}(2), \operatorname{Sl}(2), P\operatorname{Gl}(2)$
associées à une $\mathcal{O}$-alg. d'Azumaya $M$ de rg. 4.

1) Appelons \emph{triplet modulaire} de $PG$ un
triplet de sections $(\tilde{\rho}, \tilde{\sigma}, \tilde{\varepsilon}_0)$ de $PG$,
satisfaisant les conditions (49) (50) (51).

$$
\tilde{\varepsilon}_0 = \tilde{\sigma}^{-1}\tilde{\rho} \quad \text{i.e.} \quad \tilde{\rho}^{-1} \tilde{\sigma} \tilde{\varepsilon}_0 = 1 \leqno{(49)}
$$
[ de sorte que
$$
\tilde{\varepsilon}_1 \defeq \tilde{\rho} \tilde{\sigma}^{-1} = \tilde{\rho}(\tilde{\varepsilon}_0) = \tilde{\sigma}(\tilde{\varepsilon}_0) \quad ] ; \leqno{(49 \text{ bis})}
$$
$$
\tilde{\rho}^3 = \tilde{\sigma}^2 = 1, \quad \tilde{\varepsilon}_0 \text{ est unipotente régulière} \leqno{(50)}
$$
($\iff \tilde{\varepsilon}_1$ est unipotente régulière);
$$
\tilde{\varepsilon}_0 \text{ et } \tilde{\varepsilon}_1 \text{ sont [unclear: localement] équivalentes \marginpar{[unclear: analogues] dans les fibres}} \leqno{(51)}
$$
(ce qui signifie, en introduisant les relèvements
$$
\varepsilon_0 = 1 + N_0, \quad \varepsilon_1 = 1 + N_1 \quad (N_0 \text{ et } N_1 \text{ strictement}
$$
$$
\text{nilpotents ... cf (42)})
$$
dans $G$, que $N_0$ et $N_1$ sont homothétiques dans les
fibres).

2) Appelons triplet modulaire de $G$ un

\newpage

%% ===== Page 334 =====

triplet de sections $(\rho, \sigma, \varepsilon_0)$ satisfaisant les conditions suivantes

$$ -\sigma^{-1} \rho = \varepsilon_0 \qquad (\iff \rho^{-1} \sigma \varepsilon_0 = -1\ .) \leqno{(52)} $$
$$ \left[ \iff -\rho \sigma^{-1} = \varepsilon_1 \quad \text{où } \varepsilon_1 \defeq \rho \varepsilon_0 \rho^{-1} \iff -\rho \sigma^{-1} = \varepsilon_1 \quad \text{et } \varepsilon_1 = \sigma \varepsilon_0 \sigma^{-1} \right] $$

$$ (\rho + \sigma)^2 = \rho + \sigma \leqno{(53)} $$

\marginpar{$\det \varepsilon_0 = 1$ \\ $\Tr \varepsilon_0 = 2$ \\ $\varepsilon_0 \neq 1$ sur tte fibre}
$$ \varepsilon_0 \text{ unipotent régulier i.e. } \varepsilon_0 = 1 + N_0 \text{ avec} \leqno{(54)} $$
$$ N_0 \text{ nilpotent régulier } (\Tr N_0 = \det N_0 = 0, N_0 \text{ non nul sur chaque fibre}) $$
$$ \Big( \iff \varepsilon_1 \defeq \rho \varepsilon_0 \rho^{-1} = \sigma \varepsilon_0 \sigma^{-1} \text{ est unipot. rég.} \Big) $$

$$ \varepsilon_0 \text{ et } \varepsilon_1 \text{ sont [unclear: inéquivalents] en chaque fibre} \leqno{(55)} $$
$$ (\iff N_0 \text{ et } N_1 \text{ non homoth. en tte fibre}) $$

$$ \text{On a une des quatre conditions équivalentes} \leqno{(56)} $$
$$ \boxed{\sigma^2+1=0} \quad \boxed{\rho^2-\rho+1=0} $$
$$ \boxed{\Tr\sigma=0, \det\sigma=1} \quad \boxed{\Tr\rho=\det\rho=1} $$

3) Appelons triplet modulaire standard dans $\operatorname{Gl}(2)_S$ (resp. dans $\operatorname{PGL}(2)_S$) le triplet
$$ \left( \rho_0 = \begin{pmatrix} 0 & 1 \\ -1 & 1 \end{pmatrix}, \sigma_0 = \begin{pmatrix} 0 & -1 \\ 1 & 0 \end{pmatrix}, \varepsilon_0 = \begin{pmatrix} 1 & -1 \\ 0 & 1 \end{pmatrix} \right) \leqno{(57)} $$
(resp. son image $(\tilde{\rho}_0, \tilde{\sigma}_0, \tilde{\varepsilon}_0)$ dans $\operatorname{PGL}(2)_S$) -
qui est bel et bien un triplet modulaire.

Ceci posé, on a le diagramme commut. d'applications canoniques

\newpage

%% ===== Page 335 =====

(58)
$$
\begin{tikzcd}
Isom(\operatorname{Gl}(2)_S, G) \ar[r, "\varphi \mapsto (\varphi(\rho_0), \varphi(\sigma_0), \varphi(\varepsilon_0))"] \ar[d, "\simeq"'] & TripMod(G) \ar[d] & (\rho, \sigma, \varepsilon_0) \ar[d, mapsto] \\
Isom(PGl(2)_S, PG) \ar[r, "\psi \mapsto (\psi(\tilde{\rho}_0), \psi(\tilde{\sigma}_0), \psi(\tilde{\varepsilon}_0))"'] & Tripmod(PG) & (\tilde{\rho}, \tilde{\sigma}, \tilde{\varepsilon}_0)
\end{tikzcd}
$$

et on a le

Théorème d'épinglage. Toutes les flèches
de (58) sont bijectives.

Ce théorème ne fait que résumer certains
des résultats principaux des reflexions
du présent § VI. Notons que l'injectivité
des flèches horizontales signifie qu'un
hom. $\operatorname{Gl}(2)_S = G_0 \to G$ (resp. $PGl(2)_S \to PG$)
est connu quand on connait ses effets sur
$\sigma$ et $\rho$ (donc aussi sur $\varepsilon_0$) (resp. quand on
connait ses effets sur $\tilde{\sigma}$ et $\tilde{\rho}$) - ce qui
signifie aussi, [unclear: à le dire] pour les
automorphismes de $\operatorname{Gl}(2)_S$ (resp. de $PGl(2)_S$)
ce qui signifie aussi que
(59) $Z_\rho \cap Z_\sigma = Cent \ G_0$
(l'intersection des centralisateurs de $\rho, \sigma$ dans
$\operatorname{Gl}(2)_S$ est [crossed out] au centre $Cent \ G_0$) -
resp. la relation plus forte
(60) $N_\rho \cap N_\sigma = Cent \ G_0$

\newpage

%% ===== Page 336 =====

où $N'_\rho, N'_\sigma$ sont les ss-schémas en groupes de $\operatorname{Gl}(2)_{\mathbf{Z}}$ définis (resp$^t$) comme stabilisateurs de $\tilde{\rho}, \tilde{\sigma}$ dans $\operatorname{Gl}(2)_{\mathbf{Z}}$, i.e.

$$
\begin{cases}
N'_\rho = \left\{ g \mid g(\rho) = \varepsilon \rho, \, \text{$\varepsilon$ scalaire} \right\} \\
N'_\sigma = \left\{ g \mid g(\sigma) = \varepsilon \sigma, \, \text{$\varepsilon$ scalaire} \right\}
\end{cases}
\leqno{(61)}
$$
\marginpar{NB on aura automat. $\varepsilon^2 = 1$}

\underline{Rmq} Il est plausible, d'après ces résultats, que $GP(1)_{\mathbf{Z}}$ - en tant que faisceau en groupes "absolu" (sur $\operatorname{Spec} \mathbf{Z}$) pour $\operatorname{fppf}$, est engendré par les hom.
$$ \mathbf{Z}/2\mathbf{Z} \xrightarrow{1 \mapsto \tilde{\sigma}} GP(1)_{\mathbf{Z}} $$
$$ \mathbf{Z}/3\mathbf{Z} \xrightarrow{1 \mapsto \tilde{\rho}} GP(1)_{\mathbf{Z}} $$
$$ \mathbb{G}_a \xrightarrow{\lambda \mapsto \tilde{\varepsilon}_0(\lambda)} GP(1)_{\mathbf{Z}} $$
satisfaisant les seules relations
$$ \tilde{\rho}^{-1} \tilde{\sigma} \tilde{\varepsilon}_0(1) = 1 \leqno{(62)} $$
- il vaudrait bien la peine de tirer au clair cette question - ainsi que la question analogue pour $\operatorname{Gl}(2)_{\mathbf{Z}}$, où les dits groupes $\mathbf{Z}/2\mathbf{Z}$, $\mathbf{Z}/3\mathbf{Z}$ doivent se remplacer par $\mathbf{Z}/4\mathbf{Z}$, $\mathbf{Z}/6\mathbf{Z}$, et la relation (62) par
$$ \tilde{\rho}^{-1} \tilde{\sigma} \tilde{\varepsilon}_0(1) = \tilde{\rho}^3 = \tilde{\sigma}^2 \leqno{(63)} $$

\newpage

%% ===== Page 337 =====

\underline{Rmq} J'ai la nette impression que le paquet de conditions (51) à (56) est surabondant, et se réduit à peu de chose si l'on exige des conditions algébriques simples\footnote{NB. Il faut peut-être préciser $\operatorname{Tr}\sigma=0, \det\sigma=1, \operatorname{Tr}\rho=1, \det\rho=1$ pour que la relation (64) soit régulière. Si on s'en tient à la relation (64), les autres conditions s'ensuivent.}
$$
\sigma^2+1 = \rho^2-\rho+1 = 0, \quad (\rho+\sigma)^2 = \rho+\sigma \leqno{(64)}
$$
(qui définit $\varepsilon_0$ par (51)).

La troisième s'exprime en posant
$$
N = \rho+\sigma \leqno{(65)}
$$
et on a
$$
N^2=N \leqno{(66)}
$$
et en termes de $N$ on aura $\sigma^{-1}N = \sigma^{-1}\rho + 1$ i.e.
$$
\begin{cases} \varepsilon_0 = -\sigma^{-1}\rho = 1 - \sigma^{-1}N = 1+ N_0 & (N_0 = -\sigma^{-1}N = \sigma N) \\ \varepsilon_1 = -\rho\sigma^{-1} = 1 - N\sigma^{-1} = 1+ N_1 & (N_1 = -N\sigma^{-1} = N\sigma) \end{cases} \leqno{(67)}
$$
et la condition $\varepsilon_0$ strictement [unclear: régulière] s'exprime par les conditions
$$
\det N = \operatorname{Tr} N = 0, \quad \sigma N \text{ non nul sur toute fibre.} \leqno{(68)}
$$
(i.e. $N$ non nul sur toute fibre).

Or la validité de ces conditions (68) résulte de ce que, si on avait $N=0$ sur une fibre, on pourrait se ramener à $S$ spectre d'un corps, $N=0$, i.e. $\rho=-\sigma$, d'où on aurait à la fois $\sigma^2+1=0$ et $\sigma^2+\sigma+1=0$ d'où $\sigma=0$, absurde.

Je dis aussi qu'on n'aura jamais $N=1$ sur une fibre, sinon on aurait en cette fibre
$$
\rho+\sigma = 1_E
$$
et en substituant dans $\sigma^2+1=0$ on trouve
$$
(1-\rho)^2+1=0 \quad \text{i.e.} \quad \rho^2-2\rho+2=0
$$

\newpage

%% ===== Page 338 =====

comparant avec $\rho^2-\rho+1=0$, on trouve par
par différence $\rho-1=0$ i.e. $\rho=1$, mais
la relation $\rho^2-\rho+1=0$ devient $1=0$, absurde.

Donc l'opérateur idempotent $N$ est partout
de rang $1$, c.à.d. que $N$ est décomposable (un tenseur loc\textsuperscript{t}
de $M \simeq \check{F} \otimes F$) décomposable sous la forme
$x' \otimes x$, avec $\langle x, x' \rangle = 1$, et
$$
\mathrm{Tr}\, N = 1, \quad \det N = 0
$$
donc la relation $\det(\sigma^{-1}N) = 0$ de (68) est
également automatique, puisque $\det(\sigma^{-1}N) = \det \sigma^{-1} \det N = 0$. Il reste la seule relation
$$
\mathrm{Tr}\,\sigma N = 0
$$
(avec \sout{$\sigma^2+\sigma \rho$} $= \sigma \rho - 1$)
qui s'écrit aussi
$$
\mathrm{Tr}\,\sigma \rho = 2 \quad (\Leftrightarrow \mathrm{Tr}\,\sigma^{-1}\rho = -2) \leqno{(69)}
$$
laquelle il faudrait voir si elle
résulte des autres relations.

Prenons
$$
\begin{cases}
\rho = \rho_0 = \begin{pmatrix} 0 & 1 \\ -1 & 1 \end{pmatrix} \\
\sigma = \underline{u} \sigma_0 \underline{u}^{-1} = \begin{pmatrix} a & b \\ c & d \end{pmatrix} \begin{pmatrix} 0 & -1 \\ 1 & 0 \end{pmatrix} \begin{pmatrix} d & -b \\ -c & a \end{pmatrix} \\
\text{où } \underline{u} = \begin{pmatrix} a & b \\ c & d \end{pmatrix} \in \operatorname{Sl}(2) \text{, i.e. } ad-bc = 1
\end{cases} \leqno{(70)}
$$
on trouve

\newpage

%% ===== Page 339 =====

$$
\sigma = \begin{pmatrix} bd+ac & -(a^2+b^2) \\ c^2+d^2 & -(bd+ac) \end{pmatrix} \leqno{(70bis)}
$$

d'où
$$
N = \rho+\sigma = \begin{pmatrix} (bd+ac) & 1-(a^2+b^2) \\ -1+(c^2+d^2) & 1-(bd+ac) \end{pmatrix}
$$
$$
N-1 = \begin{pmatrix} (bd+ac)-1 & 1-(a^2+b^2) \\ -1+(c^2+d^2) & -(bd+ac) \end{pmatrix} = -\det(N) N^{-1} \quad \text{si } \det N \text{ inv.}
$$

et on trouve
$$
N^2 - N = N(N-1) = - \det(N) \cdot 1 \leqno{(71)}
$$

donc la condition
$$
N^2 = N \iff \det N = 0 \quad , \quad (\implies \Tr N = 1) \leqno{(72)}
$$

or on a
\begin{align*}
\det N &= bd+ac - [ -1 - (a^2+b^2)(c^2+d^2) + a^2+b^2 + c^2+d^2 ] \\
&= bd+ac - b^2d^2 - a^2c^2 - 2abcd + 1 + a^2c^2 + a^2d^2 + b^2c^2 + b^2d^2 - a^2 - b^2 - c^2 - d^2 \\
&= (bd+ac) - (a^2+b^2+c^2+d^2) + \underbrace{(ad-bc)^2}_{1} + 1 \\
&= 2 + (bd+ac) - (a^2+b^2+c^2+d^2)
\end{align*}

donc la condition (71) équivaut à
$$
a^2+b^2+c^2+d^2 = 2 + (bd+ac) \leqno{(73)}
$$

par ailleurs on a par (70), (70bis)
$$
\rho\sigma = \begin{pmatrix} c^2+d^2 & -(bd+ac) \\ c^2+d^2 - (bd+ac) & a^2+b^2 - (bd+ac) \end{pmatrix}
$$

donc
$$
\Tr \rho\sigma = \Tr \rho^{-1}\sigma = a^2+b^2+c^2+d^2 - (bd+ac) = 2 \leqno{(74)}
$$

\newpage

%% ===== Page 340 =====

Ce calcul suggère fortement que
les relations (64) impliquent déjà (68)..
c'est-à-peu de choses près une démonstration....
Si on admet ce point, il
reste seulement à vérifier que les
conditions précédentes (64) (i.e.\ constituées des
(51) à (54) et de (55')) impliquent aussi (55),
i.e.\ l'indépendance de $\varepsilon_0, \varepsilon_1$ en toute fibre,
i.e.\ que $N_0, N_1$ ne sont pas homothétiques en
toute fibre. Or (67)
$$
N_1 = \tilde{\sigma}(N_0) \leqno{(*)}
$$
Si on avait sur une fibre
$$
N_1 = \lambda N_0 \leqno{(74)}
$$
on en conclurait, en appliquant $\tilde{\sigma}$ et utilisant
$\tilde{\sigma}^2=1$, les relations
$\tilde{\sigma}(N_1) = N_0, \tilde{\sigma}(N_0) = N_1$.

Mais un calcul immédiat donne aussi
$$
N_0 = \tilde{\rho}^{-1} N_1 \leqno{(**)}
$$
$\lambda^2=1$, $\lambda^3=1$ \quad d'où $\lambda=1$

i.e.\ $N_0 = N_1$ i.e.\ $N_0$ commute à $\sigma, \rho$

donc aussi $N = -\sigma N_0$ commute à $\sigma, \rho$,

et comme $N = \sigma + \rho$, $\sigma$ et $\rho$ commutent.

Mais alors

\newpage

%% ===== Page 341 =====

$$
(\sigma+\rho)^2 = \underset{-1}{\sigma^2} + \underset{\rho-1}{\rho^2} + 2\sigma\rho = \rho - 2 + 2\sigma\rho =
$$
donc la relation
$$
(\sigma+\rho)^2 = \sigma+\rho \quad \text{s'écrit}
$$
$$
\cancel{\rho} - 2 + 2\sigma\rho = \sigma + \cancel{\rho} \quad \text{i.e.}
$$
$$
2\sigma\rho = \sigma + 2
$$
$$
\text{soit} \quad \underbrace{2\sigma^2}_{-2}\rho = \underbrace{\sigma^2}_{-1} + 2\sigma \quad \text{i.e.} \quad 2(\sigma\rho) = 1
$$

Prenant les traces des deux membres, et utilisant $\Tr \rho = 2$ ([unclear: cf 4]), $\Tr \sigma = 0$, $\Tr 1 = 2$, on trouve $4 = 4$, pas de contradiction, mais prenant les déterminants (en tenant compte, ce qui est une conséquence de (64) je pense, que $\sigma, \rho$ soient réguliers i.e. ne soient des scalaires sur aucune fibre) on trouve
$$
\cancel{4} = \underset{4}{2^2} + \underset{0}{2 \Tr \sigma} + \underset{1}{(\det \sigma)} \quad \text{i.e.} \quad \cancel{1} = 0 \leqno{(**)}
$$
ce qui est absurde.

De $(*)$ $2N=1$, en prenant les déterminants
$$
4 \det N = 1 \quad \text{i.e.} \quad 0 = 1, \quad \text{absurde.}
$$

Je ne serais pas étonné qu'on ait aussi des conditions équivalentes à (51), (56) sous la forme
$$
\sigma^2 + 1 = \rho^2 - \rho + 1 = 0, \quad \varepsilon_0 \defeq - \sigma \rho \text{ est un projecteur régulier} \leqno{(75)}
$$

\newpage

%% ===== Page 342 =====

en exigeant de plus, au besoin, que $\sigma, \rho$ soient
également réguliers i.e.\ ne soient des scalaires
sur aucune fibre. Il faudrait se
convaincre que cela implique que
$$N \defeq \rho + \sigma$$
satisfait bien la relation d'idempotence
$$N^2 - N = 0 \leqno{(75)}$$
(ce qui équivaut, si $N$ régulier, à
$$\operatorname{Tr} N = 1, \det N = 0 \quad (?) \leqno{(77)}$$
Mais $N$ est régulier, car si sur une fibre
il était scalaire on aurait que (car
se restreignant à la fibre) $\rho$ et $\sigma$
commutent, on aurait
$$-\sigma^{-1}\rho = \sigma\rho = \varepsilon_0 = 1 + N_0$$
avec
$$(1 - \sigma\rho)^2 = 0 \quad \text{i.e.} \quad 1 - 2\sigma\rho + \underbrace{\sigma^2}_{-1}\underbrace{\rho^2}_{\rho-1} = 0 \quad \text{i.e.}$$
$$1 - 2\sigma\rho + 1 - \rho = 0 \quad \text{i.e.} \quad 2 = \rho + 2\sigma\rho$$
soit, en multipliant par $\sigma$
$$2\sigma = \rho\sigma + 2\underbrace{\sigma^2}_{-1}\rho \quad \text{i.e.} \quad 2(\underbrace{\sigma + \rho}_{N}) = \rho\sigma$$

Mais par hypothèse on a $N = \lambda \cdot 1$, $\lambda$ scalaire,
d'où $\rho\sigma = 2\lambda$, et prenant les
déterminants
$$\det(\rho\sigma) = 4\lambda^2$$
i.e.
$$1 = 4\lambda^2 \quad \text{i.e.} \quad \lambda = \frac{\varepsilon}{2}, \text{ avec } \varepsilon^2 = 1,$$
i.e.
$$\rho\sigma = \varepsilon \qquad (\varepsilon^2 = 1)$$
$$\rho^2\sigma^2 = 1 \qquad \underbrace{\rho^2}_{\rho-1} = -1, \quad \text{i.e.} \quad \rho = 0, \text{ absurde.}$$

\newpage

%% ===== Page 343 =====

Mais $\operatorname{Tr} N = 0$ [unclear: rayé]

Donc $N$ est régulière, il reste à prouver (77).

Mais $\operatorname{Tr} N = \underset{0}{\operatorname{Tr} \sigma} + \underset{1}{\operatorname{Tr} \rho} = 1$, OK, reste

à prouver $\det(\rho+\sigma) = 0$ mais
$$
\det(\rho+\sigma) = \underset{1}{\det \rho} + \underset{1}{\det \sigma} + \operatorname{Tr}(\sigma \rho) \quad \text{est la relation}
$$
$$
\det(\sigma(1+\sigma^{-1}\rho)) = \det \sigma \cdot \det(1+\sigma^{-1}\rho) = \underset{1}{\det \sigma} \cdot (\underset{\det \sigma^{-1} \det \rho = 1}{\det \sigma^{-1}\rho} + \operatorname{Tr}(\sigma^{-1}\rho) + 1)
$$
$$
\det(\rho+\sigma)\footnote{calcul déjà fait ! i.e. $\det(\rho+\sigma) \dots$ (cf formule (69))} = 2 + \operatorname{Tr}(\sigma^{-1}\rho) = 2 - \operatorname{Tr} \sigma \rho , \quad \text{donc}
$$
la condition $\det N = 0$ équivaut à :
$$
\operatorname{Tr} \sigma \rho = 2 \leqno{(69)}
$$
(c'est au moins testé !)

Pour la prouver (cette relation) on peut songer à procéder comme tantôt, où je suis parti de (68) au lieu de (75). Ce n'est pas entièrement élucidé - à revoir : le danger !

\newpage

%% ===== Page 344 =====

\section*{VII. Récapitulations sur les épinglages.}

On désigne par $G, SG, PG$ des formes de $\operatorname{Gl}(2)$, $\operatorname{Sl}(2)$, $P\operatorname{Gl}(2)$ associés à une $\mathcal{O}_S$-Algèbre d'Azumaya $\mathcal{M}$, i.e.\ à un fibré en droites projectives $P$ sur $S$.

1) \underline{Sections régulières de $G$}. (pour mémoire)

Soit $\rho$ une section de $\mathcal{M}$, posons
$$
v = \det \rho, \quad w = \mathrm{Tr} \rho \leqno{(1)}
$$
de sorte qu'on aura l'équation d'Hamilton-Cayley
$$
\rho^2 - w\rho + v = 0 \qquad \text{i.e.} \quad \rho^2 = w\rho - v \leqno{(2)}
$$

\underline{Proposition}. Conditions équivalentes :
\begin{enumerate}
    \item[a)] $\rho$ n'est un scalaire sur aucune fibre.
    \item[a')] \marginpar{NB \`A cause de (2), $A_\rho$ est une sous-Alg\`ebre de $\mathcal{M}$.} $A_\rho = \underline{1}_S + \mathcal{O}_S \rho$ est un sous-module localement facteur direct libre de rang 2 de $\mathcal{M}$.
    \item[b)] (Si $\mathcal{M} = \underline{\mathrm{End}}_{\mathcal{O}_S}(F)$) Il existe loc. (Zar) une section $e_0$ de $F$ telle que $(e_0, \rho e_0 \defeq \rho(e_0))$ forment une base de $F$.
    \item[b')] (Si $\mathcal{M} = \underline{\mathrm{End}}_{\mathcal{O}_S}(F)$) Il existe loc. Zar un ``sous-fibré en droites'' $D_0 \subset F$, tel que $F = D_0 \oplus D_1$, où $D_1 = \rho(D_0)$.
    \item[c)] (Si $\mathcal{M} = \underline{\mathrm{End}}_{\mathcal{O}_S}(F)$) Le $\mathcal{O}_S[T]$-module $F_\rho$ ($T$ [unclear: désigne] une indéterminée, qui opère sur $F$ par $\rho$) est [unclear: monogène] loc! (Zar).
    \item[c')] \marginpar{(car $\mathcal{M} = \underline{\mathrm{End}}(F)$)} (Si $\mathcal{M} = \underline{\mathrm{End}}_{\mathcal{O}_S}(F)$) Le $\mathcal{O}_S[T]$-module $F_\rho$ est isomorphe à $\mathcal{O}_S[T] / (T^2 - w T + v)$
    \item[c'')] Il existe loc Zar une base $e_0, e_1$ de $F$, telle que la matrice de $\rho$ p.r.\ à cette base...
\end{enumerate}

\newpage

%% ===== Page 345 =====

soit
$$
\rho_{e_0, e_1} = \begin{pmatrix} 0 & -v \\ 1 & w \end{pmatrix} \leqno (3)
$$
(i.e.\ $\rho e_0 = e_1$, d'où $\rho e_1 = \rho^2 e_0 = -v e_0 + w e_1$)

(Cas $\mathcal{M} = \underline{\mathrm{End}}(\mathcal{F})$)

c'') $\exists$ loc.\ Zar.\ une base $(e_0, e_1)$ de $\mathcal{F}$ telle que la matrice de $\rho$ p.r.\ à celle-ci soit
$$
\rho_{e_0, e_1} = \begin{pmatrix} 0 & 1 \\ -v & w \end{pmatrix} \leqno (4)
$$

d) Le commutant de $\rho$ dans $\mathcal{M}$ est réduit à : $A_\rho = \underline{\mathcal{O}}_S + \underline{\mathcal{O}}_S \rho$, et il en reste ainsi après chaque changement de base $S' \to S$.

d') Idem sur chaque fibre.

Dem : $a \iff a'$ trivial (algèbre linéaire des modules $\mathcal{M}$ et de son sous-module [rayé: illisible] loc.\ facteur direct)

$b \iff b' \iff b''$ trivial, $b) \implies a) \implies a'$

$a \implies b$ il suffit de choisir $e_0$ dans une fibre de façon à ne pas être vecteur propre pour $\rho$ (possible sur un corps, pour un op.\ non scalaire) puis de relever au voisinage.

$c) \iff c')$ trivial, les bases pour lesquelles $\rho_{e_0,e_1}$ soit de la forme $\begin{pmatrix} 0 & -v \\ 1 & w \end{pmatrix}$ étant justement les bases de la forme $(e_0, e_1 = \rho e_0)$.

$c') \iff c'')$ signifie que si $\rho$ satisfait aux conditions envisagées, ${}^t\rho$ aussi. [rayé: Et] et inversement il en résulte que la condition pour ${}^t\rho$ l'implique pour $\rho = {}^t({}^t\rho)$, i.e.\ $c'') \implies c')$ -- mais c'est immédiat sans ce [unclear: fourbi] --).

Ainsi les conditions a) à c'') sont équivalentes. Par ailleurs $d) \implies d') \implies a)$ est trivial, reste à prouver que $b \text{ (p.\ ex.)} \implies d)$. Comme $A_\rho$ et le Commutant se conservent par changement de base, il suffit de voir que

\newpage

%% ===== Page 346 =====

974

\noindent
NB En fait, $\underline{G} = \underline{M}^\times$ opère sur $M$ (par autom.
intérieurs) ou sur $\underline{G}$, $S\underline{G}$, via $P\underline{G}$
$= \underline{M}^\times / \mathbf{G}_m$, et $P\underline{G}$ est quotient [unclear: fppf]
de $S\underline{G}$ (et $=$ pour la top. ét., si 2 inv.).
Donc on voit que (la notion de conjugaison
précédente, dans le [unclear: formalisme])
$$
\rho' = \operatorname{int}(\alpha) \rho = \alpha \rho \alpha^{-1} \leqno (6)
$$
où $\alpha \in \underline{G}(S)$.

\noindent
Cor. 2 Soient $\rho, \xi$ deux sections de $\underline{G}$,
$\rho$ étant régulière. Pour qu'on puisse
trouver loc. une $\alpha \in P\underline{G}$ telle que
$$
(\alpha, \rho) = \xi \quad \text{i.e. } \alpha \rho \alpha^{-1} = \xi \rho \leqno (7)
$$
il faut et il suffit qu'on ait
\begin{enumerate}
\item[a)] $\xi \rho$ est régulière, i.e. [unclear: n'admettant pas la] valeur propre
$= \rho^{-1}$ sur aucune fibre
\item[b)] $\det \xi = 1$, $\operatorname{Tr} \xi \rho = \operatorname{Tr} \rho$
\end{enumerate}
(On peut alors prendre $\alpha$ dans $S\underline{G}$...)
Comme [unclear: par] les conditions du cor. 1 [unclear: s'agissant]
de conjugaison de $\xi \rho$ à $\rho$, la condition
$\det \xi \rho = \det \rho$ s'écrit aussi $\det \xi \det \rho = \det \rho$
i.e. $\det \xi = 1$.

\noindent
Corollaire 3 L'indétermination de $\alpha$ dans la
résolution de (7) est de la forme $\alpha \mapsto \alpha \underline{t}$,
où $\underline{t}$ section de $A_\rho^\times$, i.e. $\underline{t} = c \rho + d$, avec
$c, d \in \Gamma(\mathcal{O}_S)$, $\det(c \rho + d) \in \Gamma(\mathcal{O}_S^\times)$, i.e. $c^2 v + w c d + d^2$ inv.
(où $v = \det \rho$, $w = \operatorname{Tr} \rho$, cf (1)).

\newpage

%% ===== Page 347 =====

si $u$ est une section de $\mathcal{M}$ telle qui commute
à $\rho$, alors $u = a + b \rho$ ($a, b \in \Gamma_{\mathcal{O}_S}$). Par a) $a$ et $b$
sont uniques, et cela montre que la [unclear: question est]
locale fppf. On peut supposer alors $\mathcal{M} = \End(F)$, $F$ muni d'une
base $(e_0, e_1 = \rho(e_0))$, on a, si
$u(e_0) = a e_0 + b e_1$, en posant $u' = a + b\rho$,
et si $u$ commute à $\rho$, on a
$u(e_1) = u\rho(e_0) = \rho u(e_0)$
$u'(e_1) = u'\rho(e_0) = \rho u'(e_0)$ car $u'$ commute à $\rho$
donc $u(e_0) = u'(e_0), u(e_1) = u'(e_1)$ donc $u = u'$,
cqfd.

Définition. Une section de $\mathcal{M}$ qui satisfait aux conditions
équivalentes de la prop. est dite \underline{régulière}. Idem
pour sections de $G, \mathcal{S}G$, interprétées comme
sections de $\mathcal{M}$ -- une telle section
est régulière ssi elle n'est
pas un scalaire, i.e. n'est pas dans le centre. Ceci est
stable par multiplication par un scalaire inver-
sible - donc par passage au quotient on déduit
la notion de section régulière de $P G$ -
ce sont les sections qui ne sont $= 1$ en
aucune fibre.

Corollaire 1. Soient $\rho, \rho'$ deux sections régulières de
$\mathcal{M}$. Pour qu'elles soient conjuguées (loc. Zar)
via $\mathcal{M}^*$, il faut et suffit que
$$
\det \rho = \det \rho' \quad , \quad \Tr \rho = \Tr \rho' \text{ (conditions numériques)} \leqno{(5)}
$$

\newpage

%% ===== Page 348 =====

Ceci résulte de l'équivalence a) $\iff$ d) de la prop.\ préc.

\underline{Proposition} Soit
$$ \varepsilon = 1 + N \leqno{(9)} $$
une section de $\mathcal{M}$\footnote{NB $\det(1+N) = 1 + \operatorname{Tr} N + \det N$, $\operatorname{Tr}(1+N) = 2 + \operatorname{Tr} N$}. Conditions équivalentes :
\begin{enumerate}
\item[a)]
$$
\begin{cases}
\det \varepsilon = \det 1 \quad (=1) \\
\operatorname{Tr} \varepsilon = \operatorname{Tr} 1 \quad (=2)
\end{cases}
\qquad
\begin{matrix}
\text{$\varepsilon$ unipotente} \\
\text{et $N$ nilpotente}
\end{matrix}
\leqno{(10)}
$$
\item[b)]
$$ \det N = \operatorname{Tr} N = 0 \leqno{(11)} $$
\item[c)] (Si $\mathcal{M} = \underline{\operatorname{End}}_{\mathcal{O}_S}(F)$) $\exists$ fibré en droites $D \subset F$ tel que l'on ait
$$ N(F) = D, \quad N(D) = 0 $$
(alors $D$ est uniquement déterminé par $\varepsilon$ i.e.\ par $N$, $D = \Ker N = \Im N$).
\item[c')] $\exists$ base $e_1, e_2$ de $F$ tel que $\varepsilon = \begin{pmatrix} 1 & 1 \\ 0 & 1 \end{pmatrix}$ i.e.\ $N = \begin{pmatrix} 0 & 1 \\ 0 & 0 \end{pmatrix}$.
\item[d)] $\exists$ Borel $B \subset G$, tel que $\varepsilon$ soit une section partout $\neq 1$ de $B_u$.
\item[d')] $\exists$ Borel $SB \subset \SL_2$, tel que $\varepsilon$ soit une section partout $\neq 1$ de $(SB)_u$.
\end{enumerate}

De plus, $D, B, SB$ sont bien déterminés par $\varepsilon$ i.e.\ par $N$, et déterminent $N$ à homothétie près.

\underline{Définition} : on dit que $\varepsilon$ est (strictement) unipotente si on a $\det \varepsilon = 1, \operatorname{Tr} \varepsilon = 2$, i.e.\ que $N$ est strictement nilpotente $\iff \det N = \operatorname{Tr} N = 0$ (d'où $N^2 = 0$ par Hamilton-Cayley). On dit que $\varepsilon$ est unipotente régulière resp.\ nilpotente régulière (on laisse tomber le qualificatif "strictement") si de plus il est régulier.

\newpage

%% ===== Page 349 =====

Proposition Soit $P$ une section de $\mathcal{M}$. Conditions équivalentes :

\begin{enumerate}
\item[a)] $P$ est régulière, et $P^2 = P$
\item[b)] $\det P = 0$, $\Tr P = 1$
\item[c)] (Si $\mathcal{M} = \underline{End}_{\mathcal{O}_S}(F)$) $\exists$ base $e_0, e_1$ de $F$ telle que la matrice de $P$ est $P_{e_0,e_1} = \begin{pmatrix} 0 & 0 \\ 0 & 1 \end{pmatrix}$ i.e.\ $P$ est la projection de $\mathcal{O}_S e_0 \oplus \mathcal{O}_S e_1$ sur le facteur $\mathcal{O}_S e_1$.
\item[d)] $\exists$ un isomorphisme $\mathcal{M} \isommap M_2(\mathcal{O}_S)$, transformant $\begin{pmatrix} 0 & 0 \\ 0 & 1 \end{pmatrix}$ en $P$.
\item[e)] [unclear: (Si $\mathcal{M} = \underline{End}_{\mathcal{O}_S}(F)$) $P$ est un projecteur dont le noyau est un sous-module de rang $1$.]
\end{enumerate}

Remarque. Pour la démonstration, on notera que la condition b) implique déjà que $P$ est régulière, car si $P$ était un scalaire $\lambda$ (dans un [unclear: anneau local]) on aurait $\det P = \lambda^2 = 0$, $\Tr P = 2\lambda = 1$, donc $\lambda$ serait inversible et aurait un carré nul, absurde. Noter que cet argument montre, plus généralement, que si $p$ est une section de $\mathcal{M}$ telle que
$$ \delta(p) = (\Tr p)^2 - 4 \det p = u^2 - 4v \leqno{(12)} $$
\marginpar{(= discriminant du polynôme caractéristique de Hamilton-Cayley de $p$)}
soit une section inversible de $\mathcal{O}_S$, alors $p$ est une section régulière de $\mathcal{M}$.

Définition. Les $P$ satisfaisant les conditions de la proposition précédente sont appelés les projecteurs élémentaires de $\mathcal{M}$.

Corollaire. Il y a une correspondance biunivoque entre les projecteurs élémentaires de $\mathcal{M}$, et les couples de Borel $(B,T)$ dans $G$ (i.e.\ $SB, ST$ dans $S_G$, i.e.\ $PB, PT$ dans $PG$), en [unclear: associant]...

\newpage

%% ===== Page 350 =====

associant à $P$ le tore $T = A_P^\ast$, et le Borel $B_0$ formé des $g$ qui stabilisent le sous-fibré $D = \Im P$ (avec le sous-groupe unipotent fermé $N$ des $g$ qui induisent $1$ sur $D$ et sur $F/D \simeq D_{1-P}$),
$$
T = \left\{ \begin{pmatrix} \mu_0 & 0 \\ 0 & \mu_1 \end{pmatrix} \mid \mu_0, \mu_1 \in \Gamma \mathbb{G}_m \right\} \leqno{(13)}
$$
$$
B_0 = \left\{ \begin{pmatrix} \mu_0 & \lambda \\ 0 & \mu_1 \end{pmatrix} \mid \mu_0, \mu_1 \in \Gamma \mathbb{G}_m, \lambda \in \Gamma \mathbb{G}_a \right\}
$$
Le couple de Killing opposé à $(T, B_0)$ est $(T, B_1)$, avec
$$
B_1 = \left\{ \begin{pmatrix} \mu_0 & 0 \\ \lambda & \mu_1 \end{pmatrix} \mid \mu_0, \mu_1 \in \Gamma \mathbb{G}_m, \lambda \in \Gamma \mathbb{G}_a \right\} \leqno{(14)}
$$
De façon un peu plus intrinsèque, $B_1$ est le stabilisateur de $D_{1-P}$. Le projecteur élémentaire associé au couple de Killing opposé $(T, B_1)$\marginpar{ou au sous-groupe de Borel $B_1$} est $1-P$.

Remarque. Le schéma des épinglages (i.e.\ un $\mathbb{G}_m$-torseur) sur le schéma des Borels $\mathcal{B}$ (de $G$, de $\mathcal{S}G$ ou de $\mathcal{P}G$, cela revient au $=$), le schéma des projecteurs élémentaires $\mathcal{M}$ est un fibré lisse de dimension relative 2. Celui des couples de Killing, fibré lisse de dimension relative 1 sur le schéma des Borels. Le produit fibré de ces deux schémas au dessus du $\mathcal{B} =$ schéma des Borels, est isomorphe au schéma des épinglages de $G$ (ou de $\mathcal{S}G$, ou de $\mathcal{P}G$) (i.e.\ des isomorphismes $\mathcal{G}_{éping} \isom \operatorname{Sl}(2)_S \hookrightarrow \mathcal{G}$).

\newpage

%% ===== Page 351 =====

(resp! $G P U_{\bar{S}} \isom P G$), qui est un torseur sous $\operatorname{Aut}_{S-\text{gr}} \operatorname{Gl}(W_S) \simeq \operatorname{Aut}_{S-\text{gr}} \operatorname{Sl}(W_S) \simeq \operatorname{Aut}_{S-\text{gr}}(G)$. On dit que $\varepsilon$ et $P$ (ou $N = \varepsilon - 1$ et $P$ associé) sont \underline{associés} s'ils définissent le \underline{même} Borel, i.e. \underline{proviennent} d'un \underline{même} épinglage.

On vérifie que ceci équivaut aux relations suivantes :
$$
P N = 0 \quad , \quad N P = N \leqno{(15)}
$$
$$
P \varepsilon = P \quad , \quad \varepsilon P = P + \underbrace{\varepsilon - 1}_{N} \leqno{(16)}
$$
\marginpar{[unclear: c'est un peu plat!]}

Nous considérons maintenant la question d'équivalence (conjugaison) de sections $\tilde{\rho}, \tilde{\rho}'$ de $P G$, définies par des sections $\rho, \rho'$ de $G$. On suppose $\tilde{\rho}$ \underline{régulière} i.e. $\neq 1$ ou ce qui revient au même $\rho$ \underline{régulière}. La pseudo-conjugaison de $\rho, \rho'$ par une section $u$ de $G$ (disons) signifie qu'on a
$$
\rho' = \lambda u \rho u^{-1} \quad , \quad \lambda \in \Gamma G_m \leqno{(17)}
$$
d'où on déduit
$$
\det \rho' = \lambda^2 \det \rho \quad \operatorname{Tr} \rho' = \lambda \operatorname{Tr} \rho
$$
et par suite
$$
\frac{\operatorname{Tr}^2 \rho'}{\det \rho'} = \frac{\operatorname{Tr}^2 \rho}{\det \rho} \leqno{(18)}
$$
Posons, pour une section $\rho$ de $G$
$$
\Delta(\rho) \defeq \frac{\operatorname{Tr}^2 \rho}{\det \rho} \leqno{(19)}
$$

\newpage

%% ===== Page 352 =====

- fait ce qu'il fallait pour que dans $G \to E'$ que passe au quotient en
$$
PG \xrightarrow{\Delta} E' \leqno{(20)}
$$
$$
\Delta(\tilde{\rho}) = \Delta(\rho) = \frac{\operatorname{Tr}^2 \rho}{\det \rho}
$$
et que soit invariant par automorphismes intérieurs.

Proposition. Pour que deux sections $\tilde{\rho}, \tilde{\rho}'$ de $PG$ soient conjuguées loc. f.p.p.f.\footnote{$\tilde{\rho}$ et $\tilde{\rho}'$ régulières}, il faut et il suffit que $\tilde{\rho}'$ soit régulière, et qu'on ait
$$
\Delta(\tilde{\rho}) = \Delta(\tilde{\rho}') \leqno{(21)}
$$
ici en termes de $\rho, \rho'$ relevant $\tilde{\rho}, \tilde{\rho}'$, que $\rho'$ soit régulière et satisfasse (18). La condition est satisfaite quant $\Delta(\tilde{\rho})$ inv. i.e. $\operatorname{Tr} \rho$ inv.

Il reste à prouver la suffisance. Notons que loc. f.p.p.f., existe $\lambda \in [unclear: \mathbf{G}_m]$ tel que
$$
\det(\rho') = \det(\lambda \rho) \quad \text{i.e.} \quad \lambda^2 = \frac{\det \rho'}{\det \rho}
$$
On peut (quitte à remplacer $\rho$ par $\lambda \rho$) supposer $\det \rho = \det \rho'$. Alors (18) devient
$$
\operatorname{Tr}^2 \rho' = \operatorname{Tr}^2 \rho \leqno{(22)}
$$
Cela la question revient à voir que $\exists \varepsilon \in \Gamma(\mu_2)$, telle que
$$
\operatorname{Tr} \rho' = \varepsilon \operatorname{Tr} \rho \quad (= \operatorname{Tr} \varepsilon \rho)
$$
ainsi
$$
\det \rho' = \det \varepsilon \rho \quad (= \varepsilon^2 \det \rho = \det \rho)
$$
d'où loc. conjugaison de $\rho'$ et $\varepsilon \rho$, [unclear: terminé].

Mais si $\tau', \tau$ sont deux sections de $G_S$, [unclear: avec $\tau \simeq \tau'$], alors que $\tau'^2 = \tau^2$, on aura bien $\tau' = \varepsilon \tau$ ($\varepsilon = \tau' / \tau$), d'où $\varepsilon^2 = 1$, OK.

\newpage

%% ===== Page 353 =====

Rmq. Si $\Delta(\rho)$ [unclear: non] inv. il y a des contre-exemples évidents, p. ex. avec $\operatorname{Tr}\rho = 0$, $(\operatorname{Tr}\rho')^2 = 0$, $\det\rho = \det\rho' = 1$, $\operatorname{Tr}\rho$ et $\operatorname{Tr}\rho'$ n'étant pas homothétiques. On peut prendre
$$ \rho = \begin{pmatrix} 0 & 1 \\ -1 & w \end{pmatrix}, \quad \rho' = \begin{pmatrix} 0 & 1 \\ -1 & w' \end{pmatrix} $$
$w^2 = w'^2 = 0$, $w,w'$ non homothétiques, p.ex. [unclear: sur] un anneau de base $A = A_0[w, w']/(w^2, ww', w'^2)$ ([unclear: le radical de $A$ étant un] module libre de rg $2$ engendré par $w, w'$).

2) Position relative de deux sections $\rho, \sigma$
de $G$, i.e. $\tilde{\rho}, \tilde{\sigma}$ de deux sections de $PG$.

On s'intéresse dans $PG$ à :
$$ \tilde{\varepsilon}_0 = \tilde{\sigma}^{-1} \tilde{\rho} \quad \text{et} \quad \tilde{\varepsilon}_1 = \tilde{\rho} \tilde{\sigma}^{-1} = \tilde{\rho}(\tilde{\varepsilon}_0) = \tilde{\sigma}(\tilde{\varepsilon}_0) \leqno{(23)} $$
et à la condition que $\tilde{\varepsilon}_0$ (donc $\tilde{\varepsilon}_1$)\marginpar{\textcircled{10}} soit unipotent régulier --- ce qui proviendra canoniquement de $\varepsilon_0, \varepsilon_1$ unipotents réguliers de $\operatorname{Sl}(2)$ satisfaisant à\marginpar{[unclear: d'après les conditions imposées]}
$$ \varepsilon_1 = \tilde{\rho}(\varepsilon_0) = \tilde{\sigma}(\varepsilon_0) \leqno{(24)} $$
$$ \begin{cases} \varepsilon_0 = (\text{scalaire inv.}) \sigma^{-1} \rho & \text{soit } \varepsilon_0 = \pm \sigma^{-1} \rho \\ \varepsilon_1 = (\text{  --- " ---  }) \rho \sigma^{-1} & \varepsilon_1 = \pm \rho \sigma^{-1} \end{cases} \leqno{(25)} $$
et par (20) $\tilde{\varepsilon}_0, \tilde{\varepsilon}_1$ soient [unclear: par suite] "unipotents réguliers", i.e. $\varepsilon_0, \varepsilon_1$ [unclear: l'étant]. Posons
$$ \varepsilon_0 = 1 + N_0, \quad \varepsilon_1 = 1 + N_1 \leqno{(26)} $$

\newpage

%% ===== Page 354 =====

que $N_0$ et $N_1$ soient [unclear: partout non] homothétiques.
Alors $\tilde{\varepsilon}_0(G_a)$, $\tilde{\varepsilon}_1(G_a)$, $\tilde{\varepsilon}_\infty(G_a)$ définissent
On aura alors, pour cet épinglage,

$$
\rho = \begin{pmatrix} 0 & u \\ v & w \end{pmatrix} \qquad \sigma = \begin{pmatrix} 0 & u' \\ v' & w' \end{pmatrix} \leqno{(26)}
$$

avec $u, v, u', v' \in \Gamma(G_m)$, $w, w' \in \Gamma G_a$,

et on aura

$$
\begin{cases}
- \sigma^{-1} \rho = \begin{pmatrix} -\frac{v}{v'} & \frac{u w' - w u'}{u' v'} \\ 0 & -\frac{u}{u'} \end{pmatrix} = \varepsilon_0 = \begin{pmatrix} 1 & [unclear: \phi] \\ 0 & 1 \end{pmatrix} \\
- \rho \sigma^{-1} = \begin{pmatrix} -\frac{u'}{u} & 0 \\ \frac{v u' - w' v}{u v} & -\frac{v'}{v} \end{pmatrix} = \varepsilon_1 = \begin{pmatrix} 1 & 0 \\ [unclear: \phi'] & 1 \end{pmatrix}
\end{cases} \leqno{(27)}
$$
\marginpar{bien définies au choix d'un épinglage}

donc il faudra

$$
u/u' = -1, \quad v/v' = -1 \quad \text{i.e.} \quad u' = -u, \quad v' = -v \leqno{(28)}
$$

donc
$$
- \rho \sigma^{-1} = \begin{pmatrix} 1 & 0 \\ \frac{1}{uv}(w+w') & 1 \end{pmatrix}
$$
et il faut

$$
w' + w = u \quad \text{i.e.} \quad w' = u - w \leqno{(29)}
$$

en d'autres termes (ré-énonçant (26)(27))

$$
\rho + \sigma = u P \qquad \text{où} \quad P = \begin{pmatrix} 0 & 0 \\ 0 & 1 \end{pmatrix} \leqno{(30)}
$$

ce qui détermine $\sigma$ en termes de $\rho$, et de la projection $P$.
où $P$ est la projection évidente associée à l'épinglage envisagé.

$$
-\sigma^{-1} \rho = \varepsilon_0 = \begin{pmatrix} 1 & u/v \\ 0 & 1 \end{pmatrix} \qquad N_0 = \begin{pmatrix} 0 & u/v \\ 0 & 0 \end{pmatrix} \leqno{(31)}
$$

\newpage

%% ===== Page 355 =====

On aimerait donner des conditions numériques ``intrinsèques'' pour que $\rho,\sigma$ aient les propriétés précédentes...

Proposition. Soient $\rho,\sigma$ deux matrices dans $\bar{G}$,
$$
\varepsilon_0 = - \sigma^{-1} \rho \quad , \quad \varepsilon_1 = - \rho \sigma^{-1}
$$
de sorte qu'on ait donc
$$
\varepsilon_1 = \sigma(\varepsilon_0) = \rho(\varepsilon_0).
$$

Conditions équivalentes :

\begin{enumerate}
\item[a)] $\varepsilon_0$ (et $\varepsilon_1$) sont unipotentes régulières, et, en posant :
$$
\varepsilon_0 = 1 + N_0, \quad \varepsilon_1 = 1 + N_1
$$
$N_0$ (resp. $N_1$) est nilpotente régulière, et $N_0$ et $N_1$ [unclear: ne sont pas homothétiques].

\item[b)] \marginpar{cf prop. p. 147} On a :
$$
\det \rho = \det \sigma \leqno{(32)}
$$
$$
\operatorname{Tr} \sigma^{-1} \rho = -2 \footnote{car ces relations expriment $\det(\varepsilon_0)=1, \operatorname{Tr} \varepsilon_0 = 2$ i.e. $\varepsilon_0$ est unipotente (et $\varepsilon_0 \neq 1$ car $\rho \neq -\sigma$ par loc. cit.)} \leqno{(33)}
$$
$$
u \defeq \operatorname{Tr}(\sigma+\rho) \text{ est inversible} \leqno{(34)}
$$

\item[c)] On a :
$$
\det(\sigma+\rho) = 0 \leqno{(32')}
$$
$$
\operatorname{Tr}(\sigma+\rho) \text{ inversible} \leqno{(33')}
$$

\item[d)] On a $\det \sigma = \det \rho$, et il existe un projecteur $P$ et un scalaire $u \in \Gamma$\footnote{NB. [unclear: Évidemment pour $u$] inversible, on a $P = \frac{1}{u}(\sigma+\rho)$} tels qu'on ait
$$
\sigma+\rho = u P \leqno{(34')}
$$
\end{enumerate}

\newpage

%% ===== Page 356 =====

\begin{enumerate}
\item[e)] $\exists$ isomorphisme $\operatorname{Gl}(2) \isommap G$ \footnote{NB cet isomorphisme est unique}, tel que l'on ait
$$
\begin{cases} 
\varphi^{-1}(\varepsilon_1) = \begin{pmatrix} 1 & 0 \\ 1 & 1 \end{pmatrix} \\ 
\varphi^{-1}(\varepsilon_0) = \begin{pmatrix} 1 & 1 \\ 0 & 1 \end{pmatrix} 
\end{cases} \leqno{(35)}
$$

\item[f)] $\exists$ isomorphisme $\operatorname{Gl}(2) \isommap G$ \footnote{cet isomorphisme $\varphi$ est unique; avec $u = \operatorname{Tr}(\sigma+\rho) (= -u')$, $w = \operatorname{Tr} \rho$ ($w' = \operatorname{Tr} \sigma$), $-uv = \det \rho = \det \sigma$ ce qui dicte $u, v, u', v'$ en termes de $\rho, \sigma$ \dots $u, v, u', v' \dots$} tel que l'on ait
$$
\begin{cases} 
\varphi^{-1}(\rho) = \begin{pmatrix} 0 & u \\ v & w \end{pmatrix} \\ 
\varphi^{-1}(\sigma) = \begin{pmatrix} 0 & u' \\ v' & w' \end{pmatrix} 
\end{cases} \leqno{(36)}
$$
avec $u+u' = v+v' = 0$, $w+w'=u$ \marginpar{NB On a $P = \begin{pmatrix} 0 & 0 \\ 0 & 1 \end{pmatrix}$}
\end{enumerate}

Démonstration facile. L'équivalence a) $\Leftrightarrow$ b) est dans [unclear: prop. p. 449]. Pour b) $\Leftrightarrow$ c) , on écrit
$$
\det(\sigma + \rho) = \det(\sigma (1 + \sigma^{-1}\rho)) = \det \sigma \underset{1 + \operatorname{Tr} \sigma^{-1}\rho + \det \sigma^{-1}\rho}{\det(1 + \sigma^{-1}\rho)}
$$
$$
= \det \sigma + \det \rho + (\det \sigma)\operatorname{Tr}(\sigma^{-1}\rho) \quad \text{(car } \det \rho = \det \sigma \det \sigma^{-1}\rho \text{)}
$$
donc en vertu de l'hypothèse
$$ \det \sigma = \det \rho $$
commune à b) et c) , on a
$$
\det(\sigma + \rho) = \det \sigma (2 + \operatorname{Tr}(\sigma^{-1}\rho)) \leqno{(40)}
$$
donc la première n'est autre que $\operatorname{Tr} \sigma^{-1}\rho = -2$, q.f.d.

\newpage

%% ===== Page 357 =====

Posant $Q = \sigma \varepsilon \rho$, l'équivalence $b \Leftrightarrow c$ provient de l'équivalence (valable pour tte section $Q$ de $\mathcal{M}$)
$$
(\det Q = 0, \operatorname{Tr} Q \text{ inv.}) \Longleftrightarrow (\exists \text{ projecteur idém. } P \text{ tel que } Q = u P) \leqno{(40)}
$$
qui résulte facilement des résultats sur les
idempotents, c'est un corollaire dans le page
précédente, que j'ai oublié d'expliciter.

L'équivalence avec $e), f)$ est aussi immédiate,
et c'est en fait que résument certains
calculs précédents.
\marginpar{Corollaire : Sous ces conditions $1, N_0, N_1, P$ forment une base de $\mathcal{M}$, et $\sigma, \rho$ se posent en combin. lin. de $1, N_0, N_1, P$ par $\rho = (N_0+N_1) + uP$, $\sigma^{-1} = u(N_0+N_1) + (u-u')P$. D'où : $\rho, \sigma$ forment une base [unclear: id.] de sous-[unclear: idéal] engendré par $(N_0+N_1)$. [unclear: equation]}

Notre point de vue ici est que nous
partons de $\tilde{\rho}, \tilde{\sigma}$, que nous voulons
fixer -- l'épinglage est en effet défini déjà
en termes de $\tilde{\rho}, \tilde{\sigma}$, mais il reste à
normaliser le relèvement de
$\tilde{\rho}, \tilde{\sigma}$ en $\rho, \sigma$. Ce sera déjà fait une
première restriction par les formules
$$
\varepsilon_1 = - \sigma^{-1} \rho \quad \text{[unclear: facteur]}
$$
en éliminant les scalaires de - en
multipliant $\tilde{\rho}, \tilde{\sigma}$ par des
facteurs $u$, de façon à obtenir la
relation simple :
$$
\operatorname{Tr}(\rho + \sigma) = 1 \text{ i.e. } \operatorname{Tr} \rho + \operatorname{Tr} \sigma = 1 \text{ i.e. } u = 1, \leqno{(41)}
$$
ce qui signifiera donc que $\rho, \sigma$ sont de la

\newpage

%% ===== Page 358 =====

formes
$$
\rho = \begin{pmatrix} 0 & 1 \\ v & w \end{pmatrix} \quad , \quad \sigma = \begin{pmatrix} 0 & -1 \\ v' & w' \end{pmatrix} \leqno{(41)}
$$
$$
v+v'=0, \quad w+w'=1 \quad \text{i.e.} \quad v'=-v, \quad w'=1-w \quad ,
$$
donc
$$
\varepsilon_0 = \begin{pmatrix} 1 & 1/v \\ 0 & 1 \end{pmatrix} \quad , \leqno{(42)}
$$
où on a identifié, par $\varphi$, $\operatorname{Gl}(2)_{\hat{\mathbf{Z}}}$ et $G$.

Y a-t-il unicité de $\rho, \sigma$ au dessus de $\tilde{\rho}, \tilde{\sigma}$ vu les conditions $\varepsilon_1 = -\rho \sigma$ (i.e. $\varepsilon_0 = -\sigma \rho^{-1}$) et (40) ?

Je voudrais d'abord vérifier une assertion faite par la bande un certain nombre de fois, savoir que l'application
$$
G \to PG \simeq G / \{\pm 1\}
$$
induit une bijection\footnote{[unclear: Donc, si l'on se donne une section de $PG$ en unipotents réguliers, elle se relève localement en une section de $G$ en unipotents réguliers.]}
$$
\text{unipotents réguliers de } G \to \text{unipotents réguliers de } PG \quad \text{[unclear: (à isomorphisme près)]} \leqno{(43)}
$$

la surjectivité en effet [unclear: résulte de la] définition des unipotents réguliers dans $PG$. Donc il s'agit de tester l'injectivité, ce qui signifie aussi ceci : si 
$\varepsilon, \varepsilon'$ sont deux sections unipot. régulières de $G$, telles que $\varepsilon' = \lambda \varepsilon$ (i.e. $\lambda = \pm 1$) ([unclear: $\lambda \in \Gamma(S, \mathbf{G}_m)$]), alors $\varepsilon = \varepsilon'$. Prenant les déterminants des deux membres, on aura
$$
\lambda^2 = 1 \quad , \leqno{(44)}
$$

\newpage

%% ===== Page 359 =====

prenant les traces on trouve
$$
2 \lambda = 2 \leqno{(45)}
$$
-- ce sont les conditions néc. et suff. sur $\lambda$ pour que ($\varepsilon$ étant unipot. régulier) $\lambda \varepsilon$ soit encore unipot. (et alors néc. unipotent régulier). Si $2$ est inv. sur $S$, plus géné. si $2$ est section régulière sur $\mathcal{O}_S$, cette condition implique $\lambda = 1$, OK, mais dans l'absence de telles hypothèses (p. ex. en car. $2$ [unclear: ou sur un] schéma non réduit) il n'en est plus néc. ainsi. Écrivons
$$
\lambda = 1 + \nu
$$
les conditions (44) (45) s'écrivent
$$
2\nu = 0, \quad \nu^2 = 0 \leqno{(46)}
$$
donc pour que sur $S$ l'application (43) soit injective, il faut et il suffit que la seule section sur $S$ qui satisfasse (46) soit la section nulle.

En général, désignant par $\mu'_2$ le sous-schéma en groupes de $\mathbf{G}_m$ défini par (45) (clair. c'est un sous-groupe fini sur $\mathbf{Z}$, non plat sur $\mathbf{Z}$, sous-groupe infinitésimal absolu, [unclear: de spectre] $\mathbf{Z}[T]/(T^2-1, 2T-2) \simeq \mathbf{Z}[S]/(S^2, 2S)$), [unclear: isomorphisme avec le] schéma des nombres duaux sur $\mathbf{Z}$, relatif au module $\mathbf{Z}/2\mathbf{Z}$ sur $\mathbf{Z}$, 
$$
\mu'_2 \simeq \operatorname{Spec}(\mathbf{Z} \oplus \mathbf{Z}/2\mathbf{Z}) \quad \text{(idéal de carré nul)}
$$

\newpage

%% ===== Page 360 =====

Il se trouve que $\mu'_{2, S}$ opère [MARGIN: librement par translations] (sur le schéma des unipotents réguliers de $G$, et on a

(47) $$ Unip.rég(G)/\mu'_{2, S} \xrightarrow{\sim} Unip.rég(PG) $$ .

Ceci semble à première vue en contradiction avec la bijection des Borels

(48)
\begin{tikzcd}
B \ar[d] & Borels(G) \ar[r, "\sim", "B \mapsto B/G_{...}"] \ar[d, "\wr"] & Borels(PG) \ar[d, "\wr"] & \tilde{B} \ar[d] \\
B_u & \text{"Sous-groupes unipots. à 1 param." de } G \ar[r, "\sim"] & \text{"S.-groupes unipots. à 1 par." de } PG & \tilde{B}_u \\
L \ar[rr, "|->"] & & p(L) & 
\end{tikzcd}

(où $p : G \to PG$ est le morphisme can. !)

La difficulté proviendrait-elle du fait que les flèches verticales (surjectives a priori par définition des s-groupes unipot. à 1 paramètre de $G$ resp de $PG$) ne seraient pas injectives, en dehors de l'hyp. 2 inv. (i.e. car. rés. $\neq 2$) ?

J'ai envie de dire que

(49) $$ B = \operatorname{Norm}(B_u) \text{ , } \tilde{B} = \operatorname{Norm}(\tilde{B}_u) $$

sont des formules qui permettent de retrouver $B, \tilde{B}$ à l'aide de $B_u, \tilde{B}_u$ - il suffirait bien sûr de vérifier la première de ces formules pour un groupe d'alg. lin. $M$ (i.e. le genre :

les sous-groupes unipot. à 1 paramètre de $G$ correspondent 1-1 aux sections de $\mathbb{P}$ sur $S$, i.e. aux s-fibrés en droites $D$ de $F$ ...) -

\newpage

%% ===== Page 361 =====

mais ça marche pareil dans $PG$, c'est à
dire $u \in [unclear: \operatorname{Aut}(M)]$ est tel que $u N u^{-1} = \mu N + \lambda$,
La difficulté doit plutôt provenir du
fait qu'un [unclear: élément] unipotent régulier
de $PG$ ne définit pas (sa restriction
sur $S$) un sous-groupe unipotent ;
-- alors qu'un tel unipotent régulier de $G$
$$ \varepsilon = 1 + N $$
en définit un par la formule
$$ \varepsilon(t) = 1 + t N $$
On note que si on remplace $\varepsilon$ par $\varepsilon' = (1+\nu)\varepsilon$
(avec $\nu$ satisfaisant (44)(45) $\nu^2=0, \nu N = N \nu$)
On aura
$$ \varepsilon' = (1+\nu)(1+N) = 1 + N' \quad N' = (1+\nu)N + \nu $$
et
$$ \varepsilon'(t) = 1 + t N' $$
je dis que si $\nu \neq 0$, alors les
homomorphismes composés
$$ \tilde{\varepsilon}, \tilde{\varepsilon}' : \mathbf{G}_a \to G \to PG $$
associés à $\varepsilon, \varepsilon'$ sont distincts,
bien qu'ils aient même valeur pour $t=1$.
En effet, en prenant les dérivées à l'origine
de $\tilde{\varepsilon}, \tilde{\varepsilon}'$ s'identifiant aux composées
$$ \mathcal{O}_S \to M \to M/M_0 \simeq [unclear: \operatorname{Lie}(PG)] $$

\newpage

%% ===== Page 362 =====

qui transforment $1$ en $\tilde{N}$ resp. en $\tilde{N}'$,
(où le $\sim$ désigne ici le passage au quotient
additif) on a
$$
\tilde{N}' = (1+\nu)\tilde{N}
$$
et $\tilde{N}$ est une section de $P G$ partout non
nulle, de sorte qu'on ne peut avoir
$$
\tilde{N}' = \tilde{N} \quad \text{i.e.} \quad (1+\nu)\tilde{N} = \tilde{N} \quad \text{i.e.} \quad \nu\tilde{N} = 0
$$
que si $\nu=0$, i.e. $\lambda = 1+\nu = 1$, cqfd.

En conclusion, il faut considérer que
l'on se donne dans $P G$, non seulement
$\tilde{\rho}$ et $\tilde{\sigma}$ (d'où $\tilde{\varepsilon}_0 = \tilde{\sigma}^{-1} \tilde{\rho}$, $\tilde{\varepsilon}_1 = \tilde{\rho} \tilde{\sigma}^{-1}$), mais
aussi des "sous-groupes à 1 paramètre"
(géométriques) [unclear: paramétrés] $\subset P G$,
$$
\begin{cases} 
\tilde{\varepsilon}_0 : \mathbf{G}_a \to P G \\ 
\tilde{\varepsilon}_1 : \mathbf{G}_a \to P G 
\end{cases} 
\leqno{(49)}
$$
reliés bien sûr par
$$
\begin{cases} 
\tilde{\varepsilon}_1 = {}^{\tilde{\rho}}(\tilde{\varepsilon}_0) = {}^{\tilde{\sigma}}(\tilde{\varepsilon}_0) \quad \text{i.e.} \\ 
\tilde{\varepsilon}_1(t) = {}^{\tilde{\rho}}(\tilde{\varepsilon}_0(t)) = {}^{\tilde{\sigma}}(\tilde{\varepsilon}_0(t)) 
\end{cases} 
\leqno{(50)}
$$
ce qui signifie justement qu'on se donne, en
plus de $\tilde{\rho}, \tilde{\sigma}$, un relèvement unipotent\footnote{On aura nécessairement un relèvement $\tilde{\varepsilon}_0$ tel que $\tilde{\sigma}^{-1} \tilde{\rho} = \tilde{\varepsilon}_0(1)$ [unclear: dont l'image $\dots$].} (rég.)
$\tilde{\varepsilon}_0$ de $\varepsilon_0$ (dont le choix n'est donc pas
\textit{a priori} unique), et on en déduit $\tilde{\varepsilon}_1$
par
$$
\tilde{\varepsilon}_1 = {}^{\tilde{\rho}}\tilde{\varepsilon}_0 = {}^{\tilde{\sigma}}(\tilde{\varepsilon}_0)
$$

\newpage

%% ===== Page 363 =====

Ainsi, on a apporté la [unclear: description] des demi-tours, pour pouvoir donner des choix de normalisations précis - savoir :

Proposition Soient $\tilde{\rho}, \tilde{\sigma}$ deux sections [unclear: additives] de $P G$, et $\tilde{\xi}_0$ un sous-groupe à 1-paramètre "strict" (i.e.\ monomorphisme et dont l'image soit un $\tilde{B}_n$) :
$$ \tilde{\xi} : G_a \to P G $$
tels que l'on ait :
$$ \tilde{\sigma}^{-1} \tilde{\rho} = \tilde{\xi}_0(1) \defeq \tilde{\varepsilon}_0 \leqno{(51)} $$
d'où
$$ \tilde{\sigma}(\tilde{\xi}_0) = \tilde{\rho}(\tilde{\xi}_0) \defeq \tilde{\xi}_1 \ , \text{ et} $$
$$ \tilde{\rho} \tilde{\sigma}^{-1} = \tilde{\xi}_1(1) \defeq \tilde{\varepsilon}_1 = \tilde{\rho}(\tilde{\varepsilon}_0) = \tilde{\sigma}(\tilde{\varepsilon}_0) \leqno{(52)} $$

Alors :

\begin{itemize}
\item[1)] $\exists$ uniques $\varepsilon_0, \varepsilon_1$ sections unipotentes dans $G$ telles que (posant [unclear: $\varepsilon_i(1) = 1+N_i$], $\varepsilon_i(t) = 1+t N_i$) on ait
$$ \tilde{\xi}_i(t) = \varepsilon_i(t)^\sim \leqno{(53)} $$

\item[3)] Soient $\rho, \sigma$ des relèvements de $\tilde{\rho}, \tilde{\sigma}$, \dots on a alors
$$ \varepsilon_0 = -\lambda_0 \sigma^{-1} \rho \ , \ \varepsilon_1 = -\lambda_0 \rho \sigma^{-1} \leqno{(55)} $$
pour $\lambda_0 \in \Gamma$ uniquement déterminé.

Pour un choix convenable de $\rho$ (quel que soit $\sigma$ et [unclear: déjà choisi], $\tilde{\rho}, \tilde{\sigma}$ étant donnés) $\lambda_0 = 1$ (ce qui détermine [unclear: uniquement $\rho$] \dots)
\end{itemize}

\newpage

%% ===== Page 364 =====

de sorte qu'on a
$$ \varepsilon_0 = -\sigma^{-1}\rho \quad (d'o\grave{u}\ \varepsilon_1 = -\rho\sigma^{-1}) \leqno{(56)} $$

4) $\rho$ et $\sigma$ \'{e}tant choisis de fa\c{c}on \`{a} satisfaire \marginpar{en chaque fibre} \`{a} (56), pour que $\varepsilon_0, \varepsilon_1$ soient (partout) ind\'ependants, i.e.\ pour que l'on ait $N_0 \oplus N_1$ un isomorphisme sur les fibres, il f.\ et i.s.\ qu'on ait
$$ \Tr(\rho+\sigma) \quad \text{inv.} \leqno{(57)} $$
(cf.\ p.\ 478' pour d'autres conditions \'{e}quivalentes).

5) Pour que $\varepsilon_0, \varepsilon_1$ soient (en chaque fibre) ind\'ependants, il f.\ et i.s.\ que l'on puisse trouver \marginpar{[unclear: loc. pour]} $\rho, \sigma$ au-dessus de $\tilde{\rho}, \tilde{\sigma}$ de fa\c{c}on \`{a} satisfaire \`{a} (56) et \`{a}
$$ \Tr(\rho\sigma) = 1 \quad \text{i.e.} \quad \Tr\rho \cdot \Tr\sigma = 1 \leqno{(58)} $$
\marginpar{[unclear: par un ajustement]} et cela d\'etermine $\rho, \sigma$ de fa\c{c}on unique (en termes de $\tilde{\rho}, \tilde{\sigma}, \tilde{\varepsilon}_0$).

Remarques 1) Les conditions sur ($\rho, \sigma$) envisag\'ees dans la p.\ 478', avec la normalisation $n=1$ i.e.\ $\Tr(\rho\sigma)=1$ i.e.\ $\sigma+\rho$ proj.\ \'el\'em., sont sym\'etriques en ($\rho, \sigma$) : elles sont satisfaites pour ($\rho, \sigma$) si elles le sont pour ($\sigma, \rho$) - [unclear: l'identit\'e de fa\c{c}on pr\'ecise, posant]
$$ \rho'=\sigma, \quad \sigma'=\rho \qquad \varepsilon'_0 = -\sigma'^{-1}\rho' = -\rho^{-1}\sigma, \quad \varepsilon'_1 = -\rho'\sigma'^{-1} = -\sigma\rho^{-1} \leqno{(58)} $$
d'o\`u
$$ \varepsilon'_0 = \varepsilon_0^{-1}, \quad \varepsilon'_1 = \varepsilon_1^{-1} \qquad \text{(de la p. 478')} \leqno{(59)} $$
De m\^eme, $(\rho^{-1}, \sigma^{-1})$ satisfont aux m\^emes conditions, et posant

\newpage

%% ===== Page 365 =====

$$ \rho'' = \rho^{-1}, \quad \sigma'' = \sigma^{-1}, \quad \varepsilon_0'' = -\sigma''^{-1}\rho'' = -\sigma\rho^{-1}, \quad \varepsilon_1'' = -\rho''\sigma''^{-1} = -\rho^{-1}\sigma \leqno(60) $$

$$ \varepsilon_0'' = \varepsilon_1^{-1}, \quad \varepsilon_1'' = \varepsilon_0^{-1} \leqno(61) $$

mais on fera Attention qu'on n'aura plus en général $\operatorname{Tr}(\rho''+\sigma'') = 1$, mais (par un calcul matriciel immédiat)
$$ \operatorname{Tr}(\rho''+\sigma'') = \operatorname{Tr}(\rho^{-1}+\sigma^{-1}) = \frac{1}{\det \rho} = (\det \rho)^{-1} = (\det \sigma)^{-1} \leqno(62) $$
(en supposant bien sûr $u=1$ i.e.\ $\operatorname{Tr}(\rho+\sigma)=1$)

- on n'a donc égal à $1$ que si
$$ \det \rho \ (= \det \sigma) = 1 \leqno(63) $$

- C'est la condition de la section 30) en [unclear: tête] (p. 463'), qui signifie en un sens que la normalisation $u=1$ est stable par la transformation $\rho \mapsto \rho''$, $\sigma \mapsto \sigma''$.

Par contre, le choix de l'"épi-conjugaison" associée à $(\rho, \sigma)$ (via $\tilde{\rho}, \tilde{\sigma}$) $\Big[$dans lequel l'accent a été mis sur $\varepsilon_1 = -\rho\sigma^{-1}$, qui reçoit la forme simple
$$ \varepsilon_1 = \begin{pmatrix} 1 & 1 \\ 0 & 1 \end{pmatrix} , $$
alors que $\varepsilon_0 = -\sigma^{-1}\rho$ prend une forme plus compliquée
$$ \varepsilon_0 = \begin{pmatrix} 1 & u/v \\ 0 & 1 \end{pmatrix} \quad \left( = \begin{pmatrix} 1 & -1 \\ 0 & 1 \end{pmatrix} \text{ dans le cas modulaire} \right) \Big] $$
fait intervenir implicitement l'ordre dans lequel on se donne $\rho, \sigma$ - si on les change, on change l'épi-conjugaison par [unclear: le changement] de $\varepsilon_1$ en $\varepsilon_1' = \varepsilon_1^{-1}$.

En fait, la relation affectant les deux autres des quatre épi-conjugaisons "modulaires", associées à la paire de

\newpage

%% ===== Page 366 =====

demi-groupes ; aux paires mixtes $\{L_0, L_1\}$, correspondant aux
quatre sections $\varepsilon_0, \varepsilon_0^{-1}, \varepsilon_1, \varepsilon_1^{-1}$ de l'un sur l'autre.
On a préféré normaliser à l'aide de $L_1$, i.e.\ ceux
de la forme la plus commode de l'équation
des lacets (à laquelle j'ai fini par parvenir d'ailleurs
après des pérégrinations assez lourdes, où
c'est $L_0$ que je m'obstinais à privilégier dans
les formules - cf.\ § 44, la notation définitive
n'apparaît qu'au § 45 -. La préférence donnée
à $-\rho \sigma^{-1}$ (noté $\varepsilon_1$) à son inverse $-\sigma \rho^{-1}$ (noté $\varepsilon_1^{-1}$)
provient de la préférence (justifiée par les faits)
donnée à $\rho$ sur $\sigma$, dans la formule des
lacets écrite sous la forme
$$
[\alpha, \rho] = [\beta, \sigma] \varepsilon_1(P)
$$
au lieu de l'écrire
$$
[\beta, \sigma] = [\alpha, \rho] \varepsilon'_1(P) \quad (\varepsilon'_1 = \varepsilon_1^{-1})
$$
Donc c'est $\varepsilon_1$ qui méritait alors d'avoir la forme
matricielle la plus simple possible - on a
pris $\varepsilon_1 = \begin{pmatrix} 1 & 0 \\ 1 & 1 \end{pmatrix}$ \dots

2) En termes du système
$$
\begin{matrix}
\tilde{l}_0, & \tilde{l}_1, & \tilde{l}'_0 \\
\parallel & \parallel & \parallel \\
\tilde{\rho} & \tilde{\sigma} & \tilde{\varepsilon}_0
\end{matrix} \quad \text{satisfaisant } \tilde{l}_0 . \tilde{l}_1 . \tilde{l}'_0 = 1 \leqno{(64)}
$$
(i.e.\ $\tilde{\rho}.\tilde{\sigma}.\tilde{\varepsilon}_0 = 1$ i.e.\ $\tilde{\varepsilon}_0 = \tilde{\sigma}^{-1}\tilde{\rho}^{-1}$), le groupe
$\overline{\mathfrak{S}}'_3$ y opère (cf.\ § 42 I) par permutations circulaires
de [unclear: ces] [unclear: générateurs] $(\tilde{\varepsilon}_0^{-1}, \tilde{\sigma}^{-1}, (\tilde{\rho}^{-1})^{-1} = \tilde{\rho})$, qui [unclear: donne] [unclear: les] [unclear: trois] permutations circulaires dans $(\tilde{\sigma}^{-1}, \tilde{\rho}, \tilde{\varepsilon}_0^{-1})$ -

\newpage

%% ===== Page 367 =====

le passage de $(\tilde{\rho}^{-1}, \tilde{\sigma}, \tilde{\varepsilon}_0)$ à $(\tilde{\sigma}^{-1}, \tilde{\rho}, \tilde{\varepsilon}_0^{-1})$ corres-
pond alors à l'échange $(\rho, \sigma) \mapsto (\sigma, \rho)$. Comment
s'applique - $(\tilde{\rho}^{-1}, \tilde{\sigma}, \tilde{\varepsilon}_0)$ interprétant la transformation
$(\rho, \sigma) \mapsto (\rho^{-1}, \sigma^{-1})$, i.e. $(\tilde{\rho}^{-1}, \tilde{\sigma}, \tilde{\varepsilon}_0) \mapsto (\tilde{\rho}, \tilde{\sigma}^{-1}, \tilde{\varepsilon}_1)$ ?

Interprétant les triples comme des hom
$$
\begin{tikzcd} \Pi_{0,3} \arrow[r, "\varphi"] & TPG \end{tikzcd}
$$
$$
\begin{cases} l_0 \mapsto \tilde{l}_0 \\ l_1 \mapsto \tilde{l}_1 \\ l_\infty \mapsto \tilde{l}_\infty \end{cases}
$$
on voit que l'échange $(\rho, \sigma) \mapsto (\sigma, \rho)$ correspond
à la composition de $\varphi$ avec l'autom. de $\Pi_{0,3}$ tel que
$$
l_0 \mapsto l_1^{-1}, \quad l_1 \mapsto l_0^{-1} \quad \text{donc } l_\infty \mapsto l_\infty^{-1}
$$
i.e. avec $\tau_\infty'$. De même la transformation
$(\rho, \sigma) \mapsto (\rho^{-1}, \sigma^{-1})$ correspond à la composition avec l'autom
$$
l_0 \mapsto l_0^{-1} \quad \text{d'où} \quad l_\infty \mapsto l_0 l_1
$$
$$
l_1 \mapsto l_1^{-1}
$$
qui n'est autre que $\tau_\infty = \sigma_{01} \tau_\infty'$. Donc le groupe
de quatre transformations [unclear: qui opère] sur l'un
des couples $(\rho, \sigma)$ qu'on a envisagé [unclear: ci-dessus]
\marginpar{10)} s'identifie [unclear: purement au] sous-groupe $\simeq \mathbf{Z}/2\mathbf{Z} \times \mathbf{Z}/2\mathbf{Z}$
engendré par $\tau_0', \tau_\infty'$ dans $\Autext(\Pi_{0,3}) = \mathfrak{S}_{0,3}$,
qui [unclear: certes] stabilise $L_0 \subset \Pi_{0,3}$ (donc son action
sur ce troisième facteur $\tilde{\varepsilon}_0$ [unclear: sera]
de le [unclear: conserver ou le remplacer par son inverse] ),
ce qui assure [unclear: de conserver la composante] 2- continue
unipotente, [unclear: auquel] nous tenons ! ).

\newpage

%% ===== Page 368 =====

\section*{VIII --- Récapitulation sur l'équation $[\alpha, \rho] = [\beta, \sigma] \varepsilon_1(\nu)$. \leqno{(1)}}

Le résultat essentiel auquel nous sommes parvenus dans (III), c'est la possibilité de paramétriser les solutions $\alpha, \beta, \nu$ d'une telle équation (sous des hypothèses convenables sur $\rho, \sigma$) par un groupe $B'_0 \times T_\rho \times T_\sigma$, où $B'_0$ est un « pré-Borel » dans $G$ (précisément, pour une base convenable, la forme $B'_0 = \left\{ \begin{pmatrix} \mu & \tau \\ 0 & 1 \end{pmatrix} \big| \begin{smallmatrix} \mu \in \hat{\mathbf{Z}}^* \\ \tau \in \hat{\mathbf{Z}} \end{smallmatrix} \right\}$) et $T_\rho, T_\sigma$ sont les centralisateurs des sections régulières $\rho, \sigma$ de $G$, par le morphisme
$$
\begin{cases}
B'_0 \times T_\rho \times T_\sigma \to \mathcal{G}_{\rho, \sigma} \\
(\underline{u}, \underline{a}, \underline{b}) \mapsto (\underline{u} \underline{a}^{-1}, \underline{u} \underline{b}^{-1}, \nu(\underline{u}))
\end{cases} \leqno{(2)}
$$
(où le fait que $\nu(\underline{u}) = \det(\underline{u}) - 1$ n'a pas pour le moment d'importance capitale).

\marginpar{[unclear: Nous partons ici de $\rho, \sigma, B'_0, \varepsilon_1$ (d'où $L_1$) et nous posons la question : quelles conditions pour paramétriser les solutions de (1)]}
Pour que ce morphisme à valeurs dans $G \times G \times E'$ donne des solutions de (1), il faut et il suffit ainsi que identiquement
$$
\underline{u} \rho \underline{u}^{-1} = \underline{u} \sigma \underline{u}^{-1} \varepsilon_1(\nu(\underline{u}))
$$
ce qui s'écrit aussi, en simplifiant par $\underline{u}$ et posant $y = \nu(\underline{u})$ :
$$
\rho(y) = \varepsilon_1(-y) \sigma(y) . \leqno{(3)}
$$
La possibilité de trouver $\underline{u} : B'_0 \to G$ (morphisme de schémas) donnant via (2), un morphisme $B'_0 \times T_\rho \times T_\sigma \to \mathcal{G}_{\rho, \sigma}$ de solutions...

\newpage

%% ===== Page 369 =====

des solutions de (1), équivaut donc à celle que $\sigma$ et $\rho$ [unclear: aient des] actions sur $B'_0$, modulo multiplication à gauche par [unclear: un facteur dans] $L_1$, qui intervient dans l'écriture de l'équation (1) que nous étudions. Cela se conçoit difficilement si l'on ne suppose que 
$L_1 \subset \rho(B'_0) \defeq B'_1$, ce qui impliquera d'ailleurs $\sigma(B'_0) \subset B'_1$ (et par un argument [unclear: analogue] $\sigma(B'_0) = B'_1$) et $L_1 = (B'_1)_u$. Inversement, si on a $\rho(B'_0) = \sigma(B'_0)$, i.e. $\sigma^{-1}\rho$ normalise $B'_0$, donc i.e. $\sigma^{-1}\rho$ section de $B_0$, alors $\sigma(u) = \sigma u \sigma^{-1}$ et $\rho(u) = \rho u \rho^{-1}$ sont des sections de $B'_1$ qui [unclear: sont déterminées à un facteur] près dans [unclear: $L_1$], qui coïnciderait d'ailleurs avec $(B'_1)_u$.

Ainsi on voit que la propriété essentielle pour obtenir des solutions de (1) via (2),\footnote{NB Cette condition d'avoir $(\sigma^{-1}\rho)(B'_0) = B'_0$ i.e. $\sigma^{-1}\rho \in B_0 \dots$ signifie que l'équation (4) (dont le second membre est défini \dots) exprime l'unicité de la représentation de \dots de type $(\alpha, \beta, \gamma)$ dans le groupe $(\underline{u}^{-1}, \underline{v}^{-1}, \underline{v}(\sim))$.} faut supposer p. ex.
$$
T_\sigma \cap B'_0 = \{1\} \leqno{(5)}
$$
$$
T_\rho \cap B'_0 = \{1\} \quad \text{et} \quad [unclear: G_{ur}] \cap B'_0 = \{1\} \leqno{(6)}
$$

\newpage

%% ===== Page 370 =====

[unclear: dans la forme] \sout{ $(\underline{a}, \underline{b}, c) \mapsto (U\underline{a}', V\underline{b}', \dots)$ } ramène le choix à une section de $T_\rho \cap T_\sigma \cap B_0'$. Reste la question de la surjectivité de (2) - toute solution de (1) s'écrit-elle bien sous la forme $(u \underline{a}', u \underline{b}', \nu(\underline{u}))$ ? Il revient au même (si $\rho,\sigma$ sont réguliers) de dire que la flèche
$$ B_0' \to \mathcal{H}_{\rho,\sigma} \leqno{(6)} $$
$$ \underline{u} \mapsto ([\underline{u},\rho], [\underline{u},\sigma], \nu(\underline{u})) $$
où $\mathcal{H}_{\rho,\sigma}$ est le schéma des solutions simultanées en $\xi, \eta, \nu$ de
$$ \det \xi = \det \eta = 1, \Tr \xi\rho = \Tr \rho, \Tr \eta\sigma = \Tr \sigma, \xi \eta = \varepsilon(\nu). \leqno{(7)} $$

C'est essentiellement pour obtenir cette surjectivité que nous avons été amenés dans la section (II) à faire les hypothèses draconiennes (p. 435, 436) qui s'expriment, au prix d'un détail, par les conditions que
$$ \begin{cases} \text{[unclear: par section] de } B_0 \to (B_0)_{red} = L_0 \text{ ([unclear: en tant que s-groupe de] } B_0) \\ B_0 \text{ et } B_1 = \rho B_0 \rho^{-1} = \sigma B_0 \sigma^{-1} \text{ sont ''en position générale''} \\ \text{i.e. } L_0 \text{ et } L_1 = \rho L_0 \rho^{-1} = \sigma L_0 \sigma^{-1} \text{ indépendants} \\ \text{(ou du moins [unclear: quasi-fixes])} \end{cases} \leqno{(8)} $$
mais il est [unclear: bien] clair que ces conditions (suggérées par des calculs antérieurs, et donnant lieu en effet à de très jolis calculs) sont vraisemblablement indispensables pour la validité du paramétrisa-

\newpage

%% ===== Page 371 =====

Dans notre contexte il est [unclear: naturel] de supposer
$\varepsilon_0 = [unclear: \delta^{-1} \rho]$ unipotent (ou au moins quasi-unipotent
i.e. une puissance [unclear: convenable] est unipotente), mais
la condition b) de (8), plus subtile, apparaît
comme [unclear: moins] impérative. La suite montrera si
cette restriction est judicieuse.

Dans le cas où on considère une équation
$$
[\alpha, \rho] = \varepsilon [\beta, \sigma] \varepsilon_1(v) \leqno{(9)}
$$
avec un paramètre $\varepsilon$ (pouvant [unclear: utilement]
prendre la valeur importante $-1$) $\in [unclear: \Gamma]$, on
nous avons pu pousser un peu plus loin ces calculs
détruisant la symétrie entre $\rho$ et $\sigma$, en
introduisant la condition $\operatorname{Tr} \sigma = 0$ comme
condition [unclear: d'orbitalité] des solutions $\mathcal{H}_{\rho, \sigma}$ en $y$
satisfaisant
$\det y = 1, \quad \operatorname{Tr} y \sigma = \operatorname{Tr} \sigma$
[unclear: par $fl_2$]. Le résultat essentiel auquel
nous tenions, était l'isomorphisme
$$
B'_0 \times T_\rho \times N'_0 \isommap \mathcal{G}^*_{\rho, \sigma} \leqno{(10)}
$$
$$
(u, \underline{a}, \underline{b}) \longmapsto (\alpha = u \underline{a}, \beta = u \underline{b}^{-1}, y = [unclear: \det y], \xi = \varepsilon(\underline{b}))
$$
où $N'_0 \subset N_\sigma$ est un [unclear: sous-groupe dans le] normalisateur
de $T_\sigma$ défini par la condition
$$
\underline{b} \sigma (\underline{b})^{-1} = \varepsilon(\underline{b}) \sigma \quad \text{i.e.} \quad \underline{b}(\sigma) = \varepsilon(\underline{b}) \sigma \leqno{(11)}
$$

\newpage

%% ===== Page 372 =====

Que le morphisme (10) s'envoie bien dans $G^*_{\rho,\sigma}$,
est trivial par définition de $N'_0$ et vu le fait
analogue pour (2) --- sans hypothèse du genre
$\Tr \sigma = 0$. Mais la possibilité d'obtenir (loc$^t$)
n'importe quel $\varepsilon$ via un $\underline{b} \in N'_0$ --- ou
simplement $\varepsilon = -1$, si $2 \neq 0$) exige justement
$\Tr \sigma = 0$.

Bien entendu, nous considérons
$G^*_{\rho,\sigma}$ comme un schéma en groupes,
par transport de structure
de groupe produit de $B'_0 \times T_\rho \times N'_0$.
On y distingue le sous-groupe
$$ G^{*(0)}_{\rho,\sigma} \text{ correspondant à } T'_0 \times T_\rho \times N'_0 \leqno{(12)} $$
où
$$ T'_0 = B'_0 \cap B'_1 \leqno{(13)} $$
est le tore maximal privilégié de $B'_0$, qui
pour notre [unclear: repère canonique] se décrit
comme
$$ T'_0 = \left\{ \begin{pmatrix} \mu & 0 \\ 0 & \mu^{-1} \end{pmatrix} \;\middle|\; \mu \in k^* \right\} \leqno{(14)} $$
On pose de même
$$ G^{* \#}_{\rho,\sigma} = G_{\rho,\sigma} \cap G^{*(0)}_{\rho,\sigma} \text{ correspondant à } T'_0 \times T_\rho \times T_0 \leqno{(15)} $$

\newpage

%% ===== Page 373 =====

On a donc des structures de produit semi-direct

$$
\mathcal{G}_{\rho, \sigma} \isom \mathcal{G}_{\rho, \sigma}^{(0)} \cdot \Lambda_0 \quad , \quad \mathcal{G}_{\rho, \sigma}^* \isom \mathcal{G}_{\rho, \sigma}^{*(0)} \cdot \Lambda_0 \quad ,
\leqno (16)
$$

$$
\left.
\begin{aligned}
\Lambda_0 \subset \mathcal{G}_{\rho, \sigma} , \quad \Lambda_0 &= \left\{ (\underset{\varepsilon_0^n}{\alpha=\lambda}, \underset{\varepsilon_0^n}{\beta=\lambda}, \mu=1) \mid \underset{(\text{ou } \in \Gamma G_0)}{\lambda \in \Gamma L_0} \right\} \\
& \text{correspondant à } \quad L_0 \times \{1\} \times \{1\} \subset B'_0 \times T_\rho \times T_\sigma
\end{aligned}
\right\}
\leqno (17)
$$

Nous allons enfin considérer un autre sous-groupe de $\mathcal{G}_{\rho, \sigma}^*$ , par les conditions

$$
\begin{cases}
\det \alpha \in \{\pm 1\} \text{ i.e. } (\det \alpha)^2 = 1 \\
\det \beta = 1
\end{cases}
\quad \text{d'où } S\mathcal{G}_{\rho, \sigma}^* \subset \mathcal{G}_{\rho, \sigma}^*
\leqno (18)
$$

[unclear: Ce s-groupe] $S\mathcal{G}_{\rho, \sigma}^*$ peut s'interpréter comme la solution du système

$$
(\alpha, \beta, \mu, \varepsilon, \varepsilon')
$$

satisfaisant à\footnote{NB $\varepsilon^2=1$ résulte déjà de la première relation.}

$$
[\alpha, \beta] = \varepsilon [\beta, \sigma] \varepsilon_1(\mu-1), \quad \det \alpha = \varepsilon', \quad
\begin{matrix}
\det \beta = 1, \\
\varepsilon^2 = {\varepsilon'}^2 = 1
\end{matrix}
\leqno (19)
$$

En termes des systèmes

$$
(\underline{x}, \underline{y}, \underline{b}) \in \Gamma B'_0 \times T_\rho \times N'_0
$$

les équations (19) (dont la première est automatique) se réduisent à

$$
\det \underline{y} = \det \underline{b} = \varepsilon \det \underline{x} \quad , \quad {\varepsilon'}^2 = 1 \quad ,
\leqno (20)
$$

\newpage

%% ===== Page 374 =====

qui montrent que $S\mathcal{G}_{p,\sigma}^*$ est bien\marginpar{app.} un sous-groupe
de $\mathcal{G}_{p,\sigma}^*$. Ce sous-groupe contient $\Lambda_0$
(correspondant à $L_0 \times \{1\} \subset B'_0 \times T_p \times \mathcal{N}'_0$ ) - qui
correspond au cas où l'on a
$$
\det \underline{y} = \det \underline{b} = \det \underline{a} = 1 \quad (\varepsilon' = 1) , \leqno{(21)}
$$
et on a une structure de produit semi-direct
$$
S\mathcal{G}_{p,\sigma}^* \simeq S\mathcal{G}_{p,\sigma}^{*(0)} \cdot \Lambda_0 \leqno{(22)}
$$
où
$$
S\mathcal{G}_{p,\sigma}^{*(0)} = S\mathcal{G}_{p,\sigma}^* \cap \mathcal{G}_{p,\sigma}^{*(0)} \leqno{(23)}
$$
correspond au sous-groupe de $B'_0 \times T_p \times \mathcal{N}'_0$
défini par les dites équations
(20). Bien sûr, on a aussi
$$
S\mathcal{G}_{p,\sigma} \ , \ S\mathcal{G}_{p,\sigma}^{(0)} \quad (\text{sans } {}^*) \leqno{(24)}
$$
obtenus comme intersections de $S\mathcal{G}_{p,\sigma}^*$, $S\mathcal{G}_{p,\sigma}^{*(0)}$ avec
$\mathcal{G}_{p,\sigma}$ - se déduit de l'expression
de ceux-ci en termes de $(B'_0, T_0), T_p, \mathcal{N}'_0$
en remplaçant $\mathcal{N}'_0$ par $T_\sigma$.

Pour expliciter la structure de $S\mathcal{G}_{p,\sigma}^*$,
on est ramené vu (17) à étudier le
groupe
$$
S\mathcal{G}_{p,\sigma}^{*(0)} \subset T_0 \times T_p \times \mathcal{N}'_0 \leqno{(25)}
$$

\newpage

%% ===== Page 375 =====

décrit par les équations (20). Mais on a
$$
T_0 \isommap \mathbf{G}_m \quad , \quad u \mapsto \det u
$$
ce qui implique que la projection
$$
S\mathcal{G}^*_{\rho, \sigma(0)} \to T_\rho \times N'_0
$$
\begin{tikzcd}
(u a^{-1}, u b^{-1}, \det u, \varepsilon=\varepsilon(b)) \arrow[r, mapsto] & (a,b) \\
\rotatebox{90}{$\simeq$} & \\
(u, a, b) & 
\end{tikzcd}
est [unclear: surj.], et induit un iso. de
$S\mathcal{G}^{*(0)}_{\rho, \sigma(0)}$ sur le s-groupe de $T_\rho \times N'_0$
défini par l'équation
$$
\det b = \varepsilon' \det a \quad , \quad \varepsilon'^2 = 1 \leqno{(21)}
$$
ou encore $(\det b)^2 = (\det a)^2$.
En termes des homomorphismes (définis [unclear: ci-dessus])
$$
\delta_\rho : T_\rho \to \mathbf{G}_m \quad , \quad \delta_0 : N'_0 \to \mathbf{G}_m
$$
cela revient à prendre l'image réciproque de la diagonale par $\delta_\rho^2, \delta_0^2$.

On voit tout de suite (comme $\delta_\rho, \delta_0$
sont épis) que l'hom.
$$
S\mathcal{G}^*_{\rho, \sigma(0)} \to \mu_2 \times \mu_2 \leqno{(22)}
$$
$$
(\alpha, \beta, \gamma, \varepsilon, \varepsilon') \mapsto (\varepsilon, \varepsilon')
$$
$$
\text{i.e.} \quad (u, a, b) \mapsto (\varepsilon(b), \det b / \det a)
$$
est épi, et le noyau [unclear: s'insère dans] une suite exacte

\newpage

%% ===== Page 376 =====

$$
1 \to S\mathcal{G}_{\rho,\sigma(\omega)}^{* 0 \natural} \to S\mathcal{G}_{\rho,\sigma}^{* \natural} \to \mu_2 \times \mu_2 \to 1 \leqno (23)
$$
où $S\mathcal{G}_{\rho,\sigma(\omega)}^{0 \natural}$ est isom. au sous-schéma de $T_{\rho} \times T_{\sigma}$ défini par l'équation
$$
\det \underline{a} = \det \underline{b} \leqno (24)
$$
et se dévisse donc :
$$
1 \to S T_{\rho}^{\natural} \times S T_{\sigma}^{\natural} \to S\mathcal{G}_{\rho,\sigma(\omega)}^{0 \natural} \to \mathbb{G}_m \to 1 \leqno (25)
$$
$$
(\underline{a}, \underline{b}) \mapsto \det \underline{a} = \det \underline{b} = \det
$$
où on a posé
$$
\begin{aligned}
S T_{\rho}^{\natural} &= \operatorname{Ker}(T_{\rho} \xrightarrow{\delta_{\rho}} \mathbb{G}_m) \\
S T_{\sigma}^{\natural} &= \operatorname{Ker}(T_{\sigma} \xrightarrow{\delta_{\sigma}} \mathbb{G}_m)
\end{aligned} \leqno (26)
$$
\marginpar{m-à-j}

Si par exemple $\rho$ est [unclear: régulier], donc $T_{\rho}$ un tore maximal, alors $ST_{\rho}^{\natural}$ est un tore.
Ainsi, si dans le "cas modulaire",
$$
\rho = \begin{pmatrix} 0 & -1 \\ 1 & 1 \end{pmatrix} , \quad \sigma = \begin{pmatrix} 0 & -1 \\ 1 & 0 \end{pmatrix} , \quad S T_{\rho}^{\natural} \text{ est un}
$$
tore de dim. $1$ (un $\mathbb{G}_m$), $S T_{\sigma}^{\natural}$ tore de dim. $0$ (un $\mu_2$), $S\mathcal{G}_{\rho,\sigma(\omega)}^{0\natural}$ est un tore de dim. $2$.

Pour terminer de déterminer la structure de $S\mathcal{G}_{\rho,\sigma}^{*\natural}$ au travers de $S\mathcal{G}_{\rho,\sigma(\omega)}^{0\natural}$ et de la décomposition (22) en produit semi-direct, il faut expliciter l'opération de

\newpage

%% ===== Page 377 =====

$SG_{\rho,\sigma}^{*(0)}$ sur $\Lambda_0$, en identifiant
$SG_{\rho,\sigma}^{*(0)}$ comme sous-groupe de $T_\rho \times N'_\sigma$,
et $\Lambda_0$ comme $\simeq \mathbb{G}_m$ par $\lambda \mapsto \varepsilon_0(\lambda)$.

Or la relation

$$ \underline{u} (\varepsilon_0(\lambda)) \underline{u}^{-1} = \varepsilon_0( \det(\underline{y}) \lambda ) \quad , \quad \det(\underline{y}) = \det(b) = \delta_\sigma(b) $$

montre que l'opération se fait via
l'application composée

(27)
$$ SG_{\rho,\sigma}^{*(0)} \hookrightarrow N'_\sigma \xrightarrow{\delta_\sigma} \mathbb{G}_m \quad , \quad OK. $$

Pour la suite, nous aurons besoin
d'un hom. canonique

(28)
\begin{tikzcd}
G_{\rho,\sigma}^* \ar[r, "t"] & G \\
B'_0 \times T_\rho \times N'_\sigma \ar[ur, dashed] &
\end{tikzcd}

qui s'identifie à la première projection
$$ (y, \underline{a}, \underline{b}) \mapsto y $$
du produit, mais considérée comme
un morphisme dans $G$ grâce à l'inclusion
$B'_0 \hookrightarrow G$.

On pose aussi
$$ u = \underline{\Theta}(\alpha, \beta, \mu, \varepsilon), $$
ce dont on aura [unclear]

\newpage

%% ===== Page 378 =====

$$
\begin{cases}
\underline{u}(\rho) = \xi \rho, & \underline{u}(\sigma) = \eta \sigma \\
\| \text{défin.} & \| \text{défin.} \\
u \rho u^{-1} \operatorname{int} h_\rho \rho & \varepsilon \operatorname{int} h_\sigma \sigma \\
\det \underline{u} = \mu
\end{cases} \leqno (29)
$$
\marginpar{$\{ \xi \alpha(\rho)^\sim, \eta \beta(\sigma)^\sim \}$ [unclear: à un] facteur $\pm 1$}

Proposition. Le morphisme de schémas $\Theta$ est
caractérisé [unclear: (pour les foncteurs $\underline{G}_{\rho,\sigma} \to \mu_2$)] par les formules (29), mod un
morphisme $\underline{G}_{\rho, \sigma} \to \mu_2$.
En effet, les premières deux formules, qui
décrivent l'action de $\underline{u}$ sur $\rho$ (resp. $\sigma$), déterminent
$\underline{u}$ mod $T_\rho$ resp. mod $T_\sigma$, donc (comme
$T_\rho \cap T_\sigma = \{ \pm 1 \}$,) la [unclear: deuxième] le détermine
mod $S G_1$, -- d'où avec la précédente,
mod $\mu_2$. Donc si on a un autre morphisme
$\Theta'$ ayant ces propriétés on a $\Theta' = \lambda \Theta$,
(i.e. $\Theta'(g) = \lambda(g) \Theta(g)$), avec $\lambda : \underline{G}_{\rho, \sigma} \to \mu_2$.

\marginpar{Si on travaille sur $S\underline{G}_{\rho,\sigma}$ [unclear: donc] l'isomorphisme [unclear: est canonique] (car [unclear: $\lambda(\pm 1) = \pm 1$])}
Remarque. Notons d'ailleurs que tout $\Theta$ est normalisé par
$$ \Theta(1,1,1,1) = 1 $$
ce qui impose une [unclear: condition supplémentaire]
$\lambda(1) = 1$, sur $\lambda$ -- mais il y a en outre une
implication déjà (car mod $\rho$, $\sigma$ sont d'ordre 2
dans $T_\rho, T_\sigma$ des [unclear: torseurs] \dots) que $\lambda$ est un hom.
[unclear: trivial sur] la composante (qui est [unclear: connexe] \dots)
multiplicative sur $\mathcal{B}'_\sigma \times T_\rho \times T_\sigma$ (de [unclear: niveau] \dots $\ne 2$)
donc un terme essentiellement [unclear: nul] modulo une
indétermination possible, la multiplication par un
$\varepsilon$ (correspondante à l'hom. [unclear: trivial]
de $N_{\sigma}/T_{\sigma}$ sur $\mu_2$). $\square$

\newpage

%% ===== Page 379 =====

\noindent
(IX) \underline{Les groupes $\mathcal{M}_{\rho, \sigma}$ et $S\mathcal{M}_{\rho, \sigma}$.}

Soient toujours $\rho, \sigma$ deux sections de $G$, satisfaisant les conditions générales (cf p. ex. page [unclear: 478]), plus la condition $\operatorname{Tr} \sigma = 0$ (cf $\overline{X}$ - p. ex. lemme p. 455 et son cor. p. 457) - nous pourrions faire la construction [unclear: omettant] ces conditions supplémentaires (en nous contentant de travailler avec $\mathcal{G}_{\rho, \sigma}$ au lieu de $\mathcal{G}_{\rho, \sigma}^*$) - mais peu importera dans la suite. Nous pourrons au besoin supposer $\rho, \sigma$ normalisés de façon qu'on ait $\operatorname{Tr}(\rho \sigma) = 1$.

On désigne par $S\mathcal{M}_{\rho, \sigma}$ le schéma produit\marginpar{Pour la suite la "proanab." du profini [unclear: des natures de groupes] $S\mathcal{M}_{\rho, \sigma}$, cf p. 559 ff.}
$$
S\mathcal{M}_{\rho, \sigma} = S\mathcal{G}_{\rho, \sigma}^* \times_S S\mathcal{G} \leqno{(1)}
$$
ses sections sont notées
$$
(\alpha, \beta, \mu, \varepsilon, f) \ ,
$$
$f$ étant donc une section arbitraire de $S\mathcal{G}$
(i.e.\ $\operatorname{Sl}(2)_S$).

Dans le cas $S = \operatorname{Spec}(A)$, $A = \hat{\mathbf{Z}}$,
$G = \operatorname{Gl}(2)_S$, $\rho = \begin{pmatrix} 0 & 1 \\ -1 & 1 \end{pmatrix}$, $\sigma = \begin{pmatrix} 0 & -1 \\ 1 & 0 \end{pmatrix}$,
et on a dans ce cas un homomorphisme canonique
$$
\hat{\mathfrak{S}}_{0,3} \to G(\hat{\mathbf{Z}})^{\pm 1} \simeq \operatorname{Sl}(2, \hat{\mathbf{Z}})^{\pm 1} / \{\pm 1\} \leqno{(2)}
$$
à valeurs dans le sous-groupe des éléments de dét $\pm 1$,

\newpage

%% ===== Page 380 =====

envoyant $\widehat{\mathfrak{S}}^+_{0,3}$ dans $\SL_2(\hat{\mathbf{Z}})' = \operatorname{Sl}(2,\hat{\mathbf{Z}})' = \SL_2(\hat{\mathbf{Z}})/\{\pm 1\}$

on va construire une variante du groupe
$\widetilde{\mathcal{M}}_{0,3}$ (cf p. 390), qui s'envoie canoniquement dans le
produit (1), ce qui nous amènera à définir
une structure de groupe sur $S\mathcal{M}_{\rho,\sigma}$, de
façon à rendre ledit morphisme multiplicatif.

Considérons d'abord la version ``rigide''
$S\widetilde{\mathcal{M}}_{0,3}$ de $\widetilde{\mathcal{M}}_{0,3}$, extension de $\widetilde{\mathcal{M}}_{0,3}$
par $\hat{\mathbf{Z}}^3$, déduite d'une extension de
$\widetilde{\mathcal{M}}_{0,\hat{3}} = \widetilde{\mathcal{M}}_{0,3} / \mathfrak{T}_{0,\hat{3}} \subset \widehat{\mathfrak{T}}_{0,3}$ , [unclear: par eh $=$]
[unclear: Nous noterons] à $\widehat{\mathfrak{T}}_{0,3}$ (via $\widetilde{\mathcal{M}}_{0,3} \to \widehat{\mathfrak{T}}_{0,3}$)

formée des
$$ (u, f_0, f_1, f_\infty) \quad (u \in \widetilde{\mathcal{M}}_{0,3}, f_0, f_1, f_\infty \in \widehat{\mathfrak{T}}_{0,3}) \leqno{(3)} $$
satisfaisant
$$ u(l_i) = \operatorname{int}(f_{u(i)}^{-1})(l_{u(i)}) \leqno{(4)} $$
où [unclear: $l_i$ est la multiplication] de $u$, et $u(i)$ l'action de $u$ sur $\{0, 1, \infty\}$
[unclear: dans $\dots$]. Le produit est défini par
$$ (u', f'_0, f'_1, f'_\infty)(u, f_0, f_1, f_\infty) = (u'u, f'_0 u'(f_0), f'_1 u'(f_1), f'_\infty u'(f_\infty)) \leqno{(5)} $$
bien entendu (disons $f'_{u(i)} u'(f_i)$ !) -- en
[unclear: $u, u' \in \widetilde{\mathcal{M}}_{0,3}$ avec $u, u' = 1$].

\newpage

%% ===== Page 381 =====

On a un hom. surjectif $S\mathcal{M}_{0,3}^{\sim} \to \mathcal{M}_{0,3}^{\sim}$,
dont le noyau est formé des
$$ (1, f_0, f_1, f_\infty) \qquad f_i \in L_i \leqno (6) $$
d'où
$$ 1 \to \hat{\mathbf{Z}}^3 \to S\mathcal{M}_{0,3}^{\sim} \to \mathcal{M}_{0,3}^{\sim} \to 1 \leqno (7) $$

Nous allons définir un sous-groupe intéressant
$$ S_0\mathcal{M}_{0,3}^{\sim} \subset S\mathcal{M}_{0,3}^{\sim} \leqno (8) $$
en utilisant maintenant la définition des
$\mathcal{M}_{0,3}^{\sim}$, à savoir la propriété des $u \in \mathcal{M}_{0,3}^{\sim}$
de commuter extérieurement à $\rho$, de sorte qu'on a
$$ u \rho u^{-1} = g \rho \qquad \text{i.e.} \quad [u, \rho] = g \quad , \quad g \in \hat{\Pi}_{0,3} \leqno (9) $$

Ceci permet, quand on a [unclear: le choix] des $f_i \in \hat{\Pi}_{0,3}$
satisfaisant
$$ u(L_i) = \operatorname{int}(f_i^{-1})(L_i) \leqno (10) $$
(en supposant pour simplifier $u \in \mathcal{M}_{0,3}^{\sim \wedge}$),
d'en déduire un [unclear: choix] fixe, fixé via
la formule
$$ f_{i+1}^{-1} = g \rho(f_i^{-1}) \qquad \left[ \text{et } f_{i-1}^{-1} = g \rho(f_{i+1}^{-1}) = g \rho(g) \rho^2(f_i^{-1}) \right] \leqno (11) $$
(où $g = [u, \rho] \in \hat{\Pi}_{0,3}$)
(on a posé $\rho(L_i) = L_{i+1}$ pour $i=0,1,\infty$
i.e. $0+1=1$, $1+1=\infty$, $\infty+1=0$ !) - on retrouve
alors
$$ f_i^{-1} = f_{(i+2)+1}^{-1} = g \rho(f_{i+2}^{-1}) = \underbrace{g \rho(g) \rho^2(g)}_{1} \rho^3(f_i^{-1}) = f_i^{-1} $$

\newpage

%% ===== Page 382 =====

492

les $(u, f_0, f_1, f_\infty) \in \mathcal{S}M_{0,3}^\sim$ qui satisfont
(11) pour tout $i \in \{0,1,\infty\}$ (il suffit qu'il le
vérifie pour deux indices, pour que ce soit
vérifié pour tous les trois) forment un
sous-groupe de $\mathcal{S}M_{0,3}^\sim$, que je note
$\mathcal{S}_0M_{0,3}^\sim$ - il s'identifie aussi au groupe
des couples $(u, f_0)$ ($u \in M$, $f_0 \in \hat{\Pi}_{0,3}^\sim$)
satisfaisant
(12) $$u(l_0) = int(f_0^{-1})(l_0^u) \quad (l_0^u = \chi(u)) .$$

On a une structure d'extensions
(13) $$1 \to \hat{\mathbb{Z}} \to \mathcal{S}_0M_{0,3}^{!\sim} \to M_{0,3}^{!\sim} \to 1 ,$$

et un hom d'extensions
(14) 
\begin{tikzcd}
1 \ar[r] & \hat{\mathbb{Z}} \ar[r] \ar[d, "diagonale"'] & \mathcal{S}_0M_{0,3}^{!\sim} \ar[r] \ar[d] & M_{0,3}^{!\sim} \ar[r] \ar[d, "\simeq"] & 1 \\
1 \ar[r] & \hat{\mathbb{Z}}^3 \ar[r] & \mathcal{S}M_{0,3}^{\sim} \ar[r] & M_{0,3}^{\sim} \ar[r] & 1
\end{tikzcd}

(où $\mathcal{S}_0M_{0,3}^{!\sim}$ est [unclear: par définition] l'image inverse de $M_{0,3}^{!\sim} \subset M_{0,3}^\sim$ dans $\mathcal{S}M_{0,3}^\sim$). On peut donc dire que l'extension $\mathcal{S}M_{0,3}^\sim$ (de $M_{0,3}^\sim$ par $\hat{\mathbb{Z}}^3$)
est déduite, par l'hom. de groupes

(15) $$\hat{\mathbb{Z}} \to \hat{\mathbb{Z}}^3$$
$$n \mapsto (n_0=n, n_1=n, n_\infty=n)$$

d'une extension "plus fine" $\mathcal{S}_0M_{0,3}^{!\sim}$ de

\newpage

%% ===== Page 383 =====

$\mathcal{M}_{0,3}^{\sim}$ par $\hat{\mathbf{Z}}$.

\marginpar{[raturé : associant à $(u,f)$ (où $u \in \hat{\mathbf{Z}}, f \in \mathcal{M}_{0,3}^{\sim [\sigma]}$) un elt de $\mathcal{M}_{0,3}^{\sim}$ de la forme $u \alpha(u) f$]}
Les calculs du § 46 V (p 407 ff) permettent de fait d'expliciter les [unclear: éléments] $(\alpha, \beta, f)$ au travers des systèmes
$$
(\alpha, \beta, f) \quad \left\{ \begin{array}{l} \alpha \in \mathfrak{S}_{0,3}^{\sim} = \hat{\Pi}_{0,3} \cdot \{1, \tau_\infty\} = \text{image inv. de } \{1, \tau\} \\ \text{par } \mathfrak{S}_{0,3}^{\sim} \to \mathcal{M}_{0,3} \\ \beta, f \in \hat{\Pi}_{0,3}^{\sim} \end{array} \right. \leqno{(16)}
$$
satisfaisant des relations (une infinité) que j'écris par extension
$$
\alpha \rho \alpha^{-1} = \beta \rho f \beta f^{-1} \varepsilon_1(\mu-1) \leqno{(17)}
$$
où $\mu \in \hat{\mathbf{Z}}^*$ est défini sans ambiguïté par $\alpha$ et ne dépend que de l'image de $\alpha$ dans $\operatorname{Gl}(2, \hat{\mathbf{Z}})$,
$$
\mu = \det(\bar{\alpha})^{-1} \text{, où } \bar{\alpha} \text{ est l'image de } \alpha \text{ dans } \operatorname{Gl}(2, \hat{\mathbf{Z}}).
$$

Le cas $f=1$ correspond au sous-groupe de $S_0 \mathcal{M}_{0,3}^{\sim}$, qui s'identifie au normalisateur de $\tau_\infty$ dans $\mathcal{M}_{0,3}^{\sim}$, que j'avais noté $\mathcal{M}_{0,3}^{\sim [\sigma]}$ ( § 46 V $\dots$ ) :
$$
\mathcal{M}_{0,3}^{\sim [\sigma]} \simeq \left\{ (\alpha, \beta, f) \middle| \begin{array}{l} f=1, \alpha, \beta \\ \text{satisfaisant (17)} \end{array} \right\} \subset S_0 \mathcal{M}_{0,3}^{\sim} \leqno{(18)}
$$
$$
= \operatorname{Norm}_{\mathcal{M}_{0,3}^{\sim}}(\langle \tau_\infty \rangle)
$$

Si on désigne par $u_{\alpha, \beta}$ l'automorphisme de $\hat{\Pi}_{0,3}^{\sim}$ défini au moyen de $\alpha, \beta$, on a :
$$
\xi = \alpha \rho \alpha^{-1}, \quad \eta = \beta \rho \beta^{-1} \leqno{(19)}
$$

\newpage

%% ===== Page 384 =====

$$
u(\rho) = \xi \rho \, , \quad u(\sigma) = \eta \sigma \leqno (20)
$$
(NB $\eta$ est défini en termes de $\xi$ grâce à (17), et $\mu = \mu(f)$ est déterminée par $f$ ), on aura par construction (loc.\ cit.)
$$
u_{\alpha, \beta, f} = \operatorname{int}(f^{-1}) u_{\alpha, \beta} \leqno (21)
$$
et un triplet $(\alpha, \beta, f)$ correspond à l'élément $(u_{\alpha, \beta, f}, f)$ de $So\tilde{\mathcal{M}}_{0,3}^\sim$ :
$$
(\alpha, \beta, f) \longmapsto (u_{\alpha, \beta, f}, f) \in So\tilde{\mathcal{M}}_{0,3}^\sim \ . \leqno (22)
$$
Ayant ainsi explicité les éléments de $So\tilde{\mathcal{M}}_{0,3}^\sim$, on définit une application
$$
\begin{tikzcd}
So\tilde{\mathcal{M}}_{0,3}^\sim \ar[r, "\isom"] & S\mathcal{M}_{\rho, \sigma}(\hat{\mathbf{Z}})' \defeq S\mathcal{G}_{\rho, \sigma}^*(\hat{\mathbf{Z}})' \times \operatorname{Sl}(2, \hat{\mathbf{Z}})' \\
(\alpha, \beta, f) \ar[r, |->] & (\underline{\alpha}, \underline{\beta}, \mu=\mu(f), \xi = \alpha\rho(\alpha)^{-1}, \underline{f})
\end{tikzcd}
\leqno (23)
$$
où les soulignés $\underline{\alpha}, \underline{\beta}, \underline{f}$ désignent les images de $\alpha, \beta, f$ dans $\operatorname{Sl}(2, \hat{\mathbf{Z}})'$, et le [unclear: terme] $S\mathcal{M}_{\rho, \sigma}(\hat{\mathbf{Z}})'$ désigne l'ensemble des systèmes $(\underline{\alpha}, \underline{\beta}, \mu, \xi, \underline{f})$ d'éléments $\in \operatorname{Sl}(2, \hat{\mathbf{Z}})' \times \operatorname{Sl}(2, \hat{\mathbf{Z}})' \times \hat{\mathbf{Z}}^\times \times \pi_1 \times \operatorname{Sl}(2, \hat{\mathbf{Z}})'$ (dans $\operatorname{Sl}(2, \hat{\mathbf{Z}})'$ (pas dans $\operatorname{Sl}(2, \hat{\mathbf{Z}})$ ! )), satisfaisant à la relation locale
$$
\underline{\alpha} \rho(\underline{\alpha})^{-1} = \varepsilon\ \underline{\beta} \sigma(\underline{\beta})^{-1}\ \varepsilon_1(\mu - 1) \leqno (24)
$$
(qui ne dépend pas du choix des signes .... )

\newpage

%% ===== Page 385 =====

On a "fait ce qu'il fallait" pour que la restriction de (23) à $S_0\mathcal{M}^\sim_{0,3}[0]$ ((18) identifiant aux $(\alpha, \beta, f)$ pour lesquels $\underline{f}=1$) soit un homomorphisme dans le groupe $S\mathcal{G}^\ast_{p,\sigma}$ -- en définissant du coup, et de façon ad hoc, la loi de composition dans $S\mathcal{G}^\ast_{p,\sigma}$. Je vais la réécrire quand même pour me rassurer, en considérant $u=u_{\alpha,\beta}, u'=u_{\alpha',\beta'}$, $u''=uu'$ sous la forme $u_{\alpha'',\beta''}$, et on va vérifier la "relation des cocycles" qui donne, à des éléments du centre près, $(\alpha'',\beta'')$ en termes de $(\alpha,\beta), (\alpha',\beta')$ (§ 46 II 5°)
$$
\alpha'' = \underline{u}'(\alpha)\alpha' \quad \beta'' = \underline{u}'(\beta)\beta' \quad \mu'' = \mu\mu' \leqno{(25)}
$$
où les $\underline{\alpha}, \underline{\beta}, \underline{\alpha}', \underline{\beta}'$ désignent des éléments de $\operatorname{Gl}(2, \hat{\mathbf{Z}})^{\pm}$ (déterminés à un signe près) qui relèvent les images de $\alpha, \alpha', \beta, \beta'$ dans $\operatorname{Gl}(2, \hat{\mathbf{Z}})'$, de sorte que $(\alpha, \beta, \mu, \varepsilon_i)$ et $(\alpha', \beta', \mu', \varepsilon_i')$ sont dans $S\mathcal{G}^\ast_{p,\sigma}(\hat{\mathbf{Z}})$, et on associe des $\underline{u}, \underline{u}' \in \operatorname{Gl}(2, \hat{\mathbf{Z}})$ (en fait $\in B'_0(\hat{\mathbf{Z}})$) tels que
$$
\begin{cases}
\underline{u}'(p) = \operatorname{int}(\underline{\alpha}') p \quad \underline{u}'(\sigma) = \varepsilon_2(\mu') \operatorname{int}(\underline{\beta}')(\sigma) \\
\det \underline{u}' = \mu'
\end{cases} \leqno{(26)}
$$
(cf relation id. pour $u$) et on aura alors, pour tout $x \in \hat{\mathfrak{T}}_{0,3}$
$$
u'(x) = \underline{u}'(x) \leqno{(27)}
$$
(cf [unclear: ...]) Donc, pour des relevés convenables $\underline{\alpha}'', \underline{\beta}''$ des images de $\alpha'', \beta''$ dans $\operatorname{Gl}(2, \hat{\mathbf{Z}})^{\pm}$, on aura
$$
\underline{\alpha}'' = \underline{u}'(\underline{\alpha}) \underline{\alpha}', \quad \underline{\beta}'' = \underline{u}'(\underline{\beta}) \underline{\beta}' \quad (\text{en plus de } \mu'' = \mu\mu' \text{ d'où } \varepsilon_i'' = \varepsilon_i\varepsilon_i') \leqno{(28)}
$$

\newpage

%% ===== Page 386 =====

et il reste à vérifier que l'élément
$$
g'' = (\underline{\alpha}'', \underline{\beta}'', \mu''=\mu\mu', \varepsilon_2''=\varepsilon_2'\varepsilon_2)
$$
est bien le produit des éléments
$g' = (\underline{\alpha}', \underline{\beta}', \mu', \varepsilon_2')$ et $g = (\underline{\alpha}, \underline{\beta}, \mu, \varepsilon_2)$.
Que ce soit ok pour $\mu''$ et $\varepsilon''$ est trivial, il reste
: vérifier que dans $\mathcal{G}_{\vec{p}, \sigma}^*$ , le produit est bien
Ceci en ce qui concerne les composantes en $\underline{\alpha}, \underline{\beta}$
donné par les formules (28) , ce qui
sera valable, bien entendu, sans se borner
pour $\underline{u}, \underline{u}'$ quelconques, et dans les conditions générales du
IV (ou du III si on se restreint à $\mathcal{G}_{\vec{p}, \sigma}^*$ au lieu
de $\mathcal{G}_{\vec{p}, \vec{\sigma}}^*$) . On aura
$$
\begin{cases} 
\underline{\alpha} = \underline{u} \underline{a}^{-1}, \quad \underline{\beta} = \underline{u} \underline{b}^{-1} \\ 
\underline{\alpha}' = \underline{u}' \underline{a}'^{-1}, \quad \underline{\beta}' = \underline{u}' \underline{b}'^{-1} 
\end{cases} 
\qquad 
\begin{matrix} 
\underline{\alpha}, \underline{\alpha}' \in T_{\vec{p}}(\hat{\mathbf{Z}}) \\ 
\underline{\beta}, \underline{\beta}' \in N_{0}(\hat{\mathbf{Z}}) 
\end{matrix} 
\leqno (29)
$$
et par définition du produit dans $\mathcal{G}_{\vec{p}, \sigma}^*$ ,
le produit $g' g$ aura pour composantes $\underline{\alpha}''$ et $\underline{\beta}''$ :
$$
\underline{\alpha}'' = (\underline{u}' \underline{u}) (\underline{a}' \underline{a})^{-1} \quad , \quad \underline{\beta}'' = (\underline{u}' \underline{u}) (\underline{b}' \underline{b})^{-1}
$$
et il reste à prouver
$$
\underline{\alpha}'' = \underline{u}'(\underline{\alpha}) \underline{\alpha}' \quad , \quad \underline{\beta}'' = \underline{u}'(\underline{\beta}) \underline{\beta}'
$$
, i.e.\ les relations
$$
\underline{u}'(\underline{\alpha}) . \underline{\alpha}' = (\underline{u}' \underline{u}) (\underline{a}' \underline{a})^{-1} \qquad \underline{u}'(\underline{\beta}) \underline{\beta}' = (\underline{u}' \underline{u}) (\underline{b}' \underline{b})^{-1} \leqno (30)
$$
où $\underline{\alpha}, \underline{\beta}, \underline{\alpha}', \underline{\beta}'$ sont donnés en fonction de
$\underline{u}, \underline{a}, \underline{b}, \underline{u}', \underline{a}', \underline{b}'$ par (29). On a
$$
\underline{u}'(\underline{\alpha}) \underline{\alpha}' = \underline{u}' (\underline{u} \underline{a}^{-1}) \underline{u}'^{-1} \underline{u}' \underline{a}'^{-1} = (\underline{u}' \underline{u}) (\underline{a}^{-1} \underline{a}'^{-1}) = (\underline{u}' \underline{u}) (\underline{a}' \underline{a})^{-1} ,
$$
ce qui prouve la 1ère des deux relations 30),
et idem pour la 2e , ce qui prouve notre assertion.

\newpage

%% ===== Page 387 =====

Notons à cette occasion que les systèmes
$$
(\alpha, \beta, 1, 1) \quad \text{avec } \alpha \in \{\pm 1\}, \beta \in \{\pm 1\} \leqno{(31)}
$$
forment en fait un s-groupe de $\mathcal{S}\mathcal{G}_{p,\sigma}$ (puisqu'ils se relèvent en fait
par multiplicateurs, [unclear: disons] $(\alpha, \beta, \mu, \varepsilon)$,
d'où $\alpha$ et $\beta$ dégénèrent par une section de $p$
ce qui en [unclear: découle] pour $\bar{\alpha}, \bar{\beta}$ ...) ils correspondent,
dans, par l'isomorphisme $\mathcal{G}_{p,\sigma}^* \simeq B'_0 \times T_{\underline{p}} \times T_0$
aux systèmes $(u, \underline{a}, b)$ avec $u=1$, $\underline{a} \in \{\pm 1\}$, $b \in \{\pm 1\}$
et font un s-groupe (avec la multiplication ordinaire des
composantes comme loi de groupe) - où $\alpha, \beta$
sont en effet, [unclear: par définition], les signes de $\underline{a}, b$)
et [unclear: est] section du s-groupe
dit "du retour trivial" de $\mathcal{S}\mathcal{G}_{p,\sigma}$. L'action
par translation de ce s-groupe [unclear: $\simeq \mathbf{Z}/2\mathbf{Z} \times \mathbf{Z}/2\mathbf{Z}$]
sur $\mathcal{G}_{p,\sigma}^*$ se fait par multiplication des
deux composantes en question. Cela montre
que le fait de travailler avec des
$\alpha, \beta$ définis seulement "au signe près"
n'est nullement gênant : au vu des
groupes quotients de $\mathcal{G}_{p,\sigma}^*(\hat{\mathbf{Z}})$ (ou de $\mathcal{S}\mathcal{G}_{p,\sigma}(\hat{\mathbf{Z}})$)
par le s-groupe $Z_{00} \simeq \mathbf{Z}/2\mathbf{Z} \times \mathbf{Z}/2\mathbf{Z}$ engendré par les
éléments $(\alpha, \beta, 1, 1)$ avec $\alpha, \beta \in \{\pm 1\}$. Ceci

\newpage

%% ===== Page 388 =====

précise en même temps le sens de la locution :
[unclear: vrai] hom. de groupes ....

Revenons maintenant à l'application (23)
avec $f$ quelconque, et précisons la loi de
composition de $S_0\mathcal{M}_{0,3}^{!\sim}$ en termes de
$(\alpha,\beta,f)$ i.e. en termes de $u_{\alpha,\beta,f}$. Considérons
le sous-groupe formé des $(1,1,f)$, ou encore
ceux qui correspondent aux couples $(u,f)$ avec
$u = \operatorname{int}(f)$, i.e. aux couples $(\operatorname{int}(f), f)$
— ils forment un ``groupe isomorphe à $\hat{\Pi}_{0,3}$''
et on aura une structure de produit semi-direct
$$
S_0\mathcal{M}_{0,3}^{!\sim} \simeq \mathcal{M}_{0,3}^{!\sim}[0] \cdot \hat{\Pi}_{0,3} \leqno{(31)}
$$

$$
\begin{cases}
\hat{\Pi}_{0,3} \longrightarrow S_0\mathcal{M}_{0,3}^{!\sim} \\
f \longmapsto (\operatorname{int}(f), f) \leftrightarrow (1,1,f^{-1}),
\end{cases} \leqno{(32)}
$$
l'opération de $\mathcal{M}_{0,3}^{!\sim}[0]$ sur $\hat{\Pi}_{0,3}$ étant
l'opération tautologique (déduite de
l'inclusion $\mathcal{M}_{0,3}^{!\sim}[0] \subset \mathcal{M}_{0,3}^{!\sim} \subset \operatorname{Aut}(\hat{\Pi}_{0,3})$).
Le sous-groupe $\hat{\mathbf{Z}}$ de $S_0\mathcal{M}_{0,3}^{!\sim}$ est
donné par les couples $(u,f) = (1, b^n) \quad (n \in \hat{\mathbf{Z}})$
qu'on met sous la forme $(\alpha, \beta, f)$
[unclear: s.p. soit tel que] $u_{\alpha,\beta} = f u = \operatorname{int}(b^n)$,
donc $\alpha = \beta = b^n$, i.e.
$$
\hat{\mathbf{Z}} \to S_0\mathcal{M}_{0,3}^{!\sim} \text{ s'identifie à } n \mapsto (\alpha=b^n, \beta=b^n, f=b^n) \leqno{(33)}
$$

\newpage

%% ===== Page 389 =====

On aura
$$
(f'^{-1} u_{\alpha',\beta'} ) (f^{-1} u_{\alpha,\beta}) = (f'^{-1} u_{\alpha',\beta'}(f^{-1})) (u_{\alpha',\beta'} u_{\alpha,\beta})
$$
où la multiplication dans $S_0 \mathcal{M}_{0,3}^{\sim}$ se fait, en termes des composantes $(\alpha',\beta',f') (\alpha,\beta,f)$, par $(\alpha'',\beta'',f'')$, avec
$$
\begin{cases}
u_{\alpha'',\beta''} = u_{\alpha',\beta'} u_{\alpha,\beta} \quad \text{i.e.} \quad \begin{cases} \alpha'' = u_{\alpha',\beta'}(\alpha)\alpha' \\ \beta'' = u_{\alpha',\beta'}(\beta)\beta' \end{cases} , \quad \text{et} \\
f'' = u_{\alpha',\beta'}(f) f'
\end{cases} \leqno (34)
$$

Revenons maintenant à la définition de $S \mathcal{M}_{g,\nu}$ (1), en vue d'y introduire la multiplication convenable, ce qui revient à définir $S \mathcal{M}_{g,\nu}$ comme un produit $1/2$ direct
$$
S \mathcal{M}_{g,\nu} \defeq S \mathcal{G}_{g,\nu}^* \cdot S G \leqno (35)
$$
(produit $1/2$ direct)
par l'opération de $S \mathcal{G}_{g,\nu}^*$ sur $G$
via l'homomorphisme $\Theta$ (p.\ 489),
et on définit l'isomorphisme
$$
\begin{tikzcd}
S \mathcal{M}_{g,\nu} \arrow[r] & S \mathcal{G}_{g,\nu}^* \times_s G \\
U \cdot f \arrow[r, maps to] & (U, f^{-1})
\end{tikzcd} \leqno (36)
$$
\marginpar{attention au passage à $f^{-1}$ !}
où $U = (\alpha,\beta,\mu,\varepsilon)$ est dans $S \mathcal{G}_{g,\nu}^*$.

\newpage

%% ===== Page 390 =====

On a donc fait ce qu'il fallait, dans les
cas $S = \operatorname{Spec}(\hat{\mathbf{Z}})$ etc., pour avoir un
homomorphisme de groupes
$$
S_0 M_{0,3}^{1\sim} \longrightarrow S\mathcal{M}_{p,\sigma}(\hat{\mathbf{Z}}) / \widetilde{\Sigma} \leqno{(37)}
$$
où $\Sigma \simeq \mathbf{Z}/2 \times \mathbf{Z}/2$ est le sous-groupe [unclear: strict]
de $S\mathcal{G}_{p,\sigma}^*$ introduit p.~494', qui opère
trivialement sur $G$ donc est central dans
$S\mathcal{M}_{p,\sigma}$, donc dans $S\mathcal{M}_{p,\sigma}(\hat{\mathbf{Z}})$, et où
$$
\widetilde{\Sigma} = \Sigma \times \{ \pm 1 \}\footnote{$=$ automorphisme identique $=$ centre de $S\operatorname{Gl}(2,\mathbf{Z})$, centre de $G(\hat{\mathbf{Z}})$} \simeq \mathbf{Z}/2 \times \mathbf{Z}/2 \times \mathbf{Z}/2 \leqno{(38)}
$$
est le groupe distingué correspondant
dans le produit $1/2$ direct -- le fait
de passer au quotient par $\widetilde{\Sigma}$ exprime
les trois ambiguïtés de signes, sur $\alpha$, sur $\beta$
et sur $f$.

Considérons maintenant l'homomorphisme
$$
\begin{tikzcd}
\Gamma_{\mathbf{Q}} \arrow[r] & S_0 M_{0,3}^{1\sim} \\
\lambda \arrow[r, maps to] & (\alpha = c(\lambda), \beta = c(\lambda), \mu = 1, \varepsilon = 1, f = \varepsilon_0(\lambda))
\end{tikzcd} \leqno{(39)}
$$
C'est un homomorphisme de groupes, et fait
de $\Gamma_{\mathbf{Q}}$ un sous-groupe distingué.
Pour ceci le plus commode est d'identifier
$S_0 M_{0,3}^{1\sim}$ au produit $1/2$-direct (35),

\newpage

%% ===== Page 391 =====

et on identifie $S\mathcal{G}_{\vec{p}, \vec{\sigma}}^\ast$ à un sous-groupe de $B_0^\ast \times T_{\vec{p}} \times N_{\vec{\sigma}}^\ast$, opérant sur $G$ par autom. int. via $B_0^\ast$.

Le morphisme précité correspond donc à :
$$ \underset{\in B_0^\ast}{\lambda} \mapsto (\lambda, 1, 1, \underset{\in G}{\lambda^{-1}}) = \lambda \cdot j(\lambda)\footnote{$(j(\lambda) \in S\mathcal{M}_{\vec{p},\vec{\sigma}}^\ast \text{ est l'élém. du groupe de ramification})$} \leqno{(40)} $$

On identifie $S\mathcal{G}_{\vec{p}, \vec{\sigma}}^\ast$ à un sous-groupe de $\mathcal{G}_{\vec{p}, \vec{\sigma}}^\ast$, et de ce fait le produit 1/2 direct $S\mathcal{M}_{\vec{p}, \vec{\sigma}}^\ast$ à un sous-groupe du produit 1/2 direct
$$ \tilde{S}\mathcal{M}_{\vec{p}, \vec{\sigma}}^\ast = \mathcal{G}_{\vec{p}, \vec{\sigma}}^\ast \cdot G \simeq (B_0^\ast \times T_{\vec{p}} \times N_{\vec{\sigma}}^\ast) \cdot G \leqno{(41)} $$

et on voit que l'image du mono $i$ (40) dans $\mathcal{G}_{\vec{p}, \vec{\sigma}}^\ast \cdot G$ est distinguée dans ce groupe plus grand. Pour un peu plus d'aisance dans les calculs, on désigne par $i_{B_0^\ast}, i_{\vec{p}}, i_{\vec{\sigma}}, i_G$ les morphismes d'inclusion
$$
\left.\begin{matrix}
B_0^\ast \to \\
T_{\vec{p}} \to \\
N_{\vec{\sigma}}^\ast \to \\
G \to
\end{matrix}\right\} \tilde{S}\mathcal{M}_{\vec{p}, \vec{\sigma}}^\ast \leqno{(42)}
$$

\newpage

%% ===== Page 392 =====

il est entendu, bien sûr, qu'on fait opérer
$G_{p,\sigma}^*$ sur $G$ intérieurement via le
facteur $B_0'$. Il faut vérifier que
$$ i_G(L_0) = \{ i_B(\lambda) i_G(\lambda) \mid \lambda \in \Gamma_{B_0} \} $$
est normalisé par \marginpar{dire tout de suite que $i(L_0)$ est invariant dans tout ...} $i_B(x)$, $i_p(x)$, $i_\sigma(x)$,
$i_B(x)$ et $i_G(x)$. Pour les deux premiers,
c'est trivial car action triviale sur $L_0$.
Si $x \in \Gamma_{B_0'}$, on aura
$$
\begin{aligned}
i_B(x)(i(\lambda)) &= i_B(x) \left(i_B(\lambda) i_G(\lambda^{-1})\right) = \\
&= \underset{i_B(x\lambda x^{-1})}{i_B(x)(i_B(\lambda))} \cdot \underset{\underset{= i_G(x\lambda^{-1}x^{-1})}{i_G(x)(i_G(\lambda^{-1}))}}{i_B(x)(i_G(\lambda^{-1}))} \\
&= i(x\lambda x^{-1})
\end{aligned}
$$
ce qui montre que $L_0 \hookrightarrow G_{p,\sigma}^* \cdot G = \widetilde{S} M_{p,\sigma}^*$
commute : l'action naturelle de $B_0'$ sur
les deux côtés (sur le deuxième via la
$i_B : B_0' \to \widetilde{S} M_{p,\sigma}^*$). Enfin
$$ i_G(x)(i(\lambda)) = i_G(x)(i_B(\lambda) i_G(\lambda^{-1})) \cdot i_G(x)^{-1} $$
or par définition, $i_B(\lambda)$ et $i_G(\lambda)$ opèrent
de la = façon sur $G$, donc $i_B(\lambda) i_G(\lambda^{-1})$
y opère trivialement, dans le dernier facteur
$i_G(x)^{-1}$ passe à travers, d'où le facteur $i(\lambda)$, on trouve
$$ i_G(x) i_G(x)^{-1} i(\lambda) = i(\lambda) $$
O.K.

\newpage

%% ===== Page 393 =====

\underline{Rmq} $\tilde{S}M_{\rho,\sigma}$ est un groupe plat de dim relative
$\underbrace{2}_{B'_0} + \underbrace{2}_{T_\rho} + \underbrace{2}_{N_0} + \underbrace{4}_{G} = 10$, $SM_{\rho,\sigma}$ est de dim relative
$4+3=7$, en passant au quotient par $i(L_0) \simeq G_a$
[unclear: crossed out text]
produit des groupes croisés $\tilde{M}_{\rho,\sigma}$, $M_{\rho,\sigma}$,
de dim. relative $9$ resp. $6$, de sorte qu'on
a un diagramme
$$
(43) \begin{tikzcd}
1 \ar[r] & G_a \ar[r, "\lambda \mapsto i(\lambda)"] \ar[d, equal] & S\tilde{M}_{\rho,\sigma} \ar[r] \ar[d, "sur"'] & \tilde{M}_{\rho,\sigma} \ar[r] \ar[d, "sur"'] & 1 \\
1 \ar[r] & G_a \ar[r] & SM_{\rho,\sigma} \ar[r] & M_{\rho,\sigma} \ar[r] & 1
\end{tikzcd}
$$
Le passage en quotient ne se fait pas sans
difficulté - les quotients sont en effet
représentables comme des produits 1/2 directs
$$
(44) \begin{cases} 
\tilde{S}M_{\rho,\sigma} \simeq S\tilde{G}_{\rho,\sigma}^* \ltimes SG \\ 
\tilde{M}_{\rho,\sigma} \simeq \tilde{G}_{\rho,\sigma}^* \ltimes G 
\end{cases}
$$
où les sous-schémas en groupes
$S\tilde{G}_{\rho,\sigma}^*$ et $\tilde{G}_{\rho,\sigma}^*$ de $S\tilde{G}_{\rho,\sigma}$, $\tilde{G}_{\rho,\sigma}$
sont définis dans (VIII), $\tilde{G}_{\rho,\sigma}^*$ s'identifiant
à $T_0 \times T_\rho \times N'_0$, et $S\tilde{G}_{\rho,\sigma}^*$ à son intersection
avec $S\tilde{G}_{\rho,\sigma}$ (défini par les conditions
simples (20) de p. 487!).

\newpage

%% ===== Page 394 =====

478

[MARGIN: compte tenu de (13)]
Revenant maintenant à la situation arithmétique (par passage au quotient via B)), on trouve un isomorphisme canonique

(45) $$ \mathcal{M}_{0,3}^{\prime \sim} \overset{\sim}{\longrightarrow} M_{\rho,\sigma}(\hat{\mathbb{Z}})/\Sigma $$
[unclear: crossed out text]
opérant sur $\operatorname{Sl}(2,\hat{\mathbb{Z}})$ via le facteur $T_{\sigma}$.

$\Sigma$ étant le sous-groupe $\simeq \{\pm 1\}^3$, produit de trois facteurs $\mathbb{Z}/2\mathbb{Z}$ inclus resp.t dans $T_{\rho}(\hat{\mathbb{Z}})$, $T_{\sigma}(\hat{\mathbb{Z}})$ et $\operatorname{Sl}(2,\hat{\mathbb{Z}})$.

Remq. Cet homomorphisme a été obtenu [unclear: crossed out text] ici, à l'aide de l'hom. de $\mathcal{M}_{0,3}^{\prime \sim}$ (ou ?)

(46) $$ \mathcal{M}_{0,3}^{\sim}(\infty) = \operatorname{Norm}_{\mathcal{M}_{0,3}^{\prime \sim}} (L_0) $$
[MARGIN: sous-groupe des autom. "unimodulaires"]
[MARGIN: (cf §46, V)]

\begin{tikzcd}
& \searrow \\
& S\mathcal{G}_{\rho,\sigma}^{**}(\hat{\mathbb{Z}})/\Sigma_0 \ar[r, hook] & S\mathcal{G}_{\rho,\sigma}^{*}(\hat{\mathbb{Z}})/\Sigma_0
\end{tikzcd}

et du scindage

$$ \hat{\Pi}_{0,3} \hookrightarrow \hat{\Gamma}_{0,3} \twoheadrightarrow \operatorname{Sl}(2,\hat{\mathbb{Z}})^{\sim} $$
[MARGIN: $\simeq \operatorname{Sl}(2,\mathbb{Z})^{\wedge}$]

en utilisant

(47) $$ \mathcal{M}_{0,3}^{\sim} \simeq \mathcal{M}_{0,3}^{\prime \sim} \cdot \hat{\Pi}_{0,3} $$

\newpage

%% ===== Page 395 =====

Mais on a - ici - un produit semi-direct
$$
\mathcal{M}_{0,3}^{\sim} \simeq \mathcal{M}_{0,3}^{!\sim}(\omega) \cdot \widehat{\mathfrak{S}}_{0,3}^+ \leqno{(48)}
$$
qui, combiné directement à
$$
\widehat{\mathfrak{S}}_{0,3}^+ \longrightarrow \operatorname{Sl}(2, \hat{\mathbf{Z}})'
$$
donne un homomorphisme
$$
\mathcal{M}_{0,3}^{\sim} \longrightarrow \big( \mathcal{M}_{\rho, \sigma}(\hat{\mathbf{Z}}) \simeq S\mathcal{G}_{\rho, \sigma}^*(\hat{\mathbf{Z}}) \cdot \operatorname{Sl}(2, \hat{\mathbf{Z}}) \big) \big/ \Sigma \leqno{(49)}
$$
qui prolonge l'homomorphisme (45), défini d'abord sur son sous-groupe distingué $\mathcal{M}_{0,3}^{!\sim}$ d'indice 6.

\newpage

%% ===== Page 396 =====

\section*{X \quad Étude du groupe $\mathcal{M}_{p,\sigma}$ et de $\mathcal{M}_{p,\sigma} \to \operatorname{Gl}(2)$}

Par construction des produits semi-directs
$$
\begin{cases}
\widetilde{\mathcal{M}}_{p,\sigma} = \mathcal{G}_{p,\sigma}^{*\#}(\omega) . G \simeq \underset{\text{via } T_0 \hookrightarrow G}{\underset{\text{opère dans } G}{\widetilde{(T_0 \times T_p \times N'_0)}}} . G \\
\mathcal{M}_{p,\sigma} = S\mathcal{G}_{p,\sigma}^{*\#}(\omega) . SG \hookrightarrow \underset{\text{via } N'_0 \xrightarrow{\det} \{\pm 1\} \xrightarrow{\sim} T_0 \hookrightarrow G}{\underset{\text{opère sur } G}{\widetilde{(T_p \times N'_0)}}} . G
\end{cases} \leqno (1)
$$

on a un homomorphisme de groupes
$$
\begin{cases}
\widetilde{\mathcal{M}}_{p,\sigma} \overset{\Theta}{\longrightarrow} G \\
\Theta( i_B(x) i_p(y) i_\sigma(z) i_G(f) ) = x . f \\
( x \in \widetilde{T}_0 , y \in \widetilde{T}_p , z \in \widetilde{N}'_0 , f \in G ) .
\end{cases} \leqno (2)
$$

L'homomorphisme composé
$$
\underset{\simeq \mathcal{G}_{p,\sigma}^{*\#} . G}{S\widetilde{\mathcal{M}}_{p,\sigma}} \longrightarrow \widetilde{\mathcal{M}}_{p,\sigma} \overset{\Theta}{\longrightarrow} G \qquad \text{([unclear: composante] intérieure } \Theta) \leqno (3)
$$
est décrit par
$$
\Theta( i_B(x) i_p(y) i_\sigma(z) i_G(f) ) = x . f
$$
et $\Theta$ s'annule trivialement pour les
$$
i(\lambda) = i_B(\lambda) i_G(\lambda^{-1}) \quad (\lambda \in \widetilde{T}_0)
$$
et passe donc au quotient en $\Theta: \mathcal{M}_{p,\sigma} \to G$.

\newpage

%% ===== Page 397 =====

Je m'intéresse surtout à l'hom. induit

(4) $$M_{\rho,\sigma} \overset{\Theta}{\longrightarrow} G \simeq \widehat{Gl}(2)_S$$

il est évident épi, car son image contient
trivialement $SG$, donc il suffit de voir
que le composé avec
$$G/SG \overset{\det}{\underset{\sim}{\longrightarrow}} \widehat{\mathbb{G}}_m$$
est épi, i.e. que le composé sus-dit l'est.
Or ce composé est aussi l'hom. composé

(5)
\begin{tikzcd}
M_{\rho,\sigma} \ar[r, "\text{projection}", "\text{canonique}"'] \ar[d, "\simeq"'] & S\mathcal{G}^{*4}_{\rho,\sigma(\varpi)} \ar[r, "*4 \text{ multiplication}"] & \widehat{\mathbb{G}}_m \\
S\mathcal{G}_{\rho,\sigma(\varpi)}^{*4}. SG & (\alpha, \beta, \mu, \varepsilon) \ar[r, mapsto] & \mu
\end{tikzcd}

qui est épi — en termes de

(6) $$S\mathcal{G}^{*4}_{\rho,\sigma(\varpi)} \simeq \text{[unclear: crossed out text]}$$
ce qui $(\underline{a}, \underline{b}) \mapsto \det \underline{b}$

en termes de

(7) $$S\mathcal{G}^{*4}_{\rho,\sigma(\varpi)} \simeq T_\rho \times N'_0$$
ce sera $(\underline{a}, \underline{b}) \mapsto \det \underline{b}$

(cf VIII, p. 189 en particulier (25)).

Définition

(8) $$M^+_{\rho,\sigma} \overset{df}{=} \operatorname{Ker}(M_{\rho,\sigma} \overset{\Phi}{\longrightarrow} \widehat{\mathbb{G}}_m) \quad ,$$

où $\Phi$ désigne l'hom. de surjection [unclear: considéré].

\newpage

%% ===== Page 398 =====

(9)
\begin{tikzcd}
1 \ar[r] & M_{\rho,\sigma}^+ \ar[r] \ar[d] & M_{\rho,\sigma} \ar[r] \ar[d, "\Theta"] & \mathbb{G}_m \ar[r] \ar[d, equal] & 1 \\
1 \ar[r] & SG \ar[r] & G \ar[r, "\det"] & \mathbb{G}_m \ar[r] & 1
\end{tikzcd}

qu'on rapprochera de §39 p.299 (1)

Comme $M_{\rho,\sigma}$ est plat de dim relative 6, et $G$ lisse de dim relative 4, le noyau de l'épi $\Theta : M_{\rho,\sigma} \to G$ est nécessairement plat de dim relative 2. Notons-le $I_{\rho,\sigma}$, il est contenu dans le noyau $J_{\rho,\sigma}$ (de dim relative 3) du composé
$$M_{\rho,\sigma} \to G \to PG \ ,$$
de sorte qu'on a
(10)
$$1 \to I_{\rho,\sigma} \to J_{\rho,\sigma} \to \mathbb{G}_m \to 1$$
et une suite de composition de
$M_{\rho,\sigma}$

(11)
\begin{tikzcd}
1 \ar[r, phantom, "\subset"] & I_{\rho,\sigma} \ar[r, phantom, "\subset"] & J_{\rho,\sigma} \ar[r, phantom, "\subset"] & M_{\rho,\sigma} \\
& & & PG \simeq SG/\mu_2 \\
& & G &
\end{tikzcd}

Le groupe $J_{\rho,\sigma}$ joue le rôle d'un radical - on prévoit qu'il est commutatif, [unclear: à isogénie près], (de dim relative 3) - et torique si...

\newpage

%% ===== Page 399 =====

$T_\rho^0$, $T_\sigma^0$ sont des tores.
Mais je vois que je me complique
l'existence pour être trop savant dans mes
déductions. En effet, par construction
du produit $\frac{1}{2}$ direct, on a :
$$M_{\rho,\sigma} = I\!I_{\rho,\sigma}^a . SG \subset (T_\sigma \times T_\rho \times N'_\sigma) . SG$$
$$\simeq (T_\rho \times N'_\sigma) \times_G (T_\sigma . SG)$$

ainsi

(12)
\begin{tikzcd}
M_{\rho,\sigma} \ar[r, hook] \ar[dr, "\Theta"'] & T_\rho \times N'_\sigma \times G \ar[d, "pr_3"] \\
& G
\end{tikzcd}

d'ailleurs $M_{\rho,\sigma}$ est défini dans ce
produit par les conditions

(13) $$M_{\rho,\sigma} \subset T_\rho \times N'_\sigma \times G$$
$$ \simeq \{ (x,y,f) \mid \det f = \det y = \varepsilon' \det x , \; {\varepsilon'}^2 = 1 \}$$

Cela montre que

(14) $$I\!I_{\rho,\sigma} \subset T_\rho \times N'_\sigma \times \mathbb{G}_m$$
$$ \simeq \{ (x,y,\lambda) \mid \det y = \varepsilon \det x = \lambda^2 , \; \varepsilon^2 = 1 \}$$

$$J_{\rho,\sigma} \subset T_\rho \times {N'_\sigma}^1$$
$$ \simeq \{ (x,y) \mid \det y = 1 , (\det x)^2 = 1 \}$$

Il me semble plus naturel de
faire le dévissage de $M_{\rho,\sigma}$ en commençant
par descendre à sa composante neutre $M_{\rho,\sigma}^0$,

\newpage

%% ===== Page 400 =====

Le groupe \marginpar{[unclear: de composantes]}
$$
\mathcal{M}_{p,\sigma} / \mathcal{M}_{p,\sigma}^\circ \isom \mathcal{G}_{p,\sigma}^*(0) / \mathcal{G}_{p,\sigma}^*(0)^\circ = \mu_2 \times \mu_2 \leqno{(15)}
$$
(cf p. 489 (23)), on aura donc
$$
\begin{array}{ccc}
\mathcal{M}_{p,\sigma}^\circ = \mathcal{G}_{p,\sigma}^*(0)^\circ \cdot SG & \subset & (T_0 \times T_p \times T_\sigma) \cdot SG \\
& & \rotatebox{90}{$\simeq$} \\
& & T_p \times T_\sigma \times \underbrace{(T_0 \cdot SG)}_{\textstyle\rotatebox{90}{$\simeq$}\atop\textstyle G} \\
& & \rotatebox{90}{$\simeq$} \\
& & T_p \times T_\sigma \times G
\end{array}
$$

$$
\begin{array}{ccl}
\mathcal{M}_{p,\sigma}^\circ & \subset & T_p \times T_\sigma \times G \\
\rotatebox{90}{$=$} & & \\
\multicolumn{3}{l}{\big\{ (\underset{\begin{subarray}{c}\in\\T_p\end{subarray}}{x}, \underset{\begin{subarray}{c}\in\\T_\sigma\end{subarray}}{y}, \underset{\begin{subarray}{c}\in\\G\end{subarray}}{f}) \mid \det x = \det y = \det f \big\}}
\end{array} \leqno{(16)}
$$

Donc le noyau
$$
\begin{array}{ccl}
\mathcal{I}_{p,\sigma}^\circ & \defeq & \Ker \big( \mathcal{M}_{p,\sigma}^\circ \to \underset{\begin{subarray}{c}\rotatebox{90}{$\simeq$}\\G/\mathbb{G}_m\end{subarray}}{P G} \big) \\
\rotatebox{90}{$\simeq$} & & \\
\multicolumn{3}{l}{ \big\{ (\underset{\begin{subarray}{c}\in\\T_p\end{subarray}}{x}, \underset{\begin{subarray}{c}\in\\T_\sigma\end{subarray}}{y}, \underset{\begin{subarray}{c}\in\\mathbb{G}_m\end{subarray}}{\lambda}) \mid \det x = \det y = \lambda^2 \big\} } \\
\rotatebox{90}{$\simeq$} & & \\
\multicolumn{3}{l}{ T_p \times_{\mathbb{G}_m} T_\sigma = \{ (x,y) \mid \det x = \det y \} } \\
\rotatebox{90}{$\simeq$} & & \\
S\mathcal{G}_{p,\sigma}(0)^\circ & &
\end{array} \leqno{(17)}
$$

Enfin, on en déduit que
$$
\begin{array}{ccl}
\mathcal{J}_{p,\sigma}^\circ & \defeq & \Ker \big( \mathcal{M}_{p,\sigma}^\circ \to G \big) \\
\rotatebox{90}{$\simeq$} & & \\
\multicolumn{3}{l}{ ST_p \times ST_\sigma = \big\{ (\underset{\begin{subarray}{c}\in\\T_p\end{subarray}}{x}, \underset{\begin{subarray}{c}\in\\T_\sigma\end{subarray}}{y}) \mid \det x = \det y = 1 \big\} }
\end{array} \leqno{(18)}
$$

Récapitulons :

\newpage

%% ===== Page 401 =====

Théorème. Le groupe
(19) $$M_{\rho,\sigma} \overset{def}{=} G_{\rho,\sigma}^*(0) \cdot SG$$
[MARGIN: $\uparrow$ cf définition p. 487' (30)]
est extension de $\mu_2 \times \mu_2$ par $M_{\rho,\sigma}^\circ$

(20) $$1 \to M_{\rho,\sigma}^\circ \to M_{\rho,\sigma} \to \mu_2 \times \mu_2 \to 1$$
déduit de $G_{\rho,\sigma}^*(0) \to \mu_2 \times \mu_2$
$(\alpha, \beta, \mu, \varepsilon) \mapsto (\det \alpha, \varepsilon \det \beta)$

Dans $M_{\rho,\sigma}^\circ$ on a
[unclear: remarquable : le sous-groupe central (de $M_{\rho,\sigma}^\circ$ + action)]
(21) $$ST_\rho \times ST_\sigma \hookrightarrow G_{\rho,\sigma}^\circ(0) \hookrightarrow M_{\rho,\sigma}^\circ \quad ,$$
et on a une suite exacte
(22)
\begin{tikzcd}
1 \ar[r] & ST_\rho \times ST_\sigma \ar[r] & M_{\rho,\sigma}^\circ \ar[r, "\Theta"] & G \ar[r] & 1 \\
& & SG \ar[u, hook] \ar[ur, "\simeq"'] & & 
\end{tikzcd}

Comme [unclear: dans $M_{\rho,\sigma}^\circ$ on a]
[unclear: $ST_\rho \times ST_\sigma \times SG \subset J_{\rho,\sigma}^\circ \times SG$]
[MARGIN: NB $J_{\rho,\sigma}^\circ$ est encore central dans $M_{\rho,\sigma}^\circ$ et $J_{\rho,\sigma}^\circ \cap SG = 1$]

(23) $$1 \to ST_\rho \times ST_\sigma \times SG \to M_{\rho,\sigma} \to \mathbb{G}_m \times \mu_2 \times \mu_2 \to 1$$
$x \mapsto (\chi(x), \varepsilon(x), \varepsilon'(x))$

(24) $$M_{\rho,\sigma}^{0+} \simeq ST_\rho \times ST_\sigma \times SG$$
provenant de la
section ($G \hookrightarrow M_{\rho,\sigma}^\circ$)
[unclear: $ST_\rho \times ST_\sigma \times G \to M_{\rho,\sigma}^\circ \to 1$]

On a un hom. canonique
(25) $$\mu_2 \times \mu_2 \times \mu_2 \hookrightarrow M_{\rho,\sigma}^{0+} \overset{def}{=} \operatorname{Ker}(M_{\rho,\sigma}^\circ \to \mathbb{G}_m)$$
$\simeq$ (24)
$$ST_\rho \times ST_\sigma \times SG$$

\newpage

%% ===== Page 402 =====

grâce aux inclusions canoniques
$$
\mu_2 \hookrightarrow S\mathcal{T}_p, \quad \mu_2 \hookrightarrow S\mathcal{T}_\sigma, \quad \mu_2 \hookrightarrow SG \leqno (26)
$$
donnant naissance à [unclear: un s-groupe] commutatif d'exposant 2, [unclear: engendré] par ailleurs
$$
\Sigma \simeq (\mathbf{Z}/2\mathbf{Z})^3 \hookrightarrow \mathcal{M}_{\vec{p},\vec{\sigma}}^{0+} \subset \mathcal{M}_{\vec{p},\vec{\sigma}} \leqno (27)
$$

Je vais introduire au sujet des groupes $\mathcal{M}_{\vec{p},\vec{\sigma}}$ des notations qui paraphrasent les notations $\mathcal{M}_{0,3}$ etc, dans le contexte des schémas en groupes. Le sous-groupe invariant $SG$ de $\mathcal{M}_{\vec{p},\vec{\sigma}}$ sera ainsi noté $\mathfrak{S}_{\vec{p},\vec{\sigma}}^+$ (l'analogue de $\mathfrak{S}_{0,3}^+$), le quotient\footnote{pour [unclear: un peu moins patiner] à la rigueur ! (cf (28)) et tomber [unclear: d'accord] même dans la suite ... c'est $\mathcal{T}_{\vec{p},\vec{\sigma}}$ qu'il conviendrait de ... noter ce groupe ! (29)}
$$
\mathcal{M}_{\vec{p},\vec{\sigma}} / \mathfrak{S}_{\vec{p},\vec{\sigma}}^+ \simeq \mathcal{T}_{\vec{p},\vec{\sigma}}^! \leqno (28)
$$
$\mathcal{T}_{0,3}^!$, ou, d'ailleurs
$$
\mathcal{T}_{\vec{p},\vec{\sigma}}^! \simeq S\mathcal{G}_{\vec{p},\vec{\sigma}}^*(0) \hookrightarrow \mathcal{T}_0 \times \mathcal{T}_p \times N'_0 \leqno (29)
$$
--- par la suite la notation plus simple et plus suggestive $\mathcal{T}_{\vec{p},\vec{\sigma}}$ sera utilisée de préférence à $S\mathcal{G}_{\vec{p},\vec{\sigma}}^*(0)$. On a donc
$$
1 \longrightarrow \mathfrak{S}_{\vec{p},\vec{\sigma}}^+ \longrightarrow \mathcal{M}_{\vec{p},\vec{\sigma}} \longrightarrow \mathcal{T}_{\vec{p},\vec{\sigma}} \longrightarrow 1 \leqno (30)
$$
et un scindage canonique en produit semi-direct
$$
\mathcal{M}_{\vec{p},\vec{\sigma}} \simeq \mathcal{M}_{\vec{p},\vec{\sigma}}(\omega) \cdot \mathfrak{S}_{\vec{p},\vec{\sigma}}^+ \leqno (31)
$$

\newpage

%% ===== Page 403 =====

(32)
\begin{tikzcd}
M_{\rho,\sigma}(\infty) \ar[r, hook] \ar[dr, "\sim"'] & M_{\rho,\sigma} \ar[d, "can"] \\
& \tau_{\rho,\sigma}^!
\end{tikzcd}

est un sous-groupe isomorphe à $\tau_{\rho,\sigma}^!$
qui correspond au sous-groupe $M_{0,3}^{!\sim}(\infty)$ dans
$M_{0,3}^\sim$ ($\S$ V).

Rem. On aurait envie d'introduire une
$\tau_{\rho,\sigma}$ extension (un produit) de $\tau_{\rho,\sigma}^!$ par un groupe
tel que $\mathfrak{S}_3$ (qui serait noté $\tau_{\rho,\sigma}^+$), en
retenant un certain sous-groupe $\Pi_{\rho,\sigma}$ de
$\mathfrak{S}_{\rho,\sigma}^+$ tel que $\mathfrak{S}_{\rho,\sigma}^+ / \Pi_{\rho,\sigma} \simeq \mathfrak{S}_3$ -
mais ce n'est pas possible dans le contexte
des schémas en groupes (ni dans le cas
particulier modulaire), vu que $\mathfrak{S}_{\rho,\sigma}^+ \simeq \operatorname{Sl}(2)_{\mathbb{Z}}$
est [unclear: connexe] sur $S$. Par contre on
peut, dans le "cas modulaire" $\rho = \begin{pmatrix} 0 & 1 \\ -1 & 0 \end{pmatrix}, \sigma = \begin{pmatrix} 0 & 1 \\ -1 & 1 \end{pmatrix}$
introduire une section canonique

(33) $$t = \bar{\tau}_\infty \in \Gamma(\tau_{\rho,\sigma}^!)$$

dont le représentant $\tau_\infty$ dans $\Gamma M_{\rho,\sigma} (\infty)$ s'explicite
comme

(34) $$\tau_\infty = \begin{pmatrix} -1 & 0 \\ 0 & +1 \end{pmatrix}$$

[unclear: Il serait intéressant de voir si, en dehors des]

\newpage

%% ===== Page 404 =====

cas modulaire, on peut introduire des «~antiinvo-
lutions~» qui généralisent les réflexions géométriques
$$
\begin{cases}
\tau_0' = \tau_{0\infty}' = \begin{pmatrix} 1 & 0 \\ 0 & -1 \end{pmatrix} \\
\tau_1' = - \tau_0' = \begin{pmatrix} 0 & 1 \\ 1 & 0 \end{pmatrix} \\
\tau_\infty' = \tau_{\infty 0}' = \begin{pmatrix} -1 & 0 \\ 0 & 1 \end{pmatrix}
\end{cases}
\leqno(35)
$$
\marginpar{(cf. p. 335 (§ 42, IV))}

Je suis convaincu que oui. D'ailleurs n'avions-
nous pas mis en évidence dans $\mathcal{M}_{p,\sigma}$ un
sous-groupe $\mu_2 \times \mu_2 \times \mu_2$, correspondant d'ailleurs aux
trois sections remarquables (correspondant respect. aux
trois sections $-1$ des trois facteurs $\mu_2$)?
Mais je me mélange là les pinceaux - ces
sections étaient des sections de $\mathcal{M}_{p,\sigma}$ - tandis
que les sections $\tau_0', \tau_1', \tau_\infty'$ donnant des multiplications
par $-1$, donc correspondaient à $\varepsilon = \varepsilon' = -1$, donc correspon-
daient à la section «~le noyau trivial~» $(-1,-1)$
dans $\mathcal{M}_{p,\sigma} / \mathcal{M}_{p,\sigma}^\circ \simeq \mu_2 \times \mu_2$.

$$
\hat{\mathfrak{T}}_{p,\sigma}^! \simeq \hat{\mathcal{G}}_{p,\sigma}(0) \simeq T_0 \times T_p \times \mathcal{N}_0' \leqno(36)
$$
$$
\simeq \{ (u, \underline{a}, \underline{b}) \mid \det u = \det \underline{b} = \varepsilon' \det \underline{a} \; , \; \varepsilon'^2 = 1 \}
$$

et on rappelle les isomorphismes
$$ \delta_0 : u \mapsto \det u $$
$$ T_0 \isommap \hat{\mathbf{G}}_m $$
et des hom. canoniques «~positifs~»
$$ \hat{\mathbf{G}}_m \overset{j_p}{\longrightarrow} T_p $$
$$ \hat{\mathbf{G}}_m \overset{j_0}{\longrightarrow} T_0 $$
dont la composée avec $\delta_p, \delta_0$ (les hom. det) est $\lambda \mapsto \lambda^2$.
Ainsi dans la deuxième [unclear: ligne] de (36) il y a un

\newpage

%% ===== Page 405 =====

[unclear: sous-groupe de certains twists]
$$ \mathbb{G}_m \times \mathbb{G}_m \times \mathbb{G}_m \to T_0 \times T_p \times \mathcal{N}_0' $$
Faisons les calculs pour déterminer les invariants $\alpha, \beta$ \dots correspondants à $r_0, r_1, r_\infty$.

i) $$ r_\infty = \tau_\infty = \begin{pmatrix} -1 & 0 \\ 0 & 1 \end{pmatrix} \in \operatorname{Gl}(2, \mathbf{Z}) \subset \operatorname{Gl}(2, \hat{\mathbf{Z}})^{\pm 1} $$
On a ici $\tau_\infty \in T_0'(\mathbf{Z})$, c'est le seul élément non trivial de $T_0'(\mathbf{Z}) \simeq \mathbb{G}_m(\mathbf{Z}) \simeq \{\pm 1\}$. L'automorphisme extérieur qu'il définit dans $\mathfrak{T}_{0,3}$ ou dans $\operatorname{Sl}(2, \mathbf{Z})$ ou ses complétés profinis stabilise le sous-groupe $L_0$ engendré par $\varepsilon_0$.
(i.e.\ $L_0$ --- $\varepsilon_0^2 = l_0$) où $\tau_\infty(\varepsilon_0) = \varepsilon_0^{-1}$, $\tau_\infty(l_0) = l_0^{-1}$
donc
$$ \tau_\infty \in \mathcal{M}_{0,3}^{+ \wedge}[0] $$
Et en fait il est de la forme $u_{\alpha, \beta}$, avec des calculs déjà faits donnant
$$ \alpha = \tau_\infty, \quad \gamma = q = \alpha \beta \alpha^{-1} = - l_1^{-1} = \begin{pmatrix} 1 & 0 \\ -2 & 1 \end{pmatrix} $$
(attention aux signes --- quand je fais les calculs dans $\operatorname{Sl}(2, \mathbf{Z})$ il s'agit de $\operatorname{Sl}(2, \mathbf{Z})'$ !)

$$ \mu = -1, \quad \varepsilon = \varepsilon_2(\mu) = -1, \quad \varepsilon' = \varepsilon_3(\mu) = -1 $$
On a ici $B = 0, D = -1$, donc
$$ u = \begin{pmatrix} \mu & 0 \\ B/D = 0 & D \end{pmatrix} = \begin{pmatrix} -1 & 0 \\ 0 & 1 \end{pmatrix} = \tau_\infty $$
$$ \underline{a} = \alpha^{-1} u = \tau_\infty^{-1} \tau_\infty = 1, \quad \underline{b} = \beta^{-1} u = \tau_\infty $$
d'où \hfill (l'élément [unclear: d'ordre 2] dans $T_0(\mathbf{Z}) = {_2T_0}(\mathbf{Z})$)
$$ r_\infty = \tau_\infty = (\underline{u} = \tau_\infty, \underline{a} = 1, \underline{b} = \tau_\infty) \in T_0 \times T_p \times \mathcal{N}_0' \leqno{(36)} $$

\underline{NB} Soit plus généralement
$$ \underline{u} = \begin{pmatrix} \mu & \lambda \\ 0 & \nu \end{pmatrix} \in \operatorname{Gl}(2, \mathbf{Z}) \subset \operatorname{Gl}(2, \hat{\mathbf{Z}}) $$
et considérons l'automorphisme intérieur de $\operatorname{Sl}(2, \hat{\mathbf{Z}})'$ qu'il induit, d'où un autom.\ de $\widehat{\mathcal{M}}_{0,3}^+$ [unclear: opéré], quels en sont les invariants $\alpha, \beta$ etc.? On pourra se réduire

\newpage

%% ===== Page 406 =====

mieux prendre $f=1$, et on aura
$$ \alpha = u, \quad \text{d'où} \quad a = 1, $$
d'autre part, notons que $\mu \in \{\pm 1\}$ (i.e.\ $\mu = (-1)^{\det u}$, car $u \in \operatorname{Gl}(2,\mathbf{Z})$),
on aura $\varepsilon = \varepsilon' = \mu$, et si $\mu = 1$, on aura $\beta = u$,
(puisque alors $b=1$), sinon $\beta = \tau_\infty u \tau_\infty$, donc $b = \tau_\infty$,
dans tous les cas on aura $\beta = u^{\tau_\infty^k}$, $b = \tau_\infty^k$,
en résumé
$$ \underline{u} = (\underset{\in B'_0}{u}, \underset{\in \Gamma'_p}{a=1}, \underset{\in N'_0}{b=\tau_\infty^k}) \in \mathcal{M}_{p,\sigma,[0]}^*(\hat{\mathbf{Z}}) \leqno{(37)} $$

[unclear: Mais ce n'est pas un élément sous forme standardisée à la p. 11, il n'est pas dans $\Gamma_{p,\sigma}^{0}(0)$, sauf si $\lambda=0$.]

2°) $u = - \tau'_0 = \begin{pmatrix} -1 & 1 \\ 0 & 1 \end{pmatrix}$, c'est justement un cas
qui tombe sous les remarques précédentes, on a :\marginpar{NB réduire modulo le $\operatorname{Cent}_{\mathcal{E}_{p,\sigma}}$ de $\Gamma_1$ dans l'image de $(1, \tau_\infty)$.}
$$ \tau_1 = \tau'_0 = (\underline{u}_1 = -\tau'_0 = \begin{pmatrix} -1 & 1 \\ 0 & 1 \end{pmatrix}, a=1, b=\tau_\infty) \in \mathcal{M}_{p,\sigma,[0]}(\mathbf{Z}) \leqno{(38)} $$

3°) $\tau_0 \tau'_0 = \begin{pmatrix} 0 & 1 \\ 1 & 0 \end{pmatrix}$, cette fois-ci l'automorphisme [unclear: qu'il induit] ne stabilise pas $L_0$ -- ça vaut mieux, si on veut que $T_0, \Gamma_1, \Gamma_\infty$ engendrent $\operatorname{Gl}(2,\mathbf{Z})$ (ou "seulement" $\operatorname{Sl}(2,\mathbf{Z})$) !
On a $\tau'_0 = \tau_\infty \tau_\alpha \tau_\infty$, donc son image dans $\mathcal{M}_{p,\sigma,[0]}(\mathbf{Z}) \subset \mathcal{M}_{p,\sigma}(\mathbf{Z})$ est contrôlée produit des images de $\tau_\infty$ et $\tau_\alpha$ -- la première a été calculée dans 1°, la deuxième est tautologique puisque $\operatorname{Sl} \subset \mathcal{E}_{p,\sigma}^+$ et $\tau_\alpha \in \operatorname{Sl}(2,\mathbf{Z})$. On n'ira donc pas [unclear: plus loin] pour évaluer cette image.

\newpage

%% ===== Page 407 =====

dans $\mathcal{G}_{p,\sigma}(\mathbf{Z})$ pour $\tau_\infty = \tau_0$ et $\tau_1 = -\tau_0'$, d'où toujours le même élément
$$
(-1, 1, \tau_\infty) \in (\underset{\simeq \mathbf{G}_m}{T_0} \times T_p \times N_{\sigma'})(\mathbf{Z}) \leqno{(39)}
$$

Calculons d'ailleurs
$$
T_0 \times T_p \times N_{\sigma'}(\mathbf{Z}) = T_0(\mathbf{Z}) \times T_p(\mathbf{Z}) \times N_{\sigma'}(\mathbf{Z})
$$
On a :
\begin{align*}
T_0 &\simeq \mathbf{G}_m, \quad T_0(\mathbf{Z}) \simeq \{\pm 1\} \\
T_p(\mathbf{Z}) &\simeq (\mathbf{Z} + \mathbf{Z} p)^* \simeq \mathbf{Z}/6\mathbf{Z} \quad \text{groupe engendré par } \rho 
\end{align*}
$$
N_{\sigma'}(\mathbf{Z}) \simeq \{1, \tau_\infty\} \cdot T_{\sigma'}(\mathbf{Z}) \quad \text{(produit semi-direct)} \leqno{(40)}
$$
\begin{align*}
T_{\sigma'}(\mathbf{Z}) &\simeq (\mathbf{Z} + \mathbf{Z} \sigma')^* \simeq \mathbf{Z}/4\mathbf{Z} \quad \text{groupe engendré par } \sigma' \\
N_{\sigma'}(\mathbf{Z}) &\simeq D_4 \quad \text{groupe diédral d'ordre 4 ("symétries du carré")}
\end{align*}

d'où
$$
(T_0 \times T_p \times N_{\sigma'})(\mathbf{Z}) \simeq \underset{\substack{\text{groupe engendré} \\ \text{par } \tau_\infty}}{(\mathbf{Z}/2\mathbf{Z})} \cdot ( \mathbf{Z}/2\mathbf{Z} \times \mathbf{Z}/6\mathbf{Z} \times \mathbf{Z}/4\mathbf{Z} ) \leqno{(41)}
$$
$\tau_\infty$ opère trivialement sur les deux premiers facteurs, et par symétrie sur le dernier.
$$
\simeq \mathbf{Z}/2\mathbf{Z} \times \mathbf{Z}/6\mathbf{Z} \times \underset{\simeq D_4}{(\mathbf{Z}/2\mathbf{Z} \cdot \mathbf{Z}/4\mathbf{Z})}
$$

$$
S\mathcal{G}_{p,\sigma'[0]}(\mathbf{Z}) = \mathcal{G}_{p,\sigma'}^+(\mathbf{Z})\footnote{NB la condition $\det a = \varepsilon \det u, \varepsilon'=1$ est automatique car $\mathbf{G}_m(\mathbf{Z}) = \{\pm 1\}$.} = \{ (u,a,b) \in (T_0 \times T_p \times N_{\sigma'})(\mathbf{Z}) \mid \det u = \det b \} \leqno{(42)}
$$
$$
\simeq (T_p \times N_{\sigma'})(\mathbf{Z}) \simeq \mathbf{Z}/6\mathbf{Z} \times D_4
$$

\newpage

%% ===== Page 408 =====

l'élément remarquable, i.e.\ qui correspond à $\tau_0 = \tau_{\infty}$, $\tau_1 = -\tau'_0$, $\tau_{\infty} = \tau_{\infty}$ l'élément $(1, \tau_{\infty})$ de $\mathbf{Z}/6\mathbf{Z} \times (D_4 \simeq \{1, \tau_{\infty}\} \cdot \mathbf{Z}/4\mathbf{Z})$.

Ceci nous donne d'ailleurs une expression pour 
$$
\begin{cases}
\mathcal{M}_{p,\sigma}(\mathbf{Z}) \simeq \underbrace{\mathcal{T}_{p,\sigma}(\mathbf{Z})}_{\mathcal{M}_{p,\sigma(0)}(\mathbf{Z})} \cdot \operatorname{Sl}(2,\mathbf{Z}) \\
\text{(produit semi-direct)} \\
\simeq (\mathbf{Z}/6\mathbf{Z} \times D_4) \cdot \operatorname{Sl}(2,\mathbf{Z})
\end{cases} \leqno{(43)}
$$
où $\mathbf{Z}/6\mathbf{Z} \times D_4$ opère sur $\operatorname{Sl}(2,\mathbf{Z})$ via l'hom.
$$
\mathbf{Z}/6\mathbf{Z} \times D_4 \xrightarrow{\text{projection canonique}} D_4 \xrightarrow{\text{hom. can.}} \{1, \tau_{\infty}\} \leqno{(44)}
$$
$\tau_{\infty}$ opère sur $\operatorname{Sl}(2,\mathbf{Z})$ par autom. int. avec la matrice
$$
\tau_{\infty} = \begin{pmatrix} -1 & 0 \\ 0 & 1 \end{pmatrix}.
$$
On note que l'homom. canonique
$$
\mathcal{M}_{p,\sigma} \to \mu_2 \times \mu_2
$$
se factorise par l'homom $\mathcal{T}_{p,\sigma}$, et induit un hom. sur les points entiers
$$
\mathcal{M}_{p,\sigma}(\mathbf{Z}) \to \mathcal{T}_{p,\sigma}(\mathbf{Z}) \simeq \mathbf{Z}/6\mathbf{Z} \times D_4 \to (\mu_2 \times \mu_2)(\mathbf{Z}) \simeq \{\pm 1\} \times \{\pm 1\} \leqno{(45)}
$$
Cet hom. est nul sur $\mathbf{Z}/6\mathbf{Z} \times \mathbf{Z}/4\mathbf{Z}$ (où $\mathbf{Z}/4\mathbf{Z} \subset D_4$) (correspondant à des déterminants égaux à $1$), et sur $\tau_{\infty}$ prend la valeur $(-1, -1)$. [unclear: Ce fait] que cet hom. sur les pts entiers n'est

\newpage

%% ===== Page 409 =====

pas surjectif.

Centre de $\mathcal{M}_{p,\sigma}$\footnote{Faux, dans un produit $1/2$ direct le centre peut être plus grand que le produit des centres.} $= \mathcal{M}_{p,\sigma(0)} . SG$. C'est le produit du sous-groupe central de $\mathcal{M}_{p,\sigma(0)}$ (qui opère trivialement sur le centre de $SG$) et du centre de $SG$. Le 2ème est bien connu, c'est $\mu_2$, et $\mathcal{M}_{p,\sigma}$ y opère trivialement.

D'autre part le sous-groupe de $\mathcal{M}_{p,\sigma(0)} \simeq \mathfrak{T}_{p,\sigma} \simeq B'_0 \times T_p \times N'_0$ qui opère trivialement sur $SG$, est $B'_0 \times N'_0$, car $T_p$ opère fidèlement, (puisque
$$
T_p \cap \underset{\text{centre de } SG}{\mu_2} = \{1\} \leqno{(46)}
$$

C'est donc le sous-groupe des
$$
\left\{ (\underset{\in B'_0}{1}, \underset{\in T_p}{\underline{a}}, \underset{\in N'_0}{\underline{b}}) \;\middle|\; (\det \underline{a})^2 = 1, \det \underline{b} = 1 \right\}
$$

\marginpar{non !}Reste à déterminer le centre de $\mathfrak{T}_{p,\sigma}$, extension de $\mu_2 \times \mu_2$ par $\mathfrak{T}_{p,\sigma}^\circ \simeq \mathcal{G}_{p,\sigma(0)}^\circ \subset B'_0 \times T_p \times N'_0$. Comme $\mathfrak{T}_{p,\sigma}^\circ$ est commutatif, il suffit de déterminer l'action de $\mu_2 \times \mu_2$ sur $\mathfrak{T}_{p,\sigma}^\circ$, et d'en déduire les invariants.\marginpar{Cela donne au moins l'intersection du centre avec $\mathfrak{T}_{p,\sigma}^\circ$.}

\newpage

%% ===== Page 410 =====

505

Notons que l'on a
$$ \mathcal{T}_{\rho,\sigma} = S G_{\rho,\sigma}^*(0) \subset G_{\rho,\sigma}^*(0) \simeq T_0 \times T_\rho \times N'_{0'} $$
et un hom. de suites exactes
$$ ((u,a,b) \mapsto (\varepsilon(b), \det\underline{a}/\det(\underline{u}))) $$
(48)
\begin{tikzcd}
1 \ar[r] & \mathcal{T}_{\rho,\sigma} \ar[r] \ar[d] & G_{\rho,\sigma}^*(0) \ar[r] \ar[d] & \mu_2 \times \mu_2 \ar[r] \ar[d, "\text{première} \atop \text{proj.}"] & 1 \\
1 \ar[r] & T_0 \times T_\rho \times \bar{N}_0 \ar[r] & T_0 \times T_\rho \times N'_{0'} \ar[r] & \mu_2 \ar[r] & 1
\end{tikzcd}

qui montre que l'action de $\mu_2 \times \mu_2$ sur $\mathcal{T}_{\rho,\sigma}$ se fait par l'intermédiaire de la projection de $\mu_2 \times \mu_2$ sur son premier facteur, et par l'opération de celui-ci sur $T_0 \times T_\rho \times \bar{T}_0$ qui opère trivialement sur $T_0 \times T_\rho$, et qui sur $\bar{T}_0$ est l'opération canonique de $\mu_2$ sur $\bar{T}_0$, $\sigma \mapsto \varepsilon \sigma$, dont les invariants se réduisent aux scalaires $\mathbb{G}_m$. Donc ce raisonnement montre, plus généralement, que la trace du centre de l'ext. $G^*_{\rho,\sigma}$ sur le sous-groupe $\mathcal{T}'_{\rho,\sigma}$ noyau de 
$$ \mathcal{T}_{\rho,\sigma} \rightarrow \mu_2 \times \mu_2 \xrightarrow{1^{\text{ère}} \text{ proj.}} \mu_2 $$
$$ (u,a,b) \mapsto \varepsilon(b) $$
est aussi la trace, sur ce sous-groupe, de $T_0 \times T_\rho \times \mathbb{G}_m$. D'ailleurs, comme $\mu_2$ opère fidèlement sur $\bar{T}_0$ via $\sigma \mapsto \varepsilon \sigma$, on voit que l'on constate qu'il opère aussi fidèlement sur $\mathcal{T}_{\rho,\sigma}$ ([unclear: qui en a en effet s'envoie épimorphiquement] sur $\bar{T}_0$) donc le
$$ \text{Centre}(\mathcal{T}_{\rho,\sigma}) \subset \mathcal{T}'_{\rho,\sigma} = \operatorname{Ker}(\mathcal{T}_{\rho,\sigma} \xrightarrow{((u,a,b) \mapsto \varepsilon(b))} \mu_2) $$

\newpage

%% ===== Page 411 =====

Donc en résumé

$$
\begin{cases}
\operatorname{Centre}(G_{\rho, \sigma}^*(\omega)) = T_{\sigma} \times T_{\rho} \times \mathbb{G}_m \\
\operatorname{Centre}(\mathcal{T}_{\rho, \sigma}) = (T_{\sigma} \times T_{\rho} \times \mathbb{G}_m) \cap \mathcal{T}_{\rho, \sigma} \\
\qquad \qquad \qquad = \left\{ (\underset{\in T_{\sigma}}{u}, \underset{\in T_{\rho}}{\underline{a}}, \underset{\in \mathbb{G}_m}{\lambda}) \mid u = \lambda^2, \det \underline{a} = \varepsilon \lambda^4, \varepsilon^{12}=1 \right\}
\end{cases} \leqno (49)
$$

Ce sous-groupe du centre, qui opère trivialement sur $S G$ est donc, en vertu de (47), donné par la restriction $u=1$, il s'identifie donc à

$$ 1 \times S T_{\rho}^* \times \mu_2 \subset \mathcal{T}_{\rho, \sigma} \subset T_{\sigma} \times T_{\rho} \times N_{\sigma}' \leqno (50) $$

où $S T_{\rho}^* \subset T_{\rho}$
$= \{ \underline{a} \mid (\det \underline{a})^3 = 1 \}$

\footnote{Faux. le centre est plus grand,}
$$ \operatorname{Centre}(\mathcal{M}_{\rho, \sigma}) \simeq \underset{\subset \mathcal{T}_0}{(1 \times S T_{\rho}^* \times \mu_2)} \times \underset{\subset S G}{\mu_2} \leqno (51) $$

Le groupe des points entiers de ce groupe, dans le cas modulaire, est réduit à :
$$ \mathbf{Z}/6\mathbf{Z} \times \mathbf{Z}/2\mathbf{Z} \times \mathbf{Z}/2\mathbf{Z} $$

\marginpar{NB Ce groupe est en fait contenu dans $\mathcal{M}_{\rho,\sigma}(\mathbf{Z})$.}
$$ \simeq \mathbf{Z}/3\mathbf{Z} \times (\mathbf{Z}/2\mathbf{Z} \times \mathbf{Z}/2\mathbf{Z} \times \mathbf{Z}/2\mathbf{Z}) \leqno (52) $$

produit d'un groupe cyclique d'ordre 3 (engendré visiblement par $-\rho$) par un groupe $(\mathbf{Z}/2\mathbf{Z})^3$, qui n'est autre que le groupe $\Sigma$ déjà dans (18).

Plus intéressant serait d'étudier le

\newpage

%% ===== Page 412 =====

\marginpar{faux}centre de $\mathcal{M}_{\bar{p},\bar{\sigma}^\circ}$, qui sera donc le produit
direct du centre de $SG$ (sur lequel le groupe
$\mathcal{T}_{\bar{p},\bar{\sigma}^\circ}$ opère trivialement) par le ss-groupe
de $\mathcal{T}_{\bar{p},\bar{\sigma}^\circ} \subset T_0 \times T_{\bar{p}} \times T_\sigma$ qui opère
trivialement sur $S\ell_1$, savoir :
$$
\mathcal{T}_{\bar{p},\bar{\sigma}^\circ} \cap 1 \times T_{\bar{p}} \times T_\sigma = \{ (1,a,b) \mid \det b = 1, \det a = 1 \} \footnote{et pas seulement $(\det a)^2=1$, car $(1,a,b) \in \mathcal{T}_{\bar{p},\bar{\sigma}^\circ}$}
$$
$$
\simeq S T_{\bar{p}} \times S T_\sigma
$$

Donc
\marginpar{faux, le centre est de dim 3}
$$
\operatorname{Centr} (\mathcal{M}_{\bar{p},\bar{\sigma}^\circ}) = (1 \times S T_{\bar{p}} \times S T_{\sigma}) \times \underset{\subset SG}{\mu_2} \leqno{(53)}
$$

d'où dans le cas modulaire :
$$
\operatorname{Centr}(\mathcal{M}_{\bar{p},\bar{\sigma}^\circ})(\mathbf{Z}) \simeq 1 \times \underset{\Gamma(S T_{\bar{p}})}{\mathbf{Z}/6\mathbf{Z}} \times \underset{\Gamma(S T_{\bar{\sigma}^\circ})}{\mathbf{Z}/4\mathbf{Z}} \times \underset{\Gamma(\mu_2 \subset SG)}{\mathbf{Z}/2\mathbf{Z}} \leqno{(54)}
$$

qui n'est que deux fois plus gros que
le groupe des pts entiers du centre
de $\mathcal{M}_{\bar{p},\bar{\sigma}}$.

Rmq. La composante neutre du centre de
$\mathcal{M}_{\bar{p},\bar{\sigma}^\circ}$ est bien sûr $S T_{\bar{p}} \times S T_\sigma$, comme
le suggérait la discussion de la p. 501.

\newpage

%% ===== Page 413 =====

\chapter*{\S \space XI. --- Récapitulation sur le groupe $\mathcal{M}_{p, \sigma}$.}\thispagestyle{empty}
\addcontentsline{toc}{section}{XI. Récapitulation sur le groupe $\mathcal{M}_{p, \sigma}$.}
\label{sec:XI}
\section*{}

\marginpar{Préliminaires} $1^\circ$ Étant donnés, sur un schéma $S$, une forme $G$ de $\operatorname{Gl}(2)$, d'où des formes $\mathfrak{S}, E$ de $\operatorname{Sl}(2), PG$, et deux sections $\bar{\rho}, \bar{\sigma}$ de $PG$, et un sous-groupe à $1$-paramètre additif $\bar{\mathcal{E}}_0 : \mathbb{G}_a \hookrightarrow PG$ (ce qui revient à une section d'éléments unipotents réguliers de $\mathcal{E}_0$ de $G$, ou une section nilpotente régulière $\mathcal{N}_0$ de la forme $M$ de $M_2(\mathcal{O}_S)$ associée à $G$) -- satisfaisant les conditions
$$
\bar{\rho}\bar{\sigma} = \bar{\mathcal{E}}_0(1), \quad \bar{\mathcal{E}}_0 \text{ et } \bar{\mathcal{E}}_1 = \bar{\rho}(\bar{\mathcal{E}}_0) \ (= \bar{\sigma}(\bar{\mathcal{E}}_0)) \text{ sont } \leqno{(1)}
$$
indépendants sur chaque fibre
$$
\Tr \sigma = 0 \quad \text{si } \sigma \text{ est un relèvement local de } \bar{\sigma} \leqno{(2)}
$$

On définit des relèvements canoniques $\rho, \sigma$ de $\bar{\rho}, \bar{\sigma}$ à $G$, satisfaits par les conditions
$$
\begin{cases} \mathcal{E}_0 = -\sigma^{-1} \rho & \text{i.e. } \mathcal{E}_1 = -\rho \sigma^{-1} \\ \Tr(\rho+\sigma)=1 & \text{i.e. } \Tr \rho=1 \end{cases} \leqno{(3)}
$$
et les constructions qui suivent ont un sens.
On construit un épinglage
$$
\operatorname{Gl}(2)_S \isommap G \leqno{(5)}
$$
par la condition que $\mathcal{E}_1$ prenne la forme
$$
\mathcal{E}_1 = \begin{pmatrix} 1 & 0 \\ 1 & 1 \end{pmatrix} \leqno{(6)}
$$
d'où
$$
\begin{cases} 
\rho = \begin{pmatrix} 0 & 1 \\ -1 & 1 \end{pmatrix}, \quad \sigma = \begin{pmatrix} 0 & -1 \\ -\nu & 0 \end{pmatrix} \quad (\nu \in \Gamma(S, \mathbb{G}_m)) \\
\det \rho = \det \sigma = -\nu 
\end{cases} \leqno{(7)}
$$
$$
\mathcal{E}_0 = \begin{pmatrix} 1 & 1/\nu \\ 0 & 1 \end{pmatrix} \leqno{(8)}
$$

\newpage

%% ===== Page 414 =====

on considère les sous-groupes

(9)
$$ \begin{cases} L_0 = \{ \varepsilon_0(\lambda) \} = \left\{ \begin{pmatrix} 1 & \lambda \\ 0 & 1 \end{pmatrix} \right\} \qquad (\lambda \in \Gamma \mathbb{G}_a) \\ L_1 = \sigma(L_0) = \rho(L_0) = \{ \varepsilon_1(\lambda) \} = \left\{ \begin{pmatrix} 1 & 0 \\ \lambda & 1 \end{pmatrix} \right\} \\ B'_0 = \left\{ \begin{pmatrix} \mu & \lambda \\ 0 & 1 \end{pmatrix} \mid \mu \in \Gamma \mathbb{G}_m, \lambda \in \Gamma \mathbb{G}_a \right\} \\ T_0 = \left\{ \begin{pmatrix} \mu & 0 \\ 0 & 1 \end{pmatrix} \mid \mu \in \Gamma \mathbb{G}_m \right\} \end{cases} $$

d'où

(10)
$$ B'_0 = T_0 \cdot L_0 \qquad (\text{produit } 1/2\text{-direct}) $$

et on a des iso.

(11)
$$ \begin{cases} \xi_0 : \mathbb{G}_a \simeq L_0 \qquad \lambda \mapsto \varepsilon_0(\lambda) = \begin{pmatrix} 1 & \lambda \\ 0 & 1 \end{pmatrix} \\ \xi_1 : \mathbb{G}_a \simeq L_1 \qquad \lambda \mapsto \varepsilon_1(\lambda) = \begin{pmatrix} 1 & 0 \\ \lambda & 1 \end{pmatrix} \end{cases} $$
$$ \text{avec } \varepsilon_1(\lambda) = \rho(\varepsilon_0(\lambda)) = \sigma^-(\varepsilon_0(\lambda)) $$

On [unclear: pose] considère

(12)
$$ \begin{cases} T_{\rho} = \text{schém. des unités de } \mathcal{A}_{\rho} = \underline{\mathcal{O}}_S + \underline{\mathcal{O}}_S \rho \simeq \underline{\mathcal{O}}_S[X]/(X^2-X-\nu) \\ T_{\sigma} = \text{schém. des unités de } \mathcal{A}_{\sigma} = \underline{\mathcal{O}}_S + \underline{\mathcal{O}}_S \sigma \simeq \underline{\mathcal{O}}_S[Y]/(Y^2-\nu) \end{cases} $$

(13)
$$ N'_{\sigma} = \{ g \in G \mid g(\sigma) = \varepsilon \sigma, \varepsilon^2 = 1 \} $$
$$ \cap $$
$$ G $$

d'où

(14)
$$ 1 \to T_{\sigma} \to N'_{\sigma} \xrightarrow{\varepsilon_{\sigma}} \mu_2 \to 1 $$

On pose

(15)
$$ \begin{cases} \mathcal{G}_{\rho,\sigma} = B'_0 \times T_{\rho} \times T_{\sigma} \\ \mathcal{G}^*_{\rho,\sigma} = B'_0 \times T_{\rho} \times N'_{\sigma} \end{cases} $$

\newpage

%% ===== Page 415 =====

2° On introduit
$$ M_{\rho,\sigma}[0] \subset G_{\rho,\sigma}^* = B_0' \times T_\rho \times N_\sigma' \leqno{(16)} $$
$$ \| \defeq $$
$$ \{ (u, a, b) \mid \det u = \det b = \varepsilon' \det a, \varepsilon'^2 = 1 \} $$

$$ M_{\rho,\sigma}(0) \subset M_{\rho,\sigma}[0] \quad \text{défini par } u=1 \text{ i.e.} \leqno{(17)} $$
$$ M_{\rho,\sigma}(0) = M_{\rho,\sigma}[0] \cap (T_0 \times T_\rho \times N_\sigma') \quad \text{dans } G_{\rho,\sigma}^* $$

$$ L_0 \hookrightarrow M_{\rho,\sigma}[0] \leqno{(18)} $$
$$ \lambda \mapsto (\lambda, 1, 1) $$

donc on a
$$ M_{\rho,\sigma}[0] \simeq M_{\rho,\sigma}(0) \cdot L_0 \leqno{(19)} $$
produit $\frac{1}{2}$ direct

On a un homomorphisme
$$ M_{\rho,\sigma}[0] \longrightarrow G_m \times \mu_2 \times \mu_2 \leqno{(20)} $$
noté $g \mapsto (\chi(g), \varepsilon_1(g), \varepsilon_2(g))$

pour $g = (u, a, b)$ on a
$$
\begin{cases}
\chi(g) = \det u = \det b \\
\varepsilon_1(g) = \varepsilon(b) \quad (cf \ (13), (14)) \\
\varepsilon_2(g) = \det(a) / \det b \quad \text{i.e.} \quad \det a = \varepsilon_2 \mu \text{ (où } \mu = \chi(g))
\end{cases}
$$

trivial sur $L_0$ (donc se factorise par le quotient). On pose ainsi
$$ \mathcal{T}_{\rho,\sigma}\footnote{[unclear: Cette notation s'impose en fait, pour des raisons évidentes. L'inconvénient qu'il y aura pour la suite en traitant le précédent comme combinaison de $G_{\rho,\sigma}$ et $\Gamma_{0,3}$]} \simeq M_{\rho,\sigma}[0] / L_0 \simeq \dots \leqno{(21)} $$

$$
\begin{tikzcd}
M_{\rho,\sigma}(0) \arrow[r, hook] \arrow[rr, bend right=20, "\text{iso. déduit}"'] & M_{\rho,\sigma}[0] \arrow[r] & \mathcal{T}_{\rho,\sigma}
\end{tikzcd} \leqno{(22)}
$$

$$ M_{\rho,\sigma}(0) \simeq \mathcal{T}_{\rho,\sigma} \hookrightarrow T_0 \times T_\rho \times N_\sigma' \leqno{(23)} $$

\newpage

%% ===== Page 416 =====

La composition avec la projection can. dans $T_{\rho} \times N'_{\sigma}$ donne un mono

[MARGIN: car $\det$ sur $T_0 \xrightarrow{\sim} \mathbb{G}_m$]
(24) $$ \mathcal{T}_{\rho, \sigma} \hookrightarrow T_{\rho} \times N'_{\sigma} $$
$$ \simeq \left\{ (\underline{a}, \underline{b}) \mid (\det \underline{a} / \det \underline{b})^2 = 1 \right\} $$
$$ \text{i.e. } \det \underline{a} = \varepsilon' \det \underline{b} \text{ avec } {\varepsilon'}^2=1 $$

et on aura dans

(25) $$ \begin{cases} \chi(\underline{a}, \underline{b}) = \det \underline{b} \\ \varepsilon_2(\underline{a}, \underline{b}) = \varepsilon_{\sigma}(\underline{b}) \\ \varepsilon_3(\underline{a}, \underline{b}) = \det(\underline{a}) / \det(\underline{b}) = \det \underline{b} / \det \underline{a} \end{cases} $$

On définit

(26) $$ \underset{\simeq T_{\rho} \times_{\mathbb{G}_m} T_{\sigma}}{\mathcal{T}_{\rho, \sigma}^{\circ}} = \text{Ker } \Big( \mathcal{T}_{\rho, \sigma} \xrightarrow{(\varepsilon_2, \varepsilon_3)} \mu_2 \times \mu_2 \Big) $$

de sorte on a

(27) $$ 1 \to \mathcal{T}_{\rho, \sigma}^{\circ} \to \mathcal{T}_{\rho, \sigma} \to \mu_2 \times \mu_2 \to 1 $$

et une suite exacte canonique

(28) $$ 1 \to S T_{\rho} \times S T_{\sigma} \to \mathcal{T}_{\rho, \sigma}^{\circ} \xrightarrow{\chi} \mathbb{G}_m \to 1 $$

où on définit

(29) $$ \begin{cases} S T_{\rho} = \{ g \in T_{\rho} \mid \det g = 1 \} \\ S T_{\sigma} = \{ g \in T_{\sigma} \mid \det g = 1 \} \end{cases} $$

[MARGIN: L'opération de $\mathcal{T}_{\rho, \sigma} \simeq \mathcal{M}_{\rho, \sigma}(0)$ sur $L_0$ est donnée par $g(\lambda, \mu) = \varepsilon_3(g)(\lambda, \mu)$, où $g \in \mathcal{T}_{\rho, \sigma}$]

On fait opérer $\mathcal{M}_{\rho, \sigma}(0) (\simeq \mathcal{T}_{\rho, \sigma})$ sur $SG$, par projection

(30)
\begin{tikzcd}
\mathcal{M}_{\rho, \sigma}(0) \ar[r, "\text{projection}"] \ar[d, hook] & T_0 \ar[r, "\text{can.}"', "\text{(mono)}"] \ar[d, hook] & P G \ar[r, "\sim"] \ar[d, equal] & \operatorname{Aut}_{gr}(SG) \\
T_0 \times T_\rho \times N'_{\sigma} & G & G / \mathbb{G}_m \simeq \operatorname{Aut}(G) &
\end{tikzcd}

\newpage

%% ===== Page 417 =====

et on définit
$$ \mathcal{M}_{\rho, \sigma} \simeq \mathcal{M}_{\rho, \sigma}(0) . SG \leqno (32) $$
produit $\frac{1}{2}$ direct.

On a donc un hom. can.
$$ SG \hookrightarrow \mathcal{M}_{\rho, \sigma} \leqno (33) $$
dont l'image est notée $\mathfrak{S}_{\rho, \sigma}^+$, d'où
$$ 1 \to \mathfrak{S}_{\rho, \sigma}^+ \to \mathcal{M}_{\rho, \sigma} \to \mathfrak{T}_{\rho, \sigma} \to 1 \leqno (34) $$
$$ \mathcal{M}_{\rho, \sigma} \simeq \mathcal{M}_{\rho, \sigma}(0) \underset{\mathfrak{T}_{\rho, \sigma}^\circ}{.} \mathfrak{S}_{\rho, \sigma}^+ \leqno (35) $$
produit $\frac{1}{2}$ direct.

\textsc{NB.} On utilise la notation $\mathfrak{T}_{\rho, \sigma}$ pour le
groupe quotient
$$ \mathfrak{T}_{\rho, \sigma} \simeq \mathcal{M}_{\rho, \sigma} / \mathfrak{S}_{\rho, \sigma}^+ \quad , \leqno (36) $$
la notation $\mathcal{M}_{\rho, \sigma}[0]$ pour le sous-groupe
de $\mathcal{M}_{\rho, \sigma}$.

On définit un plongement canonique
$$ \mathcal{M}_{\rho, \sigma}[0] \hookrightarrow \mathcal{M}_{\rho, \sigma} \leqno (37) $$
qui prolonge l'indentification des sous-groupes
$\mathcal{M}_{\rho, \sigma}(0)$ communs de $\mathcal{M}_{\rho, \sigma}[0]$ et de $\mathcal{M}_{\rho, \sigma}$,
et qui sur $L_0$ est donné par
$$ i(\lambda) = j(\lambda) \quad (cf (33)) \leqno (38) $$
(ce qui est bien compatible avec les opérations
de $\mathcal{M}_{\rho, \sigma}(0)$ sur $L_0$ d'une part, sur $SG \simeq \mathfrak{S}_{\rho, \sigma}^+$
de l'autre).

\newpage

%% ===== Page 418 =====

Une autre façon de le dire, c'est de constater
que $\mathcal{M}_{\rho,\sigma(0)}$ opérant sur $SG \simeq \mathfrak{S}_{\rho,\sigma}^+$
normalise le sous-groupe $j(L_0)$ - que nous
identifions désormais à $L_0$), de sorte qu'on
aura un sous-groupe
$$
\mathcal{M}_{\rho,\sigma(0)} . L_0 \subset \mathcal{M}_{\rho,\sigma}
$$
qui sera canoniquement isom. à $\mathcal{M}_{\rho,\sigma}[0]$
défini plus haut, grâce à [unclear: l'inclusion (37) de]
celui-ci comme produit $1/2$-direct.

Remarque. On a
$$
\mathcal{M}_{\rho,\sigma}[0] = \operatorname{Norm}_{\mathcal{M}_{\rho,\sigma}}(L_0)
$$
Par construction, on a $\subset$, l'inclusion
en sens inverse reviendrait à dire que $L_0$
est son propre normalisateur dans $SG \simeq \mathfrak{S}_{\rho,\sigma}^+$
-- ce qui bien sûr n'est pas vrai du tout
-- bien que le fait analogue soit vrai
dans le groupe profini $\hat{\mathfrak{S}}_{0,3}^+$, pour le s-.
groupe engendré par $s_0$, ou dans $\hat{\mathfrak{T}}_{0,3}$ pour
celui engendré par $l_0$ (une [unclear: simple] vérification
d'ailleurs que je n'ai pas faite).

\newpage

%% ===== Page 419 =====

On désigne par $\mathcal{M}_{\rho,\sigma}^\circ$, $\mathcal{M}_{\rho,\sigma}[\circ]^\circ$, $\mathcal{M}_{\rho,\sigma}(\circ)^\circ$\marginpar{$\mathcal{M}_{\rho,\sigma}^\circ, \mathcal{M}_{\rho,\sigma}^{\circ +}, \mathcal{M}_{\rho,\sigma}^{+ \circ}, \mathfrak{S}_{\rho,\sigma}^{+ \circ}, \mathfrak{T}_{\rho,\sigma}^\circ \dots$ quelques sous-groupes calqués sur celui des pts rationnels...} les composantes neutres, on aura p.ex.
$$
1 \to \underbrace{\mathcal{M}_{\rho,\sigma}^\circ}_{\substack{\simeq \mathcal{M}_{\rho,\sigma}(\circ) \cdot \mathfrak{S}_{\rho,\sigma}^+\\ \text{produit 1/2 direct}}} \to \mathcal{M}_{\rho,\sigma} \to \mu_3 \times \mu_2 \to 1 \leqno (39)
$$

$$
\mathcal{M}_{\rho,\sigma}(\circ)^\circ \simeq \mathfrak{T}_{\rho,\sigma}^\circ \simeq \mathfrak{T}_\rho \times_{\mathfrak{S}_m} \mathfrak{T}_\sigma \leqno (40)
$$

$$
\mathcal{M}_{\rho,\sigma}^\circ \text{ est l'image inverse de } \mathfrak{T}_{\rho,\sigma}^\circ \text{ par } \mathcal{M}_{\rho,\sigma} \xrightarrow{\text{can}} \mathfrak{T}_{\rho,\sigma} \leqno (41)
$$

On a un hom canonique surjectif
$$
\mathcal{M}_{\rho,\sigma} \to \mathfrak{T}_{\rho,\sigma} \xrightarrow{\chi} \mathfrak{S}_m \leqno (42)
$$

également noté $\chi$, son noyau est noté $\mathcal{M}_{\rho,\sigma}^+$, l'on définit de même $\mathcal{M}_{\rho,\sigma}^{\circ +}$, etc.
$$
\begin{cases}
1 \to \mathcal{M}_{\rho,\sigma}^+ \to \mathcal{M}_{\rho,\sigma} \xrightarrow{\chi} \mathfrak{S}_m \to 1 \\
1 \to \mathcal{M}_{\rho,\sigma}^{\circ +} \to \mathcal{M}_{\rho,\sigma}^\circ \xrightarrow{\chi} \mathfrak{S}_m \to 1
\end{cases} \leqno (43)
$$

$$
\begin{cases}
\mathcal{M}_{\rho,\sigma}^+ \simeq \mathcal{M}_{\rho,\sigma}^+(\circ) \cdot \mathfrak{S}_{\rho,\sigma}^+ \\
\mathcal{M}_{\rho,\sigma}^{\circ +} \simeq \mathcal{M}_{\rho,\sigma}^{\circ +}(\circ) \cdot \mathfrak{S}_{\rho,\sigma}^+
\end{cases} \leqno (44)
$$
\begin{center}
$\mathfrak{T}_{\rho,\sigma}^{\circ +} \simeq ST_\rho \times ST_\sigma \quad \text{cf (28)}$
\end{center}

Donc on a
$$
\mathcal{M}_{\rho,\sigma}^{\circ +}(0) \defeq \Ker \left( \mathcal{M}_{\rho,\sigma}(0) \xrightarrow{\chi,\varepsilon_2,\varepsilon_3} \mathfrak{S}_m \times \mu_3 \times \mu_2 \right) \leqno (45)
$$
$$
\simeq ST_\rho \times ST_\sigma
$$

et ce sous-groupe est central dans $\mathcal{M}_{\rho,\sigma}^{\circ}$ ([unclear: près de] $\mathcal{M}_{\rho,\sigma}$), de façon précise

\newpage

%% ===== Page 420 =====

\footnote{Toute la fin du n$^\circ$ devrait venir après la [unclear: vue] du n$^\circ$ [unclear: 36] après la "vérification" + attention, [unclear: avec] [unclear: l'idée] [unclear: de] l'égalité y plus bas...}
$$ 
\mathcal{M}_{\rho, \sigma}^{0 +} \defeq \Ker \left( \mathcal{M}_{\rho, \sigma}^0 \xrightarrow{(\chi, \varepsilon_1, \varepsilon_2)} \mathbb{G}_m \times \mu_2 \times \mu_2 \right) \leqno{(46)}
$$
$$
\simeq \underset{\subset \mathcal{M}_{\rho, \sigma}(0)}{S T_{\rho} \times S T_{\sigma}} \times \mathfrak{S}_{\rho, \sigma}^+
$$

$$ 
\operatorname{Centr}(\mathcal{M}_{\rho, \sigma}^0) \supset \underset{\subset \mathcal{M}_{\rho, \sigma}(0)^0}{1 \times S T_{\rho} \times S T_{\sigma}} \times \underset{\subset \mathfrak{S}_{\rho, \sigma}^+ = S G}{\mu_2} \leqno{(47)} 
$$

D'où un sous-groupe

$$ 
\mu_2 \times \mu_2 \times \mu_2 \subset S T_{\rho} \times S T_{\sigma} \times \mu_2 \leqno{(48)} 
$$

qui est en fait contenu dans le centre de $\mathcal{M}_{\rho, \sigma}^0$ tout entier, car

$$ 
\operatorname{Centr}(\mathcal{M}_{\rho, \sigma}) \supset \underset{\cap \mathcal{M}_{\rho, \sigma}(0)}{1 \times S' T_{\rho}} \times \underset{\subset \mathfrak{S}_{\rho, \sigma}^+ = S G}{\mu_2 \times \mu_2} \leqno{(49)} 
$$

avec

$$ 
S' T_{\rho} \defeq \{ g \in S T_{\rho} \mid (\det g)^2 = 1 \} \leqno{(50)} 
$$

de sorte qu'on a

$$ 
1 \to S T_{\rho} \to S' T_{\rho} \xrightarrow{\varepsilon_3} \mu_2 \to 1 \leqno{(51)} 
$$

On a donc

$$ 
\operatorname{Centr}(\mathcal{M}_{\rho, \sigma}) \cap \mathcal{M}_{\rho, \sigma}^0 \supset \underset{T_0}{1} \times \underset{T_\rho}{S T_{\rho}} \times \underset{T_\sigma}{\mu_2} \times \underset{\mathfrak{S}_{\rho, \sigma}^+}{\mu_2} \leqno{(52)} 
$$

Le "cas modulaire" est celui où

$$ 
\begin{cases} 
\det \rho = 1 \text{ i.e. } \det \sigma = 1 \text{ i.e. } \nu = -1 \\ 
\text{i.e. } \rho = \begin{pmatrix} 0 & 1 \\ -1 & 1 \end{pmatrix} \text{ i.e. } \sigma = \begin{pmatrix} 0 & 1 \\ -1 & 0 \end{pmatrix} 
\end{cases} \leqno{(53)} 
$$

(cf VI pour d'autres conditions, telles $\sigma^2 = -1$, ou $\rho^2 - \rho + 1 = 0$, ou $\tilde{\sigma}^2 = \tilde{\rho}^3 = 1$ ... ), cela signifie aussi [unclear: qui] [unclear: plus] [unclear: est] [unclear: que] l'image de $T_\rho \isommap S T_\rho$ est dans le centre de $\mathcal{M}_{\rho, \sigma}^0$.

\newpage

%% ===== Page 421 =====

de même que $\sigma$ [unclear: est un élément] qui est dans $T_{\sigma}$ et dans $S T_{\sigma}$ i.e. dans le centre de $M_{\rho, \sigma}$. Donc dans ce cas on trouve un homomorphisme
$$
\begin{cases}
\mathbf{Z}/6\mathbf{Z} \times \mathbf{Z}/4\mathbf{Z} \to \operatorname{Centr}(M_{\rho, \sigma}) \\
(i, j) \mapsto i_{\rho}(\rho^i) i_{\sigma}(\sigma^j)
\end{cases} \leqno{(54)}
$$
(où on définit $i_{\rho}, i_{\sigma}$ par les homomorphismes canoniques
$$
\begin{cases}
i_{\rho} : S T_{\rho} \to M_{\rho, \sigma}^\circ \hookrightarrow M_{\rho, \sigma} \\
i_{\sigma} : S T_{\sigma} \to M_{\rho, \sigma}^\circ \hookrightarrow M_{\rho, \sigma} \\
i_{G} : SG \simeq \hat{\mathfrak{S}}^+_{\rho, \sigma} \to M_{\rho, \sigma} \text{,}
\end{cases} \leqno{(55)}
$$
dont l'image pour $2$ inversible dans $S$, contient l'image de $\mu_2 \times \mu_2 \times 1$ dans (48). Pour $2$ inv. dans $S$, l'hom. (54) est un mono (mais pas un iso. bien sûr). Sur le bon absolu $\mathbf{Z}$, on a par ailleurs, dans le cas modulaire
$$
\begin{cases}
\tau_{\rho, \sigma}(\mathbf{Z}) \simeq T_{\rho}(\mathbf{Z}) \times \operatorname{N}^1_{\sigma}(\mathbf{Z}) \simeq \mathbf{Z}/6\mathbf{Z} \times D_4 \\
M_{\rho, \sigma}(\mathbf{Z}) \simeq \tau_{\rho, \sigma}(\mathbf{Z}) \cdot \operatorname{Sl}(2, \mathbf{Z}) \\
\qquad \begin{matrix}
\text{produit } 1/2 \text{ direct} \\
\tau_{\rho, \sigma} \simeq \mathbf{Z}/6\mathbf{Z} \times D_4 \text{ opérant} \\
\text{via le caractère } D_4 \xrightarrow{\varepsilon_2} \mathbf{Z}/2\mathbf{Z} \simeq \{\pm 1\} \\
\tau_{\sigma} \text{ opérant sur } \operatorname{Sl}(2, \mathbf{Z}) \text{ par} \\
\text{automo induit par la matrice} \\
\tau_{\sigma} = \begin{pmatrix} -1 & 0 \\ 0 & 1 \end{pmatrix}
\end{matrix}
\end{cases} \leqno{(56)}
$$

\newpage

%% ===== Page 422 =====

[MARGIN: Les homs $\Theta, \Theta^\circ$]
[MARGIN: 40]

On définit un hom.
(57) $$M_{\rho,\sigma} \overset{\Theta}{\longrightarrow} G$$
$\Theta | \underline{G}_{\rho,\sigma}^+ \simeq SG$ est l'identité
$\Theta | M_{\rho,\sigma}(0) \simeq T_0 \times \Gamma_\rho \times N_\sigma'$ donné par $\Theta(\underline{u}, \underline{a}, \underline{b}) = \underline{u}$

d'où
(58) $\Theta | M_{\rho,\sigma}[0] \simeq B_0' \times \Gamma_\rho \times N_\sigma'$ donné par $\Theta(\underline{u}, \underline{a}, \underline{b}) = \underline{u}$

On a
(59)
\begin{tikzcd}
M_{\rho,\sigma}(0) \ar[r, equal] & T_0 \times \Gamma_\rho \times N_\sigma' \\
\operatorname{Ker} \Theta \ar[u, hook] \ar[r, "\subset"] & 1 \times \Gamma_\rho \times N_\sigma' \ar[u, hook, "cart."'] \\
& \{ (1, \underline{a}, \underline{b}) \mid \det \underline{b} = 1, (\det \underline{a})^2 = 1 \} \ar[u, "\simeq"'] \\
& \Gamma T_\rho \times \Gamma N_\sigma' \ar[u, "\subset"'] \\
& S T_\rho \times S N_\sigma' \ar[u, "\simeq"']
\end{tikzcd}
où $SN_\sigma' \overset{d\acute{e}f}{=} \operatorname{Ker}(N_\sigma' \to G_m)$

Soit
(60) $\Theta^\circ = \Theta | M_{\rho,\sigma}^\circ$ , on trouve

(61)
\begin{tikzcd}
1 \ar[r] & \operatorname{Ker} \Theta^\circ \ar[r] \ar[d, "\simeq"'] & \operatorname{Ker} \Theta \ar[r, "(\varepsilon_1{,}\varepsilon_2)"] & \mu_2 \times \mu_2 \ar[r] & 1 \\
& S T_\rho \times S T_\sigma \ar[d, "\simeq"'] & & & \\
& \operatorname{Rad}(M_{\rho,\sigma}^\circ) \ar[d, "\simeq"'] & & & \\
& \operatorname{Centr}(M_{\rho,\sigma}^\circ) & & &
\end{tikzcd}

Posons
(62) $SR_{\rho,\sigma}^\circ = S T_\rho \times S T_\sigma = \operatorname{Ker} \Theta^\circ = (\operatorname{Ker} \Theta) \cap M_{\rho,\sigma}^\circ$
qui est un sous-groupe central de $M_{\rho,\sigma}^\circ$ (il est

\newpage

%% ===== Page 423 =====

distingué [crossed out: invariant] mais pas central dans $M_{\rho,\sigma}$ tout entier)

le groupe
(63) $$M_{\rho,\sigma} / S R_{\rho,\sigma}^0 = M_{\rho,\sigma}^\natural$$
est donc une extension de $\mu_2 \times \mu_2$ par
$M_{\rho,\sigma}^0 / S R_{\rho,\sigma}^0 \underset{\text{déf}}{=} M_{\rho,\sigma}^{0\natural}$, isomorphe à $G$ par
$\Theta^0$ (puisque $S R_{\rho,\sigma}^0 = \text{Ker } \Theta^0$), d'où
(64) $$1 \to G \to M_{\rho,\sigma}^\natural \to \mu_2 \times \mu_2 \to 1,$$
d'ailleurs l'hom.
$$\Theta : M_{\rho,\sigma} \to G$$
étant trivial sur $S R_{\rho,\sigma}^0$ passe au quotient en
(65) $$\Theta^\natural : M_{\rho,\sigma}^\natural \to G,$$
dont la restriction à $G$ (via (64)) est l'
identité, donc on trouve, via (64)
(66) $$M_{\rho,\sigma}^\natural \simeq G \times \mu_2 \times \mu_2.$$
On trouve donc
(67)
\begin{tikzcd}
1 \ar[r] & S R_{\rho,\sigma}^0 \ar[r] \ar[d, "\simeq"'] & M_{\rho,\sigma} \ar[r, "((\Theta), \varepsilon_3, \varepsilon_2)"] & M_{\rho,\sigma}^\natural \ar[r] \ar[d, "\simeq"] & 1 \\
& S T_\rho \times S T_\sigma & & G \times \mu_2 \times \mu_2 &
\end{tikzcd}

Pour voir comment $M_{\rho,\sigma}^\natural \simeq G \times \mu_2 \times \mu_2$
opère sur $S T_\rho \times S T_\sigma$, on [unclear: note que le s-groupe]
$M_{\rho,\sigma}^0$ [crossed out: est l'] image inverse de $G \simeq M_{\rho,\sigma}^{0\natural}$
dans l'extension (67), or $S R_{\rho,\sigma}^0$ est central
dans $M_{\rho,\sigma}^0$, donc $G$ y opère trivialement.

\newpage

%% ===== Page 424 =====

D'autre part on a
$$
\begin{tikzcd}
S\mathcal{R}^0_{\rho, \sigma} \ar[d, "\simeq"'] \ar[r, hook] & \mathcal{M}_{\rho, \sigma}(0) \ar[d, "\cap"'] \ar[r, twoheadrightarrow, "\text{épi}", "(\varepsilon_3, \varepsilon_2)"'] & \mu_2 \times \mu_2 \\
S\mathfrak{T}_{\rho} \times S\mathfrak{T}_{\sigma} & \mathfrak{T}_0 \times \mathfrak{T}_{\rho} \times \mathcal{N}_{\sigma} &
\end{tikzcd}
$$
et on trouve que $\mu_2 \times \mu_2$ opère sur $S\mathcal{R}_{\rho,\sigma}$ via ses deux projections (correspondant à $\varepsilon_2(u,v,k) = \varepsilon_v(k)$), un $\varepsilon$ de celui-ci ($\varepsilon = \pm 1$) opère sur $S\mathcal{R}_{\rho,\sigma} \simeq S\mathfrak{T}_{\rho} \times S\mathfrak{T}_{\sigma}$ en laissant inchangé le premier facteur, et en opérant sur $\mathfrak{T}_{\sigma}$ donc sur $S\mathfrak{T}_{\sigma}$ via
$$
\sigma \mapsto \varepsilon \sigma \leqno{(68)}
$$
(compatible avec l'équation de Cayley $\sigma^2 - v = 0$ de $\sigma$).

Soit\footnote{Il y aurait intérêt de suite à [unclear: l'étude] de l'extension $\mathcal{R}_{\rho,\sigma}$ (voir (69) plus bas) et de la relation de $\mathcal{M}_{\rho,\sigma}$, - les deux s'obtenant de [unclear: façon] voir... p. (76).}
$$
\mathcal{R}^0_{\rho,\sigma} \quad \text{image inverse de } G_{\text{un}} \subset G \text{ dans } \mathcal{M}_{\rho,\sigma} \text{ par } \Theta: \mathcal{M}_{\rho,\sigma} \to G \leqno{(69)}
$$
d'où une structure d'extension
$$
1 \to S\mathcal{R}^0_{\rho,\sigma} \to \mathcal{R}^0_{\rho,\sigma} \to G_{\text{un}} \to 1 \leqno{(70)}
$$
On aura
$$
\mathcal{R}^0_{\rho,\sigma} = \left\{ (\underset{\in \mathfrak{T}_0}{\underline{u}}, \underset{\in \mathfrak{T}_\rho}{\underline{\sigma}}, \underset{\in \mathcal{N}_\sigma}{\underline{k}}), f \ \left| \ \det u = \det \sigma = \det k \ , \ u, f \in G_{\text{un}} \atop{\text{i.e. } f = \lambda \cdot u^\sim \atop{\text{où } \lambda \in G_{\text{un}}, \ \lambda^2 = \det u}} \right. \right\}
$$

\newpage

%% ===== Page 425 =====

514

$$ (71) \quad R_{\rho,\sigma}^0 \simeq \left\{ (\mu, \underset{\in \Gamma_{T_\rho}}{\underline{a}}, \underset{\in \Gamma_{T_\sigma}}{\underline{b}}) \mid \mu \in \mathbb{G}_m, \text{det } a = \text{det } b = \mu^{-2} \right\} $$
[MARGIN: NB. le $\text{det}$ est $1$ car $\text{det } j_{\mathbb{G}}(\mu) = \mu^{-2}$]

où
$$ (72) \quad \begin{cases} j_{+}(\mu) = \begin{pmatrix} \mu & 0 \\ 0 & 1 \end{pmatrix} \in S\mathbb{G} \quad (\mu \in \mathbb{G}_m) \\ j_{T_0} : \mathbb{G}_m \to T_0 \quad \text{défini plus haut} \\ j_{\mathbb{G}} : \mathbb{G}_m \to \mathbb{G} \quad \text{défini plus haut} \end{cases} $$

On trouve ainsi que la composition
$$ R_{\rho,\sigma}^0 \hookrightarrow M_{\rho,\sigma}^0 \twoheadrightarrow \tau_{\rho,\sigma}^0 \simeq T_\rho \times_{\mathbb{G}_m} T_\sigma $$
$$ M_{\rho,\sigma}^0 / S P_{\rho,\sigma}^0 $$

$$ (73) \quad \begin{cases} \text{fait de } R_{\rho,\sigma}^0 \text{ un [unclear: revêtement] de groupes} \\ \text{de } \tau_{\rho,\sigma}^0 \end{cases} $$

et plus précisément qu'on a un carré cartésien

(74)
\begin{tikzcd}
R_{\rho,\sigma}^0 \ar[r, "(73)"] \ar[d, "\text{l'hom. can.}"', "\text{de (70)}"] & T_\rho \times_{\mathbb{G}_m} T_\sigma \simeq \tau_{\rho,\sigma}^0 \ar[d, "\text{can}"] \\
\mathbb{G}_m \ar[r, "\mu \mapsto \mu^{-2}"'] & \mathbb{G}_m
\end{tikzcd}

(75) Proposition : $R_{\rho,\sigma}^0 = \text{Centre} (M_{\rho,\sigma}^0)$

Évidemment le centre de $M_{\rho,\sigma}^0$ est contenu dans $R_{\rho,\sigma}^0$, il s'agit inversement de [unclear: justifier] des assertions...

\newpage

%% ===== Page 426 =====

de $M_{\rho,\sigma}^\circ = \mathcal{M}_{\rho,\sigma}^\circ/\mathcal{S}_{\rho,\sigma} \simeq G$, il reste à prouver que $R_{\rho,\sigma}^\circ$ est central dans $\mathcal{M}_{\rho,\sigma}^\circ$. Mais [unclear: puisque] $\mathcal{M}_{\rho,\sigma}^\circ$ est engendré par $R_{\rho,\sigma}^\circ$ et par $\mathcal{S}_{\rho,\sigma}^+ \simeq \SL_2$, puisque $R_{\rho,\sigma}^\circ \to C_{\rho,\sigma}^\circ \simeq \mathcal{M}_{\rho,\sigma}^\circ / \mathcal{S}_{\rho,\sigma}^+$ est épi (73)(74), et comme $R_{\rho,\sigma}^\circ$ est commutatif (74), il reste à montrer qu'il commute à $\mathcal{S}_{\rho,\sigma}^+$ i.e.\ opère trivialement sur $\mathcal{S}_{\rho,\sigma}^+$. Or une section de $R_{\rho,\sigma}^\circ$ est de la forme :
$$ (j_{\Gamma_0}(\mu^2), \underline{a}, \underline{b}) \cdot (\mathcal{J}_{j_{\Gamma_0}(\mu^{-2})}) $$
le premier facteur opérant comme $\operatorname{int} j_G(\mu^2)$, le deuxième comme $\operatorname{int} j_G(\mu^{-2})$, donc leur produit opère trivialement cqfd.

Notons maintenant que
$$ R_{\rho,\sigma}^\circ \cap \underset{\substack{\parallel \\ S G}}{\mathcal{S}_{\rho,\sigma}^+} = \operatorname{Ker}(R_{\rho,\sigma}^\circ \to C_{\rho,\sigma}^\circ) = \mu_2 \leqno{(75)} $$
\marginpar{cf (74)}

(on peut dire aussi plus simplement, que

\newpage

%% ===== Page 427 =====

$R_{p,\sigma}^\circ \cap \mathfrak{S}_{p,\sigma}^+$, étant contenu dans le centre de $\mathfrak{S}_{p,\sigma}^+ \simeq \mathfrak{S} G$ qui est $\mu_2$, est $\subset \mu_2$ --- mais cela ne montre pas l'égalité.)

\medskip

\noindent
Remarque. Cela donne une description plus simple de $\mathcal{M}_{p,\sigma}^\circ$ que comme produit $\frac{1}{2}$ direct, de la façon suivante. On définit $R_{p,\sigma}^\circ$ par le diagramme (74), d'où une suite exacte
$$
1 \to \mu_2 \xrightarrow{j'} R_{p,\sigma}^\circ \to T_p \times_{G_{ini}} T_\sigma \to 1 \leqno{(77)}
$$
et on définit $\mathcal{M}_{p,\sigma}^\circ$ comme somme amalgamée centrale
$$
\begin{tikzcd}[row sep=tiny]
1 \arrow[r] & \mu_2 \arrow[r] \arrow[d, maps to] & R_{p,\sigma}^\circ \times \underset{\displaystyle \shortparallel \vspace{1ex}}{\mathfrak{S}_{p,\sigma}^+} \arrow[r] & \mathcal{M}_{p,\sigma}^\circ \arrow[r] & 1 \\[-1ex]
& \varepsilon \arrow[r, maps to] & (j'(\varepsilon), j''(\varepsilon^{-1})) & & \\[-3ex]
& & \mathfrak{S} G & & 
\end{tikzcd} \leqno{(78)}
$$
(où $j'$ est donné dans (77) et $j'' \colon \mu_2 \to \mathfrak{S} G = \mathfrak{S}_{p,\sigma}^+$ l'homomorphisme can.).

Notons qu'on a de même une description bien connue de $G$ en termes de $\mathfrak{S} G$ :
$$
\begin{tikzcd}[row sep=tiny]
1 \arrow[r] & \mu_2 \arrow[r] \arrow[d, maps to] & G_{ini} \times \mathfrak{S} G \arrow[r] & G \arrow[r] & 1 \\[-1ex]
& \varepsilon \arrow[r, maps to] & (\varepsilon, j''(\varepsilon^{-1})) & &
\end{tikzcd} \leqno{(79)}
$$
et on a un hom. de (78) dans (79), grâce à l'hom. can.
$$
R_{p,\sigma}^\circ \to G_{ini} \leqno{(80)}
$$
de (74).

\newpage

%% ===== Page 428 =====

ce qui par passage aux quotients des
$$ R^\circ_{\rho,\sigma} \times S \mathfrak{S} \xrightarrow{\widetilde{\Theta}^\circ} G_{can} \times S \mathfrak{S} $$
permet de reconstituer
$$ \Theta^\circ \colon \mathcal{M}^\circ_{\rho,\sigma} \longrightarrow G , $$
tout comme le groupe $\mathcal{M}^\circ_{\rho,\sigma}$ lui-même,
en termes des groupes ($\simeq$ toriques) "en
l'infini" $T_\rho, T_\sigma$ et des épimorphismes
$$ T_\rho \xrightarrow{\delta_\rho} G_{can} \quad , \quad T_\sigma \xrightarrow{\delta_\sigma} G_{can} \ , $$
et de la forme $b$ de $\operatorname{Sl}(2)_{\mathbf{Z}}$ bien sûr.
Ces [unclear: degrés d'] arbitraire = des plongements
$T_\rho \hookrightarrow G$, $T_\sigma \hookrightarrow G$, et un choix [unclear: continu]
des sections $\rho, \sigma$ de $G \dots$). On
trouve $\mathfrak{S}^+_{\rho,\sigma}$ comme image de $S\mathfrak{S}$ dans $\mathcal{M}^\circ_{\rho,\sigma}$,
Le quotient
$$ \mathcal{T}^\circ_{\rho,\sigma} \defeq \mathcal{M}^\circ_{\rho,\sigma} / \mathfrak{S}^+_{\rho,\sigma} \simeq R^\circ_{\rho,\sigma} / \mu_2 \simeq T_\rho \times_{G_{can}} T_\sigma $$

Mais pour notre propos, le
scindage (dit "distingué") de l'extension ainsi construite
de $\mathcal{T}^\circ_{\rho,\sigma}$ par $\mathfrak{S}^+_{\rho,\sigma} \simeq S\mathfrak{S}$, correspondant
à un sous-groupe section $\mathcal{M}^\circ_{\rho,\sigma}(b)$, jouera un
rôle essentiel - en comparant ces différentes sections.

\newpage

%% ===== Page 429 =====

\marginpar{516}
$$ \mathcal{T}_{\rho,\sigma}^\circ \simeq \mathcal{T}_\rho \times_{G_-} \mathcal{T}_\sigma \longrightarrow \mathcal{M}_{\rho,\sigma}^\circ $$
avec
$$ \mathcal{M}_{\rho,\sigma}^\circ \longrightarrow G $$
on trouve un homomorphisme qui se factorise ainsi :
$$ 
\begin{tikzcd}
\mathcal{T}_{\rho,\sigma}^\circ \simeq \mathcal{T}_\rho \times_{G_-} \mathcal{T}_\sigma \arrow[r, "\text{épi}", "\text{can}"'] & G_- \arrow[r, hook, "\text{mono}"] & G
\end{tikzcd}
$$
dont l'image est le tore $T_-$...

Remarque. On peut [unclear: de façon analogue] construire $\mathcal{M}_{\rho,\sigma}$ de façon analogue, en introduisant l'image inverse --- notons le $\mathcal{R}_{\rho,\sigma}$\footnote{NB $\mathcal{R}_{\rho,\sigma}$ est [unclear: défini] plus simplement, comme noyau du composé $\overline{\Theta} : \mathcal{M}_{\rho,\sigma} \to G \to \operatorname{PGL}(2)$ dans la suite exacte $1 \to \mathcal{R}_{\rho,\sigma} \to \mathcal{M}_{\rho,\sigma} \to \operatorname{PGL}(2) \to 1$\leqno{(81)}} --- de $G_m \times [unclear: fl] \times [unclear: fl]$ par $\mathcal{M}_{\rho,\sigma} \to G_{aff} \times [unclear: fl]$ dans le [unclear: dictionnaire] avec $\mathfrak{S}_{\rho,\sigma}^+$ évidemment égale à $[unclear: fl]$, puisque (en fait on a)
$$ \mathcal{R}_{\rho,\sigma}^\circ = \mathcal{R}_{\rho,\sigma} \cap \mathcal{M}_{\rho,\sigma}^\circ = (\mathcal{R}_{\rho,\sigma})^\circ \quad \text{(composante neutre)} $$
(il faudrait peut-être noter $\mathcal{R}_{\rho,\sigma}^\sim$ plutôt $\mathcal{R}_{\rho,\sigma}$, et l'ancien $\mathcal{R}_{\rho,\sigma}$ deviendrait $\mathcal{R}_{\rho,\sigma}^\circ$, de sorte que $\mathfrak{S}_{\rho,\sigma}$ deviendrait $\mathfrak{S}_{\rho,\sigma}^\circ$ ...).
Je ne vais pas expliciter ici la construction de $\mathcal{R}_{\rho,\sigma}^\circ$ en termes de $\mathcal{T}_\rho, N_{\rho}^\circ$ etc., n'étant pas sûr qu'il y en aura besoin par la suite.

\newpage

%% ===== Page 430 =====

Dans le cas modulaire, sur la base absolue $\mathbf{Z}$, on a

$$
\begin{tikzcd}[row sep=tiny]
S\mathfrak{P}_{p,\sigma}^\circ(\mathbf{Z}) \ar[r, "\isom"] \ar[d, equal] & \mathfrak{T}_{p,\sigma}^\circ(\mathbf{Z}) \simeq \mathfrak{T}_p(\mathbf{Z}) \times \mathfrak{T}_\sigma(\mathbf{Z}) \simeq \mathbf{Z}/6\mathbf{Z} \times \mathbf{Z}/4\mathbf{Z} \\
S\mathfrak{T}_p(\mathbf{Z}) \times S\mathfrak{T}_\sigma(\mathbf{Z}) & \text{car } \mathfrak{T}_p(\mathbf{Z}) = S\mathfrak{T}_p(\mathbf{Z}), \mathfrak{T}_\sigma(\mathbf{Z}) = S\mathfrak{T}_\sigma(\mathbf{Z})
\end{tikzcd} \leqno{(81)}
$$
d'où, via (74)

$$
\mathcal{R}_{p,\sigma}^\circ(\mathbf{Z}) \simeq (\mathbf{Z}/6\mathbf{Z} \times \mathbf{Z}/4\mathbf{Z}) \times \{\pm 1\} \leqno{(82)}
$$
les deux projections sur les deux facteurs du second membre correspondent justement aux hom.
$$
\mathcal{R}_{p,\sigma}^\circ \to \mathfrak{T}_{p,\sigma}^\circ, \quad \mathcal{R}_{p,\sigma}^\circ \to G_m
$$
de (74), de sorte que l'inclusion

$$
\mu_2(\mathbf{Z}) = \{\pm 1\} \hookrightarrow \mathcal{R}_{p,\sigma}^\circ(\mathbf{Z}) \simeq (\mathbf{Z}/6\mathbf{Z} \times \mathbf{Z}/4\mathbf{Z}) \times \{\pm 1\} \leqno{(83)}
$$
est l'inclusion évidente $\varepsilon \mapsto (0) \times \varepsilon$.

Une façon plus parlante de dire ceci est que l'inclusion canonique

$$
\begin{tikzcd}[row sep=small, column sep=small]
S\mathfrak{P}_{p,\sigma}^\circ \ar[d, equal] &[-1em] \times &[-1em] \mu_2 \ar[r, hook] \ar[d, equal] & \mathcal{R}_{p,\sigma}^\circ \\
S\mathfrak{T}_p \times S\mathfrak{T}_\sigma \ar[d, hook] & & \operatorname{Centr}(S\mathcal{G}) \\
\mathcal{M}_{p,\sigma}^\circ(0)
\end{tikzcd} \leqno{(84)}
$$
induit un isomorphisme

$$
\begin{tikzcd}[row sep=tiny, column sep=tiny]
S\mathfrak{P}_{p,\sigma}^\circ(\mathbf{Z}) \ar[d, phantom, "\shortparallel"] & \times & \mu_2(\mathbf{Z}) \ar[r, "\isom"] \ar[d, phantom, "\shortparallel"] & \mathcal{R}_{p,\sigma}^\circ(\mathbf{Z}) \\
\mathbf{Z}/6\mathbf{Z} \times \mathbf{Z}/4\mathbf{Z} & & \{\pm 1\} &
\end{tikzcd} \leqno{(85)}
$$

\newpage

%% ===== Page 431 =====

[MARGIN: Commençons l'étude du radical $C R_{\rho,\sigma}$, $S R_{\rho,\sigma}$, $R_{\rho,\sigma}$, $S \dots$]

$6^o$) Je vais revenir sur

(86) $$R_{\rho,\sigma} \overset{\text{df}}{=} \operatorname{Ker}(\mathcal{M}_{\rho,\sigma} \overset{\overline{\Theta}}{\longrightarrow} \mathbb{P} G)$$

donnant lieu à la suite exacte

(87) $$1 \longrightarrow R_{\rho,\sigma} \longrightarrow \mathcal{M}_{\rho,\sigma} \longrightarrow \mathbb{P} G \longrightarrow 1 \ .$$

On a aussi

(88) $$R_{\rho,\sigma}^o \underset{\substack{\text{défini} \\ \text{dans 69}}}{=} R_{\rho,\sigma} \cap \mathcal{M}_{\rho,\sigma}^o = \substack{\text{comp. neutre} \\ \text{de } R_{\rho,\sigma}}$$

D'ailleurs on a vu

(89) $$R_{\rho,\sigma} = \Theta^{-1}(\mathbb{G}_m) \quad (\Theta: \mathcal{M}_{\rho,\sigma} \to G)$$

d'où une suite exacte

(90) 
$$1 \longrightarrow S R_{\rho,\sigma} \longrightarrow R_{\rho,\sigma} \overset{\mathbf{k}}{\longrightarrow} \mathbb{G}_m \longrightarrow 1$$
$$ \Vert \text{déf}$$
$$\text{Ker } \mathbf{k} \, (R_{\rho,\sigma} \overset{\mathbf{k}}{\longrightarrow} \mathbb{G}_m)$$
$$ \Vert$$
$$\text{Ker } \Theta \, (\mathcal{M}_{\rho,\sigma} \overset{\Theta}{\longrightarrow} G)$$

où $\mathbf{k}$ est défini par la condition que

(91)
\begin{tikzcd}
R_{\rho,\sigma} \ar[r, "\mathbf{k}"] \ar[drrr, "\Theta"', bend right] & \mathbb{G}_m \ar[r, "\simeq"] & \text{Centr}(G) \ar[r, hook] & G
\end{tikzcd}

soit commutatif.

Notons en passant que l'hom. $\chi$ de (42) $\mathcal{M}_{\rho,\sigma} \overset{\chi}{\longrightarrow} \mathbb{G}_m$ peut se déduire aussi de $\Theta$ (cf (57) par composition) :

(92)
\begin{tikzcd}
\mathcal{M}_{\rho,\sigma} \ar[r, "\Theta"] \ar[dr, "\chi"'] & G \ar[d, "\text{det}"] \\
& \mathbb{G}_m
\end{tikzcd}

et on a donc

(93) $$\chi(g) = \mathbf{k}(g)^2 \quad \text{si } g \in R_{\rho,\sigma}$$

i.e.

\begin{tikzcd}
R_{\rho,\sigma} \ar[r, "\mathbf{k}"] \ar[dr, "\chi | R_{\rho,\sigma}"'] & \mathbb{G}_m \ar[d, "\lambda \mapsto \lambda^2"] \\
& \mathbb{G}_m
\end{tikzcd}
commutatif.

\newpage

%% ===== Page 432 =====

On a le diagramme d'inclusions
(94)
\begin{tikzcd}
\underset{\text{"Ker }\kappa\text{"}}{SR_{\rho,\sigma}} \ar[r, hook, "\mathbb{G}_m \ (90)"] \ar[d, "\mu_2 \times \mu_2"'] & \underset{\text{"Ker }\Theta\text{"}}{R_{\rho,\sigma}} \ar[r, hook, "PG \ (86)"] \ar[d, "\text{cart et cocart}"'] & M_{\rho,\sigma} \ar[d, "\mu_2 \times \mu_2"] \\
\underset{\text{"Ker }\kappa^\circ\text{" \ (62)}}{\underset{\simeq S T_\rho \times S T_\sigma}{SR_{\rho,\sigma}^\circ}} \ar[r, hook, "\mathbb{G}_m"] & \underset{\text{"Ker }\Theta^\circ\text{" \ (69)}}{R_{\rho,\sigma}^\circ} \ar[r, hook, "PG"] & M_{\rho,\sigma}^\circ
\end{tikzcd}

[MARGIN: NB il y a une vérification à faire, qui se réduit à voir que les flèches $SR_{\rho,\sigma}^+ \to SR_{\rho,\sigma}$, $R_{\rho,\sigma}^+ \to R_{\rho,\sigma}$ sont des isomorph. — oui, cf plus bas en (100) où la flèche est rendue explicite par ex.]

La deuxième ligne est formée des composantes neutres de la première, qu'on peut réécrire plus joliment encore

(95)
\begin{tikzcd}
& R_{\rho,\sigma} \ar[dr, "PG"] \ar[dd] \\
SR_{\rho,\sigma} \ar[ur, "\mathbb{G}_m"] \ar[rr, dashed, "G"] \ar[dd, "\mu_2 \times \mu_2"'] & & M_{\rho,\sigma} \ar[dd, "\mu_2 \times \mu_2"] \\
& R_{\rho,\sigma}^\circ \ar[dr, "PG" near start] \\
\underset{(S T_\rho \times S T_\sigma)^-}{\underset{\parallel}{SR_{\rho,\sigma}^\circ}} \ar[ur, "\mathbb{G}_m"] \ar[rr, dashed, "G"'] & & M_{\rho,\sigma}^\circ
\end{tikzcd}

sous forme d'un prisme dont les trois carrés latéraux sont à la fois cartésiens et cocartésiens — donc donnent lieu à des quotients (sur des cotés parallèles) isomorphes — les quotients qui interviennent étant $\mu_2 \times \mu_2$, $\mathbb{G}_m$, $PG$ et l'extension $G$ de $

\newpage

%% ===== Page 433 =====

510

qui explicite la structure de $R_{p,\sigma}$ en termes
des groupes très simples $ST_p, ST_\sigma, \mathbb{G}_m, \mu_2$.

On a un cube, construit sur la face
verticale postérieure gauche de (95) (qui devient la
face verticale postérieure)

(95)
\begin{tikzcd}
S R_{p,\sigma} \ar[rr, "\mathbb{G}_m", hook] \ar[dd, "\mu_2 \times \mu_2"'] \ar[dr, "\mu_2"] & & R_{p,\sigma} \ar[rr, dashed] \ar[dd, "\mu_2 \times \mu_2"'] \ar[dr, "\mu_2"] & & \mathbb{G}_m \ar[dd, equal] \ar[dr, dashed, "\lambda \mapsto \lambda^2"] \\
& T_{p,\sigma}^+ \ar[rr, "\mathbb{G}_m", hook, crossing over] \ar[dd] & & T_{p,\sigma} \ar[rr, dashed, crossing over] \ar[dd] & & \mathbb{G}_m \ar[dd, equal] \\
S R_{p,\tilde{\sigma}} \ar[rr, "\mathbb{G}_m"', hook] \ar[dr, "\mu_2"] & & R_{p,\tilde{\sigma}} \ar[rr, dashed] \ar[dr, "\mu_2"] & & \mathbb{G}_m \ar[dr, dashed, "\lambda \mapsto \lambda^2"'] \\
& T_{p,\tilde{\sigma}}^+ \ar[rr, "\mathbb{G}_m"', hook] & & T_{p,\tilde{\sigma}} \ar[rr, dashed] & & \mathbb{G}_m
\end{tikzcd}

dont les faces verticales (antérieure et
postérieure) sont cartésiennes et cocartésiennes
et les flèches horizontales qui vont de l'une à
l'autre sont des morphismes de passage
au quotient, ayant tous des noyaux isomorphes
soit à $\mu_2$ (noyaux des flèches), soit à 1
(noyaux des autres flèches), ce qui
explicite les relations :

(97)
$\begin{cases}
R_{p,\sigma} \cap S G = R_{p,\sigma} \cap S h = \mu_2 \\
S R_{p,\sigma} \cap S h = S R_{p,\sigma} \cap S G = \{1\} \quad \text{explicite bien ça}
\end{cases}$

Pour la première ligne il s'agit de voir que
[unclear] est vrai dans la première ligne
ce qui équivaut à dire que l'on a une inclusion
$$\mu_2 = Sh_{p,\sigma} \cap S G \subset R_{p,\sigma} \cap S G$$

\newpage

%% ===== Page 434 =====

on note
$$ \underset{\Ker \tilde{\Theta}}{R_{\rho,\sigma} \cap SG} = \operatorname{Ker}(\tilde{\Theta} \mid SG), \quad \underset{\Ker \tilde{\Theta}^\circ}{R_{\rho,\sigma}^\circ \cap SG} = \operatorname{Ker}(\tilde{\Theta}^\circ \mid SG), $$
où $\tilde{\Theta} \mid SG = \tilde{\Theta}^\circ \mid SG = \text{morphisme can } SG \to PG$
dont le noyau est bien $\mu_2$.

D'autre part
$$ \underset{\Ker \Theta}{SR_{\rho,\sigma} \cap S\ell} = \operatorname{Ker}(\Theta \mid S\ell) = \{1\} $$
puisque $\Theta \mid S\ell$ est l'inclusion can.\ $S\ell \hookrightarrow G$,
a fortiori $SR_{\rho,\sigma}^\circ \cap S\ell = \{1\}$. Le fait qu'[unclear: il y a] des isomorphismes
$$ \begin{cases} SR_{\rho,\sigma} \isommap \tau_{\rho,\sigma}^+ \simeq S'T_\rho \times W_\sigma' \\ SR_{\rho,\sigma}^\circ \isommap \tau_{\rho,\sigma}^{+ \circ} \simeq ST_\rho \times ST_\sigma \end{cases} \leqno{(98)} $$
résulte, pour la deuxième, de la comparaison de (62) et de (28'), pour la première elle s'en déduit par le lemme du cinq, puisque
$$ SR_{\rho,\sigma} / SR_{\rho,\sigma}^\circ \isommap \tau_{\rho,\sigma}^+ / \tau_{\rho,\sigma}^{+ \circ} \simeq \mu_2 \times \mu_2 . $$

Le fait que l'un et l'autre des iso (98)
impliquent que dans
$$ M_{\rho,\sigma}^+ \simeq \underset{\operatorname{Ker}(M_{\rho,\sigma}^+ \to \tau_{\rho,\sigma}^+)}{\underline{M}_{\rho,\sigma}^+(0)} \underset{SG}{\cdot} \tau_{\rho,\sigma}^+ \quad \text{([unclear: produit fibré] déduit de (49))} $$
regardé comme extension de $\tau_{\rho,\sigma}^+$ par $SG$,
il y a, à côté du scindage par $\underline{M}_{\rho,\sigma}^+(0) \subset M_{\rho,\sigma}^+$
(fourni par $\underline{M}_{\rho,\sigma}(0)$), un autre scindage,

\newpage

%% ===== Page 435 =====

liée de façon beaucoup plus naturelle à la
structure du groupe $M_{\rho,\sigma}$ et à ses hom.
[crossed out: D...s] $G$ et $\gamma = Out \Theta$ dans $\hat{\Gamma}_{n,1}$, savoir le
sous-groupe section $S R_{\rho,\sigma}$, permettant de l'écrire
dans $M_{\rho,\sigma}^+$ comme un produit direct
$$ (99) \quad M_{\rho,\sigma}^+ \simeq \underset{\parallel S' T_{\rho} \times S N'_{\sigma}}{S R_{\rho,\sigma}} \times \underset{\parallel S G}{\mathfrak{S}_{\rho,\sigma}^+} $$

induisant une décomposition analogue

[MARGIN: c'est (45), par le biais d'un autre dévissage immédiat]

$$ (100) \quad M_{\rho,\sigma}^{+0} \simeq \underset{S T_{\rho} \times S T_{\sigma}}{S R_{\rho,\sigma}^0} \times \underset{S G}{\mathfrak{S}_{\rho,\sigma}^{+0}} $$
$$ \simeq S T_{\rho} \times S T_{\sigma} \times S G $$

qui dit mieux !. D'autre part, le fait
que $R_{\rho,\sigma} \to \Gamma_{\rho,\sigma}$ soit un isomorphisme
nous montre que l'on a, à isomorphisme
près, une "section" [crossed out text] de $M_{\rho,\sigma}$ sur $\Gamma_{\rho,\sigma}$ en tant
qu'extension de $\Gamma_{\rho,\sigma}$ par $\mathfrak{S}_{\rho,\sigma}^+ \simeq S G$.

En fait, on a
$$ R_{\rho,\sigma} \cap S G = \{1\} $$
[crossed out text]
$R_{\rho,\sigma}$ commute à $S G$
(trivial par définition de
$R_{\rho,\sigma}$ comme $\ker \tilde{\Theta}$)
donc on a une suite exacte

\newpage

%% ===== Page 436 =====

510

(102)
$$ 1 \longrightarrow \mu_2 \longrightarrow R_{\rho,\sigma} \times S G \longrightarrow M_{\rho,\sigma} \longrightarrow 1 $$
$$ \varepsilon \longmapsto (\varepsilon, \varepsilon^{-1}) $$

et de même

(103)
$$ 1 \longrightarrow \mu_2 \longrightarrow R_{\rho,\sigma}^+ \times S G \longrightarrow M_{\rho,\sigma}^+ \longrightarrow 1 $$

déjà envisagée (un peu timidement) dans (78).
C'est le moment décidément d'expliciter
entièrement la description de $R_{\rho,\sigma}$ en
termes de $T_\rho, N_\sigma'...$ -- on l'obtient, dans
le contexte de la p. 510, comme

(104)
$$ R_{\rho,\sigma} = \left\{ (\underset{\in}{\underset{\Gamma_{T_0}}{u}}, \underset{\in}{\underset{\Gamma_{T_\rho}}{\underline{a}}}, \underset{\in}{\underset{\Gamma_{N'_\sigma}}{\underline{b}}}, \underset{\in}{\underset{SG}{f}}) \left| \begin{array}{l} \det u = \det \underline{b} = \varepsilon' \det \underline{a},\ \varepsilon'^2=1 \\ u f = \lambda \\ (\text{d'où } \lambda^2 = \det u) \end{array} \quad (\varepsilon',\lambda \in \Gamma_{\mathbb{G}_m}) \right. \right\} $$
$$ = \left\{ (\underset{\in}{\underset{T_0}{u}}, \underset{\in}{\underset{T_\rho}{\underline{a}}}, \underset{\in}{\underset{N'_\sigma}{\underline{b}}}, \underset{\in}{\underset{\mathbb{G}_m}{\lambda}}) \left| \det u = \det \underline{b} = \varepsilon' \det \underline{a} = \lambda^2 \right. \right\} $$
$$ \simeq \left\{ (\underset{\in}{\underset{T_\rho}{\underline{a}}}, \underset{\in}{\underset{N'_\sigma}{\underline{b}}}, \underset{\in}{\underset{\mathbb{G}_m}{\lambda}}) \left| \det \underline{b} = \lambda^2,\ \det \underline{a}^2 = \det \underline{b}^2 \right. \right\} $$
$$ \simeq (\widetilde{T}_\rho, \delta_\rho^2) \times_{\mathbb{G}_m} (\widetilde{N}'_\sigma, \delta_\sigma^2) \times_{\mathbb{G}_m (\delta^2)} \mathbb{G}_m (\delta^2) $$

ou plus simplement ,

[raturé]
$$ R_{\rho,\sigma} \simeq \left\{ (\underset{\in}{\underset{T_\rho}{\underline{a}}}, \underset{\in}{\underset{N'_\sigma}{\underline{b}}}, \underset{\in}{\underset{\mathbb{G}_m}{\lambda}}) \left| \det \underline{a} = \det \underline{b} = \lambda^2 \right. \right\} $$

(105)
\begin{tikzcd}
R_{\rho,\sigma} \ar[r] \ar[d] & \widetilde{C}_{\rho,\sigma} \ar[d] \\
\mathbb{G}_m \ar[r, "\lambda \mapsto \lambda^2"'] & \mathbb{G}_m
\end{tikzcd}
diag. cartésien

\newpage

%% ===== Page 437 =====

s'insère dans un diagramme à suites exactes

\begin{tikzcd}
& & & 1 \ar[d] & 1 \ar[d] & \\
& & & S R_{\rho,\sigma} \ar[r, "\sim"] \ar[d] & \mathcal{C}^+_{\rho,\sigma} \ar[d] & \\
(105) & 1 \ar[r] & \mu_2 \ar[r] \ar[d, equal] & R_{\rho,\sigma} \ar[r, "can."] \ar[d, "\pi"] & \mathcal{C}_{\rho,\sigma} \ar[r] \ar[d, "\chi"] & 1 \\
& 1 \ar[r] & \mu_2 \ar[r] & \mathbb{G}_m \ar[r, "\lambda \mapsto \lambda^2"] \ar[d] & \mathbb{G}_m \ar[r] \ar[d] & 1 \\
& & & 1 & 1 & 
\end{tikzcd}

qui n'est qu'une explicitation de la
fin supérieure (*) [MARGIN: (*) plus au-dessus] des diagrammes analogues
(95) -- faire attention que celle-ci s'est
construite [plus bas] sur l'hom.
(107) groupe $\mathcal{C}_{\rho,\sigma}$, et $(\mathcal{C}_{\rho,\sigma} \xrightarrow{\chi} \mathbb{G}_m, \mathcal{C}_{\rho,\sigma} \xrightarrow{(\varepsilon_g, \varepsilon_d)} \mu_2 \times \mu_2)$
[ et par là exprime : (105). - (106) $R_{\rho,\sigma}$,
$S R_{\rho,\sigma} \simeq \mathcal{C}^+_{\rho,\sigma}$, et aussi (en [unclear: remplaçant]
$\mathcal{C}^0_{\rho,\sigma} = \operatorname{Ker}(\mathcal{C}_{\rho,\sigma} \to \mu_2 \times \mu_2)$
$R^0_{\rho,\sigma}, S R^0_{\rho,\sigma} \simeq \mathcal{C}^0_{\rho,\sigma}$ ] -- en termes des
groupes
(108) $T_{\rho}, N'_{\sigma}$ et de $T_{\rho} \xrightarrow{\delta_{\rho}} \mathbb{G}_m, N'_{\sigma} \xrightarrow{\delta_{\sigma}} \mathbb{G}_m$
$N'_{\sigma} \xrightarrow{\varepsilon_{\sigma}} \mu_2$

\newpage

%% ===== Page 438 =====

511

comme sous-groupe

(109)
\begin{tikzcd}
\mathcal{T}_{\rho, \sigma} \ar[r, hook] \ar[d, equal] & T_0 \times T_\rho \times N'_\sigma \\
\{(\underline{u}, \underline{a}, \underline{b}) \mid \delta_0(\underline{u}) = \delta_\rho(\underline{a}) = \varepsilon' \delta_\sigma(\underline{b}) , {\varepsilon'}^2 = 1 \} & T_0 \ar[d, "\delta_0"] \\
\simeq \{(\underline{a}, \underline{b}) \in T_\rho \times N'_\sigma \mid \delta_\rho(\underline{a}) = \varepsilon' \delta_\sigma(\underline{b}) , {\varepsilon'}^2 = 1 \} & \mathcal{G}_m
\end{tikzcd}

ou

(110)
$$
\begin{cases}
\chi(\underline{a}, \underline{b}) = \delta_\rho(\underline{a}) \\
\varepsilon_3(\underline{a}, \underline{b}) = \delta_\rho(\underline{a}) / \delta_\sigma(\underline{b}) \\
\varepsilon_1(\underline{a}, \underline{b}) = \varepsilon_\sigma(\underline{b})
\end{cases}
$$

A l'aide de ceci on peut construire $M_{\rho, \sigma}$
par (102), et bien entendu l'hom.

(111) $$ \underline{\theta} : M_{\rho, \sigma} \longrightarrow \underline{G} $$

se déduit par passage au quotient des hom.
évidents

(112)
\begin{tikzcd}
S\underline{G} \ar[r, hook] & \underline{G} , & R_{\rho, \sigma} \ar[r] \ar[dr, "\kappa"'] & \underline{G} \\
& & & \underline{G}_m \ar[u, "\simeq"'] \ar[r, "\simeq"] & \underline{C}ent(\underline{G})
\end{tikzcd}

On vérifie que on en déduit les compatibilités

(113)
\begin{tikzcd}
M_{\rho, \sigma} \ar[r, "\underline{\theta}"] \ar[dr, "\underline{\tilde{\theta}}"'] & \underline{G} \ar[d] \\
& P\underline{G}
\end{tikzcd}

$$ M_{\rho, \sigma} \overset{\underline{\theta}}{\longrightarrow} \underline{G} \overset{det}{\longrightarrow} \underline{G}_m $$

et on en déduit la suite exacte

(114)
\begin{tikzcd}
1 \ar[r] & S\underline{G} \ar[r] \ar[dr, "can"'] & M_{\rho, \sigma} \ar[r] & \mathcal{T}_{\rho, \sigma} \ar[r] \ar[d, equal] & 1 \\
& & R_{\rho, \sigma} \times S\underline{G} \ar[u] & R_{\rho, \sigma} / \mu_2
\end{tikzcd}

\newpage

%% ===== Page 439 =====

de la description (102), d'où des composés

(115)
\begin{tikzcd}
& & T_\rho \\
M_{\rho, \sigma} \ar[r] & T_{\rho, \sigma} \ar[ur, "pr_1"] \ar[r, "pr_2"'] \ar[dr] & N'_\sigma \ar[r, "det"] & \mathbb{G}_m \\
& & \mu \times \mu_2
\end{tikzcd}

Dans tout ceci, on n'a pas eu besoin de la
donnée de sections $\rho$, $\sigma$ de $G$, ni de
connaître quoi que ce soit sur la façon
dont $T_\rho$, $N'_\sigma$ sont [unclear: obtenus] à partir de telles
sections. Mais des structures supplémentaires
plus délicates dans $M_{\rho,\sigma}$ - telle la
donnée d'un sous-groupe section $M_{\rho,\sigma}(s)$ -
en dépendent, et permettront peut-être de
les retrouver.

Rmq. Le cadre (96) est le "compagnon"
d'un diagramme analogue, que je
ne sais expliciter, dont le [unclear]
complète (95) - qui peut être considéré
comme un diagramme partiel de
celui-ci :

\newpage

%% ===== Page 440 =====

Comme d. (95) (96), 6 faces inf. 
et la face sup. du prisme ne sont pas des
carrés cartésiens (ni cocartésiens), les quatre groupes
quotients s'identifient canoniquement entre eux par les
flèches en jeu, ils sont isomorphes à 
$M_{\rho,\sigma}^0$ --

(116)
\begin{tikzcd}[row sep=1.5em, column sep=1.5em]
G_m \ar[ddd, "\lambda \mapsto \lambda^2"'] \ar[drr, dashed] & & & G_m \ar[ddd, "\lambda \mapsto \lambda^2"'] \ar[drr, dashed] & & & G_m \ar[ddd, "\lambda \mapsto \lambda^2"'] \ar[ddr, dashed] & & \\
& & S R_{\rho,\sigma} \ar[rrr, hook] \ar[ddd, "pr_1 \times pr_2" near start] \ar[dl] & & & R_{\rho,\sigma} \ar[rrr] \ar[ddd, "pr_1 \times pr_2" near start] \ar[dl] & & & P G \ar[ddd, "id"] \ar[dl] \\
& S G \ar[rrr, crossing over, hook] \ar[ddd, "id"'] & & & G \ar[rrr, crossing over] \ar[ddd, "id"'] & & & M_{\rho,\sigma} \ar[ddd, "pr_1"'] & \\
G_m \ar[drr, dashed] & & & G_m \ar[drr, dashed] & & & G_m \ar[ddr, dashed] & & \\
& & S R_{\rho,\sigma}^0 \ar[rrr, hook] \ar[dl] & & & R_{\rho,\sigma}^0 \ar[rrr] \ar[dl] & & & P G \ar[dl] \\
& S G \ar[rrr, hook] & & & G \ar[rrr] & & & M_{\rho,\sigma}^0 & 
\end{tikzcd}

Dans le cube ci-dessous (117), les faces latérales du cube sont
des carrés cartésiens, mais pas les faces
inf. et sup. en général.

\newpage

%% ===== Page 441 =====

(117)
\begin{tikzcd}[row sep=2em, column sep=2em]
& R_{\rho, \sigma} \ar[dl] \ar[d, "\mathbb{G}_m"', dashed] \ar[dr] \ar[dddd, "pr"'] & \\
S R_{\rho, \sigma} \ar[d] \ar[drr, "\simeq" near start] \ar[dddd, "pr"'] & G \ar[dl] \ar[r] \ar[dr] \ar[dddd, "pr"'] & P G \ar[dd] \\
S G \ar[dr] \ar[dddd, "pr"'] & & M_{\rho, \sigma} \ar[dl, "+"] \ar[dr] \ar[dddd, "pr"'] & \\
& \mathcal{M}_{\rho, \sigma}^+ \ar[dddd, "pr"'] & & C_{\rho, \sigma} \ar[dddd, "pr"'] \\
& R_{\rho, \sigma}^0 \ar[dl] \ar[d, "\mathbb{G}_m"', dashed] & \\
S R_{\rho, \sigma}^0 \ar[d] \ar[drr, "\simeq" near start] & G^0 \ar[dl, dashed] \ar[dr] & \\
S G \ar[dr] & & M_{\rho, \sigma}^0 \ar[dl, "+ 0"] \ar[dr] & \\
& \mathcal{M}_{\rho, \sigma}^{+ 0} & & C_{\rho, \sigma}^0
\end{tikzcd}

La face comprise entièrement à gauche de ce cube est de nature particulièrement simple, puisque les termes médians $\mathcal{M}_{\rho, \sigma}^+$, $\mathcal{M}_{\rho, \sigma}^{+ 0}$ s'identifient explicitement aux produits directs des tronçons $S R_{\rho, \sigma} \simeq \mathcal{T}_{\rho, \sigma}^+$ et $S R_{\rho, \sigma}^0 \simeq \mathcal{T}_{\rho, \sigma}^{+ 0}$
$$ \simeq S T_{\rho} \times S N_{\sigma}' \quad \text{et} \quad \simeq S T_{\rho} \times S T_{\sigma} $$
avec le groupe $S G$. C'est pourquoi le [unclear: as-is] pour la face parallèle (face postérieure droite), on gagne [unclear: le prix] - [unclear: à condition] de remplacer les groupes
médians $M_{\rho, \sigma}$, $M_{\rho, \sigma}^0$ par les pr-revêtements

(118) $$ \mathcal{M}_{\rho, \sigma} \simeq R_{\rho, \sigma} \times S G, \quad \mathcal{M}_{\rho, \sigma}^0 \simeq R_{\rho, \sigma}^0 \times S G $$

cf (102), (103)) , ce qui [unclear: assure] la [unclear: suite].

\newpage

%% ===== Page 442 =====

[unclear: discutable] de rendre les quatre carrés des faces supérieures et inférieures des deux cubes accolés cartésiens et cocartésiens, ce qu'ils n'étaient pas avant, regrettablement.

Pour terminer ce grimoire sur le radical, je vais introduire un homomorphisme canonique
$$
G_m \overset{i_R}{\hookrightarrow} R_{p,\sigma}^\circ \leqno{(118)}
$$
via un homomorphisme canonique
$$
\begin{aligned}
G_m &\overset{i_{\mathfrak{T}}}{\hookrightarrow} \mathfrak{T}_{p,\sigma}^\circ \simeq \mathfrak{T}_p \times_{G_m} \mathfrak{T}_\sigma \\
\lambda &\longmapsto (\lambda, \lambda)
\end{aligned}
\leqno{(119)}
$$
dont la composée avec
$$
y: \mathfrak{T}_{p,\sigma}^\circ \longrightarrow G_m
$$
qui donne en effet $\lambda \longmapsto \lambda^2$, de sorte que grâce à la description (104) ou (105) de $R_{p,\sigma}$ en termes de $\mathfrak{T}_{p,\sigma}$ comme produit fibré (paraphrasé en un isomorphisme de $R_{p,\sigma}^\circ$ au dessus de $\mathfrak{T}_{p,\sigma}^\circ$) on a un unique homomorphisme (118), rendant commutatif le diagramme

\newpage

%% ===== Page 443 =====

52-9

(120)
\begin{tikzcd}
\mathbb{G}_m \ar[r, "i_\tau"] \ar[dr, "i_R"'] & \mathcal{T}_{\rho,\sigma}^\circ \\
& R_{\rho,\sigma}^\circ \ar[u, "can"']
\end{tikzcd}

-, plus joliment, le diagramme en
suite exacte

(121) $$1 \longrightarrow \mu_2 \longrightarrow \mathbb{G}_m \xrightarrow{\lambda \mapsto \lambda^2} \mathbb{G}_m \longrightarrow 1$$ .

Par construction, on trouve aussitôt

(122)
\begin{tikzcd}
\mathbb{G}_m \ar[r, "i_R"] \ar[dr, "id"'] & R_{\rho,\sigma}^\circ \ar[d, "k"] \\
& \mathbb{G}_m
\end{tikzcd}

d'où des décompositions en produit direct

(123)
$$
\begin{cases}
R_{\rho,\sigma} \simeq S R_{\rho,\sigma} \times \mathbb{G}_m \simeq S' \mathcal{T}_\rho \times S \mathcal{N}_\sigma' \times \mathbb{G}_m \\
R_{\rho,\sigma}^\circ \simeq S R_{\rho,\sigma}^\circ \times \mathbb{G}_m \\
\quad \simeq S \mathcal{T}_\rho \times S \mathcal{T}_\sigma \times \mathbb{G}_m
\end{cases}
$$

qui précise considérablement (96) dont j'étais
fort peu satisfait !

En termes de cette décomposition, on
explicite l'hom.
$$ \mu_2 \hookrightarrow R_{\rho,\sigma}^\circ \subset R_{\rho,\sigma} $$ (et $R_{\rho,\sigma} \twoheadrightarrow \mathcal{T}_\sigma$, $R_{\rho,\sigma}^\circ \twoheadrightarrow \mathcal{T}_\sigma^\circ$)
qui intervient dans (102), (103), $\mu_2 = R_{\rho,\sigma}^\circ \cap S R_{\rho,\sigma}^\circ = R_{\rho,\sigma} \cap S R_{\rho,\sigma}$
~~[unclear crossed out text]~~

\newpage

%% ===== Page 444 =====

Pour l'hom. 
$$R_{\rho,\sigma} \to T_{\rho,\sigma} \subset T_\rho \times N_\sigma$$
il suffit d'expliciter les composées
(124)
$$
\begin{tikzcd}
R_{\rho,\sigma} \ar[r] \ar[d, "\simeq"'] & T_{\rho,\sigma} \ar[r] \ar[dr] & T_\rho \\
S' T_\rho \times S N'_\sigma \times \mathbb{G}_m & & N_\sigma
\end{tikzcd}
$$
$$(a,b,\lambda) \mapsto (\lambda a, \lambda b)$$
[MARGIN: on a en effet $\det(\lambda a) = \det(\lambda b)^2$, $\lambda^2 = \lambda^2$ car $\det(a)=1, \det \underline{b}=1$]

Le noyau de cet hom. est donc
formé des $(\lambda^{-1}, \lambda^{-1}, \lambda)$, avec $\lambda \in \mathbb{G}_m$ tel que
$\lambda^{-1} \in S'T_\rho$, $\lambda^{-1} \in S N'_\sigma$, c.a.d. i.e. $(\lambda^2)^? = 1$ et $\lambda^2 = 1$,
soit $\lambda \in \mu_2$. Donc

(115) [unclear: (125)]
$$
\begin{tikzcd}
\mu_2 \ar[r, "can"] & R_{\rho,\sigma}^\circ \ar[r, hook] & R_{\rho,\sigma} \\
R_{\rho,\sigma} \cap SG \ar[u, phantom, sloped, "\simeq"] & S T_\rho \times S T_\sigma \times \mathbb{G}_m \ar[u, phantom, sloped, "\simeq"] & S' T_\rho \times S N'_\sigma \times \mathbb{G}_m \ar[u, phantom, sloped, "\simeq"]
\end{tikzcd}
$$
$$\lambda \mapsto (\lambda, \lambda, \lambda)$$

On trouve donc l'explicitation particulièrement
simple
(126)  $$T_{\rho,\sigma} \simeq (S' T_\rho \times S N'_\sigma \times \mathbb{G}_m) / \mu_2$$
par envoi "diagonal" dans le
produit triple

On notera qu'on a des hom. canoniques
$$\mu_2 \to S T_\rho \subset S' T_\rho, \quad \mu_2 \to S T_\sigma \subset S N'_\sigma, \quad \mu_2 \to \mathbb{G}_m$$
(qui sont utilisés dans (125)) d'où
(127)
$$
\begin{tikzcd}
\mu_2 \times \mu_2 \times \mu_2 \ar[r] & S T_\rho \times S T_\sigma \times \mathbb{G}_m \ar[r, hook] & S' T_\rho \times S N'_\sigma \times \mathbb{G}_m \\
& R_{\rho,\sigma}^\circ \ar[u, phantom, sloped, "\simeq"'] & R_{\rho,\sigma} \ar[u, phantom, sloped, "\simeq"']
\end{tikzcd}
$$
$$(\lambda', \lambda'', \lambda) \mapsto (\lambda', \lambda'', \lambda)$$
qui n'est autre que (98)).

\newpage

%% ===== Page 445 =====

Le $\operatorname{Centr}$ de $M_{\rho,\sigma}$. Il s'envoie par $\widetilde{\Theta}$ dans le centre de $PG$ qui est trivial, donc
$$ \operatorname{Centr}(M_{\rho,\sigma}) \subset \Ker \widetilde{\Theta} = R_{\rho,\sigma} \simeq S' T_{\rho} \times S N'_{\sigma} \times \mathbf{G}_m $$
ceci évidemment. Comme $R_{\rho,\sigma}$ commute à $SG$, et que $R_{\rho,\sigma} \cdot SG = M_{\rho,\sigma}$, on en conclut
$$ \operatorname{Centr}(M_{\rho,\sigma}) = \operatorname{Centr}(R_{\rho,\sigma}) = S' T_{\rho} \times \operatorname{Centr}(S N'_{\sigma}) \times \mathbf{G}_m \leqno{(128)} $$
déterminons donc $\operatorname{Centr}(S N'_{\sigma})$. Notons que
$$
\begin{cases}
\mu_2 \simeq N'_{\sigma} / T_{\sigma} \text{ opère fidèlement sur } T_{\sigma} \\
\text{et } T_{\sigma} = \mathbf{G}_m \cdot S T_{\sigma}
\end{cases}
\text{ donc } \mu_2 \text{ opère fidèlement sur } S T_{\sigma}
$$
d'où
$$ \operatorname{Centr}(S N'_{\sigma}) \subset S T_{\sigma} \quad \text{d'où } \operatorname{Centr}(S N'_{\sigma}) \simeq (S T_{\sigma})^{\mu_2} \leqno{(129)} $$

\marginpar{crossed out in MS}Où on constate que le facteur $-1$ de $\mu_2$ opère sur $S T_{\sigma}$ par $x \mapsto x^{-1}$, ce qui montre que
$$ \operatorname{Centr}(S N'_{\sigma}) \subset (S T_{\sigma})_2 = \left\{ g = c \sigma + d \Biggm| \begin{array}{l} -u c^2 + d^2 = 1 \\ u c^2 + d^2 + 2\sigma c d = 1 \end{array} \right\} $$
On a pour une matrice $g \in S T_{\sigma}$, $g = c \sigma + d$
$$ \det g = -u c^2 + d^2 = 1 $$
$$ g^2 = u c^2 + d^2 + 2cd \sigma = 1 $$
i.e. 
$$ \begin{cases} -u c^2 + d^2 = 1 \\ u c^2 + d^2 = 1 \\ 2cd = 0 \end{cases} $$
d'où i.e. 
$$ \begin{cases} 2u c^2 = 0 \text{ donc } 2c^2 = 0 \text{ (car } u \text{ invers.)} \\ 2cd = 0 \\ -u c^2 + d^2 = 1 \end{cases} $$
si $2$ inv. cela revient à $c^2=0$, $d^2=1$ (donc $d$ invers.) donc $c=0$, $d = \pm 1$ i.e. $g \in \mu_2$.

d'où
$$ \operatorname{Centr}(S N'_{\sigma}) = \mu_2 = S T_{\sigma} \cap \mathbf{G}_m $$
Dans le cas général on aura simplement
$$ \mu_2 \subset \operatorname{Centr}(S N'_{\sigma}) \leqno{(130)} $$
et cette inclusion induit un iso. sur les [unclear: 2-groupes] de sections dans le cas où $2$ est inversible dans l'anneau de base.

\newpage

%% ===== Page 446 =====

[unclear: $\underline{C}_{es}$. Diagrammes 
(131) $\mu'_2 \subset S T_\sigma$
sous-groupe de $S T_\sigma$ d'où finales, qu'on a :]

(128) $$\underline{Cent} (\underline{M}_{\rho,\sigma}) \simeq S' T_\rho \times \mu'_2 \times \mathbb{G}_m \subset \underline{R}_{\rho,\sigma} \simeq S' T_\rho \times S N'_\sigma \times \mathbb{G}_m$$
d'où
(129) $$\underline{Cent} (\underline{M}_{\rho,\sigma}) \cap \underline{M}_{\rho,\sigma}^\circ = S T_\rho \times \mu'_2 \times \mathbb{G}_m \subset \underline{R}_{\rho,\sigma}^\circ \simeq S T_\rho \times S T_\sigma \times \mathbb{G}_m$$

Revenons aux sections entières dans le cas
modulaire absolu

(130) $$\underline{R}_{\rho,\sigma} (\mathbb{Z}) \simeq \underbrace{S' T_\rho (\mathbb{Z})}_{T_\rho(\mathbb{Z}) \simeq \mathbb{Z}/6\mathbb{Z}} \times \underbrace{S N'_\sigma (\mathbb{Z})}_{T_\sigma(\mathbb{Z}) \simeq \mathbb{Z}/4\mathbb{Z}} \times \underbrace{\mathbb{G}_m (\mathbb{Z})}_{\mu(\mathbb{Z}) \simeq \pm 1}$$

(131) $$\underline{Cent} (\underline{M}_{\rho,\sigma}) (\mathbb{Z}) \simeq S' (T_\rho [\mathbb{Z}]) \times \mu_2 (\mathbb{Z}) \times \mathbb{G}_m (\mathbb{Z})$$
$$\simeq \mathbb{Z}/6\mathbb{Z} \times \{\pm 1\} \times \{\pm 1\}$$
et pour mémoire (cf (75), p. 514)

(132) $$\underline{Cent} (\underline{M}_{\rho,\sigma}^\circ) (\mathbb{Z}) \simeq S T_\rho (\mathbb{Z}) \times S T_\sigma (\mathbb{Z}) \times \mathbb{G}_m (\mathbb{Z})$$
$$\underline{R}_{\rho,\sigma}^\circ \simeq \mathbb{Z}/6\mathbb{Z} \times \mathbb{Z}/4\mathbb{Z} \times \{\pm 1\}$$

En résumé, on a le diagramme d'inclusions (strictes)
de schémas en groupes

(133)
\begin{tikzcd}
\text{dim 3} & \begin{matrix} \underline{Cent} (\underline{M}_{\rho,\sigma}^\circ) = \underline{R}_{\rho,\sigma}^\circ \\ \simeq S T_\rho \times S T_\sigma \times \mathbb{G}_m \end{matrix} \ar[r, hook, "\text{(quotient } \mu_2 \times \mu_2)"'] & \begin{matrix} \underline{R}_{\rho,\sigma} \\ \simeq S' T_\rho \times S N'_\sigma \times \mathbb{G}_m \end{matrix} & \text{dim 3} \\
\text{dim 2} & \begin{matrix} \underline{Cent} (\underline{M}_{\rho,\sigma}) \cap \underline{M}_{\rho,\sigma}^\circ \\ = S T_\rho \times \mu'_2 \times \mathbb{G}_m \end{matrix} \ar[u, hook] \ar[r, hook, "\text{(quotient } \mu_2)"'] & \begin{matrix} \underline{Cent} (\underline{M}_{\rho,\sigma}) \\ \simeq S' T_\rho \times \mu'_2 \times \mathbb{G}_m \end{matrix} \ar[u, hook] & \text{dim 2}
\end{tikzcd}

\newpage

%% ===== Page 447 =====

valable dans le cas général - et qui, dans le cas modulaire absolu, induit dans les points entiers des isomorphismes par les deux flèches horizontales, correspondant resp. aux deux groupes
$$
\begin{cases} 
\operatorname{Cent}(\mathcal{M}_{p,\sigma}^\circ)(\mathbf{Z}) \simeq \mathbf{Z}/6\mathbf{Z} \times \mathbf{Z}/4\mathbf{Z} \times \{\pm 1\} \\ 
\operatorname{Cent}(\mathcal{M}_{p,\sigma})(\mathbf{Z}) \simeq \mathbf{Z}/6\mathbf{Z} \times \{\pm 1\} \times \{\pm 1\} 
\end{cases}
\leqno{(134)}
$$
qui en est un sous-groupe d'indice 2.

Enfin, pour mémoire signalons à nouveau le groupe
$$
\Sigma = \{\pm 1\} \times \{\pm 1\} \times \{\pm 1\}
\leqno{(135)}
$$
contenu dans $\operatorname{Cent}(\mathcal{M}_{p,\sigma})(\mathbf{Z})$ (d'ordre 8 et d'indice 3), qui joue un rôle particulier dans l'application arithmétique, puisqu'on aura un homomorphisme canonique
$$
\boxed{ \mathcal{M}_{0,3}^\sim \longrightarrow \mathcal{M}_{p,\sigma}(\hat{\mathbf{Z}})/\Sigma }
\leqno{(136)}
$$

\newpage

%% ===== Page 448 =====

[MARGIN: Bien décrire $\mathcal{M}_{\rho,\sigma}$]

7° Faisons une liste récapitulative des principaux homomorphismes de groupes envisagés au cours de la section 6°, impliquant le groupe $\mathcal{M}_{\rho,\sigma}$

(134)
\begin{tikzcd}[row sep=1.5em, column sep=3em]
& & & & PG \\
& SG \ar[dr, "j_G"] & & G \ar[ur] \ar[dr, "\xi = \det_G"'] \\
& S'T_\rho \ar[dr, "j'_\rho"] & & & \mathbb{G}_m \\
\mu_2 \ar[uuuur] \ar[uur] \ar[ddr] \ar[ddddr] & & \mathcal{M}_{\rho,\sigma} \ar[uuur, "\Theta = \Theta_G" near start] \ar[ur, "\Theta_\rho" near start] \ar[dr, "\Theta_\sigma" near start] \ar[dddr, "\varepsilon_3" near start] & T_\rho \ar[r, "\delta_\rho = \det_{T_\rho}"] & \mathbb{G}_m \\
& SN'_\sigma \ar[ur, "j'_\sigma"'] & & N'_\sigma \ar[r, "\delta'_\sigma = \det_{N'_\sigma}"] \ar[dr, "\varepsilon_\sigma"'] & \mathbb{G}_m \\
& & & & \mu_2 \\
& \mathbb{G}_m \ar[uuur, "j_\mu"'] & & \mu_2
\end{tikzcd}

Notons que ce diagramme se reconstruit (avec les propriétés essentielles que nous venons de repasser en revue) : à partir des données

$$
(135) \qquad
\begin{cases}
\mathbb{G}_m \xrightarrow{i_\rho} T_\rho \xrightarrow{\delta_\rho} \mathbb{G}_m & (\delta_\rho i_\rho(\lambda) = \lambda^2) \\
\mathbb{G}_m \xrightarrow{i_\sigma} N'_\sigma \xrightarrow{\delta'_\sigma} \mathbb{G}_m & (\delta'_\sigma i_\sigma(\lambda) = \lambda^2) \\
N'_\sigma \xrightarrow{\varepsilon_\sigma} \mu_2 & \varepsilon_\sigma i_\sigma(\lambda) = 1
\end{cases}
$$

[diagramme rayé]
\begin{tikzcd}
1 \ar[r] & \mu_2 \ar[r] \ar[d] & SG \ar[r] \ar[d] & PG \ar[r] \ar[d, equal] & 1 \\
1 \ar[r] & \mathbb{G}_m \ar[r] \ar[d, "\lambda \mapsto \lambda^2"'] & G \ar[r] \ar[d] & PG \ar[r] & 1 \\
& \mathbb{G}_m \ar[r, equal] & \mathbb{G}_m
\end{tikzcd}

plus de la donnée familière

(136)
\begin{tikzcd}
& \mathbb{G}_m \ar[r, equal] & \mathbb{G}_m \\
1 \ar[r] & \mathbb{G}_m \ar[r, "i_G"] \ar[u, "\lambda \mapsto \lambda^2"'] & G \ar[r] \ar[u, "\delta_G"'] & PG \ar[r] & 1 \\
1 \ar[r] & \mu_2 \ar[r] \ar[u, hook] & SG \ar[r] \ar[u, hook] & PG \ar[r] \ar[u, equal] & 1 \\
& 1 \ar[u] & 1 \ar[u]
\end{tikzcd}

\newpage

%% ===== Page 449 =====

Rappelons qu'on a de fait, en termes de (135)
$$
\begin{cases}
ST_p = \Ker \delta_p, \quad S'T_p = \Ker {\delta_p}^2 \\
T_\sigma = \Ker \varepsilon_\sigma, \quad SN'_\sigma = \Ker \delta'_\sigma, \quad ST_\sigma = SN'_\sigma \cap T_\sigma = \Ker \delta_\sigma \\
\hfill (\delta_\sigma = \delta'_\sigma \mid T_\sigma) \\
\mu_2 \hookrightarrow ST_p \subset S'T_p \quad \text{comme comp. de} \quad \mu_2 \subset \mathbb{G}_m \xrightarrow{j_p} T_p \\
\mu_2 \hookrightarrow ST_\sigma \subset SN'_\sigma \quad \text{------------------} \quad \mu_2 \subset \mathbb{G}_m \xrightarrow{j_\sigma} T_\sigma
\end{cases}
\leqno{(137)}
$$
ce qui détermine la partie gauche des diagrammes (134), où $\mathcal{M}_{g,\sigma}$ peut être considéré comme quotient de la somme amalgamée des quatre facteurs $SG, S'T_p, SN'_\sigma, \mathbb{G}_m$, par $\mu_2$ envoyé diagonalement.

Ceci permet en principe de réduire la description des quatre morphismes fondamentaux de $\mathcal{M}_{g,\sigma}$ dans (134) — $\Theta_G, \Theta_p, \Theta_\sigma$ et $\varepsilon_3$ — à leurs composés avec les quatre hom. $j_G, j_p, j_\sigma, j_\mu$ — quitte à vérifier que le produit des valeurs sur les images d'un même $\varepsilon \in \mu_2$ est trivial.
$$
\begin{cases}
\Theta_G j_G : SG \hookrightarrow G \quad \text{inclusion dans (136)} \\
\Theta_G j_p, \Theta_G j_\sigma \quad \text{triviaux} \\
\Theta_G j_\mu : \mathbb{G}_m \longrightarrow \mathbb{G}_m \quad \text{incl. can.}
\end{cases}
\leqno{(138)}
$$
(conditions de compatibilité satisfaites !)

\newpage

%% ===== Page 450 =====

(139)
$\begin{cases}
\Theta_P j_G, \Theta_P j_\sigma \text{ triviaux} \\
\Theta_P j_\rho : S' T_P \hookrightarrow T_P \text{ inclusion} \\
\Theta_P j_\mu : G_m \xrightarrow{i_P} T_P
\end{cases}$

compatibilités triviales,

(140)
$\begin{cases}
\Theta_\sigma j_G, \Theta_\sigma j_\rho \text{ triviaux} \\
\Theta_\sigma j_\sigma : S N'_\sigma \hookrightarrow N'_\sigma \text{ inclusion} \\
\Theta_\sigma j_\mu : G_m \xrightarrow{i_\sigma} N'_\sigma
\end{cases}$

compatibilités triviales,

(141)
$\begin{cases}
\varepsilon_3 j_G, \varepsilon_3 j_\sigma, \varepsilon_3 j_\mu \text{ triviaux} \\
\varepsilon_3 j_\rho : S' T_P \xrightarrow{\delta_P} \mu_2
\end{cases}$

On définit donc, en termes de (134)

(142)
$$
\begin{tikzcd}
M_{\rho, \sigma} \ar[r, "\varepsilon_2"] \ar[d, "\Theta_\sigma"'] & \mu_2 \\
N'_\sigma \ar[ur, "\varepsilon_\sigma"'] &
\end{tikzcd}
\quad \varepsilon_2 = \varepsilon_\sigma \circ \Theta_\sigma
$$

$$
\begin{tikzcd}
M_{\rho, \sigma} \ar[r, "\chi"] \ar[d, "\Theta_G"'] & G_m \\
G \ar[ur, "\delta_G=\det"'] &
\end{tikzcd}
\quad \chi = \delta_G \circ \Theta_G
$$

$$
\begin{tikzcd}
M_{\rho, \sigma} \ar[r, "\widetilde{G}"] \ar[d, "\Theta_G"'] & P G \\
G \ar[ur, "can"'] &
\end{tikzcd}
\quad \widetilde{G} = can \circ \Theta_G
$$

on aura donc
(143)
$$
M_{\rho, \sigma} \xrightarrow{(\chi, \varepsilon_3, \varepsilon_2)} G_m \times \mu_2 \times \mu_2
$$
) épimorphisme

\newpage

%% ===== Page 451 =====

et les autres deux composés $\delta_p \Theta_p$ et $\delta_0' \Theta_\sigma$ de (134) s'explicitent donc par les formules
$$
\begin{cases}
\delta_0' \Theta_\sigma = \chi \quad (\overset{dfn}{=} \delta_G \Theta_G) \\
\delta_p \Theta_p(g) = \varepsilon_\sigma(g) \chi(g)
\end{cases}
\leqno(144)
$$

En termes de (143) on définit
$$
\begin{cases}
\mathcal{M}_{p,\sigma}^+ = \Ker (\mathcal{M}_{p,\sigma} \xrightarrow{\chi} G_m) \\
\mathcal{M}_{p,\sigma}^\circ = \Ker (\mathcal{M}_{p,\sigma} \xrightarrow{(\varepsilon_3, \varepsilon_2)} \mu_2 \times \mu_2) \\
\mathcal{M}_{p,\sigma}^{+\circ} = \mathcal{M}_{p,\sigma}^+ \cap \mathcal{M}_{p,\sigma}^\circ = \Ker (\mathcal{M}_{p,\sigma} \xrightarrow{(\chi, \varepsilon_3, \varepsilon_2)} G_m \times \mu_2 \times \mu_2)
\end{cases}
\leqno(145)
$$

on définit
$$
\begin{cases}
R_{p,\sigma} = \Ker (\mathcal{M}_{p,\sigma} \xrightarrow{\widetilde{\Theta}} P G) \\
R_{p,\sigma}^\circ = \mathcal{M}_{p,\sigma}^\circ \cap R_{p,\sigma} = \Ker (\mathcal{M}_{p,\sigma}^\circ \xrightarrow{\widetilde{\Theta}^\circ} P G) \\
SR_{p,\sigma} = \Ker (\mathcal{M}_{p,\sigma} \xrightarrow{\Theta} G) \\
SR_{p,\sigma}^{+\circ} = \Ker (\mathcal{M}_{p,\sigma} \xrightarrow{(\Theta, \varepsilon_3, \varepsilon_2)} G \times \mu_2 \times \mu_2) \\
\quad = SR_{p,\sigma} \cap R_{p,\sigma}^\circ
\end{cases}
\leqno(146)
$$

d'où les carrés cartésiens et cocartésiens (94), et plus richement, les diagrammes commutatifs (96) (116) (117), et les dévissages\footnote{dont le sens est évident sur le diagramme (134)}
$$
\begin{cases}
R_{p,\sigma} \simeq S' T_p \times S N_\sigma' \times G_m \\
R_{p,\sigma}^\circ \simeq S T_p \times S T_\sigma \times G_m \\
SR_{p,\sigma} \simeq S' T_p \times S N_\sigma' \times 1 \\
SR_{p,\sigma}^\circ \simeq S T_p \times S T_\sigma \times 1
\end{cases}
\leqno(147)
$$

\newpage

%% ===== Page 452 =====

(148)
\begin{tikzcd}
1 \ar[r] & \mu_2 \ar[r] \ar[d, equal] & R_{\rho,\sigma} \times S G \ar[r] \ar[d, hook, "\mu_2 \times \mu_2"'] & M_{\rho,\sigma} \ar[r] \ar[d, hook, "\mu_2 \times \mu_2"'] & 1 \\
1 \ar[r] & \mu_2 \ar[r] & R_{\rho,\mathring{\sigma}} \times S G \ar[r] & M_{\rho,\mathring{\sigma}} \ar[r] & 1
\end{tikzcd}
[MARGIN: c'est ici la définition de $M_{\rho,\mathring{\sigma}}$]

$$
(149) \qquad
\begin{cases}
M_{\rho,\sigma}^+ \simeq S R_{\rho,\sigma} \times S G \\
M_{\rho,\mathring{\sigma}}^{+0} \simeq S R_{\rho,\mathring{\sigma}} \times S G
\end{cases}
$$

Enfin, définissons

$$
(150) \qquad
\begin{cases}
\mathcal{T}_{\rho,\sigma} \simeq M_{\rho,\sigma} / S G \\
\mathcal{T}_{\rho,\mathring{\sigma}} = M_{\rho,\mathring{\sigma}} / S G = \operatorname{Ker}(\mathcal{T}_{\rho,\sigma} \to \mu_2 \times \mu_2) \\
\mathcal{T}_{\rho,\sigma}^+ = M_{\rho,\sigma}^+ / S G = \operatorname{Ker}(\mathcal{T}_{\rho,\sigma} \to \mathbb{G}_m) \\
\mathcal{T}_{\rho,\mathring{\sigma}}^{+0} = M_{\rho,\mathring{\sigma}}^{+0} / S G = \operatorname{Ker}(\mathcal{T}_{\rho,\sigma} \to \mathbb{G}_m \times \mu_2 \times \mu_2) \\
\quad = \mathcal{T}_{\rho,\sigma}^+ \cap \mathcal{T}_{\rho,\mathring{\sigma}} \quad \text{dans } \mathcal{T}_{\rho,\sigma}
\end{cases}
$$

d'où des structures d'extensions

$$ 1 \to S G \to M_{\rho,\sigma} \to \mathcal{T}_{\rho,\sigma} \to 1 $$

etc (idem avec 
$M_{\rho,\mathring{\sigma}}, \mathcal{T}_{\rho,\mathring{\sigma}}$
$M_{\rho,\sigma}^+, \mathcal{T}_{\rho,\sigma}^+$
$M_{\rho,\mathring{\sigma}}^{+0}, \mathcal{T}_{\rho,\mathring{\sigma}}^{+0}$)

et des suites exactes

$$ 1 \to \mathcal{T}_{\rho,\sigma}^+ \to \mathcal{T}_{\rho,\sigma} \to \mathbb{G}_m \to 1 $$

etc (idem pour 
$\mathcal{T}_{\rho,\mathring{\sigma}} \subset \mathcal{T}_{\rho,\sigma}$
$\mathcal{T}_{\rho,\mathring{\sigma}}^{+0} \subset \mathcal{T}_{\rho,\mathring{\sigma}}$
$\mathcal{T}_{\rho,\sigma}^{+0} \subset \mathcal{T}_{\rho,\sigma}^+$
$\mathcal{T}_{\rho,\mathring{\sigma}}^{+0} \subset \mathcal{T}_{\rho,\sigma}$)

\newpage

%% ===== Page 453 =====

529

enfin des iso
$$ (151) \begin{cases}
R_{\rho,\sigma} / \mu_2 \simeq \tau_{\rho,\sigma} \\
R_{\rho,\sigma}^{\circ} / \mu_2 \simeq \tau_{\rho,\sigma}^{\circ} \\
S R_{\rho,\sigma} \simeq \tau_{\rho,\sigma}^{+} \\
S R_{\rho,\sigma}^{\circ} \simeq \tau_{\rho,\sigma}^{+\circ}
\end{cases} $$

80
Si $\underline{U}$ est un élément de $\mathcal{M}_{\rho,\sigma}$, on lui
associe ses images par $\theta_{\underline{a}}, \theta_{\rho}, \theta_{\sigma}$, et par
$\varepsilon, \varepsilon_1, \varepsilon_2$, on les notera respectivement

[MARGIN: Il faillait expliciter que les conditions (152)' conduisent $\mathcal{M}_{\rho,\sigma}$ comme un p. fibré principal sur $G \times \Gamma_{\underline{a}} \times \mathcal{N}_{\sigma}$, ... il faudrait le vérifier... mais peut on le faire ici ??? (153)' ne suffit pas vrai-]

$$ (152) \begin{cases}
u \in \Gamma_{\underline{a}} \\
\underline{a} \in \Gamma_{\rho} \\
\underline{b} \in \Gamma_{\sigma} \\
(\varepsilon, \varepsilon' \in \Gamma_{ph}) \\
\mu \in \Gamma_{\underline{a}}
\end{cases} \quad \text{satisfaisant} \quad (152)' \begin{cases}
\mu = \delta_{\underline{a}}(u) = \delta_{\rho}(\underline{a}) = \varepsilon \delta_{\sigma}(\underline{b}) \\
\varepsilon' = \varepsilon_{\sigma}(\underline{b})
\end{cases} $$
donc ces quantités sont
déterminées par la seule
conn. de $u, \underline{a}, \underline{b}$

au lieu des notations des § 45 et suivants, on reprendra, pour l'optique de la présente section, des éléments notés

$$ u, \alpha, \beta, \alpha_0, \beta_0, f \quad \text{avec} \quad \begin{cases} u \in \Gamma_{\underline{a}} \subset \Gamma_G \\ \alpha, \alpha_0 \in \Gamma G^{\pm 1} = \Gamma_G \\ \beta, \beta_0, f \in \Gamma S_G \end{cases} $$

[MARGIN: NB des cond (152)' ou (153) on déduit des identités autres... On pourrait introduire $g = \alpha \rho(\alpha_0)^{-1} (= a \rho(a)^{-1})$, $h = \beta \sigma(\beta_0)^{-1}$ mais je ne vois pas à quoi ça sert]

(153) reliés par des relations
$$ \begin{cases}
\alpha = u \underline{a} \\
\beta = u \underline{b}
\end{cases} \text{ i.e. } u = \alpha \underline{a}^{-} = \beta \underline{b}^{-} \qquad \text{[MARGIN: formules dans } G \text{]} $$
$$ \underline{U} = f^{-1} u $$
$$ \alpha_0 = f^{-1} \alpha \text{ i.e. } \alpha = f \alpha_0 $$
$$ \beta_0 = f^{-1} \beta \text{ i.e. } \beta = f \beta_0 $$

\newpage

%% ===== Page 454 =====

formules qui ont un sens seulement si
on se donne des plongements (ou
au moins morphismes)
$$
\begin{cases}
\Gamma_\rho \to G & \text{et un sous-groupe} \\
\Gamma_\sigma \to G & \Gamma_0 \hookrightarrow G
\end{cases} \leqno{(155)}
$$

Quand on a de tels plongements, les formules
(154) ont un sens, et définissent
(pour les quantités (152) connues) les
quantités (153) en termes de l'une d'elles,
disons $\underline{u}$, ou $f$, ou $\alpha$, ou $\beta$ (mais
pas en termes de $\alpha_0$ ou de $\beta_0$).

En fait, on dispose d'un sous-groupe
$$
\Gamma_0 \subset G \leqno{(156)}
$$
tel que
$$
\begin{cases}
\Gamma_0 \xrightarrow{\delta \text{ et } \dots} G \\
\text{i.e.} \quad G \simeq \Gamma_0 \cdot SG & \text{produit } 1/2 \text{ direct} \\
\text{pour } \underline{u} \in \Gamma G ,
\end{cases}
$$
de sorte que (la quantité existe un
unique $f \in \Gamma SG$ tel que
$$
\underline{U} = f^{-1} \underline{u}
$$
d'où des uniques $\alpha, \beta, \alpha_0, \beta_0 \in \Gamma G$ satisfai-
sant (avec $f$, $u$ précédents) à (154), i.e.
$$
\alpha = \underline{u} \underline{a}^{-1} , \quad \beta = \underline{u} \underline{b}^{-1} , \quad \alpha_0 = f^{-1} \alpha = f^{-1} \underline{u} \underline{a}^{-1} , \quad \beta_0 = f^{-1} \beta = f^{-1} \underline{u} \underline{b}^{-1} \leqno{(157)}
$$

\newpage

%% ===== Page 455 =====

On aura d'ailleurs
$$ \det \underline{u} = \det(\underline{f}) = \mu $$
$$ \det \alpha = \det(\underline{u}) \det(\underline{a})^{-1} = \varepsilon' \marginpar{par (152')} $$
$$ \det \beta = \det(\underline{u}) \det(\underline{b})^{-1} = 1 $$
et $\det \alpha_0 = \det \alpha = \varepsilon'$

$\det \beta_0 = \det \beta = 1$

On a donc bien $\alpha, \alpha_0 \in \Gamma G^{\pm 1} \defeq \{ g \mid (\det g)^2 = 1 \}$

$\beta, \beta_0 \in \Gamma S G$

- admis, bien sûr, que les
homomorphismes $\delta_{\hat{\mathfrak{T}}}, \delta_0$ de (155) sont déduits, via les
plongements (155), par le déterminant $\delta_G$.

Pour parvenir à la relation
$$ [\alpha, \rho] = \varepsilon [\beta, \sigma] \varepsilon_1 f(\mu-1) \leqno{(157')} $$
il faut se donner au préalable des sections $\rho, \sigma$ de $G$
(plus précisément des éléments de $\hat{\mathfrak{T}}_g$ et $\mathfrak{T}_0$
eux-mêmes), et un représentant strict $\varepsilon_1 \in \Gamma S G$,
pour que la relation (157') prenne un sens.

Remplaçant $\alpha, \beta$ par les valeurs $\underline{u} \underline{a}^{-1}, \underline{u} \underline{b}^{-1}$,
la relation devient (utilisant la commutativité
de $\hat{\mathfrak{T}}_g$, qui n'avait pas même
fait irruption jusqu'ici)
$$ [\underline{u}, \rho] = \varepsilon' \underline{u} (\underline{b}^{-1} \sigma \underline{b})^{-1} \sigma(\underline{u})^{-1} \varepsilon_1 f(\mu-1) \leqno{(*)} $$

\newpage

%% ===== Page 456 =====

et résultera de l'ensemble des deux conditions
$$
\begin{cases}
a) & \underline{b}^{-1} \sigma(b) = \varepsilon_1(b) \text{ identiquement } b \in \Gamma \mathcal{N}_0 \\
& \text{i.e. } \underline{b}(\sigma) = \varepsilon_1(b) \sigma \text{ identiquement} \\
b) & [\underline{u}, \rho] = [\underline{u}, \sigma] \varepsilon_1 (\mu - 1) \text{ identiquement en } \underline{u} \in \Gamma_{T_0} \\
& \text{pour } \mu = \delta_{T_0}(u) = \det(u) \\
& \text{i.e. } \rho(u) = \varepsilon_1(1-\mu) \sigma(u) \text{ identiquement en } u
\end{cases}
\leqno (158)
$$
ce qui s'écrit aussi (en appliquant $\sigma^{-1}$ aux deux membres)
$$
\sigma^{-1} \rho(u) = \varepsilon_0(1-\mu) \cdot u \quad \text{où} \quad \varepsilon_0 = \sigma^{-1}(\varepsilon_1)
$$
et n'est guère commode que si $T_0$ normalise $L_0$, et si $\sigma \in \Gamma_{T_0}$. Là on a de la forme
$$
\begin{cases}
a) & \sigma^{-1}\rho = \varepsilon_0 \underline{u}_0 \quad u_0 \in \Gamma_{T_0} \\
b) & T_0 \text{ opère sur } L_0 \text{ par la condition induite} \\
& \text{par le déterminant}
\end{cases}
\leqno (159)
$$

[unclear: Il y a quand même encore de la marge, par rapport aux conditions du commencement ainsi qu'à des VII ?)]

Mais ici l'équation 157) venait un peu comme un deus ex machina --- il doit y avoir des façons plus convaincantes et expéditives de s'en assurer, donnant par elles-mêmes des formules

\marginpar{attention \\ au signe}
$$
\begin{cases}
\underline{u}(\rho) = \operatorname{int}_{\alpha}(\rho) \quad ( = \alpha \rho \alpha^{-1}) & \text{i.e. } [\underline{u}, \rho] = \alpha \rho(\alpha)^{-1} \\
\underline{v}(\sigma) = \operatorname{int}_{\varepsilon\beta}(\sigma) \quad ( = \varepsilon\beta \sigma (\varepsilon\beta)^{-1}) & [\underline{v}, \sigma] = \varepsilon\beta \sigma(\beta)^{-1} \\
\underline{u}(\varepsilon_0) = \operatorname{int}_{\alpha}(\varepsilon_0) \text{ ($=1$ !)} &
\end{cases}
\leqno{(160)}
$$

La troisième relation résulte (159 b) du fait que par hypothèse $\underline{u} \in \Gamma_{T_0}$ opère sur $L_0$ (identiquement) par la multiplication par $\det(u)$.

\newpage

%% ===== Page 457 =====

Les deux premières formules s'explicitent successive-
ment, en écrivant $\alpha = \underline{u} a^{-1}$, $\beta = \underline{u} b^{-1}$, $y = f^{-1} \underline{u}$
$$ \alpha_0 = f^{-1} \alpha, \quad \beta_0 = f^{-1} \beta $$

$$
\begin{cases}
\underline{u}(\rho) = \operatorname{int}(\alpha)(\rho) \\
\underline{u}(\sigma) = \varepsilon \operatorname{int}(\beta)(\sigma)
\end{cases}
\quad \text{où } f \text{ ne figure plus}
\leqno{(161)}
$$
i.e.
$$
\begin{aligned}
\underline{u}(\rho) &= \operatorname{int}(\underline{u})(\rho) \quad \text{puisque } \operatorname{int}(a^{-1})\rho = \rho \\
&\text{ok!} \qquad \text{(comme } \Gamma_{\underline{\rho}} \text{ commut.)} \\
\underline{u}(\sigma) &= \varepsilon \operatorname{int}(\underline{u})(\underbrace{b^{-1}(\sigma)}_{= \varepsilon_0 \sigma}) = \operatorname{int}(\underline{u})\sigma \quad \text{OK!} \\
&\qquad \text{cf (158)(a)}
\end{aligned}
$$

Pour prouver (160) on n'a utilisé que (158 a)
et (159) (b). Essayons d'en déduire mainte-
nant la "relation des lacets" (157) - ça doit être
un calcul formel, que j'ai déjà dû faire une
douzaine de fois, dans des contextes qui ne
différaient entre eux que d'un chouia ! On s'y ramène !
en écrivant les relations (160) sous la forme
équivalente (161) et
$$
\text{(161 bis)} \quad \underline{u}(\varepsilon_0) = \varepsilon_0(\mu) \quad \text{i.e.} \quad \underbrace{[\underline{u},\varepsilon_0]}_{\underline{u} \varepsilon_0 \underline{u}^{-1}} = \varepsilon_0(\mu^{-1}) \quad \text{soit}
\leqno{(\ast)}
$$

$$
\begin{array}{ll}
\text{ou} \quad [\underline{u},\varepsilon_0] = \underbrace{\underline{u} \varepsilon_0 (\underline{u})^{-1}}_{= \sigma^{-1} \rho(\underline{u}) \sigma}
&
\text{or } \begin{cases} \rho(\underline{u}) = \rho \underline{u} \rho^{-1} = \xi \underline{u} \\ \sigma \underline{u} \sigma^{-1} = \eta \underline{u} \end{cases} \text{ d'où } \begin{cases} \xi = \alpha \rho(\alpha)^{-1} \\ \eta = \varepsilon \beta \sigma(\beta)^{-1} \end{cases} \\
& \\
\text{i.e. } \sigma([\underline{u},\varepsilon_0]) = \sigma(\underline{u}) \rho(\underline{u})^{-1}
&
\begin{cases} \rho(\underline{u}^{-1}) = \underline{u}^{-1} \xi^{-1} \\ \rho(\underline{u}) = \xi \underline{u} \\ \sigma(\underline{u}) = \eta \underline{u} \end{cases}
\end{array}
$$

d'où $\quad \sigma([\underline{u},\varepsilon_0]) = \xi^{-1} \eta$

et $(\ast)$ devient, vu que $\sigma(\varepsilon_0) = \varepsilon_1$
$$
\xi^{-1} \eta = \varepsilon_1(1-\mu) \quad \text{i.e.} \quad \xi = \eta \varepsilon_1(\mu-1) \quad \text{i.e. (157)}
$$
OK !

\newpage

%% ===== Page 458 =====

En résumé:

Proposition Supposons donnés, en plus des données
de la section précédente $7^\circ$ (cf.\ p.\ 526' H 3, notamment
(135) (136) ), ceci :
$$
\begin{cases}
\text{a) sections } \rho \in \Gamma_{T_S}, \sigma \in \Gamma_{T_{S'}} \text{ sections \underline{centrales}} \\
\text{b) plongements } \Gamma_{T_S} \hookrightarrow G, \Gamma_{T_{S'}} \hookrightarrow G \\
\text{c) tore } T_0 \subset G \text{ tel que } \delta_G : T_0 \isommap G_m \\
\quad \text{i.e. } G \simeq T_0 \cdot SG \text{ (semi-direct)} \\
\text{d) } \varepsilon_0 \in \Gamma_G \text{ unipotent strict. régulier, définissant } \\
\quad \lambda \mapsto \varepsilon_0(\lambda) : G_m \hookrightarrow SG
\end{cases} \leqno (162)
$$

et supposons :
\marginpar{relie $L_0, T_0$ \\ relie $L_0, T_0 / d)$ \\ $\rho, \sigma$}
$$
\begin{cases}
\text{a) } \delta_G | \Gamma_{T_S} = \delta_S, \quad \delta_G | \Gamma_{T_{S'}} = \delta_{S'} \\
\text{b) } \operatorname{int}(k)(\sigma) = \varepsilon_0(\lambda) \sigma \text{ dans } G \text{ [unclear: s'identifiant] } \\
\quad \text{pour } k \in \Gamma_{T_{S'}} \\
\text{c) } T_0 \text{ normalise } L_0, \text{ et on a } : \\
\quad u(\varepsilon_0(\lambda)) = \varepsilon_0(\mu \lambda) \quad \text{où } \mu = \delta_G(u) \\
\text{d) } \sigma \rho^{-1} \in L_0 \varepsilon_0 L_0 \text{ i.e. } \exists \text{ loc.t sect.n } \lambda \in \Gamma_{G_m} \\
\quad \text{telle que } \lambda_0 \sigma \rho^{-1} = \varepsilon_0 u_0, \text{ avec } u_0 \in \Gamma_{T_0}
\end{cases} \leqno (163)
$$

Sous ces conditions, on a pour le système
(152) satisfaisant les relations (152') (i.e.\ en
fait pour le système de $\mathcal{M}_{g, \nu}$ !) un
unique système (153), satisfaisant les relations
(154) . Soit $U$ mis de façon unique
sous la forme
$$
U = f^{-1} u, \quad f \in \Gamma_{SG}, u \in \Gamma_{T_0} \leqno (164)
$$
et on pose
$$
\alpha = u a^{-1}, \beta = u b^{-1}, \alpha_0 = f^{-1} \alpha, \beta_0 = f^{-1} \beta = f^{-1} u b^{-1} . \leqno (165)
$$

\newpage

%% ===== Page 459 =====

Tout ceci posé, on a alors les formules (160), (161), (161 bis), et la formule "des bouts" (157).

Nous nous posons maintenant la question sous quelles conditions la connaissance des $(\alpha, \beta, f)$ (p. ex.) parmi les quantités (153), permet de récupérer les quantités (153) - i.e.\ une section de $\mathcal{M}_{g,r}$ - et quelles conditions il faut imposer à $(\alpha, \beta, f)$. Il suffirait de (pour l'unicité) pouvoir construire $\underline{u}$, d'où $\underline{U}, \underline{a}, \underline{b}$ par les trois premières formules (154). Il [unclear: s'agira] de trouver $\underline{a} \in \Gamma T_\rho$ tel que $\alpha \underline{a}$ soit section de $T_0$, - indétermination sur $\underline{a}$ est dans $\Gamma T_0 \cap T_\rho$, et $\underline{b} \in \Gamma N'_0$ tel que $\beta \underline{b}$ soit section de $T_0$ - indétermination sur $\underline{b}$ [unclear: est dans] $\Gamma T_0 \cap N'_0$.

L'indétermination sur $\underline{u}$ est la multiplication à droite par une section de $T_0 \cap T_\rho \cap N'_0$, il suffira donc qu'on ait\footnote{En fait, on a $T_0 \cap T_\rho = \{ 1 \}$ dans les cas envisagés dans II, III et suite...}
$$
T_0 \cap T_\rho \cap N'_0 = \{ 1 \} \leqno{(164)}
$$
Supposant cette condition satisfaite, quelles conditions sur les [unclear: seuls] $(\alpha, \beta, f)$ pour provenir d'un syst.\ (152) - i.e.\ d'une section de $\mathcal{M}_{g,r}$ ?

\newpage

%% ===== Page 460 =====

On voit que si une section $\varphi_0$ de $\mathcal{M}_{g,r}^\sim$, de [unclear: composantes] correspondant à $\alpha_0, \beta_0, f_0$, est modifiée en $\varphi_1 = f^{-1} \varphi_0$, correspondant à $\alpha_1, \beta_1, f_1$, alors
$$
\alpha_1 = \alpha_0, \beta_1 = \beta_0, f_1 = f_0 f \leqno{(165')}
$$
On voit que si $\varphi_0$ correspond à $(U_0, \underline{a}_0, \underline{b}_0)$, et $\varphi_1$ à $(U_1, \underline{a}_1, \underline{b}_1)$, alors on a
$$
U_1 = f^{-1} U_0, \underline{a}_1 = \underline{a}_0, \underline{b}_1 = \underline{b}_0 \leqno{(166)}
$$
(car $\theta_{\mathcal{G}} | S\mathcal{G} = \text{l'inclusion canonique } S\mathcal{G} \subset \mathcal{G}$, et $\theta_{\mathcal{G}}$ [unclear: est triviale] sur $S\mathcal{G}$).

D'où $\underline{u}_1 = \underline{u}_0$ (car $\det U_1 = \det U_0$)

d'où (165'). Cela montre que la condition cherchée sur $(\alpha, \beta, f)$ ne dépend [unclear: que des] $\alpha, \beta$, pas de $f$.

On voudrait qu'elle s'exprime par la validité de la formule (157)\footnote{A vrai dire, cette condition = des relations (157) avec $f=1$ (et bien entendu $u \in \Gamma_0$) est en fait toujours vérifiée (qui [unclear: fait exiger] que dans ce cas $B=0$, ou $u=1$ \dots ce qu'on a supposé [unclear: plus haut]!). L'intérêt de [unclear: substituer] la validité de (157) pour le couple $(\alpha, \beta)$ qu'on obtient en "multipliant" l'action sur $(a_i)$ par $f$. (On remarquera que dans ce cas $U=1, \dots$).} pour un $\varepsilon$ précis ($\varepsilon = 1$) (qui sera donc [unclear: d'ailleurs] une identité - et par suite [unclear: induira] une relation locale\dots)

Je me contenterai de cette condition pour l'objet des [unclear: n\textsuperscript{os}] de (II) (III), des conditions (plus draconiennes que celles [unclear: de la] prop. précédente (p. 531') qui nous [unclear: sont] imposées au triplet $(\varepsilon_0, f, \sigma)$. Il

\newpage

%% ===== Page 461 =====

est possible que les conditions auxquelles nous
nous sommes [unclear: astreints] étaient plus restrictives
que nécessaires, et que je sois obligé
d'y revenir par la suite, dans un contexte
plus général que celui des [unclear: formes] et
[unclear: arbres], $\operatorname{Gl}(1), \operatorname{Sl}(2)$, mais pour le moment
l'étude détaillée faite en nous inspirant
des hypothèses dégagées dans II, III nous
servira de déblayage, pour nous éclaircir
les idées sur un cas particulier
important.

Remarque. Rappelons que nous avons vu que, en
fait $p$ est déjà défini en termes de $\alpha$ seul,
il en est de même de $\alpha p(\alpha)^{-1} = [unclear: \rho]$
Mais il n'est nullement déterminé en termes de $p$
— en fait, le $\beta$ n'apparaît que modulo
le [unclear: commutant] de $\varepsilon$ (qui d'ailleurs, modulo
l'identité \dots). Mais si on impose (par exemple) $\varepsilon = 1$,
la donnée de $\alpha$ ne détermine $\beta$ que
modulo $\beta \mapsto \beta [unclear: \gamma_\alpha]$ \dots

\newpage

%% ===== Page 462 =====

La structure de $\mathcal{M}_{p,\sigma}$ étant bien comprise (cf p. 495!), revenons sur le groupe $S\mathcal{M}_{p,\sigma}$, considéré comme produit semi-direct de $\mathcal{M}_{p,\sigma}$ par le sous-groupe $L_0$ de $SG$, sur lequel on fait opérer $\mathcal{M}_{p,\sigma}$ via $\chi$. Désignant par
$$ 
\mathcal{E}_{L_0} = \operatorname{Ker}(S\mathcal{M}_{p,\sigma} \to \mathcal{M}_{p,\sigma}) 
\leqno{(165)} 
$$
(correspondant à $\varepsilon_0$, par déf. et construction)
$$ 
u(\mathcal{E}_{L_0}(\lambda)) = \mathcal{E}_{L_0}(\mu \lambda) \quad \text{où } \mu = \chi(u). 
\leqno{(166)} 
$$

\newpage

%% ===== Page 463 =====

\marginpar{533}
g) Les sous-groupes $M_{p,\sigma}(\alpha)$, $M_{p,\sigma}[J]$, et le groupe $SM_{p,\sigma}$, l'isom.\ $SM_{p,\sigma} \isommap \Sigma_{p,\sigma}^+$.

Le groupe $M_{p,\sigma}$ est extension de $\Gamma_{p,\sigma}$ par $\mathfrak{S}_{p,\sigma}^+ \simeq SG$ et les trois sous-groupes qu'on a vus plus haut s'explicitent à propos des sections, partielles --- [unclear: le fait est que] pour en trouver, il suffira d'une section partielle seulement, au-dessus du sous-groupe (d'indice $2$ de $\Gamma_{p,\sigma}$) $\Gamma_{p,\sigma}^+$, image de $E_{\sigma}$.

\marginpar{finalement je ne reviens pas à \dots}Déterminons de façon générale les sections de $M_{p,\sigma}$ sur $\Gamma_{p,\sigma}$.
$$
Q_{p,\sigma} \xrightarrow{u} \Gamma_{p,\sigma} \leqno{(1)}
$$
(épi.\ de noyau [unclear: profini])

[unclear: Si \dots]

Il revient au même de se donner une section, ou un hom.
$$
Q_{p,\sigma} \xrightarrow{\Psi} M_{p,\sigma} \leqno{(2)}
$$
qui soit
$$
\begin{cases}
a) \text{ trivial sur le noyau \quad ("remonte" (1))} \\
b) \text{ [unclear: se factorise] \dots (identique \dots)}
\end{cases}
$$

Les hom.\ $\Psi$ qui satisfont à b) correspondent dans le scindage de l'extension triviale de $Q_{p,\sigma}$ par $SG$ (extension produit), aux homomorphismes
$$
Q_{p,\sigma} \xrightarrow{\Psi^0} SG \leqno{(4)}
$$

\newpage

%% ===== Page 464 =====

en [unclear: posant] : $\Psi^\circ$ le hom.
$$ x \mapsto \Psi(x) = x \Psi^\circ(x) \leqno (5) $$
Donc la condition a prend la forme
$$ \Psi^\circ(x) = x \quad \text{sur} \quad pr_\nu = R_{g,\nu} \cap SG . \leqno (6) $$
Les $\Psi, \Psi^\circ$ correspondants aux sections $\mathcal{M}_{g,\nu}(0), \mathcal{M}_{g,\nu}(1/2), \mathcal{M}_{g,\nu}(-J)$ seront notés resp.\ $(\Psi_0, \Psi_0^\circ), (\Psi_{1/2}, \Psi_{1/2}^\circ), (\Psi_{-J}, \Psi_{-J}^\circ)$.

1$^\circ$) Cas de $\mathcal{M}_{g,\nu}(0)$. Avec les notations de la p.\ 529, il correspondra [unclear: à des] $\mathcal{M}_{g,\nu}$ satisfaisant à : l'une des conditions équivalentes $Y=y$,
$$ f=1, \quad (U \in \Gamma_{T_0}, \alpha_0=\alpha, \beta_0=\beta) \leqno (7) $$
Donc on a
$$ \mathcal{M}_{g,\nu}(0) = \Theta^{-1}(T_0) \leqno (8) $$
(en termes de la description de $\mathcal{M}_{g,\nu}$ de la p.\ 526).

Vérifions que c'est un sous-groupe-section :
On a $\mathcal{M}_{g,\nu}(0) \cap \mathfrak{S}_{g,\nu}^+ = \{1\}$ car $\Theta | (\mathfrak{S}_{g,\nu}^+ = SG)$ est l'inclusion canonique $SG \hookrightarrow \mathcal{G}$, $\Rightarrow SG \cap T_0 = \{1\}$.
Reste à vérifier que $\mathcal{M}_{g,\nu}(0) \to \mathfrak{T}_{g,\nu}$ (qui est mono par (a)) est épi. Mais en

\newpage

%% ===== Page 465 =====

termes de la [unclear: relation] can.
$$
\mathcal{M}_{p,\sigma} \simeq \underbrace{R_{p,\sigma} \times S\mathcal{G}}_{\widetilde{\mathcal{M}}_{p,\sigma}} / \mu_2 \text{ diag.} \leqno{(9)}
$$

$$
i_{T_0} : \mathbb{G}_m \isommap T_0 \subset \mathcal{G} \footnote{$k : R_{p,\sigma} \to \mathbb{G}_m$ l'hom. canonique provenant de l'hom. $\Theta \mid R_{p,\sigma}$} \leqno{(10)}
$$
iso inverse de l'isom.
$$
\delta_{T_0} = \det \mid T_0 : T_0 \isommap \mathbb{G}_m ,
$$
on a, pour un élément $(\xi, U_0)$ de $\widetilde{\mathcal{M}}_{p,\sigma}$
$$
\Theta(\xi, U_0) = k(\xi) U_0 , \text{ déterminant } k(\xi)^2
$$
dire que c'est $\in T_0$ signifie donc que
$$
k(\xi) U_0 = i_{T_0}(k(\xi)^2) \quad \text{i.e.}
$$
$$
U_0 = k(\xi)^{-1} i_{T_0}(k(\xi)^2) \leqno{(11)}
$$
donc l'image inverse $\mathcal{M}_{p,\sigma}^{(0)}$ de $T_0$
dans $\widetilde{\mathcal{M}}_{p,\sigma}$ est formée des
$$
(\xi, k(\xi)^{-1} i_{T_0}(k(\xi)^2)) \leqno{(12)}
$$
en d'autres termes cela correspond à un
sous-groupe section, avec
$$
\psi_0^\circ(\xi) = k(\xi)^{-1} i_{T_0}(k(\xi)^2) \leqno{(13)}
$$
Comme il se doit, c'est bien un hom.
$$
\psi_0^\circ : R_{p,\sigma} \to S\mathcal{G}
$$
et sa valeur sur $\mu_2$ est [unclear: l'identité]
car $k(\xi)^2 = 1$ et $\psi_0^\circ(\xi) = k(\xi)^{-1} = \xi$ .

\underline{Rem} Pour définir cette section, il suffisait

\newpage

%% ===== Page 466 =====

[unclear: l'on prenne un au pur hasard]
ou [unclear: hom.] quelconque, noté $p_{T_0}$, de $\Gamma$ dans
$SG$, et de poser (13).

Avant de quitter $\mathcal{M}_{\rho,\sigma}(0)$, je vais
introduire un groupe un peu plus grand
(de dim 4, au lieu de 3)
$$
\mathcal{M}_{\rho,\sigma}[0] \defeq \varpi^{-1}(B'_0) \leqno{(14)}
$$
où, je le rappelle, $B'_0$ est défini comme le
produit $\frac{1}{2}$ direct
$$
B'_0 = T_0 . L_0 \leqno{(15)}
$$
Cette fois-ci, $U = U_0 k(\Gamma) \in \varpi^{-1}(B'_0)$ n'est pas déterminé
par la seule connaissance de son déterminant
$k(\Gamma)^2$, qui est le déterminant modulo
$L_0$. Cela correspond au fait que
$\mathcal{M}_{\rho,\sigma}[0] \cap \widehat{SG} = B'_0 \cap \widehat{SG} = L_0$ (qui n'est pas réduit à $e$ ! (1)),
mais : $L_0$ apparaît donc comme
un sous-groupe invariant de $\mathcal{M}_{\rho,\sigma}[0]$, et
ce manifestement. Compte tenu que
$\mathcal{M}_{\rho,\sigma}(0) \subset \mathcal{M}_{\rho,\sigma}[0]$ et $\mathcal{M}_{\rho,\sigma}(0) \isommap \tilde{\Gamma}_{\rho,\sigma}$
puis
$$
\mathcal{M}_{\rho,\sigma}[0] = \mathcal{M}_{\rho,\sigma}(0) . L_0 \leqno{(16)}
$$
produit $\frac{1}{2}$ direct.

On a encore une constatation qui a
son sens dès que l'on a un $T_0$ comme dessus,
soit un $T_0$ et un sous-groupe $L_0$ de $SG$,

\newpage

%% ===== Page 467 =====

invarié par $T_0$.

L'opération de $\mathcal{M}_{g,\nu}[0]$ sur $SG$ se fait par l'intermédiaire de $\theta$, [unclear: définie par ailleurs] comme un ép.
$$
\mathcal{M}_{g,\nu}[0] \overset{\theta | \mathcal{M}_{g,\nu}[0]}{\longrightarrow} B'_0 \ ;
$$

\marginpar{[unclear: fait pb] pour une déf.\ plus théo.\ cf p 538, 17)}
Cette opération étant précisée, on peut introduire
$$
S\mathcal{M}_{g,\nu} \defeq \mathcal{M}_{g,\nu}[0] \cdot SG \leqno{(17)}
$$
(produit 1/2 direct) dans lequel $SG$ est invariant, et $\mathcal{M}_{g,\nu}[0]$ s'identifie à un sous-groupe de $S\mathcal{M}_{g,\nu}$. Comme $\mathcal{M}_{g,\nu}(0) \subset \mathcal{M}_{g,\nu}[0]$, on a une inclusion can.
$$
\underset{= \mathcal{S}\mathcal{M}_{g,\nu}}{\mathcal{M}_{g,\nu}(0) \cdot SG} \subset \underset{= S\mathcal{M}_{g,\nu}}{\mathcal{M}_{g,\nu}[0] \cdot SG} \leqno{(18)}
$$

Pour ne pas commencer à m'embrouiller i. e.\ m'y perdre, je vais m'accrocher - pour ma compréhension - à la signification ``géom.'' des $S\mathcal{M}_{g,\nu}$, $\mathcal{M}_{g,\nu}[0]$, aux notations de la p.\ 529, qui mérite d'être explicitée par un

\newpage

%% ===== Page 468 =====

Proposition Considérons l'homomorphisme
$$ \mathcal{M}_{p,\sigma} \longrightarrow C \times T_p \times N'_{\sigma} \times \dots \leqno{(19)} $$
$$ x \longmapsto (\Theta_a(x), \Theta_p(x), \Theta_\sigma(x), \mu(x), \varepsilon(x), \varepsilon'(x)) $$
C'est un monomorphisme, dont l'image est formée des multiplets $(U, \underline{a}, \underline{b}, \mu, \varepsilon, \varepsilon')$ satisfaisant
$$ \begin{cases} \mu = \delta_a(U) = \delta_\sigma(\underline{b}) = \varepsilon' \delta_p(\underline{a}) \\ \varepsilon = \varepsilon_\sigma(\underline{b}) \end{cases} \leqno{(20)} $$

Corollaire Le morphisme
$$ \mathcal{M}_{p,\sigma} \longrightarrow C \times T_p \times N'_{\sigma} \leqno{(21)} $$
$$ x \longmapsto (\Theta_a(x), \Theta_p(x), \Theta_\sigma(x)) $$
est un monomorphisme, dont l'image est formée des triplets $(U, \underline{a}, \underline{b})$ satisfaisant\marginpar{NB on a $\begin{cases} \varepsilon = \varepsilon_\sigma(\underline{b}) \\ \varepsilon' = \delta_p(\underline{a})/\delta_\sigma(\underline{b}) \\ \mu = \delta_\sigma(\underline{b}) = \delta_a(U) \end{cases}$}
$$ \begin{cases} \delta_a(U) = \delta_\sigma(\underline{b}) \\ (\delta_p(\underline{a}) / \delta_\sigma(\underline{b}))^2 = 1 \end{cases} \leqno{(22)} $$

Démonstration de la proposition : On a 
$$ \operatorname{Ker}(\Theta_a, \mu, \varepsilon, \varepsilon') = S \mathcal{R}_{p,\sigma} \simeq S T_p \times S T_\sigma $$
et $\Theta_p, \Theta_\sigma$ donnés sur $S \mathcal{R}_{p,\sigma}$ s'identifient aux projections, donc
$$ \operatorname{Ker}( (\Theta_p, \Theta_\sigma) | S \mathcal{R}_{p,\sigma} ) = (1), \text{ donc } $$
$$ \operatorname{Ker}(\Theta_a, \mu, \varepsilon, \varepsilon', \Theta_p, \Theta_\sigma) = (1) $$

\newpage

%% ===== Page 469 =====

donc monomorphisme. Pour l'épimorphisme, on est ramené à la forme
équivalente du Corollaire. On peut supposer qu'on a un $\lambda$ tel
que\footnote{NB Où l'on a vu $\mathcal{M}_{\rho,\sigma} = R_{\rho,\sigma} \rtimes G$ que j'avais noté $\Theta_G$.}
$$
\lambda^2 = \mu \defeq \det \underline{U}
$$
on veut écrire
$$
\underset{\in \mathcal{M}_{\rho,\sigma}}{\underline{U}} = \text{[unclear: }\Theta\text{]} (\underset{\in R_\rho}{r}, \underset{\in S\mathcal{G}}{\underline{U}_0}) = k(r)\underline{U}_0 \quad k(r)^2 = \det \underline{U} \text{ i.e. } k(r)^2 = \mu,
$$
on prendra $r$ avec $k(r) = \lambda$,
et on pose donc
$$
\underline{U}_0 = \lambda^{-1}\underline{U}, \quad \underline{a}_0 = \lambda^{-1}\underline{a}, \quad \underline{b}_0 = \lambda^{-1}\underline{b}
$$
et on est ramené à voir que $\underline{U}_0, \underline{a}_0, \underline{b}_0$ est dans
l'image de $\mathcal{M}_{\rho,\sigma}$ (on vient de "corriger"
$(\underline{U}, \underline{a}, \underline{b})$ par l'image dans $G \times T_\rho \times T_\sigma$ de
la section $\lambda$ de $G \hookrightarrow R_{\rho,\sigma}$) i.e. ramené
à voir que $(\underline{a}_0, \underline{b}_0)$ est dans l'image de
$R_{\rho,\sigma} \to T_\rho \times \mathcal{N}'_\sigma$. Mais maintenant
$\underline{a}_0$ est dans $S^1 T_\rho$, $\underline{b}_0$ dans $S \mathcal{N}'_\sigma$, et justement
$R_{\rho,\sigma} \simeq S^1 T_\rho \times S \mathcal{N}'_\sigma$, ce qui gagne \dots

Remarque. On ne peut remplacer $G$ par $PG$,
i.e. $\mathcal{M}_{\rho,\sigma} \to PG \times T_\rho \times \mathcal{N}'_\sigma$ (ici on peut rajouter
encore $\epsilon_1, \epsilon_2, \epsilon_3$ si on veut!) n'est pas monomorphisme -
son noyau est justement le sous-groupe d'ordre
2, le centre de $\mathcal{G}^+_{\rho,\sigma} = S\mathcal{G}$.

\newpage

%% ===== Page 470 =====

Considérons maintenant que les quantités $\underline{U}, \underline{a}, \underline{b}$ (et $\mu, \varepsilon, \varepsilon'$) de p. 539 (152) sont comprises, les quantités (153) en loc. cit. $\underline{u}, f, \alpha, \beta, \alpha_0, \beta_0$ ($g, h, s, y$ n'y étaient, en plus..) Elles sont déterminées en termes des $\underline{U}, \underline{a}, \underline{b}$, dès qu'on connait soit $\underline{u}$, soit $f$, par les formules (154), $\underline{u}$ et $f$ étant reliées par la seule relation
$$ \underline{U} = f^{-1} \underline{u} \leqno{(23)} $$
avec en plus sur $f$ la condition
$$ f \in \Gamma S\underline{G} = \widehat{\mathfrak{S}}_{p,\sigma}^+ \leqno{(24)} $$
plus une condition sur $\underline{u}$ qui, dans loc. cit., s'écrivait
$$ \underline{u} \in \Gamma T_0 \leqno{(25)} $$
et était suggérée par les calculs de § 46 V (repris ensuite dans § 48 II, III etc.). Cette condition assure l'unicité de $f, \underline{u}$ satisfaisant à (23) -- c'est cela, [unclear: au plan des vertus] de calculs, son principal mérite ! -- mais dans certains contextes il est préférable de la remplacer par la condition plus large
$$ \underline{u} \in \Gamma B'_0 \qquad \therefore B'_0 = T_0 \cdot L_0 \text{ (produit semi-direct)} \leqno{(26)} $$

\newpage

%% ===== Page 471 =====

On n'a plus de résultat d'unicité pour l'écriture de $\underline{u}$ sous forme (23) -- par contre on a un théorème d'existence sous forme simple :

Si on part de $(\alpha, \beta, \rho) \in$ [unclear: certaines] de $G$ telles que
$$
(\det \alpha)^2 \det \rho = \det f = 1, \leqno{(27)}
$$
pour que [unclear: un tel triplet] soit par conséquent les quantités qui s'en déduisent :
$$
\begin{cases} \alpha_0 = f^{-1} \alpha, \beta_0 = f^{-1} \rho \beta \\ y = \alpha \rho (\alpha \rho)^{-1}, y' = \beta \sigma (\rho \beta)^{-1} \\ g = \alpha_0 \rho (\alpha_0 \rho)^{-1}, h = \beta_0 \sigma (\rho \beta_0)^{-1} \end{cases} \leqno{(28)}
$$
proviennent d'un système $(x, \underline{u})$, $x \in \Gamma \mathcal{M}_{p, \sigma}$, $\underline{u} \in \Gamma \mathcal{B}^1_0$ tel que
$$
f \defeq x \underline{u} \in \Gamma \mathfrak{S} \mathcal{G} \quad \text{i.e.} \quad \det \underline{u} = \det \mathcal{U}, \leqno{(29)}
$$
par les formules
$$
\alpha = \underline{u} \alpha_0, \quad \beta = \underline{u} \beta_0, \quad f = \underline{u} \mathcal{U}^{-1} \leqno{(30)}
$$
il faut et il suffit que l'on ait la relation (qui en fait par l'inversion de $f$, entraîne $\alpha_0, \beta_0$) :
$$
\alpha \rho (\alpha^{-1} \sigma) = \varepsilon \beta \sigma (\beta^{-1}) \varepsilon_1 (p^{-1}) \qquad \text{pour } \varepsilon \in \Gamma \mathcal{M}_{p, \sigma}, \mu \in \Gamma_{\mathcal{B}_0} \text{ convenables.} \leqno{(31)}
$$

Alors le couple $(x, \underline{u})$ est uniquement déterminé à [unclear: un facteur près]. Il correspond à 1-1 entre les $\rho$-classes de $\mathcal{M}_{p, \sigma} [0]$ et $\mathcal{M}_{p, \sigma}$.

\newpage

%% ===== Page 472 =====

des couples $\alpha, \beta$ satisfaisant aux deux relations $(\det \alpha)^2 = \det \beta = 1$ et (27), et la relation des lacets. Pour cela, on n'a [unclear: d'utilité] pour le groupe $S\mathcal{M}_{g,\sigma}$.

Mais [unclear: prenons] un cas $f$ quelconque, pour interpréter de façon algébrique simple les triplets $(\alpha, p, f)$ satisfaisant (27) (31), (ayant la structure de [unclear: groupoïde] via relations de cocycles, groupoïde étrange dessus (sur laquelle il me faudra revenir un peu plus longuement...) où d'ailleurs ces trois quantités $\alpha, p, f$ se mélangent intimement --- il y a lieu d'introduire $S\mathcal{M}_{g,\sigma}$ --- qui sera isomorphe à [unclear: chacun] (des groupes) de solutions communes de (27) (31).

De ce point de vue, $S\mathcal{M}_{g,\sigma}$ se définit bien comme le sous-groupe des triplets
$$ S\mathcal{M}_{g,\sigma} = \text{ ensemble des triplets } (x, \underline{u}, f) \dots \leqno{(32)} $$
$$ \left\{ \begin{array}{l} x \in \mathcal{M}_{g,\sigma}, \\ \underline{u} \in B'_0 \\ f \in S\mathcal{G} \end{array} \right\} \quad f^{\underline{u}} = \underset{\defeq \Theta(x)}{f^x} $$

D'où
$$ S\mathcal{M}_{g,\sigma} \hookrightarrow \mathcal{M}_{g,\sigma} \times B'_0 \leqno{(33)} $$
$$ \simeq \left\{ (x, \underline{u}) \mid f(x) = [unclear: \det(\underline{u})] \right\} $$

on [unclear: aura], plus simple [unclear: $\det u$]

\newpage

%% ===== Page 473 =====

$$
\begin{aligned}
& S\mathcal{M}_{g,\sigma} \isommap \mathcal{M}_{g,\sigma} \times_{G_m} B'_0 \\
& \text{où } \begin{cases}
\mathcal{M}_{g,\sigma} \xrightarrow{\chi} G_m \\
B'_0 \xrightarrow{\delta_{B'} = \overline{\varepsilon}_G | B'_0} G_m
\end{cases} \text{, ouf !}
\end{aligned} \leqno (34)
$$

Cette formule va nous servir de définition de $S\mathcal{M}_{g,\sigma}$ dorénavant - elle définit du même coup une loi de groupe sur $S\mathcal{M}_{g,\sigma}$, et on trouve sur $S\mathcal{M}_{g,\sigma}$ deux structures d'extensions

$$ 
1 \to \underset{\substack{\parallel \\ \operatorname{Ker}(B'_0 \xrightarrow{\delta_{B'}} G_m)}}{L_\bullet} \to S\mathcal{M}_{g,\sigma} \xrightarrow{\tilde{x}=(x,y) \mapsto x} \mathcal{M}_{g,\sigma} \to 1 \leqno (35) 
$$

$$ 
1 \to \underset{\substack{\parallel \\ \operatorname{Ker}(\mathcal{M}_{g,\sigma} \xrightarrow{\chi} G_m)}}{\mathcal{M}_{g,\sigma}^+} \to S\mathcal{M}_{g,\sigma} \xrightarrow{\underset{\substack{\parallel \\ (x,u)}}{\tilde{x}} \mapsto u} B'_0 \to 1 \leqno (36) 
$$

Le sous-groupe $S\mathcal{M}_{g,\sigma}[0]$ de $S\mathcal{M}_{g,\sigma}$ correspond à la condition $f=1$ i.e.\ $u = \underset{\substack{\parallel \\ \Theta_G(x)}}{\underline{u}}$ (qui implique déjà $\det u = \underset{\substack{\parallel \\ \chi(x)}}{\det \underline{u}}$), c'est-à-dire le sous-groupe des $(x,u)$ de $\mathcal{M}_{g,\sigma} \times B'_0$ tels que l'on ait $u = \Theta_G(x)$. Donc $S\mathcal{M}_{g,\sigma}[0]$ est le sous-groupe en $S\mathcal{M}_{g,\sigma}$ image de $\mathcal{M}_{g,\sigma}[0]$ (cf (14) p. 536) par $x \mapsto (x, \Theta_G(x))$ :

$$ 
S\mathcal{M}_{g,\sigma}[0] = \operatorname{Im} \big( \mathcal{M}_{g,\sigma}[0] \underset{x \mapsto (x, \Theta_G(x))}{\hookrightarrow} S\mathcal{M}_{g,\sigma} \big) \leqno (37) 
$$

$$ 
( \simeq \mathcal{M}_{g,\sigma}[0] = \Theta_G^{-1}(B'_0) ) 
$$

\newpage

%% ===== Page 474 =====

La donnée d'un scindage de (35) revient
à un scindage de
$$
1 \longrightarrow L_\circ \longrightarrow B'_\circ \longrightarrow G_m \longrightarrow 1 \quad ,
$$
or il y en a un tout prêt par le sous-groupe
$T_\circ$. Donc on en déduit un plongement de
$\mathcal{M}_{g,\nu}$ dans $S\mathcal{M}_{g,\nu}$, qui en fait
bien un produit semi-direct
$$
S\mathcal{M}_{g,\nu} \simeq \mathcal{M}_{g,\nu} \cdot L_\circ \leqno{(38)}
$$
$\mathcal{M}_{g,\nu}$ opérant d'ailleurs sur $L_\circ$ via le
caractère $\chi : \mathcal{M}_{g,\nu} \longrightarrow G_m$ --- mais
cette description me semble
beaucoup moins \og parlante \fg{} que (34),
car elle ne met pas en évidence l'homomorphisme can.\ 
$$
S\mathcal{M}_{g,\nu} \longrightarrow B'_\circ
$$
qui paraîtra comme la chose essentielle
pour l'introduction de $S\mathcal{M}_{g,\nu}$. Une fois
ceci posé, on trouve (par [unclear: inversion] ?) le

\underline{Théorème} --- un isom.\ can.\ 
$$
S\mathcal{M}_{g,\nu} \isommap \mathfrak{S}^*_{g,\nu}\footnote{sous-groupe des systèmes (cf. \S f) satisfaisant (17) et (31)} \leqno{(39)}
$$

\newpage

%% ===== Page 475 =====

donné par
$$
(x, y) \longmapsto (\alpha, \beta, f) \quad \text{où} \quad \alpha = \underline{u}\underline{a}^{-1}, \; \beta = \underline{u}\underline{b}^{-1}, \; f = \underline{u}\underline{U}^{-1} \leqno{(40)}
$$
(comme c'est bizarre !)
$$
\text{où } \underline{a} = \theta_p(x), \; \underline{b} = \theta_\sigma(x), \; \underline{U} = \theta_G(x)
$$

L'opération de $L_0$ dans $\Sigma_{p,\sigma}$ par
$$
(\alpha, \beta, f) \longrightarrow (l\alpha, l\beta, lf) \leqno{(41)}
$$
correspond à l'opération de translation à gauche du sous-groupe $L_0$ de $S\mathcal{M}_{p,\sigma}$ (cf (35)), pour $l \in L_0$. Donc on récupère $\mathcal{M}_{p,\sigma} \isom S\mathcal{M}_{p,\sigma} / L_0$ comme le schéma (quotient) de $\Sigma_{p,\sigma}^*$ par l'opération (41) de $L_0$.

Dans cette description des $(\alpha, \beta, f)$ interviennent les quantités indésirables $\underline{u}, \underline{a}, \underline{b}, \underline{U}$, ce qui conduit à considérer l'homomorphisme
$$
S\mathcal{M}_{p,\sigma} \longrightarrow B'_0 \times \Gamma'_p \times \Gamma'_\sigma \times G \leqno{(42)}
$$
$$
(x, u) \longmapsto (\underline{u}, \underline{a}, \underline{b}, \underline{U}) \quad \text{avec} \quad \begin{cases} \underline{a} = \theta_p(x) \\ \underline{b} = \theta_\sigma(x) \\ \underline{U} = \theta_G(x) \end{cases}
$$

On a à ce sujet :

\textbf{Proposition} L'homomorphisme (42) est un [unclear: isomorphisme] sur le fermé des systèmes satisfaisant
$$
\det \underline{u} = \det \underline{b} = \det \underline{U} \text{ soit } \mu, \text{ et } \det \underline{a} / \mu \in \Gamma_{[unclear: f/\mu]} \leqno{(43)}
$$

C'est une conséquence immédiate des corollaires à la fin de la p. 536', et de la définition de $S\mathcal{M}_{p,\sigma}$ comme produit fibré.

\newpage

%% ===== Page 476 =====

10°) Les sous-groupes-sections $M_{p, \sigma}(\sigma)$, $M_{p, \sigma}(p)$.

On définit $M_{p, \sigma}(\sigma)$ par la condition
$\beta_0 = 1$ i.e. $\beta = f$ i.e. (cf (40))
$$ \underline{U} = \underline{b} \leqno{(43)} $$
il est évident par prop. p. 536' que le sous-groupe
$$ M_{p,\sigma}(\sigma) \subset M_{p,\sigma} \leqno{(44)} $$
$$ \defeq \{ (\underline{a}, \underline{b}, \underline{U}) \in M_{p,\sigma} \mid \underline{U} = \underline{b} \} $$
forme une section — pour $\underline{a}, \underline{b}$ dans $T_{p,\sigma}$ satisfaisant la condition
$$ \det(\underline{a}) / \det(\underline{b}) \in \Gamma \mu_r $$
(la façon la plus simple de le relever en un élém. $(\underline{a}, \underline{b}, \underline{U})$ de $M_{p,\sigma}$ n'est-il pas par $(\underline{a}, \underline{b}, \underline{U} = \underline{b})$ !)

(Et on aura bien $\det(\underline{U}) = \det(\underline{b})$ ). -- On définit
de même $M_{p, \sigma}(p)$ par $\alpha_0 = 1$ i.e. $\alpha = f$ i.e.

$$ M_{p,\sigma}(p) \subset M_{p,\sigma} \leqno{(45)} $$
$$ \defeq \{ (\underline{a}, \underline{b}, \underline{U}) \in M_{p,\sigma} \mid \underline{U} = \underline{a} \} $$

ce qui fait une section partielle -- au-dessus du
sous-groupe $T'_{p,\sigma}$ de $T_{p,\sigma}$ défini comme

$$ T'_{p,\sigma} \subset T_{p,\sigma} \leqno{(46)} $$
$$ \defeq \{ (\underline{a}, \underline{b}) \in T_{p} \times \mathcal{N}'_0 \mid \det(\underline{a}) = \det(\underline{b}) \} $$
$$ = \operatorname{Ker}( T_{p,\sigma} \xrightarrow{[unclear: \varepsilon_p]} \mu_r ) $$

d'où
$$ T'_{p,\sigma} = T_{p} \times_{G_m} \mathcal{N}'_0 \quad \text{, [unclear: clair] .} \leqno{(47)} $$

\newpage

%% ===== Page 477 =====

donc
$$
\mathfrak{T}_{\dot{p},\sigma} / \mathfrak{T}_{\dot{p},\sigma}^\circ \xrightarrow{\epsilon_2} \mu_2 . \leqno{(48)}
$$

Remq. Ces sous-groupes peuvent se définir, et fournissent une section resp une section au dessus de $\mathfrak{T}_{\dot{p},\sigma}^\circ$, dès qu'on a (dans la situation générale de $\mathcal{P}$, p.\ 526' ff) des $p$-groupes $\mathcal{T}_{\dot{p}} \hookrightarrow G$, $\mathcal{N}'_\sigma \subset G$, tels que $\delta_0, \delta_\sigma$ soient induits par $\delta_G$.

Dans le cas modulaire, on écrit ainsi
$$
\mathcal{M}_{\dot{p},\sigma}(1/2) = \mathcal{M}_{\dot{p},\sigma}(\sigma), \quad \mathcal{M}_{\dot{p},\sigma}(-\rho) = \mathcal{M}_{\dot{p},\sigma}(\rho) , \leqno{(49)}
$$
et on définit plus généralement
$$
\mathcal{M}_{\dot{p},\sigma}(q) , \quad q \in \{ \rho, 1, \infty, 1/2, 2, -1 \} \leqno{(50)}
$$
à partir de $\mathcal{M}_{\dot{p},\sigma}(0)$ et $\mathcal{M}_{\dot{p},\sigma}(1/2)$ par\marginpar{[unclear: $\rho_0$ est $\rho$ ou tout autre él. du s-groupe $\simeq \mathfrak{S}_3$ de $\mathcal{M}_{\dot{p},\sigma}$]}
$$
\operatorname{int}(\rho) \mathcal{M}_{\dot{p},\sigma}(q) = \mathcal{M}_{\dot{p},\sigma}(\tilde{\rho}(q)) \leqno{(51)}
$$
où $\tilde{\rho}$ est la transformation homographique
$$
\tilde{\rho} : q \mapsto \frac{1}{1-q} \leqno{(52)}
$$
(dont l'itéré troisième est l'identité, tout comme l'itéré troisième de $\rho$).

Il resterait à définir $\mathcal{M}_{\dot{p},\sigma}(-1)$, et le choix de l'élément d'ordre deux [unclear: dans $\operatorname{Sl}(2, \hat{\mathbf{Z}})$ n'a guère d'influence, le groupe est bien déterminé, au choix près].

\newpage

%% ===== Page 478 =====

naturel serait entre
$$
\operatorname{int}(\sigma_i)(\mathcal{M}_{p,\sigma}(0)) \qquad i \in \{0,1,\infty\} \leqno (53)
$$

\newpage

%% ===== Page 479 =====

\chapter*{\S \space XII. --- RÉTROSPECTIVE}\thispagestyle{empty}
\addcontentsline{toc}{section}{XII. Rétrospective}
\label{sec:XII}
\section*{}

Par bien des détours inutilement compliqués, j'ai fini de façon bien simple par construire deux schémas en groupes sur $\mathbf{Z}$
$$
\mathcal{M}_{\rho,\sigma} \subset G \times T_\rho \times N'_\sigma \times G \leqno{(1)}
$$
$$
S\mathcal{M}_{\rho,\sigma} \subset B'_0 \times T_\rho \times N'_\sigma \times B'_0 \leqno{(2)}
$$
par
$$
\mathcal{M}_{\rho,\sigma} = \left\{ x = (\overset{x}{\underline{u}}, \underline{a}, \underline{b}, \underline{v}) \mid \overset{\chi(x)}{\det \underline{u}} = \det \underline{b} = \overset{\varepsilon_2(x)}{\varepsilon'} \det \underline{v}, \varepsilon'^2 = 1 \right\} \leqno{(3)}
$$
$$
S\mathcal{M}_{\rho,\sigma} = \left\{ (\underline{u}, \underline{a}, \underline{b}, \underline{v}) \mid \det \underline{u} = \det \underline{a} = \det \underline{b} = \varepsilon' \det \underline{v}, \varepsilon'^2 = 1 \right\} \leqno{(4)}
$$
avec
$$
\mathcal{M}_{\rho,\sigma} \simeq S\mathcal{M}_{\rho,\sigma} / L_0 \leqno{(5)}
$$
où $L_0 \hookrightarrow S\mathcal{M}_{\rho,\sigma}$ \newline
$\lambda \mapsto (1, 1, 1, \lambda)$

et
$$
S\mathcal{M}_{\rho,\sigma} \simeq \mathcal{M}_{\rho,\sigma} \cdot L_0 \quad \text{produit } \frac{1}{2} \text{ direct} \leqno{(6)}
$$
où
$$
\begin{cases}
\mathcal{M}_{\rho,\sigma} \hookrightarrow S\mathcal{M}_{\rho,\sigma} \\
x = (\underline{u}, \underline{a}, \underline{b}, \underline{v}) \mapsto (\underline{u}, \underline{a}, \underline{b}, i_{T_0}(\mu) \underline{v}) \text{ où } \mu = \chi(x)
\end{cases} \leqno{(7)}
$$

Ici (pour mémoire) dans le cas modulaire
$$
\rho = \begin{pmatrix} 0 & 1 \\ -1 & 1 \end{pmatrix} \quad \sigma = \begin{pmatrix} 0 & 1 \\ -1 & 0 \end{pmatrix} \leqno{(8)}
$$

$$
\begin{cases}
T_\rho = \{ (c\rho + d) \mid \underbrace{\det(c\rho + d)}_{c^2 + cd + d^2} \text{ inv.} \} \\
T_\sigma = \{ (t\sigma + u) \mid \underbrace{\det(t\sigma + u)}_{t^2 + u^2} \text{ inv.} \} 
\end{cases} \quad \begin{cases} B'_0 = \left\{ \begin{pmatrix} \mu & \lambda \\ 0 & 1 \end{pmatrix} \mid \mu \text{ inv.} \right\} \\ T_0 = \left\{ \begin{pmatrix} \mu & 0 \\ 0 & 1 \end{pmatrix} \mid \mu \text{ inv.} \right\} \\ i_{T_0}(\mu) = \begin{pmatrix} \mu & 0 \\ 0 & 1 \end{pmatrix} \\ L_0 = \left\{ \begin{pmatrix} \lambda & 0 \\ 0 & \lambda \end{pmatrix} \right\} \end{cases} \leqno{(9)}
$$

$$
N'_\sigma = \left\{ \underline{b} \text{ dans } G \mid \underline{b}(\sigma) = \underbrace{\varepsilon}_{\varepsilon(\underline{b})} \sigma, \ \varepsilon^2 = 1 \right\} \leqno{(10)}
$$
d'où suite exacte
$$
1 \to T_\sigma \to N'_\sigma \xrightarrow{\bar{\varepsilon}_\sigma} \mu_2 \to 1 \ , \leqno{(11)}
$$

\newpage

%% ===== Page 480 =====

et les trois homomorphismes
$$
\mathcal{M}_{\rho,\sigma} \xrightarrow{\Theta_G} G, \quad \mathcal{M}_{\rho,\sigma} \xrightarrow{\Theta_\rho} T_\rho, \quad \mathcal{M}_{\rho,\sigma} \xrightarrow{\Theta_\sigma} N_\sigma' \leqno{(12)}
$$
donnés dans la définition (1)(3) se complètent
par
$$
\mathcal{M}_{\rho,\sigma} \xrightarrow{\chi = \det \circ \Theta_G} G_m, \quad \mathcal{M}_{\rho,\sigma} \xrightarrow{\varepsilon_2 = \varepsilon_\sigma \circ \Theta_\sigma} \mu_2, \quad \mathcal{M}_{\rho,\sigma} \xrightarrow{\varepsilon_3(x) = \frac{\det \underline{a}}{\det \underline{b}}} \mu_2 \leqno{(12 \text{ bis})}
$$
de sorte qu'on a l'homomorphisme
$$
\mathcal{M}_{\rho,\sigma} \xrightarrow{(\Theta_G, \Theta_\rho, \Theta_\sigma, \chi, \varepsilon_2, \varepsilon_3)} G \times T_\rho \times N_\sigma' \times G_m \times \mu_2 \times \mu_2 \leqno{(13)}
$$
$$
x \longmapsto (\underline{U}, \underline{a}, \underline{b}, \mu, \varepsilon, \varepsilon')
$$
dont l'image est formée des systèmes satisfaisant à
$$
\begin{cases}
\det \underline{U} = \det \underline{b} = \varepsilon' \det \underline{a} = \mu \\
\varepsilon = \varepsilon_\sigma(\underline{b}) \text{ i.e. } \quad \underline{b}(\sigma) = \varepsilon \sigma \text{ i.e. } \underset{\varepsilon \underline{b}^{-1}}{\sigma \underline{U}_b} = \varepsilon \underline{b} \text{ i.e. } [unclear: \underline{b} \,^t\underline{b} = \varepsilon]
\end{cases} \leqno{(14)}
$$
De même on a un homomorphisme
$$
S\mathcal{M}_{\rho,\sigma} \subset \underline{B}_0' \times T_\sigma \times N_\sigma' \times \underline{B}_0 \times G_m \times \mu_2 \times \mu_2 \leqno{(15)}
$$
$$
x \longmapsto (\underline{U}, \underline{a}, \underline{b}, \underline{u}, \mu, \varepsilon, \varepsilon')
$$
dont l'image est formée des systèmes satisfaisant à
$$
\begin{cases}
\mu = \det \underline{u} = \det \underline{U} = \det \underline{b} = \varepsilon' \det \underline{a} \\
\varepsilon = \varepsilon_\sigma(\underline{b}) \text{ i.e. } \quad \underline{b}(\sigma) = \varepsilon(\sigma)
\end{cases} \leqno{(16)}
$$
Soit d'autre part
$$
S\Sigma_{\rho,\sigma}^\ast \hookrightarrow G \times G \times G \times G_m \times \mu_2 \leqno{(17)}
$$
$$
= \left\{ (\alpha, \beta, f, \mu, \varepsilon) \;\middle|\;
\begin{aligned}
&(\det \alpha)^2 = \det \beta = \det f = 1 \\
&\underbrace{\alpha \rho \alpha^{-1}}_{\xi} = \varepsilon \underbrace{\beta \sigma f \sigma^{-1}}_{\eta} \varepsilon_1(\mu-1)
\end{aligned}
\right\}
$$
où on a défini
$$
\varepsilon_1 = \begin{pmatrix} 1 & 0 \\ 1 & 1 \end{pmatrix} , \quad \varepsilon_1(\mu) = \varepsilon_1^\mu \defeq \begin{pmatrix} 1 & 0 \\ \mu & 1 \end{pmatrix} \leqno{(18)}
$$

\newpage

%% ===== Page 481 =====

en fait la projection
$$
\begin{cases}
S\Sigma_{\rho,\sigma}^* \hookrightarrow G \times G \times G \\
y \longmapsto (\alpha, \beta, f)
\end{cases} \leqno{(19)}
$$
est déjà un monomorphisme. On fait opérer \marginpar{à gauche $L_0$} $L_0 = \operatorname{Ker}(B'_0 \xrightarrow{\det} \mathbb{G}_m) = \left\{ \begin{pmatrix} l & 0 \\ 0 & l^{-1} \end{pmatrix} \right\}$ sur $S\Sigma_{\rho,\sigma}^*$ par
$$
(\alpha, \beta, f) \longmapsto (l\alpha, l\beta, lf) \leqno{(20)}
$$
et on pose
$$
\Sigma_{\rho,\sigma}^* \defeq L_0 \backslash S\Sigma_{\rho,\sigma}^* \leqno{(21)}
$$

Ceci posé, on a un isomorphisme
$$
\begin{array}{ccc}
S\mathcal{M}_{\rho,\sigma} & \isommap & S\Sigma_{\rho,\sigma}^* \\
(\underline{U}, \underline{a}, \underline{b}, \underline{U}') & \longmapsto & (\alpha, \beta, f)
\end{array}
\quad (\alpha = \underline{U}\underline{a}^{-1}, \beta = \underline{U}\underline{b}^{-1}, f = \underline{U}\underline{U}'^{-1}) \leqno{(22)}
$$
transformant l'action de $L_0$ sur $S\mathcal{M}_{\rho,\sigma}$ [unclear: naturelle] en l'action de $L_0$ décrite dans (20), et induisant un isomorphisme
$$
\mathcal{M}_{\rho,\sigma} \isommap \Sigma_{\rho,\sigma}^* \leqno{(23)}
$$
par passage au quotient.

On a des homomorphismes canoniques
$$
S\mathcal{M}_{0,3}^\sim \longrightarrow S\mathcal{M}_{\rho,\sigma}(\hat{\mathbf{Z}}) / \Sigma \leqno{(24)}
$$
$$
\mathcal{M}_{0,3}^\sim \longrightarrow \mathcal{M}_{\rho,\sigma}(\hat{\mathbf{Z}}) / \Sigma \leqno{(25)}
$$
le deuxième déduit du premier par passage au quotient, et coïncident sur $\mathfrak{S}_{0,3}^{+\wedge}$.

\newpage

%% ===== Page 482 =====

l'hom composé § 42 IV

(26) $$ \hat{G}_{0,3}^+ \overset{\sim}{\longrightarrow} S\ell_2(\hat{\mathbb{Z}})' \hookrightarrow M_{\rho,\sigma}(\hat{\mathbb{Z}})/\Sigma $$
$$ \parallel $$
$$ SG(\hat{\mathbb{Z}})/SG(\hat{\mathbb{Z}}) \cap \Sigma $$

grâce à l'hom canonique

(27) $$ SG \hookrightarrow M_{\rho,\sigma} $$
$$ S\ell(2) \cong \left\{ (\underline{u}, \underline{a}, \underline{b}) \mid \underline{a} = \underline{b} = 1 \right\} $$

On a défini

(28) $$ \Sigma = (\mathbb{Z}/2\mathbb{Z})^3 \hookrightarrow M_{\rho,\sigma}(\mathbb{Z}) \subset M_{\rho,\sigma}(\hat{\mathbb{Z}}) $$
$$ \parallel $$
$$ (\mu_2)^3(\mathbb{Z}) $$

via l'hom canonique (central)

(29) $$ \mu_2^3 \longrightarrow M_{\rho,\sigma} \subset T_{\rho} \times N'_{\sigma} \times G $$

défini via les composantes canoniques

(30) 
\begin{tikzcd}
& T_\rho \\
\mu_2 \ar[r] & \mathbb{G}_m \ar[u] \ar[r] \ar[d] & T_\sigma \ar[r, hook] & N'_\sigma \\
& G
\end{tikzcd}

Par passage au quotient, l'hom (25)
induit

(31) $$ \hat{\Gamma}_{0,3} \longrightarrow \Gamma_{\rho,\sigma}(\hat{\mathbb{Z}})/ \underbrace{\{1 \times \pm 1 \}}_{\Sigma_0} $$

où

(32) $$ \Gamma_{\rho,\sigma} \subset T_{\rho} \times N'_{\sigma} $$
$$ = \left\{ (\underline{a}, \underline{b}) \mid \det \underline{a} / \det \underline{b} \in \mu_2 \right\} $$

La donnée - plutôt la construction faite
plus haut - donne un hom de groupes (24)

\newpage

%% ===== Page 483 =====

revient à la construction des quatre homomorphismes de groupes

$$
S_0\mathcal{M}_{0,3}^\sim \xrightarrow[\theta_b^\pi]{\sim} B'_0(\hat{\mathbf{Z}}) = \hat{\mathbf{Z}}^* \cdot \hat{\mathbf{Z}} \quad \text{(donnant $u$)} \leqno{(33)}
$$
$$
\mathcal{M}_{0,3}^\sim \xrightarrow[\theta_G^\pi]{\sim} G(\hat{\mathbf{Z}})/\pm 1 = \operatorname{Gl}(2, \hat{\mathbf{Z}})/\pm 1 \quad \text{(donnant $U$ (mod signe) (produit $1/2$ direct))} \leqno{(34)}
$$
$$
\Gamma_{0,3}^\sim \xrightarrow[\theta_P^\pi]{\sim} T_P(\hat{\mathbf{Z}})/\pm 1 \quad \text{(donnant $a$ (mod signe))} \quad (\text{unités adéliques (mod } \pm 1) \text{ du corps } \mathbf{Q}(j)) \leqno{(35)}
$$
$$
\overline{\Gamma}_{0,3}^\sim \xrightarrow[\theta_{\overline{P}}^\pi]{\sim} T_0(\hat{\mathbf{Z}})/\pm 1 \quad \text{(donnant $b$ (mod signe))} \quad (\text{unités adéliques (mod } \pm 1) \text{ du corps } \mathbf{Q}(i)) \leqno{(36)}
$$

reliés par les relations (4), que nous allons identifier. On utilise le sous-groupe
$$
\mathcal{M}_{0,3}^\sim [0] = \text{sous-groupe de } \mathcal{M}_{0,3}^\sim \text{ qui} \leqno{(37)}
$$
(invarie les [unclear: éléments de base] de $\mathfrak{T}_{0,3}^\sim$ , et qui) normalise $L_{0 \pi}$ (sous-groupe fini engendré par $l_0$)

et on note que
$$
1 \to L_{0\pi} \to \mathcal{M}_{0,3}^\sim [0] \to \Gamma_{0,3}^\sim \to 1 \leqno{(38)}
$$
et que
$$
\begin{cases}
\mathcal{M}_{0,3}^\sim = \mathcal{M}_{0,3}^\sim[0] \cdot \widehat{\mathfrak{S}}_{0,3}^{+\wedge} \quad \text{(sous-groupe invariant de } \mathcal{M}_{0,3}^\sim) \\
\mathcal{M}_{0,3}^\sim[0] \cap \widehat{\mathfrak{S}}_{0,3}^{+\wedge} = L_{0\pi}
\end{cases} \leqno{(39)}
$$

--- bien sûr, (38) est conséquence des (39).

enfin
$$
S_0\mathcal{M}_{0,3}^\sim \simeq \mathcal{M}_{0,3}^\sim[0] \cdot \widehat{\mathfrak{S}}_{0,3}^{+\wedge} \quad \text{(produit $1/2$ direct)} \leqno{(40)}
$$

1º) L'homomorphisme (33) est trivial sur le facteur $\widehat{\mathfrak{S}}_{0,3}^{+\wedge}$

\newpage

%% ===== Page 484 =====

\marginpar{La dépendance de $\underline{u}$ (description de (33) restreinte à $\mathcal{M}_{0,3}^{\sim}[0] \to B'_0(\hat{\mathbf{Z}})$)}

$\underline{u}$ ne dépend que du couple $(\alpha, \beta)$ (ou du triplet $(\alpha, \beta, \gamma)$ qui décrit un élément de $\mathcal{M}_{0,3}^{\sim}$), et la [unclear: pré]définition de $\underline{u}$ en termes de $\alpha, \beta$ -- ou seulement en termes de l'image de $\xi = \alpha\beta$ dans $\operatorname{Sl}(2, \hat{\mathbf{Z}})$ [unclear: en omettant l'image ch $\dots$] a été faite essentiellement dans le § 45 ou tout au moins dans § 45 III, -- on a
$$
\underline{u} = \begin{pmatrix} 1/D & B/D \\ 0 & 1 \end{pmatrix} \sim \xi = \begin{pmatrix} A & B \\ C & D \end{pmatrix} \quad (\text{dans } \operatorname{Sl}(2, \hat{\mathbf{Z}})) \leqno{(41)}
$$
ou encore, en termes de
$$
\underline{\alpha} = \begin{pmatrix} a & b \\ c & d \end{pmatrix} \dots \text{ avec } ad-bc \in \{\pm 1\}
$$
$$
\begin{cases}
\varepsilon' D = c^2 + cd + d^2 \\
\varepsilon' B = ac + bd + bc
\end{cases} \leqno{(42)}
$$

Cela ne dépend pas de fondements non faits sur ces groupes profinis -- les ``calculs en $S,T$'' nous en [unclear: affranchissent]. C'est grâce à la description de cet homomorphisme qu'on peut définir
$$
\mathcal{M}_{0,3}^{\sim}(0) \subset \mathcal{M}_{0,3}^{\sim}[0] \leqno{(43)}
$$
$$
= \{ \underline{u} \mid \underline{u} \in T_0(\hat{\mathbf{Z}}) \text{ i.e. } B=0 \}
$$
$$
(\text{condition de normalisation, i.e. } ac+bd+bc=0)
$$

\newpage

%% ===== Page 485 =====

et obtenir des décompositions canoniques
en produits semi-directs
$$ \mathcal{M}_{0,3}^{\sim}[0] \simeq \mathcal{M}_{0,3}^{\sim}(0) . L_0 \pi \leqno{(44)} $$
$$ \mathcal{M}_{0,3}^{\sim} \simeq \mathcal{M}_{0,3}^{\sim}(0) . \mathfrak{S}_{0,3}^{+\wedge} \leqno{(45)} $$

2º) Pour décrire (34), comme il est déjà
connu sur $\mathfrak{S}_{0,3}^{+\wedge}$, il suffit de le décrire
sur $\mathcal{M}_{0,3}^{\sim}[0]$ ou seulement sur
$\mathcal{M}_{0,3}^{\sim}(0)$ (moyennant bien sûr des compati-
bilités...). Et sur $\mathcal{M}_{0,3}^{\sim}[0]$ il coïncide
avec (33) --- tandis que sur $\mathcal{M}_{0,3}^{\sim}(0)$ il
est pratiquement trivial et se réduit au
multiplicateur $\mu$
$$ \Theta_{\underline{u}}^\pi = \Theta_{B'_0}^\pi \quad \text{sur } \mathcal{M}_{0,3}^{\sim}[0] \leqno{(46)} $$
$$
\begin{cases}
\Theta_{\underline{u}}^\pi = i_{T_0}(\mu) & \text{où } \mu \text{ est le} \\
& \text{multiplicateur de } \underline{u} \\
\text{si } \underline{u} \in \mathcal{M}_{0,3}^{\sim}(0)
\end{cases} \leqno{(47)}
$$
Cela = [unclear: une définition possible de] $\mathcal{M}_{0,3}^{\sim}(0)$.

Remarque : Comme $\Theta_{\underline{u}}^\pi$ a été construit
"à la force des poignets" par l'intermédiaire
de $\Theta_{B'_0}^\pi$ i.e. $\underline{U}$ a été construit par $\underline{u}$,
il y a eu longtemps une [unclear: espèce] de

\newpage

%% ===== Page 486 =====

confusion dans mon esprit entre $\underline{u}$ et $\underline{v}$,
d'autant plus que la plus grande partie
du temps je me réduisais à travailler
dans $\mathcal{M}_{0,3}^{\sim}[0]$ (après avoir pas mal
cafouillé avec des points de vue pure-
ment [unclear: globaux], avant de me rendre compte
que techniquement c'était la façon
la plus élégante de procéder --- à
partir de \S~45). Cette confusion ne
s'est définitivement clarifiée que hier
et aujourd'hui, à force de mettre et
remettre les pieds dans mes propres
plats...

3º) Les homomorphismes (35)(36) n'ont été obtenus
qu'en travaillant à fond sur les
paramètres $\underline{a},\underline{b}$ pour $\underline{u}$ --- bien entendu,
il suffit de travailler dans $\mathcal{M}_{0,3}^{\sim}[0]$,
et de définir $\underline{u} \mapsto \underline{a}$
$$
\mathcal{M}_{0,3}^{\sim}[0] \longrightarrow T_l(\hat{\mathbf{Z}}) / \pm 1
$$
$\underline{u} \mapsto \underline{b}$
$$
\longrightarrow \text{[unclear: raturé]} / \pm 1
$$

\newpage

%% ===== Page 487 =====

de façon que $a$ soit trivial sur $L_0^{\boldsymbol{\Gamma}}$.
Les quantités $a, b$ sont apparues dans les
calculs progressivement à partir du
§ 44 (p. 366 ff), la lumière sur leur
rôle a commencé à se faire jour seulement
à partir des § 48 III, IV, V - et une compré-
hension complète de leur nature, du point
de vue des schémas en groupes qui inter-
viennent dans la théorie, s'est fait égale-
ment seulement en ces tout derniers
jours. Cette compréhension reste d'
ailleurs pour le moment assez
formelle, et il faudrait approfondir
dans plusieurs directions :

\begin{enumerate}
\item[a)] Expliciter la signification des points
de vue des corps de classes de
$\mathbf{Q}(j)$ et $\mathbf{Q}(\sqrt[3]{j})$.

\item[b)] Expliciter la signification pour
la théorie des courbes "[unclear: c-quantifiées]"
(i.e. sous-entendu : des [unclear: représ.]
du quotient $\operatorname{Sl}(2, \hat{\mathbf{Z}})$ ou $\operatorname{Sl}(2, \mathbf{Z}/N)$)
des résultats obtenus,
\end{enumerate}

\newpage

%% ===== Page 488 =====

\begin{enumerate}
\item[c)] Introduire sur $\mathcal{M}_{g,\nu}$ les éléments de structure caractéristiques des groupes de Galois motiviques - et faire le lien avec les constructions typiques de cet ordre en théorie des corps de classes.

\item[d)] Examiner dans quelle mesure dans ce contexte on peut faire intervenir les conditions de compatibilité diophantiennes draconiennes sus-dites du § 47 liées à la compatibilité avec
$$ \varphi : \hat{\Pi}_{0,3} \isommap \mathfrak{S}^{+\wedge}_{0,3} $$
- je prévois d'ailleurs que le cadre des calculs précédents est trop étroit pour espérer y faire intervenir ces conditions supplémentaires.

\item[e)] Ce qui m'intéresse le plus pour le moment, c'est une systématisation des calculs précédents dans un cadre beaucoup plus vaste que celui du seul groupe $\mathfrak{S}_g$ et de l'homomorphisme de Torelli.
\end{enumerate}

\newpage

%% ===== Page 489 =====

$$
\mathfrak{S}_{0,3}^{+\wedge} \to \operatorname{Gl}(\hat{\mathbf{Z}})/\pm 1
$$

- j'ai le sentiment d'avoir mis le doigt sur un procédé systématique pour construire des représentations "motiviques" de $\Gamma_{\mathbf{Q}}$, via des représentations de la partie "géométrique" $\mathfrak{S}_{0,3}^{+\wedge}$ de $\mathcal{M}_{0,3} \simeq \pi_1(M_{0,3 \mathbf{Q}}, [unclear: \vec{v}_0]) \simeq \pi_1(M_{1,1 \mathbf{Q}})$ - ou mieux, via des représentations du groupe $\mathfrak{S}_{0,3}^+ \simeq \operatorname{Sl}(2, \mathbf{Z})'$ discret (ou de variantes telles $\operatorname{Sl}(2, \mathbf{Z})$, $\Pi_{0,3} \dots$ ) à valeurs dans les points entiers de schémas en groupes (réductifs ??) sur $\mathbf{Z}$.

\begin{enumerate}
\item[f)] Il y a pourtant une autre direction d'investigation qui m'attire beaucoup, pas nécessairement liée à l'étude exclusive des schémas en groupes, et qui pourrait se poursuivre sur la base des seuls "gros" groupes fondamentaux profinis -
\end{enumerate}

\newpage

%% ===== Page 490 =====

savoir l'étude de certains $\mathbf{Q}$-revêtements particuliers de $M_{0,3\mathbf{Q}} = \mathbf{P}_{1,\mathbf{Q}} - \{0,1,\infty\}$, et des renseignements qu'on doit en tirer sur des propriétés supplémentaires de l'action ext.\ "arithmétique" de $\pi_{0,3}$ (du type des résultats obtenus dans [unclear: § 4]), et sur des calculs de relevages de l'extension $M_{0,3}$ ou $\pi_{0,3}$ par $\mathfrak{S}_{0,3}^{+\wedge}$ en d'autres points que les points "arithmétiques" $0, 1, \infty, 1/2, 2, -1, -j, \bar{j}$.

(Sans compter la question d'asseoir les fondements, via la méthode suggérée par Deligne - et d'étudier l'équation des droites dans la compactific.\ profinie [unclear: $L_4$] et des coniques : investigations dans le discret. - des questions fondamentales, intimement liées à l'arithmétique des points complexes - etc etc.).

\newpage

%% ===== Page 491 =====

Revenons pour l'instant à la question d'identifier les homomorphismes obtenus (33) \& (36), en termes de choses classiques. On peut considérer l'homomorphisme (33), $\Theta_{B'_4}^F$, comme de nature essentielle-ment technique, ou du moins comme étant subordonné : l'homomorphisme $\Theta_G^F$ de (34). D'ailleurs, (35) et (36) s'en déduisent, en composant avec les deux homomorphismes
$$
\Pi_{0,3}^\sim \longrightarrow \mathcal{M}_{0,3}^\sim
$$
(morphismes de groupes, pour le moment !) associés aux deux points $1/2$ et $-j$ de $\mathcal{M}_{0,3 \mathbf{Q}}(\mathbf{C})$, et qui correspondent, au niveau des schémas en groupes, aux deux scindages $\underset{(3)}{y_\rho}, \underset{(31)}{y_\sigma}$ de $\mathcal{M}_{\rho,\sigma}$ sur $\mathfrak{T}_{\rho,\sigma}$, donnés respectivement par
$$ 
\begin{cases} 
y_\rho : (\underline{a}, \underline{b}) \longmapsto (\underline{a}, \underline{b}, \underline{a}) \\ 
y_\sigma : (\underline{a}, \underline{b}) \longmapsto (\underline{c} = \underline{a}\underline{b}, \underline{b}) 
\end{cases} \leqno{(38)} 
$$
avec un grain de sel que le premier n'est défini que sur le sous-groupe $\mathfrak{T}'_{\rho,\sigma}$ défini par
$$ 
\begin{cases} 
\mathfrak{T}'_{\rho,\sigma} \subset \mathfrak{T}_{\rho,\sigma} \\ 
\text{''} = \{ (\underline{a}, \underline{b}) \in \mathfrak{T}_{\rho,\sigma} \mid \det \underline{a} = \det \underline{b} \} \\ 
\text{''} \isom \mathfrak{T}_\rho \times N'_0 
\end{cases} \leqno{(49)} 
$$

\newpage

%% ===== Page 492 =====

et de
$$
\begin{cases}
\hat{\Gamma}_{\mathbf{Q}}^{\pi} : \hat{\Gamma}_{\mathbf{Q}}^{1 \sim} \simeq \hat{C}_3^{1\sim \prime} \longrightarrow T_g(\hat{\mathbf{Z}})/\pm 1 \\
\tilde{\hat{\Gamma}}_{\mathbf{Q}}^{\pi} : \tilde{\hat{\Gamma}}_{\mathbf{Q}}^{1 \sim} = \left\{ \gamma \in \hat{\Gamma}_{\mathbf{Q}}^{1 \sim \prime} \mid \varepsilon_3(\mu) = 1 , \dots \mu = \chi(\gamma) \right\}
\end{cases}
$$

On voit donc\footnote{Au fait l'hom. $\Theta_{\mathbf{Q}}^\pi$ s'identifie à l'hom. continu $\Gamma_{\mathbf{Q}} \to \Out(\hat{\pi}_{1,1})$ à déf. d'une section près.} que les signes pour $\hat{\Theta}_{\mathbf{Q}}^\pi$, $\tilde{\hat{\Theta}}_{\mathbf{Q}}^\pi$ correspondent à des indéterminations pour $\Theta_{\mathbf{Q}}^\pi$ lui-même - tout se tient ! Cette dernière est levée par [unclear: le remplacement] en remplaçant $\mathcal{M}_{0,3}^\sim$ par une extension de degré 2 $\tilde{\mathfrak{T}}_{1,1}^\sim \to \pm 1$, de sorte qu'on trouve
$$
\tilde{\mathfrak{T}}_{1,1}^\sim \longrightarrow \operatorname{Gl}(2, \hat{\mathbf{Z}}) \leqno{(50)}
$$
et il faudrait [unclear: prendre] cet homomorphisme comme l'homomorphisme fondamental "anabélien" de la théorie des courbes elliptiques (où l'image inverse de $\Gamma_{\mathbf{Q}} \subset \tilde{\Gamma}_{\mathbf{Q}}^\sim$ dans $\tilde{\mathfrak{T}}_{1,1}^\sim$ peut s'interpréter comme $\tilde{\mathfrak{T}}_{1,1} \simeq \pi_1(\mathcal{M}_{1/\mathbf{Q}})$ dont on fera plus un usage [unclear: topologique] par la suite).

Je ne vois clair en tous cas que sur les points de $\mathcal{M}_{0,3}^\sim$, un [unclear: choix] de deux [unclear: coïncidant aux] éléments de
$$
\mu_2(\hat{\mathbf{Z}}) = \{\pm 1\}^\mathcal{P} \quad (\mathcal{P} = \text{ens des nbs premiers}) \leqno{(51)}
$$

\newpage

%% ===== Page 493 =====

i.e. ils coïncident modulo un hom.
$$
\widetilde{\boldsymbol{\Gamma}}_{\mathbf{Q}} \longrightarrow \mu_2(\hat{\mathbf{Z}}) \simeq \{\pm 1\}^{\mathbb{P}}
$$
i.e. un hom.
$$
\mathbb{G}_m(\hat{\mathbf{Z}})_2 = (\hat{\mathbf{Z}}^\ast)_2 \simeq \{\pm 1\}^{\mathbb{P}} \longrightarrow \mu_2(G_m(\hat{\mathbf{Z}})) = \mu_2(\hat{\mathbf{Z}}) \simeq \{\pm 1\}^{\mathbb{P}}
$$

Je [unclear: donne ça] pour un [unclear: calcul] ind. qui est un [unclear: casse-tête] ! C'est moralement sûr que les deux coïncident -- je n'ai pas essayé de le prouver.

Quant à (35) (36), ils doivent provenir du résumé : l'action de $\boldsymbol{\Gamma}_{\mathbf{Q}}$ sur les modules de Tate des courbes elliptiques à isomorphismes exceptionnels, correspondant justement aux pts $(1/2, -1)$ et $(-j, -\bar{j})$ de $U_{0,3 \mathbf{Q}} = M_{1,1}[2]^\prime_{\mathbf{Q}}$. L'indétermination de signes dont [unclear: j'ai parlé] correspond à l'indétermination à remonter un point de $M_{1,1}[2]^\prime_{\mathbf{Q}}$ en un point de $M_{1,1}[2]_{\mathbf{Q}}$ , i.e. de choisir une courbe elliptique (avec "rigidification" d'échelon 2) correspondant aux données de Legendre $(0, 1, \infty, \rho)$ avec $\rho = 1/2$ ou $\rho = -j$.

\newpage

%% ===== Page 494 =====

Mieux une fois que l'ambiguïté de signes pour $\Theta_{\mathbf{Q}}^\pi$ est levée en prenant : l'extension $\widetilde{\mathfrak{T}}_{1,1}$ et des feuillets sur $\widetilde{\mathcal{M}}_{0,3}$, l'indétermination de signes pour $\Theta_I^\pi$, $\Theta_J^\pi$ n'est pas levée pour autant, car les demi-tours $\psi_{1/2}^\pi, \psi_{-j}^\pi$ ne se remontent pas [unclear: en] des demi-tours (ni p. strictes) des extensions $\widetilde{\mathfrak{T}}_{1,1}$, mais seulement de $\widetilde{\mathcal{M}}_{0,3}$ --- et des relèvements usuels de ces demi-tours de $\widetilde{\mathfrak{T}}_{1,1}$ sur [unclear: $\Gamma_{0,3}$] revient très précisément, à choisir justement des courbes elliptiques à rigidification de Legendre au-dessus des pts $1/2, -j$. Ainsi, des calculs faits en termes de l'action du groupe à lacets $\widehat{\Pi}_{0,3}\tilde{\ }$ (de $\widehat{U}_{0,3}^+$), (tous les soins bien choisis!), donnent des résultats qui collent de très près : des réalités arithmético-géométriques particulièrement cruciales.

\newpage

%% ===== Page 495 =====

\chapter*{\S \space 49. --- HOMOMORPHISMES DE $\mathcal{M}_{0,3}$, LES GROUPES $\mathcal{M}_{g,\nu}$ GÉNÉRAUX}\thispagestyle{empty}
\addcontentsline{toc}{section}{49. Homomorphismes de $\mathcal{M}_{0,3}$, les groupes $\mathcal{M}_{g,\nu}$ généraux}
\label{sec:49}
\section*{}

I) Considérons un sous-groupe fermé $G$ de $\mathcal{M}_{0,3}$, contenant $\mathfrak{S}_{0,3}^{+\wedge}$, donc image inverse d'un sous-groupe fermé de $\Gamma_{0,3}^\sim$\footnote{Sauf dit le contraire. Toute cette section semble inutile si on remplace la [unclear: note] ci-contre dans le texte [unclear: p. 533'] par l'énoncé préliminaire de la débandade jetée !! p. 533'}. On suppose que celui-ci est «~suffisamment gros~» --- disons qu'il contienne un sous-groupe d'indice fini de $\Gamma_{\mathbf{Q}} \subset \Gamma_{0,3}^\sim$. Pour le moment, on supposera pour simplifier que $G \subset \Ker \varepsilon_3$. Cela revient à dire que
$$
G \subset \mathfrak{S}_{0,3}^{+\wedge} \mathcal{M}_{0,3}^\sim \qquad \text{où } g \mapsto \operatorname{perm} \equiv 1 \ (3) \leqno{(1)}
$$
Cette dernière hypothèse implique donc que pour tout $u \in G$,
$$
u \xrightarrow{\operatorname{int}(u)} \mathfrak{S}_{0,3}^{+\wedge} \leqno{(2)}
$$
s'exprime par
$$
\begin{cases}
u(\rho) = \operatorname{int}(\alpha_0) \rho \\
u(\sigma) = \operatorname{int}(\beta_0) \sigma
\end{cases} \leqno{(2')}
$$
On considère
$$
\begin{cases}
G_{(0)} = G \cap \mathcal{M}_{0,3}^\sim(0) \\
G_{[0]} = G \cap \mathcal{M}_{0,3}^\sim[0] = \operatorname{Norm}_G(L_0^\oplus)
\end{cases} \leqno{(3)}
$$
de sorte que
$$
G = G_{(0)} \cdot \mathfrak{S}_{0,3}^{+\wedge} \qquad \text{(produit semi-direct)} \leqno{(4)}
$$
$$
G_{[0]} = G_{(0)} \cdot L_0^\oplus \qquad \text{id} \leqno{(5)}
$$
$$
G_{[0]} \cap \mathfrak{S}_{0,3}^{+\wedge} = L_0^\oplus \leqno{(6)}
$$
On a posé
$$
L_0^\oplus = \text{sous-groupe fermé de } \mathfrak{S}_{0,3}^{+\wedge} \text{ engendré} \leqno{(7)}
$$
par $\tilde{L}_0, \hat{L}_0$.

\newpage

%% ===== Page 496 =====

Considérons un groupe pro-fini $\mathcal{M}$, et un homomorphisme surjectif continu
$$ \varphi : \mathcal{G} \to \mathcal{M}. \leqno{(8)} $$
On pose
$$ \begin{cases} \varphi(\mathfrak{S}^{+\wedge}_{0,3}) = \mathfrak{S} \quad \text{sous-groupe distingué de } \mathcal{M} \\ \varphi(\mathcal{G}_{(\infty)}) = \mathcal{M}_{(\infty)} \\ \varphi(\mathcal{G}_{[0]}) = \mathcal{M}_{[0]} \\ \varphi(L^\sim_0) = L^\circ_0 \end{cases} \leqno{(9)} $$
Donc ce sont des sous-groupes fermés de $\mathcal{M}$, satisfaisant
$$ \mathcal{M}_{(\infty)} \subset \mathcal{M}_{[0]} \subset \operatorname{Norm}_{\mathcal{M}}(L^\circ_0) \leqno{(10)} $$
$$ L^\circ_0 \subset \mathfrak{S} \cap \mathcal{M}_{[0]} \leqno{(11)} $$
Par abus de langage, on désigne encore par $\rho, \sigma, \varepsilon_0, \varepsilon_1$
les images de ces éléments dans $\mathfrak{S}$, qui y
satisfont donc les relations connues dans $\mathfrak{S}^{+\wedge}_{0,3}$ ,
en particulier
$$ \sigma^2 = \rho^3 = 1, \quad \sigma\rho = \varepsilon_0, \quad \rho\sigma = \varepsilon_1, \quad \varepsilon_1 = \sigma(\varepsilon_0) = \rho(\varepsilon_0). \leqno{(12)} $$
Les restrictions des opérations $\alpha_0, \beta_0$ de (2, 2')
à $\mathcal{M}_{0,3}[0]$ seront plutôt notées $\alpha, \beta$, où $\alpha, \beta$ sont
en fait des opérations
$$ \begin{gathered} \mathcal{S}\mathcal{G} \underset{\beta}{\overset{\alpha}{\rightrightarrows}} \mathfrak{S}^{+\wedge}_{0,3} \\ \text{défini sous-groupe de } \mathcal{G} \times \mathfrak{S}^{+\wedge}_{0,3} \\ = \{ (u,f) \in \mathcal{G} \times \mathfrak{S}^{+\wedge}_{0,3} \mid f u \in \mathcal{G}_{[0]} \} \end{gathered} \leqno{(13)} $$
dont on repère les restrictions au sous-groupe ($\isom \mathcal{G}_{[0]}$) de $\mathcal{S}\mathcal{G}$ défini par la relation $f=1$.
Comme on a [unclear: $\alpha, \beta : \mathcal{G}_{[0]} \rightrightarrows \mathfrak{S}^{+\wedge}_{0,3}$] ou $\varphi(\mathfrak{S}^{+\wedge}_{0,3})$ ,
donnent des opérations, encore notées $\alpha, \beta$ [unclear: (ici à une confusion possible)]
$$ \mathcal{G}_{[0]} \underset{\beta}{\overset{\alpha}{\rightrightarrows}} \mathfrak{S} \leqno{(14)} $$
satisfaisant les
$$ \begin{cases} [unclear: \alpha(x \rho) = \alpha_0(x) \rho] \\ [unclear: \beta(x \sigma) = \beta_0(x) \sigma] \end{cases} \quad \text{pour } x \in \mathcal{M}_{[0]} \text{ [unclear: formés] } \varphi(u_0), \alpha = \alpha(u_0), \beta = \beta(u_0). \leqno{(15)} $$

\newpage

%% ===== Page 497 =====

A priori, $\alpha$ et $\beta$ dépendent des choix de $u_0$ dans ce procés
pas seulement de $u$ - mais
$$
\begin{cases}
\alpha \rho(\alpha)^{-1} = [u, \rho] \\
\beta \sigma(\beta)^{-1} = [u, \sigma]
\end{cases}
\leqno{(16)}
$$
qui n'en dépendent pas. Si on pose
$$
\Gamma_\rho = \operatorname{Centr}_{\mathcal{M}}(\rho) \qquad \Gamma_\sigma = \operatorname{Centr}_{\mathcal{M}}(\sigma)
\leqno{(17)}
$$
on voit donc que, modulo $\Gamma_\rho$ resp.\ $\Gamma_\sigma$, $\alpha, \beta$ ne dépendent
que de $u$.

\underline{Hypothèse} $\alpha_{(= \alpha(u_0))}, \beta_{(= \beta(u_0))}$ ne dépendent que de $u = \varphi(u_0)$ (pour $u_0 \in \mathcal{G}[0]$).

On exprime cette hypothèse de façon un peu
différente, en introduisant le
$$
J = \operatorname{Ker}(\mathcal{G}[0] \to \mathcal{M}[0]) ,
\leqno{(18)}
$$
on veut donc que pour
$u'_0 = u_0 \delta_0$ avec $u_0 \in \mathcal{G}[0], \delta_0 \in J$, on ait
$\alpha(u'_0) = \alpha(u_0)$, $\beta(u'_0) = \beta(u_0)$

Soit $\alpha_0 = \underline{\alpha}(u_0)$, $\beta_0 = \underline{\beta}(u_0)$, $\alpha'_0 = \underline{\alpha}(u'_0)$, $\beta'_0 = \underline{\beta}(u'_0)$, $\overline{\alpha}_0 = \underline{\alpha}(\delta_0)$, $\overline{\beta}_0 = \underline{\beta}(\delta_0)$, dans $\mathcal{G}$,
de sorte qu'on aura, par les formules de cocycles dans $\mathcal{G}$
$\alpha'_0 = {}^{\delta_0}(\alpha_0) \overline{\alpha}_0$, $\beta'_0 = {}^{\delta_0}(\beta_0) \overline{\beta}_0$
d'où, en appliquant $\varphi$ et notant que
$\varphi({}^{\delta_0}(\alpha_0)) = {}^{\varphi(\delta_0)}(\varphi(\alpha_0)) = {}^1(\varphi(\alpha_0)) = \varphi(\alpha_0) \defeq \alpha$ et de même
$\varphi({}^{\delta_0}(\beta_0)) = {}^1\varphi(\beta_0) \defeq \beta$
$$
\alpha' = \alpha \overline{\alpha}, \quad \beta' = \beta \overline{\beta} \qquad \text{dans } \mathfrak{S}
\leqno{(19)}
$$
où $\alpha, \beta, \alpha', \beta', \overline{\alpha}, \overline{\beta}$ sont les images par $\varphi$ des quantités
$\alpha_0$ etc de $\mathfrak{S}_{0,3}^{+\wedge}$, d'où il ressort que l'hypothèse
se quantifie par :
$$
\forall \delta_0 \in J = \operatorname{Ker}(\mathcal{G}[0] \to \mathcal{M}[0]), \quad \underline{\alpha}(\delta_0), \underline{\beta}(\delta_0) \in \Ker \varphi
\leqno{(20)}
$$
Quand cette hypothèse est satisfaite, on a dans des
notations

\newpage

%% ===== Page 498 =====

$$
\mathcal{M}[\sigma_0] \overset{\alpha}{\underset{\beta}{\rightrightarrows}} \mathfrak{S} \leqno{(21)}
$$
factorisant (14), et satisfaisant aux
$$
\begin{cases} 
u \rho u^{-1} = i(\alpha) \rho \\ 
u \sigma u^{-1} = i(\beta) \sigma 
\end{cases} \quad \text{pour } u \in \mathcal{M}[\sigma_0], \ \alpha = \alpha(u), \ \beta = \beta(u) \leqno{(22)}
$$
qui généralisent (2').

D'ailleurs, posons
$$
\xi = \alpha \rho \alpha^{-1}, \quad \eta = \beta \sigma \beta^{-1} \leqno{(23)}
$$
les relations (22) s'écrivent aussi
$$
u \rho u^{-1} = \xi \rho, \quad u \sigma u^{-1} = \eta \sigma \leqno{(24)}
$$
et on a
$$
\xi = [u, \rho], \quad \eta = [u, \sigma] \leqno{(25)}
$$
et de plus
$$
\xi = \eta l_1^\nu \quad \text{où } \nu = \frac{\mu - 1}{2} \in \hat{\mathbf{Z}} \text{ (cf le multiplicateur de } u) \leqno{(26)}
$$
qui proviennent de la relation analogue dans $\mathfrak{S}_{0,3}^+$.
(On suppose pour fixer les idées que l'homomorphisme se factorise par
$$
\mathcal{G} \to \mathcal{M} \xrightarrow{\chi^\mathcal{M}} \hat{\mathbf{Z}}^\times, \leqno{(27)}
$$
de sorte qu'on a une action par multiplicateurs ($\in \hat{\mathbf{Z}}^\times$) pour les $u \in \mathcal{M}$.)

Considérons l'homomorphisme naturel
$$
\begin{tikzcd}
\mathcal{M}[\sigma_0] \arrow[r] & \operatorname{Aut}(\mathfrak{S}) \\
u \arrow[r, mapsto] & \underset{\text{via (24)}}{i_u}
\end{tikzcd} \leqno{(28)}
$$
on voit sur (22) qu'il est connu quand on connait

\newpage

%% ===== Page 499 =====

\marginpar{au point de vue topol. de $\mathfrak{S}$}les fonctions $\alpha, \beta : \mathcal{M}[\sigma] \longrightarrow \mathfrak{S}$, et $u \longmapsto$ [unclear: en fait] seulement $\xi, \eta : \mathcal{M}[\sigma] \longrightarrow \mathfrak{S}$. Si
D'ailleurs on pose
$$
\mathcal{M}'[\sigma] = \operatorname{Im}(\mathcal{M}[\sigma] \longrightarrow \operatorname{Aut}(\mathfrak{S})) \leqno{(29)}
$$
ces mêmes fonctions $\xi, \eta$ sont déjà définies sur $\mathcal{M}'[\sigma]$ - alors qu'il n'est pas clair qu'il en soit de même de $\alpha, \beta$, quand $\mathcal{M}[\sigma]$ n'opère pas fidèlement sur $\mathfrak{S}$. On a ainsi une application $(\xi, \eta)$ d'un certain sous-groupe $\mathcal{M}'[\sigma] \subset \operatorname{Norm}_{\operatorname{Aut}(\mathfrak{S})}(L_0)$ dans l'ensemble des solutions de l'équation en $\xi, \eta$ (26) qui s'écrit simplement
$$
\eta^{-1} \xi \in L_1 \leqno{(30)}
$$
relation qui [unclear: rend] que (26) (avec la forme précise de l'exposant $\nu$) n'a de sens (pas [unclear: plus] qu'on n'a que pour que $\mathfrak{S}_{0,3}^{+\wedge} \longrightarrow \operatorname{Sl}(2, \hat{\mathbf{Z}})$ se factorise par $\mathfrak{S}$ - ce qu'on ne suppose pas, pour nous simplifier la vie.

On va faire sur $\mathfrak{S}$, muni de $\rho, \sigma$, l'hypothèse plus précise, soit
$$
(\xi, \eta)
\begin{cases}
\xi, \eta \text{ conjugués à } \rho \text{ et } \sigma \\
\eta^{-1} \xi \in L_1, \text{ i.e. } \exists \nu' \in \hat{\mathbf{Z}} \text{ tel que } \xi = \eta \varepsilon_1^{\nu'}
\end{cases} \leqno{(32)}
$$
Pour $u \in \operatorname{Aut}(\mathfrak{S})$ [unclear: annulant] $L_0$, [unclear: les elements]
$\xi$ et $\eta$ [unclear: faisant la classe de conjugaison] de $\rho$
et celle de $\sigma$, i.e. tels que
$$
\xi = [u, \rho], \quad \eta = [u, \sigma]
$$
satisfont à la première condition de (32), et

\newpage

%% ===== Page 500 =====

et soit $u \in \operatorname{Aut}(\bar{\mathfrak{S}})$ un automorphisme de $\bar{\mathfrak{S}}$, et posons
$$
\begin{cases}
\xi = [u, \rho] = u \rho u^{-1} = u(\rho) \rho^{-1} & \text{i.e. } u(\rho) = \xi \rho \\
\eta = [u, \sigma] = u(\sigma) \sigma^{-1} & \text{i.e. } u(\sigma) = \eta \sigma
\end{cases} \marginpar{cf (24), (25)} \leqno (34)
$$
alors dire que $\xi \rho$ ($= u(\rho)$) est conjugué à $\rho$ signifie que $u$ fixe la classe de conj. de $\rho$, dire que $\eta \sigma$ ($=u(\sigma)$) est conjugué de $\sigma$ signifie que $u$ fixe la classe de conj. de $\sigma$. Enfin la condition $\eta^{-1} \xi \in L_1$ s'écrit
$$
\eta^{-1} \xi = (\sigma u(\sigma)^{-1}) (u(\rho) \rho^{-1}) = \sigma u(\sigma^{-1} \rho) \rho^{-1} = \sigma \left( u( \underbrace{\sigma^{-1} \rho}_{\varepsilon_0} ) \underbrace{\rho^{-1} \sigma}_{\varepsilon_0^{-1}} \right) \sigma^{-1}
$$
$$
= \sigma ( u(\varepsilon_0) \varepsilon_0^{-1} ) \sigma^{-1}
$$
donc de la forme\footnote{Dans notre [unclear: cas] où on utilise $\sigma \rho = \varepsilon_0 \in L_0$, on a plutôt [unclear: besoin de] $\sigma^2 = 1$, $\rho^3 = 1$.}
$$
\sigma ( u(\varepsilon_0) \varepsilon_0^{-1} ) \sigma^{-1} \in L_1 \quad \text{i.e.} \quad u(\varepsilon_0)\varepsilon_0^{-1} \in L_0
$$
ou encore (puisque $\varepsilon_0 \in L_0$) $u(\varepsilon_0) \in L_0$, soit (puisque $\varepsilon_0$ engendre $L_0$)
$$
u(L_0) = L_0
$$
i.e. que $u$ normalise $L_0$.

Soit donc
$$
\tilde{\mathcal{B}}'_0 \defeq \left\{ u \in \operatorname{Aut}(\bar{\mathfrak{S}}) \ \middle| \ \begin{array}{l} u \text{ fixe classes de conj. de } \rho \text{ et } \sigma \\ \text{et } u \text{ normalise } L_0 \end{array} \right\} \leqno(35)
$$
De là [unclear: par définition de $Z$ on a] une application [unclear: naturelle donnée] par
$$
\tilde{\mathcal{B}}'_0 \longrightarrow Z \leqno(36)
$$
$$
u \longmapsto ( [u, \rho] = u(\rho)\rho^{-1}, [u, \sigma] = u(\sigma)\sigma^{-1} )
$$

\newpage

%% ===== Page 501 =====

\underline{Hypothèse fondamentale sur $\mathfrak{S}$, muni de $\rho, \sigma$, d'où $\varepsilon_0 = \sigma^{-1} \rho$,}
$$
\begin{array}{l}
\varepsilon_1 = \rho \sigma^{-1} = \sigma(\varepsilon_0) = \rho(\varepsilon_0) : \\
a) \quad \mathfrak{S} \text{ est engendré topologiquement par } \rho, \sigma \\
b) \quad \text{L'application (36) (injective par } a)) \text{ est bijective.}
\end{array} \leqno{(37)}
$$

La condition b) signifie donc que pour $(\alpha, \beta) \in \mathfrak{S} \times \mathfrak{S}$, satisfaisant une relation
$$
\alpha \rho \alpha^{-1} \rho^{-1} = \beta \sigma \beta^{-1} \sigma^{-1} \varepsilon_1^{\nu'} \quad (\nu' \text{ convenable}),
$$
il existe $u \in \operatorname{Aut}(\mathfrak{S})$ tel qu'on ait
$$
\begin{cases}
[u, \rho] = \alpha \rho \alpha^{-1} \rho^{-1} & \text{i.e.} \quad u(\rho) = \operatorname{int}(\alpha)(\rho) \quad \text{(i.e. } u(\rho) = \xi \rho) \\
[u, \sigma] = \beta \sigma \beta^{-1} \sigma^{-1} & \text{i.e.} \quad u(\sigma) = \operatorname{int}(\beta)(\sigma) \quad \text{(i.e. } u(\sigma) = \eta \sigma)
\end{cases}
$$
(on aura donc nécessairement $u \in \widetilde{B}'_0$ !, et $u$ sera unique par a)).

Ainsi, sur l'ensemble $R_{\rho, \sigma}$ des $(\xi, \eta)$ satisfaisant a) b) c) , on aura une structure de groupe, déduite par transport de structure de celle de $\widetilde{B}'_0$ etc.

\newpage

%% ===== Page 502 =====

\marginpar{Les groupes $\mathcal{M}_{\rho,\sigma}$ etc.}
II Reprenons la situation de [unclear: base], en oubliant (au moins pour le moment) $\mathcal{G}$, $\varphi$, etc.

Soit $\mathfrak{S}$ un groupe\footnote{Pour une [unclear: justification] des notations suivantes cf p. 564, remarque.} --- on se placera indifféremment, soit dans la catégorie des groupes discrets, soit dans celle des groupes profinis, --- le choix n'a pas ici d'incidence. On supposera pour fixer les idées de se placer dans les groupes profinis. On se donne dans $\mathfrak{S}$
$$
\rho, \sigma \in \mathfrak{S} \leqno{(1)}
$$

\marginpar{[unclear: Lier en outre $\epsilon_0 = \rho_0^2, \sigma = \rho_0^3$ et $\rho^2 = \dots$ aux relations (2) de l'algèbre]}
$$
\epsilon_0 = \sigma^{-1} \rho , \quad \epsilon_1 = \rho \sigma^{-1} = \sigma(\epsilon_0) = \rho(\epsilon_0)
$$

$$
\begin{aligned}
L_0 &= \text{sous-groupe fermé engendré par } \epsilon_0 \\
L_1 &= \qquad \qquad \dots \qquad \qquad \qquad \qquad \epsilon_1
\end{aligned} \leqno{(3)}
$$
donc $L_1 = \rho(L_0) = \sigma(L_0)$.

On considère
$$
\operatorname{Aut}(\mathfrak{S}) \xrightarrow{\Phi = (\xi, \eta)} \mathfrak{S} \times \mathfrak{S} \leqno{(4)}
$$
défini par
$$
\begin{cases}
\xi(u) = [u, \rho] = u(\rho) \rho^{-1} & \text{donc } u(\rho) = \xi \rho \quad (\xi = \xi(u)) \\
\eta(u) = [u, \sigma] = u(\sigma) \sigma^{-1} & \text{donc } u(\sigma) = \eta \sigma \quad (\eta = \eta(u))
\end{cases} \leqno{(5)}
$$

Donc pour $u \in \operatorname{Aut}(\mathfrak{S})$, et $\xi = \xi(u), \eta = \eta(u)$, on a :
$$
\begin{cases}
u \text{ fixe la classe de conjugaison de } \rho \iff \xi \rho \sim \rho \iff \xi = \alpha \rho \alpha^{-1} \rho^{-1} \ (\alpha \in \mathfrak{S}) \\
u \text{ fixe la classe de conjugaison de } \sigma \iff \eta \sigma \sim \sigma \iff \eta = \beta \sigma \beta^{-1} \sigma^{-1} \ (\beta \in \mathfrak{S}) \\
u \text{ normalise } L_0 \iff \eta^{-1} \xi \in L_1 \iff \xi \in \eta L_1 \ \text{[unclear: (c.a.d. } \xi \equiv \eta \pmod{L_1}\text{)]}
\end{cases} \leqno{(6)}
$$

\newpage

%% ===== Page 503 =====

$$
B_{\rho, \sigma} \subset \operatorname{Aut}(G) \leqno{(7)}
$$
$$
= \{ u \mid \text{satisfaisant les trois conditions } (5) \}
$$
$$
\begin{aligned}
\Sigma_{\rho, \sigma} &\subset G \times G \\
&= \left\{ (\xi, \eta) \left| 
\begin{array}{l}
\xi = \alpha \rho \alpha^{-1} \rho^{-1} \text{ (i.e. } \xi \rho \sim \rho), \quad \eta = \beta \sigma \beta^{-1} \sigma^{-1} \text{ (i.e. } \eta \sigma \sim \sigma) \quad (\alpha, \beta \in G) \\
\eta^{-1} \xi \in L_1 \text{ i.e. } \xi = \eta \xi_1^\nu \quad (\nu \in \hat{\mathbf{Z}})
\end{array}
\right. \right\} ,
\end{aligned}
\leqno{(8)}
$$
de sorte que $B_{\rho, \sigma}$ est l'image inverse de $\Sigma_{\rho, \sigma}$ par (4), et on a l'application induite
$$
B_{\rho, \sigma} \xrightarrow{\varphi = \varphi_{\rho, \sigma}} \Sigma_{\rho, \sigma} \leqno{(9)}
$$
On suppose :\footnote{Attention : cette condition est satisfaite pour $G = \operatorname{Sl}(2, \hat{\mathbf{Z}})$ !! (et $\rho, \sigma$ les générateurs canoniques $S, T$\dots)}
$$
\begin{cases}
\text{a) } \rho, \sigma \text{ engendrent } G \text{ topologiquement} \\
\text{b) l'application (9) (injective par a)) est surjective (donc bijective).}
\end{cases} \leqno{(10)}
$$
On transporte alors par $\varphi$ la structure de groupe de $B_{\rho, \sigma}$ sur $\Sigma_{\rho, \sigma}$. On désigne par
$$
u_{\xi, \eta} = \varphi^{-1}(\xi, \eta) \in B_{\rho, \sigma} \quad ((\xi, \eta) \in \Sigma_{\rho, \sigma}) \leqno{(11)}
$$
de sorte que l'autom.\ $u_{\xi, \eta}$ est caractérisé par
$$
\begin{cases}
u_{\xi, \eta}(\rho) = \xi \rho \\
u_{\xi, \eta}(\sigma) = \eta \sigma
\end{cases} \leqno{(12)}
$$
L'hypothèse b) signifie donc que, dès que $(\xi, \eta) \in \Sigma_{\rho, \sigma}$, i.e.\ $\xi \rho \sim \rho$, $\eta \sigma \sim \sigma$ dans $G$ et $\eta^{-1} \xi \in L_1$, il existe

\newpage

%% ===== Page 504 =====

un automorphisme $u_{\xi,\eta}$ de $\widehat{\mathcal{G}}$ satisfaisant (12)
Cet $u_{\xi,\eta}$ sera unique par (10) a) ). La loi de
groupe sur $\Sigma_{\rho,\sigma}$ s'explicite ainsi :

$$
(\xi',\eta')(\xi,\eta) = (\xi'',\eta'') \iff \begin{cases} \xi'' = u_{\xi,\eta}(\xi') \xi \\ \eta'' = u_{\xi,\eta}(\eta') \eta \end{cases} \leqno{(13)}
$$

Considérons l'homomorphisme [unclear: \sout{donné par}]
$$ \widehat{\mathcal{G}} \to \operatorname{Aut}(\widehat{\mathcal{G}}) $$
$$ g \mapsto \operatorname{int}(g) $$
et l'image inverse de $B_{\rho,\sigma}$ \sout{dans} [unclear: \dots] le normalisateur de $L_0$ dans $\widehat{\mathcal{G}}$, soit $N_0$. On veut
construire \sout{dans} un groupe $G_{\rho,\sigma}$, contenant $\widehat{\mathcal{G}}$
et dans lequel s'envoie $B_{\rho,\sigma}$, de façon que
l'action tautologique de $B_{\rho,\sigma}$ sur $\widehat{\mathcal{G}}$ soit
induite par cet homomorphisme, et que
$G_{\rho,\sigma}$ soit engendré par $\widehat{\mathcal{G}}$ et l'image de $B_{\rho,\sigma}$.

\sout{soit $B'_{\rho,\sigma}$, qui contienne $N_0$.}

Si le centre de $\widehat{\mathcal{G}}$ est trivial i.e. $\widehat{\mathcal{G}} \hookrightarrow \operatorname{Aut}(\widehat{\mathcal{G}})$, \sout{la construction} [unclear: \dots]
($\widehat{\mathcal{G}} \subset \text{sous-groupe distingué}$)
[ $B_{\rho,\sigma} \cdot \widehat{\mathcal{G}}$ (produit semi-direct) il suffit de
considérer l'application \sout{pour le sous-groupe}
$N_0 \to B_{\rho,\sigma} \cdot \widehat{\mathcal{G}}$ \quad ([unclear: \sout{du sous-groupe \dots}])
$g \mapsto (\operatorname{int}(g), g^{-1})$ ]

Mais l'hypothèse
que le centre de $\widehat{\mathcal{G}}$ est trivial n'est pas heureuse -
par exemple elle n'est pas satisfaite si $\mathcal{G} = \operatorname{Sl}(2,\hat{\mathbf{Z}})/\pm 1$.
Il me semble que l'introduction de $N_0$ (sur
lequel on \sout{pratiq} pratiquera une certaine torsion)

\newpage

%% ===== Page 505 =====

n'était pas non plus judicieuse, et que la
condition $B_{\rho,\sigma} \supset N_0$ - il suffisait de $B_{\rho,\sigma} \supset L_0$.

Je vais donc considerer
$$ S_0 G_{\rho,\sigma} \overset{dfn}{=} B_{\rho,\sigma} . \underline{C} \quad \text{le produit 1/2 direct, et} $$
$$ (14) \quad \begin{cases} L_0 \overset{i}{\hookrightarrow} S_0 G_{\rho,\sigma} \\ g \longmapsto i(g).g^{-1} \end{cases} $$
l'image de cet homomorphisme est un sous-groupe
distingué, car a) il commute à $\underline{C}$ (trivial) et
b) il est invariant par $B_{\rho,\sigma}$, car pour $u \in B_{\rho,\sigma}$
$$ int(u)(i(g)) = int(u) (i(g).g^{-1}) = int(u)(i(g)) . \underbrace{int(u)(g^{-1})}_{=u(g)^{-1}} = i(u(g)) $$
(i.e. $i$ commute aux actions de $B_{\rho,\sigma}$ , agissant sur
$L_0$ par restriction à $L_0$ de l'op. de $B_{\rho,\sigma}$ sur
$\underline{C}$, et sur $S_0 G_{\rho,\sigma}$ par autom. intérieurs ).

On peut donc considerer
$$ (15) \quad G_{\rho,\sigma} \overset{dfn}{=} S_0 G_{\rho,\sigma} / i(L_0) $$

Alors les hom. d'inclusion
$$ B_{\rho,\sigma} \hookrightarrow S_0 G_{\rho,\sigma} , \quad \underline{C} \hookrightarrow S_0 G_{\rho,\sigma} $$
donnent, par composition avec $S_0 G_{\rho,\sigma} \twoheadrightarrow G_{\rho,\sigma}$,
des hom.
$$ (16) \quad \underline{C} \to G_{\rho,\sigma} , \quad B_{\rho,\sigma} \to G_{\rho,\sigma} $$
rendant commutatif le diagramme
(17)
\begin{tikzcd}
L_0 \ar[r, hook] \ar[d, hook] & \underline{C} \ar[d] \\
B_{\rho,\sigma} \ar[r] & G_{\rho,\sigma}
\end{tikzcd}

On note :
$$ (18) \quad \operatorname{Ker}(\underline{C} \to G_{\rho,\sigma}) = L_0 \cap \text{Centr}(\underline{C}) = \operatorname{Ker}(L_0 \to B_{\rho,\sigma}) $$

\newpage

%% ===== Page 506 =====

Dans le cas $\underline{G} = S \ell(2, \hat{\mathbb{Z}})'$, (où $\rho,\sigma$ comme à l'ordinaire
l'on a $L_0 \cap Cent(\underline{G}) = \{1\}$. Donc il
ne semble pas a priori qu'il y ait des inconvé-
nients majeurs à supposer que

$$ (18') \quad L_0 \cap Cent(\underline{G}) = \{1\} \text{ i.e. } \underline{G} \hookrightarrow G_{\rho,\sigma} \text{ inj.} $$

[MARGIN: i.e. l'un des 2 diagrammes exacts
(10bis)
\begin{tikzcd}
& & 1 \ar[d] & 1 \ar[d] \\
& & L_0 \ar[d] \ar[r, equal] & L_0 \ar[d] \\
1 \ar[r] & \underline{G} \ar[r] \ar[d, equal] & S G_{\rho,\sigma} \ar[r] \ar[d] & B_{\rho,\sigma} \ar[r] \ar[d] & 1 \\
1 \ar[r] & \underline{G} \ar[r] & G_{\rho,\sigma} \ar[r] \ar[d] & \Gamma_{\rho,\sigma} \ar[r] \ar[d] & 1 \\
& & 1 & 1
\end{tikzcd}
]

Notons que $L_0$ est l'image de
$$ L_0 \to B_{\rho,\sigma} $$
$$ \ell \mapsto i(\ell) $$
Composée de (14) et de la projection $S G_{\rho,\sigma} \to B_{\rho,\sigma}$)
et évidemment $L_0 \cap \underline{G} = 1$, et posant

$$ (19) \quad \Gamma_{\rho,\sigma} \stackrel{dfn}{=} B_{\rho,\sigma} / \text{Im } L_0 $$

on trouve

$$ (20) \quad G_{\rho,\sigma} / \underline{G} \simeq B_{\rho,\sigma} / L_0 \simeq \Gamma_{\rho,\sigma} $$

NB Dans le cas $\underline{G} = S \ell(2, \hat{\mathbb{Z}})'$, $B_{\rho,\sigma}$ est,
(à [unclear: $C P(\hat{\mathbb{Z}})$] près le groupe noté précédemment $\tilde{B}_0'$
ou plutôt un groupe isomorphe $\tilde{B}_0'$), et
l'on a $det \mid \tilde{B}_0' = B_{\rho,\sigma}$ isom. (trivial sur $L_0$)
induisant par passage au quotient un isomorphisme

$$ \Gamma_{\rho,\sigma} \buildrel \sim \over \to G_m $$

et on a donc, pour utiliser l'isomorphisme
$$ \tilde{B}_0' = \tilde{\Gamma}_0 \cdot L_0 \quad (\text{1/2 direct}) \quad \text{où } \tilde{\Gamma}_0 \to \tilde{B}_0' \simeq G_m $$
$$ \Gamma_{\rho,\sigma} = \tilde{\Gamma}_{\rho,\sigma} \cdot L_0 $$
pour donner

$$ (22) \quad G_{\rho,\sigma} = S \ell(2, \hat{\mathbb{Z}})' . $$

[MARGIN: dans le cas de $\underline{G} = \widehat{S \ell}(2, \mathbb{Z}), \bar{\rho}, \bar{\sigma}$ (21)
tout se transpose avec
$\underline{G} = S \ell(2, \hat{\mathbb{Z}})' \dots$]

\newpage

%% ===== Page 507 =====

Mais dans des cas plus généraux auxquels nous
aspirons, il ne faut pas dans H s'attendre à ce que
$\Gamma_{\rho, \sigma}$ soit aussi petit - p. ex. qu'il soit commutatif.
Par exemple dans le cas "universel", $\mathcal{G} = \mathfrak{S}_{0,3}^{+\wedge}$,
on aura
$$ B_{\rho, \sigma} = \hat{\mathcal{M}}_{0,3}^\sim[0]' = \operatorname{Ker}(\hat{\mathcal{M}}_{0,3}^\sim[0] \xrightarrow{\varepsilon_3} \{ \pm 1 \} ) \leqno{(23)} $$
donc
$$ \Gamma_{\rho, \sigma} \simeq \Gamma_{0,3}^{\sim 1} \quad (\supset \boldsymbol{\Gamma} \ ! ) \leqno{(24)} $$
Par contre, on doit pouvoir dans ce cas sans doute
(conjecture)
postuler l'existence d'un scindage de l'extension
$B_{\rho, \sigma}$ de $\Gamma_{\rho, \sigma}$ par $L_0$. En fait, si on suppose,
que $\rho^2 = \sigma^3 = 1$, i.e. $\mathcal{G}$
quotient de $\mathfrak{S}_{0,3}^{+\wedge}$, que le quotient $\mathcal{G}$ de
$\mathfrak{S}_{0,3}^{+\wedge}$ est sans quotient fini s'envoyant dans
$\subset \operatorname{Sl}(2, \hat{\mathbf{Z}})$, alors l'hom. $B_{\rho, \sigma} \to \mathfrak{S}_{\rho, \sigma}$ s'envoie dans
$\tilde{B}_{0,3}'$, et l'image inverse de $\tilde{\Gamma}_0$ est un
sous-groupe (qu'on pourrait noter $T_{\rho, \sigma}$)
$$ T_{\rho, \sigma} \subset B_{\rho, \sigma} \quad T_{\rho, \sigma} \isommap \Gamma_{\rho, \sigma} \leqno{(25)} $$
qui est un sous-groupe section, pour $B_{\rho, \sigma} \to \Gamma_{\rho, \sigma}$,
il en est de même (sous-groupe de $G_{\rho, \sigma}$)
pour $G_{\rho, \sigma} \to \Gamma_{\rho, \sigma}$, de sorte qu'on a
$$ B_{\rho, \sigma} \simeq T_{\rho, \sigma} \cdot L_0, \quad G_{\rho, \sigma} \simeq T_{\rho, \sigma} \cdot \mathfrak{S}_{\rho, \sigma} \quad (\text{produits } 1/2 \text{ directs}) \leqno{(26)} $$

\newpage

%% ===== Page 508 =====

Soit maintenant
$$
\mathcal{S}\mathcal{G}_{\rho,\sigma} \subset \mathfrak{S} \times \mathfrak{S} \leqno{(27)}
$$
$$
\defeq \left\{ (\alpha, \beta) \mid \alpha \rho (\alpha)^{-1} = \beta \sigma (\beta)^{-1} \varepsilon^\nu, \text{ pour } \nu \in \hat{\mathbf{Z}} \text{ convenable} \right\}
$$
(i.e.\ [unclear: $\xi \eta^{-1} \in L_1$])

On a une application canonique
$$
\mathcal{S}\mathcal{G}_{\rho,\sigma} \longrightarrow \Sigma_{\rho,\sigma} \leqno{(28)}
$$
$$
(\alpha, \beta) \longmapsto (\xi, \eta), \begin{cases} \xi = \alpha \rho \alpha^{-1} \\ \eta = \beta \sigma \beta^{-1} \end{cases}
$$
qui est surjective par définition de $\Sigma_{\rho,\sigma}$, et en fait par l'hypothèse (10) a une section canonique
$$
\Sigma_{\rho,\sigma} \longrightarrow \mathcal{S}\mathcal{G}_{\rho,\sigma}
$$
$$
(\xi, \eta) \longmapsto u_{\xi,\eta}
$$
à cela près que $u_{\xi,\eta}$ n'est pas dans $\mathfrak{S}$, mais seulement dans $G_{\rho,\sigma}$. Nous nous arrêterons dans l'ensemble plus gros
$$
\mathcal{G}_{\rho,\sigma} \subset G_{\rho,\sigma} \times G_{\rho,\sigma} \leqno{(29)}
$$
$$
\defeq \left\{ (\alpha, \beta) \mid \alpha \rho (\alpha)^{-1} = \beta \sigma (\beta)^{-1} \varepsilon^\nu \text{ pour } \nu \in \hat{\mathbf{Z}} \text{ convenable} \right\}
$$
\footnote{NB On a bien $\xi = \alpha (\rho) \alpha^{-1}$, $\eta = \beta (\sigma) \beta^{-1}$ d'où $\xi, \eta \in \mathfrak{S}$ (puisque $\mathfrak{S}$ inv. dans $G_{\rho,\sigma}$).}
de sorte que l'on a
$$
\mathcal{S}\mathcal{G}_{\rho,\sigma} = \mathcal{G}_{\rho,\sigma} \cap (\mathfrak{S} \times \mathfrak{S}) \leqno{(30)}
$$
(sans supposer (18) pour $\varepsilon$ ou $L_1$).

\newpage

%% ===== Page 509 =====

Ainsi on trouve encore
$$
\begin{cases}
\mathcal{G}_{\rho,\sigma} \longrightarrow \Sigma_{\rho,\sigma} \\
(\alpha,\beta) \longmapsto (\xi,\eta), \quad \xi = \alpha\rho\alpha^{-1} = [\alpha,\rho] = \alpha(\rho)\rho^{-1} \\
\phantom{(\alpha,\beta) \longmapsto (\xi,\eta), \quad} \eta = \beta\sigma\beta^{-1} = [\beta,\sigma] = \beta(\sigma)\sigma^{-1}
\end{cases} \leqno (31)
$$
(Par définition de $B_{\rho,\sigma}$ on trouve bien que les $(\xi,\eta)$ associés à $(\alpha,\beta)$ sont dans $\Sigma_{\rho,\sigma}$). Cette fois-ci l'hyp. (10) implique qu'on a une section canonique
$$
\begin{cases}
\Sigma_{\rho,\sigma} \longrightarrow \mathcal{G}_{\rho,\sigma} \\
(\xi,\eta) \longmapsto (u_{\xi,\eta}, u_{\xi,\eta})
\end{cases} \leqno (32)
$$
et on trouve [unclear: le scindage de] (10).

Théorème Considérons les sous-groupes de $\mathcal{G}$\footnote{NB dans le cas $\mathcal{G} = \operatorname{Sl}(2, \hat{\mathbf{Z}})$ et [unclear: $\dots$], on a $Z_{\rho} = T_{\rho}$, $Z_{\sigma} = N_{\sigma}'$}
$$
\begin{cases}
Z_{\rho} = \operatorname{Centr}_{\mathcal{G}_{\rho,\sigma}}(\rho) \\
Z_{\sigma} = \operatorname{Centr}_{\mathcal{G}_{\rho,\sigma}}(\sigma)
\end{cases} \leqno (33)
$$
et le groupe produit
$$
B_{\rho,\sigma} \times Z_{\rho} \times Z_{\sigma}.
$$
Alors l'application naturelle est bijective.
$$
B_{\rho,\sigma} \times Z_{\rho} \times Z_{\sigma} \isommap \mathcal{G}_{\rho,\sigma} \leqno (34)
$$
$$
(u, a, b) \longmapsto (u a^{-1}, u b^{-1})
$$

Corollaire\footnote{Par définition de $B_{\rho,\sigma}$ on a [unclear: un] hom. $Z_{\rho} \to B_{\rho,\sigma} \times Z_{\sigma}$ donné par les équations (35) qui [unclear: donne] l'iso en question.} Considérons les hom. de groupes
$$
\delta_{B} : B_{\rho,\sigma} \to \mathcal{G}_{\rho,\sigma}, \quad \delta_{\rho} : Z_{\rho} \to \mathcal{G}_{\rho,\sigma}, \quad \delta_{\sigma} : Z_{\sigma} \to \mathcal{G}_{\rho,\sigma} \leqno (35)
$$
où $\delta_{B}$ est l'hom. canonique de $B_{\rho,\sigma}$ dans $\mathcal{G}_{\rho,\sigma} = \mathcal{G}$, et $\delta_{\rho}, \delta_{\sigma}$ sont les composés de $\delta_{B}$ avec les inclusions $Z_{\rho} \subset \mathcal{G}_{\rho,\sigma}$, $Z_{\sigma} \subset \mathcal{G}_{\rho,\sigma}$.

\newpage

%% ===== Page 510 =====

\marginpar{[unclear: notations de la § 48, IX H suit]}
$$
S_{\rho,\sigma}^* \defeq \{ (u, a, b) \in B_{\rho,\sigma} \times Z_\rho \times Z_\sigma \mid \delta_{\rho,\sigma}(u) = \delta_\rho(a) = \delta_\sigma(b) \} \leqno{(36)}
$$
i.e.\ $S_{\rho,\sigma}^* = B_{\rho,\sigma} \times_{\overline{\Gamma}_{\rho,\sigma}} Z_\rho \times_{\overline{\Gamma}_{\rho,\sigma}} Z_\sigma$.

une bijection\footnote{NB la lettre $S$ ici est inspirée de l'usage des \emph{systèmes} dans (38) un peu plus bas où l'on se propose de munir d'une structure de \emph{groupe} (par l'intermédiaire de bijections \dots)}
$$
S_{\rho,\sigma}^* \isommap S\mathcal{G}_{\rho,\sigma} \quad \text{(défini dans (27))} \leqno{(37)}
$$
$$
(u, a, b) \longmapsto (\xi, \eta) , \quad \text{où } \xi = u a^{-1} , \, \eta = u b^{-1}
$$

\section*{Scholie}

Par les bijections (34) (38), on transporte sur $\mathcal{G}_{\rho,\sigma}$, $S\mathcal{G}_{\rho,\sigma}$ les structures de groupes naturelles de $\Gamma_{\rho,\sigma}$, $S_{\rho,\sigma}^*$. Expliciter cette structure de groupe --- on trouve notamment pour $S\mathcal{G}_{\rho,\sigma}$ :

$$
\begin{cases}
(\alpha', \beta')(\alpha, \beta) = (\alpha'', \beta'') , \quad \text{avec } \alpha'' = u'(\alpha)\alpha' , \quad \beta'' = u'(\beta)\beta' \\
u' = u_{\xi', \eta'} : \mathcal{C} \to \mathcal{C}, \text{défini par } \begin{cases} u'(\rho) = \xi' \rho = \alpha' \rho {\alpha'}^{-1} \\ u'(\sigma) = \eta' \sigma = \beta' \sigma {\beta'}^{-1} \end{cases}
\end{cases} \leqno{(39)}
$$

[unclear: Dans $S\Gamma_{\rho,\sigma}^*$ on a le sous-groupe $L_0 \to S\Gamma_{\rho,\sigma}^*$ et la suite exacte
$$
1 \to L_0 \to S\Gamma_{\rho,\sigma}^* \to \overline{\Gamma}_{\rho,\sigma} \to 1 \leqno{(40)}
$$
$$
\overline{\Gamma}_{\rho,\sigma} = Z_\rho \times_{\overline{\Gamma}_{\rho,\sigma}} Z_\sigma \leqno{(41)}
$$
$$
1 \to S Z_\rho \times S Z_\sigma \to S\Gamma_{\rho,\sigma}^* \to \overline{\Gamma}_{\rho,\sigma} \to 1 \leqno{(42)}
$$
\emph{(Note : La fin de la page, à partir de « Dans $S\Gamma_{\rho,\sigma}^*$ », est barrée d'un grand trait diagonal, incluant un bloc à gauche concernant $M_{\rho,\sigma} \subset \operatorname{Aut}(\mathcal{C})$.)}]

\newpage

%% ===== Page 511 =====

$$
L_0 \longrightarrow S\mathbf{T}_{\rho,\sigma}^0 \subset B_{\rho,\sigma} \times Z_{\bar{\rho}} \times Z_{\bar{\sigma}} \leqno{(40)}
$$
$$
\ell \longmapsto (\ell, 1, 1)
$$
et l'action de $L_0$ par translation à g. sur $S\mathbf{T}_{\rho,\sigma}^0$ devient, par (39), l'action sur $S\mathcal{G}_{\rho,\sigma}$ donnée par
$$
(\alpha,\beta) \longmapsto (\ell\alpha, \ell\beta) \leqno{(41)}
$$
Posons\footnote{intro $\mathbf{T}_{\rho,\sigma}^0$ ds § 48, IX H}
$$
\mathbf{T}_{\rho,\sigma}^0 \defeq S\mathbf{T}_{\rho,\sigma}^0 / \Im L_0 \simeq Z_{\bar{\rho}} \times_{\Gamma_{\rho,\sigma}} Z_{\bar{\sigma}} \leqno{(42)}
$$
on trouve que (38) induit une bijection canonique
$$
\mathbf{T}_{\rho,\sigma}^0 \isommap L_0 \setminus S\mathcal{G}_{\rho,\sigma} \leqno{(43)}
$$
Posant\footnote{introduisons [unclear: aussi] $S\mathbf{T}_{\rho,\sigma} = \operatorname{Ker}(\mathbf{T}_{\rho,\sigma} \to \Gamma_{\rho,\sigma})$, $S\mathcal{G}_{\rho,\sigma} = \operatorname{Ker}(\mathcal{G}_{\rho,\sigma} \to \Gamma_{\rho,\sigma}) = \mathcal{G}_{\rho,\sigma}^0$}
$$
\begin{cases}
S Z_{\bar{\rho}} = \operatorname{Ker}(Z_{\bar{\rho}} \to \Gamma_{\rho,\sigma}) = \operatorname{Centr}_{Z_{\bar{\rho}}}(\rho) \\
S Z_{\bar{\sigma}} = \operatorname{Ker}(Z_{\bar{\sigma}} \to \Gamma_{\rho,\sigma}) = \operatorname{Centr}_{Z_{\bar{\sigma}}}(\sigma)
\end{cases} \leqno{(44)}
$$
on a des suites exactes
$$
\begin{cases}
1 \to S Z_{\bar{\rho}} \to Z_{\bar{\rho}} \to \Gamma_{\rho,\sigma} \to 1 \\
1 \to S Z_{\bar{\sigma}} \to Z_{\bar{\sigma}} \to \Gamma_{\rho,\sigma} \to 1
\end{cases} \leqno{(45)}
$$
et les expressions (37), (42) de $S\mathbf{T}_{\rho,\sigma}^0$, $\mathbf{T}_{\rho,\sigma}^0$ donnent
$$
1 \to L_0 \times S Z_{\bar{\rho}} \times S Z_{\bar{\sigma}} \to S\mathbf{T}_{\rho,\sigma}^0 \to \Gamma_{\rho,\sigma} \to 1 \leqno{(46)}
$$
$$
1 \to S Z_{\bar{\rho}} \times S Z_{\bar{\sigma}} \to \mathbf{T}_{\rho,\sigma}^0 \to \Gamma_{\rho,\sigma} \to 1 \leqno{(47)}
$$

NB Dans le cas $\mathcal{C} = \operatorname{Sl}(2, \hat{\mathbf{Z}})$, $\rho$ et $\sigma$ "[unclear: modulaires]", on retrouve les $S\mathbf{T}_{\rho,\sigma}^0$ et $\mathbf{T}_{\rho,\sigma}^0$ de § 48 \underline{IX}.

\newpage

%% ===== Page 512 =====

Soit maintenant

$$
S\mathcal{M}_{\rho,\sigma}^0 \subset S\mathfrak{T}_{\rho,\sigma}^0 \times G_{\rho,\sigma} \quad (\subset B_{\rho,\sigma} \times Z_\rho \times Z_\sigma \times G_{\rho,\sigma}) \leqno{(49)}
$$
$$
\defeq \{ (u,a,b,U) \mid \delta_B(u) = \delta_\rho(a) = \delta_\sigma(b) = \delta_G(U) \}
$$
(relations dans $\Gamma_{\rho,\sigma}$ cf (19))

$$
\mathcal{M}_{\rho,\sigma}^0 \subset \mathfrak{T}_{\rho,\sigma}^0 \times G_{\rho,\sigma} \leqno{(50)}
$$
$$
\defeq \{ (a,b,U) \mid \delta_\rho(a) = \delta_\sigma(b) = \delta_G(U) \}
$$
(relations dans $\Gamma_{\rho,\sigma}$)

On a donc un homomorphisme canonique
$$
\begin{array}{ccc}
S\mathcal{M}_{\rho,\sigma}^0 & \longrightarrow & \mathcal{M}_{\rho,\sigma}^0 \\
(u,a,b,U) & \longmapsto & (a,b,U)
\end{array} \leqno{(51)}
$$
dont le noyau est formé des $(u,1,1,1)$ où $u \in B_{\rho,\sigma}$ tel que $\delta_B(u) = 1$ i.e.\ $u \in L_0$, d'où une suite exacte

$$
1 \longrightarrow L_0 \longrightarrow S\mathcal{M}_{\rho,\sigma}^0 \longrightarrow \mathcal{M}_{\rho,\sigma}^0 \longrightarrow 1 \leqno{(52)}
$$
et des homomorphismes canoniques

$$
\begin{tikzcd}
& G_{\rho,\sigma} \\
\mathcal{M}_{\rho,\sigma}^0 \arrow[ru, "\Theta_G"] \arrow[r, "\Theta_\rho"] \arrow[rd, "\Theta_\sigma"'] & Z_\rho \\
& Z_\sigma
\end{tikzcd} \leqno{(53)}
$$
(comp.\ $(a,b,U) \mapsto \begin{cases} U \\ a \\ b \end{cases}$)

reliant les homomorphismes analogues sur $S\mathcal{M}_{\rho,\sigma}^0$ par composition avec (51), et de plus on a un homomorphisme canonique
$$
\begin{array}{ccc}
S\mathcal{M}_{\rho,\sigma}^0 & \overset{\Theta_B}{\longrightarrow} & B_{\rho,\sigma} \\
(u,a,b,U) & \longmapsto & u
\end{array} \leqno{(54)}
$$

\newpage

%% ===== Page 513 =====

On a ainsi l'isomorphisme
$$
S\mathcal{M}_{\rho,\sigma}^0 \simeq \mathcal{M}_{\rho,\sigma}^0 \times_{T_{\rho,\sigma}} B_{\rho,\sigma} ; \leqno{(55)}
$$
et $S\mathcal{M}_{\rho,\sigma}^0$ s'exprime à l'aide de $\mathcal{M}_{\rho,\sigma}^0$ et $B_{\rho,\sigma}$,
est un torseur qui s'obtient de $\mathcal{M}_{\rho,\sigma}^0$ par $L_0$
comme tirage inverse de l'extension $B_{\rho,\sigma}$ de
$T_{\rho,\sigma}$ par $L_0$.
De même on a
$$
\mathcal{M}_{\rho,\sigma}^0 \simeq Z_\rho \times_{T_{\rho,\sigma}} Z_\sigma \times_{T_{\rho,\sigma}} G_{\rho,\sigma} \leqno{(56)}
$$
d'où on déduit [unclear: (via $\delta_\rho, \delta_\sigma$ sur $T_{\rho,\sigma}$)] :
$$
1 \to \underset{\defeq \mathcal{M}_{\rho,\sigma}^{0+}}{SZ_\rho \times SZ_\sigma \times \mathcal{G}} \to \mathcal{M}_{\rho,\sigma}^0 \overset{\delta_{\mathcal{M}}^\ast}{\to} T_{\rho,\sigma} \to 1 \leqno{(57)}
$$
où on a posé (cf (44))
$$
\begin{cases} 
SZ_\rho = Z_\rho \cap \mathcal{G} = \operatorname{Centr}_{\mathcal{G}}(\rho) \\ 
SZ_\sigma = Z_\sigma \cap \mathcal{G} = \operatorname{Centr}_{\mathcal{G}}(\sigma) 
\end{cases}
$$
De même on a
$$
S\mathcal{M}_{\rho,\sigma}^0 \simeq B_{\rho,\sigma} \times_{T_{\rho,\sigma}} Z_\rho \times_{T_{\rho,\sigma}} Z_\sigma \times_{T_{\rho,\sigma}} G_{\rho,\sigma} \leqno{(49)}
$$
d'où une suite exacte canonique :
$$
1 \to L_0 \times SZ_\rho \times SZ_\sigma \times \mathcal{G} \to S\mathcal{M}_{\rho,\sigma}^0 \overset{\delta_{S\mathcal{M}^0}^\ast}{\to} T_{\rho,\sigma} \to 1 \leqno{(50)}
$$
Considérons le hom. can.
$$
\mathcal{G} \hookrightarrow \mathcal{M}_{\rho,\sigma}^0 \leqno{(51)}
$$
$$
u \mapsto (1, 1, u)
$$

\newpage

%% ===== Page 514 =====

dont l'image est le troisième facteur dans le noyau de (57), et on a une suite exacte canonique (déduite de (56) et (57))
\marginpar{(53') $1 \to S Z_\rho \times S Z_\sigma \to \mathcal{M}_{\rho,\sigma}^\circ \to \text{[unclear: G]} \to 1$}
$$ 
1 \to \underbrace{\mathfrak{S}}_{\text{parfois noté } \mathfrak{S}_{\rho,\sigma} ?} \to \mathcal{M}_{\rho,\sigma}^\circ \to \Gamma_{\rho,\sigma}^\circ \to 1 \leqno{(52)} 
$$
et de même (via (49) ou (50))
$$ 
1 \to \mathfrak{S} \to S_0 \mathcal{M}_{\rho,\sigma}^\circ \to S_0 \Gamma_{\rho,\sigma}^\circ \to 1 \leqno{(53)} 
$$
Considérons l'application
$$
S_0 \mathcal{M}_{\rho,\sigma}^\circ \longrightarrow \mathfrak{S} \times \mathfrak{S} \times \mathfrak{S}
$$
$$ 
\{u, a, b, U\} \longmapsto (u a^{-1}, u b^{-1}, u U^{-1}) = (\alpha, \beta, f) \leqno{(54)} 
$$
On voit par le corollaire p.\ 557 (cf.\ (37)) que cela induit une bijection
$$ 
S_0 \mathcal{M}_{\rho,\sigma}^\circ \simeq (S \mathcal{G}_{\rho,\sigma}) \times \mathfrak{S} \leqno{(55)} 
$$
On a ainsi pour un $x \in S_0 \mathcal{M}_{\rho,\sigma}^\circ$, d'une part les quantités
$$
\begin{matrix}
u, & a, & b, & U \\
\in & \in & \in & \in \\
B_{\rho,\sigma} & Z_\rho & Z_\sigma & G_{\rho,\sigma}
\end{matrix} \leqno{(56)}
$$
associés, d'autre part les éléments
$$ 
\underbrace{\alpha, \beta, f}_{\substack{\text{définis en termes} \\ \text{de (56), dans (54)}}}, \xi, \eta, \alpha_0, \beta_0, f_0, g, h \quad \text{dans } \mathfrak{S} \leqno{(57)} 
$$
et on a, posant
$$
\begin{cases}
\xi = a \rho(a)^{-1}, \, \eta = b \sigma(b)^{-1}, & \alpha_0 = f^{-1} \alpha f, \, \beta_0 = f^{-1} \beta f \\
g = \alpha_0 \rho(\alpha_0)^{-1} = f^{-1} \xi \rho(f), & h = \beta_0 \sigma(\beta_0)^{-1} = f^{-1} \eta \sigma(f)
\end{cases} \leqno{(58)}
$$

\newpage

%% ===== Page 515 =====

On aura, en termes de $\alpha, \beta, f$, les expressions suivantes de l'action de $U$ sur $\hat{C}$ :
$$
\begin{cases}
U(\rho) = \operatorname{int}(f^{-1}\alpha)(\rho) = \operatorname{int}(f^{-1}) (\xi \rho) \\
U(\sigma) = \operatorname{int}(f^{-1}\beta)(\sigma) = \operatorname{int}(f^{-1}) (\eta \sigma)
\end{cases} \leqno (59)
$$

Remarques. Finalement, dans tous ces développements, on n'a pas vraiment utilisé la soi-disante "hyp. fondamentale" (20), sauf la partie a). En l'absence de b), il y a lieu d'introduire
$$ \Sigma_{\rho,\sigma}^{\text{eff}} \subset \Sigma_{\rho,\sigma} \subset \hat{C} \times \hat{C} \leqno (60) $$
$$ \Sigma_{\rho,\sigma}^{\text{eff}} = \{ (\xi,\eta) \in \hat{C} \times \hat{C} \mid \exists u \in B_{\rho,\sigma} , u(\rho) = \xi \rho , u(\sigma) = \eta \sigma \} $$
qui n'est autre que l'image de l'[unclear: application] (1) : couples $(\xi,\eta)$ satisfaisant les conditions définissant $\Sigma_{\rho,\sigma}$ ((28)).
L'application $S\mathcal{G}_{\rho,\sigma} \xrightarrow{(28)} \Sigma_{\rho,\sigma}$ [unclear: induit une surjection] sur son image $\Sigma_{\rho,\sigma}^{\text{eff}} \simeq B_{\rho,\sigma}$.
D'ailleurs il y a [unclear: lieu de noter]
$$ S\mathcal{G}_{\rho,\sigma}^{\text{eff}} = \text{Image inv. de } \Sigma_{\rho,\sigma}^{\text{eff}} \text{ dans } S\mathcal{G}_{\rho,\sigma} \to \Sigma_{\rho,\sigma} \text{ (28)} \leqno (61) $$
[unclear: bien qu'isom.], et ce sera ce $S\mathcal{G}_{\rho,\sigma}^{\text{eff}}$ qu'on [unclear: sera amené à substituer à] $S\mathcal{G}_{\rho,\sigma}$.

\newpage

%% ===== Page 516 =====

de groupes isomorphes : $S_0 \mathcal{G}_{\rho,\sigma}$, par l'isomorphisme (qui remplace (38))
$$ S_0 \mathcal{G}_{\rho,\sigma} \isommap S_0 \mathcal{G}_{\rho,\sigma}^{\text{eff}} \leqno{(62)} $$
$$ (u, \underline{a}, \underline{b}) \longmapsto (u \underline{a}^{-1}, u \underline{b}^{-1}) $$

Ceci posé, revenons à $\mathcal{M}_{0,3}^\sim$, et la suite exacte
$$ 1 \to \mathfrak{S}_{0,3}^{+\wedge} \to \mathcal{M}_{0,3}^\sim \to \underset{\simeq \Gamma_{\mathbf{Q}}}{\Gamma_{0,3}^\sim} \to 1 \leqno{(63)} $$

Considérons le sous-groupe
$$ \Gamma_{\mathbf{Q}}^{\prime \sim} = \operatorname{Ker}( \Gamma_{0,3}^\sim \xrightarrow{(43)} \hat{\mathbf{Z}}^\times ) = \{ \gamma \in \Gamma_{0,3}^\sim \mid \mu(\gamma) = 1 \} \leqno{(64)} $$

et soit
$$ \mathcal{M}_{0,3}^{\prime \sim} = \operatorname{Im. inv.} \text{ de } \Gamma_{0,3}^{\prime \sim} \text{ dans } \mathcal{M}_{0,3}^\sim \leqno{(65)} $$

Soit $G$ un groupe profini, et
$$ \psi : \mathfrak{S}_{0,3}^{+\wedge} \to G \leqno{(66)} $$
un homomorphisme continu surjectif, ce qui revient à la donnée de deux éléments
$$ \rho, \sigma \in G \leqno{(67)} $$
(les images de $\rho, \sigma \in \mathfrak{S}_{0,3}^+$) satisfaisant
$$ \rho^3 = \sigma^2 = 1 \leqno{(68)} $$

\newpage

%% ===== Page 517 =====

[MARGIN: l'introduire
ces conditions
suivantes équiv. à
a) $M_{0,3}^\sim [0] \to \Sigma_{\rho,\sigma}^{eff}$
b) $\forall u \in M_{0,3}^\sim, \exists y$
tel que $uy \in M_{0,3}^\sim [0]$
b') Idem avec $u\Gamma = y\Gamma$
c) $\exists \Psi_u$ qui prolonge $\psi$
c') $\exists \Psi_G$ qui prolonge $u$
On dira dans
la suite (?)
que $\psi$ est effective
si (70) a) satisfait]

C'est-à-dire que $\psi$ est effectif si pour
tout $u \in M_{0,3}^\sim$, posant $u(\rho) = \xi \rho$, $u(\sigma) = \eta \sigma$,

$$(\xi_{\tilde{\mathbb{G}}}, \eta_{\tilde{\mathbb{G}}}) = (u(\xi), u(\eta)) \in \Sigma_{\rho,\sigma}$$

est effectif (i.e. $\in \Sigma_{\rho,\sigma}^{eff}$, i.e. $\exists\ u_{\tilde{\mathbb{G}}}$ autom.
de $\mathbb{G}$ tel qu'on ait

(69) $$u_{\tilde{\mathbb{G}}}(\rho_{\tilde{\mathbb{G}}}) = \xi_{\tilde{\mathbb{G}}} \rho_{\tilde{\mathbb{G}}} \quad , \quad u_{\tilde{\mathbb{G}}}(\sigma_{\tilde{\mathbb{G}}}) = \eta_{\tilde{\mathbb{G}}} \sigma_{\tilde{\mathbb{G}}}$$

ce qui revient encore, par commutativité dans

(70)
\begin{tikzcd}
\mathbb{G}_{0,3}^{+\wedge} \ar[r, "u"] \ar[d, "\psi"'] & \mathbb{G}_{0,3}^{+\wedge} \ar[d, "\psi"] \\
\mathbb{G} \ar[r, "u_{\tilde{\mathbb{G}}}"] & \mathbb{G}
\end{tikzcd}

(Dans le cas général où on ne l'a pas, le $u \in M_{0,3}^\sim$ effectif signifie que cet $u$ transforme
$\operatorname{Ker}\ \psi$ en lui-même, ou encore que
$M_{0,3}^\sim [0] \quad (= M_{0,3}^\sim [0] \cdot \mathbb{G}_{0,3}^{+\wedge})$ tout entier
stabilise $\operatorname{Ker}\ \psi$. Donc cette condition
est automatiquement remplie
quand $\mathbb{G}, \rho_{\tilde{\mathbb{G}}}, \sigma_{\tilde{\mathbb{G}}}$ satisfait la condition
b) de (20).

On a fait ce qu'il fallait pour avoir...

\newpage

%% ===== Page 518 =====

562

Théorème Soit $\psi : \overline{G}_{0,3}^{+ \wedge} \to \overline{G}$ un
hom surjectif effectif (cf ci-dessus).
Alors il existe un hom. unique

(71) $$ \Psi_\psi : M_{0,3}^{\prime \sim} \to M_\psi $$

[MARGIN: En fait, je crains bien - on aura $S_0 M_{0,3}^\sim \to S_0 M_\psi$]

[unclear: crossed out diagram]
(72)
\begin{tikzcd}
\overline{G}_{0,3}^{+ \wedge} \ar[r, "*"] \ar[d, "\psi"'] & M_{0,3}^{\prime \sim} \ar[d, "\Psi_\psi"] \\
\overline{G} \ar[r, "*"] & M_{\bar{\rho}, \bar{\sigma}} \quad (= M_\psi)
\end{tikzcd}

$\Psi_\psi$ a la propriété suivante: Si 
$u \in M_{0,3}^{\prime \sim}$ est de la forme $u = \alpha \beta \gamma$
($\alpha, \beta, \gamma \in \overline{G}_{0,3}^{+ \wedge}$, satisfaisant aux équations
des lacets [unclear: appelés = polygonaux] ), alors
$\Psi_\psi(u)$ est l'élément de $M_\psi$
défini par les invariants $\psi(\alpha), \psi(\beta), \psi(\gamma)$.

Corollaire On a un diag. de suites exactes

(72)
\begin{tikzcd}
1 \ar[r] & \overline{G}_{0,3}^{+ \wedge} \ar[r] \ar[d, "\psi"'] & M_{0,3}^{\prime \sim} \ar[r] \ar[d, "\Psi_\psi"'] & \Gamma_{0,3}^{\prime \sim} \ar[r] \ar[d, "\overline{\psi}"'] & 1 \\
1 \ar[r] & \overline{G} \ar[r] & M_\psi \ar[r] & \Gamma_\psi \ar[r] & 1
\end{tikzcd}

L'hom. $\Psi_\psi$ s'explicite à l'aide
de trois homomorphismes

517

\newpage

%% ===== Page 519 =====

(73)
$$ \psi_{T_\rho} : \Gamma_Q^{\prime \sim} \to T_\rho $$
$$ \psi_{T_\sigma} : \Gamma_Q^{\prime \sim} \to T_\sigma $$

[MARGIN: NB on écrit pour simplifier $\rho, \sigma$ au lieu de $\rho_{\overline{Q}}, \sigma_{\overline{Q}}$ en indices]

(74)
$$ \psi_G : M_{0,3}^{\prime \sim} \to G_{\rho,\sigma} $$

satisfaisant les conditions : pour $x \in M_{0,3}^{\prime \sim}$, $\dot{x}$ son image dans $\Gamma_Q^{\prime \sim}$,
~~$a = \psi_{T_\rho}(\dot{x}), b = \psi_{T_\sigma}(\dot{x}), U = \psi_G(x)$~~

(75)
$$ a = \psi_{T_\rho}(\dot{x}), \quad b = \psi_{T_\sigma}(\dot{x}), \quad U = \psi_G(x) \quad , $$

(76)
$$ \delta_\rho(a) = \delta_\sigma(b) = \delta_G(U) \quad \text{dans} \quad \Phi_{\rho,\sigma} $$

i.e. les trois homs. composés

(77)
\begin{tikzcd}
& \Gamma_Q^{\prime \sim} \ar[r, "\psi_{T_\rho}"] & T_\rho \ar[dr, "\delta_\rho"] \\
M_{0,3}^{\prime \sim} \ar[ur, "can"] \ar[r, "can"] \ar[dr, "\psi_G"'] & \Gamma_Q^{\prime \sim} \ar[r, "\psi_{T_\sigma}"] & T_\sigma \ar[r, "\delta_\sigma"] & \Phi_{\rho,\sigma} \\
& G_{\rho,\sigma} \ar[urr, "\delta_G"']
\end{tikzcd}

sont égaux. De plus, on aura d'après (73)

(78)
$$ \psi_G | \mathbb{G}_{0,3}^{+ \wedge} = \Psi_{\mathfrak{S}} \quad , $$

et enfin ceci :
(79) Si $x \in M_{0,3}^{\prime \sim}$, $U = \psi_G(x)$ , on a 
commutativité dans

(80)
\begin{tikzcd}
\mathbb{G}_{0,3}^{+ \wedge} \ar[r, "\psi"] \ar[d, "int(x) | \mathbb{G}_{0,3}^{+ \wedge}"'] & \mathfrak{S} \ar[d, "int(U) | \mathfrak{S}"] \\
\mathbb{G}_{0,3}^{+ \wedge} \ar[r, "\psi"'] & \mathfrak{S}
\end{tikzcd}

\newpage

%% ===== Page 520 =====

563

On peut reprendre tout ceci comme
énoncé de fonctorialité des constructions
faites, p. ex. si un hom. surjectif

[MARGIN: 
il aurait
mieux valu
prendre
$\underline{G}' \xrightarrow{\varphi} \underline{G}$
pour
spécialiser
$\underline{G}$
]

(81) $$ \underline{G} \xrightarrow{\varphi = \varphi_{\underline{G}}} \underline{G}' $$

transformant le couple $(\rho, \sigma)$ de générateurs
topologiques de $\underline{G}$ en un couple $(\rho', \sigma')$.
Il y a lors un unique homomorphisme

(82) $$ S_0\mathcal{M}_{\underline{G}} \xrightarrow{\varphi_{S_0\mathcal{M}}} S_0\mathcal{M}_{\underline{G}'} $$

satisfaisant les conditions suivantes :

[MARGIN:
ceci même sans
supposer que
$\varphi_{S_0 \mathcal{M}}$ se déduise
d'un hom.
$\varphi_{\mathcal{M}}$,
pourvu qu'on
se donne les
$\varphi_B, \varphi_\rho, \varphi_\sigma, \varphi_G$
]

a) qui [unclear] sur les quotients [unclear] les hom. :
(83)
$$ \varphi_B : B_{\rho, \sigma} \to B_{\rho', \sigma'} $$
$$ \varphi_\rho : Z \to Z $$
$$ \varphi_\sigma : Z \to Z $$
$$ \varphi_G : G_{\rho, \sigma} \to G_{\rho', \sigma'} ; $$

b) Compatibilité avec $\varphi_G$ (ce qui montre que la connais-
sance de $\varphi_{S_0 \mathcal{M}}$ en revient à celle de $\varphi_G$) ;

c) On a commutativité dans

(85)
\begin{tikzcd}
\underline{G} \ar[r] \ar[d, "\varphi = \varphi_{\underline{G}}"'] & G_{\rho, \sigma} \ar[d, "\varphi_G"] \\
\underline{G}' \ar[r] & G_{\rho', \sigma'}
\end{tikzcd}

(ce qui montre, puisque $G_{\rho, \sigma} = B_{\rho, \sigma} \cdot \underline{G}$, que
$\varphi_G$ est [unclear] quand on connait $\varphi_B$ ;)

d) Pour $u \in B_{\rho, \sigma}$, désignant par $u' = \varphi_B(u)$, on a
commutativité dans
(86)
\begin{tikzcd}
\underline{G} \ar[r, "u"] \ar[d, "\varphi"'] & \underline{G} \ar[d, "\varphi"] \\
\underline{G}' \ar[r, "u'"] & \underline{G}'
\end{tikzcd}

519

\newpage

%% ===== Page 521 =====

-(ce qui détermine $\varphi_B$ en termes de $\varphi_{\mathfrak{G}}$,
donc cela détermine $\varphi_{\mathfrak{S}M}$, et par passage
aux quotients, un homomorphisme
$$
\varphi_M: M_{p,\sigma} \to M_{p',\sigma'} \leqno{(87)}
$$

Une autre façon, plus directe, de décrire
$\varphi_{\mathfrak{S}_0 M}$ et $\varphi_M$, en termes des coordonnées
$(\alpha, \beta, f)$ au lieu de $(u, a, b, U)$, est de
dire qu'elle [unclear: s'exprime] pour un
$u_{\alpha, \beta, f}$ sur le $u_{\alpha', \beta', f'}$, où $\alpha', \beta', f'$ sont les
images de $\alpha, \beta, f$ par $\varphi$ :
$$
\varphi_{\mathfrak{S}_0 M_{p,\sigma}}(u_{\alpha, \beta, f}) = u_{\alpha', \beta', f'}, \quad \begin{cases} \alpha' = \varphi(\alpha) \\ \beta' = \varphi(\beta) \\ f' = \varphi(f) \end{cases} \leqno{(88)}
$$

On retrouve, ainsi, sous une forme
un peu plus précise, le th. précédent,
auquel le "th." général se ramène
($\mathfrak{G} = \widehat{\mathfrak{G}}_{0,3}^{+ \wedge}$ (le "cas universel"\footnote{tout-au-moins si $\mathfrak{G}$ "universel" pour un tel $p, \sigma$, i.e. $p=1, \sigma^2=1$}), $\mathfrak{G}'$
quelconque). Il [unclear: y a lieu même de le noter]
(89) indépendamment du rôle de $\mathfrak{G}$ et $\mathfrak{G}'$
([unclear: en explicitant] les invariants attachés à $\mathfrak{G}$)
Il est bon d'expliciter dans ce cas "universel" :

$$
\begin{aligned}
\mathfrak{G} &= \widehat{\mathfrak{G}}_{0,3}^{+ \wedge}, \quad B_{p,\sigma} = \mathcal{M}_{0,3}^{1 \sim}[0], \quad \mathcal{P}_{p,\sigma} = \widehat{\Gamma}_{\mathbf{Q}}^{1 \sim} \\
G_{p,\sigma} &= \mathcal{M}_{0,3}^{1 \sim}, \quad Z_p = \mathcal{M}_{0,3}^{1 \sim}(T_p), \quad Z_{\sigma} = \mathcal{M}_{0,3}^{1 \sim}(\sigma/2) \\
\Gamma_{p,\sigma} &= \widehat{\Gamma}_{\mathbf{Q}}
\end{aligned} \leqno{(90)}
$$

\newpage

%% ===== Page 522 =====

[MARGIN: donc $S Z_{\rho'} \simeq S Z_{\rho,\sigma}'$, $G(\hat{\mathbb{Z}}) \simeq S \Gamma_{\rho',\sigma'}$]
$$ \Gamma_{\rho',\sigma'} \simeq \Gamma_{\rho,\sigma}', \quad Z_{\rho'} \simeq Z_{\rho,\sigma}', \quad Z_{\sigma'} \simeq Z_{\rho,\sigma}', \quad \text{donc} $$
(92) $$ S_0 M_{\rho',\sigma'} \simeq S_0 M_{\rho,\sigma}', \quad M_{\rho',\sigma'} \simeq M_{\rho,\sigma}' $$

[MARGIN: Ce n'est pas si simple ! par ex. pour $G = \operatorname{Sl}(2,\hat{\mathbb{Z}})$ $\rho=3$, $\sigma=2$ (ou plus exactement leurs images dans $G/\{\pm 1\}$) on verrait qu'on ne peut se servir des anciennes notations.]
Remarque. En écrivant les relations (80), il devenait évident que les notations utilisées dans cette section II (depuis p. 80 jusqu'à 519) étaient inadéquates, et qu'il convient de les modifier de la façon suivante :
$$ \Sigma_{\rho,\sigma}, B_{\rho,\sigma}, G_{\rho,\sigma}, \gamma_{\rho,\sigma}, Z_{\rho}, Z_{\sigma}, \Gamma_{\rho,\sigma}, S G_{\rho,\sigma} $$
(93) $$ Z_{\rho}, Z_{\sigma}, S_0\Gamma_{\rho,\sigma}, \Gamma_{\rho,\sigma}, S Z_{\rho}, S Z_{\sigma}, S_0 M_{\rho,\sigma}, M_{\rho,\sigma} $$
On met partout un index $'$ pour en faire $\Sigma_{\rho,\sigma}', B_{\rho,\sigma}'$ etc, et on oublie l'exposant $0$ de $S_0\Gamma_{\rho,\sigma}$, $\Gamma_{\rho,\sigma}$, $S_0 M_{\rho,\sigma}$, $M_{\rho,\sigma}$ (qui se trouve donc remplacé par un $'$).

Autres cas particuliers importants explicite-
ment :

2) $G = \operatorname{Sl}(2,\hat{\mathbb{Z}})/\pm 1$, $(\rho, \sigma)$ couple modulaire.

$$ B_{\rho,\sigma}' = \left\{ \begin{pmatrix} \mu & \lambda \\ 0 & 1 \end{pmatrix} \mid \mu \in \hat{\mathbb{Z}}^*, \mu \equiv 1 \, (3), \lambda \in \hat{\mathbb{Z}} \right\} $$
$$ L_0 = \left\{ \begin{pmatrix} 1 & \lambda \\ 0 & 1 \end{pmatrix} \mid \lambda \in \hat{\mathbb{Z}} \right\} $$

[MARGIN: c'est $G(\hat{\mathbb{Z}})/ \mu_2(\hat{\mathbb{Z}})$ avec notations du § 48]
$$ G_{\rho,\sigma}' = Gl'(2, \hat{\mathbb{Z}})/ \pm 1, \quad \text{où } Gl'(2,\hat{\mathbb{Z}}) \overset{\text{def}}{=} \{ U \in \operatorname{Gl}(2,\hat{\mathbb{Z}}) \mid \det(U) \equiv 1 \, (3) \} $$

$$ Z_{\rho}' = \{ c\rho+d \mid c^2+cd+d^2 \in \hat{\mathbb{Z}}^* \} / \pm 1 $$

(94) $$ Z_{\sigma}' = \text{[unclear crossed out text]} \quad \text{avec notation du § 48} $$

$$ S_0 M_{\rho,\sigma}' \simeq S M_{\rho,\sigma}(\hat{\mathbb{Z}})/\Sigma \quad \text{avec notations du § 48} $$

$$ M_{\rho,\sigma}' \simeq M_{\rho,\sigma}(\hat{\mathbb{Z}})/\Sigma \quad id. $$

$$ \Gamma_{\rho,\sigma}' \simeq \Gamma_{\rho,\sigma}(\hat{\mathbb{Z}})/\Sigma_0 \quad \text{avec notations du § 48.} $$

\newpage

%% ===== Page 523 =====

3^\circ) Prenons $\mathcal{G} = \operatorname{Sl}(2, \hat{\mathbf{Z}})$, [unclear: un] couple modulaire
$$
\begin{cases}
B'_{\rho, \sigma} = \left\{ \left. \begin{pmatrix} \mu & \lambda \\ 0 & 1 \end{pmatrix} \right| \mu \in \hat{\mathbf{Z}}^*,\ \mu \equiv 1\ (3),\ \lambda \in \hat{\mathbf{Z}} \right\} \\
G'_{\rho, \sigma} = Gl'(2, \hat{\mathbf{Z}}) = \{ g \in \operatorname{Gl}(2, \hat{\mathbf{Z}}) \mid \det g \equiv 1\ (3) \} \\
Z'_{\rho} = \left\{ c\rho + d \mid c, d \in \hat{\mathbf{Z}}, c^2+cd+d^2 \in \hat{\mathbf{Z}}^* \right\} \\
Z'_{\sigma} = T'_{\sigma} \cdot \{ 1, \tau_{\sigma} \}, \text{ où } T'_{\sigma} = \left\{ t \sigma + u \mid \begin{smallmatrix} t,u \in \hat{\mathbf{Z}} \\ t^2+u^2 \in \hat{\mathbf{Z}}^* \end{smallmatrix} \right\}, \tau_{\sigma} = \begin{pmatrix} -1 & 0 \\ 0 & 1 \end{pmatrix} \\
S_0 M'_{\rho, \sigma} = S_0 M'_{\rho, \sigma}(\hat{\mathbf{Z}}) / \Sigma_0 \text{ avec notations } § 48 \\
M'_{\rho, \sigma} = M'_{\rho, \sigma}(\hat{\mathbf{Z}}) / \Sigma_0 \text{ \hspace{1cm} \hrulefill} \\
\Gamma'_{\rho, \sigma} = \Gamma'_{\rho, \sigma}(\hat{\mathbf{Z}}) / \Sigma_0 \text{ \hspace{1cm} \hrulefill}
\end{cases} \leqno{(95)}
$$

\vspace{0.5cm}
\noindent \underline{III \quad Le groupe $M_{\rho, \sigma}$}

Considérons l'homomorphisme
$$
\begin{aligned}
\hat{\mathbf{Z}} &\xrightarrow{u} \mathcal{G} \\
n &\mapsto \xi_0^n \quad \text{où } \xi_0 = \sigma^{-1}\rho
\end{aligned} \leqno{(1)}
$$
[unclear: (dont on note $L_0$ l'image)], et soit $J$ [unclear: l'idéal noyau de $u$], d'où
$$
(\hat{\mathbf{Z}}/J)^* \isommap \operatorname{Aut}(L_0) \leqno{(2)}
$$
On a un homomorphisme induit
$$
B'_{\rho, \sigma} \to \operatorname{Aut}(L_0) \isommap (\hat{\mathbf{Z}}/J)^* \leqno{(3)}
$$
[unclear: qui sur $L_0$ s'annule, d'où]
$$
\Gamma'_{\rho, \sigma} \to (\hat{\mathbf{Z}}/J)^* \leqno{(4)}
$$
[unclear: d'où par composition]
$$
G'_{\rho, \sigma} \xrightarrow{\chi_G} (\hat{\mathbf{Z}}/J)^* \leqno{(4 \text{ bis})}
$$

\newpage

%% ===== Page 524 =====

[III Considérations de...]

Considérons l'extension $G\ell(2, \hat{\mathbb{Z}})$ de $\mathcal{G}_{0,3}^\wedge$ par $\{\pm 1\}$ :

(1)
\begin{tikzcd}
& & 1 \ar[d] & 1 \ar[d] \\
1 \ar[r] & \{\pm 1\} \ar[r] \ar[d, equal] & S\ell(2, \hat{\mathbb{Z}}) \ar[r] \ar[d] & S_0 \mathcal{G}_{0,3}^\wedge \ar[r] \ar[d] & 1 \\
1 \ar[r] & \{\pm 1\} \ar[r] & G\ell(2, \hat{\mathbb{Z}}) \ar[r] \ar[d] & \mathcal{G}_{0,3}^\wedge \ar[r] \ar[d] & 1 \\
& & \{\pm 1\} \ar[r] \ar[d] & \{\pm 1\} \ar[d] \\
& & 1 & 1
\end{tikzcd}

On s'intéresse au $\tilde{N}_{1,1}$, défini comme l'image inverse de $G\ell(2, \hat{\mathbb{Z}})$ dans $G\ell(2, \hat{\mathbb{Z}})/\{\pm 1\}$ [unclear] via l'hom canonique

(2)
$$ \tilde{M}_{0,3} \longrightarrow G\ell(2, \hat{\mathbb{Z}})' = G\ell(2, \hat{\mathbb{Z}})/\{\pm 1\} $$

(cf § 4 [unclear]...), de sorte qu'on trouve un diagramme de suites exactes

(3)
\begin{tikzcd}
1 \ar[r] & \{\pm 1\} \ar[r] \ar[d, equal] & \tilde{N}_{1,1} \ar[r] \ar[d] & \tilde{M}_{0,3} \ar[r] \ar[d] & 1 \\
1 \ar[r] & \{\pm 1\} \ar[r] & G\ell(2, \hat{\mathbb{Z}}) \ar[r] & G\ell(2, \hat{\mathbb{Z}})' \ar[r] & 1
\end{tikzcd}

On trouve, pour l'image inverse dans $S_0 \mathcal{G}_{0,3}^\wedge \subset G\ell(2, \hat{\mathbb{Z}})'$ de $\tilde{M}_{0,3}$ dans $\tilde{N}_{1,1}$ , ... [unclear]

[MARGIN: NB $\tau_{1,1} \simeq G\ell(2, \hat{\mathbb{Z}})$]

\newpage

%% ===== Page 525 =====

donnant un hom.

(4) $\operatorname{Sl}(2,\Z)^\wedge \to \widetilde{\mathcal{N}}_{1,1}$

dont la restriction à $\operatorname{Sl}(2,\Z)^\wedge$ s'insère dans une suite exacte canoni-
que

(5) $$1 \to \underset{\simeq \widetilde{C}_{1,1}^{+ \wedge}}{\operatorname{Sl}(2,\Z)^\wedge} \to \widetilde{\mathcal{N}}_{1,1} \to \underset{\simeq \widetilde{\Gamma}_{0,3}}{\widetilde{\Gamma}_Q} \to 1$$

qui à son tour s'insère dans le diagramme

(6)
\begin{tikzcd}
1 \ar[r] & \operatorname{Sl}(2,\Z)^\wedge \ar[r] \ar[d, "\text{épi}"'] & \widetilde{\mathcal{N}}_{1,1} \ar[r] \ar[d, "\text{épi}"'] & \widetilde{\Gamma}_Q \ar[r] \ar[d, "\text{épi}"'] & 1 \\
1 \ar[r] & \operatorname{Sl}(2,\widehat{\Z}) \ar[r] & \operatorname{Gl}(2,\widehat{\Z}) \ar[r] & \widehat{\Z}^\ast \ar[r] & 1
\end{tikzcd}

On peut écrire 
$\simeq \widetilde{\Gamma}_Q \simeq \operatorname{Sl}(2,\Z)^\wedge \simeq \widehat{Sl(2,\Z)}/\pm 1$
(7) $\widetilde{\mathcal{M}}_{0,3} \simeq \widetilde{\mathcal{M}}_{0,3}(0) \cdot \widetilde{\Gamma}_{0,3}^{+ \wedge}$
(semi-direct)

en termes de cette décomposition, et
compte tenu qu'on a un hom.
canonique
(8) $\widetilde{\mathcal{M}}_{0,3}(0) \to \operatorname{Gl}(2,\widehat{\Z})$

qui relève l'hom. composé de (2)
avec l'iso. can.
$\widetilde{\mathcal{M}}_{0,3}(0) \xrightarrow{\sim} \dots$

\newpage

%% ===== Page 526 =====

525

(9) $$ \mathcal{N}_{1,1}^\sim \simeq \mathcal{N}_{1,1}^\sim(0) . \operatorname{Sl}(2, \mathbb{Z})^\wedge $$
$$ \mathcal{M}_{0,3}^\sim(0) \simeq \Gamma_Q^\sim \qquad \mathcal{N}_{1,1}^\sim(0) \simeq \dots $$

Ainsi on a une opération de $\Gamma_Q^\sim$
sur $\operatorname{Sl}(2, \mathbb{Z})^\wedge$
que je voudrais expliciter,
utilisant le diagramme cartésien

(10)
\begin{tikzcd}
\operatorname{Sl}(2, \mathbb{Z})^\wedge \ar[r] \ar[d, "can"'] & \widehat{\mathfrak{S}}_{0,3}^+ \simeq \operatorname{Sl}(2, \mathbb{Z})^\wedge / \pm 1 \ar[d, "induit \text{ par } (2)"] \\
\operatorname{Sl}(2, \hat{\mathbb{Z}}) \ar[r] & \operatorname{Sl}(2, \hat{\mathbb{Z}}) / \pm 1
\end{tikzcd}

Un élément $\tilde{v}$ de $\Gamma_Q^\sim$ correspond à
un couple $(\alpha, \beta)$ d'élts

(11) $$ (\alpha, \beta) \in \widehat{\Pi}_{0,3} \times \widehat{\Pi}_{0,3} $$

satisfaisant l'équation des tresses

(12) $$ \alpha \rho \alpha^{-1} = \beta \sigma \beta^{-1} l^v \quad (v \in \hat{\mathbb{Z}}) $$

qui opère sur $\widehat{\mathfrak{S}}_{0,3}^+ \simeq \operatorname{Sl}(2, \mathbb{Z})^\wedge$ par

(13)
$$ \begin{cases} u(\rho) = int_\alpha \rho = \xi \rho \quad & \xi = \alpha \rho \alpha^{-1} \\ u(\sigma) = int_\beta \sigma = \eta \sigma \quad & \eta = \beta \sigma \beta^{-1} \end{cases} $$

ce qui est compatible avec une opération
sur $\operatorname{Sl}(2, \hat{\mathbb{Z}})$ induite
par $x \mapsto u x u^{-1}$, $u \in \operatorname{Gl}(2, \hat{\mathbb{Z}})$ où l'élé-

\newpage

%% ===== Page 527 =====

défini [unclear: comme] l'image de $u$ par (8), de
façon précise
$$
\tilde{u} = i_{\mathfrak{T}_0}(\mu) , \quad \text{où } \mu = \chi(u) \text{ est le multiplicateur, et}
\leqno{(14)}
$$
$$
\begin{cases}
i_{\mathfrak{T}_0} : \mathbf{G}_m \to \operatorname{Gl}(2,\hat{\mathbf{Z}}) \\
\mu \mapsto \begin{pmatrix} \mu & 0 \\ 0 & 1 \end{pmatrix}
\end{cases}
\leqno{(15)}
$$
On aura , de façon précise

\marginpar{les [unclear: notations] soulignées désignent des éléments de $\operatorname{Gl}(2,\hat{\mathbf{Z}})$ dét- ... par $\underline{\alpha}$ et $\underline{\beta}$ ... $\xi, \eta$ ...}
$$
\begin{cases}
u(\rho) = \xi & \quad ; \quad \xi = \underline{\alpha} \rho (\underline{\alpha})^{-1} \\
u(\sigma) = \eta & \quad ; \quad \eta = \varepsilon_\mu \underline{\beta} \sigma (\underline{\beta})^{-1}
\end{cases}
\leqno{(16)}
$$
où, pour mémoire
$$
\varepsilon = \varepsilon_2(\mu) \in \{\pm 1\} \quad \begin{cases} \varepsilon = + 1 & \text{si } \mu \equiv 1\ (4) \\ \varepsilon = - 1 & \text{si } \mu \equiv -1\ (4) \end{cases}
\leqno{(17)}
$$

Bien entendu, $\operatorname{Sl}(2,\mathbf{Z})^\wedge$ est engendré
top[ologiquement] par $\rho$ et $\sigma$, et l'on identifie
$\pm 1$ au sous-groupe central de $\operatorname{Sl}(2,\mathbf{Z})$
dans $\operatorname{Sl}(2,\mathbf{Z})^\wedge$, on écrira
$-g$ pour le produit de $g \in \operatorname{Sl}(2,\mathbf{Z})^\wedge$
par $-1$. On a donc dans $\operatorname{Sl}(2,\mathbf{Z})^\wedge$
$$
\begin{cases}
u(\rho) = \operatorname{int}(\tilde{\alpha})(\rho) = \xi \\
u(\sigma) = \varepsilon \operatorname{int}(\tilde{\beta})(\sigma) = \eta
\end{cases}
\leqno{(18)}
$$
où
$$
\begin{cases}
\xi = \tilde{\alpha} \rho (\tilde{\alpha})^{-1}, \quad \eta = \varepsilon \tilde{\beta} \sigma (\tilde{\beta})^{-1} \\
\varepsilon = \varepsilon_2(\mu) \in \{\pm 1\} \subset \operatorname{Sl}(2,\mathbf{Z})^\wedge
\end{cases}
\leqno{(19)}
$$

\newpage

%% ===== Page 528 =====

où maintenant $\tilde{\alpha}$ et $\tilde{\beta}$ désignent des relèvements arbitraires de $\alpha$ et $\beta$ dans $\operatorname{Sl}(2, \hat{\mathbf{Z}})^{\wedge}$ - ils ne sont déterminés qu'au signe près (i.e. $-\tilde{\alpha}, -\tilde{\beta}$), mais $\xi$ et $\eta$ sont déterminés sans ambiguité par $u$. Notons aussi
$$
u(\varepsilon_0) = \varepsilon_0^\mu \qquad \text{[unclear: avec \dots]} \leqno{(20)}
$$
Si on se donne un isomorphisme
$$
N_{1,1}^\sim \xrightarrow{\sim \Psi} \mathcal{G} \leqno{(20)}
$$
déduisant un homomorphisme
$$
\operatorname{Sl}(2, \hat{\mathbf{Z}})^{\wedge} \xrightarrow{\tilde{\Psi}} \mathfrak{S} \subset G \leqno{(21)}
$$
les images de $\rho, \sigma$ dans $\mathfrak{S}$ (notées par les mêmes lettres) satisferont\marginpar{dans le cas où il y a un [unclear: centre] canonique involutif, noté $-1 \dots$}
$$
\rho^6 = \sigma^4 = 1, \qquad \rho^3 = \sigma^2 \text{ et qui soit \dots} \leqno{(22)}
$$
et il se peut, pour $u \in N_{1,1}^\sim$, correspondant à des éléments $\tilde{\alpha}, \tilde{\beta}$ de $\operatorname{Sl}(2, \hat{\mathbf{Z}})^{\wedge}$ (déterminés au signe près), d'en déduire des images de ceux-ci dans $\mathfrak{S}$ et $\mathcal{G}$, satisfaisant les relations idoines
$$
u(\tilde{\rho}) = \tilde{\xi} \tilde{\rho} \tilde{\xi}^{-1}, \qquad u(\tilde{\sigma}) = \tilde{\eta} \tilde{\sigma} \tilde{\eta}^{-1} \leqno{(23)}
$$
et faire jouer par $\tilde{\alpha}$ et $\tilde{\beta}$. Mais il ne

\newpage

%% ===== Page 529 =====

obligés, comme [unclear: en général] on n'a pas \marginpar{[unclear: noté !]}
$\tilde{\alpha} \in \operatorname{Sl}(2,\hat{\mathbf{Z}})^{\wedge}$, seulement $\tilde{\alpha} \in \operatorname{Gl}(2,\hat{\mathbf{Z}})^{\wedge}$, de
faire intervenir l'élément $\tau$ de $\operatorname{Gl}(2,\hat{\mathbf{Z}})^{\wedge}$,
et [unclear: par suite] de $G$ sur $\mathfrak{S}$. On
aura, en posant maintenant\footnote{Ce qui est la convention habituelle je pense...}

$$
\begin{aligned}
\varepsilon_0 &= \tilde{\rho}\tilde{\sigma}, & \varepsilon_1 &= \tilde{\rho}\tilde{\sigma}\tilde{\rho} \\
&= -\tilde{\sigma}\tilde{\rho} & &= -\tilde{\rho}\tilde{\sigma}^{-1}
\end{aligned} \leqno{(24)} \marginpar{je laisse tomber en redac}
$$

$$
\tau(\sigma) = \sigma^{-1}, \quad \tau(\rho) = \rho^{-1} \leqno{(25)}
$$
qui implique, via (24) et (22)
$$
\tau(\varepsilon_0) = \varepsilon_0^{-1}, \quad \tau(\varepsilon_1) = \varepsilon_1^{-1} \leqno{(26)}
$$

et on introduit le groupe
$$
\mathfrak{S}^{\pm} \defeq \mathfrak{S} \rtimes \{1,\tau\} \leqno{(27)}
$$

où $\{1,\tau\}$ opère sur $\mathfrak{S}$ par les formules (25).

Ceci dit, on a [unclear: un moyen] de pour [unclear: passer aux] constructions déjà faites par ailleurs.

\newpage

%% ===== Page 530 =====

IV.

Soit donc $\mathfrak{S}$ un groupe profini, muni de $2$ générateurs\footnote{Ceci signifie simplement que $\mathfrak{S}$ est un quotient de $\operatorname{Sl}(2,\hat{\mathbf{Z}})^\sim$.} $\rho, \sigma \in \mathfrak{S}$ satisfaisant les relations
$$
\rho^6 = \sigma^4 = 1, \quad \rho^3 = \sigma^2 \quad (\defeq "-1") \leqno{(1)}
$$
un cas particulièrement intéressant est celui où $\rho^3 = \sigma^2 = 1$, au quel nous nous restreindrons à la limite. Soit $\tau$ un automorphisme de $\mathfrak{S}$ satisfaisant\marginpar{Introduisons de suite l'invol. $\tau' = \tau c$ (ce qui signifie que $\tau'(\rho) = \rho^{-1}$, $\tau'(\sigma) = \sigma$, $\tau'(\varepsilon_0) = \varepsilon_1^{-1}$, $\tau'(\varepsilon_1) = \varepsilon_0^{-1}$)}
$$
\tau(\sigma) = \sigma^{-1} \ (= -\sigma), \quad \tau(\rho) = \sigma(\rho^{-1}) \leqno{(2)}
$$
(ce qui signifie que $\tau$ "renverse" l'orientation de $\mathfrak{S}$), ce qui implique
$$
\tau^2 = 1 \leqno{(3)}
$$
$$
\tau(\varepsilon_0) = \varepsilon_0^{-1}, \quad \tau(\varepsilon_1) = \varepsilon_1^{-1} \leqno{(4)}
$$
\marginpar{[unclear: De fa. analogue l'invol. $\tau'$ \dots (5) demande $\varepsilon_0 \leftrightarrow \varepsilon_1$]}
$$
\varepsilon_0 = \sigma\rho \ (= -\sigma^{-1}\rho), \quad \varepsilon_1 = \rho\sigma = \sigma(\varepsilon_0) = \sigma^{-1}(\varepsilon_0) = \rho(\varepsilon_0) \leqno{(5)}
$$
$$
\mathfrak{S}' = \{1, \tau\} \cdot \mathfrak{S}
$$
donc
$$
1 \to \mathfrak{S} \to \mathfrak{S}' \to \{1, \tau\} \to 1 \leqno{(6)}
$$
Soit
$$
B_{\rho, \sigma} \subset \operatorname{Aut}(\mathfrak{S}) \times \hat{\mathbf{Z}}^* \leqno{(7)}
$$
$$
= \left\{ (u, \mu) \ \Big|\  \begin{array}{l} u(\varepsilon_0) = \varepsilon_0^\mu \\ u(\sigma) \text{ conjugué dans } \mathfrak{S} \text{ à } \sigma^\mu \\ u(\rho) \text{ conjugué dans } \mathfrak{S} \text{ à } \rho^\mu \end{array} \right\}
$$
qui est évidemment un sous-groupe, la troisième condition signifie d'ailleurs,

\newpage

%% ===== Page 531 =====

puisque $\rho^6 = 1$ i.e.\ $\rho^\mu$ ne dépend que de l'
image de $\mu$ dans $(\mathbf{Z}/6\mathbf{Z})^* \simeq (\mathbf{Z}/3\mathbf{Z})^* \simeq \{\pm 1\}$
on a donc
$$
\rho^\mu = \rho^{\varepsilon_3(\mu)} \leqno{(8)}
$$
où pour $\varepsilon_3(\mu)=1$ on a $u(\rho)$ est conjugué à
$\rho$ dans $\mathfrak{S}$, et pour $\varepsilon_3(\mu)=-1$ il est conjugué
à $\rho^{-1}$ dans $\mathfrak{S}$, c'est-à-dire à $\rho$ par un
élément de $\mathfrak{S}' \setminus \mathfrak{S}$, i.e.\ un élément de la
forme $\alpha \tau$ avec $\alpha \in \mathfrak{S}$ (puisque
$\tau \rho \tau^{-1} = \rho^{-1}$ justement\dots). On voit de même qu'on a
$$
\sigma^\mu = \sigma^{\varepsilon_2(\mu)} = \varepsilon_2(\mu)\sigma =
\begin{cases}
= \sigma & \text{ si } \mu \equiv 1 \ (4) \\
= \sigma^{-1} = -\sigma & \text{ si } \mu \equiv -1 \ (4)
\end{cases} \leqno{(9)}
$$
Ainsi :
Pour qu'un automorphisme $(u,\mu) \in \operatorname{Aut}(\mathfrak{S}) \times \hat{\mathbf{Z}}^*$
soit dans $B_{\rho,\sigma}$ il faut et il suffit qu'il
existe $\alpha \in \mathfrak{S}'$, $\beta \in \mathfrak{S}$ tels qu'on ait

\newpage

%% ===== Page 532 =====

$$
u(\rho) = \operatorname{int}(\alpha)(\rho) , \quad u(\sigma) = \varepsilon_2(\mu) \operatorname{int}(\beta)(\sigma)
\leqno{(10)}
$$
et enfin
$$
u(\varepsilon_0) = {\varepsilon_0}^\mu ,
\leqno{(11)}
$$
où de plus \marginpar{p.e. $\alpha \in \mathfrak{T}$}
$$
\mu \equiv 1 \ (3)
\leqno{(12)}
$$
On notera que $u$ (ou $(u,\mu)$) est uniquement déterminé par la connaissance de $\alpha$, de $\beta$, et de $\varepsilon_2(\mu)$ et de $\mu$ :
$$
\xi = \alpha \rho (\alpha)^{-1}, \quad \eta = \varepsilon_2(\mu) \beta \sigma (\beta)^{-1}
\leqno{(13)}
$$
puisqu'on aura alors $u$ par
$$
u(\rho) = \xi, \quad u(\sigma) = \eta .
\leqno{(14)}
$$
D'ailleurs, dans le cas où
$$
\hat{\mathbf{Z}} \longrightarrow \widetilde{\mathcal{G}}
\leqno{(15)}
$$
$$
1 \longmapsto \varepsilon_0
$$
est injectif, alors $\mu$ est déterminé [unclear: univoquement] en fonction de $u$ comme restriction de $u$ sur $L_0 \subset \widetilde{\mathcal{G}}$, image de (15), donc $(u,\mu)$ est déterminé par $u$, donc par $(\alpha,\beta, \varepsilon_2(\mu))$. D'ailleurs, si $\widetilde{\mathcal{G}}$ est sans [unclear: points] pour [unclear: d'ordre 2] dans $\operatorname{Sl}(2,\hat{\mathbf{Z}})'$ (cela fera [unclear: moyen] [unclear: point] sur les éléments [unclear: habituels]), on aura $\alpha$ [unclear: à lui-même] seul.

\newpage

%% ===== Page 533 =====

détermine déjà $\mu$, donc $\varepsilon$, donc [unclear]
cet élément de $B_{\rho,\sigma}$ est déterminé
à l'aide de $\alpha, \beta$ seuls.

Mais quelles conditions doivent satisfaire
$(\alpha, \beta, \mu)$ pour définir un autom. de $B_{\rho,\sigma}$ ?
Tout d'abord, il faut qu'il existe bien un
homom. (nécess. unique) satisfais. à
(14) — c'est une condition sur $\xi, \eta$ seuls.
Il faut ensuite exprimer (11), qui s'écrit, puisque

$$ \varepsilon_0 = \sigma p $$
$$ u(\varepsilon_0) = u(\sigma p) = [unclear] = \eta \sigma \xi p = \eta \sigma \xi \sigma^{-1} (\overbrace{\sigma p}^{= \varepsilon_0}) $$
$$ \eta (\sigma \xi \sigma^{-1}) \varepsilon_0 = \varepsilon_0^\mu \text{ i.e.} $$
$$ \eta (\sigma \xi \sigma^{-1}) = \varepsilon_0^{\mu-1} [unclear] $$

Or on a,
$$ \sigma(\eta) = [unclear] = \varepsilon \sigma(\beta) \sigma^2(\beta)^{-1} = [unclear] $$
i.e.
(15) $$ \eta \sigma(\eta) = 1 $$

(*)
$$ \varepsilon_0' = \sigma^{-1} p = - \varepsilon_0 $$
d'où
$$ (\varepsilon'_0)^\mu = (-1)^\mu \varepsilon_0^\mu = - \varepsilon_0^\mu \quad (\text{car } \mu \equiv 1 \, (2)) $$
$$ u(\varepsilon'_0) = - u(\varepsilon_0) $$ car $u(-1) = -1$ cf ci-dessous.
donc la relation (11) équivaut à :
(10) $$ u(\varepsilon'_0) = (\varepsilon'_0)^\mu $$
qui, pour la susdite valeur, s'écrit

\newpage

%% ===== Page 534 =====

$$
\eta = \varepsilon_1^{\prime 2\nu} y \qquad \nu = (\mu-1)/2, \quad \varepsilon'_1 = -\varepsilon_1 = \rho\sigma^{-1}
$$
$$
\varepsilon_1^{\prime 2\nu} = (-1)^{2\nu} \varepsilon_1^{2\nu} = \varepsilon_1^{2\nu} = l_1^\nu,
$$
d'où la relation du haut prend la forme
$$
\eta = y l_1^\nu \leqno (16)
$$
avec $\nu = (\mu-1)/2$.

On doit encore prouver le

Lemme. Soit $u$ un automorphisme de $\mathcal{G}$ tel que $u(\rho)$ soit conjugué à $\rho$, $u(\sigma)$ conjugué à $\sigma^{\pm 1}$. Alors
$$
u(-1) = -1 \leqno (17)
$$
Dém. $u(\sigma) = \operatorname{int}(\alpha)(\sigma^{\pm 1})$ d'où $u(\sigma^2) = \operatorname{int}(\alpha)(\sigma^{\pm 2}) = \operatorname{int}(\alpha)(-1) = -1$, OK.

Reste : à exprimer que $\xi$ est conjugué à $\rho^\mu$, et $u(\sigma)$ conjugué à $\sigma$.
et celà est tautologique pour $u(\sigma)$ puisque $u(\sigma) = \gamma_3(\mu)(\varepsilon_2(\mu)\sigma)\gamma_3(\mu)^{-1}$ (voir (19)),
et pour $\rho$ il signifie que l'on a \marginpar{[unclear: où $\gamma_3(\mu) \in \mathcal{G}$ l'élément de $\mathcal{G}$ défini aux \dots $\varepsilon_3(\mu)$]}
$$
\alpha \equiv \tau \gamma_3(\mu) \pmod{\operatorname{Centr}(\rho)}
$$
qui s'écrit aussi $\tau^{-1} \gamma_3(\mu) \alpha \in \operatorname{Centr}(\rho)$.

Conclusion pour l'application
$$
B_{\rho, \sigma} \longrightarrow \underline{\mathcal{G}} \times \underline{\mathcal{G}} \times \hat{\mathbf{Z}}^* \leqno (18)
$$
$$
(u, \mu) \longmapsto ([u(\rho)], [u(\sigma)], \mu)
$$
est injective, et elle a comme image l'ensemble des triplets $(\xi, \eta, \mu)$ satisfaisant les conditions suivantes :
\begin{enumerate}
\item[1º)] $\xi \sim \rho^\mu$ \quad ($\xi \sim \varepsilon_1(\mu)$ dans $\mathcal{G}$)
\item[2º)] $\eta \sim \sigma$ \quad ($\eta \sim \varepsilon_2(\mu)$ dans $\mathcal{G}$)
\end{enumerate}

\newpage

%% ===== Page 535 =====

3°) On a la relation des lacets (16).

4°) $\exists$ autom. $u$ de $\mathcal{G}$, satisfaisant $u(p) = \xi p$, $u(\sigma) = \eta \sigma$.

Les conditions 1'), 2'), 3') sont équivalentes à
celle-ci :

4') $\exists (\alpha, \mu) \in \mathcal{G}' \times \widehat{\mathbf{Z}}^*$, avec $\alpha \in \mathcal{G}$ et $\mu \equiv 1\ (3)$
i.e $\alpha \bmod{\widehat{\mathcal{G}}} = \tau^{\nu_3(\mu)}$, et tels qu'on ait
$\xi = \alpha p (\alpha)^{-1}$, $\eta = \varepsilon_2(\mu)$.

Considérons l'hom. de groupes
$$
\begin{tikzcd}
L_0 \arrow[r, "i_{L_0,B}"] & B_{p,\sigma} \\
\ell \arrow[r, maps to] & (\operatorname{int}(\ell), 1)
\end{tikzcd} \leqno{(19)}
$$
Il est injectif si on suppose toujours
$$
L_0 \cap \operatorname{Centr}(\mathcal{G}) = \{ 1 \} , \leqno{(20)}
$$
son image est distinguée, et on a
$$
u(i_{L_0,B}(\ell)) = i_{L_0,B}(u(\ell)) = i_{L_0,B}(\ell^\mu) \leqno{(21)}
$$
où $\mu$ est le multiplicateur. On pose
$$
\gamma_{p,\sigma} = B_{p,\sigma} / \operatorname{Im} L_0 , \leqno{(22)}
$$
on construit alors
la suite exacte canonique
$$
1 \to L_0 \to B_{p,\sigma} \to \gamma_{p,\sigma} \to 1 \leqno{(23)}
$$
et on a un hom. ([unclear: déduit de $f$ défini sur] $B_{p,\sigma}$)
$$
\gamma_{p,\sigma} \xrightarrow{f_\sigma} \widehat{\mathbf{Z}}^* \leqno{(24)}
$$
dont le noyau est noté $\gamma_{p,\sigma}^+$ :
$$
1 \to \gamma_{p,\sigma}^+ \to \gamma_{p,\sigma} \to \widehat{\mathbf{Z}}^* \to 1 \leqno{(25)}
$$

\newpage

%% ===== Page 536 =====

571

Considérons les hom. composés 
(26)
$$ \gamma_{\rho,\sigma} \xrightarrow{\chi_\gamma} \hat{\mathbb{Z}}^\times \xrightarrow{\varepsilon_2} \pm 1 $$
$$ \gamma_{\rho,\sigma} \xrightarrow{\chi_\gamma} \hat{\mathbb{Z}}^\times \xrightarrow{\varepsilon_3} \pm 1 $$
et soient 
(27)
$$ \gamma_{\rho,\sigma}', \gamma_{\rho,\sigma}'' \subset \gamma_{\rho,\sigma} \text{ sous-groupes d'indice 2, définis par} $$
$$ \gamma_{\rho,\sigma}' = \operatorname{Ker}(\varepsilon_3 \circ \chi_\gamma), \quad \gamma_{\rho,\sigma}'' = \operatorname{Ker}(\varepsilon_2 \circ \chi_\gamma) $$

Soit maintenant
(28)
$$ G_{\rho,\sigma} = (B_{\rho,\sigma} \cdot \widehat{\mathfrak{S}}) / L_0 $$
(1/2 direct)
$$ \simeq S_0 G_{\rho,\sigma} $$
où 
(29)
$$ L_0 \xrightarrow{i_{L_0, B}} B_{\rho,\sigma} \cdot \widehat{\mathfrak{S}} $$
$$ \ell \mapsto (i(\ell), \ell^{-1}) $$
dont l'image est invariante. On a un 
diagramme commutatif 
(30)
\begin{tikzcd}
L_0 \ar[r, hook] \ar[dr] & \widehat{\mathfrak{S}} \ar[dr] \\
& B_{\rho,\sigma} \ar[r] & G_{\rho,\sigma}
\end{tikzcd}
et 
$$ \operatorname{Ker}(\widehat{\mathfrak{S}} \to G_{\rho,\sigma}) \simeq \operatorname{Ker}(L_0 \to B_{\rho,\sigma}) $$
$$ \simeq L_0 \cap \text{Centre}(\widehat{\mathfrak{S}}) \text{ est trivial vu } \dots (20), $$
ce qui [unclear] justifie les 
suites exactes 
(31)
\begin{tikzcd}
1 \ar[r] & L_0 \ar[r] \ar[d, hook] & B_{\rho,\sigma} \ar[r, "\delta_B"] \ar[d] & \gamma_{\rho,\sigma} \ar[r] \ar[d, equal] & 1 \\
1 \ar[r] & \widehat{\mathfrak{S}} \ar[r] & G_{\rho,\sigma} \ar[r, "\delta_G"'] & \gamma_{\rho,\sigma} 
\end{tikzcd}
$$ [ \xrightarrow{\chi_G} \hat{\mathbb{Z}}^\times ] $$

\newpage

%% ===== Page 537 =====

Posons
$$
\begin{cases} 
Z_{\rho} = \operatorname{Centr}_{C_{\rho, \sigma}}(\rho) \\ 
Z_{\sigma} = \operatorname{Centr}_{C_{\rho, \sigma}}(\sigma) 
\end{cases} \leqno{(32)}
$$
alors les hom. induits par $\delta_C$
$$
Z_{\rho} \xrightarrow{\delta_{\rho}} \gamma_{\rho, \sigma} \, , \quad Z_{\sigma} \xrightarrow{\delta_{\sigma}} \gamma_{\rho, \sigma} \leqno{(33)}
$$
ont des images qui contiennent $\gamma'_{\rho, \sigma}$ et $\gamma''_{\rho, \sigma}$ respectivement. D'ailleurs l'image de $Z_{\rho}$ est contenue dans $\gamma'_{\rho, \sigma}$ car $\rho$ et $\rho^{-1}$ sont conjugués dans $\overline{\mathfrak{S}}$. L'image de $Z_{\sigma}$ est évidente : $\gamma''_{\rho, \sigma}$ car $\sigma$ et $\sigma^{-1} = -\sigma$ sont conjugués dans $\overline{\mathfrak{S}}$.\footnote{La condition pour que l'image [unclear: de $Z_{\rho}$] soit [unclear: d'indice $2$] n'est d'ailleurs pas satisfaite [unclear: dans $C_{\rho, \sigma}$]. Celui-ci pour $\overline{\mathfrak{S}} = \operatorname{Sl}(2, \hat{\mathbf{Z}})$, l'image de $\delta_{\rho}$ est $\gamma'_{\rho, \sigma}$... Mais d'après [unclear: (1)] c'est un autom. intérieur de $\overline{\mathfrak{S}}$.}

Notons p.ex.\ pour $\overline{\mathfrak{S}} = \operatorname{Sl}(2, \hat{\mathbf{Z}})$, $\sigma^2=1$, la dernière condition n'est pas satisfaite, donc $Z_{\sigma} \to \gamma_{\rho, \sigma}$ surjective, équivalente à $\sigma$ et $-\sigma$ conjugués dans $\operatorname{Sl}(2, \hat{\mathbf{Z}})$ (cf.\ Steinberg p.ex.) --- a-t-on pareille propriété dans le $\operatorname{Sl}(2, \mathbf{Z})^{\wedge} = \mathfrak{S}_{1,1}^{\wedge +}$ ?

\newpage

%% ===== Page 538 =====

Proposition
\begin{enumerate}
\item[1^\circ)] L'image $\delta_{\rho}(Z_{\rho})$ contient $\gamma_{\rho,\sigma}'$, et est réduite à $\gamma_{\rho,\sigma}'$ ssi $\rho$ et $\rho^{-1}$ non conjugués dans $\mathfrak{S}$ --- donc elle est égale à $\gamma_{\rho,\sigma}$ i.e.\ $\delta_{\rho} \colon Z_{\rho} \to \gamma_{\rho,\sigma}$ épi ssi $\rho$ et $\rho^{-1}$ sont conjugués dans $\mathfrak{S}$.
\item[2^\circ)] L'image $\delta_{\sigma}(Z_{\sigma})$ contient $\gamma_{\rho,\sigma}''$, et elle est réduite à $\gamma_{\rho,\sigma}''$ ssi $\sigma$ et $\sigma^{-1} = -\sigma$ non conjugués dans $\mathfrak{S}$ --- donc elle est égale à $\gamma_{\rho,\sigma}$ i.e.\ $\delta_{\sigma} \colon Z_{\sigma} \to \gamma_{\rho,\sigma}$ épi ssi $\sigma$ et $\sigma^{-1} = -\sigma$ sont conjugués dans $\mathfrak{S}$.
\end{enumerate}

Corollaire. Si $\mathfrak{S}$ contient comme quotient $\operatorname{Sl}(2,\mathbf{Z}/3\mathbf{Z})'$ ; avec $\rho = \begin{pmatrix} 0 & -1 \\ 1 & 1 \end{pmatrix}$, $\sigma = \begin{pmatrix} 0 & -1 \\ 1 & 0 \end{pmatrix}$ $(\text{mod } \pm 1)$, alors l'image de $\delta_{\rho}$ est $\gamma_{\rho,\sigma}'$.

En effet, on voit que $\rho$ et $\rho^{-1}$ ne sont pas conjugués dans $\mathfrak{S}$, car ils ne le sont pas dans $\operatorname{Sl}(2,\mathbf{Z}/3\mathbf{Z})'$ --- On a $(\tau_{\infty})(\rho) = \rho^{-1}$, donc tous les éléments de $\operatorname{Gl}(2,\mathbf{Z}/3\mathbf{Z})$ qui conjuguent $\rho$ en $\rho^{-1}$ sont de la forme $f \tau_{\infty}$, où $f = c\rho + d$, on a $\det(f \tau_{\infty}) = 1$ i.e.\ $\det(f) = -1$ car $\det(\tau_{\infty}) = -1$ i.e.\ $c^2 + cd + d^2 = -1$, impossible.

\newpage

%% ===== Page 539 =====

\noindent
Corollaire 2 Si $\sigma^2=1$, alors $Z_\sigma \to \gamma_{\rho,\sigma}$ est surjectif. Si par contre $\mathcal{G}$ contient comme quotient $\operatorname{Sl}(2,\mathbf{Z}/4\mathbf{Z})$ (sans \og $'$ \fg{} !), avec $\rho,\sigma$ comme d'habitude, alors $\operatorname{Im}(\delta_\sigma) = \gamma''_{\rho,\sigma}$.

Il faut voir que $\sigma$ et $\sigma^{-1}=-\sigma$ non conjugués dans $\operatorname{Sl}(2,\mathbf{Z}/4\mathbf{Z})$, en utilisant que $\tau_\infty(\sigma) = \sigma^{-1}$, donc il faut les élts de $\operatorname{Gl}(2,\mathbf{Z}/4\mathbf{Z})$ qui conjuguent $\sigma$ en $\sigma^{-1}$ (avec $g = f \tau_\infty \dots$) on a $f = t \sigma + u$, et dire que $\det(g) = 1$ signifie $\det f = -1$ i.e.\ $t^2+u^2=-1$, impossible.

Posons
$$
\begin{cases}
S_{Z_\rho} = Z_\rho \cap \mathcal{G} = \operatorname{Centr}_{\mathcal{G}}(\rho) \\
S_{Z_\sigma} = Z_\sigma \cap \mathcal{G} = \operatorname{Centr}_{\mathcal{G}}(\sigma) \\
S_{B_{\rho,\sigma}} = B_{\rho,\sigma} \cap \mathcal{G} = L_0
\end{cases} \leqno{(34)}
$$

$$
S_{\mathcal{M}_{\rho,\sigma}} \subset B_{\rho,\sigma} \times_{Z_\rho \times Z_\sigma} \mathcal{G}_{\rho,\sigma} \leqno{(35)}
$$

\newpage

%% ===== Page 540 =====

Notons que d'après (4), on a
$$
\tau_B \defeq (\tau, -1) \in B_{\rho, \sigma} \leqno{(35)}
$$
Et évidemment $\tau_B \notin L_0 = B_{\rho, \sigma} \cap \widehat{\mathcal{G}}$. Donc on a un plongement\footnote{NB $\mathcal{G}'$ est un quotient de $\operatorname{Gl}(2, \hat{\mathbf{Z}})$ et le $\tau$ s'identifie à l'image $\tau_\infty = \begin{pmatrix} -1 & 0 \\ 0 & 1 \end{pmatrix}$ de $\operatorname{Gl}(2, \hat{\mathbf{Z}})$.}
$$
\begin{tikzcd}[row sep=tiny]
\mathcal{G}' = \{1, \tau\} \cdot \widehat{\mathcal{G}} \ar[r, hook] & G_{\rho, \sigma} \\
\tau \ar[r, |->] & \tau_B \\
g \ar[r, |->] & g \in \widehat{\mathcal{G}}
\end{tikzcd} \leqno{(36)}
$$
donc l'action $\tau$ de $\widehat{\mathcal{G}}$ [unclear: doit coïncider avec l'action] induite par $\operatorname{int}(\tau_B)$, i.e.\ $\tau_B$ se conduit bien comme un élément de $G_{\rho, \sigma}$. En effet [unclear: on a] l'identité
$$
\tau_B(\sigma) = \sigma^{-1} \ (= -\sigma)
$$
qui montre que $\tau_B$ normalise le sous-groupe $L_0$ engendré par $\sigma$ (isomorphe au groupe $\mathbf{Z}/4\mathbf{Z}$). Posons $Z_\sigma = \operatorname{Centr}_{G_{\rho, \sigma}}(L_0)$, introduisons
$$
N_\sigma = \{1, \tau_B\} \cdot Z_\sigma \quad \text{([unclear: produit semi-direct])} \leqno{(37)}
$$
et le hom. canonique
$$
\begin{cases} 
N_\sigma \overset{i_\sigma}{\longrightarrow} G_{\rho, \sigma} \\ 
\tau_B \mapsto \tau_B, \ x \mapsto x \ (x \in Z_\sigma) 
\end{cases} \leqno{(38)}
$$
L'image de ce hom. n'est autre que le normalisateur de $\sigma$ dans $G_{\rho, \sigma}$, car un élément $g \in G_{\rho, \sigma}$ qui normalise $L_0$ doit transformer $\sigma$ en un générateur de $L_0$, i.e.\ en $\sigma$ ou $\sigma^{-1}$; dans le premier cas on a $g \in Z_\sigma$, dans le deuxième $g(\sigma) = \tau_B(\sigma)$ i.e.\ $\tau_B^{-1} g \in Z_\sigma$ i.e.\ $g \in \tau_B Z_\sigma$. L'hom. (38) est injectif.

\newpage

%% ===== Page 541 =====

$$
\tau_B \notin Z_\sigma \text{ i.e. } \sigma \neq \sigma^{-1} \text{ i.e. } \sigma^2 \neq 1
$$
i.e.\ $"-1" \neq 1$.

Puisqu'on y est, profitons de la formule
$$
\tau_B(\rho) = \sigma(\rho^{-1})
$$
pour introduire
$$
\tau' = \sigma^{-1} \tau \in \widetilde{\mathfrak{S}}' \subset G_{\rho, \sigma} \leqno{(39)}
$$
(qui correspond à l'élément
$$
\tau'_{\infty} = \sigma_{\infty}^{-1} \tau_{\infty} = - \sigma_{\infty} \tau_{\infty} = \begin{pmatrix} 0 & 1 \\ 1 & 0 \end{pmatrix} \in \operatorname{Gl}(2, \mathbf{Z})) \leqno{(40)}
$$
on a donc\footnote{NB. On a $\tau'^2 = 1$, donc $\tau'(\sigma) = \sigma^{-1}$ etc. On peut considérer $\tau'$ comme élément de $N_\rho$, i.e.\ $N_\rho = \{1, \tau'\} \cdot Z_\rho$ (produit semi-direct) dans $G_{\rho, \sigma}$.}
$$
\tau'^2 = 1, \quad \tau'(\rho) = \rho^{-1} \leqno{(40)}
$$
Ainsi $\tau'$ normalise le sous-groupe $L_\rho$ (isomorphe à un quotient de $\hat{\mathbf{Z}}$ engendré par $\rho$), donc il normalise le centralisateur de $\rho$, et on peut construire
$$
N_\rho \defeq \{1, \tau'\} \cdot Z_\rho \quad \text{(semi-direct)} \leqno{(41)}
$$
et on a l'hom.\ canonique
$$
\begin{cases} N_\rho \xrightarrow{i_\rho} G_{\rho, \sigma} \\ \tau' \mapsto \tau', \quad x \mapsto x \text{ pour } x \in Z_\rho \end{cases} \leqno{(42)}
$$

L'image en est le normalisateur de $L_\rho$ dans $G_{\rho, \sigma}$ - on le voit comme plus haut pour $\sigma$, $\rho, \rho^{-1}$ étant les seuls générateurs de $L_\rho$. L'hom.\ $i_\rho$ est injectif ssi $\tau' \notin Z_\rho$ i.e.\ $\rho \neq \rho^{-1}$, i.e.\ $\rho^2 \neq 1$, i.e.\ $\rho \neq -1$ (NB $\rho = -1$ implique\footnote{car $\rho = -1$ implique que $\rho$ commute à $\sigma$, et que $\widetilde{\mathfrak{S}}'$ soit un quotient de $\mathbf{Z}/2\mathbf{Z} \times \mathbf{Z}/2\mathbf{Z}$ (ce qu'il est en effet pour $2\rho = 0$, c'est à dire en $\mathbf{Z}/4\mathbf{Z}$ avec $\rho$ correspondant à $2$, $\sigma$ à $-1$ \dots). On a fait ce qu'il fallait pour exclure de tels épimorphismes.} que $\rho$ commute à $\sigma$, et que $\widetilde{\mathfrak{S}}'$ \dots).

\marginpar{épimorphismes}
$$
\begin{cases}
N_\sigma \xrightarrow{j_\sigma = \delta_\sigma \circ i_\sigma} G_{\rho,\sigma} \\
N_\rho \xrightarrow{j_\rho = \delta_\rho \circ i_\rho} G_{\rho,\sigma}
\end{cases} \leqno{(43)}
$$

\newpage

%% ===== Page 542 =====

$$
\mathcal{S}\mathcal{G}_{\rho,\sigma}^\# \subset \mathfrak{S}' \times \mathfrak{S} \times \hat{\mathbf{Z}}^\times \leqno{(44)}
$$
$$
= \left\{ (\alpha, \beta, \mu) \mid \begin{array}{l} 
\alpha \equiv \tau^{\nu_3(\mu)} \bmod \mathfrak{L} \text{ i.e. } \chi(\alpha) = \varepsilon_3(\mu) \\
\alpha \rho(\alpha)^{-1} = \varepsilon_2(\mu) \beta \sigma(\beta)^{-1} l_1^\nu, \text{ où } \nu = (\mu-1)/2 \\
\hspace{5cm} (l_1 = \varepsilon_1^2)
\end{array} \right\}
$$

d'où
\marginpar{ (46) $\mathcal{S}\mathcal{G}_{\rho,\sigma}^\# \xrightarrow{\text{épi.}} \Sigma_{\rho,\sigma}^\#$ }
$$
\mathcal{S}\mathcal{G}_{\rho,\sigma}^\# \longrightarrow \Sigma_{\rho,\sigma} \subset \mathfrak{S}' \times \mathfrak{S} \times \hat{\mathbf{Z}}^\times \leqno{(45)}
$$
$$
\begin{array}{l}
\hspace{2.5cm} \text{cl. conj. de } \rho, \sigma \\
(\alpha, \beta, \mu) \mapsto (\alpha \rho \alpha^{-1}, \beta \sigma \beta^{-1}, \mu)
\end{array}
$$

On désigne par $\mathcal{G}_{\rho,\sigma}^\#$ l'image inverse de $\Sigma_{\rho,\sigma}^\#$ (cf. [unclear: corr.] p. 540) $\equiv \mathcal{B}_{\rho,\sigma}$. Par construction on a un [unclear: homomorphisme] surjectif $\mathcal{S}\mathcal{G}_{\rho,\sigma}^\# \twoheadrightarrow \mathcal{G}_{\rho,\sigma}^\#$. 

$$
\mathcal{G}_{\rho,\sigma}^\# \subset \mathcal{G}_{\rho,\sigma} \times \mathcal{G}_{\rho,\sigma} \times \hat{\mathbf{Z}}^\times \leqno{(47)}
$$
$$
= \left\{ (\alpha, \beta, \mu) \mid \begin{array}{l} 
\chi(\alpha) = \varepsilon_3(\mu) \\
\alpha \rho(\alpha)^{-1} = \varepsilon_2(\mu) \beta \sigma(\beta)^{-1} l_1^\nu \\
= \alpha (\rho \alpha^{-1} \rho^{-1}) \qquad = \varepsilon_2(\mu) \beta (\sigma \beta^{-1} \sigma^{-1}) 
\end{array} \right\}
$$

de sorte que l'on a
$$
\mathcal{S}\mathcal{G}_{\rho,\sigma}^\# = \mathcal{G}_{\rho,\sigma}^\# \cap (\mathfrak{S}' \times \mathfrak{S} \times \hat{\mathbf{Z}}^\times) \leqno{(48)}
$$
$$
= \left\{ (\alpha, \beta, \mu) \in \mathcal{G}_{\rho,\sigma}^\# \mid \begin{array}{l} \delta(\beta) = 1 \\ \delta_\rho(\alpha \tau^{\nu_3(\mu)}) = 1 \end{array} \right\}
$$

\newpage

%% ===== Page 543 =====

et on a encore une application canonique

(49)
$$ g_{\rho,\sigma}^\# \longrightarrow \Sigma_{\rho,\sigma} $$
$$ (\alpha, \beta, \mu) \longmapsto ([\alpha, \rho], \varepsilon_2(\mu)[\beta, \sigma], \mu) \ . $$

On note $g_{\rho,\sigma}^{\sharp\sharp}$ l'image inverse de $\Sigma_{\rho,\sigma}^{eff}$ , et on trouve

(50)
$$ Sg_{\rho,\sigma}^{\sharp\sharp} = g_{\rho,\sigma}^{\sharp\sharp} \cap (\underline{\mathfrak{S}}' \times \underline{\mathfrak{S}} \times \hat{\mathbb{Z}}^*) $$
$$ = \left\{ (\alpha, \beta, \mu) \in g_{\rho,\sigma}^{\sharp\sharp} \mid \begin{matrix} \delta_\sigma(\beta) = 1 \\ \delta_\rho(\alpha \tau^{\nu_3(\mu)}) = 1 \end{matrix} \right\} $$

On a une section de

(51)
\begin{tikzcd}
g_{\rho,\sigma}^{\sharp\sharp} \ar[r] & \Sigma_{\rho,\sigma}^{eff} \\
& B_{\rho,\sigma} \ar[u, "\simeq"'] \ar[ul, "section"]
\end{tikzcd}

[unclear: marginal text to the right of diagram]

donnée par

(52)
$$ u \longmapsto (\varepsilon u, u \tau^{\nu_2(\mu)}, \mu = \chi_D(u)) $$

puisqu'on a (posant $\varepsilon = \varepsilon_2(\mu)$ [unclear])

$$ [u,\rho] = u \rho u^{-1} \rho^{-1} = \xi $$
$$ \varepsilon [u\tau^{\nu_2(\mu)}, \sigma] = \varepsilon u \tau^{\nu_2} \sigma (u\tau^{\nu_2})^{-1} \sigma^{-1} = \varepsilon u (\tau^{\nu_2} \sigma \tau^{-\nu_2} \sigma^{-1}) \sigma(u^{-1}) \sigma^{-1} = \varepsilon u \sigma(u^{-1}) = \eta $$

OK. Donc à $(\alpha, \beta, \mu)$ correspond dans $\Sigma_{\rho,\sigma}^{eff}$

$$ (y_1, y_2, \mu) = ([u,\rho], [u,\sigma], \mu) $$ exactement [unclear]

\newpage

%% ===== Page 544 =====

$G_{\rho, \sigma}(\mu) =$
$$
\left\{ ( \overbrace{u a_0^{-1}}^{\alpha}, \overbrace{u \tau^{1/2} b_0^{-1}}^{\beta}, \mu ) \mid a_0 \in Z_\rho, b_0 \in Z_\sigma, \nu_2 = \nu_2(\mu) \right\} \leqno{(53)}
$$
\begin{center}
$(= (u a_0^{-1}, u b^{-1}, \mu))$ avec $b = b_0 \tau^{-1/2}$
\end{center}
qu'on écrit aussi
$$
G_{\rho, \sigma}(\mu) = \left\{ ( \overbrace{u a_0^{-1}}^{\alpha}, \overbrace{u i_\sigma(b)}^{\beta}, \mu ) \mid a_0 \in Z_\rho, b \in N_\sigma, \underbrace{\varepsilon_\sigma(b) = \varepsilon_2(\mu)}_{\text{cf } (37) (38)} \right\} \leqno{(54)}
$$
où $\varepsilon_\sigma : N_\sigma \to \{\pm 1\}$ est défini sur (37) comme l'hom.\ can.\ de la suite exacte can.
$$
1 \to Z_\sigma \to N_\sigma \xrightarrow{\varepsilon_\sigma} \underbrace{\{\pm 1\}}_{\simeq \{1, \tau\}} \to 1 \leqno{(55)}
$$
Si on cherche les éléments de $SG_{\rho, \sigma}^{\text{eff}}$ au-dessus de $(u, \eta, \mu)$, il faut de plus imposer les conditions
$$
u a_0^{-1} \tau^{1/2} \in \widehat{G} , \quad u i_\sigma(b) \in \widehat{G} \leqno{(56)}
$$
i.e.
$$
\delta_G(u a_0^{-1} \tau^{1/2}) = 1 , \quad \delta_G(u i_\sigma(b)) = 1
$$
soit encore
$$
\delta_G(u) = \underbrace{\delta_\sigma}_{\text{cf } (43)}(b) = \underbrace{\delta_\rho}_{\text{cf } (43)}(\tau^{1/2} a_0) \leqno{(57)}
$$
Posons donc $a = \tau^{1/2} \cdot a_0$ dans $N_\rho$ (donc $a_0 = \tau^{-1/2} a$) on voit que
$$
G_{\rho, \sigma}(u, \mu) = \left\{ (u a^{-1} \tau^{1/2}, u b^{-1}, \mu) \;\middle|\; \substack{a \in N_\rho, b \in N_\sigma \text{ tels que} \\ \varepsilon_\rho(a) = \varepsilon_1(\mu), \varepsilon_\sigma(b) = \varepsilon_2(\mu) \\ \delta_\rho(a) = \delta_\sigma(b) = \delta_G(u)} \right\} \leqno{(58)}
$$
Introduisons donc le groupe
$$
M_{\rho, \sigma} = N_\rho \times_{\widehat{\Gamma}} N_\sigma \leqno{(59)}
$$
et les homomorphismes composés
$$
\begin{tikzcd}[row sep=tiny]
M_{\rho, \sigma} \arrow[r] & N_\rho \arrow[r, "\varepsilon_\rho"] & \{\pm 1\} \\
M_{\rho, \sigma} \arrow[r] & N_\sigma \arrow[r, "\varepsilon_\sigma"] & \{\pm 1\}
\end{tikzcd} \leqno{(60)}
$$
(encore notés $\varepsilon_\rho, \varepsilon_\sigma$ par abus de notation)

\newpage

%% ===== Page 545 =====

et
$$
\mathcal{M}_{\rho, [\sigma]} \xrightarrow{\delta_{\mathcal{M}}} \Gamma_{\rho,\sigma} \xrightarrow{\chi_{\Gamma}} \hat{\mathbf{Z}}^* \quad \text{soit } \chi_{\mathcal{M}} \quad (\text{ou simplement } \chi)
\leqno{(57)}
$$

considérons
$$
\mathcal{M}_{\rho, [\sigma]} \xrightarrow{(\varepsilon_{\rho}, \varepsilon_{\sigma})} \{\pm 1\} \times \{\pm 1\}
\leqno{(58)}
$$
et
$$
(\varepsilon_{2,3} : x \mapsto (\varepsilon_2(\chi(x)), \varepsilon_3(\chi(x)))
$$
$$
\mathcal{M}_{\rho, [\sigma]} \to \{\pm 1\} \times \{\pm 1\}
\leqno{(59)}
$$

et considérons le sous-groupe de coïncidence des deux hom. précédents
$$
\mathcal{M}_{\rho, [\sigma]}^{\text{eff}} \defeq \{ x \in \mathcal{M}_{\rho, [\sigma]} \mid \varepsilon_{\rho,\sigma}(x) = \varepsilon_{2,3}(x) \}
\leqno{(60)}
$$
$$
= \left\{ (u,a,b) \in \mathcal{B}_{\rho,\sigma} \times \mathcal{N}_{\rho} \times \mathcal{N}_{\sigma} \;\middle|\; \begin{array}{l} \delta_{\mathcal{B}}(u) = \delta_{\rho}(a) = \delta_{\sigma}(b) \text{ soit } \Delta \in \Gamma_{\rho,\sigma} \\ \varepsilon_{\rho}(a) = \varepsilon_2(\chi_{\Gamma}(\Delta)), \varepsilon_{\sigma}(b) = \varepsilon_3(\chi_{\Gamma}(\Delta)) \end{array} \right\}
$$

On a une application naturelle
$$
\mathcal{M}_{\rho, [\sigma]}^{\text{eff}} \to \mathfrak{S}' \times \mathfrak{S} \times \hat{\mathbf{Z}}^*
\leqno{(61)}
$$
$$
(u,a,b) \mapsto (\alpha, \beta, \mu), \quad \text{où } \alpha = u a^{-1} \tau^{\prime \nu_{\rho}(a)}\footnote{NB On a posé $\nu_{\rho}(a) \in \hat{\mathbf{Z}}/2\hat{\mathbf{Z}}$ tel que $(-1)^{\nu_{\rho}(a)} = \varepsilon_{\rho}(a)$. On identifie $a \in \mathcal{N}_{\rho}, b \in \mathcal{N}_{\sigma}$ avec $\delta_{\mathcal{B}}(u) = \delta_{\rho}(a) = \delta_{\sigma}(b)$ dans $\Gamma_{\rho,\sigma}$.} \quad \beta = u b^{-1}
$$

[unclear: Le] fait [unclear: ce] qu'il [unclear: a] fallu poser pour que celle-ci soit [unclear: légitime] et ait comme image $\mathcal{S}\mathcal{G}_{\rho, \sigma}^{\text{eff}}$, donne une bijection
$$
\mathcal{M}_{\rho, [\sigma]}^{\text{eff}} \xrightarrow{\sim} \mathcal{S}\mathcal{G}_{\rho, \sigma}^{\text{eff}}
\leqno{(62)}
$$

En fait (61) définissait une application
$$
\mathcal{M}_{\rho, [\sigma]} \to \mathfrak{S}' \times \mathfrak{S} \times \hat{\mathbf{Z}}^*
\leqno{(63)}
$$
et pour un $(u, a, b) \in \mathcal{M}_{\rho, [\sigma]}$, son image $(\alpha, \beta, \mu)$ est liée à $u \in \mathcal{B}_{\rho, \sigma}$ par les formules

\newpage

%% ===== Page 546 =====

$$
\begin{cases}
u(\rho) = \operatorname{int}(\alpha)(\rho) & (\text{à conjuguer : } \rho^{\varepsilon'}, \text{ où } \varepsilon' = \varepsilon_{\rho}(a)) \\
u(\sigma) = \varepsilon \operatorname{int}(\beta)(\sigma) & (\text{à conjuguer : } \sigma^{\varepsilon}, \text{ où } \varepsilon = \varepsilon_{\sigma}(b))
\end{cases} \leqno (63)
$$
(mais sans qu'on ait nécessairement $\varepsilon' = \varepsilon_3(\mu), \varepsilon = \varepsilon_2(\mu)$).

Ceci suggère une définition d'un groupe un peu plus gros que $B_{\rho,\sigma}$, dans lequel $B_{\rho,\sigma}$ sera comme sous-groupe d'indice 1, 2, ou 4 - et qui aura l'avantage, dans les "bons cas", d'être le groupe des points rationnels d'un schéma en groupes [unclear: lisse] sur $\mathbf{Z}$. D'autre part, pour éliminer le signe $\varepsilon$ dans (63), (qui se présente si $(u,a,b) \in \mathcal{M}_{\rho,\sigma}[0]^{\text{eff}}$), l'astuce de remplacer $\beta$ par $\beta' = \tau'^{\nu} \beta$ (où $\nu \in \hat{\mathbf{Z}}/2\mathbf{Z}$ est tel que $(-1)^{\nu} = \varepsilon$)
$$
\beta' = \tau'^{\nu} \beta \leqno (64)
$$
de sorte qu'on aura
$$
u(\sigma) = \operatorname{int}(\beta')(\sigma) \quad \text{i.e. } \quad [u, \sigma] = \beta' \sigma (\beta')^{-1} \ . \leqno (65)
$$

Ceci est lié à la convention qui nous a été [unclear: utile], que nous répétons pour le couple $(\alpha, \beta) \in \hat{\Pi}_{0,3} \times \hat{\Pi}_{0,3}$, remplacée désormais opérer par le couple $(\alpha, \beta') \in \hat{\Pi}_{0,3} \times \hat{\Pi}_{0,3}$, avec
$$
\chi(\alpha) = \varepsilon_3(\mu), \chi(\beta') = \varepsilon_2(\mu) \quad (\text{où } \mu = \chi(u) \text{ par le multiplicateur}) \leqno (66)
$$
ce qui donne pour l'opération de $u$ sur 
$$
\mathfrak{S}_{0,3}^{+\wedge} = \operatorname{Sl}(2, \hat{\mathbf{Z}})^{(1)} = \operatorname{Sl}(2, \hat{\mathbf{Z}})/ \pm 1
$$
un [unclear: dévissage pour en tirer] les formules partantes :
$$
\begin{cases}
u(\rho) = \operatorname{int}(\alpha) \cdot \rho, \quad u(\sigma) = \operatorname{int}(\beta') \cdot \sigma \quad (= \alpha \rho \alpha^{-1}, \beta' \sigma (\beta')^{-1}) \\
\alpha \rho \alpha^{-1} \cdot \beta' \sigma (\beta')^{-1} \dots \quad (= \beta' \sigma \tau (\beta')^{-1}) \quad \nu = \frac{\mu-1}{2}
\end{cases} \leqno (67)
$$

\newpage

%% ===== Page 547 =====

mais pour l'opération dans
$\mathfrak{T}_{1,1}^{+\wedge} \simeq \operatorname{Sl}(2, \hat{\mathbf{Z}})^{\wedge}$ lui-même
donne des formules plus simples, à savoir,-
au lieu de
$$
\begin{cases}
u(\rho) = \operatorname{int}(\alpha)\rho \quad u(\sigma) = \varepsilon \cdot \operatorname{int}(\beta)\sigma \\
\alpha \rho (\alpha)^{-1} = \varepsilon \beta \sigma (\beta)^{-1} l_1^\nu
\end{cases} \leqno (68)
$$
les formules plus symétriques
$$
\begin{cases}
u(\rho) = \operatorname{int}(\alpha)\rho, \quad u(\sigma) = \operatorname{int}(\beta)\sigma \\
\alpha \rho (\alpha)^{-1} = \beta \sigma (\beta)^{-1} l_1^\nu
\end{cases} \leqno (69)
$$
où le signe $\varepsilon = \varepsilon_2(\mu)$ disparait. Il
apparait dans les formules
$$
\chi(\alpha) = \varepsilon_3(\mu), \quad \chi(\beta) = \varepsilon_2(\mu) . \leqno (70)
$$

Il me semble en tous cas qu'on pourra
reprendre "à main levée" les
constructions précédentes et les achever avec
des notations plus adéquates, partant de $\mathfrak{C}_{1,1}^\wedge \simeq \operatorname{Gl}(2, \hat{\mathbf{Z}})^\wedge$.
Le groupe $\mathfrak{S}_{1,1}^\wedge$ joue d'ailleurs une
importance croissante dans nos calculs, c'est
lui qui mérite [unclear: vraiment la] notation $\mathfrak{S}_{1,1}^\wedge$,
ou $\mathfrak{S}$, le sous-groupe fermé engendré par $\rho, \sigma$
devenant $\mathfrak{S}^+$ -- il est isomorphe à $\mathfrak{T}_{1,1}^{+\wedge} \simeq \operatorname{Sl}(2, \hat{\mathbf{Z}})^{\wedge}$.
On suppose [unclear: d'ailleurs] que $\tau \notin \mathfrak{S}^+$ (c.à.d. l'image
de $\tau_\infty$), ce qui [unclear: veut dire], que
$$
\tau' \notin \mathfrak{S}^+ \leqno (71)
$$

\newpage

%% ===== Page 548 =====

C'est p.-e. ce relèvement $\tau'$ lui-même\footnote{Néanmoins $\tau'$ a sur $\tau$ l'avantage sérieux d'être un "vrai" élément d'ordre 2 de $\hat{\mathfrak{T}}_{0,3} \simeq \hat{\mathfrak{T}}_{1,1}$. On a la relation : $\tau'\tau \equiv 1$ (mod le centre de $\hat{\mathfrak{T}}_{0,3} \simeq \hat{\mathbf{Z}}$), ou encore $\tau = \tau' \dots$, ce qui détruit [unclear: d'un coup la] rigidité de $\tau$. -- Alors que de l'autre côté on a [unclear: construit] $\tau'$, mais qui ne relève pas de $\tau$ ! De plus $\tau'$ a le désavantage de [unclear: mélanger] les $M_i^o$ $\dots$}, qui normalise à la fois $L_\rho$ et $L_\sigma$ par les formules remarquables
$$
\tau'(\rho) = \rho^{-1}, \quad \tau'(\sigma) = \sigma^{-1} \leqno (*)
$$
que ne possède pas un $\tau$, qui satisfait à
$$
\tau(\varepsilon_0) = \varepsilon_0^{-1}, \quad \tau(\varepsilon_1) = \varepsilon_1^{-1} \quad \text{et même } \tau(\sigma) = \sigma^{-1} \leqno (40)
$$
(qui sont évidemment très vertueux ! ),
d'où l'action sur $\rho$
$$
\tau(\rho) = \sigma(\rho^{-1})
$$
est moins commode.

\marginpar{\underline{Données} \underline{préliminaires}}
10) Soit donc $\mathcal{G}$ un groupe profini, muni d'un homomorphisme surjectif
$$
\hat{\mathfrak{T}}_{1,1}^\wedge \simeq \operatorname{Gl}(2,\hat{\mathbf{Z}})^\wedge \xrightarrow{\text{épi.}} \mathcal{G} \leqno (1)
$$
ce que nous exprimerons par la donnée de trois générateurs topologiques $\rho, \sigma, \tau' \in \mathcal{G}$, satisfaisant
$$
\rho^6 = \sigma^4 = {\tau'}^2 = 1, \quad \rho^3 = \sigma^2 \text{ (noté } -1), \quad \tau' \rho {\tau'}^{-1} = \rho^{-1}, \quad \tau' \sigma {\tau'}^{-1} = \sigma^{-1}. \leqno (2)
$$
On désigne par $\mathcal{G}^+$ le sous-groupe fermé engendré par $\rho, \sigma$, qui est donc l'image de
$$
\hat{\mathfrak{T}}_{1,1}^+ \simeq \operatorname{Sl}(2,\hat{\mathbf{Z}})^\wedge
$$

\newpage

%% ===== Page 549 =====

et on suppose que $\mathfrak{S}^+ \neq \mathfrak{S}$, donc
$$ \mathfrak{S} \simeq \underset{(\text{1/2 direct})}{\{1, \tau\} \cdot \mathfrak{S}^+} \leqno{(3)} $$
et on a une suite exacte canonique
$$ 1 \longrightarrow \mathfrak{S}^+ \longrightarrow \mathfrak{S} \overset{\varepsilon}{\longrightarrow} \{\pm 1\} \longrightarrow 1 \leqno{(4)} $$
On désignera aussi par $\nu_0$ la composition
$$
\begin{tikzcd}[row sep=tiny]
\mathfrak{S} \arrow[r, "\varepsilon"] & \{\pm 1\} & \mathbf{Z}/2\mathbf{Z} \arrow[l, "\sim"'] \\
& (-1)^\nu & \nu \arrow[l, maps to]
\end{tikzcd} \leqno{(5)}
$$

De façon générale, on désignera souvent par une même lettre des éléments de $\operatorname{Gl}(2, \mathbf{Z})$ - tels $\varepsilon_0, \varepsilon_1, \tau=\tau_\infty, \tau'=\tau'_\infty$, $\sigma=\sigma_\infty, \rho=\rho_\infty$ etc - et leurs images dans $\mathfrak{S}$.

\marginpar{les groupes $Z_{\rho,\sigma}$, $Z_{\rho,\tau}$}
2º) Soit
$$ Z_{\rho,\sigma} \subset \operatorname{Aut}(\mathfrak{S}^+) \times \hat{\mathbf{Z}}^\times \times \{\pm 1\} \times \{\pm 1\} \leqno{(6)} $$
$$ \defeq \left\{ (u, \mu, \varepsilon, \varepsilon') \;\middle|\; \begin{array}{l} u(\varepsilon_0) = \varepsilon_0^\mu \\ u(\rho) \text{ conjugué dans } \mathfrak{S}^+ \text{ à } \rho^{\varepsilon'} \\ u(\sigma) \hspace{1.5cm} \text{''} \hspace{1.5cm} \text{ à } \sigma^\varepsilon \end{array} \right\} $$

C'est un sous-groupe du groupe produit.

Posons
$$ \begin{cases} L_\rho = \text{sous-groupe de } \mathfrak{S} \text{ engendré par } \rho \\ L_\sigma = \hspace{1.5cm} \text{''} \hspace{1.5cm} \text{''} \hspace{1.5cm} \sigma \end{cases} \leqno{(7)} $$
de sorte qu'on a des épimorphismes
$$ \begin{cases} \mathbf{Z}/6\mathbf{Z} \overset{\text{épi}}{\longrightarrow} L_\rho \qquad n \mapsto \rho^n \\ \mathbf{Z}/4\mathbf{Z} \longrightarrow L_\sigma \qquad n \mapsto \sigma^n \end{cases} \leqno{(8)} $$

\newpage

%% ===== Page 550 =====

un diagramme commutatif
et des homomorphismes

(9)
\begin{tikzcd}
& L_\rho \ar[dr] & \\
\{\pm 1\} \ar[ur] \ar[dr] & & \check{\mathbb{G}}^+ \\
& L_\sigma \ar[ur] &
\end{tikzcd}

- Si $\check{\mathbb{G}}$ est "assez gros", on aura d'ailleurs
que $\{\pm 1\} \hookrightarrow \check{\mathbb{G}}^+$ i.e. "$-1$" $\neq \mathbb{1}$ dans $\check{\mathbb{G}}^+$, et
même
$$L_\rho \cap L_\sigma = \{\pm 1\}$$

Les conditions (6) impliquent donc que

(10)
$$
\begin{cases}
u(L_\rho) \text{ } \check{\mathbb{G}}^+ \text{-conjugué à } L_\rho \\
u(L_\sigma) \text{ } \check{\mathbb{G}}^+ \text{-conjugué à } L_\sigma
\end{cases}
$$

Si $\check{\mathbb{G}}^+$ est "assez gros" pour que $\rho$ et $\rho^{-1}$, $\sigma$ et $\sigma^{-1}$
n'y soient pas conjugués, alors un élément
$(u,\mu,\varepsilon,\varepsilon')$ de $Z_{\rho,\sigma}$ est connu par la seule
connaissance de $(u,\mu)$. Si d'autre part l'hom.

(11)
\begin{tikzcd}
\hat{\mathbb{Z}} \ar[r] & \check{\mathbb{G}}^+ \\
1 \ar[r, mapsto] & \varepsilon_0
\end{tikzcd}

est injectif, alors $\mu$ est connu quand on connait
$u$. Donc si $\check{\mathbb{G}}^+$ est assez gros pour satisfaire
aux trois conditions envisagées (p.ex. pour
qu'on ait un "homomorphisme de bases"

(12)
$$\check{\mathbb{G}} \longrightarrow \operatorname{Gl}(2, \hat{\mathbb{Z}})$$

(transformant $\rho, \sigma, \tau$ en les éléments de ce nom),
alors un élément de $Z_{\rho,\sigma}$ est connu par la seule

\newpage

%% ===== Page 551 =====

connaissance de $u \in \operatorname{Aut}(\mathfrak{S}^+)$. 

On a un homomorphisme
$$
\begin{aligned}
L_0 &\xrightarrow{i_{L_0,Z}} Z_{\rho,\sigma} \\
l &\longmapsto (\operatorname{int}(l), 1, 1, 1)
\end{aligned} \leqno{(13)}
$$
dont le noyau est manifestement l'intersection $L_0 \cap \operatorname{Centr}(\mathfrak{S}^+)$. Notons que l'on a 
$$
L_0 \cap \operatorname{Centr}(\mathfrak{S}^+) \subset L_0 \cap \operatorname{Centr}_{\mathfrak{S}^+}(\rho) = L_0 \cap L_1 \leqno{(14)}
$$
car si $l \in L_0$ et $\rho$ commutent, i.e.\ $l = \rho(l)$, comme $\rho(L_0) = L_1$, on a $l \in L_1$. Donc si $L_0 \cap L_1 = \{\pm 1\}$ (ce qui est le cas si (12) existe), et même s'il existe seulement 
$$
\mathfrak{S}^+ \to \operatorname{Gl}(2, \hat{\mathbf{Z}})' \quad \text{[unclear: modulaire]} \leqno{(11')}
$$
alors on a 
$$
L_0 \cap \operatorname{Centr}(\mathfrak{S}^+) = \{ 1 \} \leqno{(15)}
$$
et (13) sera \emph{injectif}. Nous allons le supposer, pour fixer les idées.

On pose
$$
\gamma_{\rho,\sigma} = Z_{\rho,\sigma} / \operatorname{Im} L_0 \leqno{(16)}
$$
de sorte qu'on aura une suite exacte canonique
$$
1 \to L_0 \xrightarrow{i_{L_0,Z}} Z_{\rho,\sigma} \xrightarrow{p_Z} \gamma_{\rho,\sigma} \to 1 \leqno{(17)}
$$
Les homomorphismes canoniques de projection $(u,\mu,\xi,\xi') \mapsto \mu, \xi, \xi'$ se factorisent par $\gamma_{\rho,\sigma}$ et fournissent des homomorphismes canoniques

\newpage

%% ===== Page 552 =====

$$
\begin{cases}
\gamma_{\rho,\sigma} \xrightarrow{\chi_\gamma} \hat{\mathbf{Z}}^* \\
\gamma_{\rho,\sigma} \xrightarrow{\varepsilon_\gamma} \{\pm 1\} \\
\gamma_{\rho,\sigma} \xrightarrow{\varepsilon'_\gamma} \{\pm 1\}
\end{cases}
\quad \text{soit} \quad \gamma_{\rho,\sigma} \xrightarrow{(\chi_\gamma, \varepsilon_\gamma, \varepsilon'_\gamma)} \hat{\mathbf{Z}}^* \times \{\pm 1\} \times \{\pm 1\} \leqno (18)
$$

les composés de ceux-ci avec $Z_{\rho,\sigma} \to \gamma_{\rho,\sigma}$
sont notés
$$
\chi_Z = \chi_\gamma \delta_Z, \quad \varepsilon_Z = \varepsilon_\gamma \delta_Z, \quad \varepsilon'_Z = \varepsilon'_\gamma \delta_Z \leqno (19)
$$
de sorte que
$$
\chi_Z(u, \mu, \varepsilon, \varepsilon') = \mu, \quad \varepsilon_Z(u, \mu, \varepsilon, \varepsilon') = \varepsilon, \quad \varepsilon'_Z(u, \mu, \varepsilon, \varepsilon') = \varepsilon' \leqno (20)
$$

On aurait envie de désigner par $\gamma_{\rho,\sigma}^\circ$ le
noyau de $(\varepsilon_\gamma, \varepsilon'_\gamma)$, par $\gamma_{\rho,\sigma}^+$ celui de $\chi_\gamma$,
par $\gamma_{\rho,\sigma}^{\circ+}$ leur intersection, et d'introduire
par les notations $Z_{\rho,\sigma}^\circ$, $Z_{\rho,\sigma}^+$, $Z_{\rho,\sigma}^{\circ+}$ leurs
images inverses dans $Z_{\rho,\sigma}$, i.e.\ les noyaux resp.\
de $(\varepsilon_Z, \varepsilon'_Z)$, de $\chi_Z$ et de $(\chi_Z, \varepsilon_Z, \varepsilon'_Z)$.

\marginpar{Les groupes $\mathfrak{S}_{G,\rho,\sigma}$, $G_{\rho,\sigma}$}

3º)
$$
Z_{\rho,\sigma} \longrightarrow \operatorname{Aut}(\mathfrak{S}^+) \leqno (21)
$$
$$
(u, \mu, \varepsilon, \varepsilon') \longmapsto u
$$
donc $Z_{\rho,\sigma}$ opère sur $\mathfrak{S}^+$, et on peut
donc considérer le produit 1/2 direct
$$
\mathfrak{S}_{G,\rho,\sigma} = Z_{\rho,\sigma} \underset{(1/2\text{-direct})}{\cdot} \mathfrak{S}^+ \leqno (22)
$$

\newpage

%% ===== Page 553 =====

On a un hom.
(23)
$$L_0 \xrightarrow{i_{L_0, SG}} S_0 G_{\rho,\sigma} = Z_{\rho,\sigma} \cdot \widehat{C}^+$$
$$l \longmapsto i_{L_0,Z}(l) \cdot l^{-1}$$
dont l'image est encore invariante, d'où
(24)
$$G_{\rho,\sigma} = S_0 G_{\rho,\sigma} / \text{Im } i_{L_0, SG} ,$$
de sorte qu'on a une suite exacte
(25)
$$1 \longrightarrow L_0 \xrightarrow{i_{L_0, SG}} S_0 G_{\rho,\sigma} \longrightarrow G_{\rho,\sigma} \longrightarrow 1$$
et un diagramme (d'homomorphismes)
(26)
\begin{tikzcd}
& Z_{\rho,\sigma} \ar[dr, "i_Z"] \\
L_0 \ar[ur, "i_{L_0,Z}"] \ar[dr, "i_{L_0,\widehat{C}}"] & & G_{\rho,\sigma} \\
& \widehat{C}^+ \ar[ur, "i_{\widehat{C}^+}"']
\end{tikzcd}
( a priori $\operatorname{Ker}(\widehat{C}^+ \to G_{\rho,\sigma}) \simeq \operatorname{Ker}(L_0 \to Z_{\rho,\sigma})$
$\simeq L_0 \cap \text{Cent}(\widehat{C}^+)$, mais par hyp. (15)
il se réduit à $\{1\}$ ), s'insérant dans un
diagramme de suites exactes
(27)
\begin{tikzcd}
1 \ar[r] & L_0 \ar[r, "i_{L_0,Z}"] \ar[d, "i_{L_0,\widehat{C}}"'] & Z_{\rho,\sigma} \ar[r, "\delta_Z"] \ar[d] & \Gamma_{\rho,\sigma} \ar[r] \ar[d, equal] & 1 \\
1 \ar[r] & \widehat{C}^+ \ar[r, "i_{\widehat{C}^+}"'] & G_{\rho,\sigma} \ar[r, "\delta_G"'] & \Gamma_{\rho,\sigma} \ar[r] & 1
\end{tikzcd}
On posera
(28)
$$\chi_G = \chi_\Gamma \delta_G , \varepsilon_G = \varepsilon_\Gamma \delta_G , \varepsilon'_G = \varepsilon'_\Gamma \delta_G$$

\newpage

%% ===== Page 554 =====

d'où un homomorphisme
$$
G_{p,\sigma} \xrightarrow{(\chi_G, \varepsilon_G, \varepsilon'_G)} \hat{\mathbf{Z}}^* \times \{\pm 1\} \times \{\pm 1\} \leqno (29)
$$
nous conduisant à introduire $G_{p,\sigma}^+, G_{p,\sigma}^-$, comme d'habitude.

L'injection $i_{\mathfrak{S}^+} : \mathfrak{S}^+ \hookrightarrow G_{p,\sigma}$ se prolonge en une [unclear: plongeance rigide] (notée $i_{\mathfrak{S}}$)
$$
i_{\mathfrak{S}} : \mathfrak{S} \hookrightarrow G_{p,\sigma} \leqno (30)
$$
que je définis par sa valeur sur $\tau = \sigma \tau'$, par
$$
i_{\mathfrak{S}}(\tau) = i_Z(\tau_Z) \footnote{où $\tau_Z$ désigne l'image de $\tau = \sigma \tau' \in \mathfrak{S} \setminus \mathfrak{S}^+$ dans $G_{p,\sigma}$ [unclear: qui s'identifie à l'élément d'ordre 2 de $\Gamma_{p,\sigma}$ d'image $+1, \varepsilon_{\Gamma_0}$ dans le groupe $\hat{\mathbf{Z}}^* \times \{\pm 1\} \times \{\pm 1\}$]} \leqno (31)
$$
$$
\begin{cases} \tau_Z \in Z_{p,\sigma} \text{ défini comme} \\ \tau_Z = (\operatorname{int}(\tau) \mid \mathfrak{S}^+ , -1 , -1 , -1 ) \\ \in \operatorname{Aut}(\mathfrak{S}^+) \times \hat{\mathbf{Z}}^* \times \{\pm 1\} \times \{\pm 1\} \end{cases} \leqno (32)
$$

Désormais\footnote{Noter que $\tau'_G \notin Z_{p,\sigma}$, par contre $\tau_G \in Z_{p,\sigma}$.} nous identifions $\mathfrak{S}$ à un sous-groupe de $G_{p,\sigma}$, et nous désignons par la même lettre des éléments de $\mathfrak{S}$ (tels que $\tau', \sigma, \tau_0, S, \dots$) et leurs images dans $G_{p,\sigma}$. On aura d'ailleurs
$$
\chi_G | \mathfrak{S} = \varepsilon_G | \mathfrak{S} = \varepsilon'_G | \mathfrak{S} = \varepsilon_{\mathfrak{S}} : \begin{cases} \text{trivial sur } \mathfrak{S}^+ \\ -1 \text{ sur } \mathfrak{S} \setminus \mathfrak{S}^+ \end{cases} \leqno (32 \text{ bis})
$$

\marginpar{Les groupes $Z_p, Z_\sigma, N_p, N_\sigma$, $S N_p, S N_\sigma \dots$}
4º) On [unclear: constate] que l'élément $\tau'$ de $G_{p,\sigma}$ normalise $L_p$ et $L_\sigma$, donc normalise $\dots$
$$
\begin{cases} Z_p \defeq \operatorname{Centr}_{G_{p,\sigma}}(p) = \operatorname{Centr}_{G_{p,\sigma}}(L_p) \\ Z_\sigma \defeq \operatorname{Centr}_{G_{p,\sigma}}(\sigma) = \operatorname{Centr}_{G_{p,\sigma}}(L_\sigma) \end{cases} \leqno (33)
$$

\newpage

%% ===== Page 555 =====

dans $\{1,\tau'\}$ opère sur ces groupes, et on
peut considérer

(34)
$$
\begin{cases}
N_\rho = \{1,\tau'\}. Z_\rho \\
N_\sigma = \{1,\tau'\}. Z_\sigma
\end{cases} \quad \text{(produits semi-directs).}
$$

On a des suites exactes canoniques

(35)
$$1 \to Z_\rho \to N_\rho \xrightarrow{\varepsilon_\rho} \{\pm 1\} \to 1$$
$$1 \to Z_\sigma \to N_\sigma \xrightarrow{\varepsilon_\sigma} \{\pm 1\} \to 1$$

et un diagramme commutatif d'hom.
canoniques

(36)
\begin{tikzcd}
& D_\rho \ar[r] & N_\rho \ar[dr, "i_\rho"] & \\
\{1,\tau'\} \times \{\pm 1\} \ar[ur] \ar[dr] & & & G_{\rho,\sigma} \\
& D_\sigma \ar[r] & N_\sigma \ar[ur, "i_\sigma"'] &
\end{tikzcd}

où

(37)
$$
\begin{cases}
D_\rho \simeq D_6 \subset \operatorname{Gl}(2, \mathbb{Z}) & \text{engendré par } \rho, \tau' \\
D_\sigma \simeq D_4 \subset \operatorname{Gl}(2, \mathbb{Z}) & \text{------- } \sigma, \tau' \\
\{1,\tau'\} \times \{\pm 1\} = D_\rho \cap D_\sigma \text{ dans } \operatorname{Gl}(2, \mathbb{Z})
\end{cases}
$$

dans lesquels

(38)
$$
\begin{cases}
D_\rho = \{1,\tau'\}.L_\rho & \text{(où } L_\rho \text{ est d'ordre } 6 \text{ dans } \bar{\mathbb{S}}^+ \text{)} \\
D_\sigma = \{1,\tau'\}.L_\sigma & \text{( ------- } \sigma \text{ d'ordre } 4 \text{ --------)}
\end{cases}
$$

et les hom.

$$D_\rho \to N_\rho, \quad D_\sigma \to N_\sigma$$

sont évidents. Les hom. $i_\rho, i_\sigma$ sont inj., bien sûr.

Rappelons que

(39)
$$
\begin{cases}
i_\rho \text{ i-jectif} \iff \tau' \notin Z_\rho \iff \rho^2 \neq 1 \iff \rho \neq \pm 1 \iff \bar{\mathbb{S}}^+ \text{ est [unclear]} \\
i_\sigma \text{ i-jectif} \iff \tau' \notin Z_\sigma \iff \sigma^2 \neq 1 \iff \text{[unclear]}
\end{cases}
$$

\newpage

%% ===== Page 556 =====

[MARGIN: Si $g \in G_{\rho,\sigma}$ et $\varepsilon = \varepsilon_G(g)$, $g^{-1} \rho g = \rho^\varepsilon$, $int(g)(\rho) = \rho^\varepsilon$ donc $int(g)(\sigma) = \sigma^\varepsilon$
[unclear]
]

Faisant opérer $G_{\rho,\sigma}$ sur $\widehat{\mathbb{C}}^+$, sous-groupe distingué, via autom. int., on voit que l'image de $G_{\rho,\sigma}$ dans le gr. des [unclear] de $\widehat{\mathbb{C}}^+$ est [unclear]
[crossed out text]
(posant
(40) $N_\rho^* = \{ x \in N_\rho \mid \varepsilon_\rho(x) = \varepsilon_G(i_\rho(x)) \}$, $N_\sigma^* = \{ x \in N_\sigma \mid \varepsilon_\sigma(x) = \varepsilon_G(i_\sigma(x)) \}$
les homomorphismes composés

$$ N_\rho^* \xrightarrow{i_\rho} G_{\rho,\sigma} \xrightarrow{\delta_a} \Gamma_{\rho,\sigma} , \delta_\rho^* = \delta_G \circ i_\rho^* $$
$$ N_\sigma^* \xrightarrow{i_\sigma} G_{\rho,\sigma} \xrightarrow{\delta_a} \Gamma_{\rho,\sigma} , \delta_\sigma^* = \delta_G \circ i_\sigma^* $$
sont des épimorphismes.

a) Si $\rho^2 \neq 1$, alors on a commut. [unclear] des
diagrammes de suites exactes
(41)
\begin{tikzcd}
1 \ar[r] & Z_\rho \ar[r] \ar[d, "j_\rho"'] & N_\rho^* \ar[r, "\varepsilon_\rho"] \ar[d, "\delta_\rho^* \text{ épi}"] & \{\pm 1\} \ar[r] \ar[d, equal] & 1 \\
1 \ar[r] & \Gamma_{\rho,\sigma}' \ar[r] & \Gamma_{\rho,\sigma} \ar[r, "\varepsilon_\Gamma'"'] & \{\pm 1\} \ar[r] & 1
\end{tikzcd}
où l'on a posé
(42) $$ \Gamma_{\rho,\sigma}' = \operatorname{Ker}(\Gamma_{\rho,\sigma} \xrightarrow{\varepsilon_\Gamma'} \{\pm 1\}) ; $$

b) Si $\sigma^2 \neq 1$ (cf (39)), alors on a [unclear] diag. comm. de suites exactes
(43)
\begin{tikzcd}
1 \ar[r] & Z_\sigma \ar[r] \ar[d, "j_\sigma"'] & N_\sigma^* \ar[r, "\varepsilon_\sigma"] \ar[d, "\delta_\sigma^* \text{ épi}"] & \{\pm 1\} \ar[r] \ar[d, equal] & 1 \\
1 \ar[r] & \Gamma_{\rho,\sigma}'' \ar[r] & \Gamma_{\rho,\sigma} \ar[r, "\varepsilon_\Gamma''"'] & \{\pm 1\} \ar[r] & 1
\end{tikzcd}
où
(44) $$ \Gamma_{\rho,\sigma}'' = \operatorname{Ker}(\Gamma_{\rho,\sigma} \xrightarrow{\varepsilon_\Gamma''} \{\pm 1\}) . $$

\newpage

%% ===== Page 557 =====

NB Les notations $\varepsilon_\gamma, \varepsilon_{\gamma'} : \Gamma_{\rho,\sigma} \to \{\pm 1\}$
pourraient prêter à confusion, la première
exprimant l'action d'un $u$ sur $L_\sigma$ (non sur
$L_\rho$). Une notation un peu plus lourde,
mais ne prêtant pas à confusion, serait
$\varepsilon_{\gamma,\rho}, \varepsilon_{\gamma,\sigma}$ - ou on écrira seulement $\varepsilon_\rho, \varepsilon_\sigma$,
la confusion avec les hom de ces noms
$N_\rho^* \to \{\pm 1\}$, $N_\sigma^* \to \{\pm 1\}$ étant exclue lorsque
$\rho^2 \neq 1, \sigma^2 \neq 1$, à cause des compatibilités (41) (43).

On a donc construit, à l'aide
de (1) i.e. de $\widehat{\mathcal{L}}$ munie de $\rho, \sigma, \tau'$, des
groupes $C^+, G_{\rho,\sigma}, \Gamma_{\rho,\sigma}, Z_{\rho,\sigma}^*, N_\rho^*, N_\sigma^*$, s'insérant
dans un diagramme commutatif de
groupes

(45)
\begin{tikzcd}
L_0 \simeq S Z_{\rho,\sigma} \ar[r, hook] \ar[d, hook] \ar[dd, hook, bend right=40] & Z_{\rho,\sigma}^* \ar[d, hook] \ar[dd, hook, bend left=40] \ar[dddr, dashed, "\delta_Z"] & & \\
S N_\rho^* \ar[r, hook] \ar[dd, hook, bend right=40] & N_\rho^* \ar[dd, hook, bend left=40] \ar[ddr, dashed, "\delta_\rho"'] & & \hat{\mathbb{Z}}^* \times \{\pm 1\} \times \{\pm 1\} \\
S N_\sigma^* \ar[r, hook] \ar[d, hook] & N_\sigma^* \ar[d, hook] \ar[dr, dashed, "\delta_\sigma"'] & & \\
1 \ar[r] & C^+ \ar[loop below, "G"'] \ar[r, "i^+"'] & G_{\rho,\sigma} \ar[r, "\delta_G"'] & \Gamma_{\rho,\sigma} \ar[r] \ar[uuu, hook, "(\chi{,} \varepsilon_{\gamma,\rho}{,} \varepsilon_{\gamma,\sigma})"'] & 1
\end{tikzcd}

où on a posé

(46)
$$ S Z_{\rho,\sigma} = i_{L_0, Z}(L_0) = \ker(\delta_Z : Z_{\rho,\sigma}^* \to \Gamma_{\rho,\sigma}) $$
$$ S N_\rho^* = \ker(\delta_\rho : N_\rho^* \to \Gamma_{\rho,\sigma}) = \text{Cent}_{C^+}(\rho) $$
$$ S N_\sigma^* = \ker(\delta_\sigma : N_\sigma^* \to \Gamma_{\rho,\sigma}) = \text{Cent}_{C^+}(\sigma) $$

[MARGIN: NB On peut d'ailleurs les noter $S Z_\rho, S Z_\sigma$ au lieu de $S N_\rho^*, S N_\sigma^*$.]

\newpage

%% ===== Page 558 =====

\marginpar{On a -- \\ $\operatorname{Im}(N_{\rho}^* \to \mathfrak{S}^+) = \operatorname{Norm}_{\mathfrak{S}^+}(\{\rho\})$ \\ $\operatorname{Im}(N_{\sigma}^* \to \mathfrak{S}^+) = \operatorname{Norm}_{\mathfrak{S}^+}(\{\sigma\})$ \\ (au moins pour $\rho^2 \neq 1$ resp.\ $\sigma^2 \neq 1$)}%
qui a des suites exactes
$$
\begin{cases} 
1 \to S N_{\rho}^* \to N_{\rho}^* \xrightarrow{\delta_\rho} \Upsilon_{\rho,\sigma} \to 1 \\ 
1 \to S N_{\sigma}^* \to N_{\sigma}^* \xrightarrow{\delta_\sigma} \Upsilon_{\rho,\sigma} \to 1 
\end{cases}
$$
s'envoyant dans la suite exacte
$$
1 \to \mathfrak{S}^+ \to G_{\rho,\sigma} \to \Upsilon_{\rho,\sigma} \to 1 , \leqno (47)
$$
de façon évidente (47). De plus, les homomorphismes $i_\rho: N_{\rho}^* \to G_{\rho,\sigma}$ et $i_\sigma: N_{\sigma}^* \to G_{\rho,\sigma}$ sont des injections (ce qui justifie (cf (39)), ce qui est indiqué dans le diagramme (45) par la ligne $\hookrightarrow$ (signe d'inclusion pointillé).

On introduit aussi
$$
\begin{cases} 
\chi_{\rho} = \chi_\Upsilon \circ \delta_\rho : N_{\rho}^* \to \hat{\mathbf{Z}}^* \\ 
\chi_{\sigma} = \chi_\Upsilon \circ \delta_\sigma : N_{\sigma}^* \to \hat{\mathbf{Z}}^* 
\end{cases} , \leqno{(48)}
$$
dont les noyaux (qui contiennent $S N_{\rho}^*$ resp.\ $S N_{\sigma}^*$, et qu'on se gardera de confondre avec eux) seront notés comme de juste
$$
N_{\rho}^{**}, N_{\sigma}^{**} .
$$

NB Il faudrait encore examiner les composées
$$
\begin{cases} 
N_{\rho}^* \to \Upsilon_{\rho,\sigma} \xrightarrow{\varepsilon_{\Upsilon,\rho}} \{ \pm 1 \} \\ 
N_{\sigma}^* \to \Upsilon_{\rho,\sigma} \xrightarrow{\varepsilon_{\Upsilon,\sigma}} \{ \pm 1 \} 
\end{cases} \leqno{(49)}
$$
on verra que dans les "bons cas", la restriction de ces hom.\ à $Z_\rho$, $Z_\sigma$ respectivement sont [unclear: surjectifs].

\newpage

%% ===== Page 559 =====

583

[MARGIN: Cancelé / En révisant (15) / dans le cas où / une certaine [unclear] / $Norm$ de $(L_\rho)$ / et idem pour $L_\sigma, L_0$]

On pose

(50)
$$
\begin{cases} 
S N_\rho^* = \operatorname{Ker}( N_\rho^* \xrightarrow{\delta_\rho} \gamma_{\rho,\sigma} ) \doteq \operatorname{Norm}_{\hat{\mathbb{Z}}}(L_\rho) \\
S N_\sigma^* = \operatorname{Ker}( N_\sigma^* \xrightarrow{\delta_\sigma} \gamma_{\rho,\sigma} ) = \operatorname{Norm}_{\hat{\mathbb{Z}}}(L_\sigma) \\
S B_{\rho,\sigma} = \operatorname{Ker}( B_{\rho,\sigma} \xrightarrow{\delta_B} \gamma_{\rho,\sigma} ) = [unclear] \cap \operatorname{Norm}_{\hat{\mathbb{Z}}}(L_0) 
\end{cases}
$$

[MARGIN: Les groupes $S_0 M_{\rho,\sigma}, M_{\rho,\sigma}, S_0 \Gamma_{\rho,\sigma}, \Gamma_{\rho,\sigma}$]

$$
\begin{cases} 
S_0 M_{\rho,\sigma} = Z_{\rho,\sigma} \times_{\gamma_{\rho,\sigma}} N_\rho^* \times_{\gamma_{\rho,\sigma}} N_\sigma^* \times_{\gamma_{\rho,\sigma}} C_{\rho,\sigma} \\
M_{\rho,\sigma} = N_\rho^* \times_{\gamma_{\rho,\sigma}} N_\sigma^* \times_{\gamma_{\rho,\sigma}} C_{\rho,\sigma} 
\end{cases}
$$

(51)
$$
\begin{cases} 
S_0 \Gamma_{\rho,\sigma} = Z_{\rho,\sigma} \times_{\gamma_{\rho,\sigma}} N_\rho^* \times_{\gamma_{\rho,\sigma}} N_\sigma^* \\
\Gamma_{\rho,\sigma} = N_\rho^* \times_{\gamma_{\rho,\sigma}} N_\sigma^* 
\end{cases}
$$

d'où un diagramme de suites ex. [unclear]

(52)
\begin{tikzcd}
1 \ar[r] & C^+ \ar[r] \ar[d, equal] & S_0 M_{\rho,\sigma} \ar[r] \ar[d] & S_0 \Gamma_{\rho,\sigma} \ar[r] \ar[d] & 1 \\
1 \ar[r] & C^+ \ar[r] & M_{\rho,\sigma} \ar[r] \ar[d] & \Gamma_{\rho,\sigma} \ar[r] \ar[d] & 1 \\
& & 1 & 1
\end{tikzcd}

et des suites ex.

(53)
$$ 1 \to S Z_{\rho,\sigma} \times S N_\rho^* \times S N_\sigma^* \times C^+ \xrightarrow{\simeq L_0} S_0 M_{\rho,\sigma} \to \gamma_{\rho,\sigma} \to 1 $$

(54)
$$ 1 \to S N_\rho^* \times S N_\sigma^* \times C^+ \to M_{\rho,\sigma} \to \gamma_{\rho,\sigma} \to 1 $$

(55)
$$ 1 \to S Z_{\rho,\sigma} \times S N_\rho^* \times S N_\sigma^* \xrightarrow{\simeq L_0} S_0 \Gamma_{\rho,\sigma} \to \gamma_{\rho,\sigma} \to 1 $$

(56)
$$ 1 \to S N_\rho^* \times S N_\sigma^* \to \Gamma_{\rho,\sigma} \to \gamma_{\rho,\sigma} \to 1 $$

\newpage

%% ===== Page 560 =====

On a aussi

(57) $$S_0 \mathcal{M}_{\rho,\sigma} \simeq \mathcal{M}_{\rho,\sigma} \times_{\gamma_{\rho,\sigma}} \mathbb{Z}_{\rho,\sigma}$$
(58) $$S_0 \Gamma_{\rho,\sigma} \simeq \Gamma_{\rho,\sigma} \times_{\gamma_{\rho,\sigma}} \mathbb{Z}_{\rho,\sigma}$$

i.e. les ext. $S_0 \mathcal{M}_{\rho,\sigma}$, $S_0 \Gamma_{\rho,\sigma}$ de $\mathcal{M}_{\rho,\sigma}, \Gamma_{\rho,\sigma}$
par $L_0$ se déduisent de l'extension $\mathbb{Z}_{\rho,\sigma}$ de
$\gamma_{\rho,\sigma}$ par $L_0$ par image inverse par
$\mathcal{M}_{\rho,\sigma} \to \gamma_{\rho,\sigma}$, $\Gamma_{\rho,\sigma} \to \gamma_{\rho,\sigma}$ - et --

aussi

(59) $$S_0 \mathcal{M}_{\rho,\sigma} \simeq \mathcal{M}_{\rho,\sigma} \times_{\Gamma_{\rho,\sigma}} S_0 \Gamma_{\rho,\sigma}$$

i.e. l'extension $S_0 \mathcal{M}_{\rho,\sigma}$ de $\mathcal{M}_{\rho,\sigma}$ par $L_0$
se déduit de l'extension $S_0 \Gamma_{\rho,\sigma}$ de
$\Gamma_{\rho,\sigma}$ par $L_0$, par image inverse par
$\mathcal{M}_{\rho,\sigma} \to \Gamma_{\rho,\sigma}$.

On introduit un élément

(60) $\tau_M \in S_0 \mathcal{M}_{\rho,\sigma} \quad (\tau_M = (\overset{a}{\tau'_Z}, \overset{b}{\tau'_N}, \overset{a}{\tau''_N}, \overset{b}{\tau'_G}))$

[unclear: et isomorphie entre les groupes introduits, je vais les indiquer ci-dessous]

où $\tau'_Z, \tau'_N, \tau''_N, \tau'_G$
"relevés" de $\tau'$, pris resp. dans $\mathbb{Z}_{\rho,\sigma}$, $N_\rho^*, N_\sigma^*, C_{\rho,\sigma}$. Ceux-
ci ont bien tous la même image dans $\gamma_{\rho,\sigma}$ et [unclear]
des $\gamma_{\rho,\sigma}$, donc définissent un él. de $S_0 \mathcal{M}_{\rho,\sigma}$, dont
l'opération sur $\widetilde{C}^+$ est celle de $\tau$. D'où un plon-
gement

(61) $\begin{tikzcd} \widetilde{C} \ar[r, hook] & S_0 \mathcal{M}_{\rho,\sigma} \end{tikzcd} \quad \begin{matrix} : \text{identifiant } \widetilde{C}^+ \\ \tau \mapsto \tau_M \end{matrix}$

\newpage

%% ===== Page 561 =====

\[
\begin{tikzcd}[row sep=1.5em, column sep=1.5em]
% Row 1
& L_0 \times S_{N_{\tilde{P}}^*} \times S_{N_\sigma^*} \arrow[rr, "pr"] \arrow[dd, "pr"'] \arrow[dr, twoheadrightarrow] & & L_0 \arrow[dd, twoheadrightarrow] & & \\
% Row 2
& & S_0 \Gamma_{\tilde{P},\sigma} \arrow[rr, "\Theta_Z"] \arrow[dd, twoheadrightarrow] & & \mathbf{Z}_{\tilde{P},\sigma} \arrow[dd, hook] \arrow[dr, "\delta_Z"] & \\
% Row 3
& L_0 \times S_{\{\pm 1\}} \times S_{\{\pm 1\}} \arrow[dr, twoheadrightarrow] \arrow[dd, "pr"'] & & & & \Gamma_{\tilde{P},\sigma} \arrow[r, shift left, "\varepsilon_{\tilde{P},\sigma} = \varepsilon_{\tilde{P}} \cdot \{\pm 1\}"] \arrow[r, shift right, "\varepsilon_{\sigma} = \varepsilon_\sigma \cdot \{\pm 1\}"'] & \hat{\mathbf{Z}}^* \\
% Row 4
& & S_0 M_{\tilde{P},\sigma} \arrow[rr, "\Theta_P" near start] \arrow[dd, twoheadrightarrow] & & N_{\tilde{P}}^* \arrow[dd, twoheadrightarrow] \arrow[ur, "\delta_P"'] & \\
% Row 5
C^+ \arrow[r, hook, "inj"] \arrow[ddrrr, dashed, "\sim"'] & S_{N_{\tilde{P}}^*} \times S_{N_\sigma^*} \times C^+ \arrow[rr, crossing over, "pr" near start] \arrow[dd, "pr"'] \arrow[dr, twoheadrightarrow] & & M_{\tilde{P},\sigma} \arrow[rr, "\Theta_\sigma" near start] \arrow[dd, twoheadrightarrow] & & N_\sigma^* \arrow[dd, twoheadrightarrow] \arrow[ul, "\delta_\sigma"'] & \\
% Row 6
& & & & & \\
% Row 7
& S_{\{\pm 1\}} \times S_{\{\pm 1\}} \times C^+ \arrow[dr, twoheadrightarrow] & & C^+ \arrow[rr, "\Theta_G" near start] & & G_{\tilde{P},\sigma} \arrow[uuul, "\delta_G"'] & \\
% Row 8
& & & & & &
\end{tikzcd}
\leqno{(62)}
\]

On a résumé certaines des relations entre les groupes introduits jusqu'à présent dans le diagramme commutatif (62), que je commente brièvement.

\begin{enumerate}
\item[a)] Toutes les flèches sont soit injectives ($\hookrightarrow$) soit surjectives ($\twoheadrightarrow$), à l'exception des deux flèches issues de $\Gamma_{\tilde{P},\sigma}$.
\item[b)] Les dix carrés (six faces d'un cube, et quatre carrés marqués I, II, III, IV) sont cartésiens.
\item[c)] Les flèches marquées $pr$ sont des projections canoniques de groupes produits sur des produits partiels. La flèche marquée $inj$ : $C^+ \hookrightarrow S_{N_{\tilde{P}}^*} \times S_{N_\sigma^*} \times C^+$ est l'injection canonique.
\item[d)] Pour ne pas détruire la commutativité du diagramme, nous nous sommes dispensés d'y faire figurer les flèches évidentes (inclusions)
$$
N_{\tilde{P},\sigma} \to G_{\tilde{P},\sigma} \quad , \quad N_{\tilde{P}}^* \to G_{\tilde{P},\sigma} \quad , \quad N_\sigma^* \to G_{\tilde{P},\sigma}
$$
qui figurent dans (45).
\end{enumerate}

\newpage

%% ===== Page 562 =====

e) Pour ne pas surcharger le diagramme,
je n'ai pas introduit dans ce diagramme des
notations pour certaines flèches composées utiles,
telles

$$ (63) \left\{ \begin{array}{l} \delta_\Gamma : \Gamma_{\rho,\sigma} \to G_{\rho,\sigma} \\ \delta_{S_0\Gamma} : S_0\Gamma_{\rho,\sigma} \to G_{\rho,\sigma} \\ \delta_M : M_{\rho,\sigma} \to G_{\rho,\sigma} \\ \delta_{S_0M} : S_0M_{\rho,\sigma} \to G_{\rho,\sigma} \end{array} \right. $$

et

$$ (64) \left\{ \begin{array}{l} \chi_\Gamma = \chi_{G_{\rho,\sigma}} \circ \delta_\Gamma : \Gamma_{\rho,\sigma} \to \hat{\mathbb{Z}}^\times \\ \chi_{S_0\Gamma} = \chi_{G_{\rho,\sigma}} \circ \delta_{S_0\Gamma} : S_0\Gamma_{\rho,\sigma} \to \hat{\mathbb{Z}}^\times \\ \chi_M = \chi_{G_{\rho,\sigma}} \circ \delta_M : M_{\rho,\sigma} \to \hat{\mathbb{Z}}^\times \\ \chi_{S_0M} = \chi_{G_{\rho,\sigma}} \circ \delta_{S_0M} : S_0M_{\rho,\sigma} \to \hat{\mathbb{Z}}^\times \end{array} \right. $$

et les composées $\varepsilon_{?,\rho}, \varepsilon_{?,\sigma} : ? \to \{\pm 1\}$, où $?$ peut
prendre toute valeur $Z, G$ (déjà introduits) $N_\rho^*, N_\sigma^*$
(on a aussi $\varepsilon_{N_\rho^*,\rho} = \varepsilon_\rho$ [unclear], $\varepsilon_{N_\sigma^*,\sigma} = \varepsilon_\sigma$ [unclear]),
$\Gamma, S_0\Gamma, M, S_0M$. Le noyau de $\chi_?$ est noté $?^+$
[MARGIN: il y a une petite impropriété de notations car $N_\rho^*, N_\sigma^*$ déjà utilisés dans un tout autre sens]
celui de $\varepsilon_{?,\rho}$ est noté $?'$, celui de $\varepsilon_{?,\sigma}$ est noté $?''$ -
on a $?^\circ = ?' \cap ?''$, $?^{+'} = ?^+ \cap ?'$, $?^{+''} = ?^+ \cap ?''$,
$?^{+\circ} = ?^+ \cap ?^\circ$, $S? = \text{Ker } \delta_?$, et donc

(65

\newpage

%% ===== Page 563 =====

585

Le noyau de $\varepsilon_{\hat{\mathbb{Z}}}$ n'est autre que $\hat{\mathbb{Z}}^+$. En tant
que sous-groupe de $M_{\rho,\sigma}$ (resp. de $\Gamma_{\rho,\sigma}$) $\hat{\mathbb{Z}}^+$
s'identifie à l'image inverse dans ce groupe \textsuperscript{d'un sous-groupe}
du sous-groupe $\{1, \tau'_\rho\}$ de $\bar{\Gamma}_{\rho,\sigma}$ (resp. $\{1, \bar{\tau}'_\rho\}$
de $\bar{\Gamma}_{\rho,\sigma}$), où $\tau'_\rho, \bar{\tau}'_\rho$ désignent les
images de $\tau' \in \widehat{G}$ dans $\bar{\Gamma}_{\rho,\sigma}$ et $\bar{\Gamma}_{\rho,\sigma}$ respectivement.

g) Le diagramme (62) a comme partie
principale

(67)
\begin{tikzcd}
& \bullet \ar[r] \ar[dr] & \bullet \ar[dr] \ar[rr, dashed] & & \\
\bullet \ar[ur] \ar[r] \ar[dr] & \bullet \ar[r] & \bullet \ar[r] \ar[rr, dashed] & \bullet & {[ \ \ ]} \\
& \bullet \ar[ur] \ar[r] & \bullet \ar[ur] \ar[rr, dashed] & &
\end{tikzcd}

Tous les groupes marqués dans ce \textsuperscript{dernier} diagramme, à
l'exception de $\hat{\mathbb{Z}}^*$, $\{\pm 1\}$, $\{\pm 1\}$ (mis entre crochets
$[ \ \ ]$) s'envoient dans $\bar{\Gamma}_{\rho,\sigma}$, qui est le terme
terminal de ce diagramme, relatif à $8$ groupes
(en ne comptant pas \sout{[unclear]} \textsuperscript{l'image inverse par une flèche pointillée}) dont $\bar{\Gamma}_{\rho,\sigma}$, déduits
de quatre d'entre eux $Z_{\rho,\sigma}$, $N_\rho^*$, $N_\sigma^*$, $G_{\rho,\sigma}$
(qui s'envoient dans $\bar{\Gamma}_{\rho,\sigma}$ par des épis) en prenant
des produits fibrés sur $\bar{\Gamma}_{\rho,\sigma}$ de certains
d'entre eux. Le terme initial étant $S_0 M_{\rho,\sigma}$,
qui est le produit fibré des quatre facteurs
précédents. En prenant les noyaux des hom.
de ces \sout{quatre} neuf groupes dans $\bar{\Gamma}_{\rho,\sigma}$, on trouve
un diagramme du même type combinatoire, relié

562

\newpage

%% ===== Page 564 =====

neuf groupes, dont $H\} = \operatorname{Ker}(\Gamma_{p,\sigma} \to \bar{\Gamma}_{p,\sigma})$, formés de
$L_0 = S Z_{p,\sigma}$, $S N_p^*, S N_\sigma^*$ et $\mathcal{G}^+ = S G_{p,\sigma}$ et de certains
produits de ces groupes (dont le produit sur l'un
d'entre eux, savoir $H\} = \operatorname{Ker}(\Gamma_{p,\sigma} \xrightarrow{id} \bar{\Gamma}_{p,\sigma})$),
avec comme flèches [unclear: induites] des projections
canoniques. Nous avons inséré dans (62) seule-
ment cinq parmi ces groupes produits, avec
les injections canoniques dans les produits
fibrés dont ils proviennent, [unclear: bon d'omettre]
d'y faire figurer également
$S N_p^*, S N_\sigma^*, S G_{p,\sigma} = \mathcal{G}^+$, et les inclusions
canoniques $S N_p^* \hookrightarrow N_p^*, S N_\sigma^* \hookrightarrow N_\sigma^*, \mathcal{G}^+ = S G_{p,\sigma} \hookrightarrow G_{p,\sigma}$,
et les [unclear: par diagramme] (ou couples) des projections des
produits partiels de ces groupes et de $L_0$,
ainsi qu'une [unclear: trentaine] de flèches
y compris les projections sur ces trois groupes, et
les hom. de ces groupes et de $L_0$ dans le
groupe unité.

\newpage

%% ===== Page 565 =====

6$^{\circ}$) Liens entre $\Sigma_{\rho,\sigma}$, $S\mathcal{G}_{\rho,\sigma}$.

On pose
$$
\Sigma_{\rho,\sigma} = \mathfrak{G}_*^+ \times \mathfrak{G}_*^+ \times \hat{\mathbf{Z}}^* \times \{\pm 1\} \times \{\pm 1\} \leqno (68)
$$
$$
\parallel
$$
$$
\left\{(\xi,\eta, \mu, \varepsilon', \varepsilon) \left| \begin{array}{l} \xi \rho \text{ conjugué à } \rho^{\varepsilon'} \\ \eta \sigma \text{ conjugué à } \sigma^\varepsilon \\ \xi = \eta s_1^{\mu-1} \end{array} \right. \right\}
$$

NB Si $\rho$ est conjugué à $\rho^{-1}$ dans $\mathfrak{G}^+$, $\varepsilon'$ est connu quand on connait $\xi,\eta$ -- donc si $\sigma$ est conjugué à $\sigma^{-1}$ (donc $\sigma^2 \neq 1$), $\varepsilon$ est déterminé par $\xi,\eta$. De même $\hat{\mathbf{Z}} \to L_0$ ($u \mapsto s_0^{-u}$) est injectif, alors $\mu$ est déterminé par $\xi,\eta$, et la relation $\xi = \eta s_1^{\mu-1}$ s'écrit aussi (sans faire intervenir $\mu$)
$$
\eta^{-1} \xi \in L_1. \leqno (69)
$$

Soit d'autre part
$$
S\mathcal{G}_{\rho,\sigma} \subset \mathfrak{G}_* \times \mathfrak{G}_* \times \hat{\mathbf{Z}}^* \leqno (70)
$$
$$
\parallel
$$
$$
\left\{ (\alpha,\beta,\mu) \mid [\alpha,\rho] = [\beta,\sigma] s_1^{\mu-1} \right\}
$$

On a une application\marginpar{[unclear: NB $\xi = \alpha \rho \alpha^{-1} \rho^{-1} \in \mathfrak{G}_*^+$, car on a $\varepsilon_{\mathfrak{G}}(\xi) = \dots = 1$, et de même $\eta = \beta \sigma \beta^{-1} \sigma^{-1} \in \mathfrak{G}_*^+$]}
$$
S\mathcal{G}_{\rho,\sigma} \longrightarrow \Sigma_{\rho,\sigma} \leqno (71)
$$
$$
(\alpha,\beta,\mu) \longmapsto (\xi,\eta,\mu,\varepsilon',\varepsilon) \quad \text{avec } \begin{cases} \xi = [\alpha,\rho] \\ \eta = [\beta,\sigma] \\ \mu \text{ le même} \\ \varepsilon' = \varepsilon_{\mathfrak{G}}(\alpha) \\ \varepsilon = \varepsilon_{\mathfrak{G}}(\beta) \end{cases}
$$
qui est \underline{évidemment} surjective.

On définit d'autre part

\newpage

%% ===== Page 566 =====

$$
\Sigma_{\rho,\sigma}^{\text{eff}} \subset \Sigma_{\rho,\sigma} \leqno (72)
$$
$$
= \left\{ (\xi,\eta,\mu,\varepsilon',\varepsilon) \in \Sigma_{\rho,\sigma} \mid \exists u \in \operatorname{Aut}(\widehat{\mathfrak{S}}^+) \text{ avec } \begin{cases} u(\rho) = \xi\rho \\ u(\sigma) = \eta\sigma \end{cases} \text{ et induisant le car. } \mu \text{ sur } \hat{\mathbf{Z}} \text{ et les sign. } \varepsilon,\varepsilon' \right\}
$$

i.e. $\xi = u(\rho)\rho^{-1}$, $\eta = u(\sigma)\sigma^{-1}$ (cet $u$ sera alors unique, cf. [unclear: la suite]).

Il y a correspondance 1-1 entre les éléments de $\Sigma_{\rho,\sigma}^{\text{eff}}$ et les
systèmes $(u, \mu, \varepsilon', \varepsilon) \in \operatorname{Aut}(\widehat{\mathfrak{S}}^+) \times \hat{\mathbf{Z}}^* \times \{\pm 1\} \times \{\pm 1\}$,
tels que l'on ait
$$
\begin{cases}
u(\rho) \text{ est } \widehat{\mathfrak{S}}^+\text{-conjugué à } \rho^{\varepsilon'} \\
u(\sigma) \text{ est } \widehat{\mathfrak{S}}^+\text{-conjugué à } \sigma^{\varepsilon} \\
u(l_1) = l_1^\mu \quad (\text{qui exprime } \xi = \eta l_1^{\mu-1})
\end{cases}
$$
or cet ensemble n'est autre que le
sous-groupe $B_{\rho,\sigma}$ de $\operatorname{Aut}(\widehat{\mathfrak{S}}^+) \times \hat{\mathbf{Z}}^* \times \{\pm 1\} \times \{\pm 1\}$ ! (cf. (16) p. 577 !)

On trouve ainsi une bijection
$$
B_{\rho,\sigma} \isommap \Sigma_{\rho,\sigma}^{\text{eff}} \leqno (73)
$$
$$
u \longmapsto (\xi,\eta,\mu,\varepsilon',\varepsilon) \quad \text{avec} \begin{cases} \xi = u(\rho)\rho^{-1} \\ \eta = u(\sigma)\sigma^{-1} \\ \mu = \mu(u) \\ \varepsilon' = \varepsilon'_\rho(u) \\ \varepsilon = \varepsilon_\sigma(u) \end{cases}
$$

On désigne par $SG_{\rho,\sigma}^{\text{eff}}$ l'image inverse de $\Sigma_{\rho,\sigma}^{\text{eff}}$ dans $SG_{\rho,\sigma}$ par (71), d'où on a
$$
SG_{\rho,\sigma}^{\text{eff}} \subset SG_{\rho,\sigma} \leqno (74)
$$
$$
= \left\{ (\alpha,\beta,\mu) \in \widehat{\mathfrak{S}}^+ \times \widehat{\mathfrak{S}}^+ \times \hat{\mathbf{Z}}^* \mid \begin{array}{l} [\alpha,\rho] = [\beta,\sigma] l_1^{\mu-1} \\ \exists u \in \operatorname{Aut}(\widehat{\mathfrak{S}}^+) \text{ tel qu'on ait } \begin{cases} u(\rho) = \operatorname{int}(\alpha)(\rho) & (= \xi\rho, \text{ avec } \xi = [\alpha,\rho]) \\ u(\sigma) = \operatorname{int}(\beta)(\sigma) & (= \eta\sigma, \text{ avec } \eta = [\beta,\sigma]) \end{cases} \end{array} \right\}
$$
de sorte qu'on a un homomorphisme surjectif induit par (71)
$$
SG_{\rho,\sigma}^{\text{eff}} \longrightarrow \Sigma_{\rho,\sigma}^{\text{eff}} \leqno (75)
$$

\newpage

%% ===== Page 567 =====

On a une application canonique\footnote{pour être plus correct, il faudrait écrire $\alpha = u \cdot i_{\rho}(a)^{-1} {\tau'}^{\nu'(x)}$, $\beta = u \cdot i_{\sigma}(b)^{-1} \tau^{\nu(x)}$}
$$
S_0\Gamma_{\rho,\sigma} = \underline{Z}_{\rho,\sigma} \times_{\underline{Z}_{\rho,\sigma}} N_{\rho}^\ast \times_{\Gamma_{\rho,\sigma}} N_{\sigma}^\ast \longrightarrow SG_{\rho,\sigma}^{\text{eff}} \leqno (76)
$$
$$
x = (u, a, b) \longmapsto (\alpha, \beta, \mu) \text{ avec }
\begin{cases}
\alpha = u a^{-1} {\tau'}^{\nu'(x)} \\
\beta = u b^{-1} \tau^{\nu(x)} \\
\mu = \chi(x)
\end{cases} \in SG_{\rho,\sigma}^{\text{eff}}
$$
où on a défini $\nu'(x), \nu(x) \in \mathbf{Z}/2\mathbf{Z}$ par
$$
(-1)^{\nu'(x)} = \varepsilon'(x) \quad , \quad (-1)^{\nu(x)} = \varepsilon(x)
$$

Par définition de $S_0\Gamma_{\rho,\sigma}$ comme produit fibré, on aura en effet $u a^{-1}, u b^{-1} \in \mathfrak{S}^+$, donc $\alpha, \beta \in \mathfrak{S}$,
on aura d'ailleurs\marginpar{NB On a $a \in N_{\rho}^\ast \Rightarrow \varepsilon'(a) = \varepsilon(u)$ et $b \in N_{\sigma}^\ast \Rightarrow \varepsilon(b) = \varepsilon(u)$. Ainsi pour $[\alpha,\rho]$}
$$
\begin{aligned}
\alpha(\rho) &= u \bigl( \underbrace{a^{-1} ( \underbrace{{\tau'}^{\nu'(x)}(\rho)}_{\rho^{\varepsilon'(x)}} )}_{(\rho^{\varepsilon'(x)})^{\varepsilon'(a)} = \rho} \bigr) = u(\rho) & \text{i.e.} \quad & [\alpha,\rho] = [u,\rho] \\
\beta(\sigma) &= u(\sigma) & \text{i.e.} \quad & [\beta,\sigma] = [u,\sigma]
\end{aligned} \leqno (77)
$$
de sorte que la condition $[\alpha,\rho] = [\beta,\sigma] \varepsilon_1^{\mu-1}$ se réduit à $[u,\rho] = [u,\sigma]$, ce qui exprime en fait la condition que $u(\varepsilon_0) = \varepsilon_0^\mu$ qui intervient dans la définition de $Z_{\rho,\sigma}$.

Je dis que l'application envisagée (76) est \emph{injective}. Car il [unclear: résulte de ce qui précède] que l'on a
$$
\nu'(x) = \nu_{\rho}(\alpha) \quad , \quad \nu(x) = \nu_{\sigma}(\beta) \quad , \quad \chi(x) = \mu \leqno (78)
$$
d'où on tire $a = {\tau'}^{\nu_{\rho}(\alpha)} \alpha^{-1} u$ , $b = \tau^{\nu_{\sigma}(\beta)} \beta^{-1} u$. Donc si $u$ est connu, on connaît également $a$ et $b$ donc $x$. D'autre part pour déterminer $u \in Z_{\rho,\sigma} \subset \operatorname{Aut}(\mathfrak{S}^+) \times \widehat{\mathbf{Z}}^\ast \times \{\pm 1\} \times \{\pm 1\}$ il suffit en vertu de (78) de le déterminer en tant qu'élément de $\operatorname{Aut}(\mathfrak{S}^+)$, ce qui est

\newpage

%% ===== Page 568 =====

immédiat sur (77). On a donc
$$
\underline{S}\Gamma_{\rho,\sigma} \hookrightarrow \underline{S}\mathcal{G}_{\rho,\sigma}^\text{eff} \leqno{(79)}
$$

il reste à voir que cette application est surjective -
on a fait ce qu'il fallait pour. Notons
d'abord que le composé
$$
\underset{(x=(u,a,b))}{\underline{S}_0 \Gamma_{\rho,\sigma}} \longrightarrow \underline{S}\mathcal{G}_{\rho,\sigma}^\text{eff} \xrightarrow{\text{eff int}} \Sigma_{\rho,\sigma}^\text{eff} \ (\xrightarrow{\sim} \underline{Z}_{\rho,\sigma})
$$
n'est autre, par (77) et (71) que
$$
x = (u,a,b) \longmapsto (\xi,\eta,\mu,\varepsilon',\varepsilon) \quad
\begin{cases}
\xi = [u,\rho] & \text{où } u(\rho) = \xi\rho \\
\eta = [u,\sigma] & \text{où } u(\sigma) = \eta\sigma \\
\mu = \chi_{\underline{S}\Gamma}(x) \\
\varepsilon' = \varepsilon'(x) \\
\varepsilon = \varepsilon(x)
\end{cases} \leqno{(80)}
$$
donc son composé avec $\Sigma_{\rho,\sigma}^\text{eff} \xrightarrow{\sim} Z_{\rho,\sigma}$
n'est autre que la projection canonique
$$
\underline{S}_0 \Gamma_{\rho,\sigma} \longrightarrow Z_{\rho,\sigma}
$$
$$
(u,a,b) \longmapsto u
$$
qui est surjective, de noyau $S\mathcal{N}_\rho^* \times S\mathcal{N}_\sigma^* = \underline{S}Z_\rho \times \underline{S}Z_\sigma$.

Considérons d'autre part
$$
\begin{array}{c}
\text{opération de } S\mathcal{N}_\rho \times S\mathcal{N}_\sigma \\
\text{sur } \underline{S}\mathcal{G}_{\rho,\sigma} \subset \mathfrak{S}^\times \times \mathfrak{S}^\times \times \hat{\mathbf{Z}}^\times \\
\shortparallel \\
\{(\alpha,\beta,\mu) \mid [\alpha,\rho] = [\beta,\sigma]\varepsilon_1^{\mu-1}\}
\end{array}
\quad
\begin{cases}
(a_0,b_0) \in S\mathcal{N}_\rho^* \times S\mathcal{N}_\sigma^* \\
\text{opère par} \\
(\alpha,\beta,\mu) \longmapsto (\alpha a_0, \beta b_0, \mu)
\end{cases} \leqno{(81)}
$$
Il est clair qu'il opère fidèlement, et
que le quotient de $\underline{S}\mathcal{G}_{\rho,\sigma}$ par cette action
s'identifie à $\Sigma_{\rho,\sigma}$ , donc on a un résultat
identique ([unclear: pour $\underline{S}\mathcal{G}_{\rho,\sigma}^\text{eff}$ et $\Sigma_{\rho,\sigma}^\text{eff}$, quotients par l'action de $D_2$]), pour l'action (81) de $\underline{S}Z_\rho^* \times \underline{S}Z_\sigma^*$ correspondant
dans $\underline{S}_0\Gamma_{\rho,\sigma}$ à la translation
$$
x=(u,a,b) \longmapsto (u, \varepsilon_\rho^{\nu'(x)}(a_0^{-1})a, \varepsilon_\sigma^{\nu(x)}(b_0^{-1})b) \leqno{(82)}
$$

\newpage

%% ===== Page 569 =====

$$ \varepsilon = (-1)^{\nu}, \quad \varepsilon' = (-1)^{\nu'} \leqno{(82)} $$
fixés, l'opération précédente de $(a_0, b_0)$ correspond,
par l'application (76), à l'opération
$$ (u, a, b) \longmapsto (u, \tau^\nu(a_0^{-1})a, \tau^{\nu'}(b_0^{-1})b) \ , \leqno{(83)} $$
i.e.\ à la translation à gauche par
$$ (\tau^\nu, \tau^{\nu'})(a_0^{-1}, b_0^{-1}) \ . $$
D'où la surjectivité, et l'isomorphisme cherché
$$ S_0\Gamma_{\rho, \sigma} \isommap S\mathcal{G}^{\text{eff}}_{\rho, \sigma} \quad \text{(cf (75))} \leqno{(84)} $$

La translation à gauche par $l \in L_{S\Gamma_{\rho, \sigma}}$ correspond,
par la bijection précédente, à l'opération
$$ (\alpha, \beta, \mu) \longmapsto (l_\alpha \alpha, l_\beta \beta, \mu) \leqno{(85)} $$
dont il était clair a priori qu'elle stabilise
$S\mathcal{G}^{\text{eff}}_{\rho, \sigma}$. Passant aux quotients, on trouve
donc
$$ \Gamma_{\rho, \sigma} \isommap L \backslash S\mathcal{G}^{\sharp}_{\rho, \sigma} \ . \leqno{(86)} $$

Enfin, considérons
$$
\begin{cases}
S_0M_{\rho, \sigma} \longrightarrow S\mathcal{G}^{\text{eff}}_{\rho, \sigma} \times \widehat{\mathfrak{S}}^+ = \left\{ (\alpha, \beta, \mu, f) \in \mathfrak{G} \times \mathfrak{G} \times \hat{\mathbf{Z}} \times \widehat{\mathfrak{S}}^+ \mid [\alpha, \rho] = [\beta, \sigma] \varepsilon_0^{H-1} \right\} \\
x = (u, a, b, U) \longmapsto
\begin{cases}
\alpha = u a^{-1} \tau^\nu(a) \\
\beta = u b^{-1} \tau^{\nu'}(b) \\
\mu = \chi(x) \ ( = \chi_\rho(u) = \chi_\rho(a) = \chi_\sigma(b) = \chi_{S_M}(U) ) \\
f = u U^{-1}
\end{cases}
\end{cases} \leqno{(87)}
$$
Cette application est bijective (trivial par (84)).

\newpage

%% ===== Page 570 =====

7) Fonctorialités

Considérons
(88) $\hat{G} \twoheadrightarrow \hat{G}'$ hom. surjectif de groupes profinis

soient $\rho', \sigma', \tau', \dots$ les images de $\rho, \sigma, \tau$,
$\tau' = \tau_{\rho'\sigma'}$ dans $\hat{G}'$.

Prop. Conditions équivalentes sur $\varphi$ :

a) L'application induite par $\varphi$
(89) $\Sigma_{\rho,\sigma} \to \Sigma_{\rho',\sigma'}$
envoie $\Sigma_{\rho,\sigma}^{eff}$ dans $\Sigma_{\rho',\sigma'}^{eff}$

b) Pour tout $u \in Aut \hat{C}_{>}^+$ qui normalise $L_0$,
et tel que $u(L_\rho) \underset{\hat{C}^+}{\sim} L_\rho$, $u(L_\sigma) \underset{\hat{C}^+}{\sim} L_\sigma$,
il existe $u' \in Aut (\hat{C}'^+_{>})$ qui rend commutatif
(90)
\begin{tikzcd}
\hat{C}^+ \ar[r, "\varphi"] \ar[d, "u"'] & \hat{C}'^+ \ar[d, "u'"] \\
\hat{C}^+ \ar[r, "\varphi"'] & \hat{C}'^+
\end{tikzcd}
en d'autres termes, $u$ respecte $Ker \varphi$.

b') idem, en supposant seulement, à la place que
$u$ normalise $L_0$, que $u(L_0) \underset{\hat{C}^+}{\sim} L_0$ (plus les autres conditions)

c) $\exists$ hom. de groupes
(91) $\varphi_Z : Z_{\rho,\sigma} \to Z_{\rho',\sigma'}$
compatible avec $j, j_\rho, j_\sigma$ et tel que $\forall u \in Z_{\rho,\sigma}$, le
$u'$ qui lui est associé soit comme dans (90).

[MARGIN: alors il induit $\Gamma_{\rho,\sigma} \to \Gamma_{\rho',\sigma'}$] (NB $\varphi_Z$ uniquement déterminé par $\varphi$, car est compatible avec $j : L_0 \to Z$)

c') $\exists$ hom de groupes
(92) $\varphi_{\hat{G}} : G_{\rho,\sigma} \to G_{\rho',\sigma'}$
compatible avec $j_a, j_b, j'_c$ et satisfaisant à la
condition de compatibilité (90). (NB $\varphi_{\hat{G}}$ a [unclear])

\newpage

%% ===== Page 571 =====

est alors unique, et compatible aux $i_{Z, G}^{\varepsilon, \alpha}$

On dira alors que $\varphi$ est "prolongeable".

Corollaire On peut alors prolonger $\varphi$ en
$$
\varphi_M : S_0 M_{\underline{\rho}, \sigma} \longrightarrow S_0 M'_{\underline{\rho}, \sigma} \leqno{(93)}
$$
compatible aux $\Theta_Z, \Theta_P, \Theta_N, \Theta_G$
compte tenu de $\varphi_Z, \varphi_P, \varphi_N, \varphi_G$, et ceci
de façon "compatible avec tout", notamment
avec les diagrammes (45) (62).

Exemple Supposons $\Sigma_{\underline{\rho}, \sigma} = \Sigma_{\underline{\rho}, \sigma}^{G'}$, alors
les conditions de la prop. sont automatiquement
satisfaites quel que soit l'homo. surjectif
$G \to G'$.

Remarque \marginpar{Il faut supposer pour être tranquille que $\varphi_Z(Z^\natural) \subset Z'^\natural$} Il se peut qu'on se donne un
sous-groupe $Z_{\underline{\rho}, \sigma}^\natural \subset Z_{\underline{\rho}, \sigma}$, et qu'on
ne se donne la condition a) de la prop.
seulement pour des $u \in Z_{\underline{\rho}, \sigma}^\natural$. Mais à l'aide
de $Z_{\underline{\rho}, \sigma}^\natural$, on construit
$$
G_{\underline{\rho}, \sigma}^\natural = Z_{\underline{\rho}, \sigma}^\natural \cdot G^+, Z_P^\natural = Z_P \cap G_{\underline{\rho}, \sigma}^\natural, Z_N^\natural = Z_N \cap G_{\underline{\rho}, \sigma}^\natural \leqno{(94)}
$$
d'où $N_G^\natural$ et $N_G^{*\natural}$, $N_P^\natural$ et $N_P^{*\natural}$, $N_{P, \sigma}^\natural$, puis
$S_0 M_{\underline{\rho}, \sigma}^\natural$ et $M_G^\natural$. On a des variantes évidentes de la

\newpage

%% ===== Page 572 =====

proposition, [unclear: fournissant] des conditions de prolongement de $\varphi$ en $S_0 \mathcal{M}_{\rho,\sigma}^\natural \to S_0 \mathcal{M}_{\rho,\sigma}^{\widetilde{\sigma}'}$.

\underline{Cas universel} $\widetilde{G} = \operatorname{Gl}(2, \mathbf{Z})^\wedge$, $\rho = \begin{pmatrix} 0 & 1 \\ -1 & 1 \end{pmatrix}$, $\sigma = \begin{pmatrix} 0 & -1 \\ 1 & 0 \end{pmatrix}$
$$
\tau' = \begin{pmatrix} 0 & 1 \\ 1 & 0 \end{pmatrix}, \quad ( \tau = \begin{pmatrix} -1 & 0 \\ 0 & +1 \end{pmatrix}, \quad \varepsilon_0 = \begin{pmatrix} 1 & -1 \\ 0 & 1 \end{pmatrix},
$$
$$
\varepsilon_1 = \begin{pmatrix} 1 & 0 \\ 1 & 1 \end{pmatrix}.
$$

\newpage

%% ===== Page 573 =====

8°) Le cas universel $\hat{G} = \operatorname{Gl}(2,\mathbb{Z})^\wedge$.  (590)

Considérons la suite exacte

(1) 
\begin{tikzcd}
1 \ar[r] & \{\pm 1\} \ar[r] & \mathcal{N}_{1,1}^\sim \ar[r] \ar[d, equal] & \mathcal{M}_{0,3}^\sim \ar[r] & 1 \\
& & \mathcal{M} & &
\end{tikzcd}

[MARGIN:
$\hat{G} = \hat{G}_{1,1} \simeq \operatorname{Gl}(2,\mathbb{Z})^\wedge$
$\hat{G}^+ = \hat{\Gamma}_{1,1} \simeq \operatorname{Sl}(2,\mathbb{Z})^\wedge$
]

$\{\pm 1\} = Centr(\mathcal{M}) = Centr(\hat{G}^+)$

d'où par (1) une application

(3) $\mathcal{M}_{0,3}^\sim \to \operatorname{Aut}(\mathcal{M}) \xrightarrow{\text{restriction à } \hat{G}^+} \operatorname{Aut}(\hat{G}^+)$

D'ailleurs, dans le cas [unclear] de $\hat{G}^+$ [unclear] 
[unclear] un centre, [unclear: par suite] d'où

(4) $\operatorname{Aut}(\hat{G}^+) \to \operatorname{Aut}(\hat{G}^+ / \pm 1 \simeq \hat{G}_0^+)$

et le composé de (3) avec (4) est l'homomorphisme canonique :

(5)
\begin{tikzcd}
\mathcal{M}_{0,3}^\sim \ar[r] \ar[dr] & \operatorname{Aut}(\hat{G}^+) \ar[d] \\
& \operatorname{Aut}(\hat{G}_0^+ = \hat{G}^+ / \pm 1)
\end{tikzcd}

Cela montre que (3) est injectif, puisque $\mathcal{M}_{0,3}^\sim \to \operatorname{Aut}(\hat{G}_0^+)$ l'est. L'application

Remarque L'application $\operatorname{Aut}(\hat{G}^+) \to \operatorname{Aut}(\hat{G}_0^+)$ n'est pas injective par contre, et les éléments du noyau correspondent aux

hom. $\hat{G}^+ \to Centr(\hat{G}^+) = \{\pm 1\}$

\newpage

%% ===== Page 574 =====

ou encore aux hom. de

(6) $$ \widehat{C}^+_{ab} \longrightarrow \{\pm 1\} \quad \text{ou encore} \quad (\widehat{C}^+_{ab})_2 \longrightarrow \{\pm 1\} $$

Or on voit que

$$ \widehat{C}^+_{ab} = \{ \rho, \sigma \mid \rho^6 = \sigma^4 = [\rho, \sigma] = 1, \rho^3 = \sigma^2 \} = \mathbb{Z}/12\mathbb{Z} $$
[MARGIN: $\rho \mapsto 2 \bmod 12$, $\sigma \mapsto 3 \bmod 12$]

donc

(7) $$ \left\{ \begin{aligned} (\widehat{C}^+_{ab})_2 &\xrightarrow{\sim} \mathbb{Z}/2\mathbb{Z} \\ \rho &\longmapsto 0 \\ \sigma &\longmapsto 1 \end{aligned} \right. $$

Cet homomorphisme correspond à l'hom.
$\widehat{C}^+ \longrightarrow \{\pm 1\}$ qui s'obtient aussi comme
composé :

(8)
\begin{tikzcd}
\widehat{C}^+ \ar[r] & \mathfrak{S}_3 = \widehat{C}^+/\widehat{\pi} \ar[r, "sgn"] & \{\pm 1\}
\end{tikzcd}
[MARGIN: sous $\mathfrak{S}_3$ : groupe symétrique à trois lettres. sous $\widehat{\pi}$ : im. inv. de $\widehat{\Pi}_{0,3}$ dans $\widehat{C}^+$]

Ainsi le seul élément $\Theta$ différent de $1$
du noyau de (4) est donné par

(9) $$ \Theta : x \longmapsto x \cdot \text{sgn}(x) $$

On notera, en posant

(10) $$ \widehat{L}_0 = \varepsilon_0^{\hat{\mathbb{Z}}}, \quad \widehat{L}''_0 = \widehat{L}_0 \cap \widehat{\pi} = \widehat{L}_0^2 $$

que le s-groupe $\Theta(\widehat{L}_0)$ de $\widehat{C}^+$
n'est pas conjugué à $\widehat{L}_0$. En effet

(11) $$ \Theta(\varepsilon_0) = -\varepsilon_0 \quad (\text{car } \varepsilon_0 = \sigma \rho, \text{ sgn}(\varepsilon_0) = -1) $$

et on vérifie que $\varepsilon_0$ n'est pas conjugué à
$(-\varepsilon_0)^k = \tau_0 \varepsilon_0^k$ pour quelque $k \in \hat{\mathbb{Z}}^*$ (car on voit

\newpage

%% ===== Page 575 =====

il s'ensuit, en [unclear: projetant sur] $\hat{\mathfrak{S}}/\hat{\Pi}_0 \simeq$
$\hat{\mathfrak{S}}_{0,3} / \hat{\Pi}_{0,3}$ et en notant que $l_0 = \varepsilon_0^2$ serait
conjugué à $(-\varepsilon_0)^2 = l_0^{\mu}$, ce qui est [unclear: absurde]
($\mu = 1$, car $\varepsilon_0$ et $-\varepsilon_0$ ne sont pas conjugués
même dans $\operatorname{Sl}(2, \hat{\mathbf{Z}})$ ! ...). Or on va
s'intéresser justement aux automorphismes $u$ de $\hat{\mathfrak{S}}^{+\wedge}$
qui transforment $L_0^{\hat{\pi}}$ en un sous-groupe
conjugué. On notera que pour un tel $u$
$$
\Theta(l_0) = l_0 \quad \text{donc} \quad \Theta(L_0^{\hat{\pi}}) = L_0^{\hat{\pi}} \leqno(12)
$$

Proposition. L'homomorphisme (3)
$$
\mathcal{M}_{0,3}^{\sim} \longrightarrow \operatorname{Aut}(\hat{\mathfrak{S}}^{+\wedge})
$$
est injectif, et son image est formée des
automorphismes $u$ de $\hat{\mathfrak{S}}^{+\wedge}$ qui satisfont les deux
conditions suivantes :
\begin{enumerate}
\item[a)] $u$ stabilise $\hat{\pi}^* = \hat{\mathfrak{S}}^{+\wedge}[2] = \operatorname{Ker}(\hat{\mathfrak{S}}^{+\wedge} \to \operatorname{Sl}(2, \mathbf{Z}/2\mathbf{Z}))$
\item[b)] $u$ transforme $L_0^{\hat{\pi}}$ en un sous-groupe de $\hat{\mathfrak{S}}^{+\wedge}$
$\hat{\mathfrak{S}}^{+\wedge}$-conjugué de $L_0$, i.e.\ $\exists \, \mu \in \hat{\mathbf{Z}}^*$ et $f_0 \in \hat{\mathfrak{S}}^{+\wedge}$
tels qu'on ait
$$
u(\varepsilon_0) = \operatorname{int}(f_0^{-1})(\varepsilon_0^\mu) \leqno (13)
$$
\end{enumerate}

Dém. L'injectivité a déjà été vue. [unclear: Pour voir]
que les conditions a) b) sont nécessaires pour
qu'un $u$ [unclear: provienne] d'un $u_0 \in \mathcal{M}_{0,3}^{\sim}$, pour a)
c'est trivial, puisqu'un $u_0$ stabilise $\hat{\Pi}_{0,3}$, donc son...

\newpage

%% ===== Page 576 =====

image inverse dans $\mathfrak{S}^{+\wedge}$. Pour la condition
b), on est ramené au cas où $u_0 \in \tilde{\mathcal{M}}_{0,3}[0]$
(quitte à conjuguer par un élément de $\widehat{\mathfrak{S}}_{0,3} \simeq \mathfrak{S}^{+\wedge}/\pm 1$)
dans ce cas on a par construction de $\tilde{N}_{1,1} \subset \tilde{\mathcal{M}}_{0,3}$
(cf. III, p. 665...)
$$
u_0(\varepsilon_0) = \varepsilon_0^\mu \leqno{(14)}
$$
où $\mu \in \hat{\mathbf{Z}}^*$ est le multiplicateur, OK.

Inversement, soit $u$ un autom. de $\mathfrak{S}^{+\wedge}$ qui
satisfait a), b), considérons l'autom. $\tilde{u}$ de
$\mathfrak{S}^{+\wedge}/\pm 1 \simeq \widehat{\mathfrak{S}}_{0,3}$ qu'il induit. On voit donc
que $\tilde{u}(\hat{\pi}_{0,3}) = \hat{\pi}_{0,3}$, et que $\tilde{u}(\varepsilon_0)$ est $\widehat{\mathfrak{S}}_{0,3}$-con-
jugué à un $\varepsilon_0^\mu$ ($\mu \in \hat{\mathbf{Z}}^*$). Donc par
définition, on a $\tilde{u} \in \tilde{\mathcal{M}}_{0,3}$ (cf. page 390 ff.).
Soit $u'$ l'autom. de $\mathfrak{S}^{+\wedge}$ induit par $\tilde{u}$
(via B)) - alors $u'$ et $u$ induisent le même
autom. de $\widehat{\mathfrak{S}}_{0,3} = \mathfrak{S}^{+\wedge}/\{\pm 1\}$, donc on a
$u' = u \circ \Theta$, $\Rightarrow \Theta = u^{-1}u'$ est un autom. qui
\begin{enumerate}
\item[1)] induit l'identité sur $\mathfrak{S}^{+\wedge}/\{\pm 1\}$, et
\item[2)] transforme $L_0^\pi$ en un sous-groupe conjugué
\end{enumerate}
--- on a vu que cela implique $\Theta = 1$, cqfd.

Considérons maintenant la suite
exacte
$$
1 \to \pm 1 \to \underset{= \mathfrak{S}^{+\wedge}[2]}{\pi^\wedge} \to \hat{\pi}_{0,3} \to 0 \leqno{(15)}
$$

\newpage

%% ===== Page 577 =====

et considérons le scindage donné par
$$
\begin{cases}
\Pi_{0,3} \longrightarrow \Pi = \mathfrak{S}^+[2] = \operatorname{Ker}(\operatorname{Sl}(2,\mathbf{Z}) \to \operatorname{Sl}(2,\mathbf{Z}/2)) \\
l_i \longmapsto \lambda_i \defeq -l_i
\end{cases} \leqno{(16)}
$$
Il y a un tel homomorphisme, en vertu de
$$
\begin{cases}
(-l_0)(-l_1)(-l_\infty) = 1 \quad \text{dans } \operatorname{Sl}(2,\mathbf{Z}) \quad \text{i.e.} \\
l_\infty l_1 l_0 = -1
\end{cases} \leqno{(17)}
$$
et il est évident que c'est un scindage. On en déduit un scindage pour $\widehat{\Pi}$ sur $\widehat{\Pi}_{0,3}$.

On désigne par $\Pi_0$ l'image de $\Pi_{0,3}$, i.e.
$$
\begin{cases}
\Pi_0 = \text{sous-groupe de } \Pi = \mathfrak{S}^+[2] \text{ engendré par} \\
\lambda_0, \lambda_1, \lambda_\infty \quad (\text{où } \lambda_i = -l_i)
\end{cases} \leqno{(18)}
$$
et on a de même un sous-groupe $\widehat{\Pi}_0$ de $\widehat{\Pi}$.

\emph{Proposition.}
\begin{enumerate}
\item[a)] Le sous-groupe $\Pi_0$ de $\Pi$ est d'indice 2, et il est distingué dans $\mathfrak{S}^+$.
\item[b)] Le sous-groupe $\widehat{\Pi}_0$ de $\widehat{\Pi}$ est d'indice 2, et il est distingué dans $\widehat{\mathcal{M}}$ --- ou, ce qui revient au même, est stable sous l'action de $\widehat{\mathcal{M}} / \pm 1 \isom \widehat{\mathcal{M}}_{0,3}$.
\end{enumerate}

Démonstration. Comme
$$
\left\{
\begin{aligned}
&\rho(\lambda_i) = \lambda_{i+1} \\
&\sigma(\lambda_0) = \lambda_1^{-1}, \quad \sigma(\lambda_1) = \lambda_0^{-1}, \quad \sigma(\lambda_\infty) = \sigma(\lambda_1 \lambda_0)^{-1} = \lambda_1^{-1} \lambda_0^{-1} \\
&\tau(\lambda_0) = \lambda_0^{-1}, \quad \tau(\lambda_1) = \lambda_1^{-1}, \quad \tau(\lambda_\infty) = \tau(\lambda_1 \lambda_0)^{-1} = \lambda_0 \lambda_1
\end{aligned}
\right. \leqno{(19)}
$$
on voit que $\Pi_0$ est normalisé par $\rho, \sigma, \tau$, donc par $\mathfrak{S}^+$. Qu'il soit d'indice 2 est évident. Donc $\widehat{\Pi}_0$ est d'indice 2 dans $\widehat{\Pi}$ et normalisé par $\widehat{\mathfrak{S}}^+$. Pour voir qu'il est

\newpage

%% ===== Page 578 =====

normalisé par $\mathcal{M}$ tout entier, comme
$$
\mathcal{M} = \mathcal{M}(0) \cdot \mathfrak{S}^{+\wedge} ,
$$
il nous revient à vérifier qu'il est normalisé par $\mathcal{M}(0) \simeq \mathcal{M}_{0,3}(0)$. Or on a, pour $u \in \mathcal{M}_{0,3}(0)$
$$
\begin{cases}
u(\rho) = \operatorname{int}(\alpha) \rho \\
u(\sigma) = \varepsilon_2 \operatorname{int}(\beta) \sigma \\
u(\varepsilon_0) = {\varepsilon_0}^{k}
\end{cases}
\quad
\begin{cases}
\varepsilon_2 = \varepsilon_1(u) \\
\alpha, \beta \in \hat{\Pi} \cdot \{1, \tau\} \\
\beta \in \hat{\pi}
\end{cases}
\leqno (20)
$$
d'où
$$
\begin{cases}
u(l_0) = {l_0}^k \\
u(-l_0) = (-l_0)^k \text{ i.e. } u(\lambda_0) = {\lambda_0}^k \quad (\text{car } (-1)^k = -1)
\end{cases}
\leqno (21)
$$
Pour montrer que $u(\lambda_0) \in \hat{\Pi}_0$, il reste à prouver que $u(\lambda_1) \in \hat{\Pi}_0$, or
$$
u(\lambda_1) = u(\rho(\lambda_0)) = (u \rho)(\lambda_0) = \operatorname{int}(\alpha)(\rho) (\underbrace{u(\lambda_0)}_{{\lambda_0}^k}) = \operatorname{int}(\alpha)(\rho({\lambda_0}^k))
$$
et comme on a vu que $\hat{\Pi}_0$ est normalisé par $\mathfrak{S}^{+\wedge}$ et en particulier par $\rho, \alpha$, on trouve bien $u(\lambda_1) \in \hat{\Pi}_0$.

Corollaire. On a
$$
\begin{cases}
\Pi \simeq \Pi_0 \times \mu \\
\hat{\Pi} \simeq \hat{\Pi}_0 \times \mu
\end{cases}
\leqno (22)
$$
où
$$
\mu = \{ \pm 1 \} = \operatorname{Centr}(\mathfrak{S}^{+\wedge})
$$
C'est trivial, puisqu'on a un scindage de l'extension centrale (15).

\newpage

%% ===== Page 579 =====

Soient
$$
\begin{cases}
\mu_\tau = \{1, \tau\} \subset \widehat{\mathfrak{S}} \\
\mu_\sigma = \{1, \sigma\} \subset \widehat{\mathfrak{S}}
\end{cases}
\qquad
\begin{aligned}
D_\tau &= \mu_\tau \times \mu \subset \widehat{\mathfrak{S}} \\
D_\sigma &= \mu_\sigma \times \mu \subset \widehat{\mathfrak{S}}
\end{aligned}
\leqno{(23)}
$$

de sorte que l'on a
\marginpar{idem en remplaçant $\widehat{\mathfrak{S}}, \pi$ par $\widehat{\mathfrak{S}}_1, \hat{\pi}$}
$$
\begin{aligned}
D_\tau \cap \pi &= \mu \qquad D_\tau \cap \pi_0 = \{1\} \\
D_\sigma \cap \pi &= \mu \qquad D_\sigma \cap \pi_0 = \{1\}
\end{aligned}
\leqno{(24)}
$$

On pose
$$
\begin{cases}
\pi_\tau = \mu_\tau . \pi = D_\tau . \pi_0 , & \pi_{0,\tau} = \mu_\tau . \pi_0 \\
\pi_\sigma = \mu_\sigma . \pi = D_\sigma . \pi_0 , & \pi_{0,\sigma} = \mu_\sigma . \pi_0
\end{cases}
\leqno{(25)}
$$

de sorte que l'on a
$$
\begin{cases}
\pi_\tau = \pi_{0,\tau} \times \mu \\
\pi_\sigma = \pi_{0,\sigma} \times \mu
\end{cases}
\leqno{(26)}
$$

et idem en mettant des $\wedge$.

Normalisation de $\alpha, \beta$ \quad Les relations (20) ne déterminent $\alpha, \beta$ qu'au "[unclear: signe]" près. On lève cette ambiguïté en prenant
$$
\alpha \in \pi_{0,\tau} \quad , \quad \beta \in \pi_{0,\sigma}
\leqno{(27)}
$$

Le groupe $\widehat{\mathfrak{S}}/\pi_0$. Il est décrit par
gén. $\rho, \sigma, \tau$ et relations
$$
\begin{cases}
\rho^6 = \sigma^4 = \tau^2 = 1, \quad \rho^3 = \sigma^2 \\
\tau(\sigma) = \sigma^{-1} \\
\tau(\rho) = \rho
\end{cases}
\leqno{(28)}
$$

$$
\tau(\rho) = (-l_1^{-1}) \rho
\leqno{(29)}
$$
donc la dernière relation (28) s'écrit $-l_1^{-1} = 1$ i.e. $l_1 = -1$
(et implique [unclear: $\lambda_0 = \bar{\lambda}_0 = 1$ en multipliant $\rho$] ).

\newpage

%% ===== Page 580 =====

En termes de $\rho, \sigma, \tau'$ on trouve [unclear: donc] :
$$
\begin{cases}
\rho^6 = \sigma^4 = \tau'^2 = 1 \quad , \quad \rho^3 = \sigma^2 \\
\tau'(\rho) = \rho^{-1}, \tau'(\sigma) = \sigma^{-1} \\
\sigma(\rho) = \rho^{-1}
\end{cases} \leqno{(30)}
$$

Le sous-groupe $\mathfrak{S}^+/\Pi_0$ a la présentation
en $\rho, \sigma$
$$
\begin{cases}
\rho^6 = \sigma^4 = 1 \quad , \quad \rho^3 = \sigma^2 \\
\sigma(\rho) = \rho^{-1}
\end{cases} \leqno{(31)}
$$

Quelle est l'opération de $\mathcal{M}_{0,\tilde{5}} \simeq \mathcal{M} / \pm 1$
sur $\mathfrak{S}^+/\Pi_0$ ?

$$
1 \to \underset{\substack{\parallel \\ \pm 1}}{\mu} \to \mathfrak{S}^+/\Pi_0 \to \underset{\substack{\parallel \\ \mathfrak{S}_3}}{\mathfrak{S}^+/\mu} \to 1 \leqno{(32)}
$$

et les actions en question conservent cette
structure d'extension. Les éléments de
$\mathcal{M}_{0,\tilde{5}} \simeq \mathcal{M}/ \pm 1$ opèrent en [unclear: induisant] l'
identité sur le quotient $\mathfrak{S}_3$, et bien sûr
aussi sur $\mu$ - or un automorphisme de
l'extension correspond à un homomorphisme
$$
\mathfrak{S}_3 \to \mu \quad \text{i.e.} \quad \mathfrak{S}_{3,ab} \to \mu \leqno{(33)}
$$
or $\mathfrak{S}_{3,ab} \xrightarrow{\sim} \{\pm 1\} = \mu$
donc le groupe des automorphismes de $\mathfrak{S}^+/\Pi_0$ indui-
sant l'identité sur $\mathfrak{S}_3$ et sur $\mu$ est $\simeq \mathbf{Z}/2\mathbf{Z}$, de
générateur $\Theta$ donné par
$$
\Theta(x) = \text{sgn}(x).x \quad \text{donc } \begin{cases} \Theta(\rho) = \rho \\ \Theta(\sigma) = -\sigma \end{cases} \leqno{(34)}
$$

\newpage

%% ===== Page 581 =====

Comparons avec les formules (20), on trouve
$$
u^{\mathfrak{S}^{+\wedge}/\Pi_0} = \varepsilon_2(u) : \begin{cases} \rho \mapsto \rho \\ \sigma \mapsto \varepsilon_2(u) \sigma \end{cases} \leqno{(35)}
$$
\footnote{Cet autom. agit au signe près sur les éléments de $\tilde{\Pi}_0/\Pi_0$, i.e.\ $u(x) = \text{sg}(x) \cdot x$ (car $\mathcal{M}_{0,3}^{\sim}$ est $\dots$).}i.e.
$$
[\dot{u}, \dot{\sigma}] = \varepsilon_2(u) \quad \text{dans } \mathfrak{S}^{+\wedge}/\Pi_0 \leqno{(35 \text{ bis})}
$$
(où $\dot{u}(\dot{\sigma}) \dot{\sigma}^{-1}$)

Reprenons la définition :
$$
Z = Z_{1, \sigma} \defeq \left\{ u \in \operatorname{Aut}(\mathfrak{S}^{+\wedge}) \mid u(S_3) = S_3, u(L_0) = L_0 \right\} \leqno{(36)}
$$
de sorte que pour $u \in Z$
$$
\begin{array}{l}
u(\rho) \text{ } \mathfrak{S}\text{-conjugué à } \rho^\varepsilon \\
u(\sigma) \text{ } \mathfrak{S}\text{-conjugué à } \sigma^{\varepsilon'} \quad (\varepsilon' = \pm 1) \\
u(L_0) = \varepsilon_u L_0
\end{array} \leqno{(37)}
$$
avec $\varepsilon \in \hat{\mathbf{Z}}^*$ et $\varepsilon', \varepsilon_u \in \{\pm 1\}$ bien déterminés.

(Car $Z \to \mathfrak{S}^{+\wedge}$ est injectif, $\rho$ et $\rho^{-1}$ ne sont pas $\mathfrak{S}$-conjugués, ni $\sigma$ et $\sigma^{-1}$). Je dis qu'un tel $u$ provient de $\mathcal{M}_{0,3}^{\sim}$ - il reste d'utiliser un scindage de $p_{0,3}$ sur le stabilisateur $\tilde{\pi}_0$, qui est engendré par $l_0, l_1, l_\infty$.
On a en effet
$$
u(l_0) = l_0^{\varepsilon_u} \in \pi^\wedge
$$
et
$$
u(l_1) = u({}^\rho(l_0)) = {}^{u(\rho)}(u(l_0)) = {}^{\rho^\varepsilon}(l_0^{\varepsilon_u})
$$
or l'action par conjugaison laisse stable $\pi^\wedge$,
donc
$$
u(l_1) \in \pi^\wedge
$$
de plus
$$
u(-1) = -1, \quad u(l_\infty) = - u(l_1 l_0)^{-1} = - u(l_0^{-1}) u(l_1)^{-1}
$$
$u(L_0) = L_0$, OK.

Si on identifie $u$ à un élément de $\mathcal{M}_{0,3}^{\sim}$, on voit que $u$ est en fait un élément [unclear: de $\mathcal{M}_{0,3}$].

\newpage

%% ===== Page 582 =====

$u \in \mathcal{M}_{0,3}^{\sim}[0] \simeq M'_0[0] \subset \mathcal{M}$\footnote{Attention: on aurait l'air dans ce qui suit d'identifier $\mathcal{M}_{0,3}^{\sim}[0] \subset \mathcal{M}$ seulement à $M'_0$. Or il s'avère dans ce qui suit que c'est $M_0$. p. ex. le degré de $\rho$ dans $M_0$ est 1, l'image de $\sigma$ (détailler ça). Corriger dans le texte qui précède ! En fait $M'_0 = M_0 \cap \mathcal{M}^!$, $Z_{\hat{\pi}^0} = M(0) =$ ss-groupe inv $\mathcal{M}!$}

donne (20), de façon plus précise
$$
\begin{cases}
u(\rho) = \operatorname{int}(\alpha) \rho \\
u(\sigma) = \varepsilon_2 \operatorname{int}(\beta) \sigma \\
u(\varepsilon_0) = \varepsilon_0^\mu
\end{cases}
\quad \text{avec} \quad
\begin{cases}
\alpha \in \hat{\pi}_{0,\tau}, \, \beta \in \hat{\pi}_0 \\
\mu \in \hat{\mathbf{Z}}^* \\
(\varepsilon_\tau = \varepsilon_\tau(\mu), \, \varepsilon_2(\alpha) = \varepsilon_3(\mu))
\end{cases} \leqno{(38)}
$$
bien tous impliquent (37). En fait les formules pour $u(\rho), u(\sigma), u(\varepsilon_0)$ déterminent de façon unique $\mu \in \hat{\mathbf{Z}}^*, \varepsilon_2 \in \{\pm 1\}, \alpha \in \hat{\pi}_{0,\tau}, \beta \in \hat{\pi}_0$, et on a nécessairement alors
$$
\varepsilon_2 = \varepsilon_2(\mu), \quad \varepsilon_\tau(\alpha) = \varepsilon_3(\mu) \leqno{(39)}
$$
où $\varepsilon_2$ est l'hom que [unclear: l'on] déduit de la suite exacte
$$
1 \to \hat{\pi}_0 \to \Pi_{0,\tau} \xrightarrow{\varepsilon_2} \{\pm 1\} \to 1 . \leqno{(40)}
$$

Ainsi on a prouvé

\underline{Proposition} Dans le cas universel ($\overline{\mathfrak{S}} = \widehat{\mathfrak{S}}^{+\wedge}$)
on a
$$
Z_{\hat{\pi}^0} \simeq \mathcal{M}'[0] \defeq \mathcal{M}_{0,3}^{\sim}[0] ,
$$
les composantes $u, \mu, \varepsilon, \varepsilon'$ d'un élément de $Z_{\hat{\pi}^0}$ correspondent resp. (en termes de $u_0 \in \mathcal{M}[0]$) :
$u_0$ lui-même : un él. de $\operatorname{Aut}(\widehat{\mathfrak{S}}^{+\wedge})$,
$\mu = \chi(u_0)$ (le multiplicateur de $u_0$), $\varepsilon = \varepsilon_3(\mu), \varepsilon' = \varepsilon_2(\mu)$.

\underline{Corollaire} On a $\mathfrak{S}_{\hat{\pi}^0} \simeq \mathcal{M} = \mathcal{N}_{1,1}^{\sim}$,
$$
\mathcal{P}_{\hat{\pi}^0} \simeq \mathcal{M}/\widehat{\mathfrak{S}}^{+\wedge} = \mathcal{N}_{1,1}^{\sim}/\widehat{\mathfrak{S}}^{+\wedge} \simeq \mathcal{M}_{0,3}^{\sim}/\widehat{\mathfrak{S}}^{+\wedge}_{0,3} \simeq \boldsymbol{\Gamma}_{\mathbf{Q}}^{\sim} .
$$

\newpage

%% ===== Page 583 =====

Remarques. Il est tentant d'écrire les (deux premières) relations (38) sous la forme
$$
\begin{cases}
(\alpha^{-1} u)(\rho) = \rho & \text{i.e. } \alpha^{-1} u \in Z_\rho^1 \\
(\beta^{-1} u)(\sigma) = \varepsilon_2 \sigma & \text{soit } \beta^{-1} u \in \mathcal{N}_\sigma'
\end{cases} \leqno{(41)}
$$
et d'introduire alors
$$
\begin{cases}
a = \alpha^{-1} u \in Z_\rho^1 \text{, d'où } u = \alpha a, \; \alpha = u a^{-1} \\
b = \beta^{-1} u \in \mathcal{N}_\sigma' \text{, d'où } u = \beta b, \; \beta = u b^{-1}
\end{cases} \leqno{(42)}
$$
et les développements du § 48 (notamment XI) peuvent faire s'attendre que les quantités $a, b$, pour $u \in \mathcal{M}[0]$ variables, dépendent multiplicativement de $u$. Nous allons voir cependant qu'il n'en est rien. Soit en effet
$$
\delta \colon \mathcal{G}_{\rho,\sigma} \isom \mathcal{M} \longrightarrow \Gamma_{\rho,\sigma} \isom \Gamma_{\mathbf{Q}}
$$
l'homomorphisme canonique, de sorte que l'on aura
$$
\delta(u) = \delta(b) = \tau_\gamma^{\nu_3} \delta(a) \leqno{(43)}
$$
en posant
$$
\tau_\gamma = \delta(\tau) \quad , \quad \varepsilon_3 = (-1)^{\nu_3} \quad (\nu_3 \in \mathbf{Z}/2\mathbf{Z})
$$
Si $u'$ correspond à $b', a', \nu_3'$, on aura
$$
\delta(u') = \delta(b') = \tau_\gamma^{\nu'_3} \delta(a') \leqno{(43 \text{ bis})}
$$
et si $u'' = u' u$ correspond à $a'', b'', \nu_3''$ :
$$
\delta(u'') = \delta(b'') = \tau_\gamma^{\nu''_3} \delta(a'') \leqno{(43 \text{ ter})}
$$
En fait, on aura
$$
b'' = b' b \quad , \quad \nu_2'' = \nu'_2 + \nu_2 \quad \text{(soit } \varepsilon_2'' = \varepsilon'_2 \varepsilon_2) \quad \text{et} \quad \nu_3'' = \nu'_3 + \nu_3 \quad (\text{soit } \varepsilon_3'' = \varepsilon'_3 \varepsilon_3) \leqno{(44)}
$$

\newpage

%% ===== Page 584 =====

On a
$$
u(\sigma) = \varepsilon_2 \operatorname{int}(\beta) \sigma
$$
$$
u'(\sigma) = \varepsilon'_2 \operatorname{int}(\beta') \sigma
$$
d'où
$$
(u'u)(\sigma) = u'(u(\sigma)) = \varepsilon_2 u'(\operatorname{int}(\beta)) u'(\sigma) = \varepsilon_2 \varepsilon'_2 \operatorname{int}(u'(\beta)\beta') \sigma
$$
avec $u'(\beta)\beta' \in \hat{\Pi}_{0'}$, ce qui par l'unicité prouve
$$
\beta'' = u'(\beta)\beta' \ , \ \ \varepsilon''_2 = \varepsilon'_2 \varepsilon_2 \leqno (45)
$$
d'où on conclut bien
$$
b'' = \underbrace{b'}_{\beta'^{-1}u'} \underbrace{b}_{\beta^{-1}u}
$$
i.e. $\beta''^{-1} u' u = \beta'^{-1} u' \beta^{-1} u$ car $\beta'' = u'(\beta)\beta'$ qui [unclear: est la rel.] (45).

La relation $\varepsilon''_3 = \varepsilon'_3 \varepsilon_3$ est aussi évidente p. ex. sur la relation (qui définit $\varepsilon_3$ et de même $\varepsilon'_3, \varepsilon''_3$)
$$
u(p) \text{ est } \widehat{\mathfrak{S}}^+\text{-conj. à } p^{\varepsilon_3} \leqno (46)
$$
Il reste donc à examiner la relation hypothétique
$$
a'' = a' a
$$
qui impliquerait
$$
\delta(a'') = \delta(a')\delta(a)
$$
mais comparant (43) (43 bis) (43 ter) on trouve la relation
$$
\delta(u')\delta(u) = \tau_{\gamma}^{\nu'_3} \delta(a') \tau_{\gamma}^{\nu_3} \delta(a)
$$
$$
\delta(u'u) = \tau_{\gamma}^{\nu''_3} \delta(a'')
$$
\hfill i.e.

\newpage

%% ===== Page 585 =====

de sorte que la relation (*) équivaut à
$$
\tau_{\gamma}^{\nu_3+\nu_3'} \delta(a'') = \tau_{\gamma}^{\nu_3'} \delta(a') \tau_{\gamma}^{\nu_3} \delta(a)
$$

i.e.
$$
\delta(a'') = \tau_{\gamma}^{\nu_3'} (\delta(a')) . \delta(a) \leqno{(47)}
$$

Dans la notation des opérateurs on a
\marginpar{puisque $\delta(u') = \tau_{\gamma}^{\nu_3'}(\delta(u'))$}
$$
\delta(a') = \tau_{\gamma}^{\nu_3'} (\delta(a')) \quad \text{i.e.} \quad [\delta(a'), \tau_{\gamma}^{\nu_3'}] = 1 \quad \text{car} \quad [\delta(u'), \tau_{\gamma}^{\nu_3'}] = 1 \quad \text{[unclear: tout court]} \leqno{(48)}
$$

Cela montre que si $\nu_3 = 1$ i.e. $\varepsilon_3 = -1$, alors la relation (*) équivaut à la commutativité de $\delta(a')$ avec $\tau_{\gamma}$, la même condition du commencement sur un élément de $\widetilde{\Gamma}_{\mathbf{Q}}$ qui, pour un élément de $\Gamma_{\mathbf{Q}}$ lui-même, implique que $w \in \{1, \tau_{\gamma}\}$ - et par des raisonnements heuristiques, on a cru qu'il devait en être ainsi pour tout élément de $\widetilde{\Gamma}_{\mathbf{Q}}$. 

\marginpar{Il semble qu'on ne puisse en tirer que $\Gamma_1' = \operatorname{Ker}(\dots)$}
Il apparaît que le succès de la présentation du § 48 \underline{XI} tient au fait que dans la situation envisagée là le groupe $\Gamma_0 \simeq \mathbf{Z}_p^*$ était commutatif, la quantité $a$ d'où $a \in \mathbf{Z}_p^*$ ou dépend multiplicativement de l'élément $u \in \mathcal{M}[0]$, sous la condition de base : $u \in \mathcal{M}[0]'$.

\newpage

%% ===== Page 586 =====

\marginpar{NB On a $S_0 \Gamma_{p,\sigma} / \tilde{\Gamma}_Q \simeq \hat{S}_0 \Gamma_{p,\sigma}$ dans $\Gamma_{p,\sigma}$}
Nous allons essayer de sauver la mise, en trouvant un homomorphisme canonique !
$$
\begin{tikzcd}
\tilde{M}[0] \arrow[r] & \hat{\mathbf{Z}}_{p,\sigma}^* \times_{\Gamma_{p,\sigma}} N_{\rho}^* \times_{\Gamma_{p,\sigma}} N_{\sigma}^* \\
& \hat{S}_0 \Gamma_{p,\sigma} \subset S_0 \Gamma_{p,\sigma} \text{ (indice 2)} \arrow[u, hookrightarrow, phantom]
\end{tikzcd}
$$
Notons que ce dernier ensemble est formé des triplets
$$
(u, a, b) \mid u = f b = f' a , \text{ avec } f, f' \in \hat{\Gamma}_Q \text{ convenables} \leqno{(49)}
$$
ce qui implique $\varepsilon_2(u) = \varepsilon_0(b)$, $\varepsilon_3(u) = \varepsilon_0(a)$
et pour un tel triplet, on aura
$$
\begin{cases}
u(\sigma) = \operatorname{int}(fb)(\sigma) = \operatorname{int}(f) \operatorname{int}(b)(\sigma) = \operatorname{int}(f)(\sigma^{\varepsilon_2(u)}) \\
u(\rho) = \operatorname{int}(f'a)(\rho) = \operatorname{int}(f') \operatorname{int}(a)(\rho) = \operatorname{int}(f')(\rho^{\varepsilon_3(u)}) = \operatorname{int}(f' \varepsilon_3(u))(\rho)
\end{cases} \leqno{(49 \text{ bis})}
$$
où on pose
$$
\varepsilon_3(u) = \begin{cases} 1 & \text{si } \varepsilon_3 = +1 \\ \rho^2 & \text{si } \varepsilon_3 = -1 \end{cases} \quad \text{i.e. } \varepsilon_3(u) = \rho^{\nu_3} \text{ avec } \varepsilon_3 = (-1)^{\nu_3}
$$
Si $u$ correspond à $\alpha, \beta$ (cf.\ (38)), ces
formules (49 bis) s'écrivent, vu les répartitions
$$
\begin{cases}
\operatorname{int}(f)(\sigma) = \operatorname{int}(\beta)(\sigma) & \text{i.e. } f = \beta \sigma^i \quad (i \in \mathbf{Z}/2\mathbf{Z}) \\
\operatorname{int}(f' \varepsilon_3(u))(\rho) = \operatorname{int}(\alpha)(\rho) & \text{i.e. } f' \varepsilon_3(u) = \alpha \rho^j \quad (j \in \mathbf{Z}/6\mathbf{Z})
\end{cases}
$$

Dans ce cas on a

Proposition. Pour un $u \in \tilde{Z}_{p,\sigma}^* = \tilde{M}[0] \simeq \text{[unclear: f \mapsto f_0]}$ fixé,
\marginpar{$\alpha, \beta, \varepsilon_3$ dépendent de $u$ par (38)}les $(u,a,b) \in S_0\Gamma_{p,\sigma}$ au dessus de $u$
sont exactement les éléments de la forme
$$
\left\{ (u,a,b) \mathrel{\Big|} \begin{matrix} a = f'^{-1} u \\ b = f^{-1} u \end{matrix} \text{ avec } \begin{matrix} f' = \alpha \rho^j \varepsilon_3(u)^{-1} & (j \in \mathbf{Z}/6\mathbf{Z}) \\ f = \beta \sigma^i & (i \in \mathbf{Z}/2\mathbf{Z}) \end{matrix} \right\} \leqno{(50)}
$$

Considérons alors l'élément $(u,a,b)$ au dessus de $u$,
l'expression la plus simple en termes
de $\alpha, \beta$, en prenant $i=0, j=0$, correspondant
disons aux choix :

\newpage

%% ===== Page 587 =====

$$
\begin{cases}
a = \varepsilon_3(\tau') \alpha^{-1} u & \text{i.e. } \alpha = u a^{-1} \varepsilon_3(\tau') \\
b = \beta^{-1} u & \text{i.e. } \beta = u b^{-1}
\end{cases}
\leqno{(51)}
$$

je me demande si cette section ensembliste de $S_{0}\Gamma'_{p,\sigma}$ au dessus de $Z'_{p,\sigma} = M'_{[\tau_{0}]}$ est multiplicative, i.e.\ ce qui revient au même, si $\alpha$ et $\beta$ définies par (51) dépendent multiplicativement de $u$.

Proposition. Soit $(u,a,b) \in S_{0}\Gamma'_{p,\sigma}$ (cf (50)). Pour qu'on ait $i=0, j=0$, i.e.\ pour qu'on ait (51) --- i.e.\ pour que $(u,a,b)$ appartienne à la section ensembliste précédente, il faut et il suffit qu'on ait les relations\footnote{On aura posé $\varepsilon_3(\sigma)=1$ si $\varepsilon_3=+1$, et $\varepsilon_3(\sigma)= \sigma^{-1}$ si $\varepsilon_3=-1$. NB $\tau'\tau = \sigma^{-1}$ ; $\tau'\sigma^{-1} = \tau$.}
$$
u \equiv b \equiv \varepsilon_3(\sigma) a \pmod{\widehat{\Pi}_0} \leqno{(52)}
$$

Démonstration. Comme $\beta$ ($= u b^{-1}$) et $\alpha \varepsilon_3(\tau)$ ($= u a^{-1} \varepsilon_3(\tau') \varepsilon_3(\tau) = u a^{-1} \varepsilon_3(\sigma^{-1})$) sont tous deux dans $\widehat{\Pi}_0$, la condition est nécessaire. Pour la suffisance, on note que sous (52), il s'ensuit que l'on a
$$
f, f' \varepsilon_3(\tau') \in \widehat{\Pi}_0 \leqno{(53)}
$$
et réduisant mod $\widehat{\Pi}_0$ les deux relations (50), la première s'écrit $\sigma^i \equiv 1$ et donne $i=0$ (car $\beta \in \widehat{\Pi}_0$), la seconde donne $f' \varepsilon_3(\tau') \equiv 1$, i.e.\ $\alpha \rho^j \varepsilon_3(\tau') \equiv 1$ i.e.\ $\varepsilon_3(\tau) \alpha \rho^j \equiv 1$, d'où $\rho^j \equiv 1$ et $j=0$ (car $\varepsilon_3(\tau) \alpha \in \widehat{\Pi}_0$). 

Corollaire. La section envisagée est multiplicative. \hfill cqfd.

\newpage

%% ===== Page 588 =====

Soient maintenant $u,u'$, correspondant à $a,b$, $\varepsilon_3$ et $u'$, correspondant à $a',b'$, $\varepsilon'_3$, soit $u''=u'u$, correspondant à $a'', b''=b'b$ et $\varepsilon''_3 = \varepsilon'_3 \varepsilon_3$, de sorte qu'on aura
$$
\begin{aligned}
u &\equiv \varepsilon_3(\sigma) a \quad (\hat{\Pi}_0) \\
u' &\equiv \varepsilon'_3(\sigma) a' \quad (\hat{\Pi}_0) \\
u'' &\equiv \varepsilon''_3(\sigma) a'' \quad (\hat{\Pi}_0) \\
|| & \\
u'u &\equiv \varepsilon \, \varepsilon'_3(\sigma) \varepsilon_3(\sigma) a'' \quad \text{où } \varepsilon \in \{\pm 1\}, \, \varepsilon = -1 \text{ ssi } \varepsilon_3 = -1 \text{ et } \varepsilon'_3 = -1 \\
|| & \\
(\varepsilon'_3(\sigma) a')(\varepsilon_3(\sigma) a) &
\end{aligned}
$$
\marginpar{De $\begin{cases} \sigma(u') \equiv u' \\ \sigma(u') \equiv -u' \end{cases}$ i.e.}

on aura donc
$$
a' \varepsilon_3(\sigma) a \equiv \varepsilon \, \varepsilon_3(\sigma) a'' \quad (\hat{\Pi}_0) \quad \text{i.e.}
$$
\marginpar{NB $\varepsilon_3(\sigma)$ et $\varepsilon'_3(\sigma^{-1})$ opèrent sur $\hat{\Pi}_0$}
$$
a'' \equiv \varepsilon \, \varepsilon_3(\sigma)(a') \cdot a \leqno{(53)}
$$

Donc la relation escomptée :
$$
a'' = a' a \quad \text{équivaut à} \quad a'' \equiv a' a \quad (\hat{\Pi}_0)
$$
soit
\marginpar{54 bis \\ 1) $\varepsilon_3=1, \varepsilon'_3=1 : a'=a'$ OK \\ 2) $\varepsilon_3=1, \varepsilon'_3=-1 : a'=a'$ OK \\ 3) $\varepsilon_3=-1, \varepsilon'_3=1 : \sigma(a')=-a'$ \\ 4) $\varepsilon_3=-1, \varepsilon'_3=-1 : \sigma(a')=-a'$}
$$
\varepsilon_3(\sigma)(a') \equiv \varepsilon a' \quad (\hat{\Pi}_0) \leqno{(54)}
$$

Cette condition est automatiquement vérifiée.

\newpage

%% ===== Page 589 =====

$$
\varepsilon_3(\sigma)(a') \equiv \varepsilon a' \quad (\hat{\Pi}_0) \leqno{(54)}
$$
$$
\begin{cases}
\varepsilon \in \{\pm 1\}, \ \varepsilon = -1 \text{ ssi } \varepsilon_3 = \varepsilon_3' = -1 \\
\varepsilon_3(\sigma) = \begin{cases} \id/1 & \text{si } \varepsilon_3 = +1 \\ \sigma & \text{si } \varepsilon_3 = -1 \end{cases}
\end{cases}
$$

Cette condition (54) s'explicite suivant les cas :

1/ si $\varepsilon_3 = 1$, alors $\varepsilon = 1$ et (54) est toujours vérifiée

2/ si $\varepsilon_3 = -1$, alors (54) s'écrit
$$
\sigma(a') \equiv \varepsilon_3' \cdot a' \quad (\hat{\Pi}_0) \leqno{(55)}
$$
i.e.
$$
[\sigma, a'] = \varepsilon_3'
$$
où les $[\dots]$ désignent les classes dans $\mathcal{M}/\hat{\Pi}_0$.

Comme on a 
$$
a' \equiv \varepsilon_3'(\sigma)^{-1} u' = \begin{cases} u' & \text{si } \varepsilon_3' = 1 \\ \sigma^{-1} u' & \text{si } \varepsilon_3' = -1 \end{cases}
$$
on trouve que la condition (55) équivaut à :
$$
[\sigma, u'] = \varepsilon_3' \leqno{(56)}
$$
(valable aussi bien pour $\varepsilon_3' = 1$ que $\varepsilon_3' = -1$).

Or on a vu en (35) que l'on a dans $\mathcal{M}/\hat{\Pi}_0$ :
$$
u'(\sigma) = \varepsilon_2(u') \sigma
$$
i.e.
$$
[u', \sigma] = \varepsilon_2(u') \quad \text{ou encore}
$$
$$
[\sigma, u'] = \varepsilon_2(u') \leqno{(57)}
$$

Comparant (56) et (57), on voit que (56) prend la forme équivalente
$$
\varepsilon_2(u') = \varepsilon_3' \leqno{(58)}
$$

On a donc prouvé :

\marginpar{NB
$\begin{cases} \varepsilon_3 = \varepsilon_2(u) = \varepsilon_0^p(a) \\ \varepsilon_3' = \varepsilon_2(u') = \varepsilon_0^p(a') \\ \varepsilon_3'' = \varepsilon_2(u'') = \varepsilon_0^p(a'') \end{cases}$}

\textbf{Proposition.} Soient $u, u'$ et $u''=u'u \in \mathcal{M}/\hat{\Pi}_0 \isommap \hat{\mathbf{Z}}'_{pair}$,
correspondant à $\varepsilon_3, \varepsilon_3', \varepsilon_3'' = \varepsilon_3' \varepsilon_3 \in \{\pm 1\}$, et :
$$
a = \varepsilon_3(\sigma) \alpha^{-1} u, \quad a' = \varepsilon_3'(\sigma) \alpha'^{-1} u', \quad a'' = \varepsilon_3''(\sigma) \alpha''^{-1} u'',
$$
avec $\alpha, \alpha', \alpha'' \in \hat{\mathcal{N}}_{\overline{P}}, \quad \alpha \equiv \varepsilon_3(\sigma)^{-1} u, \alpha' \equiv \dots, \alpha'' \equiv \dots$
les conditions suivantes sont équivalentes :

\newpage

%% ===== Page 590 =====

\begin{enumerate}
\item[a)] On a \qquad $a'' = a' a$
\item[a')] On a \qquad $a'' \equiv a'a \quad (\hat{\Pi}_0)$ \marginpar{i.e. $a(\rho)=\rho$}
\item[b)] On a \qquad $\varepsilon_3 = 1$ ou $\sigma(a') \equiv \varepsilon'_3 a' \quad (\hat{\Pi}_0)$
\item[b')] On a \qquad $\varepsilon_3 = 1$ ou $a'(\sigma) \equiv \varepsilon'_3 \sigma \quad (\hat{\Pi}_0)$ \marginpar{i.e. $a(\rho)=\rho$}
\item[c)] On a \qquad $\varepsilon_3 = 1$ ou $\sigma(u') \equiv \varepsilon'_3 u' \quad (\hat{\Pi}_0)$
\item[c')] On a \qquad $\varepsilon_3 = 1$ ou $u'(\sigma) \equiv \varepsilon'_3 \sigma \quad (\hat{\Pi}_0)$ \marginpar{i.e. $a(\rho)=\rho$}
\item[d)] On a \qquad \fbox{$\varepsilon_3 = 1$ ou $\varepsilon'_2 = \varepsilon'_3 \text{ i.e. } \varepsilon'_2 \varepsilon'_3 = 1$}
\end{enumerate}

Corollaire. Au dessus des deux sous-groupes
$$
Z^{!\sharp}_{\rho,\sigma} = M'^{!}_{[\tau_0]} \defeq \{ u \in Z^!_{\rho,\sigma} \mid \varepsilon_3(u) = 1 \}
$$
$$
Z^{!\flat}_{\rho,\sigma} = M''^{!}_{[\tau_0]} \defeq \{ u \in Z^!_{\rho,\sigma} \mid \varepsilon_2(u) = \varepsilon_3(u) \} \leqno{(59)}
$$
l'application quotient ensembliste
$$
\begin{tikzcd}
A \colon Z^!_{\rho,\sigma} \arrow[r] & S_0 \Gamma^!_{\rho,\sigma} \\
u \arrow[r, maps to, "\text{cf prop p. 597}"'] & (u, a = \varepsilon_3(\tau) \alpha^{-1} u, b = \beta^{-1} u)
\end{tikzcd} \leqno{(60)}
$$
est une section multiplicative. Mais elle n'est pas multiplicative au dessus de $\mathcal{M}^!_{[\tau_0]}$ tout entier.

\newpage

%% ===== Page 591 =====

\noindent
\underline{Remarque} Les éléments de l'im.\ [unclear: par l'isom.] $S_0\Gamma^!_{\rho,\sigma}$ de $Z^!_{\rho,\sigma}$ dans $S_0\Gamma_{\rho,\sigma}$ qui appartiennent à la section envisagée sont caractérisés (par vertu de la prop.\ p.\ 598) par la relation
$$
u \equiv b \equiv a \mod \widehat{\Pi}_0 \qquad \left( \begin{matrix} \text{en plus de } \varepsilon_3(u)=1 \\ \text{qui caractérise les} \\ \text{éléments de } S_0\Gamma^!_{\rho,\sigma} \end{matrix} \right) \leqno{(61)}
$$
qui rend évident que c'est un sous-groupe. Je n'ai pas trouvé de jolie expression analogue pour l'appartenance d'un élément de $S_0\Gamma^{!!}_{\rho,\sigma}$ (caractérisé par la relation $\varepsilon_2(u)\varepsilon_3(u)=1$) à la section correspondante.

\vspace{0.5cm}
\noindent
\underline{Proposition} Soit $S^{\text{II}}_0\Gamma^!_{\rho,\sigma} \subset S_0\Gamma_{\rho,\sigma}$ le sous-ensemble de $S_0\Gamma_{\rho,\sigma}$
$$
S^{\text{II}}_0\Gamma^!_{\rho,\sigma} \subset S_0\Gamma_{\rho,\sigma} \leqno{(62)}
$$
$$
\|
$$
$$
\left\{ (u,a,b) \in \underset{M[\Gamma_0]}{Z_{\rho,\sigma} \times \hat{\pi} \times \hat{\pi}} \;\middle|\; \begin{matrix} u \equiv b \ (\widehat{\Pi}_0) \\ u \equiv \varepsilon_3(\sigma) a \quad (\widehat{\Pi} = \widehat{\Pi}_0 \rtimes \{\pm 1\}) \end{matrix} \right\}
$$

Alors
\begin{enumerate}
\item[a)] $S^{\text{II}}_0\Gamma^!_{\rho,\sigma}$ est un sous-groupe de $S_0\Gamma_{\rho,\sigma}$
\item[b)] L'homomorphisme canonique $S^{\text{II}}_0\Gamma^!_{\rho,\sigma} \to Z^!_{\rho,\sigma} = \mathcal{M}^{!}_0[\Gamma_0]$,
$(u,a,b) \mapsto u$, est surjectif, et son noyau
est le sous-groupe isomorphe à $\mu = \{\pm 1\}$
des éléments $(1, \varepsilon, 1)$ où $\varepsilon \in \mu$
\end{enumerate}

\newpage

%% ===== Page 592 =====

Dém. Le fait qu'on a un sous-groupe
résulte des calculs déjà faits p. 597, 598,
la condition que le produit de deux
éléments $(u',a',b')$ et $(u,a,b)$ de $S^{\mathbf{P} !}_{\Gamma,\sigma}$ y soit
s'exprime encore par la formule
$$
[\dot{\sigma},\dot{u}'] = \dot{\varepsilon}'_3 \quad \text{dans } M/\hat{\pi} \quad \text{(au lieu de } M/\hat{\pi}_0)
\leqno{(56)}
$$
$$
[\dot{u}',\dot{\sigma}] = \dot{\varepsilon}'_3 \qquad \text{où } \quad \dot{\varepsilon}'_3 = \dot{\varepsilon}_3
\leqno{(56 \text{bis})}
$$
avec $[\dot{u}',\dot{\sigma}] = \dot{u}'(\dot{\sigma})\dot{\sigma}^{-1}$.

Or dans $\mathfrak{S}/\hat{\pi}_0$, [unclear: les images] de $s_1$ et $s_3$
sont contenues dans le sous-groupe
central $\hat{\pi}/\hat{\pi}_0$, dont ils sont triviaux dans
$\mathfrak{S}/\hat{\pi}$, donc la relation (56 bis) est trivialement
satisfaite.

La surjectivité de l'application $S^{\mathbf{P} !}_{\Gamma,\sigma} \to Z^{! \sigma}_{\Gamma}$
provient du fait que $S^{\mathbf{P} !}_{\Gamma,\sigma}$ contient la
section ensembliste $A(Z^{! \sigma}_{\Gamma})$ de la prop.
p. 597. Le noyau est formé des
$(1, a, b)$ avec $a \in N^*_{\Gamma} \cap \hat{\pi}$, $b \in N^*_{\sigma} \cap \hat{\pi}_0$.
Or $N^*_{\Gamma} \cap \hat{\pi} = Z_{\Gamma}$, $N^*_{\sigma} \cap \hat{\pi} = E_{\sigma}$, et
$$
L_{\Gamma} \cap \hat{\pi} = L_{\sigma} \cap \hat{\pi} = \mu, \quad L_{\Gamma} \cap \hat{\pi}_0 = L_{\sigma} \cap \hat{\pi}_0 = \{1\}
\leqno{(57)}
$$

\newpage

%% ===== Page 593 =====

d'où la caractérisation du noyau.

Corollaire Considérons $\mathfrak{S}^{\pi !}_{\rho,\sigma}$ comme
une extension de $Z^!_{\rho,\sigma} = M'_{[\sigma]}$ par
$L_\rho \times L_\sigma \simeq \mathbf{Z}/6\mathbf{Z} \times \mathbf{Z}/4\mathbf{Z}$, donc
(en écrivant [unclear: $\mathbf{Z}/6\mathbf{Z}$] $\simeq \mathbf{Z}/3\mathbf{Z} \times \mathbf{Z}/2\mathbf{Z}$)
une extension de $\underset{\mu}{\mathbf{Z}/2\mathbf{Z}} \times \underset{L_\rho}{\mathbf{Z}/3\mathbf{Z}} \times \underset{L_\sigma}{\mathbf{Z}/4\mathbf{Z}}$.

Cette extension correspond donc canoniquement
à trois extensions, par $\mu$, $L_\rho$, $L_\sigma$
respectivement. Les extensions par $L_\rho \times L_\sigma$
admettent des sections canoniques.

Il reste une extension (c'est juste $\mathfrak{S}^{\pi !}_{\rho,\sigma}$) de $Z^!_{\rho,\sigma}$ par $\mu$.

Cette extension correspond justement à
l'ambiguïté de signe pour choisir
un $\alpha \in \hat{\pi}$ au dessus de $\tilde{\alpha} \in \hat{\pi}/\pm 1$
défini par $u \in Z^!_{\rho,\sigma} = M'_{[\sigma]}$, qui se reflète
par l'ambiguïté à choisir le signe
de $a \in N^*_\rho$, lié à $\alpha$ par (51)
$$
a = \varepsilon_3(2i) \alpha^{-1} u, \quad \alpha = u a^{-1} \varepsilon_3(2i) \leqno{(58)}
$$
Nous avons pu lever cette ambiguïté
par l'exigence $\alpha \in \hat{\pi}_{0 \pm} = \hat{\pi}_{0, \{1, 2\}} \cdot \pi_3$, et nous

\newpage

%% ===== Page 594 =====

venons de constater que ce choix, pour
canonique qu'il soit, n'est pas multiplica-
tif. En même temps, on trouve
l'explicitation complète des $\alpha''$ (normalisé
canoniquement par la condition $\alpha'' \in \hat{\Pi}_{0\overline{2}}$)
associé à $u''=u'u$, en termes de
$\alpha'$ (associé à $u'$) et $\alpha$ (associé à $u$)
(en termes de $a', a$).
(En effet, la page
p 591 nous donne dans quels cas on a
$$
a'' = a'a \leqno{(59)}
$$
savoir ssi on a
$$
\varepsilon_3(u) = 1 \quad \text{ou} \quad \varepsilon_3(u')\varepsilon_3(u)=1 \leqno{(59\text{ bis})}
$$
Mais on sait maintenant qu'en tous
les cas doit valoir
$$
a'' = \varepsilon a'a \qquad \varepsilon \in \{\pm 1\} \leqno{(60)}
$$
et on trouve maintenant que $\varepsilon = +1$
ssi on est dans le cas (59 bis), $-1$
dans le cas contraire i.e.\ si $\varepsilon_3(u)=-1$
et $\varepsilon_3(u')\varepsilon_3(u) = +1$. En résumé -

\underline{Corollaire} Soient $u, u' \in M_{1 \to 0}^!$, d'où $u'' = u'u$,
soient $\alpha, \alpha', \alpha'' \in \hat{\Pi}_{0\overline{2}}$ associés (définis
par [unclear: $a(\rho)=\dots$ etc]) et posons

\newpage

%% ===== Page 595 =====

$$
a = \varepsilon_3(\tau') \alpha^{-1} u, \quad a' = \varepsilon'_3(\tau') \alpha'^{-1} u', \quad a'' = \varepsilon''_3(\tau') \alpha''^{-1} u'' \leqno (61)
$$
où
$$
\varepsilon_3(\tau') = \begin{cases} 1 & \text{si } \varepsilon_3 = 1 \\ \tau' & \text{si } \varepsilon_3 = -1 \end{cases} \quad \text{et de même pour } \varepsilon'_3(\tau'), \varepsilon''_3(\tau'),
$$
et où $\varepsilon_3 = \varepsilon_3(u) = \varepsilon_3(X(u))$, et de même pour $\varepsilon'_3, \varepsilon''_3$
on a $\varepsilon''_3 = \varepsilon_3 \varepsilon'_3$. On a donc
$$
a'' = \varepsilon a' a, \quad \text{avec } \varepsilon \in \{\pm 1\} \leqno (62)
$$
et
$$
\begin{cases} \varepsilon = +1 & \text{ssi } \varepsilon_3 = 1 \text{ ou } \varepsilon'_3 = 1 \\ \varepsilon = -1 & \text{ssi } \varepsilon_3 = -1 \text{ et } \varepsilon'_3 = -1 \end{cases} \leqno (63)
$$

$$
\alpha'' = \varepsilon u'(\alpha) \alpha'\footnote{pas vrai en général} \quad \text{où } \varepsilon \text{ donné par } (63) \leqno (64)
$$
En effet
$$
\begin{cases} \alpha = u a^{-1} \varepsilon_3(\tau') \\ \alpha' = u' a'^{-1} \varepsilon'_3(\tau') \end{cases} \quad \text{d'où}
$$
$$
u'(\alpha) = u' u a^{-1} \varepsilon_3(\tau') u'^{-1}
$$
$$
u'(\alpha)\alpha' = u' u a^{-1} \varepsilon_3(\tau') \cancel{u'^{-1} u'} a'^{-1} \varepsilon'_3(\tau')
$$

La formule (64), en vertu de $\ast)$, prend donc la forme
$$
a'^{-1} \varepsilon_3(\tau') = \varepsilon_3(\tau') a'^{-1}
$$
i.e.
$$
[\bar{\varepsilon}_3(\tau'), a'^{-1}] = 1
$$
elle n'est pas toujours vérifiée si $\varepsilon_3 = -1$
i.e. $\varepsilon_3(\tau') = \tau'$ : elle équivaut à une restriction draconienne sur $a'$. La formule toujours correcte est
$$
\alpha'' = (\varepsilon u'(\alpha)\alpha') \cdot \underbrace{\varepsilon'_3(\tau') \big([a', \varepsilon_3(\tau')]\big)}_{[ \varepsilon'_3(\tau')(a'), \varepsilon_3(\tau') ]} \leqno (65)
$$

\newpage

%% ===== Page 596 =====

Comme $\varepsilon_3'(\tau') a' = \alpha'{}^{-1} u'$ (61), la formule un peu alambiquée :\marginpar{Sic ! Voilà le monstre qui m'a fait froid dans le dos l'autre mois -- mais il n'en peut plus, le pauvre vieux !}
$$
\alpha'' = (\varepsilon u'(\alpha)\alpha') \cdot \left[ \underset{\overset{\uparrow}{\varepsilon_3'(\tau')a'}}{\alpha'{}^{-1}u'}, \varepsilon_3(\tau') \right] \leqno{(65)}
$$
(où $\varepsilon \in \{\pm 1\}$ défini par (63))

Il reste à examiner la question si l'extension de $S_0^\pi \Gamma_{p,\sigma}$ par $Z^1_{p,\sigma} \simeq M (\simeq \mu = \{\pm 1\})$ est triviale ou non, compte tenu du fait que ses restrictions à $Z'_{p,\sigma}$ et à $Z'''_{p,\sigma}$ sont triviales.
Soient $S_0^\pi \Gamma'_{p,\sigma}$, $S_0^\pi \Gamma'''_{p,\sigma}$ les images inverses de $Z'_{p,\sigma}$, $Z'''_{p,\sigma}$ dans $S_0^\pi \Gamma_{p,\sigma}$, donc
$$
\begin{cases}
S_0^\pi \Gamma'_{p,\sigma} = S_0^\pi \Gamma_{p,\sigma} \cap \Gamma'_{p,\sigma} \\
S_0^\pi \Gamma'''_{p,\sigma} = S_0^\pi \Gamma_{p,\sigma} \cap \Gamma'''_{p,\sigma}
\end{cases} \leqno{(66)}
$$
et on a des suites exactes d'extensions centrales
$$
\begin{cases}
1 \to \mu \to S_0^\pi \Gamma'_{p,\sigma} \to Z'_{p,\sigma} \to 1 \\
1 \to \mu \to S_0^\pi \Gamma'''_{p,\sigma} \to Z'''_{p,\sigma} \to 1
\end{cases} \leqno{(67)}
$$
qui sont scindées par des sous-groupes déjà répertoriés $S_0^{\pi_0} \Gamma'_{p,\sigma}$, $S_0^{\pi_0} \Gamma'''_{p,\sigma}$, de sorte que
$$
\begin{cases}
S_0^\pi \Gamma'_{p,\sigma} \simeq S_0^{\pi_0} \Gamma'_{p,\sigma} \times \mu \\
S_0^\pi \Gamma'''_{p,\sigma} \simeq S_0^{\pi_0} \Gamma'''_{p,\sigma} \times \mu
\end{cases} \leqno{(68)}
$$

\newpage

%% ===== Page 597 =====

Je me pose la question si ces sous-groupes\marginpar{ss-groupes d'indices 4 dans $\mathfrak{S}^\Pi \Gamma_{\rho, \sigma}^!$}
$$
\begin{cases}
\mathfrak{S}_0^\Pi \Gamma_{\rho, \sigma}^{!!} \subset \mathfrak{S}^\Pi \Gamma_{\rho, \sigma}^{!!} \subset \mathfrak{S}^\Pi \Gamma_{\rho, \sigma}^! \\
\mathfrak{S}_0^\Pi \Gamma_{\rho, \sigma}^{!!!} \subset \mathfrak{S}^\Pi \Gamma_{\rho, \sigma}^{!!!} \subset \mathfrak{S}^\Pi \Gamma_{\rho, \sigma}^!
\end{cases} \leqno{(69)}
$$
sont invariants dans $\mathfrak{S}^\Pi \Gamma_{\rho, \sigma}^!$, ce qui
signifierait (en comparant avec les quotients
par ces sous-groupes) que l'extension
$\mathfrak{S}^\Pi \Gamma_{\rho, \sigma}^!$ de $Z_{\rho, \sigma}^!$ par $\mu$, scindée sur
$Z_{\rho, \sigma}^{!!}$ et sur $Z_{\rho, \sigma}^{!!!}$, proviennent en fait
(avec lesdits scindages) d'extensions de
$Z_{\rho, \sigma}^! / Z_{\rho, \sigma}^{!!} \ (\simeq \pm 1)$ et de $Z_{\rho, \sigma}^! / Z_{\rho, \sigma}^{!!!} \ (\simeq \pm 1)$
par $\mu$, qu'il faudrait ensuite
caractériser (comme triviales ou non triviales).

Comme les sous-groupes en question
sont distingués dans $\mathfrak{S}^\Pi \Gamma^{!!}$ resp. $\mathfrak{S}^\Pi \Gamma^{!!!}$, et
si $x$ est un élément de $\mathfrak{S}^\Pi \Gamma^{!}$ qui n'est
pas dans $\mathfrak{S}^\Pi \Gamma^{!!}$ resp. $\mathfrak{S}^\Pi \Gamma^{!!!}$, dire que
$\mathfrak{S}_0^\Pi \Gamma^{!!}$ (resp. $\mathfrak{S}_0^\Pi \Gamma^{!!!}$) est distingué dans
$\mathfrak{S}^\Pi \Gamma^{!}$ et [unclear: cela] signifie simplement qu'il
est normalisé par $x$.

1º) Cas de $\mathfrak{S}_0^\Pi \Gamma^{!!}$ dans $\mathfrak{S}^\Pi \Gamma^{!}$. On prendra
$$
\begin{cases}
x = (u_0, a_0, b_0) \text{ avec } u_0 = \tau, \ d\text{où } \varepsilon_3(u_0) = \varepsilon_4(u_0) = \cancel{\chi(u_0)} = -1 \\
( \cancel{x = \tau, \text{ qui est}} \text{ dans } \mathfrak{S}_0^\Pi \Gamma^! ) \\
\cancel{a_0 = \tau', b_0 = 1} \ a'_0 = \varepsilon'_3(\tau') a_0^{-1} u_0 = \tau', \ b_0 = \beta''_0 u_0 = \tau \\
\text{donc } \underline{x = (u_0 = \tau, a_0 = \tau', b_0 = \tau)}
\end{cases} \leqno{(70)}
$$

\newpage

%% ===== Page 598 =====

$x \in \mathfrak{S}_0^{\pi_0} \Gamma^{!\prime\prime}$ et $x \notin \mathfrak{S}_0^{\pi_0} \Gamma^{!\prime}$
et il faut voir si $x$ normalise $\mathfrak{S}_0^{\pi_0} \Gamma^{!\prime}$.

Soit
$$
y = (u,a,b) \in \mathfrak{S}_0^{\pi_0} \Gamma^{!\prime} \quad \text{i.e.} \quad \begin{cases} u \equiv a \equiv b \pmod{\Pi_0} \\ \varepsilon_3(u) = 1 \end{cases} \leqno{(71)}
$$
et calculons
$$
x y x^{-1} = x y x = (\tau(u), \tau'(a), \tau(b))
$$
est-il dans $\mathfrak{S}_0^{\pi_0} \Gamma^{!\prime}$ ? [unclear: C'est un élément de] $\mathfrak{S}_0^{\pi} \Gamma^{!\prime}$ i.e. a-t-on
$$
\tau(u) \equiv \tau'(a) \equiv \tau(b) \pmod{\Pi_0} \quad ? \leqno{(72)}
$$

La relation $\tau(u) \equiv \tau(b) (\equiv \tau(a))$ est conséquence immédiate de 71), et 72) équivaut à
$$
\tau'(a) \equiv \tau(a) \pmod{\Pi_0} \quad \text{i.e.} \quad \sigma^{-1}(\tau(a)) \equiv \tau(a)
$$
\marginpar{NB}
$$
\text{i.e.} \quad [\dot{\sigma}, \tau(a)^\cdot] = 1 \quad \text{dans } \mathfrak{S}^{\pi} / \Pi_0 \leqno{(73)}
$$

i.e. $\varepsilon_2(\tau(a)) = 1$ i.e. $\varepsilon_2(a) = 1$
on a enfin (comme $a \equiv u \pmod{\Pi_0}$)
$$
\varepsilon_2(u) = 1 . \leqno{(74)}
$$

Cette condition n'est pas conséquence de la condition $\varepsilon_3(u) = 1$, de sorte que $\mathfrak{S}_0^{\pi_0} \Gamma^{!\prime}$ n'est pas invariant dans $\mathfrak{S}_0^{\pi} \Gamma^{!}$.

$\mathfrak{S}_0^{\pi_0} \Gamma^{!\prime\prime}$ est-il invariant dans $\mathfrak{S}_0^{\pi} \Gamma^{!}$ ?

Soit
$$
x = (u_0, a_0, b_0) \in \mathfrak{S}_0^{\pi_0} \Gamma^{!\prime} \setminus \mathfrak{S}_0^{\pi_0} \Gamma^{!\prime\prime} \quad \text{donc} \quad \begin{cases} \varepsilon_3(u_0) = 1 \\ \varepsilon_2(u_0) = -1 \end{cases} \leqno{(75)}
$$
avec $b_0 \equiv u_0 \equiv a_0 \pmod{\Pi_0}$.

\newpage

%% ===== Page 599 =====

et étudions un tel
$$
y = (u,a,b) \in S_0^{\Pi_0} \Gamma'''_! \quad
\begin{cases}
1) & \varepsilon_3(u) \varepsilon_3(u_0) = 1 \\
b) & u \equiv b \equiv \varepsilon_3(\sigma) a \pmod{\Pi_0}
\end{cases} \leqno{(76)}
$$
On a
$$
x y x^{-1} = (u_0(u), a_0(a), b_0(b)) \leqno{(77)}
$$
\marginpar{NB $\varepsilon_3(u_0(u)) = \varepsilon_3(u)$}et la relation $x y x^{-1} \in S_0^{\Pi_0} \Gamma'''_!$ s'écrit
$$
u_0(u) \equiv b_0(b) \equiv \varepsilon_3(\sigma) \cdot a_0(a) \leqno{(78)}
$$
Or les relations 75) b) et 76) b) impliquent aussitôt
$$
u_0(u) \equiv b_0(b) \equiv a_0(\varepsilon_3(\sigma) \cdot a)
$$
$$
\equiv \varepsilon_3[a_0(\sigma)] \cdot a_0(a)
$$
donc la relation (78) équivaut à :
$$
\varepsilon_3[\sigma] \equiv \varepsilon_3[a_0(\sigma)] \leqno{(79)}
$$
ce qui est évidemment toujours vérifié si $\varepsilon_3 = 1$
et pour $\varepsilon_3 = -1$ s'écrit :
$$
\sigma \equiv a_0(\sigma) \pmod{\hat{\Pi}_0} \leqno{(80)}
$$
i.e.\ $\varepsilon_2(a_0) = 1$ i.e.\ $\varepsilon_2(u_0) = 1$,
relation qui est fausse par l'hyp.\ 75 c).
Donc on trouve encore $S_0^{\Pi_0}\Gamma'''_!$ non
invariant dans $S_0^\Pi \Gamma'_!$. Par contre, définissons
$$
\begin{cases}
S_0 \Gamma^\circ = \Ker \big( \mathfrak{S}_0^\circ \xrightarrow{(\varepsilon_2, \varepsilon_3)} \mu \times \mu \big) = S_0 \Gamma'_! \cap S_0 \Gamma''_! \\
S_0^\Pi \Gamma^\circ = S_0 \Gamma^\circ \cap S_0^\Pi \Gamma'_! = S_0^\Pi \Gamma'_! \cap S_0 \Gamma''_! \\
S_0^{\Pi_0} \Gamma^\circ = S_0 \Gamma^\circ \cap S_0^{\Pi_0} \Gamma'_!
\end{cases} \leqno{(81)}
$$
de sorte qu'on a une suite exacte
$$
1 \to \mu \to S_0^{\Pi_0} \Gamma^\circ \to S_0 \Pi^\circ \to 1 \leqno{(82)}
$$

\newpage

%% ===== Page 600 =====

[unclear: canoniquement] scindée par $\mathfrak{S}_0^{\pi_0} \Gamma^\circ$ -- 

[unclear: On trouve donc] :
$$
\mathfrak{S}_0^\pi \Gamma^\circ = \mathfrak{S}_0^{\pi_0} \Gamma^\circ \times \mu \leqno{(82)}
$$

Ceci joint aux calculs précédents montre que $\mathfrak{S}_0^{\pi_0} \Gamma^\circ$, qui est d'indice 8 dans $\mathfrak{S}_0^\pi \Gamma^!$, est un sous-groupe invariant. En résumé :

Proposition Considérons les sous-groupes (d'indice 4, 4, 8) $\mathfrak{S}_0^\pi \Gamma^{\prime !}, \mathfrak{S}_0^{\pi_0} \Gamma^{\prime \prime !}, \mathfrak{S}_0^{\pi_0} \Gamma^{\circ !}$ de $\mathfrak{S}_0^\pi \Gamma^!$ qui sont les images inverses respectives des sous-groupes $Z', Z'', Z^\circ$ de la section involutive $A$ \marginpar{Ce qui est curieux c'est qu'il n'existe pas de section globale de $\mathfrak{S}_0^\pi \Gamma^!$ au dessus de $Z$, qui coïncide avec $A$ sur $Z^{!!!}$ !} définie par (51) (52) dans $\mathfrak{S}_0^\pi \Gamma^!$ sur $Z = \mathcal{M}[0]$ -- [unclear: qui, par extension à] ces sous-groupes, se déduisent des [unclear: dits] sous-groupes [unclear: respectifs] $\mathfrak{S}_0^\pi \Gamma', \mathfrak{S}_0^{\pi_0} \Gamma'', \mathfrak{S}_0^{\pi_0} \Gamma^\circ$ de $\mathfrak{S}_0^\pi \Gamma$ par [unclear: produit direct avec] $\mu$ dans $\mathfrak{S}_0^\pi \Gamma^!$ :
$$
\mathfrak{S}_0^\pi \Gamma^{\prime !} \simeq \mathfrak{S}_0^\pi \Gamma^{\prime} \times \mu, \quad \mathfrak{S}_0^{\pi_0} \Gamma^{\prime \prime !} \simeq \mathfrak{S}_0^{\pi_0} \Gamma^{\prime \prime} \times \mu, \quad \mathfrak{S}_0^{\pi_0} \Gamma^{\circ !} \simeq \mathfrak{S}_0^{\pi_0} \Gamma^{\circ} \times \mu \leqno{(83)}
$$

On a de plus ceci.
\begin{enumerate}
\item[a)] $\mathfrak{S}_0^\pi \Gamma^{\prime !}, \mathfrak{S}_0^{\pi_0} \Gamma^{\prime \prime !}$ ne sont pas distingués dans $\mathfrak{S}_0^\pi \Gamma^!$
\item[b)] $\mathfrak{S}_0^{\pi_0} \Gamma^\circ$ est distingué dans $\mathfrak{S}_0^\pi \Gamma^!$, de sorte que l'ext. $\mathfrak{S}_0^\pi \Gamma^!$ de $Z^! = \mathcal{M}^![0]$ par $\mu$ est [unclear: canoniquement] isom. à l'image inv. d'une ext. cent. $\mathfrak{S}_0^\pi \Gamma^! / (\mathfrak{S}_0^{\pi_0} \Gamma^\circ)$ de $Z^!/Z^\circ \simeq \mu \times \mu$ par $\mu$.
\end{enumerate}

\newpage

%% ===== Page 601 =====

604

Examinons cette extension centrale
$$ (84) \quad 1 \to \mu \to \varepsilon \xrightarrow{(\varepsilon_3, \varepsilon_2)} \mu \times \mu \to 1 $$
$$ \simeq S_0^\pi \Gamma_{\rho,\sigma}^! / S_0^{\pi_0} \Gamma_{\rho,\sigma}^0 $$
s'insérant dans le diagramme d'extensions,

$$ (85) $$
\begin{tikzcd}
1 \ar[r] & \mu \ar[r] \ar[d, equal] & \varepsilon \ar[r] \ar[d, hook] & \mu \times \mu \ar[r] & 1 \\
1 \ar[r] & \mu \ar[r] \ar[d] & S_0^\pi \Gamma_{\rho,\sigma}^! \ar[r] \ar[d] & Z_{\rho,\sigma}^! \simeq M[0] \ar[u, "(\varepsilon_3{,}\varepsilon_2)"'] \ar[d] & \\
1 \ar[r] & L_\rho \times L_\sigma \ar[r] & S_0 \Gamma_{\rho,\sigma} \ar[r] & Z_{\rho,\sigma} \ar[r] & 1
\end{tikzcd}

L'intérêt de (85) est qu'il permet de récupérer l'extension $S_0^\pi \Gamma_{\rho,\sigma}^!$ de $Z_{\rho,\sigma}^!$ par $L_\rho \times L_\sigma$, par la loi usuelle - contravariante dans le foncteur "extensions", via

$$ (85') \quad Z_{\rho,\sigma}^! \xrightarrow{(\varepsilon_3, \varepsilon_2)} \mu \times \mu \quad , \quad \mu \hookrightarrow L_\rho \times L_\sigma $$
$$ -1 \mapsto (\rho^3 \sigma^2) $$
("hom. diagonal")

ainsi que la sous-extension $S_0^\pi \Gamma_{\rho,\sigma}^!$ de celle ci (de $Z_{\rho,\sigma}^!$ par $\mu$).

Etudions d'abord l'ordre des éléments de $\varepsilon$ - ils sont évidemment diviseurs de 4, sont ils tous d'ordre 2 (... 1)? Soit

[MARGIN: NB Pour l'intuition dans ce qui suit, $\underline{u} = (u,a,b) \in S_0^\pi \Gamma^!$ (dans $Z, \hat{N}_\rho^\pi, \hat{N}_\sigma^\pi$) i.e. $u \equiv b \, (\hat{\Pi}_0)$ et $u \equiv \varepsilon_3[\sigma] a \, (\hat{\Pi})$ (87). On a $u^2 \in S_0^{\pi_0} \Gamma_0$. Evidemment $\underline{u} = (u^2, a', b') \in S_0^\pi \Gamma_0$]

\newpage

%% ===== Page 602 =====

il [unclear: reste à] voir si $u^2 \in S_0^\Pi \Gamma!$ i.e. si on a seulement $u^2 \equiv b^2 \pmod{\hat{\Pi}_0}$
--- ce qui est clair et exprime simplement $u^2 \in S_0^\Pi \Gamma$ --- mais même
$$
u^2 \equiv a^2 \pmod{\hat{\Pi}_0} \quad (?) \leqno{(88)}
$$
La relation
$$
u \equiv \varepsilon_3[\sigma].a \pmod{\hat{\Pi}}
$$
s'écrit aussi
$$
u \equiv \pm \varepsilon_3[\sigma].a \pmod{\hat{\Pi}_0}
$$
et implique
$$
u^2 \equiv \varepsilon_3[\sigma] a \varepsilon_3[\sigma] a \pmod{\hat{\Pi}_0}
$$
ce qu'on peut écrire
$$
u^2 \equiv \varepsilon_3 \sigma(a) . a = \begin{cases} a^2 & \text{si } \varepsilon_3 = 1 \\ -\sigma(a) . a & \text{si } \varepsilon_3 = -1 \end{cases} \pmod{\hat{\Pi}_0} \leqno{(89)}
$$
Donc (88) est vraie si $\varepsilon_3 = 1$, tandis que pour $\varepsilon_3 = -1$, elle équivaut à :
$$
\sigma(a) \equiv -a \pmod{\hat{\Pi}_0}
$$
i.e.
$$
[\sigma, a] = -1 \quad \text{dans } \mathcal{E} / \hat{\Pi}_0
$$
i.e.
$$
\underset{\varepsilon_2(u)}{\varepsilon_2(a)} = -1 \quad \text{i.e.} \quad \varepsilon_2(a)\varepsilon_3(a) = 1 .
$$

Donc on trouve :

Proposition. Soit $x \in \mathcal{E}$ (cf (84)). Pour que l'on ait $x^2 = 1$, il faut et il suffit qu'on ait $\varepsilon_3(x) = 1$ ou $\varepsilon_2(x)\varepsilon_3(x) = 1$, i.e. que
$$
x \in \mathcal{E}^+ \cup [unclear: \mathcal{E}'] \leqno{(90)}
$$

\newpage

%% ===== Page 603 =====

où
$$ \mathcal{E}' = \{x \in \mathcal{E} \mid \varepsilon_3(x) = 1\}, \quad \mathcal{E}''' = \{x \in \mathcal{E} \mid \varepsilon_2(x)\varepsilon_3(x) = 1\} \leqno{(91)} $$

\underline{Corollaire 1.} $\mathcal{E}$ contient deux éléments d'ordre 4, qui sont les deux éléments de $\mathcal{E}$ au dessus de l'élément $(-1, 1)$ de $\mu \times \mu$. Ces deux éléments sont inverses l'un de l'autre, et sont les deux générateurs du seul sous-groupe cyclique d'ordre 4 de $\mathcal{E}$, qui n'est autre que
$$ \mathcal{E}'' = \{x \in \mathcal{E} \mid \varepsilon_2(x) = 1\} \leqno{(92)} $$

\underline{Corollaire 2.} L'extension $\mathcal{E}$ n'est pas scindée.
En étant centrale, elle serait isomorphe à $(\mu \times \mu) \times \mu$, tous ses éléments seraient d'ordre qui divise 2.

Nous pouvons regarder maintenant $\mathcal{E}$ comme extension de $\mathcal{E}/\mathcal{E}'' \simeq [unclear: F \beta] \simeq \mu$ par $\mathcal{E}'' (\simeq \mathbf{Z}/4\mathbf{Z})$. Mais il faudrait d'abord expliciter l'opération de $\mathcal{E}/\mathcal{E}'' \simeq \mu$ sur $\mathcal{E}''$ --- l'élément $-1$ d'ordre 2 de $\mu$ induit un automorphisme involutif de $\mathcal{E}''$ (d'ailleurs $\operatorname{Aut}(\mathcal{E}'') \simeq (\mathbf{Z}/4\mathbf{Z})^* \simeq \{\pm 1\}$, donc l'automorphisme en question se réduit à la multiplication par un signe $\pm 1$ qu'il faut déterminer). On va prendre le représentant de $-1 \in \mathcal{E}/\mathcal{E}''$ dans $\mathcal{E}$ provenant de
$$ x = (u_0 = \tau, a_0 = \tau', b_0 = \tau) \in \mathfrak{S}_0^{\Gamma!} \quad \text{(cf (70))} $$
pour en faire opérer sur l'image d'un
$$ u = (u, a, b) \in \mathfrak{S}_0^{\Gamma!} \quad \text{avec} \quad \varepsilon_3(u) = -1, \varepsilon_2(u) = 1 \leqno{(93)} $$
(qui reste un générateur de $\mathcal{E}''$). ($u \equiv [unclear: b] \pmod{\hat{\Pi}_0}$, $u \equiv \sigma a \pmod{\hat{\Pi}}$)

\newpage

%% ===== Page 604 =====

On aura donc
$$ \underline{x}(u) = (\tau(u), \tau'(a), \tau(b)) \equiv (u, \varepsilon a, \varepsilon b) $$
mod $\mathfrak{S}_0^{\pi_0} \Gamma^{00}$, avec un signe $\varepsilon \in \{\pm 1\}$ à déterminer.

La condition de congruence s'écrit
$$ (u \tau(u)^{-1}, \varepsilon a \tau'(a)^{-1}, \varepsilon b \tau(b)^{-1}) \in \mathfrak{S}_0^{\pi_0} \Gamma^{00} $$
qui - pour un [unclear: élément] qui soit dans $\mathfrak{S}_0^{\pi_0} \Gamma^0$ est évidente [unclear: de toute façon] quel que soit le choix du signe $\varepsilon$, c'est la condition de congruence qui distingue $\mathfrak{S}_0^{\pi_0} \Gamma^!$ de $\mathfrak{S}_0^{\pi} \Gamma^!$ qui est essentielle, savoir
$$ u \tau(u)^{-1} \equiv \varepsilon a \tau'(a)^{-1} \pmod{\hat{\pi}_0} \leqno{(94)} $$

La difficulté vient de ce que les seuls éléments de $\mathcal{M}$ qu'on sache expliciter sont ceux qui proviennent de $\mathcal{G} \simeq \operatorname{Gl}(2, \hat{\mathbf{Z}})$ lui-même, or pour ceux-ci on a $\varepsilon_2(u) \varepsilon_3(u) = 1$, alors que le $u$ ci-dessus doit vérifier $\varepsilon_3(u) = 1$, $\varepsilon_2(u) = -1$ donc $\varepsilon_2(u)\varepsilon_3(u) = -1$. On ne peut donc pas travailler dans le $\operatorname{Gl}(2, \hat{\mathbf{Z}})$ de $\mathcal{M}$, et y prendre
$$ u = \begin{pmatrix} \mu & 0 \\ 0 & 1 \end{pmatrix} \quad \text{avec } \mu \in \hat{\mathbf{Z}}^*, \begin{cases} \mu \equiv -1 \ (3) \\ \mu \equiv 1 \ (4) \end{cases} \leqno{(95)} $$
tandis que $\tau$ correspond à $c$ i.e. $\mu \equiv 5 \ (12)$
$$ \tau = \begin{pmatrix} -1 & 0 \\ 0 & 1 \end{pmatrix} $$

\newpage

%% ===== Page 605 =====

On aura de toutes façons
$$
u \tau u^{-1} = 1
$$
soit d'autre part dans $\operatorname{Gl}(2, \hat{\mathbf{Z}})$
$$
\begin{cases}
a = \tau' a_0 \qquad a_0 = c \rho + d \qquad c^2+cd+d^2 \in \hat{\mathbf{Z}}^* \\
\tau' = \left( \begin{matrix} 0 & 1 \\ 1 & 0 \end{matrix} \right) , \rho = \left( \begin{matrix} 0 & -1 \\ 1 & 1 \end{matrix} \right)
\end{cases} \leqno{(95)}
$$
d'où
$$
a = \left( \begin{matrix} 0 & 1 \\ 1 & 0 \end{matrix} \right) \left( \begin{matrix} d & -c \\ c & c+d \end{matrix} \right) = \left( \begin{matrix} c & c+d \\ d & -c \end{matrix} \right) \leqno{(97)}
$$
et \marginpar{[unclear: un peu bof... vu plus bas...]}
$$
\mu = - (c^2+cd+d^2) \leqno{(98)}
$$
On va travailler dans le groupe fuchsien $\hat{\Pi}_0$ dans $\operatorname{Sl}(2, \hat{\mathbf{Z}})$ engendré par
$$
- l_0 = \left( \begin{matrix} 1 & -2 \\ 0 & 1 \end{matrix} \right), - l_1 = - \left( \begin{matrix} 1 & 0 \\ 1 & 1 \end{matrix} \right), - l_\infty = - \left( \begin{matrix} 3 & -2 \\ 2 & -1 \end{matrix} \right)
$$
et dans le groupe quotient de $\operatorname{Gl}(2, \hat{\mathbf{Z}})$ par $\hat{\Pi}_0$, qui s'identifie à
$$
\hat{\mathbf{Z}}^* \times \operatorname{Sl}(2, \hat{\mathbf{Z}})/\hat{\Pi}_0 \simeq \operatorname{Sl}(2, \mathbf{Z})/\Pi_0 \simeq \left\{ (\rho, \sigma) \, \middle| \, \begin{matrix} \rho^6 = \sigma^4 = 1 \\ \rho^3 = \sigma^2 \\ \sigma \rho \sigma = \rho^{-1} \end{matrix} \right\} ,
$$
avec le prod. semi-dir \marginpar{[unclear: le direct quotient par dét]}
$$
T_0 = \left\{ \left( \begin{matrix} \zeta & 0 \\ 0 & 1 \end{matrix} \right) \mid \zeta \in \hat{\mathbf{Z}}^* \right\} \isom \hat{\mathbf{Z}}^*
$$
opérant par l'intermédiaire de $\varepsilon_2(\zeta)$ \marginpar{dét}
cf p. 593', 599 - formules (34) (35) (35 bis) ). On a
$$
\sigma a \equiv u \pmod{\hat{\Pi}_0} \qquad \text{où } \sigma = \left( \begin{matrix} 0 & -1 \\ 1 & 0 \end{matrix} \right) \leqno{(99)}
$$
d'où
$$
\sigma a = \left( \begin{matrix} -d & -c \\ -c & c+d \end{matrix} \right) \leqno{(100)}
$$

\newpage

%% ===== Page 606 =====

La relation (99), il suffit d'avoir $\det a = \det u$ i.e. (98), s'écrit
$$
\sigma a \equiv u \quad (2) \leqno{(102)}
$$
soit en d'autres termes la congruence s'écrit
$$
c \equiv 0 \quad (2), \quad d \equiv 1 \quad (2) \qquad (-d \equiv \mu \quad (2) \text{ n'a aucune conséquence car } \mu \equiv 1 \equiv -1 \quad (2)) \leqno{(103)}
$$

Revenons à la relation (97), et notons ceci :

\underline{Lemme.} Soient $c, d \in \hat{\mathbf{Z}}$ satisfaisant
$$
\mu = -(c^2 + cd + d^2) \in \hat{\mathbf{Z}}^* \footnote{NB On a du reste $\mu \equiv -1 \quad (3)$.} \leqno{(104)}
$$
$$
\mu \equiv 1 \quad (4) \quad \text{(i.e. } c^2 + cd + d^2 \equiv -1 \quad (4)\text{)} \leqno{(105)}
$$
$$
c \equiv 0 \quad (2), \quad d \equiv 1 \quad (2) \leqno{(106)}
$$
Soit
$$
a = \tau'(c \rho + d) = \begin{pmatrix} -c & c+d \\ d & c \end{pmatrix} \in \operatorname{Gl}(2, \hat{\mathbf{Z}}) \leqno{(107)}
$$
(noter son déterminant $\mu$). Alors on a une congruence
$$
a \tau'(a)^{-1} \equiv \varepsilon \pmod{\hat{\Pi}_0} \quad (\subset \operatorname{Gl}(2, \hat{\mathbf{Z}})) \leqno{(108)}
$$
où le signe $\varepsilon \in \{\pm 1\}$ ne dépend pas du choix fait pour $c,d$.

\newpage

%% ===== Page 607 =====

c'est le moment de les calculer ! Reprenons
$$ a_0 = c \rho + d \in \mathbf{Z}_{\rho} $$
on a
$$ a = \tau' a_0 , \quad \tau'(a) = a_0 \tau', \quad a \tau'(a)^{-1} = \tau' a_0 \tau' a_0^{-1} = \tau'(a_0) a_0^{-1} $$
où\footnote{On pose dans $M_2(\hat{\mathbf{Z}})$, ${}^t u = (\text{Tr} u) - u$ (i.e.\ $u^{-1}$ si $\det u = 1$) et on note ${}^t \rho = \rho^{-1}$ (plus généralement ${}^t u = u^{-1}$ si $\det u = 1$).}
$$ \begin{cases} \tau'(a_0) = c \rho^{-1} + d = {}^t a_0 \\ a_0^{-1} = {}^t a_0 / \det a_0 = - \frac{1}{\mu}(c \rho^{-1} + d) \end{cases} $$
donc
$$ a \tau'(a)^{-1} = \tau'(a_0) a_0^{-1} = - \frac{1}{\mu}(c \rho^{-1} + d)^2 = \frac{1}{c^2 + c d + d^2}(c \rho^{-1} + d)^2 $$
Or,
$$ (c \rho^{-1} + d)^2 = c^2 \underbrace{\rho^{-2}}_{-\rho} + 2 c d \underbrace{\rho^{-1}}_{1-\rho} + d^2 = - (c^2 + 2 c d) \rho + (2 c d + d^2) $$
Or, comme
$$ \begin{cases} c \equiv 0 \pmod 2 \\ d \equiv 1 \pmod 2 \end{cases} \implies \begin{cases} c^2, 2 c d \equiv 0 \pmod 4 \\ d^2 \equiv 1 \pmod 4 \end{cases} $$
donc $(c \rho^{-1} + d)^2 \equiv 1 \pmod 4$.

donc le signe $\varepsilon$ est égal : $+1$

et comme par hypothèse
$$ c^2 + c d + d^2 \equiv -1 \pmod 4 $$
on aura
$$ a \tau'(a)^{-1} = \frac{1}{c^2 + c d + d^2} (c \rho^{-1} + d)^2 \equiv -1 \pmod 4 \leqno (109) $$
Donc on a prouvé ce fait

Proposition. Soient $c, d \in \hat{\mathbf{Z}}$ tels que
$c^2 + c d + d^2 \in \hat{\mathbf{Z}}^*$, soit
$$ a = \tau'(c \rho + d) = \begin{pmatrix} -c & c + d \\ d & c \end{pmatrix} \in \operatorname{Gl}(2, \hat{\mathbf{Z}}), $$
\marginpar{NB $\sigma = \begin{pmatrix} 0 & -1 \\ 1 & 0 \end{pmatrix}$}
avec
$$ c \equiv 0 \pmod 2, \quad d \equiv 1 \pmod 2 . $$

\newpage

%% ===== Page 608 =====

Alors on a (pour $\tau' = \begin{pmatrix} -1 & 0 \\ 0 & 1 \end{pmatrix}$)
$$
a \tau'(a)^{-1} \equiv \varepsilon_2 \left( \frac{c}{2} \right) \pmod{4}
$$
(= fonction [unclear: de] la congruence mod $\hat{\pi}_0$).

\underline{Corollaire} L'extension $\mathcal{E}$ de la page préc. (84), où on a $\mu$ pour $\mathcal{E}'' \simeq \mathbf{Z}/4\mathbf{Z}$, est non commutative.

Mais voici une façon de le voir sans calcul. Si l'extension était commutative, à un isomorphisme près elle serait [unclear: donnée] par un $\operatorname{Ext}^1_{\mathbf{Z}/4\mathbf{Z}}(\dots)$. Or le foncteur $\operatorname{Ext}^1(-, \mu)$ est additif. Comme la restriction de l'extension aux sous-groupes $\mathcal{E}'/\mu$ et $\mathcal{E}''/\mu$ de $\mu \times \mu$ est triviale, et que $\mu \times \mu$ est la somme directe de ses sous-groupes (qui sont resp. [unclear: l'un] et le sous-groupe diagonal ?), il s'ensuivrait que l'extension serait triviale, ce qui n'est pas le cas.

La structure du groupe $\mathcal{E}$ s'explicite aisément, par ex. en tant qu'extension de $\mathcal{E}/\mathcal{E}'' \simeq \mu$ par $\mathcal{E}'' \simeq \mathbf{Z}/4\mathbf{Z}$. Soit [unclear: $x$] un él. quelconque de $\mathcal{E} - \mathcal{E}''$ ; il est d'ordre 2 (les seuls éléments d'ordre 4 étant dans $\mathcal{E}''$) de sorte que l'

\newpage

%% ===== Page 609 =====

$$
\mathcal{E} \simeq \{1, x\} \underset{(\text{1/2 direct})}{\cdot} \mathcal{E}'' \leqno (110)
$$

Notons que si $x$ opérait trivialement sur $\mathcal{E}''$, c.à.d. si $\mathcal{E}$ était isomorphe au produit, alors pour un élément $u$ de $\mathcal{E}''$ d'ordre $4$, $x.u$ serait d'ordre $4$, ce qui est absurde (troisième démonstration du fait que $\mu$ opère non trivialement sur $\mathcal{E}''$). Or du seul fait que l'op. de $\{1, x\}$ sur $\mathcal{E}''$ est non triviale, elle est connue, la structure de $\mathcal{E}$ en est déterminée en fait -
$$
\mathcal{E} \simeq D_4 \leqno (111)
$$

Il y a un choix privilégié de $x$, en prenant pour $x$ l'élt $\tau_\xi$
$$
\begin{cases}
\tau_\xi \equiv \tau_{S_{0P}} \pmod{S_0^{P_0}} \quad \text{(où bien sûr} \\
\tau_{S_{0P}} = (\tau, \tau', \tau) )
\end{cases} \leqno (112)
$$
de sorte qu'on a là une représentation canonique
$$
\mathcal{E} \simeq \{1, \tau_\xi\} \underset{\text{produit 1/2 direct}}{\cdot} \mathcal{E}'' \qquad \begin{matrix} \mathcal{E}'' \simeq \mathbf{Z}/4\mathbf{Z} \\ \tau_\xi(u) = u^{-1} \\ \text{pour } u \in \mathcal{E}'' \end{matrix} \leqno (113)
$$

L'hom. $\mathcal{E} \to \mu \times \mu$ se récupère par ses restrictions aux $2$ facteurs :
$$
\begin{aligned}
\mu \simeq \{1, \tau_\xi\} &\longrightarrow \mu \times \mu \qquad \tau_\xi \mapsto (-1, -1) \quad \text{(hom. diagonal)} \\
\mathcal{E}'' &\longrightarrow \mu \times \{1\} \subset \mu \times \mu
\end{aligned} \leqno (114)
$$
\begin{center}
\small (l'unique hom. non trivial $\mathcal{E}'' \to \mu$ déduit de l'isom. $\mathcal{E}''_2 = \mathcal{E}''/\mathcal{E}''^2 \simeq \mu$)
\end{center}

\newpage

%% ===== Page 610 =====

Le fait que $\mathcal{E}$ ne soit pas commutatif s'exprime aussi par le fait que l'homomorphisme
$$
\mathcal{E}_{ab} \xrightarrow{\sim} \mu \times \mu \leqno{(115)}
$$
est un isomorphisme, i.e.\ que le sous-groupe [unclear: noyau] de $\mathcal{E} \to \mu \times \mu$ est aussi le sous-groupe des commutateurs.

On a vu que l'extension $\mathcal{E}$ de $\mu \times \mu$ par $\mu$ n'est pas scindée, mais qu'en est-il de son image réciproque, l'extension $\hat{\mathfrak{S}}_0^{\pi !}$ de $\mathbf{Z}^! \simeq \widehat{\mathcal{M}}(0)$ par $\mu$ ? A priori, on serait tenté de croire qu'il n'y a pas de scindage de cette extension qui prolonge le scindage canonique $A$ dont on dispose au-dessus de $\mathbf{Z}^\circ = \operatorname{Ker}(\mathbf{Z}^! \xrightarrow{(\varepsilon_3, \varepsilon_4)} \mu \times \mu)$. Je présume que $\hat{\mathfrak{S}}_0^{\pi !}$, extension de $\mathbf{Z}^{!\wedge}$ par $\mu$, est non triviale, et que l'extension qu'elle induit sur son sous-groupe
$$
\Gamma_{\mathbf{Q}}' = \Gamma_{\mathbf{Q}} \times_{\widehat{\mathcal{M}}(0)'} \widehat{\mathcal{M}}(0) \subset \widehat{\mathcal{M}}(0) \subset \widehat{\mathcal{M}}'[0] = \mathbf{Z}^!
$$
Il serait intéressant d'identifier plus précisément l'extension sur $\Gamma_{\mathbf{Q}}'$ qu'elle induit, qui correspond à un élément canonique
$$
c \in H^2(\operatorname{Spec} \mathbf{Q}, \mu_2) \subset \operatorname{Br}(\mathbf{Q}) \leqno{(115')}
$$

\newpage

%% ===== Page 611 =====

sauf erreur on aura
$$ c = "\xi_3" \cup "\xi_4" \in H^2(\mathbf{Q}, \mu_2) \leqno{(115)} $$
(cup-produits)
sont les éléments qui décrivent respectivement
les extensions $\mathbf{Q}(\sqrt[3]{1}) = \mathbf{Q}(\sqrt{-3})$, et
$\mathbf{Q}(\sqrt{3})$ -- contenues l'une et l'autre,
ainsi que $\mathbf{Q}(\sqrt{-1}) = \mathbf{Q}(i)$, dans l'extension
biquadratique $\mathbf{Q}(i, j) = \mathbf{Q}(\sqrt[12]{1})$, le corps
des racines $12^{\text{ièmes}}$ de l'unité (dont le groupe
de Galois est $(\mathbf{Z}/12\mathbf{Z})^\times \simeq (\mathbf{Z}/3\mathbf{Z})^\times \times (\mathbf{Z}/4\mathbf{Z})^\times \simeq \mu \times \mu$).
Il doit être trivial -- quand on connaît un
peu les fondements -- que $c$ ainsi
explicité est $\neq 0$, et n'est pas non plus
scindé en passant à $\mathbf{Q}(i)$, la troisième
larronne.

Considérons maintenant le
sous-groupe central
\marginpar{NB $Z' \simeq \mathcal{M}_{1,0}^\sim$ \\ $\Gamma \simeq \Gamma_{\mathbf{Q}}$}
$$ \mu \xrightarrow{i_{\text{sin}}} \mathfrak{S}_0 \Gamma_{3,4}^! = Z' \times_\Gamma N_3^\Delta \times N_4^\Delta \leqno{(116)} $$
$$ \varepsilon \longmapsto (1, \varepsilon, 1) $$
à valeurs dans
$$ \operatorname{Ker}([unclear: \mathfrak{S}_0 \Gamma_{3,4}^!] \to Z') \simeq \mu \times \mu , $$
donne un hom. canonique de relèvement
$$ \mathcal{M}'_{1,0} \simeq Z' \simeq \mathfrak{S}''_0 \Gamma_{3,4}^! / \mu \longrightarrow \mathfrak{S}_0 \Gamma_{3,4}^! (\mu) , \leqno{(117)} $$

\newpage

%% ===== Page 612 =====

qui\footnote{Ceci méritait d'être explicité dans un théorème récapitulatif} s'explicite en termes de deux hom. canoniques fondamentaux
$$
\begin{matrix}
Z'_{ \hat{p},\sigma } \\
\wr \parallel \\
\mathcal{M}_{\Gamma_{0,3}} \\
\wr \parallel \\
\Gamma^\sim_{\mathbf{Q}}
\end{matrix}
\xrightarrow{\quad a \quad} N^*_{\hat{p}}/\mu \quad , \quad Z'_{ \hat{p},\sigma } \xrightarrow{\quad b \quad} N^*_\sigma \leqno (118)
$$
où, je rappelle
$$
\begin{cases}
N^*_{\hat{p}} = \mathcal{M} = \mathcal{N}_{1,1}^\sim \simeq \Gamma^\sim_{\mathbf{Q}} \cdot \underset{\subset \operatorname{Sl}(2,\mathbf{Z})^\wedge}{\mathfrak{S}^+} \\
\quad \| \\
\{ u \in \mathcal{M} \mid u(\hat{p}) = \hat{p}^{\varepsilon_3(u)} \} \\
\\
N^*_\sigma \subset \mathcal{M} \\
\quad \| \\
\{ u \in \mathcal{M} \mid u(\sigma) = \sigma^{\varepsilon_2(u)} \quad (= \varepsilon_2(u)\sigma) \}
\end{cases} \leqno (119)
$$
Les applications $a, b$ sont données par les formules (cf 51)
$$
\begin{matrix} \underline{a}(u) = a = \varepsilon_3[\tau'] \alpha^{-1} u \\ \alpha = u a^{\varepsilon_3(\tau')} u^{-1} \end{matrix} \quad , \quad \begin{matrix} \underline{b}(u) = b = \beta^{-1} u \\ \beta = u b u^{-1} \end{matrix} \quad \text{i.e.} \leqno (120)
$$
où $\alpha \in \hat{\Pi}_{0\tau'}$, $\beta \in \hat{\Pi}_0^\wedge$ sont caractérisés par
$$
\begin{matrix} u(\hat{p}) = \operatorname{int}(\alpha)(\hat{p}) \\ \in N^*_{\hat{p}}/\mu \end{matrix} \quad , \quad \begin{matrix} u(\sigma) = \varepsilon_2(\sigma) \operatorname{int}(\beta)(\sigma) \\ \in N^*_\sigma \end{matrix} \leqno (121)
$$
Les quantités $a=a(u)$, $b=b(u)$ se caractérisent en termes de $u \in Z'$ (ou de $u \in Z!$ ca-ractérisé par ces mêmes relations) par les équations
$$
u \equiv b \pmod{\hat{\Pi}_0} \quad , \quad u \equiv \varepsilon_3[\tau'] a \pmod{\hat{\Pi}_0} \leqno (122)
$$
qui impliquent a fortiori (vu que $u \equiv 1 \pmod{\hat{\Pi}_0}$) que $a \equiv b \pmod{\hat{\Pi}_0}$, d'où

\newpage

%% ===== Page 613 =====

(123)
\begin{tikzcd}
& N^*_\rho \ar[dr, "\delta_\rho"] \\
Z^!_{\rho,\sigma} \ar[ur, dashed] \ar[r, dashed, "\delta_Z"] \ar[dr, dashed] & \dots & \gamma_{\rho,\sigma} \simeq \tilde{\Gamma}_Q \\
& N^*_\sigma \ar[ur, "\delta_\sigma"']
\end{tikzcd}

Considérons maintenant les deux
[unclear: crossed out text]
Il serait temps
de rectifier la
confusion (p. 594') entre $Z = Z_{\rho,\sigma} \subset \tilde{M}_{0,3}^\sim$
et $Z^! = Z^!_{\rho,\sigma} = Z \cap \tilde{M}_{0,3}^!$. On a en fait
$$ Z_{\rho,\sigma} \simeq \tilde{M}(0) \cdot \tilde{L}_0^\otimes \subset \mathcal{M} $$
(124) \hfill (où $\tilde{L}_0^\otimes$ est un $\frac{1}{2}$-tour)
$$ \searrow \tilde{M}_{0,3}^\sim(0) \cdot \tilde{L}_0^{\otimes_{0,3}} \subset \tilde{M}_{0,3}^\sim $$

(125) $$ Z^!_{\rho,\sigma} \simeq Z_{\rho,\sigma} \cap M^! = M(0) \cdot L_0^\pi \subset M $$
\hfill (où $L_0^\pi$ est un $\pi$-tour)
$$ \searrow M_{0,3}(0) \cdot L_0^{\pi_{0,3}} \subset M_{0,3} $$

Le lien entre les deux est donné
par l' e.g. [unclear]

(126)
\begin{tikzcd}
1 \ar[r] & L_0^\pi \ar[r] \ar[d, hook] & Z^!_{\rho,\sigma} \ar[r] \ar[d, hook, "p"] & \gamma_{\rho,\sigma} \ar[r] \ar[d, equal] & 1 \\
1 \ar[r] & \tilde{L}_0^\otimes \ar[r] & Z_{\rho,\sigma} \ar[r] & \gamma_{\rho,\sigma} \ar[r] \ar[d, equal] & 1 \\
& & & \simeq M(0) \simeq \tilde{\Gamma}_Q
\end{tikzcd}

où $L_0^\pi \subset \tilde{L}_0^\otimes$
L'inclusion d'indice 2 de
(127) $$ L_0^\pi = l_0^{\hat{\mathbb{Z}}} \subset l_0^{\frac{1}{2}\hat{\mathbb{Z}}} = \tilde{L}_0^\otimes $$

Bien entendu, c'est bien sur $Z^!_{\rho,\sigma}$, non sur
$Z_{\rho,\sigma}$, qu'est défini (117), (118) d'où (123).

\newpage

%% ===== Page 614 =====

Je m'intéresse à la restriction des
hom. (123) au sous-groupe $L_0^\pi$ de $Z_{\rho,\sigma}^!$
la concrétisation (122) montre que
l'hom. induit
$$L_0^\pi \longrightarrow N_\rho^*/\{\pm 1\}$$
est trivial (car son terme en $m$ ne
dépend que de $m \bmod \pi^\rho$), tandis que l'hom.
induit
$$L_0^\pi \longrightarrow N_\sigma^*$$
se factorise par $L_0^\pi/L_0^{\pi_0} \simeq \mathbb{Z}/2\mathbb{Z}$ (car
le terme en $m$ ne dépend que de
$m \bmod \pi_0^\sigma$). Sa valeur $b$ pour $l_0$ est
déterminée par les conditions
$$b \in N_\sigma^*, \quad b \equiv l_0 \pmod{\pi_0}$$
qui sont satisfaites évidemment par $b=-1$. Donc
il y a lieu d'introduire aussi, quitte
[unclear: à modifier un peu le diagr.] central
canonique
$$(128) \quad \mu \hookrightarrow N_\sigma^*$$
$$-1 \longmapsto -1$$
et une commutativité dans

(129)
\begin{tikzcd}
L_0^\pi \ar[r, hook] \ar[d] & Z_{\rho,\sigma}^! \ar[d] \\
L_0^\pi/L_0^{\pi_0} \simeq \mu \ar[r, hook] & N_\sigma^*
\end{tikzcd}

[unclear: Donnons un tronçon], à partir des données (123)

\newpage

%% ===== Page 615 =====

611

de $Z_{\rho,\sigma}^!$, par passage au quotient deux hom.
canoniques
$$ (130) \qquad \begin{cases} \gamma_{\rho,\sigma} \longrightarrow N_\rho^* / \{\pm 1\} \\ \Vert \ \tilde{\Pi}_Q \\ \gamma_{\rho,\sigma} \longrightarrow N_\sigma^* / \{\pm 1\} \end{cases} $$
qui sont en fait des sections des
$$ \begin{matrix} N_\rho^* / \{\pm 1\} \longrightarrow \gamma_{\rho,\sigma} \\ \Vert \ \tilde{\Gamma}_Q \simeq M(0) \\ N_\sigma^* / \{\pm 1\} \longrightarrow \gamma_{\rho,\sigma} \end{matrix} $$
de sorte qu'on trouve des décompositions
$$ (131) \qquad \begin{cases} N_\rho^* / \pm 1 \simeq \gamma_{\rho,\sigma} \ltimes S Z_\rho / \{\pm 1\} \\ \qquad \qquad \quad = L_\rho / \{\pm 1\} \simeq \mathbb{Z}/3\mathbb{Z} \\ N_\sigma^* / \pm 1 \simeq \gamma_{\rho,\sigma} \cdot S Z_\sigma / \pm 1 \\ \qquad \qquad \quad = L_\sigma / \pm 1 \simeq \mathbb{Z}/2\mathbb{Z} \end{cases} $$
(produits semi-directs). L'opération de
$\gamma_{\rho,\sigma}$ sur $L_\rho / \{\pm 1\} \simeq \mathbb{Z}/3\mathbb{Z}$ est donnée par le
caractère quadratique $\varepsilon_\rho$, celle sur $L_\sigma / \{\pm 1\} \simeq \mathbb{Z}/2\mathbb{Z}$
est donnée par le caractère quadratique $\varepsilon_\sigma$.

La symétrie de présentation des deux
lignes (130) est un peu bidon - p. ex.
l'hom. $\tilde{\Pi}_Q = \gamma_{\rho,\sigma} \longrightarrow N_\sigma^* / \{\pm 1\}$ se remonte
de façon évidente en
(132)
\begin{tikzcd}
\tilde{\Pi}_Q = \gamma_{\rho,\sigma} \ar[r] \ar[dr, "\sim"'] & N_\sigma^* \\
& M(0) \subset M[0] = Z^! \ar[u, hook]
\end{tikzcd}

\newpage

%% ===== Page 616 =====

Une autre façon de le voir est de noter que
$$
\gamma_{\rho,\sigma} \isom \mathcal{M}_0^![0] / L_0^{\pi_0} \leqno{(132)}
$$
$$
\mathcal{M}_0^![0] \defeq \mathcal{M}(0) . L_0^{\pi_0} \underset{\text{sous-groupe d'indice } 2}{\subset} \mathcal{M}^![0] = \mathcal{M}(0) . L_0^\pi
$$
or sur $L_0^{\pi_0}$, l'homomorphisme induit par
$$
\mathcal{M}^![0] = Z_{\rho,\sigma}^! \to N_\sigma^{\prime *}
$$
est trivial. Ainsi on a un isomorphisme scindé de
$$
1 \to L_\sigma \to N_\sigma^{\prime *} \to \underset{\isom \Pi_Q^\sim}{\gamma_{\rho,\sigma}} \to 1 \leqno{(133)}
$$
en un produit $\frac{1}{2}$-direct
$$
N_\sigma^{\prime *} \isom \gamma_{\rho,\sigma} . L_\sigma
$$
qu'il vaut mieux écrire en [unclear: considérant] la [unclear: sous-suite exacte] relative au sous-groupe $\gamma_{\rho,\sigma} \isom \mathcal{M}$\footnote{NB On définit $\mathcal{M}_{\mathrm{can}}^!$ l'image inverse de $\mathcal{M}_{0,3}^!$ dans le noyau de la projection canonique $\mathcal{M} \to \hat{\mathfrak{S}}^+/\hat{\pi} \isom \operatorname{Aut}(\hat{\mathfrak{S}}^+/\hat{\pi})$}, soit donc
$$
\begin{aligned}
N_0^{\prime *} &\defeq \operatorname{Im}(\mathcal{M}_0^![0] \to N_\sigma^{\prime *}) \\
&= \operatorname{Im}(\mathcal{M}[0] \hookrightarrow N_\sigma^{\prime *}) \\
&= \{ \beta^n u \mid u \in \mathcal{M}(0), \beta \in \Pi_0^\sigma \text{ tels que } u(\sigma) = \varepsilon(u) \beta^2(\sigma) \} \\
&= \{ v \in N_\sigma^{\prime *} \mid \exists \beta \in \Pi_0^\sigma \text{ tel que } \beta v \in \mathcal{M}[0] = \operatorname{Norm}_{\mathcal{M}}(L_0^\pi) \}
\end{aligned} \leqno{(134)}
$$
Posant
$$
\mathcal{M}_0^! \defeq \mathcal{M}(0) . \Pi_0 \underset{\text{indice } 2}{\subset} \mathcal{M}^! = \mathcal{M}(0) . \pi \leqno{(135)}
$$

\newpage

%% ===== Page 617 =====

de sorte qu'on a
$$
\mathcal{M}^!\footnote{NB on a $\mathcal{M}_0^! \simeq \mathcal{M}_{0,3}$} \simeq \mathcal{M}_0^! \times \mu , \leqno{(136)}
$$
on a donc
$$
\begin{cases}
\mathcal{N}_0^{*!} = \mathcal{N}_0^* \cap \mathcal{M}_0^! \\
\mathcal{N}^{*!} \defeq \mathcal{N}^* \cap \mathcal{M}^! = \mathcal{N}_0^{*!} \times \mu ,
\end{cases} \leqno{(137)}
$$
et on trouve que
$$
\mathcal{N}_0^{*!} \isommap \Gamma_{\mathfrak{P},\sigma} = \widetilde{\Gamma}_{\mathbf{Q}} \leqno{(138)}
$$
Posons
$$
\mathcal{M}(1/2) = \mathcal{N}_0^{*!} \quad \text{cf } (137) \leqno{(139)}
$$
on aura donc que $\mathcal{M}(1/2)$ est une
section pour $\mathcal{N}_0^*$ extension de $\Gamma_{\mathfrak{P},\sigma}$ par
$L_0$, et a fortiori pour $\mathcal{M}$ ext. de $\Gamma_{\mathfrak{P},\sigma}$
par $\widehat{\mathfrak{S}}^{+\wedge}$ etc ... de sorte qu'on a\marginpar{attention, $\mathcal{M}'_0$ est inv. dans $\mathcal{M}$ et $\mathcal{M}_0$ inv. dans $\mathcal{M}'$, mais non dans $\mathcal{M}$ cf plus bas}
$$
\begin{cases}
\mathcal{M}_0^! \simeq \mathcal{M}(1/2) \cdot \widehat{\Pi}_0 \\
\mathcal{M}^! \simeq \mathcal{M}(1/2) \cdot \widehat{\Pi} \\
\mathcal{M} \simeq \mathcal{M}(1/2) \cdot \widehat{\mathfrak{S}}^{+\wedge} \\
\mathcal{N}_0^* \simeq \mathcal{M}(1/2) \cdot L_0
\end{cases} \leqno{(140)}
$$
[Si on veut paraphraser en termes d'invariants
$$
\mathcal{N}_0^{*!} \defeq \mathcal{N}_0^* \cap \mathcal{M}_0^! = \left\{ u \in \mathcal{N}_0^* \mid \exists \alpha \in \widehat{\Pi}_0 \text{ t.q. } \underbrace{\alpha u}_{u} \in \mathcal{M}(0) \right\} \leqno{(141)}
$$
[unclear: c'est dire l'une des propriétés qu'on peut énoncer :]
on a $u \in \mathcal{M}(0) \subset \mathcal{M}_0^![0] = Z^1_{\Gamma_{\mathfrak{P},\sigma}}$, à l'aide de
$\alpha \in \widehat{\Pi}_0$ tel que $\alpha^{-1} u \in \mathcal{N}_0^*$ i.e.
$u(p) = \operatorname{int}(\alpha)(p^{\varepsilon_3})$ , $\varepsilon_3 = \varepsilon_3(u) = \varepsilon_3(u)$.
Mais comme $u(p)$ doit opérer trivialement dans $\mathfrak{S}_3 = \widehat{\mathfrak{S}}^{+\wedge}/\widehat{\Pi}$
pour $p$, il faut qu'on ait $\varepsilon_3=1$, et alors
$\alpha$ est l'él. unité $\alpha=1$, d'où $u \in \mathcal{N}_0^*$. Ainsi
on a l'image dans de $u$ dans $\mathcal{N}_0^*$.]

\newpage

%% ===== Page 618 =====

\underline{Proposition} L'homomorphisme $\mathcal{N}_{\hat{P}}^*! \to \Gamma_{\hat{P},\sigma} = \Gamma_{\mathbf{Q}}^\sim$ est injectif, et a pour image $\Gamma'_{\hat{P},\sigma} = \{ \gamma \in \Gamma_{\hat{P},\sigma} \mid \varepsilon_3(\gamma) = 1 \}$.

\underline{Corollaire} $\mathcal{N}_{\hat{P}}^*! = Z_{\hat{P}}^*!$ i.e.\ $\mathcal{N}_{\hat{P}}^*! = Z_{\hat{P}}^*! \defeq Z_{\hat{P}}^* \cap \mathcal{M}!$.

Il est naturel de poser
$$
\mathcal{N}_{\hat{P}}^*! = Z_{\hat{P}0}^* = \mathcal{M}(-j) . \leqno{(142)}
$$

Comme $\mathcal{N}_{\hat{P}}^* \supset Z_{\hat{P}}^* \supset \mu$, on voit que l'on a
$$
\begin{cases}
\mathcal{N}_{\hat{P}}^* = \mathcal{N}_{\hat{P}}^*! \times \mu \quad (\text{donc } \mathcal{N}_{\hat{P}}^*! = Z_{\hat{P}}^*!) \\
Z_{\hat{P}}^* = Z_{\hat{P}}^*! \times \mu
\end{cases} \leqno{(142)}
$$

Le sous-groupe $Z_{\hat{P}0}^* = \mathcal{M}(-j)$ est une section de $\mathcal{M}'$ sur $\Gamma_{\hat{P},\sigma} \simeq \Gamma_{\mathbf{Q}}^\sim$, et détermine de même des scissions dans $\mathcal{M}''$, $\mathcal{M}_1'$, $Z_{\hat{P}}'$, $Z_{\hat{P}}$, avec des décompositions en produits 1/2 directs
$$
\begin{cases}
\mathcal{M}'_1 \simeq \mathcal{M}(-j) . \Pi_0 \\
\mathcal{M}'! \simeq \mathcal{M}(-j) . \pi \\
\mathcal{M}'' \simeq \mathcal{M}(-j) . \mathfrak{S}^{+\wedge} \\
Z_{\hat{P}}' \simeq \mathcal{M}(-j) \times \mu \\
Z_{\hat{P}} \simeq \mathcal{M}(-j) \times L_{\hat{P}}
\end{cases} \leqno{(143)}
$$

Notons maintenant que $\mu_{z'} \defeq \{1, z'\} \subset \mathfrak{S}^{+\wedge}$ normalise $L_{\hat{P}}$, donc $Z_{\hat{P}}! = Z_{\hat{P}} \cap \mathcal{M}!$

\newpage

%% ===== Page 619 =====

613

Revenons au diagramme (130), pour
dire que l'hom
$$ \overline{\gamma}_{\rho,\sigma} \longrightarrow \mathcal{N}_{\rho}^*/\{\pm 1\} $$
$$ \| \wr $$
$$ \overline{\Gamma}_{\tilde{Q}}^\sim $$
ne se relève pas en un hom de $\overline{\gamma}_{\rho,\sigma}$
dans $\mathcal{N}_{\rho}^*$ — en d'autres termes, que
l'extension de $\overline{\gamma}_{\rho,\sigma}$ par $\{\pm 1\} = \mu$ qu'elle
définit n'est pas triviale. Pour identifier
cette extension, considérons le diagramme

[MARGIN: NB]
(144)
\begin{tikzcd}
\mathcal{M}'[0] = \mathcal{Z}_{\rho,\sigma}^! \ar[r, "can"] & \overline{\gamma}_{\rho,\sigma} = \mathcal{Z}_{\rho,\sigma}^! / L_0^\pi \ar[r, "\text{hom (130)}"] & \mathcal{N}_{\rho}^* / \{ \pm 1 \} \\
S_0 \Gamma_{\rho,\sigma}^! \ar[u, "can \ (u,v,b) \mapsto b"] \ar[ur, "can \atop {(\text{passage au} \atop \text{quotient})}"'] & & \\
S_0'' \Gamma_{\rho,\sigma}^! \ar[u, "incl."] \ar[ur, "incl."'] \ar[rr] & & \mathcal{N}_{\rho}^* \ar[uu]
\end{tikzcd}

qui montre que l'on a un diagramme
de carrés cartésiens

(145)
\begin{tikzcd}
\mathcal{M}(0) \ar[r, hook] \ar[dr, "\text{I}" description, phantom] & \mathcal{M}'[0] = \mathcal{Z}_{\rho,\sigma}^! \ar[r, "\text{épi}"] \ar[dr, "\text{II}" description, phantom] & \overline{\gamma}_{\rho,\sigma} \ar[r, hook] \ar[dr, "\text{III}" description, phantom] & \mathcal{N}_{\rho}^* / \{ \pm 1 \} \\
\mathcal{M}(0)^\sim \ar[u] \ar[r, hook] & S_0'' \Gamma_{\rho,\sigma}^! \ar[u, "2"'] \ar[r] & \overline{\gamma}_{\rho,\sigma}^\sim \ar[u] \ar[r, hook] & \mathcal{N}_{\rho}^* \ar[u, "2"']
\end{tikzcd}

où $\mathcal{M}(0)^\sim$ est défini comme l'extension de
$\mathcal{M}(0)$ par $\mu$ induite par l'ext. $S_0'' \Gamma_{\rho,\sigma}^!$ de $\mathcal{M}'[0]$,
qui s'identifie de même à l'image inv. de

\newpage

%% ===== Page 620 =====

l'extension $\widetilde{\gamma}_{p,\sigma}$ de $\gamma_{p,\sigma}$ par $\mu$, par l'isomorphisme composé
$\mathcal{M}(0) \isommap \gamma_{p,\sigma}$. Donc l'extension $\widetilde{\gamma}_{p,\sigma}$ de $\gamma_{p,\sigma}$
par $\mu$ s'identifie naturellement à l'extension $\mathcal{M}(0)^{\sim}$ de
$\mathcal{M}(0)$ par $\mu$ dont on s'est convaincu plus
haut qu'elle n'était pas triviale.

Voici une autre façon d'obtenir
l'homomorphisme canonique $\gamma_{p,\sigma} \to N_{\hat{P}}^*/\{\pm 1\}$
et l'extension $\widetilde{\gamma}_{p,\sigma}$ de $\gamma_{p,\sigma}$. Considérons
$$ N_{\hat{P}}^{*!} = N_{\hat{P}}^* \cap \mathcal{M}^! \leqno{(146)} $$
On a alors
$$ N_{\hat{P}}^{*!} \cap \widehat{\mathfrak{S}}^+ = \operatorname{Ker}(N_{\hat{P}}^{*!} \to \gamma_{p,\sigma}) = L_{\hat{P}} \cap \Pi = \mu \leqno{(147)} $$
d'autre part l'homomorphisme
$$ N_{\hat{P}}^{*!} \to \gamma_{p,\sigma} \quad \text{est épi.} \leqno{(148)} $$
Pour [unclear: cela] il suffit de voir que $\forall$
$u \in \mathcal{M}(0)$, il existe un élément de
$N_{\hat{P}}^{*!}$ qui lui est conjugué, mod $\mathfrak{S}$.
Or on a
$u(p) = \operatorname{int}(\alpha)(p)$ avec $\alpha \in E \defeq \widehat{\Pi}_{\tau} = \widehat{\Pi}_{0,3} . \widehat{\Pi}$,
[unclear: (où $\operatorname{int}(\alpha) \circ \varepsilon_3(\tau)(p)$)]
Si $\varepsilon_3(u) = 1$ i.e. $\alpha \in \widehat{\Pi}$, $\alpha^{-1} u \in Z_{\hat{P}} \subset N_{\hat{P}}^*$
fait l'affaire. Si $\varepsilon_3(u) = -1$ on a
$\alpha = \alpha_0 \tau$, $\alpha_0 \in \widehat{\Pi}$
et on a
$u(p) = \alpha(p) = \alpha_0(\tau(p)) = \alpha_0(\sigma(p^{-1}))$ donc
$(\sigma^{-1}\alpha_0^{-1}u)(p) = p^{-1}$ donc $\sigma^{-1}\alpha_0^{-1}u \in N_{\hat{P}}^*$ OK.

\newpage

%% ===== Page 621 =====

Donc on a une suite exacte
$$
1 \to \mu \to Z_{\rho}^{*!} \to \gamma_{\rho,\sigma} \to 1 \leqno (149)
$$
(avec $Z_{\rho}^{*!} = \tilde{\gamma}_{\rho,\sigma}$ et $\gamma_{\rho,\sigma} \simeq \widetilde{\Pi}_Q$)

qui fait de $Z_{\rho}^{*!}$ une $\mu$-ext. de $\gamma_{\rho,\sigma}$, d'où un hom. can.
$$
\gamma_{\rho,\sigma} \simeq Z_{\rho}^{*!}/\mu \to N_{\rho}^*/\mu
$$
qui, bien entendu, n'est autre que le premier hom. (140).

\section*{90) Récapitulation sur le cas universel.}

Finalement, j'ai l'impression que j'ai fait encore des longs détours et contorsions, pour finalement un peu me perdre dans les sables, et perdre un peu le fil. Je vais donc récapituler ce qui me semble essentiel.

\begin{enumerate}
\item[a)] (1) $Z_{\rho,\sigma} \simeq \mathcal{M}(0) . L_0^{\pm} = \mathcal{M}(0)^{\sim}$ \footnote{opérant sur $\hat{\mathfrak{S}}^+$ par op. induite par int (dans $\mathcal{M}$).}
on distinguera donc un sous-groupe d'indice 2
(2) $Z_{\rho,\sigma}^+ = \mathcal{M}(0) . L_0^\pi$
\end{enumerate}

\newpage

%% ===== Page 622 =====

On a

(3) $$ \widehat{G}_{\rho,\sigma} \simeq \hat{M} $$

et, plus précisément, le diagramme

(4)
\begin{tikzcd}
1 \ar[r] & L_0 \ar[r] \ar[d, hook] & Z_{\rho,\sigma} \ar[r] \ar[d, hook] & \Gamma_{\rho,\sigma} \ar[r] \ar[d, equal] & 1 \\
1 \ar[r] & \widehat{C}^+_{\rho,\sigma} \ar[r] & \widehat{G}_{\rho,\sigma} \ar[r] & \Gamma_{\rho,\sigma} \ar[r] & 1
\end{tikzcd}

s'identifie au diagramme

(4 bis)
\begin{tikzcd}
1 \ar[r] & L_0 \ar[r] \ar[d, hook] & \tilde{M}[0] \ar[r] \ar[d, hook] & \tilde{\Gamma}_Q \ar[r] \ar[d, "\simeq"'] & 1 \\
1 \ar[r] & \widehat{C}^+ \ar[r] & \hat{M} \ar[r] & \tilde{\Gamma}_Q \ar[r] & 1
\end{tikzcd}

D'ailleurs, les hom.

(5)
\begin{tikzcd}
& \hat{\mathbb{Z}}^* \\
\Gamma_{\rho,\sigma} \ar[ur, "\varepsilon_{\rho} = \varepsilon_f"] \ar[dr, "\varepsilon_{\sigma}"'] \\
& \mu
\end{tikzcd}

s'identifient à :

(5 bis)
\begin{tikzcd}
& \mu \\
\tilde{\Gamma}_Q \ar[ur, "\varepsilon_3"] \ar[r, "\chi"] \ar[dr, "\varepsilon_2"'] & \hat{\mathbb{Z}}^* \\
& \mu
\end{tikzcd}

b) On a

(6) $$ N_0^* = \{ u \in \hat{M} \mid u(\sigma) = \varepsilon_3(u) \sigma^{\varepsilon_2(u)} \} \quad \begin{array}{l} (= \varepsilon_3(u) \cdot \sigma) \\ = \varepsilon_3(\tau)(\sigma) \\ = \varepsilon_3(\tau')(\sigma) \end{array} $$

[unclear]

(7)
$$
\begin{cases}
N_0^* \cap \widehat{G}^{\wedge} = D_0 = \{1, \tau \} \cdot L_0 = \{1, \tau' \} \cdot L_0 \\
N_0^* \cap \widehat{G}^{+\wedge} = L_0 \quad (\simeq \hat{\mathbb{Z}}) \\
N_0^* \cap \Pi = \mu \\
N_0^* \cap \Pi_+ = \{ 1 \}
\end{cases}
$$

\newpage

%% ===== Page 623 =====

Posons
$$
\begin{cases}
\mathcal{M}_0^! = \mathcal{M}(0) . \pi_0 \subset \mathcal{M}^! = \mathcal{M}(0) . \pi = \operatorname{Ker}(\mathcal{M} \to \mathfrak{S}_3) \\
\text{de sorte que } \mathcal{M}^! \simeq \mathcal{M}_0^! \times \mu
\end{cases} \leqno{(8)}
$$

et
$$
\begin{cases}
\mathcal{M}(1/2) \defeq \mathcal{N}^*_\sigma \cap \mathcal{M}_0^! \\
\mathcal{N}^{*!}_\sigma = \mathcal{N}^*_\sigma \cap \mathcal{M}^!
\end{cases} \leqno{(9)}
$$

alors l'homomorphisme $\mathcal{M} \to \Gamma_{\rho,\sigma}$ induit un iso
$$
\mathcal{M}(1/2) = \mathcal{N}^{*!}_{\sigma_0} \isommap \Gamma_{\rho,\sigma} \leqno{(10)}
$$

d'où des scindages en produits semi-directs
$$
\begin{cases}
\mathcal{M}^\wedge \simeq \mathcal{M}(1/2) . \widehat{\mathfrak{S}}^{+\wedge} \\
\mathcal{M}^! \simeq \mathcal{M}(1/2) . \pi \\
\mathcal{M}_0^! \simeq \mathcal{M}(1/2) . \pi_0 \\
\mathcal{N}^*_\sigma \simeq \mathcal{M}(1/2) \times \mu \\
\mathcal{N}^\wedge_{\sigma_0} \simeq \mathcal{M}(1/2) . L_\sigma \quad \text{, et}
\end{cases} \leqno{(11)}
$$

$$
\Gamma_{\rho,\sigma} \xleftarrow{\sim} \mathcal{M}(1/2) = \mathcal{N}^{*!}_{\sigma_0} \hookrightarrow \mathcal{N}^*_\sigma . \leqno{(12)}
$$

\begin{enumerate}
\item[c)] On a
$$
\mathcal{N}^*_\rho = \{ u \in \mathcal{M} \mid u(\rho) = \rho^{\varepsilon_3(u)} \} \quad (= \mathcal{E}_3[\tau'](\rho)) \leqno{(13)}
$$

et
$$
\begin{cases}
\mathcal{N}^*_\rho \cap \widehat{\mathfrak{S}} = D_\rho = \{1, \tau'\} . L_\rho \\
\mathcal{N}^\wedge_\rho \cap \widehat{\mathfrak{S}} = L_\rho \quad (\simeq \hat{\mathbf{Z}} / 6 \hat{\mathbf{Z}}) \\
\mathcal{N}^*_\rho \cap \pi = \mu \\
\mathcal{N}^\wedge_\rho \cap \pi_0 = \{1\}
\end{cases} \leqno{(14)}
$$
\end{enumerate}

\newpage

%% ===== Page 624 =====

Soit
$$
N_{\rho}^{*!} \defeq N_{\rho}^* \cap \mathcal{M}! , \quad N_{\rho}^{\wedge!} \defeq N_{\rho}^\wedge \cap \mathcal{M}! \leqno{(15)}
$$

on trouve
$$
N_{\rho}^{*!} \subset Z_{\rho} = \operatorname{Centr}_{\mathcal{M}}(\rho) = \operatorname{Centr}_{\mathcal{M}}(L_{\rho}) \leqno{(16)}
$$

et on note aussi
$$
N_{\rho,0}^{*!} = Z_{\rho,0}^! = \mathcal{M}'(-J). \leqno{(17)}
$$

On voit que l'homomorphisme $\mathcal{M} \to \gamma_{\rho,\sigma} = \Gamma_{\mathbf{Q}}^\sim$ induit un isomorphisme
$$
\mathcal{M}'(-J) = Z_{\rho,0}^! \isommap \gamma_{\rho,\sigma}^! \defeq \{ \gamma \in \gamma_{\rho,\sigma} \mid \varepsilon_3(\gamma) = 1 \} \leqno{(18)}
$$

d'où l'on déduit des scindages en produits semi-directs (voire produits directs)
$$
\begin{cases}
\mathcal{M} \simeq \mathcal{M}'(-J) . \mathfrak{S}^{+\wedge} \\
\mathcal{M}' \simeq \mathcal{M}'(-J) . \pi \\
\mathcal{M}'_0 \simeq \mathcal{M}'(-J) . \pi_0 \\
Z_{\rho}^! \simeq \mathcal{M}'(-J) \times \mu \\
Z_{\rho} \simeq \mathcal{M}'(-J) \times L_{\rho}
\end{cases} \leqno{(19)}
$$

\newpage

%% ===== Page 625 =====

Soit maintenant
$$
\mu_{\tau'} = \{1, \tau'\} \subset \hat{\mathfrak{S}}^\wedge \leqno{(20)}
$$
alors $\mu_{\tau'}$ normalise $L_\rho \simeq \hat{\mathbf{Z}}$ car
$\tau'(\rho) = \rho^{-1}$ d'où $\tau'(L_\rho) = L_\rho$
et il normalise $Z_\rho = \operatorname{Centr}_{\mathcal{M}}(L_\rho)$, donc
aussi $Z_\rho^! = Z_\rho \cap \mathcal{M}^!$ (puisque $\mathcal{M}^!$ est
invariant dans $\mathcal{M}$) $= Z_{\rho,0}^! \times \mu = \mathcal{M}^![-j] \times \mu$.
D'ailleurs $\mu_{\tau'} \cap Z_\rho^! = \mu_{\tau'} \cap \mathcal{M}^! = \{1\}$.
$$
N_\rho^? = \mu_{\tau'} \cdot Z_\rho^! = \mu_{\tau'} \cdot (Z_{\rho,0}^! \times \mu) \leqno{(21)}
$$
$$
(\text{produit } 1/2 \text{ direct}) \quad \mathcal{M}[-j] \rtimes \mu \simeq \Gamma_{\mathbf{Q}}^\sim \times \mu
$$
On fera attention que contrairement à
ce que pourrait suggérer cette écriture,
$\mu_{\tau'}$ ne normalise pas $Z_{\rho,0}^!$ - sinon
on aurait une décomposition en produit
$N_\rho^? = N_{\rho,0}^? \times \mu$, où $N_{\rho,0}^? = \mu_{\tau'} \cdot Z_{\rho,0}^?$, or
il n'en est rien. On a une suite exacte
$$
\begin{tikzcd}
1 \arrow[r] & \mu \arrow[r] & N_\rho^? \arrow[r] & \gamma_{\rho,\sigma}^? \arrow[r] \arrow[d, equal] & 1 \\
& & & \Gamma_{\mathbf{Q}}^\sim & 
\end{tikzcd} \leqno{(22)}
$$
de sorte que $N_\rho^?$ est une quasi-section
du multiplicateur de l'extension
$N_\rho^\wedge \to \gamma_{\rho,\sigma}$, de $\mathcal{M} \to \gamma_{\rho,\sigma} = \Gamma_{\mathbf{Q}}^\sim$.

Ainsi on trouve

\newpage

%% ===== Page 626 =====

$$
\begin{cases}
\mathcal{M} = N_{\rho}^? \wedge_{\mu_2} \mathfrak{S}^{+\wedge} \\
M_{\rho}^? = N_{\rho}^? \wedge_{\mu_2} \mathbb{L}_{\rho} \simeq N_{\rho}^? \times \underset{\simeq \mathbf{Z} / 3 \mathbf{Z}}{L_{\rho}}
\end{cases} \leqno{(23)}
$$
L'[unclear: extension] (22) de $\Gamma_{\rho, \sigma}$ par $\mu \hookrightarrow$ l'extension [unclear: étudiée] $\Gamma_{\rho, \sigma}$ de la section précédente, elle est [unclear: canoniquement] scindée.
Elle donne naissance à un isomorphisme canonique
$$
\Gamma_{\rho, \sigma} \simeq N_{\rho}^? / \mu \longrightarrow N_{\rho}^* / \mu . \leqno{(24)}
$$

Remarque. Pour bien faire il faudrait expliciter l'opération de $\tau'$ sur $Z_{\rho}^? \simeq Z_{\rho}^! \times \mu \simeq \widetilde{\Gamma_{\mathbf{Q}}} \times \mu$, en termes de l'opération de l'image $\bar{\tau}'$ de $\tau'$ sur $\widetilde{\Gamma_{\mathbf{Q}}}$ par automorphisme intérieur (NB $\tau'$ n'est défini par l'élément de $\Gamma_{\mathbf{Q}} \subset \widetilde{\Gamma_{\mathbf{Q}}}$, "à conjugaison près") et d'un "caractère quadratique" de $\widetilde{\Gamma_{\mathbf{Q}}}$ : il y a fort à parier que celui-ci est juste $\varepsilon_2$ - de sorte que $\widetilde{\Gamma_{\mathbf{Q}}}^0 = \widetilde{\Gamma_{\mathbf{Q}}} \cap \Gamma_{\mathbf{Q}}^0$ apparaîtrait comme invariante par $\tau'$...
Il y aurait lieu de reprendre l'étude de l'extension faite dans la section précédente au point de vue des pôles en un des $\alpha, \beta, \gamma$.

\newpage

%% ===== Page 627 =====

de l'optique actuelle, dans des termes plus
simples. Il y aurait lieu en même temps
d'expliciter une autre section partielle,
qui serait au-dessus de $\Gamma_{\mathbf{Q}}^\sim$,
$$
\Gamma_{\mathbf{Q}}^{\sim \sim} = \{ x \in \Gamma_{\mathbf{Q}}^\sim \mid \varepsilon_1(x) \varepsilon_3(x) = 1 \}
$$
et qui mériterait peut-être la notation
\marginpar{? \\ Mieux vaut \\ s'en \\ méfier !} $\mathcal{M}(\sqrt{3})$ ou au moins $\mathcal{M}(\rho)$ , pour
un générateur $\rho$ convenable de $\mathbf{Q}(\sqrt{3})$ ???
Je me rends compte que je suis en train
de déconner - ou à peu près perdre contact
avec le contenu géométrique de mes
calculs. Il faut garder à l'esprit que
$$
\mathcal{M}_0^1 \isommap \mathcal{M}_{0,3}^{1 \sim} \leqno{(25)}
$$
correspond - à $\sim$ près, au $\pi_1$ de $U_{0,3 \mathbf{Q}}$ ,
les sections de $\mathcal{M}_0^1 \to \Gamma_{\mathbf{Q}}^\sim \simeq \overline{\Gamma}_{\sigma}$ correspondent
donc (à $\sim$ près) : celles de $\pi_1(U_{0,3 \mathbf{Q}})$ sur $\Gamma_{\mathbf{Q}}$ ,
donc ( "inutilement" ) aux points de $U_{0,3 \mathbf{Q}}$
points de $U_{0,3 \mathbf{Q}}$ : à [unclear: chacune des deux catégories]
rationnels sur $\mathbf{Q}$. Ainsi, les sections
dont b) partielle $\mathcal{M}(-1)$ correspond bien au pt $1/2$ de $U_{0,3 \mathbf{Q}}$ ,
la section envisagée ici dans c) correspond -
au point $-j$ de $U_{0,3 \mathbf{Q}}$ $\mathcal{M}(j)$ , qui est dans $U_{0,3 \mathbf{Q}}(\mathbf{Q}[\sqrt{-1}])$ ,
d'où la notation $\mathcal{M}(-j)$ raisonnable. [unclear: Comme dans] $\mathcal{M}_0^\sim$
elle n'est pas $\subset \mathcal{M}_0^1$ ! on y divise par $\mu$ ;
tout au moins diviser par $\mu$ [unclear: pour une] section

\newpage

%% ===== Page 628 =====

de $M_{0,3}^\sim$ sur $\Gamma_{\mathbb{Q}}^\sim$, mais on n'a pas une
section de $M_{0,3}^\sim$. Or à un $\sim$ près, $M_{0,3}^\sim$
est le $\pi_1$ de $M_{1,1 \mathbb{Q}}$, schéma modulaire sur $\mathbb{Q}$
des courbes elliptiques mod. symétrie.
Donc la donnée de cette section correspond à
la donnée d'une telle "courbe elliptique
mod. symétrie" sur $\mathbb{Q}$. Passant au sous-
groupe $\Gamma_{\mathbb{Q}}'^\sim$ - ce qui correspond à l'ex-
tension du corps de base de $\mathbb{Q}$ à $\mathbb{Q}(j) = \mathbb{Q}(\sqrt[3]{1})$ -
on retrouve l'image dans $M_{0,3}'^\sim$ de la quasi
section pointée $M(-j)$ - cela montre
donc qu'il s'agit toujours de la "même" courbe
elliptique. J'ai bien l'impression cependant
que sur $\Gamma_{\mathbb{Q}}^\sim$ la section, cette section de
$M_{0,3}^\sim$ ne provient pas d'une section
de $M_{1,1}^\sim$ sur $\Gamma_{\mathbb{Q}}^\sim$ - i.e. justement que l'ext. $N_{\mathbb{Q}}^*$
de $\Gamma_{\mathbb{Q}}^\sim$ par $\mu_2$ n'est pas triviale. Cela signifierait
il qu'il n'existe pas de courbe elliptique
déf / $\mathbb{Q}$ d'invariant $0$ ? (Chose qui semble absurde
- car doit pouvoir sur un corps tr. donner n'importe
[MARGIN: ?] quel invariant ...). De façon précise, on a le
diagramme de schémas modulaires sur $\mathbb{Q}$

(26)
\begin{tikzcd}
\text{sch. mod. } M_{1,1} \text{ des courbes ell.} \ar[d, "\text{gerbe [unclear]}"'] \ar[dr] & \\
\text{sch. mod. des } M'_{1,1} \text{ courbes ellipt. mod. symétrie} \ar[r] & E' \text{ schéma modulaire grossier}
\end{tikzcd}

Le fait qu'un point de $M'_{1,1}(\mathbb{Q})$ ne se re-

\newpage

%% ===== Page 629 =====

monte par un [unclear: cas] pour le $M_{1,1}$, ça signifie
que son image dans $E'$ ne se relève pas
à $M_{1,1}$. De façon précise, la digression
s'identifie à :
$$
\begin{tikzcd}
M_{1,1} \simeq (\mathcal{U}_{0,3}, \widetilde{\mathfrak{S}}_3) \arrow[d] \arrow[dr] & \\
M'_{1,1} \simeq (\mathcal{U}_{0,3}, \mathfrak{S}_3) \arrow[r] & \mathcal{U}_{0,3}/\mathfrak{S}_3 \simeq \mathcal{U}_{0,3}/\widetilde{\mathfrak{S}}_3 = E'
\end{tikzcd}
\leqno{(27)}
$$
où $\widetilde{\mathfrak{S}}_3$ est l'extension de $\mathfrak{S}_3$ par $\mu$, définie par
$$
\widetilde{\mathfrak{S}}_3 = \widetilde{\mathfrak{S}}^+ / \Pi_0 \simeq \operatorname{Sl}(2,\hat{\mathbf{Z}}) / \text{sous-groupe invariant engendré par } -e_0
\leqno{(28)}
$$
$$
\simeq \operatorname{Sl}(2,\mathbf{Z}/4\mathbf{Z}) / \text{sous-groupe (d'ordre 4) engendré par } -e_0, e_1 \dots
$$
on a une action évidente de $\widetilde{\mathfrak{S}}_3$ sur $\mathcal{U}_{0,3}$ par l'interm. de $\mathfrak{S}_3$. Un point de $M_{1,1}$\footnote{Tout le fatras des catégories fibrées de Giraud, et celui des "champs" de Deligne, me semblent d'inutiles complications pour définir des problèmes de "modules" (tels $M_{1,1}$), ou des "champs de modules", ou exprimer (comme je le fais plus bas) le fait que tel problème sur $\mathbf{Q}$ provient du problème des "courbes elliptiques".} s'identifie à
un $\widetilde{\mathfrak{S}}_3$-torseur $T$ sur $\mathbf{Q}$, avec
un morphisme
$$
T \to \mathcal{U}_{0,3}
$$
compatible avec les actions de $\widetilde{\mathfrak{S}}_3$, les
deux torseurs $M_{1,1}$ s'identifiant : un $\widetilde{\mathfrak{S}}_3$-torseur
$\widetilde{T}$ sur $\mathbf{Q}$, muni d'un morphisme
$$
\widetilde{T} \to \mathcal{U}_{0,3}
\leqno{(30)}
$$
compatible avec l'opération de $\widetilde{\mathfrak{S}}_3$. Désignons par
$$
\mu \hookrightarrow \widetilde{T}
$$
$$
T = \widetilde{T} / \mu
$$
le $\mathfrak{S}_3$-torseur correspondant, le donné de
(30) équivaut à la donnée de (29). Dans la

\newpage

%% ===== Page 630 =====

relèvements à $M_{1,1}$ d'un point de $M_{1,1}'$ à / 0
s'identifient aux $\widetilde{\mathfrak{S}}_3$-torseurs $\widetilde{T}$ qui relèvent le
$\mathfrak{S}_3$-torseur $T$, i.e. (à iso près) ils correspondent
[MARGIN: en supp. choisi un pt base de $T(\bar{\mathbb{Q}})$]
aux hom. $\Gamma_{\mathbb{Q}} \to \widetilde{\mathfrak{S}}_3$ qui relèvent l'hom
$\Gamma_{\mathbb{Q}} \to \mathfrak{S}_3$ qui définit $T$ :

\begin{tikzcd}
& \widetilde{\mathfrak{S}}_3 \ar[d] \\
\Gamma_{\mathbb{Q}} \ar[r] \ar[ur, dashed] & \mathfrak{S}_3
\end{tikzcd}

l'obstruction est, comme il est d'usage, une
extension de $\Gamma_{\mathbb{Q}}$ par $\mu_2$, i.e. un élément
de $H^2(\mathbb{Q}, \mu_2) \simeq {}_2\text{Br}(\mathbb{Q})$. Dans le cas
qui nous occupe l'image de $T$ dans $U_{0,3}$
est formé de $\{ -j, -j^2 \}$ qui est connexe sur $\mathbb{Q}$,
[MARGIN: condition d'[unclear] plus bas]
ce qui entraine $\Gamma_{\mathbb{Q}} \to \mathfrak{S}_3$ est surjectif. De façon
générale, en [unclear] le th. de [unclear]
galoisienne en termes de scindages
de $M_{0,3}$ sur $\Gamma_{\mathbb{Q}}$, l'hom. $\Gamma_{\mathbb{Q}} \to \mathfrak{S}_3$
est le composé
$$ \Gamma_{\mathbb{Q}} \xrightarrow{\kappa} M_{0,3} \to \mathfrak{S}_3 ,$$
où $\kappa$ est une section de
$$ M_{0,3}^1 = \operatorname{Ker}(M_{0,3} \to \mathfrak{S}_3) $$
sur $\Gamma_{\mathbb{Q}}$. Ici une section de
$M_{1,1}'$ au dessus de $\Gamma_{\mathbb{Q}}$ donne une section de
$M_{1,1}' \xrightarrow{\text{au dessus de}} \dots \Gamma_{\mathbb{Q}}'$ d'ordre 2,
donc $\kappa(\Gamma_{\mathbb{Q}}) = 1$, on voit que

\newpage

%% ===== Page 631 =====

c'est le sous-groupe $\{1, \tilde{\sigma}_\infty\}$ de $\widetilde{\mathfrak{S}}_3$. (Car l'image de
$\tau' \in M_{0,3}$ dans $\widetilde{\mathfrak{S}}_3$ est (comme celle de $\tau = 1$)
est égale à celle de $\sigma = \sigma_\infty$, i.e. $\tilde{\sigma}_\infty$. Ce signifie
que l'hom.
$$T \simeq T_0 \wedge_{\{1, \tilde{\sigma}_\infty\}} \widetilde{\mathfrak{S}}_3$$
i.e. $T$ se déduit d'un $\mu$-torseur (i.e. d'une
ext. quadratique de $\mathbb{Q}$, savoir justement
$\mathbb{Q}(\sqrt{j}) = \mathbb{Q}(j)$) par extension du groupe
$$ \mu \hookrightarrow \widetilde{\mathfrak{S}}_3 \quad (-1 \mapsto \tilde{\sigma}_\infty)$$
d'opérateurs. Donc, prenant l'image inverse
de $\mu = \{1, \tilde{\sigma}_\infty\}$ dans $\widetilde{\mathfrak{S}}_3$, on trouve une
extension de $\mu$ par $\mu$ - ça ne peut guère être

[MARGIN: oui, c'est clair puisqu'on a $\tilde{\sigma}_\infty$ dans $\widetilde{\mathfrak{S}}_3 = [unclear]$ est bien d'ordre 4 donc $\mathcal{E}'' \simeq \mathbb{Z}/4\mathbb{Z}$ canoniquement - ]

que $\mathbb{Z}/4\mathbb{Z}$ ! - que j'ai envie de noter
$\mathcal{E}''$ :
$$1 \to \mu \to \overset{\mathbb{Z}/4\mathbb{Z}}{\mathcal{E}''} \to \mu \to 1$$
et le pb d'obstruction est de
trouver un remontage

\begin{tikzcd}
\mathcal{E}'' \ar[r] & \mu \\
\Gamma_{\mathbb{Q}} \ar[u, dashed] \ar[ur] &
\end{tikzcd}

[MARGIN: C'est évident - [unclear: il n'y a qu'à voir] $(\Gamma_{\mathbb{Q}})_{ab} \simeq \hat{\mathbb{Z}}^\times \supset \mathbb{Z}_3^\times$, or on a $(\mathbb{Z}_3^\times)_4 \simeq \mathbb{Z}/2\mathbb{Z}$ (pas de r. q. dans $\mathbb{Z}_3$ [unclear]) - ]

de l'obstruction - s'interprète comme l'extension
de $\Gamma_{\mathbb{Q}}$ par $\mu$, image inverse de l'extension
$\mathcal{E}''$. (Il faudrait regarder pourquoi, en termes

\newpage

%% ===== Page 632 =====

classiques, cette obstruction est trivialisée, quand on passe au sous-groupe $\Gamma_{\mathbf{Q}}'''$ correspondant au corps $\mathbf{Q}(\sqrt{3})$)

Ainsi, le point $\mathbf{Q}$-rationnel envisagé de $M_{1,1}$ ne se remonte pas en un pt $\mathbf{Q}$-rationnel de $M_{1,1}$. Par contre, [unclear: son image] dans $E' \simeq U_{0,3}/\mathfrak{S}_3$ doit bien se remonter - i.e. il doit bien exister un $\widetilde{\mathfrak{S}}_3$-torseur $\widetilde{P}$, et un $\widetilde{\mathfrak{S}}_3$-morphisme
$$ \widetilde{P} \to U_{0,3} \quad \text{i.e} \quad \begin{tikzcd} P \arrow[r] \arrow[d, equal] & U_{0,3} \\ \widetilde{P}/\mu & \end{tikzcd} $$

Considérons en effet le [unclear: schéma au-dessus] de $Q = U_{0,3} \times_{E'} \{0\}$
l'image inverse de $0$, soit :
$$ P_0 = Q_{\mathrm{red}} $$
c'est un sous-schéma fini, qui sur un corps $K \supset \mathbf{Q}(j)$, se scinde en exactement deux pts, et qui sur $\mathbf{Q}$ est connexe - c'est un schéma isomorphe à :
$\operatorname{Spec}\ \mathbf{Q}(j)$ . Il est stable par $\mathfrak{S}_3$, qui y opère via
$$ \mathfrak{S}_3 \xrightarrow{sgn} \mu $$
et l'opération galoisienne de $\mu$ sur le $\mu$-torseur $P_0$.

\newpage

%% ===== Page 633 =====

Bref ... le diagramme

(31)
\begin{tikzcd}
\tilde{\mathfrak{S}}_3 \ar[r, "can"] & \mathfrak{S}_3 \ar[r, "sgn"] & \mu \\
& \Gamma_{\mathbb{Q}} \ar[u] \ar[ur] &
\end{tikzcd}

[MARGIN: NB le noyau de $\tilde{\mathfrak{S}}_3 \to \mu$ est une extension centrale de $\mathbb{Z}/3\mathbb{Z}$ par $\mu = \mathbb{Z}/2\mathbb{Z}$, elle est donc triviale et isom. : $\mathbb{Z}/6\mathbb{Z}$ (unique [unclear])]

On voit donc que la catégorie des relèvements
d'un pt $\mathbb{Q}$-rationnel $e$ de $E'$ en un
pt $\mathbb{Q}$-rationnel de $M_{1,1}^\sim$ resp. $M_{1,1}$, est
équivalente à : celle des $\tilde{\mathfrak{S}}_3$-torseurs $\tilde{P}$, resp.
des $\mathfrak{S}_3$-torseurs $P$ sur $\mathbb{Q}$, qui relèvent le
$\mu$-torseur $P_0$. On avait choisi un
$P=T$ de façon particulièrement triviale,
en utilisant le scindage
$$ \mu \xrightarrow{\sim} \{1, \sigma_3\} \subset \mathfrak{S}_3 $$
de $\mathfrak{S}_3 \to \mu$ — mais ce torseur $T$ ne se
remonte pas en un $\tilde{T}$. Les classes d'isomor-
phie de pts de $M_{1,1}^\sim$ resp. de $M_{1,1}$ au-dessus
du pt $e$ de $E'$, correspondent, à [unclear] près
(mod noyaux de $\tilde{\mathfrak{S}}_3 \to \mu$ et $\mathfrak{S}_3 \to \mu$),
aux classes de conjugaisons (des relèvements de
$\Gamma_{\mathbb{Q}} \to \mu$ en $\Gamma_{\mathbb{Q}} \to \mathfrak{S}_3$ resp. en $\Gamma_{\mathbb{Q}} \to \tilde{\mathfrak{S}}_3$).

Considérons un relèvement
$$ \Gamma_{\mathbb{Q}} \to \tilde{\mathfrak{S}}_3 $$
de ce type... un sous-groupe de $\tilde{\mathfrak{S}}_3$, qui n'est

\newpage

%% ===== Page 634 =====

$\{1, \sigma_3\}$ --- ou un de ses conjugués $\{1, \sigma_0\}, \{1, \sigma_1\}$
ou c'est $\mathfrak{S}_3$ tout entier. Dans le premier
cas, on voit que $p_0$ ne se remonte pas
en $\widetilde{\mathfrak{S}}_3$. Donc pour pouvoir remonter
il faut prendre un épi $\boldsymbol{\Gamma}_{\mathbf{Q}} \to \mathfrak{S}_3$ ---
i.e. [unclear: sortir d'une] extension abélienne sur
$\mathbf{Q}$, p. ex. entrer dans le domaine des
extensions abéliennes de $\mathbf{Q}(j)$.

Examiner en termes des groupes de
Galois s'il existe un tel relèvement genre
$\boldsymbol{\Gamma}_{\mathbf{Q}} \to \widetilde{\mathfrak{S}}_3$, et comment les obtenir,
apparaît comme un exercice plaisant
et délectable de théorie du corps de
classes, que je ne vais pas poursuivre
pour le moment, pour ne pas éterniser
cette digression. Revenant par la suite
sur ce point, il me faudrait en même
temps expliciter les relations entre les
points de vue remontages d'homomorphisme de $\boldsymbol{\Gamma}_{\mathbf{Q}}$,
et scindages de l'extension $\mathcal{M}_{0,3}$ de $\boldsymbol{\Gamma}_{\mathbf{Q}}$
par $\mathfrak{S}_{0,3}^{+\wedge}$, ou $\mathcal{M}$ de $\boldsymbol{\Gamma}_{\mathbf{Q}}$ par $\mathfrak{S}^{+\wedge} = \operatorname{Sl}(2, \hat{\mathbf{Z}})^\wedge$.

\newpage

%% ===== Page 635 =====

Revenons au sous-groupe
$$
\begin{tikzcd}
\mathcal{N}_{\rho}^? \arrow[r, hook, "3"] \arrow[d, dash, "4"] & \mathcal{N}_{\rho}^* \arrow[r, hook] & \mathcal{M} \\
Z_{\rho}^! = \mathcal{M}(-\overline{j}) & &
\end{tikzcd}
$$
contenant la section partielle $\mathcal{M}[-j] \isom \Gamma_{\mathbf{Q}}^\sim$.

Soit
$$
\mathcal{M}^\circ[-j] \overset{2}{\subset} \mathcal{M}'[-j]
$$
$$
\Ker ( \mathcal{M}[-j] \overset{\varepsilon_2}{\longrightarrow} \mu ) = \mathcal{N}_{\rho}^* \cap \mathcal{M}^\circ ! \leqno{(32)}
$$
qui est donc une section partielle au dessus
de
$$
\Gamma_{\mathbf{Q}}^{\circ \sim} = \Ker ( \Gamma_{\mathbf{Q}}^\sim \overset{(\varepsilon_2, \varepsilon_3)}{\longrightarrow} \mu \times \mu ) \leqno{(33)}
$$
$$
= \Ker ( \Gamma_{\mathbf{Q}}^\sim \longrightarrow \hat{\mathbf{Z}}^* \longrightarrow (\mathbf{Z}/12\mathbf{Z})^* )
$$

Modulo une vérification que je vais faire plus bas, $\tau'$ normalise $\mathcal{M}^\circ[-j]$, considérons dès lors
$$
\mathcal{M}''[-j] \defeq \mu_2 \cdot \mathcal{M}^\circ[-j] \quad (\text{semi-direct}) \quad , \mu_2 = \{1, \tau'\} \leqno{(34)}
$$
c'est un sous-groupe d'indice $4$ de $\mathcal{N}_{\rho}^?$,
et on a maintenant
$$
\mathcal{M}''[-j] \isom \Gamma_{\mathbf{Q}}^\sim \ , \leqno{(35)}
$$
c'est une section partielle au dessus
de $\Gamma_{\mathbf{Q}}^{\sim}$, d'où des isomorphismes

\newpage

%% ===== Page 636 =====

(36)
$$
\begin{cases}
\mathcal{M}''' \simeq \mathcal{M}'''(-\overline{j}) . C^{\pm} \\
\mathcal{N}_p^* \simeq \mathcal{M}'''(-\overline{j}) . L_p
\end{cases}
$$

On a le diagramme récapitulatif
d'inclusions de sous-groupes de $\mathcal{N}_p^*$

(37)
\begin{tikzcd}
& & Z_p \ar[ddr, "2"] \\
& \mathcal{M}'(-\overline{j}) = Z_{\hat{p}}! \ar[ur, "6"] \ar[dr, "4"] \\
\mathcal{M}^0(-\overline{j}) \ar[ur, "2"] \ar[dr, "2"'] & & \mathcal{N}_p^? \ar[r, "3"] & \mathcal{N}_p^* \\
& \mathcal{M}'''(-\overline{j}) \ar[ur, "4"']
\end{tikzcd}

Notons que l'inclusion

(38) $$ \underset{\simeq \mathbb{D}_6}{\mathbb{D}_p} \overset{df}{=} \underset{\simeq \mathbb{Z}/2 \cdot \mathbb{Z}/6\mathbb{Z}}{\mu_{\pm} \cdot L_p} \subset \mathcal{N}_p^* $$

induit un isomorphisme d'

(39) $$ \mathbb{D}_p \overset{\sim}{\longrightarrow} \mathcal{N}_p^* / \text{[unclear]} $$

\newpage

%% ===== Page 637 =====

(36) 
$$ \begin{cases} M''' \simeq M'''(-\bar{J}) . C_2^{+\wedge} \\ N_{\rho}^* \simeq M'''(-\bar{J}) . L_{\rho} \end{cases} $$

Diagramme récapitulatif des sous-groupes de $N_{\rho}^*$ obtenus jusqu'à présent :

(37)
\begin{tikzcd}
& M^0(-\bar{J}) \ar[rr, "2"] \ar[dl, "2"'] && \underset{=Z^0_{\rho}}{M'(-\bar{J})} \ar[rr, "2"] \ar[dl, "2"'] && M'(\bar{J})_{xt} \ar[rr, "3"] \ar[dl, "2"'] && Z_{\rho} \ar[dl, "2"'] \ar[dd, "2"] \\
M'''(-\bar{J}) \ar[rr, "2"] && M''(-\bar{J}) \times \mu_C \ar[rr, "2"] && N_{\rho}^{\pm C} \ar[rr, "3"] && N_{\rho}^* \ar[dr] \\
&& \mu_{\pm 1 C} \ar[ul, "1"] \ar[rr, "2"] && D_{\pm 1 C} \ar[ul, "\mu_C"'] \ar[rr, "3"] && D_{\rho} \ar[ul, "2"'] \ar[ur, "2"'] & L_{\rho}
\end{tikzcd}

où nous omettons les inclusions relatives
à 2 "faux sous-groupes", comme les
sous-groupes d'indice fini de $N_{\rho}^*$,
intermédiaires entre celui-ci et $M^0(-\bar{J})$
(d'indice 24 dans $N_{\rho}^*$). Je vais examiner
les questions d'isomorphismes de ces sous-groupes
les uns dans les autres, et la structure
des groupes quotients.

Je vais faire un diagramme qui
montrera bien les conglomérations entre
les groupes en question.

\newpage

%% ===== Page 638 =====

$$
\begin{tikzcd}[row sep=2.5em, column sep=1.5em]
& & Z_{\widehat{\rho}} \arrow[rr, "2"] \arrow[ddd, dashed] & & \mathcal{N}_{\widehat{\rho}}^\ast \arrow[ddd, dashed] & \\
& Z_{\widehat{\rho}}' = \mathcal{M}'(\overline{j}) \times \mu \arrow[ur, "3"] \arrow[rr, "2" near start] \arrow[ddd, dashed] & & \mathcal{N}_{\widehat{\rho}}' \arrow[ur, "3"] \arrow[rr, "2"] \arrow[ddd, dashed] & & \mathcal{M}'(\overline{j}) \cdot \mu_{T^2}^\ast\footnote{* Attention, $\mathcal{M}'(\overline{j}) \cdot \mu_{T^2}$ n'est pas un sous-groupe car $\mu_{T^2}$ ne normalise pas $\mathcal{M}'(\overline{j})$.} \arrow[ddd, dashed] \\
Z_{\widehat{\rho}}^\circ = \mathcal{M}^\circ(\overline{j}) \times \mu \arrow[ur, "2"] \arrow[rr, "2" near start] \arrow[ddd, dashed] & & \mathcal{N}_{\widehat{\rho}}^\circ \arrow[ur, "1"] \arrow[rr, "2"] \arrow[ddd, dashed] & & \mathcal{M}''(\overline{j}) \arrow[ur, "2"] \arrow[ddd, dashed] & \\
& & L_{\widehat{\rho}} \arrow[rr, dashed] & & D_{\widehat{\rho}} & \\
& {} \arrow[ur, dashed] \arrow[rr, dashed] & & D_{T^2}^! \arrow[ur, dashed] \arrow[rr, dashed] & & {} \\
{} \arrow[ur, dashed] \arrow[rr, dashed] & & {} \arrow[ur, dashed] \arrow[rr, dashed] & & {} \arrow[ur, dashed] &
\end{tikzcd} \leqno{(38)}
$$

a) Sous-groupes d'indice fini dans $Z_{\widehat{\rho}}$ -- dans le diagramme il y en a une [unclear: partie]

$$
\begin{tikzcd}
Z_{\widehat{\rho}}' = \mathcal{M}'(\overline{j}) \times \mu \arrow[r, "2"] \arrow[d] & \mathcal{M}'(\overline{j}) = Z_{\widehat{\rho}}^! = \mathcal{N}_{\widehat{\rho}}^! \arrow[d] \\
\mathcal{M}^\circ(\overline{j}) \times \mu \arrow[r, "2"] & \mathcal{M}^\circ(\overline{j}) = Z_{\widehat{\rho}}^{\circ !}
\end{tikzcd} \leqno{(39)}
$$

sont-ils invariants dans $Z_{\widehat{\rho}}$ ? C'est clair pour les sous-groupes [unclear: rayés]
$$
\begin{cases}
\mathcal{M}'(\overline{j}) \times \mu = Z_{\widehat{\rho}} \cap \mathcal{M}'! = Z_{\widehat{\rho}}', \text{ et pour} \\
\mathcal{M}^\circ(\overline{j}) \times \mu = Z_{\widehat{\rho}} \cap \mathcal{M}' \cap \operatorname{Ker}(\varepsilon_2 : \mathcal{M} \to \mu)
\end{cases} \leqno{(40)}
$$
- cela montre que ces groupes sont [unclear: normaux]

\newpage

%% ===== Page 639 =====

même invariants dans $N_\rho^*$, puisque $Z_\rho, \mathcal{M}', \Ker \varepsilon_2$ sont stabilisés par $N_\rho^*$.
D'ailleurs, $Z_\rho^! = \mathcal{M}'(-j) \times \mu$, pour montrer qu'il est
dans $Z_\rho$, il suffit de prouver qu'il est
normalisé par $\rho$ (puisque $Z_\rho = \langle \rho \rangle \cdot Z_\rho^!$ - c'est
$\pm$ trivial), on examine de façon générale
pour $x \in \mathcal{M}$ si il normalise $\mathcal{M}_0^!$.
$$
\mathcal{M}_0^! = \mathcal{M}(0) \cdot \Pi_0 \ ,
$$
comme il normalise $\Pi_0$, il suffit de
voir si pour $u \in \mathcal{M}(0)$, on a
$$
x u x^{-1} \in \mathcal{M}_0^!
$$
Ce qui équivaut encore à
$$
x u x^{-1} \equiv u \pmod{\Pi_0}
$$
i.e.
$$
[x, u] \in \Pi_0
$$
Or on sait que pour $x \in \widehat{\mathfrak{S}}$, $u \in \mathcal{M}'$, on a
$$
[x, u] \equiv (sg(x))^{\varepsilon_2(u)} \pmod{\Pi_0} \leqno{(5)}
$$
$$
sg : \widehat{\mathfrak{S}} \xrightarrow{can} \mathfrak{S}_3 \xrightarrow{sg} \mu
$$
Cela montre le

\underline{Lemme} Si $x \in \widehat{\mathfrak{S}}^+, u \in \mathcal{M}_0^!$, alors conditions équivalentes :
\begin{enumerate}
\item[a)] $x u x^{-1} \in \mathcal{M}_0^!$
\item[b)] $x u x^{-1} \equiv u \pmod{\Pi_0}$ i.e. $[x, u] \in \Pi_0$
\item[c)] $sg(x)^{\varepsilon_2(u)} = 1$ i.e. $sg(x) = 1$ ou $\varepsilon_2(u) = 1$
\end{enumerate}

\newpage

%% ===== Page 640 =====

Corollaire 1 Soit $x \in \widehat{\mathfrak{S}}^{+\wedge}$. Pour que $x$ normalise $M_0^!$, il faut et il suffit qu'on ait $sg(x) = +1$.

Corollaire 2 Le sous-groupe $M''^!_0 = M''^{(0)} \cdot \Pi_0$ est normalisé par $\widehat{\mathfrak{S}}^{+\wedge}$.

Cela montre en particulier que $\rho$ normalise $M_0^!$, donc son intersection $Z_p^! = Z_p \cap M_0^!$, qui est donc invariant dans $Z_p$. Il en est de même de l'intersection avec $M'(-j) \times \mu$, soit $M'_0(-j)$. (Chacun de ces groupes est égal : $Z_p \cap M''^!_0$, et $M''^!_0$ est normalisé par $\rho$\dots). Ainsi les quatre groupes $(3_1)$ sont distingués dans $Z_p$.

Pour voir s'ils sont distingués dans $N_p^\ast = H_{\tau'} \cdot Z_p$, il suffit de voir s'ils sont normalisés par $\tau'$. C'est clair pour $Z_p^! = Z_p \cap M_0^!$ --- d'où le fait que $N_p^! = H_{\tau'} \cdot Z_p^!$ est un sous-groupe de $N_p^\ast$. Pour $M'_0(-j) = Z_p \cap M''^!_0$, il faudrait donc expliciter l'action de $\tau'$ sur $Z'_p = M'(-j) \times \mu$.

\newpage

%% ===== Page 641 =====

Notons que $\tau \in \mathcal{M}(\omega) \subset \mathcal{M}_0'$ normalise $\mathcal{M}_0^!$, d'où
il opère sur $\mathcal{M}^! = \mathcal{M}_0^! \times \mu$\footnote{$\simeq \Gamma_{\mathbf{Q}}^\sim \cdot \Pi_0$} par l'opération
produit de la conjugaison par $\tau$ dans $\mathcal{M}_0^!$,
et l'identité dans $\mu$. On a d'autre part
$\tau' = \tau \sigma$, on a vu que
$$
\sigma(x) \equiv \varepsilon_2(x) x \pmod{\Pi_0} \quad (x \in \mathcal{M}) \leqno{(40)}
$$
donc
$$
\tau'(x) \equiv \varepsilon_2(x) \tau(x) \pmod{\Pi_0} \quad (x \in \mathcal{M}) \leqno{(41)}
$$
d'où il s'ensuit
$$
\text{si } x \in \mathcal{M}_0^!, \text{ alors } \tau'(x) \equiv x \pmod{\Pi_0}\footnote{$\tau'(x) \equiv x \pmod{\Pi_0} \iff \tau'(x) \in \mathcal{M}_0^!$} \iff \varepsilon_2(x) = 1 \leqno{(42)}
$$
en particulier, $\mathcal{M}_0^{!!}$ est normalisé par $\tau'$.

Donc si on applique ceci aux éléments
de $Z_\rho^! = \mathcal{M}^!(-j)$ dans $Z_\rho^! = \mathcal{M}'(-j) \times \mu$,
on trouve
$$
\text{Soit } x \in Z_\rho^! = \mathcal{M}^!(-j), \text{ alors } \tau'(x) \in Z_\rho^! \iff \varepsilon_2(x) = 1 \leqno{(43)}
$$

Corollaire. $Z_\rho^! = \mathcal{M}^!(-j)$ n'est pas normalisé par $\tau'$,
donc n'est pas invariant dans [unclear: $\hat{\Delta}_\rho^*$],
mais $\mathcal{M}^0(-j) = Z_\rho \cap \mathcal{M}_0^{!!}$ est
normalisé par $\tau'$, donc invariant
par [unclear: pr. direct par] $\mu$.

Donc le seul des quatre sous-groupes
(39) de $Z_\rho$ qui un ne soit distingué dans
$N_\rho$ est $\mathcal{M}'(-j) = Z_\rho^!$. Les trois autres

\newpage

%% ===== Page 642 =====

maintenant, par produit $\frac{1}{2}$ direct avec $\mu$,
des sous-groupes de $N_{\hat{\rho}}$, invariants dans $N_{\hat{\rho}}$.
Le plus petit de tous ces sous-groupes
est $\mathcal{M}^\circ[-j]$, qui est d'indice 24
dans $N_{\hat{\rho}}^*$. Le groupe quotient
$$
N_{\hat{\rho}}^* / \mathcal{M}^\circ[-j]
$$
est donc un groupe d'ordre 24\footnote{=[unclear: H_{24}]} qui
reçoit d'ailleurs le groupe $D_{\hat{\rho}} \subset N_{\hat{\rho}}^*$,
normalisateur de $L_{\hat{\rho}}$ dans $\operatorname{Sl}(2, \mathbf{Z})$,
groupe isomorphe à $D_6$, le dièdre d'ordre
12. Le noyau de
$$
D_{\hat{\rho}} \hookrightarrow N_{\hat{\rho}}^* / \mathcal{M}^\circ[-j] \leqno (44)
$$
est évidemment contenu dans $L_{\hat{\rho}} = D_{\hat{\rho}} \cap Z_{\hat{\rho}}$
et il est réduit à $1$ car
puisqu'on a
$$
L_{\hat{\rho}} \cap Z_{\hat{\rho}}^! \subset L_{\hat{\rho}} \cap \mathcal{M}^\circ = L_{\hat{\rho}} \cap \Pi_\circ = \{1\} .
$$
(d'après 12)
Ainsi $D_{\hat{\rho}}$ apparaît comme un sous-
groupe d'indice 2 de $N_{\hat{\rho}}^* / \mathcal{M}^\circ[-j]$,
tout comme il est sous-groupe d'indice
2 de $\mathfrak{S} / \Pi_\circ$. Il serait tentant

\newpage

%% ===== Page 643 =====

alors d'établir un isomorphisme
$$ N_{\hat{\Gamma}}^* / M^\circ(-J) \overset{?}{\simeq} \hat{\mathbb{S}}^* / \Pi_0 $$
compatible avec un plongement de $D_f$.
Essayons-le aussi :

Proposition Considérons l'hom. composé
(45)
\begin{tikzcd}
M \ar[r, "\theta_M"] & \operatorname{Gl}(2, \hat{\mathbb{Z}}) \ar[r] & \operatorname{Gl}(2, \mathbb{Z}/4\mathbb{Z}) \ar[r, "\text{épi. de}", "\text{degré 4}"'] & \hat{\mathbb{S}}^* / \Pi_0 \\
& & \operatorname{Gl}(2, \mathbb{Z}) / \mathfrak{S}[4] \ar[u, "\simeq"'] \ar[ur, "\simeq \hat{\mathbb{S}} / \hat{\Pi}_0"'] & 
\end{tikzcd}
qui induit sur $\hat{\mathbb{S}}^* \subset M$ l'homomorphisme canonique $\hat{\mathbb{S}}^* \hookrightarrow \hat{\mathbb{S}}^* / \Pi_0$. Alors l'hom. induit
$$ N_{\hat{\Gamma}}^* \longrightarrow \hat{\mathbb{S}} / \Pi_0 $$
a comme noyau $M^\circ(-J)$, et induit par passage au quotient un isomorphisme
(46)
$$ N_{\hat{\Gamma}}^* / M^\circ(-J) \overset{\sim}{\longrightarrow} \hat{\mathbb{S}} / \Pi_0 $$

Démonstration. Les éléments de $M^\circ(-J)$ sont les éléments de la forme
$$ \left\{ v = \alpha^{-1} u \ \left| \begin{array}{l} u \in M(0) \ ,\ \alpha \in \Pi_0 \\ u(p) = \alpha(p) \quad (\text{donc } \varepsilon_3(u) = 1) \\ \varepsilon_2(u) = 1 \end{array} \right. \right\} $$
ils sont en corr. 1-1 avec les éléments de
$$ M \hat{(0)} = \{ u \in M(0) \mid \varepsilon_2(u) = \varepsilon_3(u) = 1 \} \ . \quad \text{L'image de} $$

\newpage

%% ===== Page 644 =====

$v$ dans $\operatorname{Gl}(2,\hat{\mathbb{Z}})/\{\pm 1\} \Pi_0$, / [unclear] dans $\hat{\mathfrak{S}}/\Pi_0$,
est égale à celle de $u$. Donc l'assertion
que $M^\circ(-j)$ est dans le noyau de
$M \to \hat{\mathfrak{S}}/\Pi_0$, équivaut au fait que
$M^\circ(0)$ est dans ce noyau, donc que
l'on a un diagramme commut.

[MARGIN: NB. Ceci achève la démonstration, $\varepsilon_3$ n'intervenant pas dans l'exp. explicite de $M \to \hat{\mathfrak{S}}/\Pi_0$, cf plus bas...]

\begin{tikzcd}
M(0) \ar[r] \ar[d, "(\varepsilon_1,\varepsilon_3)"'] & M \ar[d, "(45)"] \\
\mu \times \mu \ar[r, "?"] & \hat{\mathfrak{S}}/\Pi_0
\end{tikzcd}

Notons que le composé
$$ M \to \hat{\mathfrak{S}}/\Pi_0 \to (\widehat{\mathbb{Z}/4\mathbb{Z}})^* \simeq \{\pm 1\} $$
n'est autre que $\varepsilon_2$. Notons aussi que
par la compatibilité

\begin{tikzcd}
\hat{\mathfrak{S}} \ar[r, hook] \ar[dr, "can"'] & M \ar[d, "(45)"] \\
& \hat{\mathfrak{S}}/\Pi_0
\end{tikzcd}

l'opération de $M$ sur $\hat{\mathfrak{S}}/\Pi_0$ déduite de
l'hom. (45) n'est autre que l'opération
déduite, par des autom. int. par passage au
quotient par $\Pi_0$, opération qui, pour $u \in M$,
est en particulier sur $M_0 \simeq M^\circ$, est
donnée par $\varepsilon_2(u)$ (cf p. 594). Ainsi

\newpage

%% ===== Page 645 =====

On voit que $\mathcal{M}'''_0$ s'envoie dans le centre de $\mathcal{G}/\Pi_0$, qui est contenu dans $\mathcal{G}^+/\Pi_0$, et en fait égal à toute l'image $\mu$ de $\Pi/\Pi_0$ dans $\mathcal{G}^+/\Pi_0$ (puisque le centre de $\mathcal{G}^+/\Pi \simeq \mathfrak{S}_3$ est trivial). Donc on a une factorisation
$$ \mathcal{M}''' \to \mu \hookrightarrow_{\text{centre}} \mathcal{G}/\Pi_0 $$
et il faudrait vérifier que cet hom. $\mathcal{M}''' \to \mu$ n'est autre que $\varepsilon_3$, qui est donc trivial sur $\mathcal{M}''!$ et a fortiori sur $\mathcal{M}''(0)$. En fait le noyau n'est autre que $\operatorname{Ker}(\varepsilon_3|\mathcal{M}''')$ (qui contient $\mathcal{M}''!$)
\marginpar{- [unclear: pas $\mu'_3$] mais $\mathcal{M}'''!$ - on y reviendra plus bas.}
Admettons provisoirement ce fait, on trouve donc une factorisation de
$$ \mathcal{N}_p^* \to \mathcal{G}/\Pi_0 $$
$$ \mathcal{N}_p^* / \mathcal{M}^\circ(-j) \to \mathcal{G}/\Pi_0 \leqno{(*)} $$
et comme l'hom. composé
$$ D_p \hookrightarrow \mathcal{N}_p^* / \mathcal{M}^\circ(-j) \to \mathcal{G}/\Pi_0 $$
est injectif, on voit que le noyau de (*) est un sous-groupe de $\mathcal{N}_p^* / \mathcal{M}^\circ(-j)$ dont l'intersection avec $D_p$ est $\{1\}$ - il est donc $1$, ou d'ordre $2$ (puisque

\newpage

%% ===== Page 646 =====

$D_{\rho}$ est d'indice 2 dans $\mathcal{N}_{\rho}^*/\mathcal{M}^{\circ}(-J)$).

Je vais mieux essayer de [unclear: l'énoncer] :

Proposition - Considérons la filtration
de $\mathfrak{S}^+/\pi_0$ par $\mathfrak{S}^+/\pi \isom \mathfrak{S}_3$, et $\pi/\pi_0 \isom \mu$,
et la filtration de $\mathcal{M}$ par $\mathcal{M}/\mathcal{M}^{\prime\prime} \overset{\varepsilon_2}{\isom} \mu$,
puis $\mathcal{M}^{\prime\prime}/\mathcal{M}^{\prime !} \isom \mathfrak{S}_3$, et $\mathcal{M}^{\prime !}/\mathcal{M}^{\circ !} \overset{\varepsilon_3}{\isom} \mu$.
Alors l'homomorphisme canonique
$$
\mathcal{M} \to \mathfrak{S}^+/\pi_0
$$
est compatible avec les filtrations,
et induit des isomorphismes évidents
sur les facteurs. En particulier,
le noyau de $\mathcal{M} \to \mathfrak{S}^+/\pi_0$ est $\mathcal{M}^{\circ !}$ !

Admettant ce point, on trouve
donc
$$
\operatorname{Ker}(\mathcal{N}_{\rho}^* \to \mathfrak{S}^+/\pi_0) = \mathcal{N}_{\rho}^* \cap \mathcal{M}^{\circ !}
$$
et comme $\mathcal{N}_{\rho}^* \cap \mathcal{M}^{\circ !} = Z_{\rho}^{\circ !} = \mathcal{M}^{\prime}(-J)$,
on trouve que le noyau cherché est
égal : $\mathcal{M}^{\prime}(-J) \cap \mathcal{M}^{\prime\prime} = \mathcal{M}^{\circ}(-J)$, cqfd.
Donc l'homomorphisme en bas est
injectif (*), et comme les groupes

\newpage

%% ===== Page 647 =====

en présence ont même ordre $24$, et l'homomorphisme \\
est bien un isomorphisme, cqfd.

\begin{enumerate}
\item[d)] Application à la structure de $\mathcal{M}_{\rho,\sigma}$ et à \\
la définition d'un homomorphisme $\mathcal{M} \to \mathcal{M}_{\rho,\sigma}/\mu_2$.
\end{enumerate}

Rappelons que
$$
\mathcal{M}_{\rho,\sigma} \defeq \mathcal{N}_\rho^* \times_{\gamma_{\rho,\sigma}} \mathcal{N}_\sigma^* \times_{\mathfrak{S}_{\rho,\sigma}} \mathfrak{S}_{\rho,\sigma} \leqno{(48)}
$$
$$
\begin{cases}
\mathfrak{S}_{\rho,\sigma} \simeq \mathcal{M} = \mathcal{M}_{0,3} \cdot \mathfrak{S}^{+\wedge} \\
\gamma_{\rho,\sigma} \defeq \mathfrak{S}_{\rho,\sigma}/\mathfrak{S}^{+\wedge} \simeq \widetilde{\Gamma}_{\mathbf{Q}} (\simeq \mathcal{M}(0)), \\
\mathcal{N}_\rho^* = \operatorname{Norm}_{\mathcal{M}_\rho^*}(L_\rho), \quad \mathcal{N}_\sigma^* = \operatorname{Norm}_{\mathcal{M}_\sigma^*}(L_\sigma)
\end{cases} \leqno{(49)}
$$

D'autre part, nous avons introduit les \\
sous-groupes remarquables
$$
\begin{cases}
Z_0^! = \mathcal{M}(1/2) \subset \mathcal{N}_\sigma^* \quad (\text{p. 615 (9)}) \\
\mathcal{N}_\rho^? = Z_\rho^! \cdot \mu_2 \subset \mathcal{N}_\sigma^* \quad (\text{p. 616 (21)}).
\end{cases} \leqno{(50)}
$$

Le premier de ces groupes ne contient \\
pas $\mu$, le deuxième le contient
$$
\begin{cases}
Z_0^! \cap \mu = \{1\} \\
\mathcal{N}_\rho^? \supset \mu
\end{cases} \leqno{(51)}
$$

et
$$
\begin{cases}
Z_0^! \isommap \gamma_{\rho,\sigma} = \widetilde{\Gamma}_{\mathbf{Q}} \\
1 \to \mu \to \mathcal{N}_\rho^? \to \widetilde{\Gamma}_{\mathbf{Q}} \to 1
\end{cases} \quad \text{exacte} \leqno{(52)}
$$

\newpage

%% ===== Page 648 =====

Soit alors
$$
\mathcal{M}_{\rho, \sigma}^{\natural} \defeq \mathcal{N}_{\rho}^{?} \times_{\gamma_{\rho, \sigma}} Z_{\sigma}^{!} \times_{\gamma_{\rho, \sigma}} G_{\rho, \sigma} \subset \mathcal{M}_{\rho, \sigma} \leqno{(53)}
$$

On a le sous-groupe central
$$
\mu \times \mu \times \mu \hookrightarrow \mathcal{M}_{\rho, \sigma} \leqno{(54)}
$$
$$
(\varepsilon, \varepsilon', \varepsilon'') \longmapsto (\varepsilon, \varepsilon', \varepsilon'')
$$
provenant des plongements évidents
$$
\mu \hookrightarrow \mathcal{N}_{\rho}^{*}, \quad \mu \hookrightarrow \mathcal{N}_{\sigma}^{*}, \quad \mu \hookrightarrow G_{\rho, \sigma} \quad (\hookrightarrow \mathfrak{S}_{\rho, \sigma}^{+} \simeq \mathfrak{S}^{+})
$$

et on a
$$
\mathcal{M}_{\rho, \sigma}^{\natural} \cap (\mu \times \mu \times \mu) = (\mu \times 1 \times \mu)
$$
d'où
$$
\mu \times \mu \hookrightarrow \mathcal{M}_{\rho, \sigma}^{\natural} \leqno{(55)}
$$

On désigne par $\mu_{\rho}$ le sous-groupe central $\mu$ de $\mathcal{N}_{\rho}^{*}$ et $\mathcal{N}_{\rho}^{?}$, par suite un sous-groupe de $\mathcal{M}_{\rho, \sigma}^{\natural}$. Il est clair alors par (52) que l'on a une suite exacte
$$
1 \to \mu_{\rho} \hookrightarrow \mathcal{M}_{\rho, \sigma}^{\natural} \to \underset{\simeq \mathcal{M}}{G_{\rho, \sigma}} \to 1 \leqno{(56)}
$$
d'où des isomorphismes canoniques

\newpage

%% ===== Page 649 =====

(56)
$$M \overset{\sim}{\longrightarrow} M_{\rho, \sigma}^\natural/\mu_{\rho} \hookrightarrow M_{\rho, \sigma}/\mu_{\rho}$$

[MARGIN: On posera donc
$M_{\rho, \sigma}^+ = M_{\rho, \sigma}^\natural/\mu_{\rho}$]

Evidemment le sous-groupe $M_{\rho, \sigma}^\natural$ de $M_{\rho, \sigma}$
contient $G_{\rho, \sigma}^+ \simeq S G_{\rho, \sigma}$ (qui en fait...
en fait on a
un diagramme de suites exactes

(57)
\begin{tikzcd}
1 \ar[r] & \mu_{\rho} \times G_{\rho, \sigma}^+ \ar[r] \ar[d] & M_{\rho, \sigma}^\natural \ar[r] \ar[d] & \Gamma_{\rho, \sigma}^\natural \ar[r] \ar[d] & 1 \\
1 \ar[r] & \underset{L_\rho}{S\hat{\mathbb{Z}}} \times \underset{L_\sigma}{S\hat{\mathbb{Z}}} \times G_{\rho, \sigma}^+ \ar[r] & M_{\rho, \sigma} \ar[r] & \Gamma_{\rho, \sigma} \ar[r] & 1
\end{tikzcd}

Dans l'isomorphisme (56), le sous-groupe
$G^{+ \wedge}$ de $M$ correspond au sous-groupe
$\mu_\rho G_{\rho, \sigma}^+/\mu_\rho \simeq G_{\rho, \sigma}^+$, il s'envoie isomorphiquement
sur le sous-groupe $G_{\rho, \sigma}^+$ de $M_{\rho, \sigma}/\mu_{\rho}$.
On trouve ainsi

(58)
\begin{tikzcd}
& & & \Gamma_{\rho, \sigma} \ar[d, "\simeq"'] & \\
1 \ar[r] & G^{+ \wedge} \ar[r] \ar[d, "\simeq"'] & M \ar[r] \ar[d, "\simeq"'] & \hat{\Gamma}_Q \ar[r] \ar[d, "\simeq"'] & 1 \\
1 \ar[r] & G_{\rho, \sigma}^+ \ar[r] \ar[d, equal] & M_{\rho, \sigma}^\natural/\mu_\rho \ar[r] \ar[d, hook] & \Gamma_{\rho, \sigma}^\natural/\mu_\rho \ar[r] \ar[d, hook] & 1 \\
1 \ar[r] & G_{\rho, \sigma}^+ \ar[r] & M_{\rho, \sigma}/\mu_\rho \ar[r] & \Gamma_{\rho, \sigma}/\mu_\rho \ar[r] & 1
\end{tikzcd}

Ainsi, se trouve prouvé que l'hom.
canonique épi.

(59)
$$1 \to \frac{S\hat{\mathbb{Z}} \times S\hat{\mathbb{Z}}}{L_\rho \times L_\sigma} \to \Gamma_{\rho, \sigma} \to \bar{\Gamma}_{\rho, \sigma} \to 1$$

admet un "bisecteur" $\Gamma_{\rho, \sigma}^\natural$ (dont
l'intersection avec $\frac{S\hat{\mathbb{Z}} \times S\hat{\mathbb{Z}}}{L_\rho \times L_\sigma}$ est $\mu_\rho \times 1$)

\newpage

%% ===== Page 650 =====

$$
\begin{cases}
\Gamma_{\rho,\sigma} \simeq \Gamma_{\rho,\sigma}^{\natural} \cdot S\Gamma_{\rho,\sigma} \quad , \quad \Gamma_{\rho,\sigma}^{\natural} \cap S\Gamma_{\rho,\sigma} = \mu_{\rho} \\
\quad\quad\quad S\mathbf{Z}_{\rho} \times S\mathbf{Z}_{\sigma} = L_{\rho} \times L_{\sigma} \\
\Gamma_{\rho,\sigma}/\mu_{\rho} \simeq \big( \Gamma_{\rho,\sigma}^{\natural}/\mu_{\rho} \big) \cdot (L_{\rho}/\mu_{\rho} \times L_{\sigma})
\end{cases} \leqno (60)
$$
$$
\Gamma_{\rho,\sigma}^{\natural} \simeq \Gamma_{\mathbf{Q}}^\sim
$$

Pour dire ce que deviennent les sous-groupes remarquables
$$
\begin{array}{cccc}
\mathcal{M}(0), & \mathcal{M}[0], & \mathcal{M}'[0], & \mathcal{M}''[0] \\
\text{"} & \mathcal{M}^\sim[0] & \mathcal{M}'^\sim[0] & \mathcal{M}''^\sim[0] \\
\text{"} & \mathcal{M}^\sim \cdot L_{\rho}^0 & \mathcal{M}^\sim \cdot L_{\sigma}^0 & \mathcal{M}^\sim \cdot L_{\rho,\sigma}^0
\end{array}
$$
$$
\begin{array}{cccc}
\text{[unclear: } \mathcal{M}(i)], & \mathcal{M}'(-j), & \mathcal{M}''(-j), & \mathcal{N}_\rho^? \\
\text{"} \quad \mathcal{Z}_0^\sim ! & & &
\end{array}
$$
de $\mathcal{M}$ par l'homomorphisme $\mathcal{M} \isommap \mathcal{M}_{\rho,\sigma}^\natural / \mu_{\rho} \subset \mathcal{M}_{\rho,\sigma}$, il faudrait d'abord introduire les groupes adéquats dans les $\mathcal{M}_{\rho,\sigma}$ généraux. Nous allons donc revenir sur ces groupes maintenant.

\newpage

%% ===== Page 651 =====

\section*{VI Les groupes $\mathcal{M}_{\rho,\sigma}$ généraux : récapitulation, raffinement.}

$1^{\circ})$ Données. Ce sont celles de V $1^{\circ})$ (p.\ 577).
En plus de $\mathfrak{S}$ (pro-fini) et du $\operatorname{Gl}(2, \mathbf{Z})^{\wedge} \to \mathfrak{S}$ homomorphisme continu surjectif, i.e.\ des générateurs $\rho, \sigma, \tau$ (ou $\rho, \sigma, \tau'$) de $\mathfrak{S}$, on se donne aussi un sous-groupe ouvert
$$
\mathfrak{S}^{\natural} \subset \mathfrak{S}^{+}, \text{ invariant dans } \mathfrak{S}. \leqno{(1)}
$$
Dans le cas universel $\mathfrak{S} = \operatorname{Gl}(2, \mathbf{Z})^{\wedge}$, la donnée de $\mathfrak{S}^{\natural}$ permet correspondre au point de vue où on regarde les sous-groupes $\mathcal{M}^{\natural}$ de $\mathcal{M}$ formés des $u$ qui opèrent trivialement dans $\mathfrak{S}^{+} / \mathfrak{S}^{\natural}$ - c'est un sous-groupe ouvert, qui définit dans $\boldsymbol{\Gamma}_{\mathbf{Q}} \simeq \mathcal{M} / \mathfrak{S}^{+}$ un sous-groupe $\boldsymbol{\Gamma}_{\mathbf{Q}}^{\natural}$ ouvert, dont on peut se proposer d'étudier les représentations dans des groupes pro-finis convenables. Un cas intéressant sera celui où
$$
\mathfrak{S}^{\natural} = \Pi_0, \text{ donc } \mathcal{M}^{\natural} = \mathcal{M}^{\prime \prime},
$$
ou
$$
\mathfrak{S}^{\natural} = \Pi, \text{ donc } \mathcal{M}^{\natural} = \mathcal{M}^{\prime}.
$$
Si
$$
\mathfrak{S}^{\natural} = \mathfrak{S}, \text{ donc } \mathcal{M}^{\natural} = \mathcal{M},
$$
on retrouve les constructions précédentes.

\newpage

%% ===== Page 652 =====

Un cas intéressant aussi est celui où $\mathfrak{G}^\natural$ est le groupe des pts.\ adéliques continus relatifs au groupe sur $\mathbf{Z}$ (dans ce cas, il est vrai que $\rho, \sigma, \tau$ n'engendrent pas [unclear: vraiment] $\mathfrak{G}^\wedge$ que modulo groupe fini i.e.\ ils n'engendrent peut-être seulement un s-groupe ouvert. Dans ce cas, il faudrait plutôt [unclear: repenser] les constructions qui suivent dans le cas schématique j'y viendrai par la suite. Un choix intéressant d'un $\mathfrak{G}^\natural$ est alors le commu-tateur [unclear: neutre] de $\mathfrak{G}$.

J'hésite à introduire des notations $Z_\mathfrak{G}^\natural, M_\mathfrak{G}^\natural$ etc --- pour ne pas surchar-ger des notations déjà lourdes. Donc je préfère utiliser les notations plus simples $Z_{\rho,\sigma}$ etc --- étant entendu que dans ces constructions le groupe $\mathfrak{G}^\natural$ intervient --- quitte à introduire la notation $\natural$ en cas de besoin.
On pose
$$
\mathfrak{G}_0^\natural = \{1, \tau\} \cdot \mathfrak{G}^\natural \subset \mathfrak{G} \leqno{(2)}
$$

\newpage

%% ===== Page 653 =====

2°) $Z_{\rho,\sigma}$, $\gamma_{\rho,\sigma}$.

Soit
$$
Z_{\rho,\sigma} \subset \operatorname{Aut}(\widehat{\mathfrak{S}}^+) \times \hat{\mathbf{Z}}^* \times \mu \times \mu \leqno{(4)}
$$
$$
(u, \mu, \varepsilon_2, \varepsilon_3)\footnote{Il s'ensuit automatiquement en partie dans d) (conjugaison canonique du quotient $\widehat{\mathfrak{S}}^+/\pi_1 \simeq \widehat{\mathfrak{S}}^+_{0,3}$ etc.) et par suite induit par $\varepsilon_3=1$, $t \mapsto \varepsilon_3(t)$ (cf [12])}
\begin{cases}
\text{a) } \exists \alpha \in \widehat{\mathfrak{S}}^+, \text{ avec } \varepsilon_2(\alpha) = \varepsilon_2, u(\rho) = \alpha \rho \\
\text{b) } \exists \beta \in \widehat{\mathfrak{S}}^+, \text{ avec } \varepsilon_3(\beta) = \varepsilon_3, u(\sigma) = \beta \sigma \\
\text{c) } u(\varepsilon_0) = \varepsilon_0^\mu \\
\text{d) } u \text{ normalise } \widehat{\mathfrak{S}}^+ / \operatorname{Im} \pi \text{ et y induit l'identification}\dots \footnote{Cette condition}
\end{cases}
$$

Proposition $Z_{\rho,\sigma}$ est un sous-groupe.

Dém. Les conditions a), d) définissent évidemment un [unclear: sous-groupe] dans ce groupe produit.

La condition b) signifie que
$$
\text{b') } u(\sigma) \text{ est } \widehat{\mathfrak{S}}^+ \text{-conjugué à } \sigma \text{ (resp. } \sigma^{\varepsilon_3} \text{)} \leqno{(5)}
$$
et sous cette forme il est clair que la condition c') définit un sous-groupe. Je ne pense pas que la condition a) définisse à elle seule un sous-groupe. Mais si $\mathcal{G}_a, \mathcal{G}_b$ désignent les sous-groupes du produit défini par ces conditions, alors $\mathcal{G}_a$ ...

Prenons $g = (u, \mu, \varepsilon_2, \varepsilon_3)$ et $g' = (u', \mu', \varepsilon_2', \varepsilon_3')$ dans $\mathcal{G}_a$, avec $\alpha, \alpha' \in \widehat{\mathfrak{S}}^+_1$, de sorte que l'on a ...

\newpage

%% ===== Page 654 =====

$$
u(\rho) = \alpha(\rho) \quad , \quad u'(\rho) = \alpha'(\rho) \quad , \quad \text{avec } \begin{cases} \varepsilon_{\overline{\tau}}(\alpha) = \varepsilon_3 \\ \varepsilon_{\overline{\tau}}(\alpha') = \varepsilon'_3 \end{cases} \leqno (6)
$$

Prouvons que $u'u \in G_{(a,b)}$ -- i.e.\ qu'il satisfait aussi à (b). Si $\varepsilon_3 = 1$ i.e.\ $\alpha \in \widehat{\mathfrak{S}}^+$
$u'(\alpha)$ est défini et on a
$$
(u'u)(\rho) = u'(u(\rho)) = u'(\alpha(\rho)) = u'(\alpha)(u'(\rho)) = u'(\alpha)(\alpha'(\rho)) = (u'(\alpha)\alpha')(\rho) = \alpha''(\rho) \quad , \quad \text{avec}
$$

$$
\alpha'' = u'(\alpha)\alpha' \in \widehat{\mathfrak{S}}_{\overline{\tau}} \qquad (\text{car } \varepsilon_3 = 1) \leqno (7)
$$
$$
\varepsilon_{\overline{\tau}}(\alpha'') = \varepsilon_{\overline{\tau}}(\alpha') = \varepsilon'_3 = \varepsilon_3 \varepsilon'_3
$$

(puisque $u'(\alpha) \in \widehat{\mathfrak{S}}_{\overline{\tau}}$ i.e.\ $\varepsilon_{\overline{\tau}}(u'(\alpha)) = 1$, et $\varepsilon_3 = 1$).
Dans ce cas on a gagné. Si $\varepsilon_3 = -1$, on écrit
$$
\alpha = \alpha_0 \tau \qquad \alpha_0 \in \widehat{\mathfrak{S}}^+
$$
et on aura
$$
u(\rho) = (\alpha_0 \tau)(\rho) = \alpha_0(\underline{\tau(\rho)}) = \alpha_0(\sigma(\rho^{-1})) = (\alpha_0 \sigma)(\rho^{-1}) ,
$$
où maintenant $\alpha_0 \sigma \in \widehat{\mathfrak{S}}^+$, donc $u'(\alpha_0 \sigma)$ est défini, et on aura
$$
(u'u)(\rho) = u'(u(\rho)) = u'(\alpha_0 \sigma) \big(\underbrace{u'(\rho^{-1})}_{\alpha'(\rho^{-1})}\big) = (u'(\alpha_0 \sigma)\alpha')(\rho^{-1})
$$
et comme $\rho^{-1} = \sigma \tau(\rho)$
on trouve
$$
(u'u)(\rho) = \alpha''(\rho) \quad \text{avec } \alpha'' \in \widehat{\mathfrak{S}}_{\overline{\tau}} \leqno (8)
$$
où
$$
\begin{cases}
\alpha'' = \varepsilon u'(\alpha_0) u'(\sigma) \alpha' \sigma \tau \\
\varepsilon \in \mu = \{\pm 1\} \subset \operatorname{Centr} \widehat{\mathfrak{S}}^+
\end{cases}
$$
est choisi de façon que $\alpha''$ (si possible) soit dans $\widehat{\mathfrak{S}}_{\overline{\tau}}$. Notons que [unclear: tous les facteurs du produit sont dans] $\widehat{\mathfrak{S}}^+$, et donc on a $\alpha'' \in \widehat{\mathfrak{S}}$.

\newpage

%% ===== Page 655 =====

Réduisons mod $\mathfrak{S}_{\widetilde{\tau}}$, la relation $\alpha'' \in \mathfrak{S}_{\widetilde{\tau}}$ s'écrit
$$
\dot{\alpha}'' \in \{1,\dot{\tau}\} \quad \text{i.e.} \quad \dot{\varepsilon} u'(\sigma)^\cdot \dot{\alpha}' \dot{\sigma} \dot{\tau} \in \{1,\dot{\tau}\} \quad \text{soit} \quad \dot{\varepsilon} u'(\sigma)^\cdot \dot{\alpha}' \dot{\sigma} \dot{\tau} \in \{1,\dot{\tau}\}
$$
On distingue deux cas :

1) $\varepsilon'_3 = 1$ i.e.\ $\dot{\alpha}' = 1$, la relation s'écrit
$$
\dot{\varepsilon} u'(\sigma)^\cdot \dot{\sigma} = 1 \leqno{(*)}
$$

2) $\varepsilon'_3 = -1$ i.e.\ $\dot{\alpha}' = \dot{\tau}$, comme $\tau \sigma = - \sigma \tau$, la relation s'écrit
$$
- \dot{\varepsilon} u'(\sigma)^\cdot \dot{\sigma} = 1 \leqno{(**)}
$$
donc dans tous les cas, la relation s'écrit
$$
\begin{cases}
u'(\sigma)^\cdot = (\varepsilon \varepsilon'_3)^\cdot \dot{\sigma}^{-1} & \text{soit encore} \\
u'(\sigma)^\cdot = (-\varepsilon \varepsilon'_3)^\cdot \dot{\sigma} &
\end{cases} \leqno{(9)}
$$
On utilise l'hyp.\ b) sur $u'$ qui s'écrit
$$
u'(\sigma)^\cdot = \dot{\varepsilon}'_2 \dot{\sigma}
$$
on voit donc qu'il suffit de choisir
$\varepsilon$ de façon que $-\varepsilon \varepsilon'_3 = \varepsilon'_2$ i.e.
$$
\varepsilon = - \varepsilon'_2 \varepsilon'_3
$$
donc on aura
$$
\alpha'' = - \varepsilon'_2 \varepsilon'_3 u'(\underbrace{\alpha_0 \sigma}_{\alpha \tau \sigma = \alpha \tau' (-\tau')}) \alpha' \sigma \tau
$$
i.e.
$$
\alpha'' = \varepsilon'_2 \varepsilon'_3 u'(\alpha \tau') (\alpha' \tau') \quad | \text{ Cas } \varepsilon_3 = -1 \leqno{(10)}
$$
formule qui a un sens, car $\alpha \tau' \in \mathfrak{S}^+$ donc
$u'(\alpha \tau')$ est défini, et l'argument précédent
prouve que l'élément en question satisfait
$$
\alpha'' \in \mathfrak{S}_{\widetilde{\tau}}, \quad \varepsilon_{\widetilde{\tau}}(\alpha'') = \varepsilon_{\widetilde{\tau}}(\alpha) \varepsilon_{\widetilde{\tau}}(\alpha') \leqno{(11)}
$$
en distinguant les cas [unclear: $\varepsilon_{\widetilde{\tau}}(\alpha)$] de (9).

Même conclusion dans le cas premier où $\varepsilon_{\widetilde{\tau}}(\alpha) = 1$ (7),
ce qui achève la démonstration.

\newpage

%% ===== Page 656 =====

\noindent
Remarques 1) Pour que la condition $b)$ soit
multiplicative, on n'a pas besoin de la condition
$a)$ dans toute sa force, il suffit que $u$ dans $\mathfrak{S}/\mathfrak{S}^+$
normalise l'image de $L_0$, i.e.\ qu'on ait
$$
u(\sigma) \equiv \pm \sigma \pmod{\mathfrak{S}^+}
$$
2) Si $\mathfrak{S}$ est ``sans centre'', l'hom.\ de
projection est injectif
$$
Z_{\bar{\rho}, \sigma} \hookrightarrow \operatorname{Aut}(\mathfrak{S}^+) \quad (cf\ p.\ 578) \leqno{(12)}
$$

Exemples dans les cas universels

\begin{enumerate}
\item[a)] $\mathfrak{S} = \mathfrak{S}^+ = \operatorname{Sl}(2, \hat{\mathbf{Z}})^\wedge$. On a vu la
condition (4) a) est vide, les conditions
$a), b)$ équivalent respectivement au fait que
$u(\rho)$ conjugué à $\rho^{\varepsilon_2}$, $u(\sigma)$ conjugué à $\sigma^{\varepsilon_3}$, i.e.\
on retombe sur la définition (6) p.\ 574, qui
donne \marginpar{NB Appliqué ici au plutôt (?) $Z_{\bar{\rho},\sigma} \subset M_{0,3} \simeq M_{0,3}[0] \simeq \hat{M}_{0,4}(0)$ etc.}
$$
Z_{\bar{\rho}, \sigma} = M(0).L_0^\pm = M^\pm[0] = M[0]. \leqno{(13)}
$$
\marginpar{[unclear: dans ce cas il est immédiat que $Z_{\bar{\rho},\sigma}$ s'identifie à un sous-groupe trivial de $M[0] \simeq M''(0)$]}

\item[b)] $\mathfrak{S} = \Pi = \Pi_0 \times \mu$. On trouve \marginpar{[unclear: explicitement... invariant par conjugaison]}
$$
Z_{\bar{\rho}, \sigma} = M(0).L_0^\Pi = M^\Pi[0] = M'[0] \leqno{(14)}
$$
\marginpar{[unclear: il s'identifie au $p(...) $ de $M(0) \times \Pi$ tiré de $Z_{\bar{\rho},\sigma}$ de $I_{0,4}$]}

\item[c)] $\mathfrak{S} = \Pi_0$. Les conditions $a)$ et $b)$
proviennent du cas précédent, puisque
$\Pi = \Pi_0 \times \mu$ et $\mu$ est central, pour évaluer les
\end{enumerate}

\newpage

%% ===== Page 657 =====

condition d) devient nettement plus stricte, on
trouve pour $Z_{\hat{\mathfrak{S}}, \sigma} \cap \hat{\mathfrak{S}}^+ = \mathbb{L}_0^{\pi_0}$ (d'indice 2 dans
le groupe correspondant $\mathbb{L}_0^{\pi}$ des cas précédents),
et comme l'opération de $\mathcal{M}(0)_{/ \pm 1} \subset \hat{\mathfrak{S}}^+ / \Pi_0$
se fait par le caractère $\varepsilon_2$, on voit que l'on a\footnote{[unclear: Mais en fait on verra d'ailleurs que $\mathcal{M}''[0] = \mathcal{M}[0]_0$, ce qui est rassurant, et avec la déf par (11) (ou (11)') on n'a plus ce pb.]}\footnote{[unclear: il aura vu avant la définition (4) ce que condition d) devient visiblement plus stricte, où il n'y figure que $\hat{\mathfrak{S}}'$ au lieu de $\hat{\mathfrak{S}}^+$ défini par (14) (ce qui me dispensera de la remarque qui suit). Donc dans le cas $\hat{\mathfrak{S}} = \Pi_0$, l'avoir remplacée par $\dots$ Mais je verrai en fait que $\mathcal{M}''[0] = \mathcal{M}[0]_0$, ce qui est rassurant. Ainsi avec cette remarque sur (4) d) devient d'ailleurs inutile.]}
maintenant
$$
Z_{\hat{\mathfrak{S}}, \sigma} = \mathcal{M}''(0) . \mathbb{L}_0^{\pi_0} = \mathcal{M}''[0]_0 \leqno{(15)}
$$
avec
$$
\begin{cases}
\mathcal{M}''(0) = \Ker ( \mathcal{M}(0) \xrightarrow{\varepsilon_2} \mu ) \\
\mathcal{M}''[0]_0 = \Ker ( \mathcal{M}[0]_0 \xrightarrow{\varepsilon_2} \mu ) \\
\mathcal{M}[0]_0 = \mathcal{M}(0) . \mathbb{L}_0^{\pi_0}
\end{cases}
$$

Ces deux exemples suggèrent l'opportunité,
dans le cas général, de disposer soit d'une sous-groupe
isotrope de $\hat{\mathfrak{S}}$, soit qu'on ait
$\hat{\mathfrak{S}}^0 \subset \hat{\mathfrak{S}}$, tel que l'on ait
$$
\hat{\mathfrak{S}} = \hat{\mathfrak{S}}^0 \times \mu \leqno{(16)}
$$
(où pour le
sous-groupe
central canonique
$\mu = \{1, \rho_{S^2}^2\} = \{1, \sigma^2\}$)
de définir $Z_{\hat{\mathfrak{S}}, \sigma}$ en termes de $\hat{\mathfrak{S}}^0$, et des
choix des $\alpha, \beta$ de la définition (4)
dans $\hat{\mathfrak{S}}^0$, d'où
$$
\hat{\mathfrak{S}}^\Delta_{\sigma} \defeq \mu_{\tau} . \hat{\mathfrak{S}}^0_\sigma . \leqno{(17)}
$$

Remarque. On est tenté de remplacer la
déf(inition) (4) (pour a), b)) par la définition d'apparence plus
simple de la p. 578, qui rend d'ailleurs
inutile le
choix \dots Mais elle n'est pas compatible
avec (le passage à $\hat{\mathfrak{S}}^+ / \Pi_0$) \dots

\newpage

%% ===== Page 658 =====

qui en se restreignant : dans $(u, \mu, \varepsilon_1, \varepsilon_3)$ ou $\varepsilon_3 = 1$,
ce $j^{-1}$ n'est pas égal à $j$ ! Il est vrai que
dans le cas universel p.ex.) [unclear: que] $u(j)$ est
conjugué dans $\underline{G}^+$ à $j^{\varepsilon_3}$, mais en général
il ne le sera pas par un élément (disons)
de $\pi$.

[MARGIN: NB On a
$L_0^{\underline{G}} \supset L_0 \cap \underline{G}^+$ et l'
inclusion peut être stricte,
(exemple p.ex. f)
ci-dessus... $\underline{G}^+ = L_0^{\pi_0}$
et $L_0^+ = L_0^\pi, L_0 \cap \underline{G}^+ = L_0^{\pi_0}$)
ce qui
s'y [unclear]]

On pose
(18) $L_0^{\underline{G}} = \{ l \in \underline{L}_0^{\underline{G}} | (int(l), 1, 1, 1) \in Z_{\rho,\sigma} \}$
(i.e. $L_0^{\underline{G}} \subset \underline{G}^+$ est l'ens
des $\varepsilon_0^n, n \in \hat{\mathbb{Z}}$)
et on a un hom. canonique

(19) $L_0^{\underline{G}} \longrightarrow Z_{\rho,\sigma}$
$$l \longmapsto (int(l), 1, 1, 1)$$

in-jectif ss.

(20) $L_0^{\underline{G}} \cap Centr(\underline{G}^+) = \{ 1 \}$

l'image invariante. On pose

(21) $\Gamma_{\rho,\sigma} = Z_{\rho,\sigma} / L_0^{\underline{G}}$ ,

d'où

(22) $$1 \longrightarrow L_0^{\underline{G}} \longrightarrow Z_{\rho,\sigma} \overset{\delta_Z}{\longrightarrow} \Gamma_{\rho,\sigma} \longrightarrow 1$$ ,

(23) 
\begin{tikzcd}
Z_{\rho,\sigma} \ar[r] & \Gamma_{\rho,\sigma} \ar[r, "\chi_\Gamma"] \ar[dr, "\varepsilon_\Gamma"'] & \hat{\mathbb{Z}}^\times \\
& & \mu
\end{tikzcd}

déduit de $(u, \mu, \varepsilon_1, \varepsilon_3) \longmapsto (\varepsilon_3 \chi, \mu)$ .

Remarque. On désignera par
(24) $\Gamma_{\rho,\sigma}' = \operatorname{Ker}(\varepsilon_\Gamma), \Gamma_{\rho,\sigma}'' = \operatorname{Ker}(\chi_\Gamma), \Gamma_{\rho,\sigma}''' = \operatorname{Ker}(\varepsilon_\Gamma, \chi_\Gamma)$
$Z_{\rho,\sigma}' = \dots, Z_{\rho,\sigma}'' = \dots, Z_{\rho,\sigma}''' = \dots$ (l'image inverse dans $Z_{\rho,\sigma}$)

\newpage

%% ===== Page 659 =====

On pose
$$
\begin{cases}
\gamma'_{p,\sigma} = \Ker \varepsilon_3 \gamma \ , \ \gamma''_{p,\sigma} = \Ker \varepsilon_1 \gamma \ , \ \gamma'''_{p,\sigma} = \Ker \varepsilon_1 \varepsilon_3 \gamma \ , \ \gamma^0_{p,\sigma} = \gamma' \cap \gamma'' = \gamma' \cap \gamma''' \\
\hspace{9cm} = \gamma'' \cap \gamma' \\
Z'_{p,\sigma}, Z''_{p,\sigma}, Z'''_{p,\sigma}, Z^0_{p,\sigma} \text{ sont les images inv. dans } Z_{p,\sigma} \text{ de } \gamma' \text{ etc.} \\
\text{dans } Z_{p,\sigma}
\end{cases} \leqno (24)
$$

Ce sont donc des sous-groupes d'indice $2$ ou quatre de $\gamma_{p,\sigma}$ et $Z_{p,\sigma}$ - l'indice $4$ ne [unclear: pouvant] se présenter que pour $\gamma^0_{p,\sigma}$, $Z^0_{p,\sigma}$.

Remarques 1) On a donc\marginpar{Plus généralement $Z''_{p,\sigma} = \left\{ (u,\mu) \in \operatorname{Aut}(12^+) \times \hat{\mathbf{Z}}^* \mid \dots \right\}$ a) $u(p) \sim p$, b) $u(\sigma) \sim \sigma$, c) $u(\varepsilon_0) = \varepsilon_0 \mu$}
$$
Z'''_{p,\sigma} = \left\{ (u,\mu) \in \operatorname{Aut}(\mathfrak{S}^+) \times \hat{\mathbf{Z}}^* \left| \begin{matrix} u(p) \text{ est } \mathfrak{S}^\natural \text{-conj. à } p \\ u(\sigma) \text{ } \dots \\ u(\varepsilon_0) = \varepsilon_0 \mu \end{matrix} \right. \right\} \leqno (25_{p,\sigma})
$$
(homomorphisme de sous-groupes). La condition d) de (4) - stabilité de $\mathfrak{S}^\natural$ par $u$ - est conséquence ici de a) b) c) , car on aura
$$
u(p) \equiv p \pmod{\mathfrak{S}^\natural}
$$
$$
u(\sigma) \equiv \sigma \pmod{\mathfrak{S}^\natural}
$$
ce qui implique\footnote{C'est [unclear: mauvais] - le sens est dans (4) d) ce qui est plus fort que (26). Ailleurs j'ai pris pour $\mathfrak{S}^\natural$ [unclear: p. 30!]}
$$
u(x) \equiv x \pmod{\mathfrak{S}^\natural} \quad \forall x \in \mathfrak{S}^+ \leqno (26bis)
$$
i.e.\ la stabilité de $\mathfrak{S}^\natural$ et son action triviale sur $\mathfrak{S}^+/\mathfrak{S}^\natural$.

2) On a encore l'élément remarquable $\tau_Z \in Z_{p,\sigma}$, $\tau_Z = (\tau, -1, -1, -1)$.

3) $\mathcal{G}_{p,\sigma}, G_{p,\sigma}, \mathcal{S}_0 \mathcal{G}_{p,\sigma}, \mathcal{G}^\natural_{p,\sigma}$.

Procédant comme page 577, on trouve
$$
Z_{p,\sigma} \to \operatorname{Aut}(\mathfrak{S}^+)
$$
d'où des produits semi-directs
$$
\mathcal{S}_0 \mathcal{G}_{p,\sigma} \defeq Z_{p,\sigma} \cdot \mathfrak{S}^+ \supset \mathcal{S}_0 \mathcal{G}^\natural_{p,\sigma} = Z_{p,\sigma} \cdot \mathfrak{S}^\natural \leqno (27)
$$
et un hom.
$$
\begin{tikzcd}
\mathcal{L}^\natural_0 \arrow[r, hook, "i_{p,\sigma}"] & \mathcal{S}_0 \mathcal{G}^\natural_{p,\sigma} \subset \mathcal{S}_0 \mathcal{G}_{p,\sigma} \\
l \arrow[r, maps to] & (\operatorname{int}(l), l^{-1})
\end{tikzcd} \leqno (28)
$$

\newpage

%% ===== Page 660 =====

d'image invariante, ce qui permet de poser

$$ (29) \quad G_{\rho, \sigma} \overset{\text{déf}}{=} S_0 G_{\rho, \sigma} / \operatorname{Im} L_0^\sigma \supset G_{\rho, \sigma}^\sigma = S_0 G_{\rho, \sigma}^\sigma / \operatorname{Im} L_0^\sigma $$
$$ Z_{\rho, \sigma} \cdot C^+ $$

d'où suites exactes

(30)
\begin{tikzcd}
1 \ar[r] & L_0^\sigma \ar[r] \ar[d, equal] & S_0 G_{\rho, \sigma}^\sigma \ar[r] \ar[d, hook] & G_{\rho, \sigma}^\sigma \ar[r] \ar[d, hook] & 1 \\
1 \ar[r] & L_0^\sigma \ar[r] & S_0 G_{\rho, \sigma} \ar[r] & G_{\rho, \sigma} \ar[r] & 1
\end{tikzcd}

et un diag. commutatif

(31)
\begin{tikzcd}
& Z_{\rho, \sigma} \ar[dr, "i_Z"] \\
L_0^\sigma \ar[ur, "i_{L_0, Z}"] \ar[dr, "i_{L_0, C^\sigma}"'] & & G_{\rho, \sigma}^\sigma \ar[r, hook] & G_{\rho, \sigma} \\
& C^\sigma \ar[ur, "i_{C^\sigma}"] \ar[rr] & & C^+ \ar[u, "i_{C^+}"']
\end{tikzcd}

s'insérant dans des diag. commutatifs et
suites exactes

(32)
\begin{tikzcd}
1 \ar[r] & L_0^\sigma \ar[r] \ar[d] & Z_{\rho, \sigma} \ar[r, "\delta_Z"] \ar[d] & \gamma_{\rho, \sigma} \ar[r] \ar[d, equal] & 1 \\
1 \ar[r] & C^\sigma \ar[r] \ar[d] & G_{\rho, \sigma}^\sigma \ar[r, "\delta_G"] \ar[d] & \gamma_{\rho, \sigma} \ar[r] \ar[d, equal] & 1 \\
1 \ar[r] & C^+ \ar[r] & G_{\rho, \sigma} \ar[r] & \gamma_{\rho, \sigma} \ar[r] & 1
\end{tikzcd}

On considère les opérations

(33)
\begin{tikzcd}
G_{\rho, \sigma} \ar[r, "\delta_G"] \ar[dr, "\chi_G"'] & \gamma_{\rho, \sigma} \ar[r, "\varepsilon_\gamma"] \ar[d, "\chi_\gamma"] & \{ \pm 1 \} \\
& \hat{\mathbb{Z}}^* \ar[ur, "\varepsilon"'] &
\end{tikzcd}

avec
$$ \chi_G = \chi_\gamma \delta_G, \quad \varepsilon_G = \varepsilon_\gamma \delta_G, \quad \varepsilon_\gamma = \varepsilon \chi_\gamma $$

On a maintenant

[MARGIN: idem pour $G_{\rho, \sigma}^+$...]
(34)
$$ G_{\rho, \sigma}', G_{\rho, \sigma}'', G_{\rho, \sigma}''', G_{\rho, \sigma}^\circ \quad \text{comme dans } (24) $$
$$ G_{\rho, \sigma}^+ = \operatorname{Ker} \chi_G, \ G_{\rho, \sigma}'^+, G_{\rho, \sigma}''^+, G_{\rho, \sigma}'''^+, G_{\rho, \sigma}^{\circ +} $$
$$ (\gamma_{\rho, \sigma}^+ = \operatorname{Ker} \chi_\gamma \text{ etc.} \dots) $$

Si
$$ (35) \quad \tau_G \in G_{\rho, \sigma}^\sigma, \tau_\gamma \in \gamma_{\rho, \sigma} $$

\newpage

%% ===== Page 661 =====

désignant les images de $\tau_Z$ (26) dans $G_{\hat{p},\sigma}, \gamma_{\hat{p},\sigma}$,

-- pour faire l'identification
$$
\begin{cases}
\mathfrak{S}^\wedge \simeq \text{image inverse de } \{1, \tau_Z\} \subset \gamma_{\hat{p},\sigma} \text{ par } G_{\hat{p},\sigma} \to \gamma_{\hat{p},\sigma} \\
\mathfrak{S}^\wedge_+ \simeq \text{image inverse de } \{1, \tau_Z\} \text{ dans } G^+_{\hat{p},\sigma} ,
\end{cases}
\leqno (36)
$$
et
$$
\varepsilon_Z = \chi_G | \mathfrak{S}^\wedge = \varepsilon_{Z\sigma} | \mathfrak{S}^\wedge = \varepsilon_\sigma | \mathfrak{S}^\wedge \quad \begin{cases} \text{trivial sur } \mathfrak{S}^\wedge_+ \\ -1 \text{ sur } \tau_Z \end{cases}
\leqno (37)
$$

Considérons l'application canonique
$$
G_{\hat{p},\sigma} \to \operatorname{Aut}(\mathfrak{S}^\wedge_+) \times \hat{\mathbf{Z}}^* \times \mu \times \mu
\leqno (38)
$$
$$
u \longmapsto ((\operatorname{int} u) | \mathfrak{S}^\wedge_+, \chi(u), \varepsilon_Z(u), \varepsilon_\sigma(u))
$$

Proposition. C'est un homomorphisme de groupes dont l'image de $G_{\hat{p},\sigma}$ est formée des
$$
(u, \mu, \varepsilon_2, \varepsilon_3) \in \operatorname{Aut}(\mathfrak{S}^\wedge_+) \times \hat{\mathbf{Z}}^* \times \mu \times \mu
$$
satisfaisant les conditions
$$
\begin{cases}
\text{-) } u(x) \text{ est } \mathfrak{S}^\wedge_+\text{-conjugué à } x^\mu \\
\text{b) } \text{[unclear: } u(\sigma) \text{ est } \mathfrak{S}^\wedge_+\text{-conjugué à } \sigma^{\varepsilon_2(\tau)\sigma} \text{]} \\
\text{c) } u(\varepsilon_0) \text{ est } \mathfrak{S}^\wedge_+\text{-conjugué : [unclear: s.k.] (si } \varepsilon_2=\varepsilon_3) \\
\text{a) } u \text{ stabilise } L^\wedge_0 \text{ [unclear: (automorphisme } \dots )]
\end{cases}
\leqno (39)
$$

Son noyau est
$$
\Ker = \text{[unclear: } \{ f \in G_{\hat{p},\sigma} \mid u = (1,1,1) \in \dots, u = \operatorname{int}(f), f \in \mathfrak{S}^\wedge_+ \} \text{]}
\leqno (40)
$$

Corollaire.
$$
Z_0 = \operatorname{Centr}_{G_{\hat{p},\sigma}}(L_0) , \quad Z^\wedge_0 = Z_0 \cap \mathfrak{S}^\wedge_+
\leqno (41)
$$
\marginpar{$(= \{ f \in \mathfrak{S}^\wedge_+ \mid \operatorname{int}(f) = 1 \} )$}
de sorte que $L^\wedge_0$ s'identifie à un sous-groupe

\newpage

%% ===== Page 662 =====

central de $Z_0^\mathfrak{S}$. Considérons l'application
$$
\begin{aligned}
Z_0^\mathfrak{S} &\longrightarrow S_0 G_{\rho,\sigma} = Z_{\rho,\sigma} \cdot \underline{\mathfrak{S}} \\
f &\longmapsto f^{-1} \cdot \underline{u}(f) \in Z_{\rho,\sigma} \cdot \underline{\mathfrak{S}} \quad \text{où } \underline{u}(f) = (\operatorname{int}(f), 1, 1, 1) \in \widetilde{G}_{\rho,\sigma}^\mathfrak{S}
\end{aligned} \leqno{(42)}
$$

Cette application est un homomorphisme de groupes, induit sur $L_0^\mathfrak{S} \subset Z_0^\mathfrak{S}$ l'homomorphisme canonique $i_{L_0 \mathfrak{S}}$ (28), et $L_0^\mathfrak{S} \subset Z_0^\mathfrak{S}$ est l'image inverse de $i_{L_0 \mathfrak{S}}(L_0^\mathfrak{S}) = \operatorname{Ker}(S_0 G_{\rho,\sigma} \longrightarrow G_{\rho,\sigma}^\mathfrak{S})$. On trouve ainsi un homomorphisme injectif
$$
Z_0^\mathfrak{S} / L_0^\mathfrak{S} \hookrightarrow G_{\rho,\sigma}^\mathfrak{S} \leqno{(43)}
$$
dont l'image est le noyau de (38), i.e.\ on a une suite exacte canonique
$$
\begin{tikzcd}
1 \arrow[r] & L_0^\mathfrak{S} \arrow[r] & Z_0^\mathfrak{S} \arrow[r] & G_{\rho,\sigma}^\mathfrak{S} \arrow[r] & \widetilde{G}_{\rho,\sigma}^\mathfrak{S} \arrow[r] \arrow[d, phantom, "\cap" description] & 1 \\
& & & & \operatorname{Aut}(\underline{\mathfrak{S}}^+) \times \hat{\mathbf{Z}}^* \times \mu \times \mu &
\end{tikzcd} \leqno{(44)}
$$
où $\widetilde{G}_{\rho,\sigma}^\mathfrak{S}$ est formé des $(u, \mu, \varepsilon_2, \varepsilon_3)$ satisfaisant les conditions (39).

Remarque. Cette proposition fait apparaître de quelle façon la \underline{construction} \underline{faible} de $G_{\rho,\sigma}^\mathfrak{S}$ est plus fine que celle en le définissant simplement comme un groupe d'automorphismes de $\underline{\mathfrak{S}}^+$ (sans avoir besoin des $\mu, \varepsilon, \varepsilon'$ par-dessus le marché...).

\newpage

%% ===== Page 663 =====

On [unclear: va considérer] l'homomorphisme (en fait effectivement surjectif par (38), même dans le cas pur dont les digressions telles que $\widetilde{\mathfrak{S}}^+ = \operatorname{Sl}(2, \hat{\mathbf{Z}})$)
$$
\begin{tikzcd}[row sep=tiny]
G_{\rho,\sigma}^\vartriangleright \arrow[r] & \operatorname{Aut}(\widetilde{\mathfrak{S}}^+) \times \hat{\mathbf{Z}}^* \times \mu \\
u \arrow[r, |->] & (\operatorname{int}(u)|\widetilde{\mathfrak{S}}^+, \varepsilon_3(u), \varepsilon_1(u))
\end{tikzcd} \leqno{(45)}
$$
induit par (38), dont l'image est formée des systèmes $(u, \varepsilon_3, \varepsilon)$ tels que
$$
\begin{array}{l}
\text{a) } u(\rho) \text{ est } \widetilde{\mathfrak{S}}^+\text{-conjugué : } \varepsilon_3[\Gamma](\rho) \\
\text{b) } u(\sigma) \text{ est } \widetilde{\mathfrak{S}}^+\text{-conjugué : } \varepsilon_3[\Gamma](\sigma) \\
\text{c) } u(L_0) \text{ est } \widetilde{\mathfrak{S}}^+\text{-conjugué : } L_0 \\
\text{d) } u \text{ stabilise } \widetilde{\mathfrak{S}}^+ \text{ (automatique si } \varepsilon_3 = \varepsilon_2 \text{, [unclear: on l'a vu])}
\end{array} \leqno{(46)}
$$

Son noyau est
$$
\left\{ f \in G_{\rho,\sigma}^\vartriangleright \;\middle|\; y = (u, h, 1, 1) \in Z_{\rho,\sigma} \text{ t.q. } u = \operatorname{int}(f) \right\} \leqno{(47)}
$$

Pour apprécier la structure du noyau, soit
$$
SN_0^\vartriangleright = \operatorname{Norm}_{G_{\rho,\sigma}^\vartriangleright}(L_0), \quad SN^\vartriangleright = SN_0 \cap \widetilde{\mathfrak{S}}^+ \leqno{(48)}
$$
$$
= \left\{ f \in \widetilde{\mathfrak{S}}^+ \;\middle|\; \exists h \in \hat{\mathbf{Z}}^* \text{ t.q. } (\operatorname{int}(f), h, 1, 1) \in Z_{\rho,\sigma} \right\}
$$

et soit
$$
\begin{array}{c}
\widetilde{SN}_0 \subset SN_0 \times \hat{\mathbf{Z}}^* \\
\| \\
\{ (u, \mu) \mid u(\xi_0) = \xi_0^\mu \}
\end{array}
\qquad , \qquad
\widetilde{SN}_0^\vartriangleright = \widetilde{SN}_0 \cap (N_0^\vartriangleright \times \hat{\mathbf{Z}}^*) \leqno{(49)}
$$

\marginpar{idem pour $\widetilde{SN}$.}
On a des [unclear: flèches] :
$$
\begin{tikzcd}[row sep=tiny]
\widetilde{SN}_0^\vartriangleright \arrow[r, "\sim"] & SN_0^\vartriangleright \\
(u, \mu) \arrow[r, |->] & u
\end{tikzcd}
\quad \text{et} \quad
\begin{tikzcd}[row sep=tiny]
\hat{\mathbf{Z}} \arrow[r, "\sim"] & L_0^\vartriangleright \\
n \arrow[r, |->] & \xi_0^n
\end{tikzcd}
\text{ est injectif.} \leqno{(50)}
$$

On a un homomorphisme canonique injectif
$$
\begin{tikzcd}[row sep=tiny]
\widetilde{SN}_0^\vartriangleright \arrow[r] & S_{\rho,\sigma} G_{\rho,\sigma}^\vartriangleright = Z_{\rho,\sigma} \cdot \widetilde{\mathfrak{S}}^+ \\
(f, \mu) \arrow[r, |->] & f^{-1}(\operatorname{int}(f), \mu, 1, 1) = u \cdot f^{-1}
\end{tikzcd} \leqno{(51)}
$$

\newpage

%% ===== Page 664 =====

qui induit sur $L_0^\mathcal{G}$ l'hom canonique $\nu^{(28)}$ (et
qui prolonge (42) ) - il est encore vrai
que l'image inverse de $L_0^\mathcal{G}$ (i.e. de $i_{L_0,\mathcal{G}}(L_0^\mathcal{G})$)
est $L_0^\mathcal{G} \subset \widetilde{SN}_0^\mathcal{G}$, et on trouve ainsi une
suite exacte canonique

(52)
\begin{tikzcd}
1 \ar[r] & L_0^\mathcal{G} \ar[r] & \widetilde{SN}_0^\mathcal{G} \ar[r] & G_{\rho,\sigma}^\mathcal{G} \ar[r] & \tilde{G}_{\rho,\sigma}^\mathcal{G} \ar[r] \ar[d, hook] & 1 \\
& & & & \operatorname{Aut}(\underline{C}^+) \times \mu \times \mu
\end{tikzcd}

[unclear: où] $\tilde{G}_{\rho,\sigma}^\mathcal{G}$ est formé des $(u,\varepsilon_2,\varepsilon_3)$ satis-
faisant (46) - i.e. ...

(52 bis) $$\widetilde{SN}_0^\mathcal{G} / L_0^\mathcal{G} \simeq \operatorname{Ker}(G_{\rho,\sigma}^\mathcal{G} \to \operatorname{Aut}(\underline{C}^+) \times \mu \times \mu)$$ .

\newpage

%% ===== Page 665 =====

4$^\circ$) $N_\rho, N_\sigma, Z_\rho, Z_\sigma, N_\rho^S, N_\sigma^S, Z_\rho^S, Z_\sigma^S$

On pose
$$
\begin{cases}
N_\rho = \{ a \in G_{\rho, \sigma} \mid a(\rho) = \rho^{\varepsilon_3}, \text{où } \varepsilon_3 = \varepsilon_3(a) \} \\
N_\sigma = \{ b \in G_{\rho, \sigma} \mid b(\sigma) = \sigma^{\varepsilon_2}, \text{où } \varepsilon_2 = \varepsilon_2(b) \}
\end{cases} \leqno (53)
$$

NB Il n'y a pas lieu [unclear: pour nous] de [unclear: donner] des définitions de $N_\rho, N_\sigma$ qui donneraient des groupes plus gros\dots (Cf p 580, 581 pour des [unclear: ratiocinations] à ce sujet\dots)

On pose
$$
\begin{cases}
Z_\rho = \operatorname{Cent}_{G_{\rho, \sigma}}(\rho) = \operatorname{Cent}_{G_{\rho, \sigma}}(U_\rho) = \Ker (N_\rho \xrightarrow{\varepsilon_3} \mu) \\
Z_\sigma = \operatorname{Cent}_{G_{\rho, \sigma}}(\sigma) = \operatorname{Cent}_{G_{\rho, \sigma}}(U_\sigma) = \Ker (N_\sigma \xrightarrow{\varepsilon_2} \mu)
\end{cases} \leqno (54)
$$

\marginpar{On a défini de façon symétrique $S N_\rho, S N_\sigma$, mais ici on a $S N_\rho = S Z_\rho$, $S N_\sigma = S Z_\sigma$}
$$
\begin{cases}
S Z_\rho = Z_\rho \cap \mathfrak{S}^+ = \Ker (Z_\rho \xrightarrow{\delta_2} \gamma_{\rho, \sigma}) \\
\qquad = \operatorname{Cent}_{\mathfrak{S}^+}(\rho) = \operatorname{Cent}_{\mathfrak{S}^+}(U_\rho) \\
S Z_\sigma = Z_\sigma \cap \mathfrak{S}^+ = \Ker (Z_\sigma \xrightarrow{\delta_3} \gamma_{\rho, \sigma}) \\
\qquad = \operatorname{Cent}_{\mathfrak{S}^+}(\sigma) = \operatorname{Cent}_{\mathfrak{S}^+}(U_\sigma)
\end{cases} \leqno (55)
$$

\marginpar{suites exactes [unclear: d'hom.] scindées par [unclear: $\mu_2$] - la deuxième scindée par $\mu_2$}
$$
\begin{matrix}
1 \to Z_\rho \to N_\rho \xrightarrow{\varepsilon_3} \mu \to 1 \\
1 \to Z_\sigma \to N_\sigma \xrightarrow{\varepsilon_2} \mu \to 1
\end{matrix} \leqno (56)
$$

Le fait qu'on a pu [unclear: mettre] des $1 \to \dots \to 1$ provient du fait que $\tau' \in N_\rho \cap N_\sigma$ satisfait
$$
\varepsilon_2(\tau') = \varepsilon_3(\tau') = \varepsilon(\tau') = -1 \leqno (57)
$$

On a un [unclear: isomorphisme] de groupes\dots (p 580 !)

\newpage

%% ===== Page 666 =====

\begin{tikzcd}
& D_\rho \ar[rr] & & \underset{= \mu_{\tau'} \cdot S Z_\rho}{N_\rho \cap \widehat{G}} \ar[rr, hook] \ar[dr, hook] & & N_\rho \ar[dd, hook] \\
(58) \quad D_{\tau'} = \mu_{\tau'} \times \mu & & & & \widehat{G} \ar[rr, hook] & & G_{\rho,\sigma} \\
& D_\sigma \ar[rr] & & \underset{= \mu_{\tau'} \cdot S Z_\sigma}{N_\sigma \cap \widehat{G}} \ar[rr, hook] \ar[ur, hook] & & N_\sigma \ar[uu, hook]
\end{tikzcd}

ou mieux

(59)
\begin{tikzcd}[row sep=1.5em, column sep=1.5em]
& & & D_\rho \ar[rr] & & S N_\rho \ar[drr, hook] \ar[rrr, hook] & & & N_\rho \ar[drr, hook] \\
\mu_{\tau'} \ar[rr] & & D_{\tau'} \ar[ur, "2"] \ar[dr, "3"'] & & & & & \widehat{G} \ar[rrr, hook] & & & G_{\rho,\sigma} \\
& & & D_\sigma \ar[rr] & & S N_\sigma \ar[urr, hook] \ar[rrr, hook] & & & N_\sigma \ar[urr, hook] \\
& & & L_\rho \ar[uuu, hook] \ar[rr] & & S Z_\rho \ar[uuu, hook] \ar[drr, hook] \ar[rrr, hook] & & & Z_\rho \ar[uuu, hook] \ar[uurr, hook, crossing over] \ar[drr, dashed] \\
1 \ar[uuuu] \ar[rr] & & \mu \ar[uuuu, hook] \ar[ur, "2"] \ar[dr, "3"'] & & & & & G^+ \ar[uuuu, hook] \ar[uuurrr, hook, crossing over] \ar[rrr, dashed] & & & G_{\rho,\sigma} \ar[uuu, hook] \\
& & & L_\sigma \ar[uuu, hook] \ar[rr] & & S Z_\sigma \ar[uuu, hook] \ar[urr, hook] \ar[rrr, hook] & & & Z_\sigma \ar[uuu, hook] \ar[uuuurr, hook, crossing over] \ar[urr, dashed]
\end{tikzcd}

On définit

\newpage

%% ===== Page 667 =====

(62) 
$$
\left\{
\begin{tikzcd}
1 \ar[r] & SZ_\rho^{\mathcal{G}} \ar[r] \ar[d] & N_\rho^{\mathcal{G}} \ar[r, "\delta_\rho^{\mathcal{G}}"] \ar[d] & \gamma_{\rho,\sigma} \ar[r] \ar[d, equal] & 1 \\
1 \ar[r] & SZ_\rho \ar[r] & N_\rho \ar[r, "\delta_\rho"'] & \gamma_{\rho,\sigma} \ar[r] & 1 \\
1 \ar[r] & SZ_\sigma^{\mathcal{G}} \ar[r] \ar[d] & N_\sigma^{\mathcal{G}} \ar[r, "\delta_\sigma^{\mathcal{G}}"] \ar[d] & \gamma_{\rho,\sigma} \ar[r] \ar[d, equal] & 1 \\
1 \ar[r] & SZ_\sigma \ar[r] & N_\sigma \ar[r, "\delta_\sigma"'] & \gamma_{\rho,\sigma} \ar[r] & 1
\end{tikzcd}
\right.
$$

(63) 
$$
\left\{
\begin{tikzcd}
1 \ar[r] & Z_\rho^{\mathcal{G}} \ar[r] \ar[d, "\text{épi}"'] & N_\rho^{\mathcal{G}} \ar[r, "\varepsilon_{3\rho}^{\mathcal{G}}"] \ar[d, "\text{épi}"'] & \mu \ar[r] \ar[d, equal] & 1 \\
1 \ar[r] & \gamma_{\rho,\sigma}^{\prime} \ar[r] & \gamma_{\rho,\sigma} \ar[r, "\varepsilon_{3\rho}"'] & \mu \ar[r] & 1 \\
1 \ar[r] & Z_\sigma^{\mathcal{G}} \ar[r] \ar[d, "\text{épi}"'] & N_\sigma^{\mathcal{G}} \ar[r, "\varepsilon_{3\sigma}^{\mathcal{G}}"] \ar[d, "\text{épi}"'] & \mu \ar[r] \ar[d, equal] & 1 \\
1 \ar[r] & \gamma_{\rho,\sigma}^{\prime\prime} \ar[r] & \gamma_{\rho,\sigma} \ar[r, "\varepsilon_{3\sigma}"'] & \mu \ar[r] & 1
\end{tikzcd}
\right.
$$

Ces diagrammes [unclear: s'intègrent] [unclear: exact] [unclear: dans] ( cf. p. [unclear] ) .
On a de même un diagramme

(64) 
\begin{tikzcd}
& L_0^\natural = S Z_{\rho,\sigma}^C \ar[rr] \ar[dl] \ar[d] & & Z_{\rho,\sigma}^C \ar[dl] \ar[d] \ar[dr, "\varepsilon_{3\gamma}"] \\
S Z_{\rho}^C \ar[d] \ar[rr] & & N_{\rho}^C \ar[d] \ar[rr] & & \mu \ar[r] & \hat{\mathbb{Z}}^* \ar[dl, "\mu"'] \\
S Z_{\sigma}^C \ar[d] \ar[rr] & & N_{\sigma}^C \ar[d] \ar[rru] \\
1 \ar[r] & C^\natural \ar[rr] \ar[d] & & G_{\rho,\sigma} \ar[rr] \ar[d] & & \gamma_{\rho,\sigma} \ar[r] \ar[d] & 1 \\
1 \ar[r] & C^+ \ar[rr] & & G \ar[rr] & & \mu_\tau \ar[r] & 1
\end{tikzcd}

( cf p. 581 (45) ) , et diagrammes analogues pour $G$.

\newpage

%% ===== Page 668 =====

Je voudrais maintenant définir l'analogue
des sous-groupes $\mathcal{M}(1/2), \mathcal{M}(-1), \mathcal{M}^{(0)}(-1), \mathcal{N}_{\mathcal{S}}^\mathcal{G}$ de p.\ 615 ff,
ce qui implique qu'on ait un sous-groupe de $G_{\bar{p}, \sigma}$
jouant le rôle de $\mathcal{M}_0$ - l'idéal serait que
ce soit $G_{\bar{p}, \sigma}$ lui-même. Mais dans le cas
universel $\mathcal{G} = \operatorname{Sl}(2, \mathbf{Z})^\wedge, \mathcal{G}^+ = \Pi_0$, d'où
$Z_{\bar{p}, \sigma} = \mathcal{M}(0)^\flat = \mathcal{M}(0) . L_0^\pi$, (où $L_0^\pi = L_0^\pi$, car $G_{\bar{p}, \sigma} \simeq \mathcal{M}$,
et $G_{\bar{p}, \sigma}^\flat =$ sous-groupe de $G_{\bar{p}, \sigma} = \mathcal{M}$ engendré par
$\mathcal{G}^+ = \Pi_0$ et par $Z_{\bar{p}, \sigma} = \mathcal{M}(0) . L_0^\pi$
$= \mathcal{M}(0) . \pi = \mathcal{M}^\flat$, et non $\mathcal{M}(0) . \Pi_0 = \mathcal{M}$ !

L'ennui provient finalement du fait
que c'est $Z_{\bar{p}, \sigma}$ qui est trop gros, on
aimerait définir dans $Z_{\bar{p}, \sigma}$ un sous-groupe
qui joue le rôle de $\mathcal{M}(0) . \Pi_0$. Il n'y
a de difficulté qu'à en définir
dans $Z_{\bar{p}, \sigma}$ un sous-groupe qui
lui corresponde : $\mathcal{M}(0)$
$$
Z_{\bar{p}, \sigma}(0) \subset Z_{\bar{p}, \sigma} \leqno{(65)}
$$
qui donne un scindage
$$
Z_{\bar{p}, \sigma} \simeq Z_{\bar{p}, \sigma}(0) . L_0^\pi \quad (\text{1/2 direct}) \leqno{(66)}
$$
et alors
$$
G_{\bar{p}, \sigma}^\flat \simeq Z_{\bar{p}, \sigma}(0) . \mathcal{G}^+ \quad (\text{1/2 direct}) \leqno{(67)}
$$

\newpage

%% ===== Page 669 =====

et on définirait
$$
(68) \qquad \begin{cases}
G_{\rho,\sigma}^{\natural\natural} \overset{dfn}{=} Z_{\rho,\sigma}(0) . \underline{S}^{\natural} \subset \underset{\underset{\textstyle Z_{\rho,\sigma}(0) . \underline{S}^+}{\Vert}}{G_{\rho,\sigma}^{\natural}} \subset G_{\rho,\sigma} \\
Z_{\rho,\sigma}^{\natural} = Z_{\rho,\sigma} \\
G_{\rho,\sigma}^{\natural\natural} \cap Z_{\rho,\sigma} = Z_{\rho,\sigma}(0) . (L_0 \cap \underline{S}^{\natural})
\end{cases}
$$

Rappelons une façon standard d'obtenir un scindage de con. de $Z_{\rho,\sigma}$ sur $\Gamma_{\rho,\sigma}$, quand on dispose d'un hom.
$$
(69) \qquad \underline{S}^+ \to \operatorname{Sl}(2,\hat{\mathbb{Z}})' = \operatorname{Sl}(2,\hat{\mathbb{Z}})/\{\pm 1\}
$$
commutant à l'action de $\pi$, ce qui induit un hom.
$$
(70) \qquad Z_{\rho,\sigma} \to B_0' = \left\{ \left(\begin{matrix} \mu & \lambda \\ 0 & 1 \end{matrix}\right) \mid \mu \in \hat{\mathbb{Z}}^\times, \lambda \in \hat{\mathbb{Z}} \right\}
$$
(qui [unclear: n' l'était au] § 48)
induisant le
$$
(71) \qquad \begin{cases}
L_0^\natural \longrightarrow L_0' = \left\{ \left(\begin{matrix} 1 & \lambda \\ 0 & 1 \end{matrix}\right) \mid \lambda \in \hat{\mathbb{Z}} \right\} \\
\rho^n \longmapsto \left(\begin{matrix} 1 & n \\ 0 & 1 \end{matrix}\right)
\end{cases}
$$
plus précisément un diagramme de suites exactes

(72)
\begin{tikzcd}
1 \ar[r] & L_0^\natural \ar[r] \ar[d, "\rho^n \mapsto n"'] & Z_{\rho,\sigma} \ar[r] \ar[d] & \Gamma_{\rho,\sigma} \ar[r] \ar[d, "\chi_\sigma"'] & 1 \\
1 \ar[r] & \hat{\mathbb{Z}} \ar[r, "\lambda \mapsto \left(\begin{smallmatrix} 1 & \lambda \\ 0 & 1 \end{smallmatrix}\right)"'] & B_0' \ar[r, "\text{det}", "\left(\begin{smallmatrix} \mu & \lambda \\ 0 & 1 \end{smallmatrix}\right) \mapsto \mu"'] & \hat{\mathbb{Z}}^\times \ar[r] & 1
\end{tikzcd}

Si l'hom. (71) a un scindage, alors tout

\newpage

%% ===== Page 670 =====

scindages de la 2$^{\text{ième}}$ suite exacte, on en déduit par
image inverse un de la première.
Mais dans les cas typiques, $L_0^{\natural}$ (qui n'
est pas un petit groupe comme "le beau"
$C_{\rho,\sigma}^{\natural}$ !) n'a pas non plus par miracle
pour donner un isomorphisme de (71) ! Mais
on peut aussi faire une hypothèse
supplémentaire - savoir que l'image
inverse par (70) du "Tore"
$$
T'_0 = \left\{ \begin{pmatrix} \mu & 0 \\ 0 & 1 \end{pmatrix} \mid \mu \in \hat{\mathbf{Z}}^\times \right\} \subset B'_0 \leqno{(73)}
$$
qui est un sous-groupe
$$
Z_{\rho,\sigma}(0) \subset Z_{\rho,\sigma} \leqno{(74)}
$$
$$
= \left\{ u \in Z_{\rho,\sigma} \mid \text{son image } \begin{pmatrix} \mu & \lambda \\ 0 & 1 \end{pmatrix} \text{ est dans } T'_0, \text{ i.e. } \lambda = 0 \right\}
$$
satisfaisant
$$
Z_{\rho,\sigma}(0) \cap L_0^{\natural} = \{ 1 \} \leqno{(75)}
$$
soit en fait un sous-groupe section,
i.e.\ que $Z_{\rho,\sigma}(0) \hookrightarrow \Gamma_{\rho,\sigma}$ soit épi i.e.\ iso,
i.e.\ que l'on ait bien
$$
Z_{\rho,\sigma} = Z_{\rho,\sigma}(0) . L_0^{\natural} . \leqno{(76)}
$$
J'ai de bonnes raisons de penser que
(dès que (69) existe) cette condition est
satisfaite dans tous les cas vraiment utiles.
quitte, s'il le faut, à [unclear: modifier] le choix d'un

\newpage

%% ===== Page 671 =====

$Z_{\rho,\sigma}(\omega)$ a la bonne propriété d'ailleurs d'être "fonctoriel" dans un sens évident (cf. plus bas la question de fonctorialité), ce qui sera essentiel pour avoir un bon formalisme. Ceci nous permet donc de définir un sous-groupe 
$$ 
(G_{\rho,\sigma})_{Top} \subset G_{\rho,\sigma} \leqno{(68)} 
$$
Si sa définition n'était ainsi un peu trop [unclear: bâtarde] \footnote{Avec "plus fine" rajouté au-dessus de bâtarde.}, c'est qu'elle serait meilleure, et mériterait en fait la notation plus simple $G_{\rho,\sigma}^\natural$.

Autre façon de procéder ? Dans le cas où on lève les restrictions odieuses dans la définition de $Z_{\rho,\sigma}$, qui (avec $L_0^\natural = L_0 \cap \mathfrak{S}^\natural$, i.e. qui) impliquent que (pour $l \in L_0$, on a
$$
( (\operatorname{int} l), 1, 1, 1 ) \in Z_{\rho,\sigma} \iff l \in \mathfrak{S}^\natural
$$

Mais dans le cas universel, où on a $\mathfrak{S}^\natural = \Pi_0$, prenant $l = l_0$, [unclear: l'image de] $\mathfrak{S}/\mathfrak{S}^\natural = \mathfrak{S}/\Pi_0$ [unclear: est un élémt central trivial de] $\mathfrak{S}/\Pi_0$, l'automorphisme qu'il définit dans $\mathfrak{S}^+ / \Pi_0$ ($cl = \text{dans } \mathfrak{S}/\Pi_0$) est trivial - donc on n'a pas pu tirer parti des propriétés de cette action pour prouver l'exercice !

Mais peut-être est-il prématuré d'essayer d'ici de "disséquer" un groupe $\hat{\mathfrak{S}}_{\rho,\sigma}^\natural, \mathcal{M}_{\rho,\sigma}^\natural$

\newpage

%% ===== Page 672 =====

"aussi fin que possible" - peut-être y a-t-il
encore des erreurs d'orientation assez grossières, qui
seront rectifiées par l'examen de cas
particuliers. Donc je vais renoncer
pour l'instant à continuer à investir de
de l'énergie dans ces digressions, et m'en
tiendrai donc là dans ce sens, quitte
à y revenir par la suite et à raffiner
les constructions présentes, qu'il faudra
cependant repenser à l'occasion, [unclear: et] celles
du § 48 \underline{XI}.

\newpage

%% ===== Page 673 =====

$5^{\circ}$) $S_0 M_{\rho,\sigma}^\natural, M_{\rho,\sigma}^\natural, S_0 \Gamma_{\rho,\sigma}^\natural, \Gamma_{\rho,\sigma}^\natural \dots$

On posera comme dans p. 583

(77)
$$ \begin{cases} S_0 M_{\rho,\sigma} = Z_{\rho,\sigma} \times_\gamma N_\rho^\natural \times_\gamma N_\sigma^\natural \times_\gamma C_{\rho,\sigma} \\ S_0 M_{\rho,\sigma}^\natural = Z_{\rho,\sigma} \times_\gamma N_\rho^\natural \times_\gamma N_\sigma^\natural \times_\gamma C_{\rho,\sigma}^\natural \end{cases} $$

[unclear: (78')]
$$ \begin{cases} S_0 \Gamma_{\rho,\sigma} = Z_{\rho,\sigma} \times_\gamma N_\rho^\natural \times N_\sigma^\natural \\ \Gamma_{\rho,\sigma} = N_\rho^\natural \times_\gamma N_\sigma^\natural \end{cases} $$

(78)
$$ \begin{cases} M_{\rho,\sigma} \simeq N_\rho^\natural \times N_\sigma^\natural \times C_{\rho,\sigma} \\ M_{\rho,\sigma}^\natural \simeq N_\rho^\natural \times N_\sigma^\natural \times C_{\rho,\sigma}^\natural \end{cases} $$

d'où un diagramme de suites exactes

(79)
\begin{tikzcd}[row sep=2em, column sep=1.5em]
& 1 \ar[rr] \ar[dl] & & C^{\natural,+} \ar[rr] \ar[dd, equal] \ar[dl] & & S_0 M_{\rho,\sigma}^\natural \ar[rr] \ar[dd] \ar[dl] & & S_0 \Gamma_{\rho,\sigma}^\natural \ar[rr] \ar[dd] \ar[dl] & & 1 \ar[dl] \\
1 \ar[rr] & & C^+ \ar[rr] \ar[dd, equal] & & S_0 M_{\rho,\sigma} \ar[rr] \ar[dd] & & S_0 \Gamma_{\rho,\sigma} \ar[rr] \ar[dd] & & 1 \\
& 1 \ar[rr] \ar[dl] & & C^{\natural,+} \ar[rr] \ar[dl] & & M_{\rho,\sigma}^\natural \ar[rr] \ar[dl] & & \Gamma_{\rho,\sigma}^\natural \ar[rr] \ar[dl] & & 1 \ar[dl] \\
1 \ar[rr] & & C^+ \ar[rr] & & M_{\rho,\sigma} \ar[rr] & & \Gamma_{\rho,\sigma} \ar[rr] & & 1
\end{tikzcd}

et des suites exactes

(80)
\begin{tikzcd}[row sep=1.5em]
1 \ar[r] & Z_{\rho,\sigma}^\natural \times S N_\rho^\natural \times S N_\sigma^\natural \times C_{\rho,\sigma}^\natural \ar[r] \ar[d] & S_0 M_{\rho,\sigma}^\natural \ar[r] \ar[d] & \overline{\Gamma}_{\rho,\sigma} \ar[r] & 1 \\
1 \ar[r] & Z_{\rho,\sigma} \times S N_\rho^\natural \times S N_\sigma^\natural \times C_{\rho,\sigma} \ar[r] & S_0 M_{\rho,\sigma} \ar[r] & \overline{\Gamma}_{\rho,\sigma} \ar[r] & 1
\end{tikzcd}

(82)
\begin{tikzcd}[row sep=1.5em]
1 \ar[r] & S N_\rho^\natural \times S N_\sigma^\natural \times C_{\rho,\sigma}^\natural \ar[r] \ar[d] & M_{\rho,\sigma}^\natural \ar[r] \ar[d] & \overline{\Gamma}_{\rho,\sigma} \ar[r] & 1 \\
1 \ar[r] & S N_\rho^\natural \times S N_\sigma^\natural \times C_{\rho,\sigma} \ar[r] & M_{\rho,\sigma} \ar[r] & \overline{\Gamma}_{\rho,\sigma} \ar[r] & 1
\end{tikzcd}

(83)
$$ 1 \to S N_\rho^\natural \times S N_\sigma^\natural \to S_0 \Gamma_{\rho,\sigma}^\natural \to \overline{\Gamma}_{\rho,\sigma} \to 1 $$

(84)
$$ 1 \to S N_\rho^\natural \times S N_\sigma^\natural \to \Gamma_{\rho,\sigma} \to \overline{\Gamma}_{\rho,\sigma} \to 1 $$

\newpage

%% ===== Page 674 =====

$$
\begin{cases}
S_0 \mathcal{M}_{\rho,\sigma} \simeq \mathcal{M}_{\rho,\sigma} \times_{\gamma_{\rho,\sigma}} Z_{\rho,\sigma} \\
S_0 \mathcal{M}_{\rho,\sigma}^\natural \simeq \mathcal{M}_{\rho,\sigma}^\natural \times_{\gamma_{\rho,\sigma}^\natural} Z_{\rho,\sigma}
\end{cases} \leqno{(85)}
$$

$$
S_0 \Gamma_{\rho,\sigma} \simeq \Gamma_{\rho,\sigma} \times_{\gamma_{\rho,\sigma}} Z_{\rho,\sigma} \leqno{(86)}
$$

$$
\begin{cases}
S_0 \mathcal{M}_{\rho,\sigma} \simeq \mathcal{M}_{\rho,\sigma} \times_{\Gamma_{\rho,\sigma}} S_0 \Gamma_{\rho,\sigma} \\
S_0 \mathcal{M}_{\rho,\sigma}^\natural \simeq \mathcal{M}_{\rho,\sigma}^\natural \times_{\Gamma_{\rho,\sigma}^\natural} S_0 \Gamma_{\rho,\sigma}
\end{cases} \leqno{(87)}
$$

Remarque. Cette kyrielle de groupes donne un peu le mal de mer. Je ne suis pas sûr que les variantes $S_0$ ($S_0 \mathcal{M}$, $S_0 \mathcal{M}^\natural$, $S_0 \Gamma$) soient très utiles. Reste la comparaison entre les variantes avec ou sans $\natural$, entre $\mathcal{M}_{\rho,\sigma}$ et $\mathcal{M}_{\rho,\sigma}^\natural$ -- dont la "différence" est "la même" qu'entre $\mathfrak{S}^+$ et $\mathfrak{S}^\natural$, et d'autre part $C_{\rho,\sigma}$ et $C_{\rho,\sigma}^\natural$, de façon plus précise $\mathcal{M}_{\rho,\sigma}$ est image inverse de $\mathcal{M}_{\rho,\sigma}^\natural$, et $C_{\rho,\sigma}$ de $C_{\rho,\sigma}^\natural$ d'où des isomorphismes

$$
\mathfrak{S}^+ / \mathfrak{S}^\natural \isommap C_{\rho,\sigma} / C_{\rho,\sigma}^\natural \isommap \mathcal{M}_{\rho,\sigma} / \mathcal{M}_{\rho,\sigma}^\natural. \leqno{(88)}
$$

Les groupes $C_{\rho,\sigma}$, $\mathcal{M}_{\rho,\sigma}$ jouent par rapport à $C_{\rho,\sigma}^\natural$, $\mathcal{M}_{\rho,\sigma}^\natural$ le rôle de fourre-touts commodes pour nous permettre de considérer les automorphismes $\rho, \sigma$ de $C_{\rho,\sigma}^\natural$, $\mathcal{M}_{\rho,\sigma}^\natural$ comme induits par des automorphismes intérieurs -- i.e.\ des générateurs.

\newpage

%% ===== Page 675 =====

environnant $G_{\rho, \sigma}, \mathcal{M}_{\rho, \sigma}$.

On a un schéma
$$
\begin{matrix}
\tau_{\mathcal{M}}' \in \mathcal{M}_{\rho, \sigma}^{\mathcal{G}} & \text{au dessus de } \tau_f' \in \Gamma_{\rho, \sigma} \\
\rotatebox{90}{$=$} & \\
(\tau_{\rho}', \tau_{\sigma}', \tau_a') & \\
\cap \quad \cap \quad \cap & \\
N_{\rho}^{\mathcal{G}} \quad N_{\sigma}^{\mathcal{G}} \quad N_a^{\mathcal{G}} \subset G_{\rho, \sigma}^{\mathcal{G}} & 
\end{matrix}
\leqno{(89)}
$$
d'où des plongements canoniques
$$
\begin{matrix}
\mathfrak{S}^{\mathcal{G}} \hookrightarrow \mathcal{M}_{\rho, \sigma}^{\mathcal{G}} & \qquad \mathfrak{S}^{\mathcal{G}} = \text{im. inv. dans } \mathcal{M}_{\rho, \sigma}^{\mathcal{G}} \text{ de } \{1, \tau_f'\} \subset \Gamma_{\rho, \sigma} \\
\mathfrak{S} \hookrightarrow \mathcal{M}_{\rho, \sigma} & \qquad \mathfrak{S} = \text{im. inv. dans } \mathcal{M}_{\rho, \sigma} \text{ de } \{1, \tau_f\} \subset \Gamma_{\rho, \sigma}
\end{matrix}
\leqno{(90)}
$$

On peut reprendre la digression bouddhique de p.\ 584, sur les $\mathcal{G}$ : le [unclear: def] \dots

Je vais ici m'en dispenser, [unclear: inutile] de redévelopper les laïus sur $Z_{\rho, \sigma}, \mathfrak{S} \mathcal{G}_{\rho, \sigma}$ etc. --- quitte à y revenir au besoin par la suite.

\newpage

%% ===== Page 676 =====

6°) Fonctorialités.

[MARGIN: NB jusqu'ici je n'avais pas utilisé l'hyp. que $\phi$ surjectif vu que dans (91)]

(cf 1 ; 5°))

cf. p. 588. En plus de l'hom.
Je n'examinerai pas [unclear] cf p. 643' ss...

(91)
\begin{tikzcd}
\mathcal{C} \ar[r, "\phi"] & \mathcal{C}'
\end{tikzcd}

supposons que $\phi$ applique $\mathcal{C}^\natural$ dans $\mathcal{C}'^\natural$.

On désigne ici par $Z_{\rho,\sigma}$, $M_{\rho,\sigma}$ etc les invariants associés à $\mathcal{C}$, $Z'_{\rho',\sigma'}$ etc ceux de $\mathcal{C}'$. On a des conditions équivalentes comme dans loc. cit - où on va laisser tomber la considération de $\bar{Z}_{\rho,\sigma}$ ... qui signifiait que pour tout
$$\underline{u} = (u, u_1, \varepsilon_2, \varepsilon_3) \in Z_{\rho,\sigma}$$
il existe un autom $u'$ de $\mathcal{C}'^+$ rendant commutatif

(92)
\begin{tikzcd}
\mathcal{C}^+ \ar[r, "\phi"] \ar[d, "u"'] & \mathcal{C}'^+ \ar[d, "u'"] \\
\mathcal{C}^+ \ar[r, "\phi"'] & \mathcal{C}'^+
\end{tikzcd}

lequel $u'$ sera unique. Je veux donc = à dire que les $u \in \operatorname{Aut}(\mathcal{C}^+)$ qui proviennent de $Z_{\rho,\sigma}$, i.e. qui normalisent $L_\sigma$, et qui transforment $k_{\rho,\sigma}$ en des $k_{\rho',\sigma'}$ [unclear] d'invariants $\mathcal{C}^\natural$ conjoints à $\rho,\sigma$, stabilisent le noyau de $\phi$ (91).

On trouve alors

\newpage

%% ===== Page 677 =====

$$ Z_{\underline{\rho},\sigma} \overset{\varphi_Z}{\longrightarrow} Z_{\underline{\rho}',\sigma'} \leqno{(93)} $$
compatible avec les actions de $Z_{\underline{\rho},\sigma}$, $Z_{\underline{\rho}',\sigma'}$ sur
$\underline{\mathfrak{S}}^+$ resp.\ $\underline{\mathfrak{S}}'^+$, et avec les hom.
$$ \underline{L}_0^\natural \subset Z_{\underline{\rho},\sigma} \to \underline{\mathfrak{S}}^+ , \quad \underline{L}'^\natural_0 \subset Z_{\underline{\rho}',\sigma'} \to \underline{\mathfrak{S}}'^+ , $$
d'où par la construction standard un hom.
$$ C_{\underline{\rho},\sigma} \overset{\varphi_C}{\longrightarrow} C_{\underline{\rho}',\sigma'} \leqno{(94)} $$
induisant
$$ \underline{C}_{\underline{\rho},\sigma}^\natural \overset{\varphi_{\underline{C}^\natural}}{\longrightarrow} \underline{C}_{\underline{\rho}',\sigma'}'^\natural \leqno{(95)} $$
d'où un hom.\ des diagrammes
(69) relatif à $(\underline{\rho},\sigma)$, dans le diagramme
analogue relatif à $(\underline{\rho}',\sigma')$, [unclear: comme on s'y attend].

Il n'y a pas à faire plus pour que toutes
les constructions des $n^{\text{os}}$ précédents, relatives
à $\underline{\rho},\sigma$, s'envoient dans celles relatives à $\underline{\rho}',\sigma'$.

Remarque: Supposons que $\underline{\mathfrak{S}}'$ satisfait à
la condition suivante : Pour tout couple
d'éléments $\alpha', \beta' \in \underline{\mathfrak{S}}'_0$, tels que l'on
ait
$$ [\alpha' \rho] = [\beta' \sigma] \varepsilon_1'^{-1} \quad \text{pour } [unclear: p \in \hat{\mathbf{Z}}] \text{ conv.} \leqno{(96)} $$
il existe un $u' \in \operatorname{Aut}(\underline{\mathfrak{S}}'^+)$ tel que
$$ u'(\rho) = \alpha'(\rho), \ u'(\sigma) = \beta'(\sigma) \leqno{(97)} $$
$$ \alpha'^{-1} u' \in \operatorname{Cent}_{\Aut}(\underline{\rho}), \ \beta'^{-1} u' \in \operatorname{Cent}_{\Aut}(\sigma) $$

\newpage

%% ===== Page 678 =====

la relation de base (96) s'exprime en termes de $u'$ par la condition
$$
u'(L_0'^{\natural}) = L_0'^{\natural} \text{ i.e. } u' \text{ normalise } \mathfrak{S}'_1 . \leqno{(99)}
$$
Alors tout homomorphisme $(\mathfrak{S}, \mathfrak{S}^{\natural}) \to (\mathfrak{S}', \mathfrak{S}'^{\natural})$ est nécessairement admissible.

Nous allons appliquer ceci au cas universel $-$ mais pour préciser ce dire, je vais d'abord parler des sous-groupes
$$
\mathcal{M}(\infty), \mathcal{M}(\rho), \mathcal{M}(\sigma) \subset \mathcal{M}_{\rho,\sigma}.
$$

\newpage

%% ===== Page 679 =====

\marginpar{Les groupes $M_{p,\sigma}[0]$, $M_{p,\sigma}(\rho)$, $M_{p,\sigma}(\sigma)$.}7º) Dans le groupe $SoM_{p,\sigma}$ -
$$
\begin{aligned}
SoM_{p,\sigma} &= Z_{p,\sigma} \times_{S_{p,\sigma}} N_p \times_{S_{p,\sigma}} N_\sigma \times_{S_{p,\sigma}} C_{p,\sigma} \\
&= \{ x = (u, a, b, U) \in Z_{p,\sigma} \times N_p \times N_\sigma \times C_{p,\sigma} \mid \delta_Z(u) = \delta_p(a) = \delta_\sigma(b) = \delta_C(U) \}
\end{aligned} \leqno{(100)}
$$
on introduit trois sous-groupes
$$
\begin{cases}
SoM_{p,\sigma}[0] \subset SoM_{p,\sigma} \quad \text{par condition} \ U = i_Z(u) \\
SoM_{p,\sigma}(\rho) \subset SoM_{p,\sigma} \quad \text{---------} \ U = i_p(a) \\
SoM_{p,\sigma}(\sigma) \subset SoM_{p,\sigma} \quad \text{---------} \ U = i_\sigma(b)
\end{cases} \leqno{(101)}
$$
en utilisant naturellement les trois hom. canon.
$$
\begin{cases}
i_Z : Z_{p,\sigma} \to C_{p,\sigma} \\
i_p : N_p \to C_{p,\sigma} \\
i_\sigma : N_\sigma \to C_{p,\sigma}
\end{cases} \leqno{(102)}
$$
Ces trois groupes sont manifestement tous
trois isomorphes : $Z_{p,\sigma} \times_{S_{p,\sigma}} N_p \times_{S_{p,\sigma}} N_\sigma = So\Gamma_{p,\sigma}$
de façon précise, dans la suite exacte (58)
$$ 1 \to \mathfrak{S}^+ \to SoM_{p,\sigma} \to So\Gamma_{p,\sigma} \to 1 $$
l'hom. canonique $SoM_{p,\sigma} \to So\Gamma_{p,\sigma}$
induit des isom. entre les groupes (101)
et $So\Gamma_{p,\sigma}$, ces trois groupes sont donc des
sous-groupes sections.\marginpar{[unclear: en $M_{p,\sigma}(k)$]}

On définit $M_{p,\sigma}[0]$, $M_{p,\sigma}(\rho)$ et $M_{p,\sigma}(\sigma)$
comme les images de
$SoM_{p,\sigma}[0]$, $SoM_{p,\sigma}(\rho)$ et $SoM_{p,\sigma}(\sigma)$ par l'hom.
canonique $SoM_{p,\sigma} \to M_{p,\sigma}$, ce qui

\newpage

%% ===== Page 680 =====

revient à poser

(103) 
$$ \begin{cases} 
M_{\rho,\sigma}[0] = \{ (a,b,U) \in M_{\rho,\sigma}^\natural \mid U \in Z_{\rho,\sigma} \} \\
M_{\rho,\sigma}(\rho) = \{ (a,b,U) \in M_{\rho,\sigma}^\natural \mid U = a \} \\
M_{\rho,\sigma}(\sigma) = \{ (a,b,U) \in M_{\rho,\sigma}^\natural \mid U = b \}
\end{cases} $$

Donc les sous-groupes $M_{\rho,\sigma}(\rho), M_{\rho,\sigma}(\sigma)$ de $M_{\rho,\sigma}^\natural$ sont des groupes entiers au dessus de $\Gamma_{\rho,\sigma}^\natural \simeq N_\rho^\natural \times_{\Gamma_{\rho,\sigma}} N_\sigma^\natural$, donc

(104) 
$$ \begin{cases} 
M_{\rho,\sigma}^\natural \simeq M_{\rho,\sigma}(\rho) \cdot C^\natural \quad (1/2 \text{ directe}) \\
M_{\rho,\sigma}^\natural \simeq M_{\rho,\sigma}(\sigma) \cdot C^\natural \quad \quad \quad )
\end{cases} $$

et de même

(105) 
$$ \begin{cases} 
M_{\rho,\sigma} \simeq M_{\rho,\sigma}(\rho) \cdot C^+ \quad (1/2 \text{ directe}) \\
M_{\rho,\sigma} \simeq M_{\rho,\sigma}(\sigma) \cdot C^+
\end{cases} $$

D'autre part, on a

(106)
\begin{tikzcd}
1 \ar[r] & L_0^\natural \ar[r] \ar[dd, equal] & M_{\rho,\sigma}[0] \ar[r] & \Gamma_{\rho,\sigma} \ar[r] \ar[dd, equal] & 1 \\
& & S_0 M_{\rho,\sigma}[0] \ar[u, "\simeq"'] \ar[d, "\simeq"] & & \\
1 \ar[r] & L_0^\natural \ar[r] & S_0 \Gamma_{\rho,\sigma} \ar[r] & \Gamma_{\rho,\sigma} \ar[r] & 1
\end{tikzcd}

donc l'extension $M_{\rho,\sigma}[0]$ de $\Gamma_{\rho,\sigma}$ par $L_0^\natural$ s'identifie à l'ext. $S_0 \Gamma_{\rho,\sigma}$ de $\Gamma_{\rho,\sigma}$ par $L_0^\natural$, i.e. à l'image inverse de l'extension $Z_{\rho,\sigma}$ de $\Gamma_{\rho,\sigma}$ par $L_0^\natural$, via l'hom. canonique

\newpage

%% ===== Page 681 =====

$$
\Gamma_{\rho,\sigma} = N_{\rho}^{\mathcal{G}} \times_{\Gamma_{\rho,\sigma}} N_{\sigma}^{\mathcal{G}} \longrightarrow \Gamma_{\rho,\sigma} \ . \leqno{(107)}
$$

Donc la donnée d'une section de $Z_{\rho,\sigma}$ sur $\Gamma_{\rho,\sigma}$ (cf à ce sujet p.\ 636! ff) en définit une de $\operatorname{Sol}_{\Gamma,\sigma}$ sur $\Gamma_{\rho,\sigma}$, donc de $\mathcal{M}_{\rho,\sigma}[0]$ sur $\Gamma_{\rho,\sigma}$. Quand une telle section est fixée, on en déduit en passant aux quotients par
$$
\mathcal{M}_{\rho,\sigma}(\rho) \subset \mathcal{M}_{\rho,\sigma}^{\mathcal{G}} \ ,
$$
d'où --
$$
\begin{cases}
\mathcal{M}_{\rho,\sigma}[0] \overset{S}{\simeq} \mathcal{M}_{\rho,\sigma}(\rho) \rtimes L_0^{\mathcal{G}} \\
\mathcal{M}_{\rho,\sigma}^{\mathcal{G}} \overset{S}{\simeq} \mathcal{M}_{\rho,\sigma}(\rho) \rtimes \mathcal{G}^{\mathcal{G}} \\
\mathcal{M}_{\rho,\sigma} \simeq \mathcal{M}_{\rho,\sigma}(\rho) \rtimes \mathcal{G}
\end{cases} \leqno{(108)}
$$

\newpage

%% ===== Page 682 =====

8\textsuperscript{o}) Retour sur le cas universel, et l'homomorphisme
$$ \mathcal{M}^\mathfrak{S} \longrightarrow \mathcal{M}_{p,\sigma}^\mathfrak{S} / \mathcal{M}_\mathfrak{S} $$

Nous supposons maintenant que $\mathfrak{S}$ correspond au cas universel $\mathfrak{S}_{0,3}$, avec un $\mathfrak{S}$ quelconque. On notera le cas universel par un indice $0$ : 
$$ \mathfrak{S}_0, \mathfrak{S}_0^+, \mathfrak{S}_0^\mathfrak{S} \text{ etc}, $$
et $\mathfrak{S}', \mathfrak{S}'^\mathfrak{S}$ etc. [unclear: le sous-groupe] $\mathfrak{S}, \mathfrak{S}^\mathfrak{S}$ etc.
\marginpar{- $\mathbb{Z}_{p,\sigma}$ et indices $0$ \dots \\ On écrira $\mathfrak{S}_0^+$ au lieu de $\mathfrak{S}_0^{+\wedge}$ \dots}

Je pose
\marginpar{NB [unclear: $\mathcal{M}_\mathfrak{S} \dots$]}
$$
\mathcal{M}^\mathfrak{S} \defeq \left\{ u \in \mathcal{M} \mathrel{\Bigg|} 
\begin{array}{l} 
\text{a) } u(p) \text{ est } \mathfrak{S}\text{-conjugué à } \varepsilon_3[\tau](p) \\ 
\text{b) } u(\sigma) \text{ est } \mathfrak{S}\text{-conjugué à } \varepsilon_2[\tau](\sigma) \\ 
\text{c) } u(\varepsilon_0) \text{ est } \mathfrak{S}\text{-conjugué à } \varepsilon_0^\mu \\ 
\text{d) } u \text{ stabilise } \mathcal{M}_\mathfrak{S} 
\end{array} 
\right\} \leqno{(107)}
$$

bien sûr
$$ \mu = \chi(u), \quad \varepsilon_2 = \varepsilon_2(u) = \varepsilon_2(1!), \quad \varepsilon_3 = \varepsilon_3(u) = \varepsilon_3(1!) $$

Je dis que $\mathcal{M}^\mathfrak{S}$ est un sous-groupe de $\mathcal{M}$. On peut le voir en reprenant les calculs de la p. 629 ff. Je préfère procéder ainsi.

\newpage

%% ===== Page 683 =====

On a [unclear: un isomorphisme] $= \mathcal{M}[0] \cdot L_0^\mu$.

$$
Z_{\mathfrak{S}_0} \hookrightarrow \mathcal{M}[0] \left[ \simeq \mathcal{M}_{0,3}^\sim [0] = \mathcal{M}_{0,3}^\sim (\infty) \cdot L_0^\mu \right] \leqno (108)
$$

On remarque que c'est une égalité si $\mathfrak{S}_0^- = \mathfrak{S}_0^+$, et que l'on identifie $Z_{\mathfrak{S}_0} : \mathcal{M}[T_0] = \mathcal{M}(0) \cdot L_0^\mu$ si $\mathfrak{S}_0 = \pi = \pi_0$.

\marginpar{cette explicitation ne servira du reste à rien.}%
[ --- en fait, si $\mathfrak{S}_0^+ \subset \pi_0$, alors

$$
Z_{\mathfrak{S}_0} = \left\{ \begin{matrix} u = u_0 \cdot l \\ \text{dans } \mathcal{M}(0) \cdot L_0^\mu \end{matrix} \middle| \begin{array}{l} \text{a) } u(p) = \alpha p \beta \text{ avec } \alpha \in \mathfrak{S}^-_0 \\ \text{b) } u(\sigma) = \beta(\sigma) \dots \beta \in \mathfrak{S}^-_0 \\ \text{c) } u \text{ stabilise } \mathcal{M}^\mathfrak{S} \end{array} \right\} \leqno (109)
$$

et on voit que $\alpha, \beta$ de a) b) sont déterminés par $u = u_0 \cdot l$ de façon unique (dès que $u \in \mathcal{M}[T_0]$ au plus). On sait que c'est un [unclear: sous-groupe] $\mathcal{M}(\mathfrak{S}_0^\pm) \simeq \hat{\mathbf{Z}}^\times \times \mu_2 \times \mathcal{M}(\mathfrak{S}_0^+)$ (cf.\ p.\ 679).

Considérons [unclear: l'image de] $\mathcal{M}_{0,3}^\sim [\mathfrak{S}]$ dans $\mathcal{M}_{0,3}^\sim [\pi_0]$, qui s'envoie dans $\mathcal{M}_{0,3}^\sim [T_0]$ comme un sous-groupe, donc dans $\mathcal{M}[\mathfrak{S}] \simeq \mathcal{M}_{0,3}^\sim [T_0]$ par un isomorphisme de groupes. ]

\newpage

%% ===== Page 684 =====

Considérons le système
$$
(Z_{0,\sigma_0}, L_0^{\S} \subset Z_{0,\sigma_0}, Z_{0,\sigma_0} \to \operatorname{Aut}(\mathfrak{S}_0^+), L_0^{\S} \hookrightarrow \mathfrak{S}_0^+) \leqno{(110)}
$$
(cf ``digression'' plus bas 90) qui donne naissance à
$$
G_{\rho_0,\sigma_0}, \mathfrak{S}_0^+ \subset G_{\rho_0,\sigma_0}, Z_{\rho_0,\sigma_0} \hookrightarrow G_{\rho_0,\sigma_0} \leqno{(110')}
$$
et considérons le système analogue à (110*) sur $\mathfrak{S}_0^+$ :
$$
(\underbrace{M[0]^{\S}}_{= \mathcal{M}_{0,3}.L_0^{\pi} \text{ (engendré par } \Sigma_0)}, L_0^{\S} \subset M[0]^{\S}, M[0]^{\S} \to \operatorname{Aut}(\mathfrak{S}_0^+), L_0^{\S} \hookrightarrow \mathfrak{S}_0^+) \leqno{(111)}
$$
qui donne naissance (loc.\ cit (90))
à la situation
$$
(\mathcal{M}, \mathfrak{S}_0^+ \subset \mathcal{M}, M[0]^{\S} \hookrightarrow \mathcal{M}). \leqno{(111')}
$$
Or grâce à (108), on a un homomorphisme de (110) dans (111), correspondant à l'identité sur $\mathfrak{S}_0^+$, (108) sur $Z_{0,\sigma_0}$ l'inclusion [unclear: can.] $L_0^{\S} \hookrightarrow L_0^{\S}$. Il en résulte un homomorphisme de (110') dans (111'), en particulier un homomorphisme canonique de suites exactes

\newpage

%% ===== Page 685 =====

(112)
\begin{tikzcd}
1 \ar[r] & \widehat{\mathfrak{S}}_0^+ \ar[r] \ar[d] & G_{\rho_0,\sigma_0} \ar[r] \ar[d] & \gamma_{\rho_0,\sigma_0} \ar[r] \ar[d] & 1 \\
1 \ar[r] & \widehat{\mathfrak{S}}_0^+ \ar[r] & \mathcal{M} \ar[r] & \tilde{\Gamma}_Q \ar[r] & 1
\end{tikzcd}

où 
\begin{tikzcd}
\gamma_{\rho_0,\sigma_0} \ar[r] \ar[d, "\simeq"'] & \tilde{\Gamma}_Q \ar[d, "\simeq"] \\
Z_{\rho_0,\sigma_0}/L_0^\natural & M[0]/L_0^\natural
\end{tikzcd}
est l'homo. induit par
(108), et s'identifie à :

(113) $Z_{\rho_0,\sigma_0}/L_0^\natural \hookrightarrow M[0]/L_0^\natural$

Par définition de $L_0^\natural$ comme inter-
sect. $Z_{\rho_0,\sigma_0} \cap \widehat{\mathfrak{S}}_0^+$ ds [unclear: $\mathcal{M}$], il résulte que
(113) est inj. donc $G_{\rho_0,\sigma_0} \to \mathcal{M}$
est toujours injectif (c'est un iso. si
$\widehat{\mathfrak{S}}_0^+ \supset \Pi_0$ ...). On déduit par
composition de (112) un hom.
de suites exactes

(114)
\begin{tikzcd}
1 \ar[r] & \widehat{\mathfrak{S}}_0^\natural \ar[r] \ar[d, hook] & G_{\rho_0,\sigma_0}^\natural \ar[r] \ar[d, hook] & \gamma_{\rho_0,\sigma_0} \ar[r] \ar[d] & 1 \\
1 \ar[r] & \widehat{\mathfrak{S}}_0^+ \ar[r] & \mathcal{M} \ar[r] & \tilde{\Gamma}_Q \ar[r] & 1
\end{tikzcd}

On voit alors immédiatement que
l'im. de $G_{\rho_0,\sigma_0}^\natural$ ds $\mathcal{M}$ n'est
autre que $\mathcal{M}^\natural$ (107). On a donc

\newpage

%% ===== Page 686 =====

\footnote{NB le [unclear: scindage] $s$ de $\gamma_{\rho_0,\sigma_0}$ (noté par la suite $s_0$) (initié par un élément $\gamma_0 = s_0 \in \gamma_{\rho_0,\sigma_0} \subset \Gamma_{\mathbf{Q}}^\sim$ en (116)) se relève en une section $u_0 \in \mathcal{M}_{\rho_0,\sigma_0}$ satisfaisant les cond. a) b) d) de (107).}
$$
\begin{cases}
\mathcal{M}^{\mathcal{G}} \isommap G_{\rho_0, \sigma_0}^{\mathcal{G}} \\
\text{définies} \\
(107) \\
\gamma_{\rho_0,\sigma_0} \subset \Gamma_{\mathbf{Q}}^\sim
\end{cases}
\leqno{(115)}
$$
qui montre en même temps dès la formule que $\mathcal{M}^{\mathcal{G}}$ est bien un sous-groupe de $\mathcal{M}$.

On a maintenant
$$
\mathcal{M}_{\rho_0,\sigma_0}^{\mathcal{G}} \defeq \mathcal{M}^{\mathcal{G}} \times_{G_{\rho_0,\sigma_0}^{\mathcal{G}}} \mathcal{N}_{\rho_0}^{\mathcal{G}} \times_{G_{\rho_0,\sigma_0}^{\mathcal{G}}} \mathcal{N}_{\sigma_0}^{\mathcal{G}} \times_{G_{\rho_0,\sigma_0}^{\mathcal{G}}} G_{\rho_0,\sigma_0}^{\mathcal{G}} \simeq \mathcal{M}^{\mathcal{G}} \leqno{(117)}
$$
$$ G_{\rho_0,\sigma_0}^{\mathcal{G}} \simeq \mathcal{M}^{\mathcal{G}} \text{ et } $$
$$ \mathcal{N}_{\rho_0}^{\mathcal{G}} = \left\{ u \in \mathcal{M}^{\mathcal{G}} \mid u(\rho) = \rho^{\varepsilon_3(u)} \right\} $$
$$ \mathcal{N}_{\sigma_0}^{\mathcal{G}} = \left\{ u \in \mathcal{M}^{\mathcal{G}} \mid u(\sigma) = \sigma^{\varepsilon_3(u)} \right\} $$
et finalement
$$
\begin{cases}
\mathcal{M}_{\rho_0,\sigma_0}^{\mathcal{G}} \simeq \mathcal{N}_{\mathcal{M}^{\mathcal{G}}}(\rho) \times_{\mathcal{M}^{\mathcal{G}}} \mathcal{N}_{\mathcal{M}^{\mathcal{G}}}(\sigma) \times_{\Gamma_{\mathbf{Q}}^\sim} \mathcal{M}^{\mathcal{G}} \\
\gamma_0 = \gamma_{\rho_0,\sigma_0} \subset \Gamma_{\mathbf{Q}}^\sim \text{ explicitées dans (116)}
\end{cases}
\leqno{(118)}
$$
d'où un homomorphisme de suites exactes
$$
\begin{tikzcd}
1 \ar[r] & S\hat{\mathbf{Z}}_{\rho} \times S\hat{\mathbf{Z}}_{\sigma} \ar[r, "L_{\rho} \times L_{\sigma}"] \ar[d, equal] & \mathcal{M}_{\rho_0,\sigma_0}^{\mathcal{G}} \ar[r] \ar[d] & \mathcal{M}^{\mathcal{G}} \ar[r] \ar[d] & 1 \\
& S\hat{\mathbf{Z}}_{\rho} \times S\hat{\mathbf{Z}}_{\sigma} \ar[r] & \dots & \dots &
\end{tikzcd}
\leqno{(119)}
$$

\newpage

%% ===== Page 687 =====

$$
\operatorname{Ker}(\mathcal{M}^\natural_{\rho_0, \sigma_0} \xrightarrow{p^\natural} \mathcal{M}^\natural) = S Z_{\rho_0} \times S Z_{\sigma_0} \simeq L_{\rho} \times L_{\sigma} \simeq \hat{\mathbf{Z}} L_{\rho} \times \hat{\mathbf{Z}} L_{\sigma} \leqno{(120)}
$$

Nous allons améliorer encore la [unclear: situation] en découvrant dans $\mathcal{M}^\natural_{\rho_0, \sigma_0}$ un sous-groupe encore plus petit, qui soit une extension de $\mathcal{M}^\natural$ par $\mu$. Pour ceci,

Notons que
$$
\mathcal{M}^\natural_{\rho_0, \sigma_0} = \mathcal{N}^\natural_{\rho_0} \times_{\Gamma^\sim_{\mathbf{Q}}} \mathcal{N}^\natural_{\sigma_0} \times_{\Gamma^\sim_{\mathbf{Q}}} \mathcal{M}^\natural \leqno{(121)}
$$
$$
\cap
$$
$$
\mathcal{N}_{\rho_0} \times_{\Gamma_{\mathbf{Q}}} \mathcal{N}_{\sigma_0} \times_{\Gamma_{\mathbf{Q}}} \mathcal{M} \quad [unclear: \to \Gamma_{\mathbf{Q}} \simeq \Gamma^\sim_{\mathbf{Q}}]
$$
où ce dernier produit fibré n'est autre que la quantité $\mathcal{M}^{(\natural')}_{\rho_0, \sigma_0}$ qui correspondrait à ce $\mathcal{M}_{\rho_0, \sigma_0}$ de $\mathfrak{S}^{\natural'} = \mathfrak{S}^+$ tout entière (noté $\mathcal{M}_{\rho, \sigma}$ à la p. 625 et suivantes). Dans ce groupe on a défini un sous-groupe, qui était noté $\mathcal{M}^\natural_{\rho, \sigma}$ à la p. 626, et que je vais plutôt noter maintenant $\mathcal{M}^\#_{\rho_0, \sigma_0}$, défini comme

\newpage

%% ===== Page 688 =====

$$
\mathcal{M}^? = \mathcal{N}_{\rho_0}^? \times_{\Gamma_0} \mathcal{Z}_{\sigma_0}^! \times_{\Gamma_0} \mathcal{M}
\leqno{(122)}
$$

où l'on a posé
$$
\begin{cases}
\mathcal{N}_{\rho_0} = \operatorname{Norm}_{\mathcal{M}}(L_{\rho_0}), \quad \mathcal{N}_{\sigma_0} = \operatorname{Norm}_{\mathcal{M}}(L_{\sigma_0}) \\
\mathcal{Z}_{\sigma_0} = \operatorname{Centr}_{\mathcal{M}}(L_{\sigma_0}) (= \operatorname{Centr}_{\mathcal{N}_{\sigma_0}}(L_{\sigma_0})) \\
\mathcal{Z}_{\sigma_0}^! = \mathcal{Z}_{\sigma_0} \cap \mathcal{M}_n^! \ ( = \mathcal{N}_{\sigma_0} \cap \mathcal{M}_n^! ) \subset \mathcal{N}_{\sigma_0} \\
\mathcal{N}_{\rho_0}^? = \mathcal{Z}_{\rho_0}^! \cdot H_\tau \ (\subset \mathcal{N}_{\rho_0} \ , \text{ où } \mathcal{Z}_{\rho_0}^! = \mathcal{Z}_{\rho_0} \cap \mathcal{M}_n^!)
\end{cases}
\leqno{(123)}
$$
de sorte qu'on a bien\marginpar{indépendant des choix divers}
$$
\mathcal{M}^? \subset \mathcal{N}_{\rho_0} \times_{\Gamma_0} \mathcal{N}_{\sigma_0} \times_{\Gamma_0} \mathcal{M}
$$

Ceci rappelé, on va poser
$$
\mathcal{M}_{\rho_0,\sigma_0}^? \subset \mathcal{M}_{\rho_0,\sigma_0}^{\mathcal{G}}
\leqno{(124)}
$$
$$
\Vert
$$
$$
\mathcal{M}_{\rho_0,\sigma_0}^{\mathcal{G}} \cap \mathcal{M}^? \quad \text{(l'intersection dans } \mathcal{N}_{\rho_0} \times_{\Gamma_0} \mathcal{N}_{\sigma_0} \times_{\Gamma_0} \mathcal{M}\text{)}
$$

et considérant le carré induit par (119)
$$
\begin{tikzcd}
\mathcal{M}_{\rho_0,\sigma_0}^? \arrow[r] \arrow[d] & \mathcal{M}^{\mathcal{G}} \arrow[d] \\
\mathcal{M}^? \arrow[r] & \mathcal{M}
\end{tikzcd}
\leqno{(125)}
$$

\newpage

%% ===== Page 689 =====

Or on a (p.~626' (55)) que $\mathcal{M}^? \to \mathcal{M}$
fait de $\mathcal{M}^?$ une extension de $\mathcal{M}$ de noyau
$\mu_{\mathbf{Q}}$. Donc le noyau de $\mathcal{M}_{\rho_0,\sigma_0}^? \to \mathcal{M}^\natural$
est égal à $\mu_{\mathbf{Q}} \cap \mathcal{M}_{\rho_0,\sigma_0}^?$, or il est clair par
l'équation (119) que l'on a $\mu_{\mathbf{Q}} \subset \mathcal{M}_{\rho_0,\sigma_0}^\natural$, donc $\mu_{\mathbf{Q}} \subset \mathcal{M}_{\rho_0,\sigma_0}^?$.
On trouve donc un morphisme de suites exactes
$$
\begin{tikzcd}
1 \arrow[r] & \mu_{\mathbf{Q}} \arrow[r] \arrow[d, equal] & \mathcal{M}_{\rho_0,\sigma_0}^? \arrow[r] \arrow[d] & \mathcal{M}^\natural \arrow[d] \\
1 \arrow[r] & \mu_{\mathbf{Q}} \arrow[r] & \mathcal{M}^? \arrow[r] & \mathcal{M} \arrow[r] & 1
\end{tikzcd} \leqno{(126)}
$$
et il reste à voir si l'homomorphisme
$$ \mathcal{M}_{\rho_0,\sigma_0}^? \to \mathcal{M}^\natural $$
est un épimorphisme, et sinon, condition(s) sur son image. \marginpar{$\gamma_0 \subset \widetilde{\Gamma_{\mathbf{Q}}}$ \\ $||$ \\ $\operatorname{Im}(\mathcal{M}^\natural)$} Notons que
la définition (124) [unclear: donne de suite une expression]
compte tenu des expressions (118) pour $\mathcal{M}_{\rho_0,\sigma_0}^\natural$ et (122) pour $\mathcal{M}^?$ :

$$ \mathcal{M}_{\rho_0,\sigma_0}^? = \mathcal{N}_{\rho_0}^? \times_{\mathcal{N}_0^?} \mathcal{M}^\natural \leqno{(127)} $$

où
$$ \begin{cases}
\mathcal{N}_{\rho_0}^? = \mathcal{N}_{\rho_0}^? \cap \mathcal{N}_{\rho_0}^? = \mathcal{N}_{\rho_0}^? \cap \mathcal{M}^? \\
\mathcal{Z}_{\sigma_0}^! = \mathcal{Z}_{\sigma_0}^! \cap \mathcal{N}_{\sigma_0}^? = \mathcal{Z}_{\sigma_0}^! \cap \mathcal{M}^?
\end{cases} \leqno{(128)} $$

\newpage

%% ===== Page 690 =====

et l'hom.\ $(*)$ n'est autre que la
troisième projection du produit fibré.
Donc l'image de $(*)$ n'est autre que l'
image inverse, dans $\mathcal{M}^\mathcal{G}$ par $\mathcal{M}^\mathcal{G} \to \Gamma_0 = \Gamma_{\sigma_0,\sigma_0}$
de l'image de
$$
\mathcal{N}_{\rho_0}^{\mathcal{G}?} \times_{\Gamma_0^\mathcal{G}} \mathcal{N}_{\sigma_0}^{\mathcal{G}!} \to \Gamma_0^\mathcal{G} \ (\subset \Gamma_0) \leqno{(**)}
$$
Il reste donc à montrer que cet hom.\
est surjectif, ou encore que les
hom.
$$ \mathcal{N}_{\rho_0}^{\mathcal{G}?} \to \Gamma_0^\mathcal{G} $$
$$ \mathcal{N}_{\sigma_0}^{\mathcal{G}!} \to \Gamma_0^\mathcal{G} $$
sont surjectifs. (Le premier sera
d'ailleurs une ext.\ de $\Gamma_0^\mathcal{G}$ par $\mu_\mathcal{G}$, le
deuxième un iso).

Soit donc $u \in \Gamma_0^\mathcal{G}$, provenant d'un
$u_\mathcal{G} \in \mathcal{M}_{\rho_0,\sigma_0}^{\mathcal{G} ?}$
satisfaisant donc aux conditions a) b) c) d)
de (107). On sait donc qu'il existe
$\beta_\mathcal{G} \in \mathcal{S}_\mathcal{M}^\mathcal{G}$ tel que $\beta^{-1} u_\mathcal{G} \in \mathcal{N}(\sigma_0)$%
\footnote{J'ai arrêté les vérifications, pensant y revenir dans les jours suivants, il s'avère que ce n'est pas urgent - il importait plutôt de dégager [unclear: les aspects] ``base'' du ``Fermat'' pour [unclear: les groupes discrets] tels que $\mathcal{M}_{\rho,\sigma}$ \dots}.

\newpage

%% ===== Page 691 =====

2$^\circ$) \underline{Digression sur le formalisme des groupes}.

$\odot$ Soient
$$
G, \quad \mathcal{G} \subset G, \quad N \xrightarrow{\varphi_N} G, \quad \text{d'où} \quad SN = \varphi^{-1}(\mathcal{G}) \subset N \leqno{(1)}
$$
la situation d'un groupe $G$, sous-groupe invariant $\mathcal{G}$, et un hom.\ $N$ tels que
$$
N \xrightarrow{epi} G/\mathcal{G} \quad \text{i.e.} \quad N/SN \isom G/\mathcal{G}
$$
on veut reconstituer cette situation à partir de la suivante :
$$
\mathcal{G}, N, \quad SN \subset N, \quad N \xrightarrow{\theta} \operatorname{Aut}(\mathcal{G}), \quad SN \xrightarrow{\varphi = \varphi_{SN}} \mathcal{G} \leqno{(2)}
$$
$\mathcal{G}, N$ deux groupes, $SN$ un sous-groupe distingué de $N$, $N \to \operatorname{Aut}(\mathcal{G})$ un hom.\ de groupes i.e.\ une action de $N$ sur $\mathcal{G}$, enfin $SN \to \mathcal{G}$ un hom.\ de groupes assujettis à :
$$
\begin{cases}
\text{a) l'action de $SN$ sur $\mathcal{G}$ induite par celle de $N$,} \\
\text{soit celle déduite de $SN \to \mathcal{G}$,} \\
\text{via l'op.\ adjointe de $\mathcal{G}$.} \\
\text{b) $SN \to \mathcal{G}$ est un hom.\ compatible avec actions de $N$}
\end{cases} \leqno{(3)}
$$
Voyons comment (1) définit (2), montrons comment un syst.\ (2) permet de reconstituer un syst.\ (1).

\newpage

%% ===== Page 692 =====

On part donc de $\mathcal{G}$, on pose
$$
S_{\theta}G = N.\mathcal{G} \quad (\text{produit semi-direct}) \leqno{(3)}
$$
on définit un hom.
$$
\begin{tikzcd}
S N \arrow[r, "i"] & S_{\theta} G \\
u \arrow[r, maps to] & u . p(u)^{-1}
\end{tikzcd}
\leqno{(5)}
$$

Vérifions que c'est un hom.
\marginpar{NB on a pour $u \in SN$, $x \in \mathcal{G}$ (dans $S_{\theta} G$)\\
(*) $u(x) = u x u^{-1} = p(u) x p(u)^{-1}$\\
dans $S_{\theta} G = N.\mathcal{G}$\\
donc si $x = p(u')$,\\
$u(x) = u x u^{-1} = x$\\
i.e. $x$ commute à $u$}

$$
i(u'u) = (u'u) p(u'u)^{-1}
$$
\begin{align*}
i(u')i(u) &= u' p(u')^{-1} u p(u)^{-1} = u' \big(p(u')^{-1} u p(u')\big) \underbrace{p(u')^{-1} p(u)^{-1}}_{= p(u'u)^{-1}} \\
&= u' u \big(u^{-1} p(u')^{-1} u p(u')\big) p(u'u)^{-1} \\
&= (u'u) p(u'u)^{-1} = i(u'u) \quad \text{OK.}
\end{align*}

On [unclear: voit] par [unclear: ex] en [unclear: utilisant] le fait que $p$ et l'hyp (b), que $SN$ est distingué dans $N$, [unclear: avec l'hyp. que] l'image de (5) est distinguée dans $S_{\theta} G$,
en effet on a pour $a \in N$, $u \in SN$
$$
a(u.p(u)^{-1})a^{-1} = a u a^{-1} \cdot a p(u)^{-1} a^{-1}
$$
$$
= \underbrace{a u a^{-1}}_{\in SN} \cdot p(a u a^{-1})^{-1}
$$
i.e. l'hom. (5) commute à l'action de $N$ sur $SN$ et sur $S_{\theta}G = N.\mathcal{G}$
[unclear: l'image de] $SN$ est normalisée par $N

\newpage

%% ===== Page 693 =====

$\forall u \in S N$, $u \varphi(u)^{-1}$ induit l'opération triviale sur $\mathfrak{S}$ --- ce qui est en effet trivial par (3).

On peut donc poser
$$
G \defeq S_0 G / i(S N) \leqno{(6)}
$$
on trouve alors un diagramme commutatif
$$
\begin{tikzcd}
& N \arrow[dr, "\varphi_N"] & \\
S N \arrow[ur, hook] \arrow[dr, "\varphi_{S N}"'] & & G \\
& \mathfrak{S} \arrow[ur, hook] &
\end{tikzcd} \leqno{(7)}
$$
on voit que l'image de $\mathfrak{S}$ dans $G$ est un sous-groupe invariant, de sorte que l'on a un isomorphisme sur les quotients :
$$
N/S N \isommap G/\mathfrak{S} \leqno{(8)}
$$
ce qui implique que l'on a un isomorphisme
$$
\operatorname{Ker}(\varphi_{S N} : S N \to \mathfrak{S}) \isommap \operatorname{Ker}(\varphi_N : N \to G) \leqno{(9)}
$$
donc
$$
\begin{gathered}
\varphi_N : N \to G \text{ est un plongement} \\
\Updownarrow \\
\varphi_{S N} : S N \to \mathfrak{S} \text{ est un plongement.}
\end{gathered} \leqno{(10)}
$$

On trouve ainsi des \underline{équivalences de catégories}\marginpar{cf. loc. cit. § 5} inverses l'une de l'autre, entre les

\newpage

%% ===== Page 694 =====

situations (4), et celle des situations (5).

\begin{enumerate}
\item[b)] Considérons maintenant une situation
$$
\begin{cases}
G, \mathfrak{G} \subset G, (\varphi_i : N_i \to G)_{i \in I} \\
\text{d'où des } SN_i = \varphi_i^{-1}(\mathfrak{G})
\end{cases} \leqno (11)
$$
avec la condition
$$
\forall i \in I, \quad \varphi_i : N_i \to G/\mathfrak{G} = \Gamma \quad \text{épi} \leqno (12)
$$
Peut-on la reconstituer à partir de
$$
\mathfrak{G}, (N_i, SN_i, (\varphi_i, \Theta_i))_{i \in I} \leqno (13)
$$
où $\forall i \in I$, $(N_i, SN_i, \varphi_i, \Theta_i)$ satisfait les conditions
(3) a) et b) --- et bien sûr $\forall i \in I$
$$
\Theta_i : N_i \to \operatorname{Aut}(\mathfrak{G}), \quad \varphi_i : SN_i \to \mathfrak{G}. \leqno (14)
$$
Je me propose de reconstituer $G$ à l'aide de $\mathfrak{G}$ et des $N_i$, comme une sorte de ``somme amalgamée'' relative :
$$
\begin{tikzcd}
& \mathfrak{G} \arrow[rd, dashed] & \\
SN_i \arrow[ru, "\varphi_i"] \arrow[r, hook, "\text{plong.}"'] & N_i \arrow[r, dashed] & G
\end{tikzcd} \leqno (15)
$$
\end{enumerate}

\newpage

%% ===== Page 695 =====

Mais il est plus vrai --- de dire, que pour tout $i$, on reconstitue via $N_i, S N_i, \varphi_i, \Theta_i$ la situation (pour $i$ fixé)
$$
\begin{tikzcd}
N_i \ar[r] & G_i (\subset G) \\
S N_i \ar[u, hook] \ar[r] & \mathcal{G} \ar[u, hook]
\end{tikzcd} \leqno{(15)}
$$
et qu'on reconstitue (11) lui-même en se donnant un syst. \underline{transitif} \underline{d'isom} \underline{entre} \underline{les} $G_i$ ($i \in I$), induisant l'identité sur $\mathcal{G}$. Une telle donnée implique qu'on [unclear: prenne tous] la donnée d'un syst. transitif d'iso entre les
$$ \gamma_i \defeq N_i / S N_i \leqno{(16)} $$
--- qu'on identifie donc à un même groupe
$$ \gamma, \text{ groupe quotient commun des } N_i \text{ (par les } S N_i) \leqno{(17)} $$
Identifions maintenant les $G_i / \mathcal{G} = \gamma_i$ à $\gamma$ :
$$ G_i / \mathcal{G} \simeq \gamma \leqno{(18)} $$
et formons le produit fibré

\newpage

%% ===== Page 696 =====

\marginpar{$g$ est extension de $\gamma$ par $\mathfrak{S}^I$}
$$
g = \prod_{\substack{i \in I \\ (\text{produit fibré sur } \gamma)}} G_i \supset \mathfrak{S}^I \leqno{(19)}
$$

La donnée d'un système transitif d'isomorphismes entre les $G_i$, induisant l'identité sur les sous-groupes $\mathfrak{S}$, et compatible avec le système transitif d'isomorphismes déjà envisagé entre les $\gamma_i = N_i / SN_i \simeq G_i / \mathfrak{S}$, équivaut alors à la donnée d'un sous-groupe $G$ de $g$, satisfaisant les conditions
$$
\begin{cases}
a) & G \cap \mathfrak{S}^I = \text{sous-groupe } \delta(\mathfrak{S}) \text{ dans } \mathfrak{S}^I \\
b) & \text{La projection } G \to \gamma \text{ est épi}
\end{cases} \leqno{(20)}
$$

En d'autres termes, on a construit, à l'aide de (13) et du système transitif d'isomorphismes entre les $\gamma_i$, une extension $g$ de leur "valeur commune" $\gamma$, par $\mathfrak{S}^I$ --- et la structure supplémentaire, correspondant à la donnée (11), est la "restriction" de cette extension à une extension de $\gamma$ par le sous-groupe diagonal de $\mathfrak{S}^I$.
